% This file is a chapter fragment, in the same style as example.tex.
% It is intended to be included by the main blueprint document.
%
% Status convention:
%   \lean{...} records a proposed declaration name.
%   A paragraph headed ``Mathlib relation'' records an audited exact,
%   equivalent, or adapter relationship with a public Mathlib declaration.
%   No \leanok or \mathlibok marker is used below, because the declarations
%   have not been checked against a concrete Lean/mathlib checkout.

\chapter{Nagata compactification}

This chapter is a formalization blueprint for Nagata's compactification
theorem for schemes.  The key result is
\begin{enumerate}
  \item \Cref{thm:nagata-compactification}, asserting that every separated
  finite type morphism $X\to S$ with $S$ quasi-compact and quasi-separated
  admits morphisms $X\xrightarrow{j}\overline X\xrightarrow{p}S$ satisfying
  $f=p\circ j$, $\operatorname{IsQuasiCompactOpenImmersion}(j)$, and
  $\operatorname{IsProper}(p)$.
\end{enumerate}

The construction refines the scheme-theoretic proof in the Stacks Project.  The
intermediate statements are deliberately finer than the printed proof: every
change of ambient scheme, every use of a scheme-theoretic image, every
admissible blow-up, and every valuative lifting step is given its own labelled
interface.  This is the level at which the objects and equalities will have to
be supplied to Lean.

Throughout this chapter, ``finite type'' includes quasi-compactness, and a
compactification always has a quasi-compact open immersion as its first map.
Fiber products and all morphisms are taken in the category of schemes.

\section{Compactifications and scheme-theoretic images}

\begin{definition}[Choice notation]
  \label{def:nagata-choice-notation}

  For every nonempty set $E$, fix a function value
  \[
    \operatorname{ch}(E)\in E.                                    \tag{C0}
  \]
  Thus every expression $\operatorname{ch}(E)$ below is defined only after
  an equality or cited theorem has established $E\neq\varnothing$.

  \paragraph{Mathlib relation.}
  This is notation for Lean's classical choice operator
  \texttt{\detokenize{Classical.choice}} applied to a term of type
  \texttt{\detokenize{Nonempty E}}.  The notation $\operatorname{ch}$ is
  local to this blueprint; the choice principle itself is in Lean core.
\end{definition}

\subsection{Foundational adapters not packaged by Mathlib}

The definitions in this subsection close representation gaps between the
mathematical notation used below and the public Mathlib interfaces.  They are
blueprint-local declarations: their constituent categories, affine
constructions, sheafification, and gluing operations are in Mathlib, but the
assembled relative constructions are not asserted to be existing Mathlib
declarations.

\begin{definition}[Quasi-coherent tensor, symmetric, and Rees interfaces]
  \label{def:nagata-qcoh-graded-algebra-interface}

  Let $X$ be a scheme.  Write $\operatorname{Mod}(X)$ for the category of
  sheaves of $\mathcal O_X$-modules, and let
  $\operatorname{QCoh}(X)\subseteq\operatorname{Mod}(X)$ be the full
  subcategory of quasi-coherent objects.  For
  $\mathcal E,\mathcal F\in\operatorname{QCoh}(X)$, define
  \[
    \mathcal E\otimes_{\mathcal O_X}\mathcal F
      :=a_X\!\left(
        U\longmapsto
        \mathcal E(U)\otimes_{\mathcal O_X(U)}\mathcal F(U)
      \right),                                                   \tag{MLQC0}
  \]
  where $a_X$ denotes sheafification.  Recursively put
  \[
    \mathcal E^{\otimes0}:=\mathcal O_X,
    \qquad
    \mathcal E^{\otimes(n+1)}
      :=\mathcal E^{\otimes n}\otimes_{\mathcal O_X}\mathcal E.
                                                                    \tag{MLQC1}
  \]
  Define the tensor and symmetric graded algebras by
  \[
  \begin{aligned}
    \operatorname T_X^\bullet(\mathcal E)
      &:=\bigoplus_{n\geq0}\mathcal E^{\otimes n},\\
    \operatorname{Sym}_X^\bullet(\mathcal E)
      &:=\operatorname T_X^\bullet(\mathcal E)/
        \langle e\otimes e'-e'\otimes e\rangle .
  \end{aligned}                                                    \tag{MLQC2}
  \]
  We also write
  $\operatorname{Sym}_{\mathcal O_X}^\bullet(\mathcal E)$ for
  $\operatorname{Sym}_X^\bullet(\mathcal E)$; this is only a change of
  subscript, not a second construction.
  For a quasi-coherent ideal sheaf
  $\mathcal I\subseteq\mathcal O_X$, first define its powers and extension
  along $f\colon Y\to X$ by
  \[
  \begin{gathered}
    \mathcal I^0:=\mathcal O_X,\qquad
    \mathcal I^{n+1}:=
      \operatorname{im}\!\left(
        \mathcal I^n\otimes_{\mathcal O_X}\mathcal I
          \longrightarrow\mathcal O_X\right),\\
    \mathcal I\mathcal O_Y:=
      \operatorname{im}\!\left(
        f^*\mathcal I\longrightarrow f^*\mathcal O_X
          \longrightarrow\mathcal O_Y\right).
  \end{gathered}                                                   \tag{MLQC2a}
  \]
  Define its sheaf Rees algebra by
  \[
    \mathcal R_X(\mathcal I)
      :=\bigoplus_{n\geq0}\mathcal I^n\mathbf u^n,
    \qquad
    (a\mathbf u^m)(b\mathbf u^n):=ab\mathbf u^{m+n}.              \tag{MLQC3}
  \]
  Here all direct sums, quotients, and multiplications are taken in sheaves
  of $\mathcal O_X$-modules and then equipped with the displayed graded
  algebra structure.

  \paragraph{Mathlib relation.}
  The ambient module-sheaf and quasi-coherence predicates are
  \texttt{\detokenize{AlgebraicGeometry.Scheme.Modules}},
  \texttt{\detokenize{SheafOfModules.IsQuasicoherent}}, and its full
  subcategory.  Pointwise presheaf tensor is supplied by
  \texttt{\detokenize{PresheafOfModules.Monoidal.tensorObj}} and
  \texttt{\detokenize{PresheafOfModules.Monoidal.tensorObj_obj}} in the
  module
  \texttt{\detokenize{Mathlib.Algebra.Category.ModuleCat.Presheaf.Monoidal}};
  the sheafification in (MLQC0) is
  \texttt{\detokenize{PresheafOfModules.sheafification}} from
  \texttt{\detokenize{Mathlib.Algebra.Category.ModuleCat.Presheaf.Sheafification}};
  affine associated
  sheaves by \texttt{\detokenize{Mathlib.AlgebraicGeometry.Modules.Tilde}}.
  Mathlib does not currently package (MLQC0)--(MLQC3) together as a
  symmetric or Rees algebra internal to $\operatorname{QCoh}(X)$.
\end{definition}

\begin{lemma}[Affine restriction and base change for the graded interfaces]
  \label{lemma:nagata-qcoh-graded-algebra-interface}
  \uses{def:nagata-qcoh-graded-algebra-interface}

  Let $j_U\colon U=\operatorname{Spec}(A)\hookrightarrow X$ be affine.  If
  $\mathcal E|_U\cong\widetilde M$,
  $\mathcal F|_U\cong\widetilde N$, and
  $\mathcal I|_U\cong\widetilde I$, then there are canonical isomorphisms
  \[
  \begin{aligned}
    (\mathcal E\otimes\mathcal F)|_U
      &\cong\widetilde{M\otimes_A N},\\
    \operatorname{Sym}_X^\bullet(\mathcal E)|_U
      &\cong\widetilde{\operatorname{Sym}_A^\bullet(M)},\\
    \mathcal R_X(\mathcal I)|_U
      &\cong\widetilde{\bigoplus_{n\geq0}I^n\mathbf u^n}.
  \end{aligned}                                                    \tag{MLQC4}
  \]
  These identifications preserve degrees, units, multiplication, restriction
  to smaller affine opens, and the canonical maps from degree zero and degree
  one.  In particular the objects in (MLQC0)--(MLQC3) are quasi-coherent.
  For every $f\colon Y\to X$ there are coherent comparison isomorphisms
  \[
  \begin{aligned}
    f^*(\mathcal E\otimes\mathcal F)
      &\cong f^*\mathcal E\otimes f^*\mathcal F,\\
    f^*\operatorname{Sym}_X^\bullet(\mathcal E)
      &\cong\operatorname{Sym}_Y^\bullet(f^*\mathcal E),\\
    f^*\mathcal R_X(\mathcal I)
      &\longrightarrow\mathcal R_Y(\mathcal I\mathcal O_Y),
  \end{aligned}                                                    \tag{MLQC5}
  \]
  and the last arrow is an isomorphism when the formation of the displayed
  ideal powers commutes with the pullback, in particular for flat $f$.
\end{lemma}

\begin{proof}
  The presheaf tensor is defined objectwise, and sheafification has its usual
  universal property.  On an affine scheme, the associated-sheaf functor
  identifies tensor products, symmetric powers, ideal powers, direct sums,
  and quotients with the corresponding module constructions.  These affine
  identifications commute with localization, so they glue and prove (MLQC4).
  The same affine-local calculation after applying $f^*$ gives (MLQC5);
  flatness makes extension of ideals commute with products and powers.
\end{proof}

\begin{definition}[Relative Proj by affine gluing]
  \label{def:nagata-relative-proj}
  \uses{def:nagata-qcoh-graded-algebra-interface,
    lemma:nagata-qcoh-graded-algebra-interface}

  Let $X$ be a scheme and let
  $\mathcal A=\bigoplus_{n\geq0}\mathcal A_n$ be a quasi-coherent graded
  $\mathcal O_X$-algebra with $\mathcal A_0=\mathcal O_X$.  Choose an affine
  open cover $j_i\colon U_i=\operatorname{Spec}(A_i)\hookrightarrow X$ and
  graded rings $B_i$ with
  $\mathcal A|_{U_i}\cong\widetilde{B_i}$.  Put
  \[
    P_i:=\operatorname{Proj}(B_i),
    \qquad p_i\colon P_i\longrightarrow U_i.                    \tag{MLRP0}
  \]
  Localization of the $B_i$ on $U_i\cap U_j$, followed by the affine
  \texttt{Proj} construction, gives canonical overlap isomorphisms
  \[
    \theta_{ij}\colon
      P_i\times_XU_j\xrightarrow{\sim}P_j\times_XU_i.             \tag{MLRP1}
  \]
  They satisfy identity, symmetry, and cocycle equations because the
  localization comparison maps do.  Define
  \[
    \operatorname{Proj}_X(\mathcal A)
      :=\operatorname{Glue}\bigl((P_i)_i,(\theta_{ij})_{i,j}\bigr),
    \qquad
    \operatorname{pr}_{X,\mathcal A}\colon
      \operatorname{Proj}_X(\mathcal A)\longrightarrow X          \tag{MLRP2}
  \]
  by scheme gluing and by gluing the compatible morphisms $p_i$.

  If $R$ is a ring and $\varphi$ is a morphism of graded $R$-algebras for
  which (MLRP3) below is defined, use the affine shorthand
  \[
    \operatorname{Proj}_R(\varphi)
      :=\operatorname{Proj}_{\operatorname{Spec}(R)}
        (\widetilde\varphi).                                      \tag{MLRP2a}
  \]

  \paragraph{Mathlib relation.}
  Each local $P_i$ is the exact declaration
  \texttt{\detokenize{AlgebraicGeometry.Proj}}.  Its standard opens and
  affine charts are \texttt{\detokenize{Proj.basicOpen}},
  \texttt{\detokenize{Proj.basicOpenIsoSpec}}, and
  \texttt{\detokenize{Proj.affineOpenCover}}.  The gluing leaves are
  \texttt{\detokenize{Scheme.GlueData}} and
  \texttt{\detokenize{Scheme.Cover.glueMorphisms}}.  The assembled relative
  object (MLRP2) is a blueprint-local construction, not an existing general
  Mathlib relative-\texttt{Proj} declaration.
\end{definition}

\begin{lemma}[Relative-Proj interface]
  \label{lemma:nagata-relative-proj-interface}
  \uses{def:nagata-relative-proj}

  The objects obtained from two choices in
  \Cref{def:nagata-relative-proj} are related by a canonical comparison
  isomorphism over $X$, compatible with all local absolute-\(\operatorname{Proj}\)
  identifications.
  It has the following interfaces.
  \begin{enumerate}
    \item If $X=\operatorname{Spec}(A)$ and
      $\mathcal A\cong\widetilde B$, then
      $\operatorname{Proj}_X(\mathcal A)\cong\operatorname{Proj}(B)$ over
      $X$.
    \item A degree-preserving morphism
      $\varphi\colon\mathcal A\to\mathcal B$ satisfying, affine-locally,
      the explicit irrelevant-ideal condition
      \[
        \sqrt{B_+}=\sqrt{B\,\varphi(A_+)}
      \]
      induces the contravariant morphism
      \[
        \operatorname{Proj}_X(\varphi)\colon
        \operatorname{Proj}_X(\mathcal B)\longrightarrow
          \operatorname{Proj}_X(\mathcal A),                     \tag{MLRP3}
      \]
      functorially in $\varphi$.  If $\varphi$ is an isomorphism, so is
      $\operatorname{Proj}_X(\varphi)$; if $\varphi$ is a surjective
      quotient, (MLRP3) is a closed immersion.
    \item For $f\colon Y\to X$, affine localization and gluing give a
      canonical base-change isomorphism
      \[
        \operatorname{Proj}_Y(f^*\mathcal A)
          \xrightarrow{\sim}
        \operatorname{Proj}_X(\mathcal A)\times_XY,               \tag{MLRP4}
      \]
      natural in $f$ and $\mathcal A$.
    \item Homogeneous sections of positive degree define standard opens;
      on every affine member of the chosen cover these are the Mathlib opens
      $D_+(s)$ and have the degree-zero homogeneous-localization chart
      $\operatorname{Proj}(B)|_{D_+(s)}\cong\operatorname{Spec}(B_{(s)})$.
    \item If $\mathcal A$ is of finite type and generated in positive
      degrees locally on $X$, the structure morphism in (MLRP2) is locally
      projective.  If a finite degree-one generating family is global, it is
      projective and hence proper.
  \end{enumerate}
\end{lemma}

\begin{proof}
  Cover every pairwise intersection of two affine presentation opens by
  affine principal opens, thereby obtaining a common affine refinement.
  The uniqueness and cocycle laws for localization identify the two glue
  data, and uniqueness in scheme gluing gives cover independence.  All five
  assertions are local on $X$.  After restriction to an affine open they are
  respectively the absolute $\operatorname{Proj}$ construction,
  contravariance of homogeneous prime spectra, the localization/base-change
  calculation, the homogeneous-localization chart theorem, and the standard
  closed immersion into projective space.  These affine maps agree on
  overlaps by their universal properties, hence glue.
\end{proof}

\begin{definition}[Projective bundles and finite relative projective space]
  \label{def:nagata-relative-projective-space}
  \uses{def:nagata-qcoh-graded-algebra-interface,
    def:nagata-relative-proj}

  For a finite-type quasi-coherent $\mathcal O_X$-module $\mathcal E$, put
  \[
    \mathbf P_X(\mathcal E)
      :=\operatorname{Proj}_X
        \bigl(\operatorname{Sym}_X^\bullet(\mathcal E)\bigr),
    \qquad
    \operatorname{pr}_{X,\mathcal E}\colon
      \mathbf P_X(\mathcal E)\longrightarrow X.                  \tag{MLPS0}
  \]
  For $n\in\mathbf N$, define
  \[
    \mathbf P_X^n:=\mathbf P_X(\mathcal O_X^{\oplus(n+1)}).
                                                                    \tag{MLPS1}
  \]
  Equivalently, using the structural map
  $s_X\colon X\to\operatorname{Spec}(\mathbf Z)$, fix the canonical
  isomorphism
  \[
    \mathbf P_X^n\xrightarrow{\sim}
      X\times_{\operatorname{Spec}(\mathbf Z)}
        \operatorname{Proj}(\mathbf Z[T_0,\ldots,T_n]).           \tag{MLPS2}
  \]
\end{definition}

\begin{lemma}[Finite relative projective-space interface]
  \label{lemma:nagata-relative-projective-space}
  \uses{def:nagata-relative-projective-space,
    lemma:nagata-relative-proj-interface}

  The projection $\mathbf P_X^n\to X$ is proper.  Its standard opens
  $D_+(T_i)$ are open affine subschemes isomorphic to $\mathbf A_X^n$, they
  cover $\mathbf P_X^n$, and all these identifications commute with arbitrary
  base change $Y\to X$.

  \paragraph{Mathlib relation.}
  The absolute projection is
  \texttt{\detokenize{AlgebraicGeometry.Proj.toSpecZero}}, whose properness
  instance is
  \texttt{\detokenize{Proj.instIsProperToSpecZeroOfFiniteTypeSubtypeMemOfNatNat}}.
  Properness after (MLPS2) is the Mathlib base-change instance for
  \texttt{\detokenize{AlgebraicGeometry.IsProper}}.  The affine-space target
  is \texttt{\detokenize{AlgebraicGeometry.AffineSpace}}.
\end{lemma}

\begin{proof}
  Formula (MLPS2) identifies the projection with the base change of the
  proper absolute $\operatorname{Proj}$ projection, so it is proper.
  The remaining assertions are the standard $D_+(T_i)$ chart calculation
  for absolute $\operatorname{Proj}$, transported through the same fibre
  product.
\end{proof}

\begin{definition}[Predicate notation for morphisms]
  \label{def:nagata-morphism-predicate-notation}
  \uses{def:nagata-relative-projective-space,
    lemma:nagata-relative-projective-space}

  For a morphism \(h\colon T\to Y\), write
  \(\Delta_h\colon T\to T\times_YT\) for its diagonal,
  \(|h|\colon|T|\to|Y|\) for its map of underlying spaces,
  \(h^\sharp\colon h^{-1}\mathcal O_Y\to\mathcal O_T\) for its
  structure-sheaf morphism, and
  \(h^\flat\colon\mathcal O_Y\to h_*\mathcal O_T\) for the adjoint of
  \(h^\sharp\).  Put
  \[
  \begin{gathered}
    \operatorname{Im}(|h|):=|h|(|T|)\subseteq|Y|
      \quad\text{with the subspace topology},\\
    |h|_{\operatorname{Im}}\colon|T|\to\operatorname{Im}(|h|),
      \qquad |h|_{\operatorname{Im}}(t):=|h|(t),\\
    \mathcal O_{Y,|h|(t)}
      :=\varinjlim_{|h|(t)\in V\subseteq|Y|\ {\rm open}}\mathcal O_Y(V),
      \qquad
    (h^{-1}\mathcal O_Y)_t
      :=\varinjlim_{|h|(t)\in V\subseteq|Y|\ {\rm open}}\mathcal O_Y(V),\\
    \lambda_{h,t}\colon
      \mathcal O_{Y,|h|(t)}
        \xrightarrow{\sim}(h^{-1}\mathcal O_Y)_t,\qquad
    \lambda_{h,t}([s]_V):=[s]_V
      \quad(t\in|T|),\\
    h_t^\sharp\colon\mathcal O_{Y,|h|(t)}\to\mathcal O_{T,t}
      \quad(t\in|T|)
      :=(h^\sharp)_t\circ\lambda_{h,t}.
  \end{gathered}                                                    \tag{P-1}
  \]
  For every morphism \(v\colon Y'\to Y\), define
  \[
    T_v:=T\times_{Y,h,v}Y',\qquad
    h_v:=\operatorname{pr}_{Y'}\colon T_v\to Y'.                 \tag{P-2}
  \]
  For a ring \(R\), an \(R\)-module \(M\), and a quasi-coherent
  \(\mathcal O_Y\)-module \(\mathcal E\), define
  \[
  \begin{aligned}
    \operatorname{Flat}_R(M)
      &:\Longleftrightarrow
        \forall\text{ \(R\)-modules }N,N'\
        \forall\text{ \(R\)-linear }u\colon N'\to N,\quad
        \ker(u)=0\Longrightarrow
        \ker(u\otimes_R\operatorname{id}_M)=0,\\
    \operatorname{FiniteTypeModule}_Y(\mathcal E)
      &:\Longleftrightarrow
        \forall y\in|Y|\ \exists(j_U\colon U\to Y)\
        \exists n\in\mathbf N\ \exists
        \bigl(a\colon\mathcal O_U^{\oplus n}\to j_U^*\mathcal E\bigr):
        \quad
        \operatorname{IsOpenImmersion}(j_U)
        \ \wedge\ y\in\operatorname{Im}(|j_U|)
        \ \wedge\ \operatorname{Coker}(a)=0.
  \end{aligned}                                                    \tag{P-3}
  \]
  For topological spaces \(A,B\), a continuous map
  \(\varphi\colon A\to B\), and a morphism
  \(a\colon\mathcal F\to\mathcal G\) of sheaves on one scheme, define
  \[
  \begin{aligned}
    \operatorname{Homeo}(\varphi)
      &:\Longleftrightarrow
        \exists\,\psi\colon B\to A\text{ continuous}:\quad
        \psi\circ\varphi=\operatorname{id}_A\ \wedge\
        \varphi\circ\psi=\operatorname{id}_B,\\
    \operatorname{IsSheafIso}(a)
      &:\Longleftrightarrow
        \exists\,b\colon\mathcal G\to\mathcal F:\quad
        b\circ a=\operatorname{id}_{\mathcal F}\ \wedge\
        a\circ b=\operatorname{id}_{\mathcal G}.
  \end{aligned}                                                    \tag{P-4}
  \]
  Define
  \[
  \begin{aligned}
    \operatorname{IsIso}(h)
      &:\Longleftrightarrow
        \exists k\colon Y\to T:\ 
        (k\circ h=\operatorname{id}_T)\wedge
        (h\circ k=\operatorname{id}_Y),\\
    \operatorname{IsSeparated}(h)
      &:\Longleftrightarrow
        \operatorname{IsClosedImmersion}(\Delta_h),\\
    \operatorname{IsFlat}(h)
      &:\Longleftrightarrow
        \forall t\in|T|,\quad
        \operatorname{Flat}_{\mathcal O_{Y,|h|(t)}}(\mathcal O_{T,t}),
        \quad
        \mathcal O_{Y,|h|(t)}\text{-linearity being induced by }h_t^\sharp,\\
    \operatorname{IsAffine}(h)
      &:\Longleftrightarrow
        \forall(A\text{ a ring})\
        \forall\bigl(v\colon\operatorname{Spec}(A)\to Y\bigr),\quad
        \operatorname{IsOpenImmersion}(v)
        \Longrightarrow T_v\text{ is an affine scheme},\\
    \operatorname{IsQuasiCompact}(h)
      &:\Longleftrightarrow
        \forall\,W\subseteq|Y|,\quad
        \bigl(W\text{ is open and quasi-compact}\bigr)
        \Longrightarrow |h|^{-1}(W)\text{ is quasi-compact},\\
    \operatorname{QuasiSeparated}(h)
      &:\Longleftrightarrow
        \operatorname{IsQuasiCompact}(\Delta_h),\\
    \operatorname{IsQCQS}(h)
      &:\Longleftrightarrow
        \operatorname{IsQuasiCompact}(h)\wedge
        \operatorname{QuasiSeparated}(h),\\
    \operatorname{IsQuasiCompactOpenImmersion}(h)
      &:\Longleftrightarrow
        \operatorname{IsOpenImmersion}(h)\wedge
        \operatorname{IsQuasiCompact}(h),\\
    \operatorname{UniversallyClosed}(h)
      &:\Longleftrightarrow
        \forall(v\colon Y'\to Y)\ \forall C\subseteq|T_v|,\quad
        C\text{ closed}\Longrightarrow |h_v|(C)\text{ closed in }|Y'|,\\
    \operatorname{IsProper}(h)
      &:\Longleftrightarrow
        \operatorname{IsSeparated}(h)\wedge
        \operatorname{FiniteType}(h)\wedge
        \operatorname{UniversallyClosed}(h),\\
    \operatorname{IsProjective}(h)
      &:\Longleftrightarrow
        \exists\,\mathcal E\in\operatorname{QCoh}(\mathcal O_Y)\ \
        \exists\,\iota\colon T\to\mathbf P_Y(\mathcal E):\
        \operatorname{FiniteTypeModule}_Y(\mathcal E)\ \wedge\
        \operatorname{IsClosedImmersion}(\iota)\ \wedge\
        \operatorname{pr}_{Y,\mathcal E}\circ\iota=h,\\
    \operatorname{LocallyProjective}(h)
      &:\Longleftrightarrow
        \exists I\ \exists(v_i\colon V_i\to Y)_{i\in I}:\
        |Y|=\bigcup_{i\in I}\operatorname{Im}(|v_i|)\ \wedge\
        (\forall i\in I)\bigl(
          \operatorname{IsOpenImmersion}(v_i)\ \wedge\
          \operatorname{IsProjective}(h_{v_i})\bigr).
  \end{aligned}                                                     \tag{P0}
  \]
  For the three immersion predicates, fix the exact meanings of
  \stacksproject{01HE}, \stacksproject{01HK}, \stacksproject{01IO}; symbolically,
  \[
  \begin{aligned}
    \operatorname{IsOpenImmersion}(h)
      &:\Longleftrightarrow
        \operatorname{Homeo}
          \bigl(|h|_{\operatorname{Im}}\bigr)
        \ \wedge\ \operatorname{Im}(|h|)\text{ is open in }|Y|
        \ \wedge\ \operatorname{IsSheafIso}(h^\sharp),\\
    \operatorname{IsClosedImmersion}(h)
      &:\Longleftrightarrow
        \operatorname{Homeo}
          \bigl(|h|_{\operatorname{Im}}\bigr)
        \ \wedge\ \operatorname{Im}(|h|)\text{ is closed in }|Y|
        \ \wedge\ \operatorname{Coker}(h^\flat)=0,\\
    \operatorname{IsImmersion}(h)
      &:\Longleftrightarrow
        \exists(W\text{ a scheme})\
        \exists(i\colon T\to W)\ \exists(j\colon W\to Y):\
        h=j\circ i\ \wedge\
        \operatorname{IsClosedImmersion}(i)\ \wedge\
        \operatorname{IsOpenImmersion}(j).
  \end{aligned}                                                     \tag{P0a}
  \]
  Define the remaining finiteness predicates by the pointwise affine
  conditions of
  \cite[\href{https://stacks.math.columbia.edu/tag/01T1}{Tag 01T1}]
  {stacks-project} and
  \cite[\href{https://stacks.math.columbia.edu/tag/01TP}{Tag 01TP}]
  {stacks-project}:
  for a ring \(A\), \(n,m\in\mathbf N\), and
  \(\mathbf F=(F_1,\ldots,F_m)\in A[x_1,\ldots,x_n]^m\), put
  \[
    Q_{A,n,\mathbf F}
      :=A[x_1,\ldots,x_n]/(F_1,\ldots,F_m),\qquad
    \overline\iota_{A,n,\mathbf F}\colon
      A\longrightarrow Q_{A,n,\mathbf F}                           \tag{P0c}
  \]
  for the composite of \(A\hookrightarrow A[x_1,\ldots,x_n]\) with the
  quotient map.  Also, for rings \(A,B\), a ring map
  \(\varphi\colon A\to B\), open
  immersions \(u\colon\operatorname{Spec}(B)\to T\) and
  \(v\colon\operatorname{Spec}(A)\to Y\), and \(t\in|T|\), define
  \[
  \begin{aligned}
    \operatorname{AffChart}_h(t;A,B,u,v,\varphi)
      :\Longleftrightarrow{}&
      \operatorname{IsOpenImmersion}(u)\ \wedge\
      \operatorname{IsOpenImmersion}(v)\ \wedge\
      t\in\operatorname{Im}(|u|)\\
      &{}\wedge\
      h\circ u=v\circ\operatorname{Spec}(\varphi).
  \end{aligned}                                                     \tag{P0b}
  \]
  \[
  \begin{aligned}
    \operatorname{LocallyFiniteType}(h)
      &:\Longleftrightarrow
        \forall t\in|T|\ \exists(A,B,u,v,\varphi)\
        \exists n\in\mathbf N\ \exists b_1,\ldots,b_n\in B:\quad
        \operatorname{AffChart}_h(t;A,B,u,v,\varphi)
        \ \wedge\
        B=\varphi(A)[b_1,\ldots,b_n],\\
    \operatorname{FiniteType}(h)
      &:\Longleftrightarrow
        \operatorname{IsQuasiCompact}(h)\wedge
        \operatorname{LocallyFiniteType}(h),\\
    \operatorname{LocallyFinitePresentation}(h)
      &:\Longleftrightarrow
        \forall t\in|T|\ \exists(A,B,u,v,\varphi)\
        \exists n,m\in\mathbf N\
        \exists\mathbf F\in A[x_1,\ldots,x_n]^m\
        \exists(\theta\colon B\to Q_{A,n,\mathbf F})\
        \exists(\theta^{-1}\colon Q_{A,n,\mathbf F}\to B):\quad
        \operatorname{AffChart}_h(t;A,B,u,v,\varphi)
        \ \wedge\
        \theta^{-1}\circ\theta=\operatorname{id}_B
        \ \wedge\
        \theta\circ\theta^{-1}=\operatorname{id}_{Q_{A,n,\mathbf F}}
        \ \wedge\
        \theta\circ\varphi=\overline\iota_{A,n,\mathbf F},
    \operatorname{FinitePresentation}(h)
      &:\Longleftrightarrow
        \operatorname{IsQuasiCompact}(h)\wedge
        \operatorname{QuasiSeparated}(h)\wedge
        \operatorname{LocallyFinitePresentation}(h),\\
    \operatorname{Proper}(h)&:\Longleftrightarrow
      \operatorname{IsProper}(h),&
    \operatorname{Separated}(h)&:\Longleftrightarrow
      \operatorname{IsSeparated}(h),\\
    \operatorname{QuasiCompact}(h)&:\Longleftrightarrow
      \operatorname{IsQuasiCompact}(h),&
    \operatorname{Projective}(h)&:\Longleftrightarrow
      \operatorname{IsProjective}(h).
  \end{aligned}                                                     \tag{P1}
  \]
  For a scheme $T$, let
  \[
    s_T\colon T\to\operatorname{Spec}(\mathbf Z),\qquad
    \Delta_T:=\Delta_{s_T}\colon
      T\to T\times_{\operatorname{Spec}(\mathbf Z)}T               \tag{P1a}
  \]
  be respectively its structural morphism and its absolute diagonal, and
  define the scheme-level predicates
  \[
    \operatorname{QuasiCompact}(T)
      :\Longleftrightarrow\operatorname{IsQuasiCompact}(s_T),
    \qquad
    \operatorname{QuasiSeparated}(T)
      :\Longleftrightarrow\operatorname{QuasiSeparated}(s_T).       \tag{P1b}
  \]
  For a ring $A$ and a scheme $T$, define
  \[
  \begin{aligned}
    \operatorname{NoetherianRing}(A)
      &:\Longleftrightarrow
        \forall(I_n\triangleleft A)_{n\in\mathbf N},\quad
        \bigl[(\forall n,\ I_n\subseteq I_{n+1})
          \Longrightarrow
          (\exists N\ \forall n\geq N,\ I_n=I_N)\bigr],\\
    \operatorname{NoetherianScheme}(T)
      &:\Longleftrightarrow
        \exists r\in\mathbf N\ \exists(A_i)_{i=1}^r\
        \exists\bigl(u_i\colon\operatorname{Spec}(A_i)\to T\bigr)_{i=1}^r:
        \quad
        \bigwedge_{i=1}^r\operatorname{IsOpenImmersion}(u_i)
        \ \wedge\
        |T|=\bigcup_{i=1}^r\operatorname{Im}(|u_i|)
        \ \wedge\
        \bigwedge_{i=1}^r\operatorname{NoetherianRing}(A_i),\\
    \operatorname{Noetherian}(T)
      &:\Longleftrightarrow\operatorname{NoetherianScheme}(T).
  \end{aligned}                                                     \tag{P2}
  \]
  Here, for every finite-type quasi-coherent $\mathcal O_Y$-module
  $\mathcal E$, the symbols in the projectivity predicates are defined by
  \[
    \mathbf P_Y(\mathcal E)
      :=\operatorname{Proj}_Y
        \bigl(\operatorname{Sym}_{\mathcal O_Y}^{\bullet}(\mathcal E)\bigr),
    \qquad
    \operatorname{pr}_{Y,\mathcal E}\colon
      \mathbf P_Y(\mathcal E)\longrightarrow Y.                    \tag{P3}
  \]
  Thus $\operatorname{IsProjective}$ and
  $\operatorname{LocallyProjective}$ are respectively the notions called
  projective and locally projective in
  \cite[\href{https://stacks.math.columbia.edu/tag/01W8}{Tag 01W8}]
  {stacks-project}.

  \paragraph{Mathlib relation.}
  The categorical isomorphism predicate is
  \texttt{\detokenize{CategoryTheory.IsIso}}.  The scheme-morphism
  predicates used here correspond to the public classes
  \texttt{\detokenize{AlgebraicGeometry.IsOpenImmersion}},
  \texttt{\detokenize{IsClosedImmersion}},
  \texttt{\detokenize{IsImmersion}},
  \texttt{\detokenize{QuasiCompact}},
  \texttt{\detokenize{QuasiSeparated}},
  \texttt{\detokenize{IsSeparated}},
  \texttt{\detokenize{Flat}},
  \texttt{\detokenize{IsAffineHom}},
  \texttt{\detokenize{LocallyOfFiniteType}},
  \texttt{\detokenize{FiniteType}},
  \texttt{\detokenize{LocallyOfFinitePresentation}},
  \texttt{\detokenize{FinitePresentation}},
  \texttt{\detokenize{UniversallyClosed}}, and
  \texttt{\detokenize{IsProper}}.  Ring/module flatness is
  \texttt{\detokenize{Module.Flat}}, and the stalk bridge is
  \texttt{\detokenize{AlgebraicGeometry.Flat.iff_flat_stalkMap}}.
  Finite-type module sheaves
  are \texttt{\detokenize{SheafOfModules.IsFiniteType}}.  The ring and
  scheme Noetherian predicates correspond to
  \texttt{\detokenize{IsNoetherianRing}} and
  \texttt{\detokenize{AlgebraicGeometry.IsNoetherian}}, with
  \texttt{\detokenize{AlgebraicGeometry.isNoetherian_iff_of_finite_iSup_eq_top}}
  as the
  finite-affine-cover bridge.  The displayed formulas deliberately remain
  the blueprint's unbundled notation; projectivity and the relative
  projective bundle use the local interfaces above because no matching
  scheme-morphism projective class is currently public in Mathlib.
\end{definition}

\begin{definition}[Closed-subscheme classes]
  \label{def:nagata-closed-subscheme-class}
  \uses{def:nagata-morphism-predicate-notation}

  For a scheme \(Y\), define
  \[
    \operatorname{ClosedEmb}(Y):=
      \{(Z,i)\mid Z\text{ is a scheme},\
        i\colon Z\to Y,\
        \operatorname{IsClosedImmersion}(i)\}.                     \tag{CSub0}
  \]
  For \((Z,i),(Z',i')\in\operatorname{ClosedEmb}(Y)\), define
  \[
    (Z,i)\sim_Y(Z',i')
      :\Longleftrightarrow
      \exists(\theta\colon Z\to Z'):\quad
      \operatorname{IsIso}(\theta)\ \wedge\ i'\circ\theta=i.       \tag{CSub1}
  \]
  The relation is reflexive because
  \(i\circ\operatorname{id}_Z=i\).  If \(\theta\) witnesses
  \((Z,i)\sim_Y(Z',i')\), choose the inverse \(\theta^{-1}\) supplied by
  \(\operatorname{IsIso}(\theta)\); then
  \[
    i\circ\theta^{-1}
      =i'\circ\theta\circ\theta^{-1}
      =i'\circ\operatorname{id}_{Z'}=i'.                           \tag{CSub1a}
  \]
  If \(\theta\) and \(\theta'\) respectively witness
  \((Z,i)\sim_Y(Z',i')\) and \((Z',i')\sim_Y(Z'',i'')\), then
  \[
  \begin{gathered}
    i''\circ(\theta'\circ\theta)
      =(i''\circ\theta')\circ\theta
      =i'\circ\theta=i,\\
    (\theta^{-1}\circ(\theta')^{-1})\circ(\theta'\circ\theta)
      =\theta^{-1}\circ\operatorname{id}_{Z'}\circ\theta
      =\operatorname{id}_Z,\\
    (\theta'\circ\theta)\circ(\theta^{-1}\circ(\theta')^{-1})
      =\theta'\circ\operatorname{id}_{Z'}\circ(\theta')^{-1}
      =\operatorname{id}_{Z''},\\
    \operatorname{IsIso}(\theta'\circ\theta).
  \end{gathered}                                                     \tag{CSub1b}
  \]
  Thus \(\sim_Y\) is an equivalence relation.  Define the set of closed
  subschemes and the arrow notation by
  \[
    \operatorname{ClosedSub}(Y):=
      \operatorname{ClosedEmb}(Y)/{\sim_Y},\qquad
    [\,Z\xhookrightarrow{i}Y\,]:=[(Z,i)]_{\sim_Y}.                 \tag{CSub2}
  \]
  Consequently,
  \[
    [\,Z\xhookrightarrow{i}Y\,]
      =[\,Z'\xhookrightarrow{i'}Y\,]
    \Longleftrightarrow
    \exists(\theta\colon Z\to Z'):\quad
      \operatorname{IsIso}(\theta)\ \wedge\ i'\circ\theta=i.       \tag{CSub3}
  \]

  \paragraph{Mathlib relation.}
  Mathlib represents a closed subscheme by
  \texttt{\detokenize{Scheme.IdealSheafData}} together with
  \texttt{\detokenize{IdealSheafData.subscheme}} and
  \texttt{\detokenize{subschemeι}}.  Closed immersions are related to that
  representation by
  \texttt{\detokenize{AlgebraicGeometry.IsClosedImmersion.overEquivIdealSheafData}}.
  The quotient by isomorphism in (CSub2) is a blueprint adapter, not
  Mathlib's definition of closed subschemes.
\end{definition}


\begin{lemma}[Morphism-property calculus]
  \label{lemma:nagata-morphism-property-calculus}
  \uses{def:nagata-morphism-predicate-notation}

  Let
  \[
    X\xrightarrow{f}Y\xrightarrow{g}Z,
    \qquad a\colon Y'\to Y,
  \]
  and define the base change of $f$ along $a$ by
  \[
    X_{Y'}:=X\times_{Y,f,a}Y',
    \qquad
    f_{Y'}:=\operatorname{pr}_{Y'}\colon X_{Y'}\to Y'.             \tag{PC0}
  \]
  Define the following set of predicates on morphisms of schemes:
  \[
  \begin{split}
    \mathfrak P:=\{&\operatorname{IsIso},
      \operatorname{IsOpenImmersion},
      \operatorname{IsClosedImmersion},
      \operatorname{IsQuasiCompact},
      \operatorname{QuasiSeparated},
      \operatorname{IsQCQS},
      \operatorname{IsSeparated},\\
      &\operatorname{FiniteType},
      \operatorname{FinitePresentation},
      \operatorname{IsProper},
      \operatorname{IsQuasiCompactOpenImmersion}\}.
  \end{split}                                                       \tag{PC1}
  \]
  Then arbitrary base change and composition are the implications
  \[
    (\forall\mathcal P\in\mathfrak P)\qquad
      \mathcal P(f)\Longrightarrow\mathcal P(f_{Y'}).             \tag{PC2}
  \]
  \[
    (\forall\mathcal P\in\mathfrak P)\qquad
      \mathcal P(f)\wedge\mathcal P(g)
        \Longrightarrow\mathcal P(g\circ f).                     \tag{PC3}
  \]
  If $(V_i)_{i\in I}$ is an open cover of $Y$, define
  \[
    v_i\colon V_i\hookrightarrow Y,
    \qquad
    X_i:=X\times_{Y,f,v_i}V_i,
    \qquad
    f_i:=\operatorname{pr}_{V_i}\colon X_i\to V_i.                \tag{PC4}
  \]
  Target locality is the equivalence
  \[
    (\forall\mathcal P\in\mathfrak P)\qquad
    \mathcal P(f)
      \Longleftrightarrow
    (\forall i\in I)\ \mathcal P(f_i).                            \tag{PC5}
  \]

  Projectivity satisfies the base-change implication
  \[
    \operatorname{IsProjective}(f)
      \Longrightarrow
    \operatorname{IsProjective}(f_{Y'}).                           \tag{PC6}
  \]
  If $s_Z\colon Z\to\operatorname{Spec}(\mathbf Z)$ denotes the structural
  morphism, then its composition implication is
  \[
    \operatorname{IsQCQS}(s_Z)\ \wedge\
    \operatorname{IsProjective}(f)\ \wedge\
    \operatorname{IsProjective}(g)
      \Longrightarrow
    \operatorname{IsProjective}(g\circ f).                        \tag{PC7}
  \]
  For isomorphisms
  $\alpha\colon X'\xrightarrow{\sim}X$ and
  $\beta\colon Y\xrightarrow{\sim}Y'$, define
  \[
    f^{\alpha,\beta}:=\beta\circ f\circ\alpha\colon X'\to Y'.
  \]
  Invariance under isomorphic source and target is
  \[
    \operatorname{IsProjective}(f^{\alpha,\beta})
      \Longleftrightarrow
    \operatorname{IsProjective}(f).                               \tag{PC8}
  \]
  Projectivity is not among the predicates in (PC5): there exist a morphism
  $h\colon T\to S$ and a two-member open cover
  $S=V_1\cup V_2$ such that, for
  $h_i:=\operatorname{pr}_{V_i}\colon T\times_SV_i\to V_i$,
  \[
    \operatorname{IsProjective}(h_1)\wedge
    \operatorname{IsProjective}(h_2)\wedge
    \neg\operatorname{IsProjective}(h).                           \tag{PC9}
  \]

  For morphisms $r\colon A\to S$ and $s\colon B\to S$, put
  \[
    P:=A\times_{S,r,s}B,
    \qquad p_A\colon P\to A,
    \qquad p_B\colon P\to B.
  \]
  Fibre-product extensionality is, for every scheme $W$ and
  $u,v\colon W\to P$, the equivalence
  \[
    u=v
      \Longleftrightarrow
    \bigl(p_A\circ u=p_A\circ v\bigr)\ \wedge\
    \bigl(p_B\circ u=p_B\circ v\bigr).                            \tag{PC10}
  \]
  Finally, let $T=\bigcup_{i\in I}U_i$ be an open cover, let
  $j_i\colon U_i\hookrightarrow T$, and define
  \[
  \begin{gathered}
    U_{ij}:=U_i\times_TU_j,
    \qquad p_{i,ij}\colon U_{ij}\to U_i,
    \qquad p_{j,ij}\colon U_{ij}\to U_j,\\
    \mathscr G_T(H,(U_i)_i):=
      \left\{(h_i)_i\in\prod_{i\in I}\operatorname{Hom}(U_i,H)
      \ \middle|\
      h_i\circ p_{i,ij}=h_j\circ p_{j,ij}
      \text{ for all }i,j\in I\right\},\\
    \operatorname{res}_{T,H}\colon\operatorname{Hom}(T,H)\to
      \mathscr G_T(H,(U_i)_i),
    \qquad
    \operatorname{res}_{T,H}(h):=(h\circ j_i)_i.
  \end{gathered}                                                   \tag{PC11}
  \]
  The Hom-gluing assertion is the bijection
  \[
    \operatorname{res}_{T,H}\colon
      \operatorname{Hom}(T,H)
      \xrightarrow{\ \sim\ }
      \mathscr G_T(H,(U_i)_i).                                    \tag{PC12}
  \]

  \paragraph{Mathlib relation.}
  This lemma aggregates Mathlib's property typeclass infrastructure:
  \texttt{\detokenize{MorphismProperty.IsStableUnderBaseChange}},
  \texttt{\detokenize{MorphismProperty.IsStableUnderComposition}},
  the concrete scheme-property instances, and
  \texttt{\detokenize{IsZariskiLocalAtTarget.iff_of_openCover}}.  Fibre
  product extensionality is \texttt{\detokenize{pullback.hom_ext}}, while
  (PC12) is implemented by
  \texttt{\detokenize{Scheme.Cover.glueMorphisms}} and
  \texttt{\detokenize{Scheme.Cover.hom_ext}}.  The projectivity clauses
  (PC6)--(PC9) remain blueprint-local because Mathlib has no matching general
  scheme-morphism projectivity predicate.
\end{lemma}

\begin{proof}
  The base-change implications (PC2) are precisely the entries of
  \cite[\href{https://stacks.math.columbia.edu/tag/02WF}{Tag 02WF}]
  {stacks-project}; the individual references for the predicates in (PC1) are
  \[
  \begin{array}{c|c}
    \text{predicate}&\text{base-change tag}\\ \hline
    \operatorname{IsOpenImmersion},\operatorname{IsClosedImmersion}
      &\text{01JY}\\
    \operatorname{IsQuasiCompact}&\text{01K5}\\
    \operatorname{QuasiSeparated},\operatorname{IsSeparated}&\text{01KU}\\
    \operatorname{FiniteType}&\text{01T4}\\
    \operatorname{FinitePresentation}&\text{01TS}\\
    \operatorname{IsProper}&\text{01W4}.
  \end{array}                                                       \tag{PC13}
  \]
  The base-change formulas for inverse morphisms from
  \cite[\href{https://stacks.math.columbia.edu/tag/0017}{Tag 0017}]
  {stacks-project} prove the $\operatorname{IsIso}$ entry.  The entries for
  $\operatorname{IsQCQS}$ and
  $\operatorname{IsQuasiCompactOpenImmersion}$ are the componentwise
  conjunctions defined in
  \Cref{def:nagata-morphism-predicate-notation} into the cited entries.

  The composition implications (PC3) are the entries of
  \cite[\href{https://stacks.math.columbia.edu/tag/02WG}{Tag 02WG}]
  {stacks-project}; individually their tags are
  \[
  \begin{array}{c|c}
    \text{predicate}&\text{composition tag}\\ \hline
    \operatorname{IsOpenImmersion},\operatorname{IsClosedImmersion}
      &\text{02V0}\\
    \operatorname{IsQuasiCompact}&\text{01K6}\\
    \operatorname{QuasiSeparated},\operatorname{IsSeparated}&\text{01KU}\\
    \operatorname{FiniteType}&\text{01T3}\\
    \operatorname{FinitePresentation}&\text{01TR}\\
    \operatorname{IsProper}&\text{01W3}.
  \end{array}                                                       \tag{PC14}
  \]
  If $f^{-1}$ and $g^{-1}$ denote the inverses supplied by
  $\operatorname{IsIso}(f)$ and $\operatorname{IsIso}(g)$, then
  \[
  \begin{aligned}
    (f^{-1}\circ g^{-1})\circ(g\circ f)
      &=f^{-1}\circ(g^{-1}\circ g)\circ f
       =f^{-1}\circ f=\operatorname{id}_X,\\
    (g\circ f)\circ(f^{-1}\circ g^{-1})
      &=g\circ(f\circ f^{-1})\circ g^{-1}
       =g\circ g^{-1}=\operatorname{id}_Z,
  \end{aligned}
  \]
  which is the $\operatorname{IsIso}$ entry of (PC3).  The two conjunction
  predicates again follow by componentwise substitution.

  In the terminology fixed by
  \cite[\href{https://stacks.math.columbia.edu/tag/02KO}{Tag 02KO}]
  {stacks-project}, (PC5) is target locality.  The individual target-locality
  tags, assembled in
  \cite[\href{https://stacks.math.columbia.edu/tag/02WH}{Tag 02WH}]
  {stacks-project}, are
  \[
  \begin{array}{c|c}
    \text{predicate}&\text{target-locality tag}\\ \hline
    \operatorname{IsOpenImmersion}&\text{02L3}\\
    \operatorname{IsClosedImmersion}&\text{02L6}\\
    \operatorname{IsIso}&\text{02L4}\\
    \operatorname{IsQuasiCompact}&\text{02KQ}\\
    \operatorname{QuasiSeparated}&\text{02KR}\\
    \operatorname{IsSeparated}&\text{02KU}\\
    \operatorname{FiniteType}&\text{02KZ}\\
    \operatorname{FinitePresentation}&\text{02L0}\\
    \operatorname{IsProper}&\text{02L1}.
  \end{array}                                                       \tag{PC15}
  \]
  Substitution of conjunction definitions proves the remaining two entries.

  Equation (PC6) is
  \cite[\href{https://stacks.math.columbia.edu/tag/02V6}{Tag 02V6}]
  {stacks-project}; (PC7) is
  \cite[\href{https://stacks.math.columbia.edu/tag/0C4P}{Tag 0C4P}]
  {stacks-project}.  To prove (PC8), define the inverse morphisms by
  \[
  \begin{gathered}
    \alpha^{-1}\circ\alpha=\operatorname{id}_{X'},\qquad
    \alpha\circ\alpha^{-1}=\operatorname{id}_{X},\\
    \beta^{-1}\circ\beta=\operatorname{id}_{Y},\qquad
    \beta\circ\beta^{-1}=\operatorname{id}_{Y'}.
  \end{gathered}                                                  \tag{PC16}
  \]
  Suppose that $(\mathcal E,\iota)$ is a witness for
  $\operatorname{IsProjective}(f)$ in (P0), and define
  \[
    \mathcal E_{\beta}:=(\beta^{-1})^*\mathcal E\in
      \operatorname{QCoh}(\mathcal O_{Y'}),
  \]
  \[
  \begin{split}
    \mathscr C_{\beta,\mathcal E}:=
      \{c\in\operatorname{Hom}(&\mathbf P_Y(\mathcal E),
          \mathbf P_{Y'}(\mathcal E_{\beta}))\mid
      \operatorname{IsIso}(c)\ \wedge{}\\
      &\operatorname{pr}_{Y',\mathcal E_{\beta}}\circ c
        =\beta\circ\operatorname{pr}_{Y,\mathcal E}\}.
  \end{split}                                                     \tag{PC17}
  \]
  Relative-Proj base change
  \cite[\href{https://stacks.math.columbia.edu/tag/01O3}{Tag 01O3}]
  {stacks-project} applied to $\beta^{-1}\colon Y'\to Y$, together with
  (PC16), is the nonemptiness assertion
  $\mathscr C_{\beta,\mathcal E}\ne\varnothing$.  Define
  \[
    c_{\beta,\mathcal E}:=
      \operatorname{ch}(\mathscr C_{\beta,\mathcal E}),
    \qquad
    \iota_{\beta,\alpha}:=
      c_{\beta,\mathcal E}\circ\iota\circ\alpha.              \tag{PC18}
  \]
  Pullback preserves finite-type modules by
  \cite[\href{https://stacks.math.columbia.edu/tag/01B6}{Tag 01B6}]
  {stacks-project}, hence
  \[
    \mathcal E\text{ finite type}
      \Longrightarrow
    \mathcal E_{\beta}\text{ finite type}.                     \tag{PC19}
  \]
  Substitution of (PC17)--(PC18) and the witness equality
  $\operatorname{pr}_{Y,\mathcal E}\circ\iota=f$ is the equality
  \[
  \begin{aligned}
    \operatorname{pr}_{Y',\mathcal E_{\beta}}
      \circ\iota_{\beta,\alpha}
      &=\operatorname{pr}_{Y',\mathcal E_{\beta}}
        \circ c_{\beta,\mathcal E}\circ\iota\circ\alpha\\
      &=\beta\circ\operatorname{pr}_{Y,\mathcal E}
        \circ\iota\circ\alpha
       =\beta\circ f\circ\alpha
       =f^{\alpha,\beta}.
  \end{aligned}                                                  \tag{PC20}
  \]
  The definitions (P0)--(P0a) are the implications
  \[
    \operatorname{IsIso}(\alpha)
      \Longrightarrow\operatorname{IsClosedImmersion}(\alpha),
    \qquad
    \operatorname{IsIso}(c_{\beta,\mathcal E})
      \Longrightarrow
      \operatorname{IsClosedImmersion}(c_{\beta,\mathcal E}). \tag{PC21}
  \]
  Applying (PC3) twice to (PC21) and
  $\operatorname{IsClosedImmersion}(\iota)$ is the predicate
  \[
    \operatorname{IsClosedImmersion}(\iota_{\beta,\alpha}).   \tag{PC22}
  \]
  Equations (PC19), (PC20), and (PC22), substituted into (P0), prove
  \[
    \operatorname{IsProjective}(f)
      \Longrightarrow
    \operatorname{IsProjective}(f^{\alpha,\beta}).             \tag{PC23}
  \]
  Applying (PC23) to
  $(f^{\alpha,\beta},\alpha^{-1},\beta^{-1})$ and substituting (PC16) is
  \[
  \begin{aligned}
    (f^{\alpha,\beta})^{\alpha^{-1},\beta^{-1}}
      &=\beta^{-1}\circ(\beta\circ f\circ\alpha)
        \circ\alpha^{-1}\\
      &=(\beta^{-1}\circ\beta)\circ f\circ
        (\alpha\circ\alpha^{-1})
       =f,
  \end{aligned}
  \]
  and therefore the reverse implication in (PC8).  The counterexample (PC9) is
  \cite[\href{https://stacks.math.columbia.edu/tag/08J1}{Tag 08J1}]
  {stacks-project}.

  The uniqueness clause in the definition of a fibre product,
  \cite[\href{https://stacks.math.columbia.edu/tag/001V}{Tag 001V}]
  {stacks-project}, is exactly (PC10).  Finally,
  \cite[\href{https://stacks.math.columbia.edu/tag/023Q}{Tag 023Q}]
  {stacks-project} states that an open cover is a universal effective
  epimorphism.  Substitution of the objects and arrows defined in (PC11) into
  its Hom equalizer is precisely the bijection (PC12).
\end{proof}


\begin{lemma}[QCQS criteria for immersions and isomorphisms]
  \label{lemma:nagata-immersion-qcqs-criteria}
  \uses{def:nagata-morphism-predicate-notation}

  For a morphism $h\colon T\to Y$, the following implications hold:
  \[
  \begin{aligned}
    \operatorname{IsImmersion}(h)\wedge
      \operatorname{IsQuasiCompact}(h)
      &\Longrightarrow\operatorname{IsQCQS}(h),\\
    \operatorname{IsClosedImmersion}(h)
      &\Longrightarrow\operatorname{IsQCQS}(h),\\
    \operatorname{IsQuasiCompactOpenImmersion}(h)
      &\Longrightarrow\operatorname{IsQCQS}(h),\\
    \operatorname{IsIso}(h)
      &\Longrightarrow\operatorname{IsQCQS}(h).
  \end{aligned}                                                     \tag{QCI0}
  \]

  \paragraph{Mathlib relation.}
  The implications are derivable from the exact immersion, quasi-compact,
  and quasi-separated instances in
  \texttt{\detokenize{Mathlib.AlgebraicGeometry.Morphisms.Immersion}} and
  the corresponding closed/open-immersion modules.  In particular Mathlib
  exposes
  \texttt{\detokenize{AlgebraicGeometry.IsImmersion.isStableUnderBaseChange}},
  \texttt{\detokenize{AlgebraicGeometry.IsImmersion.isImmersion_iff_exists}},
  and
  \texttt{\detokenize{AlgebraicGeometry.IsImmersion.isImmersion_iff_exists_of_quasiCompact}};
  the four-way
  aggregate (QCI0) is local to the blueprint.
\end{lemma}

\begin{proof}
  By
  \cite[\href{https://stacks.math.columbia.edu/tag/01L7}{Tag 01L7}]
  {stacks-project}, the diagonal definition in
  \Cref{def:nagata-morphism-predicate-notation}, and
  \cite[\href{https://stacks.math.columbia.edu/tag/01K2}{Tag 01K2}]
  {stacks-project}, the separation chain is
  \[
  \begin{aligned}
    \operatorname{IsImmersion}(h)
      &\Longrightarrow\operatorname{IsSeparated}(h)\\
      &\Longleftrightarrow
        \operatorname{IsClosedImmersion}(\Delta_h)\\
      &\Longrightarrow
        \operatorname{IsQuasiCompact}(\Delta_h)\\
      &\Longleftrightarrow\operatorname{QuasiSeparated}(h).
  \end{aligned}                                                     \tag{QCI1}
  \]
  Substitution into the definition of $\operatorname{IsQCQS}$ gives
  \[
  \begin{aligned}
    \operatorname{IsImmersion}(h)\wedge
      \operatorname{IsQuasiCompact}(h)
      &\Longrightarrow
        \operatorname{IsQuasiCompact}(h)\wedge
        \operatorname{QuasiSeparated}(h)\\
      &\Longleftrightarrow\operatorname{IsQCQS}(h).
  \end{aligned}                                                     \tag{QCI2}
  \]
  For a closed immersion, choose the immersion factorization
  $(W,i,j)=(Y,h,\operatorname{id}_Y)$ in (P0a); for a quasi-compact open
  immersion, choose $(W,i,j)=(T,\operatorname{id}_T,h)$.  The corresponding
  predicate chains are
  \[
  \begin{aligned}
    \operatorname{IsClosedImmersion}(h)
      &\Longrightarrow
        \operatorname{IsImmersion}(h)\wedge
        \operatorname{IsQuasiCompact}(h)
       \Longrightarrow\operatorname{IsQCQS}(h),\\
    \operatorname{IsQuasiCompactOpenImmersion}(h)
      &\Longleftrightarrow
        \operatorname{IsOpenImmersion}(h)\wedge
        \operatorname{IsQuasiCompact}(h)\\
      &\Longrightarrow
        \operatorname{IsImmersion}(h)\wedge
        \operatorname{IsQuasiCompact}(h)
       \Longrightarrow\operatorname{IsQCQS}(h).
  \end{aligned}                                                     \tag{QCI3}
  \]
  Finally, the definitions of an isomorphism and of a closed immersion give
  \[
    \operatorname{IsIso}(h)
      \Longrightarrow\operatorname{IsClosedImmersion}(h)
      \Longrightarrow\operatorname{IsQCQS}(h).                    \tag{QCI4}
  \]
\end{proof}


\begin{definition}[Compactification] \checked
  \label{def:nagata-compactification}
  \uses{def:nagata-morphism-predicate-notation}
  \lean{AlgebraicGeometry.Compactification}

  Let $f\colon X\to S$ be a morphism of schemes.  A \emph{compactification}
  of $f$ is a quadruple $C=(\overline X,j,p,e)$ consisting of a scheme
  $\overline X$, morphisms
  \[
    X\xrightarrow{j}\overline X\xrightarrow{p}S,
  \]
  an equality $e\colon p\circ j=f$, a quasi-compact open immersion $j$, and
  a proper morphism $p$.
  Write $\mathbf{Sch}$ for the category specified by
  \[
    \operatorname{Ob}(\mathbf{Sch})
      :=\{T\mid T\text{ is a scheme}\},\qquad
    \operatorname{Hom}_{\mathbf{Sch}}(T_1,T_2)
      :=\{h\colon T_1\to T_2\mid h\text{ is a morphism of schemes}\},
  \]
  with composition and identities equal to composition and identities of
  morphisms of schemes.  Equivalently, define the type
  \[
    \operatorname{Compactification}(f)
      :=\left\{(\overline X,j,p,e)\ \middle|\;
      \begin{array}{l}
        \overline X\in\operatorname{Ob}(\mathbf{Sch}),\quad
        j\in\operatorname{Hom}_{\mathbf{Sch}}(X,\overline X),\\
        p\in\operatorname{Hom}_{\mathbf{Sch}}(\overline X,S),\quad
        e\colon p\circ j=f,\\
        \operatorname{IsQuasiCompactOpenImmersion}(j),\quad
        \operatorname{IsProper}(p)
      \end{array}\right\}.
  \]

  \paragraph{Formalization.}
  The bundled record is the project declaration
  \texttt{\detokenize{AlgebraicGeometry.Compactification}}
  (\texttt{\detokenize{WeilConjectures/Compactification/Defs.lean}}), with fields
  \texttt{\detokenize{Z}}, \texttt{\detokenize{g}}, \texttt{\detokenize{i}},
  \texttt{\detokenize{comm}}, \texttt{\detokenize{isProper}} and
  \texttt{\detokenize{isOpenImmersion}} corresponding to
  $\overline X$, $p$, $j$, $e$, $\operatorname{IsProper}(p)$ and the open-immersion half of
  $\operatorname{IsQuasiCompactOpenImmersion}(j)$.  The component predicates are the exact
  Mathlib interfaces \texttt{\detokenize{IsOpenImmersion}},
  \texttt{\detokenize{QuasiCompact}} and \texttt{\detokenize{IsProper}}.

  \textbf{This entry is deliberately not marked as formalized, because the Lean structure is
  strictly weaker than the definition above: it has no \texttt{\detokenize{QuasiCompact i}}
  field.}  The two notions agree exactly when $f$ is quasi-compact, by
  \texttt{\detokenize{Compactification.quasiCompact_openImmersion_iff}}, which states
  $\operatorname{QuasiCompact}(j)\iff\operatorname{QuasiCompact}(f)$ for any $C$,
  and they genuinely differ otherwise: an open immersion into a proper scheme need not be
  quasi-compact, so the Lean type is a strictly larger set of quadruples for general $f$.
  The field is deliberately absent rather than overlooked --- the class
  \texttt{\detokenize{HasCompactification}} below is consumed as a hypothesis by the
  six-functor formalism (\texttt{\detokenize{Etale.lowerShriek}} and the comparison
  $f_!\to f_*$), so adding a field would make that hypothesis harder to discharge for no gain
  at the use sites, which carry \texttt{\detokenize{QuasiCompact}} separately when they need
  it.  Every application of this definition inside this chapter is at a morphism of finite
  type, hence quasi-compact, where the two notions coincide.
\end{definition}

\begin{definition}[Compactifiability]
  \label{def:nagata-compactifiable}
  \uses{def:nagata-compactification}
  \lean{AlgebraicGeometry.HasCompactification}

  Let \(f\colon X\to S\) be a morphism of schemes.  Define the predicate
  \[
    \operatorname{Compactifiable}(f)
      \quad:\Longleftrightarrow\quad
    \operatorname{Compactification}(f)\neq\varnothing.
  \]
  The phrase ``$f$ is compactifiable'' abbreviates
  $\operatorname{Compactifiable}(f)$.

  \paragraph{Formalization.}
  This is the project class
  \texttt{\detokenize{AlgebraicGeometry.HasCompactification}}, defined as
  \texttt{\detokenize{Nonempty (Compactification f)}} --- literally the displayed
  nonemptiness predicate.  A choice of compactification is extracted by
  \texttt{\detokenize{HasCompactification.compactification}}, and
  \texttt{\detokenize{Compactification.hasCompactification}} is the converse packaging.  No
  declaration named \texttt{\detokenize{AlgebraicGeometry.Compactifiable}} exists; the
  \texttt{\detokenize{Has}}-prefixed class name is the Mathlib convention for this shape.

  Not marked as formalized, for the reason given in
  \Cref{def:nagata-compactification}: the underlying Lean structure omits quasi-compactness
  of $j$, so \texttt{\detokenize{HasCompactification f}} is implied by, but does not imply,
  $\operatorname{Compactifiable}(f)$ for a general $f$.  The two agree whenever $f$ is
  quasi-compact.
\end{definition}

\begin{definition}[Direct-image structure morphism]
  \label{def:nagata-direct-image-structure-map}
  \lean{AlgebraicGeometry.Scheme.Hom.app,
    AlgebraicGeometry.Scheme.Hom.appTop,
    AlgebraicGeometry.Scheme.Hom.comp_app,
    AlgebraicGeometry.IsSchemeTheoreticallyDominant}
  \leanok

  Let $a\colon A\to B$ be a morphism of schemes and let
  \(a^\flat\colon\mathcal O_B\to a_*\mathcal O_A\) be the adjoint
  structure-sheaf morphism of
  \Cref{def:nagata-morphism-predicate-notation}.  For every open subset
  $V\subseteq|B|$, define
  \[
    \varepsilon_{a,V}:=(a^\flat)_V\colon
      \mathcal O_B(V)\longrightarrow
      \mathcal O_A(|a|^{-1}(V)).
  \]
  Define the morphism of sheaves of rings $\varepsilon_a$ by the symbolic
  equality
  \[
    \varepsilon_a:=a^\flat\colon
      \mathcal O_B\longrightarrow a_*\mathcal O_A,
    \qquad
    (a_*\mathcal O_A)(V):=\mathcal O_A(|a|^{-1}(V)),
    \qquad
    (\varepsilon_a)_V:=\varepsilon_{a,V}.
  \]
  For $b\colon B\to C$, the identity $(b\circ a)_*=b_*a_*$ and the
  composition identity for locally ringed-space morphisms are
  \[
    \varepsilon_{b\circ a,V}
      =\varepsilon_{a,|b|^{-1}(V)}\circ\varepsilon_{b,V}
      \quad(V\subseteq|C|\text{ open}),
    \qquad
    \varepsilon_{b\circ a}
      =(b_*\varepsilon_a)\circ\varepsilon_b.
  \]
  Define
  \[
    \operatorname{SchematicallyDense}(a)
      \quad\Longleftrightarrow\quad
    \ker(\varepsilon_a)=0_{\mathcal O_B}.
  \]
  If \(A=\operatorname{Spec}(R_A)\) and
  \(B=\operatorname{Spec}(R_B)\) are affine, define the induced ring map
  \[
    a^\#:=\varepsilon_{a,|B|}\colon
      R_B=\mathcal O_B(|B|)
      \longrightarrow
      \mathcal O_A(|A|)=R_A.                                      \tag{DIM0}
  \]

  \paragraph{Formalization.}
  This definition introduces notation rather than a new object, so it has no single Lean
  counterpart and none should be invented: there is no declaration named
  \texttt{\detokenize{Scheme.Hom.directImageStructureMap}}, and $\varepsilon_a$ is the
  blueprint's adapter notation for an API that is already complete in Mathlib.

  The component $\varepsilon_{a,V}$ is \texttt{\detokenize{Scheme.Hom.app a V}}, of type
  $\Gamma(B,V)\to\Gamma(A,|a|^{-1}(V))$; it is an abbreviation for the application of the
  structure field \texttt{\detokenize{AlgebraicGeometry.PresheafedSpace.Hom.c}}, written
  \texttt{\detokenize{a.c}}, which is $\varepsilon_a=a^\flat$ itself.  There is no
  declaration named \texttt{\detokenize{Scheme.Hom.c}}.  Naturality of $\varepsilon_a$ in $V$
  is \texttt{\detokenize{Scheme.Hom.naturality}}, and the induced maps on stalks are
  \texttt{\detokenize{Scheme.Hom.stalkMap}}.

  The two composition identities are \texttt{\detokenize{Scheme.Hom.comp_app}} and
  \texttt{\detokenize{Scheme.Hom.comp_appTop}}.  Both are written in diagrammatic order in
  Lean, so the first reads ``the section map of $f$ then $g$ at $U$ is the section map of $g$
  at $U$ followed by that of $f$ at $|g|^{-1}(U)$'', which is the displayed
  $\varepsilon_{b\circ a,V}=\varepsilon_{a,|b|^{-1}(V)}\circ\varepsilon_{b,V}$ with $a=f$ and
  $b=g$.

  $\operatorname{SchematicallyDense}(a)$ is exactly
  \texttt{\detokenize{AlgebraicGeometry.IsSchemeTheoreticallyDominant a}}, and the displayed
  characterization $\ker(\varepsilon_a)=0$ is
  \texttt{\detokenize{isSchemeTheoreticallyDominant_iff}}, phrased as the vanishing of the
  ideal sheaf \texttt{\detokenize{Scheme.Hom.ker a}} rather than of a literal kernel subsheaf.
  Finally (DIM0), the induced ring map $a^{\#}$ on global sections, is
  \texttt{\detokenize{Scheme.Hom.appTop}}, which is defined for arbitrary schemes and
  specializes to the affine case.
\end{definition}

\begin{definition}[Literal ideal sheaves and Mathlib ideal-sheaf data]
  \label{def:nagata-ideal-sheaf-data-adapter}
  \uses{def:nagata-direct-image-structure-map,
    def:nagata-qcoh-graded-algebra-interface}

  Let $X$ be a scheme.  If
  $\mathcal I\subseteq\mathcal O_X$ is a quasi-coherent sheaf of ideals,
  define its affine ideal datum by
  \[
    \mathsf D_X(\mathcal I).\operatorname{ideal}(U)
      :=\Gamma(U,\mathcal I|_U)
      \triangleleft\Gamma(U,\mathcal O_X)                         \tag{MLID0}
  \]
  for each affine open $U\subseteq X$, with localization compatibility
  furnished by quasi-coherence.  Conversely, for Mathlib ideal-sheaf data
  $I\colon X.\operatorname{IdealSheafData}$, define the associated literal
  ideal sheaf by
  \[
    i_I\colon I.\operatorname{subscheme}\hookrightarrow X
      \quad\text{for the canonical closed immersion},
    \qquad
    \mathsf S_X(I)
      :=\ker\!\left(
        \varepsilon_{i_I}\colon
        \mathcal O_X\longrightarrow
        (i_I)_*
          \mathcal O_{I.\operatorname{subscheme}}
      \right).                                                    \tag{MLID1}
  \]
  Fix the canonical affine-local identifications
  \[
    \mathsf D_X(\mathsf S_X(I))=I,
    \qquad
    \mathsf S_X(\mathsf D_X(\mathcal I))\cong\mathcal I.          \tag{MLID2}
  \]
  Henceforth a literal quasi-coherent ideal and its datum are silently
  transported through (MLID2).  In particular, define
  \[
  \begin{aligned}
    V(\mathcal I)&:=\mathsf D_X(\mathcal I).\operatorname{support},\\
    \mathcal I\mathcal O_Y
      &\cong\mathsf S_Y\bigl(
        \mathsf D_X(\mathcal I).\operatorname{comap}(f)\bigr)
        &&(f\colon Y\to X),\\
    \operatorname{KerIdeal}_X(f)
      &:=f.\operatorname{ker}
        &&(f\colon Y\to X),\\
    \mathsf D_X\!\left(\ker(\varepsilon_f)\right)
      &=\operatorname{KerIdeal}_X(f)
        &&(f\colon Y\to X\text{ quasi-compact}),\\
    \operatorname{FiniteTypeIdeal}_X(\mathcal I)
      &:\Longleftrightarrow
        \forall U\subseteq X\text{ affine},\quad
        \mathsf D_X(\mathcal I).\operatorname{ideal}(U)
          \text{ is finitely generated}.
  \end{aligned}                                                    \tag{MLID3}
  \]

  \paragraph{Mathlib relation.}
  The exact Mathlib representation is
  \texttt{\detokenize{AlgebraicGeometry.Scheme.IdealSheafData}} with
  \texttt{\detokenize{IdealSheafData.ideal}},
  \texttt{\detokenize{support}},
  \texttt{\detokenize{comap}},
  \texttt{\detokenize{map}},
  \texttt{\detokenize{radical}}, and its lattice, multiplication, and power
  instances.  The kernel datum in (MLID3) is the exact declaration
  \texttt{\detokenize{AlgebraicGeometry.Scheme.Hom.ker}}.  Its equality with
  the componentwise literal kernel under the displayed quasi-compactness
  hypothesis is
  \texttt{\detokenize{AlgebraicGeometry.Scheme.Hom.ker_apply}}.  Mathlib
  intentionally does not represent these data as literal
  subsheaves of $\mathcal O_X$; (MLID0)--(MLID3) are the adapter used by this
  blueprint.
\end{definition}

\begin{lemma}[Compatibility of the ideal-sheaf data adapter]
  \label{lemma:nagata-ideal-sheaf-data-adapter}
  \uses{def:nagata-ideal-sheaf-data-adapter,
    def:nagata-morphism-predicate-notation}

  The transports $\mathsf D_X$ and $\mathsf S_X$ preserve restriction to
  opens, inclusion, sums, intersections, products, powers, radicals,
  supports, and map/comap along scheme morphisms.  More precisely, all
  corresponding comparison arrows are isomorphisms and, on every affine open
  $U=\operatorname{Spec}(A)$, reduce to the identity on the associated ideal
  of $A$.  Moreover
  \[
  \begin{aligned}
    \operatorname{FiniteTypeIdeal}_X(\mathcal I)
      &\Longleftrightarrow
        \mathsf D_X(\mathcal I).\operatorname{ideal}(U)
          \text{ is finitely generated for every member of one affine cover},\\
    \operatorname{FiniteTypeIdeal}_X(\mathcal I)
      &\Longleftrightarrow
        \operatorname{FiniteTypeModule}_X(\mathcal I).
  \end{aligned}
                                                                    \tag{MLID4}
  \]
  Henceforth the phrase ``finite-type quasi-coherent ideal'' abbreviates
  either equivalent predicate in (MLID4).
\end{lemma}

\begin{proof}
  Each assertion is local on $X$.  On an affine open, (MLID0)--(MLID2)
  identify both sides with the same ordinary ideal and the statement becomes
  the corresponding ideal identity.  Compatibility with localization makes
  the affine identities agree on overlaps, and quasi-coherent sheaf gluing
  gives the global comparisons.  The first final equivalence follows by
  refining intersections by principal affine opens and using that
  localization preserves finite generation.  The second is the affine-local
  finite-generation criterion for a quasi-coherent module sheaf.
\end{proof}

\begin{definition}[Relative spectrum of a quotient algebra]
  \label{def:nagata-quotient-relative-spec}
  \uses{def:nagata-ideal-sheaf-data-adapter}

  Let $\mathcal I\subseteq\mathcal O_X$ be a quasi-coherent ideal sheaf and
  let $I:=\mathsf D_X(\mathcal I)$.  Define the quotient-relative spectrum
  and its structure morphism by the exact Mathlib subscheme construction
  \[
    \operatorname{Spec}_X(\mathcal O_X/\mathcal I)
      :=I.\operatorname{subscheme},
    \qquad
    \operatorname{Spec}_X(q_{\mathcal I})
      :=i_I\colon
      I.\operatorname{subscheme}\longrightarrow X.              \tag{MLRS0}
  \]
  If $\mathcal I\subseteq\mathcal J$, the quotient map
  $q_{\mathcal I\mathcal J}\colon
    \mathcal O_X/\mathcal I\twoheadrightarrow
    \mathcal O_X/\mathcal J$
  induces the contravariant morphism
  \[
    \operatorname{Spec}_X(q_{\mathcal I\mathcal J})\colon
      \operatorname{Spec}_X(\mathcal O_X/\mathcal J)
        \longrightarrow
      \operatorname{Spec}_X(\mathcal O_X/\mathcal I),            \tag{MLRS1}
  \]
  defined by
  \texttt{\detokenize{Scheme.IdealSheafData.inclusion}} after transport
  through \Cref{def:nagata-ideal-sheaf-data-adapter}.

  \paragraph{Mathlib relation.}
  Formula (MLRS0) is exactly
  \texttt{\detokenize{Scheme.IdealSheafData.subscheme}} and
  \texttt{\detokenize{Scheme.IdealSheafData.subschemeι}}.  The adjunction
  used below is
  \texttt{\detokenize{AlgebraicGeometry.Scheme.kerAdjunction}}.  This
  quotient adapter is sufficient for every relative-\texttt{Spec} expression
  in the chapter; no general relative spectrum of an arbitrary
  quasi-coherent algebra is assumed.
\end{definition}

\begin{lemma}[Quotient-relative-spectrum interface]
  \label{lemma:nagata-quotient-relative-spec}
  \uses{def:nagata-quotient-relative-spec,
    lemma:nagata-ideal-sheaf-data-adapter}

  The construction in \Cref{def:nagata-quotient-relative-spec} satisfies:
  \begin{enumerate}
    \item on $U=\operatorname{Spec}(A)$, if
      $\mathcal I|_U=\widetilde I$, then
      \[
        \operatorname{Spec}_X(\mathcal O_X/\mathcal I)\times_XU
          \cong\operatorname{Spec}(A/I);                         \tag{MLRS2}
      \]
    \item for $h\colon T\to X$, a factorization through (MLRS0) exists if
      and only if
      $\mathsf D_X(\mathcal I)\leq\operatorname{KerIdeal}_X(h)$, and it is
      unique because
      $i_I$ is a monomorphism;
    \item the maps (MLRS1) preserve identities and composition;
    \item for $f\colon Y\to X$ there is a natural cartesian isomorphism
      \[
        \operatorname{Spec}_Y
          (\mathcal O_Y/\mathcal I\mathcal O_Y)
          \xrightarrow{\sim}
        \operatorname{Spec}_X(\mathcal O_X/\mathcal I)\times_XY. \tag{MLRS3}
      \]
  \end{enumerate}
\end{lemma}

\begin{proof}
  The affine description is the defining chart of
  $I.\operatorname{subscheme}$.  The factorization statement is the Hom-set
  form of $X.\operatorname{kerAdjunction}$, with uniqueness supplied by the
  closed immersion $i_I$.  Functoriality is the
  corresponding identity and composition law for
  $\operatorname{IdealSheafData.inclusion}$.  Finally, both sides of (MLRS3)
  restrict on affine pullback charts to the spectrum of the same extended
  quotient ideal, and these affine isomorphisms glue.
\end{proof}

\begin{definition}[Relative spectrum with a quotient presentation]
  \label{def:nagata-presented-relative-spec}
  \uses{def:nagata-quotient-relative-spec,
    lemma:nagata-quotient-relative-spec}

  Let $\mathcal A$ be a quasi-coherent $\mathcal O_X$-algebra.  A
  \emph{quotient presentation} of $\mathcal A$ is a pair
  \[
    (\mathcal I,\chi),\qquad
    \mathcal I\subseteq\mathcal O_X\text{ a quasi-coherent ideal},qquad
    \chi\colon\mathcal O_X/\mathcal I\xrightarrow{\sim}\mathcal A.
                                                                    \tag{MLPSP0}
  \]
  For such a presentation define
  \[
    \operatorname{Spec}_X(\mathcal A;\mathcal I,\chi)
      :=\operatorname{Spec}_X(\mathcal O_X/\mathcal I),           \tag{MLPSP1}
  \]
  with structure morphism (MLRS0).  A morphism between two
  quotient-presented algebras is transported across the chosen presentation
  isomorphisms and then interpreted by (MLRS1).  If the presentation is
  clear from the preceding display, the final two arguments in (MLPSP1) are
  suppressed.

  For $f\colon Y\to X$, let
  \[
    \chi_f\colon
      \mathcal O_Y/\mathcal I\mathcal O_Y
        \xrightarrow{\sim}f^*\mathcal A
  \]
  be the quotient comparison transported through $f^*\chi$.  Define the
  cartesian comparison and the morphism associated to the adjunction map
  $f^\sharp_{\mathcal A}\colon\mathcal A\to f_*f^*\mathcal A$ by
  \[
  \begin{gathered}
    \operatorname{bc}_{f,\mathcal A}\colon
      \operatorname{Spec}_Y
        (f^*\mathcal A;\mathcal I\mathcal O_Y,\chi_f)
        \xrightarrow{\sim}
      \operatorname{Spec}_X(\mathcal A;\mathcal I,\chi)\times_XY,\\
    \operatorname{Spec}(f^\sharp_{\mathcal A})
      :=\operatorname{pr}_1\circ\operatorname{bc}_{f,\mathcal A}\colon
      \operatorname{Spec}_Y
        (f^*\mathcal A;\mathcal I\mathcal O_Y,\chi_f)
        \longrightarrow
      \operatorname{Spec}_X(\mathcal A;\mathcal I,\chi).
                                                                    \tag{MLPSP1a}
  \end{gathered}
  \]
  Thus this cross-base notation is also defined only for a displayed
  quotient presentation; (MLPSP1a) is (MLRS3) followed by a fibre-product
  projection.

  \paragraph{Mathlib relation.}
  This is not a general Mathlib relative-spectrum construction.  It is
  transport along algebra isomorphisms followed by the exact declarations
  \texttt{\detokenize{Scheme.IdealSheafData.subscheme}} and
  \texttt{\detokenize{Scheme.IdealSheafData.inclusion}}; its base-change
  map is the blueprint interface (MLRS3).  It covers every non-literal
  quotient-relative-spectrum notation used in this chapter.
\end{definition}

\begin{lemma}[Recovery of a closed immersion from its direct image algebra]
  \label{lemma:nagata-closed-immersion-relative-spec-recovery}
  \uses{def:nagata-direct-image-structure-map,
    def:nagata-presented-relative-spec,
    lemma:nagata-quotient-relative-spec}

  Let $i\colon Z\hookrightarrow X$ be a closed immersion, put
  \[
    \mathcal K_i:=\ker(\varepsilon_i),
    \qquad
    \chi_i\colon
      \mathcal O_X/\mathcal K_i\xrightarrow{\sim}i_*\mathcal O_Z,
    \qquad
    \chi_i\circ q_{\mathcal K_i}=\varepsilon_i.                 \tag{MLPSP2}
  \]
  Then $(\mathcal K_i,\chi_i)$ is a quotient presentation and there is a
  canonical isomorphism over $X$
  \[
    \rho_i\colon Z\xrightarrow{\sim}
      \operatorname{Spec}_X
        (i_*\mathcal O_Z;\mathcal K_i,\chi_i),
    \qquad
    \operatorname{Spec}_X(\varepsilon_i)\circ\rho_i=i.          \tag{MLPSP3}
  \]
  Under (MLPSP1), transport along $\chi_i$ is the canonical isomorphism from
  the right side of (MLPSP3) to
  $\operatorname{Spec}_X(\mathcal O_X/\mathcal K_i)$.

  \paragraph{Mathlib relation.}
  The quotient presentation and recovery isomorphism are supplied by
  \texttt{\detokenize{AlgebraicGeometry.IsClosedImmersion.overEquivIdealSheafData}}
  together with \texttt{\detokenize{AlgebraicGeometry.Scheme.kerAdjunction}}.
  The notation in (MLPSP3) is the blueprint's transport wrapper.
\end{lemma}

\begin{proof}
  On each affine open of $X$, a closed immersion is induced by a surjective
  ring homomorphism.  The first isomorphism theorem gives (MLPSP2), and its
  spectrum identifies $Z$ with the quotient spectrum.  These affine
  isomorphisms commute with localization and glue.  The structure-map
  equation in (MLPSP3) holds on every affine chart and therefore globally.
  Equivalently, it is the counit equation of the kernel/subscheme
  adjunction.
\end{proof}

\begin{definition}[Reduced closed subscheme supported on a closed subset]
  \label{def:nagata-reduced-closed-subscheme}
  \uses{def:nagata-ideal-sheaf-data-adapter,
    def:nagata-quotient-relative-spec,
    lemma:nagata-quotient-relative-spec}
  \lean{}

  Let $X$ be a scheme and let $F\subseteq |X|$ be closed.  For an open
  subset $V\subseteq X$, $s\in\mathcal O_X(V)$, and $x\in V$, define
  \[
    s(x):=\bigl(\mathcal O_X(V)\longrightarrow
      \mathcal O_{X,x}\longrightarrow\kappa(x)\bigr)(s).
  \]
  Define the sheaf of ideals $\mathcal I_{X,F}\subseteq\mathcal O_X$ and
  its quotient map by
  \[
  \begin{aligned}
    \mathcal I_{X,F}(V)
      &:=\{s\in\mathcal O_X(V)\mid
          (\forall x\in F\cap V)\ s(x)=0\},\\
    \pi_{X,F}&\colon\mathcal O_X\twoheadrightarrow
      \mathcal O_X/\mathcal I_{X,F}.
  \end{aligned}                                                     \tag{R0}
  \]
  For every affine open $V=\operatorname{Spec}(A)\subseteq X$, (R0) is
  the equality
  \[
    \mathcal I_{X,F}|_V
      =\widetilde{\bigcap_{\mathfrak p\in F\cap\operatorname{Spec}(A)}
          \mathfrak p},                                           \tag{R1}
  \]
  where the intersection over the empty set is $A$.  In particular,
  $\mathcal I_{X,F}$ is a radical quasi-coherent ideal sheaf.  Define
  \[
    F_{\mathrm{red}}^X
      :=\operatorname{Spec}_X(\mathcal O_X/\mathcal I_{X,F}),
    \qquad
    i_{X,F}\colon F_{\mathrm{red}}^X\hookrightarrow X              \tag{R2}
  \]
  to be the closed immersion classified by $\pi_{X,F}$.  Thus
  \[
    |F_{\mathrm{red}}^X|=F,
    \qquad
    \ker(\varepsilon_{i_{X,F}})=\mathcal I_{X,F}.
  \]
  For $F=|X|$, define
  \[
    X_{\mathrm{red}}:=|X|_{\mathrm{red}}^X,
    \qquad
    \mathcal N_X:=\mathcal I_{X,|X|}=\sqrt{0_{\mathcal O_X}}.
  \]
  For a subset $E\subseteq|X|$, define
  \[
    \overline E_{\mathrm{red}}^X
      :=(\overline E^{\,|X|})_{\mathrm{red}}^X,
  \]
  whenever the topological closure $\overline E^{\,|X|}$ is used below.

  \paragraph{Mathlib relation.}
  Under \Cref{def:nagata-ideal-sheaf-data-adapter},
  $\mathcal I_{X,F}$ is
  \texttt{\detokenize{Scheme.IdealSheafData.vanishingIdeal F}}, and (R2) is
  its exact \texttt{\detokenize{subscheme}} with inclusion
  \texttt{\detokenize{subschemeι}}.  Support is identified by
  \texttt{\detokenize{IdealSheafData.coe_support_vanishingIdeal}}; the case
  $F=|X|$ is related to \texttt{\detokenize{Scheme.nilradical}} and
  \texttt{\detokenize{IdealSheafData.vanishingIdeal_top}}.  The proposed
  bundled name \texttt{\detokenize{Scheme.reducedClosedSubschemeOfClosedSet}}
  is not a current Mathlib declaration.
\end{definition}

\begin{definition}[Scheme-theoretically dense compactification]
  \label{def:nagata-dense-compactification}
  \uses{def:nagata-compactification,
    def:nagata-direct-image-structure-map}
  \lean{AlgebraicGeometry.IsSchemeTheoreticallyDominant,
    AlgebraicGeometry.isSchemeTheoreticallyDominant_iff}
  \leanok

  Let \(f\colon X\to S\) be a morphism of schemes, and let
  $C=(\overline X,j,p,e)\in
  \operatorname{Compactification}(f)$ be as in
  \Cref{def:nagata-compactification}.  By
  \Cref{def:nagata-direct-image-structure-map}, define
  \[
    \varepsilon_j:=j^\flat\colon\mathcal O_{\overline X}\longrightarrow
      j_*\mathcal O_X,
    \qquad
    (\varepsilon_j)_V=\varepsilon_{j,V}\colon
      \mathcal O_{\overline X}(V)\longrightarrow
      \mathcal O_X(|j|^{-1}(V)).
  \]
  Define
  \[
    \operatorname{Dense}(C)
      \quad\Longleftrightarrow\quad
    \operatorname{SchematicallyDense}(j)
      \quad\Longleftrightarrow\quad
    \ker(\varepsilon_j)=0_{\mathcal O_{\overline X}}.
  \]

  \paragraph{Formalization.}
  The density condition itself is formalized exactly: $\operatorname{Dense}(C)$ is
  \texttt{\detokenize{AlgebraicGeometry.IsSchemeTheoreticallyDominant c.i}} applied to the
  open immersion field of the Lean \texttt{\detokenize{Compactification}} record, and the
  displayed kernel characterization $\ker(\varepsilon_j)=0$ is
  \texttt{\detokenize{isSchemeTheoreticallyDominant_iff}}, stated as the vanishing of the
  ideal sheaf \texttt{\detokenize{Scheme.Hom.ker j}}.  Both are Mathlib declarations; no wrapper
  predicate on \texttt{\detokenize{Compactification}} is defined, and none is needed, since
  the condition is a predicate on $j$ alone.

  The quasi-compactness caveat of \Cref{def:nagata-compactification} applies to the ambient
  quadruple $C$, not to this predicate.
\end{definition}

\begin{definition}[Scheme-theoretic image]
  \label{def:nagata-scheme-theoretic-image}
  \uses{def:nagata-direct-image-structure-map,
    def:nagata-quotient-relative-spec}
  \lean{AlgebraicGeometry.Scheme.Hom.ker,
    AlgebraicGeometry.Scheme.Hom.image,
    AlgebraicGeometry.Scheme.Hom.imageι}
  \leanok

  Let \(g\colon T\to Y\) satisfy
  \(\operatorname{IsQCQS}(g)\).  By
  \Cref{def:nagata-direct-image-structure-map}, define
  \[
    \varepsilon_g:=g^\flat\colon\mathcal O_Y\longrightarrow g_*\mathcal O_T,
    \qquad
    (\varepsilon_g)_V=\varepsilon_{g,V},
    \qquad
    \mathcal K_g:=\ker(\varepsilon_g)\subseteq\mathcal O_Y,
    \qquad
    \pi_g\colon\mathcal O_Y\twoheadrightarrow
      \mathcal O_Y/\mathcal K_g.
  \]
  For every open $V\subseteq|Y|$, the definitions of kernel and quotient are
  the component equalities
  \[
    \mathcal K_g(V)
      =\ker\!\left(
        \varepsilon_{g,V}\colon\mathcal O_Y(V)\to
        \mathcal O_T(|g|^{-1}(V))
      \right),
    \qquad
    (\pi_g)_V(s)=[s].
  \]
  The conclusion of
  \cite[\href{https://stacks.math.columbia.edu/tag/01R5}{Tag 01R5}]
  {stacks-project}, applied to the quasi-compact and quasi-separated morphism
  $g$, is
  \[
    \mathcal K_g\text{ is quasi-coherent}.
  \]
  Define
  \[
    \operatorname{SIm}_Y(g):=
      \operatorname{Spec}_Y(\mathcal O_Y/\mathcal K_g),
    \qquad
    i_g\colon\operatorname{SIm}_Y(g)\hookrightarrow Y
  \]
  where $i_g$ is the closed immersion classified by $\pi_g$.  Define
  $\chi_g$ to be its structure isomorphism:
  \[
    \chi_g\colon\mathcal O_Y/\mathcal K_g
      \xrightarrow{\ \sim\ }(i_g)_*
      \mathcal O_{\operatorname{SIm}_Y(g)},
    \qquad
    \varepsilon_{i_g}=\chi_g\circ\pi_g.
  \]
  When the displayed codomain of $g$ is $Y$, define the abbreviated notation
  \[
    \operatorname{SIm}(g):=\operatorname{SIm}_Y(g).
  \]

  \paragraph{Formalization.}
  $\operatorname{SIm}_Y(g)$ is \texttt{\detokenize{AlgebraicGeometry.Scheme.Hom.image g}},
  reducibly \texttt{\detokenize{(Scheme.Hom.ker g).subscheme}}, and $i_g$ is
  \texttt{\detokenize{Scheme.Hom.imageι g}}.  The proposed name
  \texttt{\detokenize{Scheme.Hom.schemeTheoreticImage}} is a blueprint alias for those, not a
  Mathlib declaration.  The ideal sheaf $\mathcal K_g$ is
  \texttt{\detokenize{Scheme.Hom.ker g}}.  The structure isomorphism $\chi_g$ has no Lean
  counterpart as an isomorphism; its content --- that the closed immersion $i_g$ cuts out
  exactly $\mathcal K_g$, i.e.\ $\ker(\varepsilon_{i_g})=\mathcal K_g$ --- is the equation
  \texttt{\detokenize{Scheme.IdealSheafData.ker_subschemeι}}, and that is the form every
  downstream argument in this chapter uses.

  Two differences of presentation are worth recording, neither of them a gap.

  First, Mathlib carries no quasi-coherence hypothesis here.  The blueprint assumes
  $\operatorname{IsQCQS}(g)$ so that Tag 01R5 makes $\mathcal K_g$ quasi-coherent, and then
  forms $\operatorname{Spec}_Y$ of the quotient.  Mathlib's
  \texttt{\detokenize{Scheme.Hom.ker}} lands in
  \texttt{\detokenize{Scheme.IdealSheafData}}, which encodes a compatible family of ideals on
  affine opens and is therefore quasi-coherent by construction, so
  \texttt{\detokenize{Scheme.Hom.image g}} is defined for \emph{every} $g$ with no
  hypothesis.  The Lean object is thus available in strictly greater generality, and agrees
  with the blueprint's under $\operatorname{IsQCQS}(g)$.

  Second, the relative spectrum does not appear: the closed subscheme cut out by an ideal
  sheaf is constructed directly as
  \texttt{\detokenize{Scheme.IdealSheafData.subscheme}}, with inclusion
  \texttt{\detokenize{subschemeι}}, rather than as
  $\operatorname{Spec}_Y(\mathcal O_Y/\mathcal K_g)$.
\end{definition}

\begin{definition}[Canonical factor through a scheme-theoretic image]
  \label{def:nagata-scheme-image-factor}
  \uses{def:nagata-direct-image-structure-map,
    def:nagata-scheme-theoretic-image,
    lemma:nagata-quotient-relative-spec}
  \lean{AlgebraicGeometry.Scheme.Hom.toImage,
    AlgebraicGeometry.Scheme.Hom.toImage_imageι}
  \leanok

  For $g\colon T\to Y$ and
  $(\mathcal K_g,\pi_g,i_g,\chi_g)$ as in
  \Cref{def:nagata-scheme-theoretic-image}, define
  $\overline\varepsilon_g$ to be the unique morphism of
  $\mathcal O_Y$-algebras satisfying
  \[
    \overline\varepsilon_g\colon
      \mathcal O_Y/\mathcal K_g\longrightarrow g_*\mathcal O_T,
    \qquad
    \overline\varepsilon_g\circ\pi_g=\varepsilon_g,
    \qquad
    \ker(\overline\varepsilon_g)=0.
  \]
  Under the relative-spectrum adjunction
  \[
    \operatorname{Hom}_Y
      \!\left(T,\operatorname{Spec}_Y(\mathcal O_Y/\mathcal K_g)\right)
      \cong
    \operatorname{Hom}_{\mathcal O_Y\text{-Alg}}
      \!\left(\mathcal O_Y/\mathcal K_g,g_*\mathcal O_T\right),
  \]
  define
  \[
    g^{\mathrm{im}}\colon T\longrightarrow\operatorname{SIm}_Y(g)
  \]
  to be the morphism corresponding to $\overline\varepsilon_g$.  Define the
  equality witness $e_g^{\mathrm{im}}$ and its sheaf equality by
  \[
    e_g^{\mathrm{im}}\colon i_g\circ g^{\mathrm{im}}=g,
    \qquad
    \varepsilon_g
      =\overline\varepsilon_g\circ\pi_g
      =\bigl((i_g)_*\varepsilon_{g^{\mathrm{im}}}\bigr)
         \circ\chi_g\circ\pi_g.
  \]

  \paragraph{Formalization.}
  $g^{\mathrm{im}}$ is \texttt{\detokenize{AlgebraicGeometry.Scheme.Hom.toImage g}}, and the
  equality witness $e_g^{\mathrm{im}}\colon i_g\circ g^{\mathrm{im}}=g$ is
  \texttt{\detokenize{Scheme.Hom.toImage_imageι}}, written in diagrammatic order as
  ``\texttt{\detokenize{g.toImage}} then \texttt{\detokenize{g.imageι}} equals
  \texttt{\detokenize{g}}''; the reassociated form
  \texttt{\detokenize{toImage_imageι_assoc}} is provided for rewriting.  The proposed name
  \texttt{\detokenize{toSchemeTheoreticImage}} is not a declaration.

  The Lean construction does not go through the relative-spectrum adjunction: it is the
  universal property of \texttt{\detokenize{Scheme.IdealSheafData.subscheme}}, packaged
  categorically by \texttt{\detokenize{Scheme.kerAdjunction}} as an adjunction over $Y$
  between the closed-subscheme functor
  \texttt{\detokenize{Scheme.IdealSheafData.subschemeFunctor}}, taking
  $\mathcal K$ to $V(\mathcal K)$, and the kernel functor
  \texttt{\detokenize{Scheme.kerFunctor}}, taking $h$ to $\ker(\varepsilon_h)$.  So
  $\overline\varepsilon_g$ is not named separately; the morphism it classifies is.
\end{definition}

\begin{definition}[Comparison map for a closed factorization]
  \label{def:nagata-scheme-image-comparison}
  \uses{def:nagata-choice-notation,
    def:nagata-closed-subscheme-class,
    def:nagata-direct-image-structure-map,
    def:nagata-scheme-theoretic-image}
  \lean{AlgebraicGeometry.IsClosedImmersion.lift,
    AlgebraicGeometry.IsClosedImmersion.lift_fac,
    AlgebraicGeometry.Scheme.Hom.ker_le_ker_imageι_comp}
  \leanok

  Let $g\colon T\to Y$ be as in
  \Cref{def:nagata-scheme-theoretic-image}.  Suppose that
  $h\colon T\to Z$, that $z\colon Z\hookrightarrow Y$ is a closed
  immersion, and that $e_{z,h}\colon g=z\circ h$ is a specified equality.
  Define
  \[
    \mathcal K_z:=\ker(\varepsilon_z),
    \qquad
    \pi_z\colon\mathcal O_Y\twoheadrightarrow
      \mathcal O_Y/\mathcal K_z.
  \]
  Let
  \[
    r_z\colon\operatorname{Spec}_Y(\mathcal O_Y/\mathcal K_z)
      \hookrightarrow Y
  \]
  be the closed immersion classified by $\pi_z$.  The equality
  $\mathcal K_z=\ker(\varepsilon_z)$ is the class equality
  \[
    [\,Z\xhookrightarrow{z}Y\,]
      =[\,\operatorname{Spec}_Y(\mathcal O_Y/\mathcal K_z)
        \xhookrightarrow{r_z}Y\,]
      \quad\text{in }\operatorname{ClosedSub}(Y).                \tag{SIC0}
  \]
  Define
  \[
    \mathscr S_z:=
      \left\{\theta\in
        \operatorname{Hom}\!\left(
          Z,\operatorname{Spec}_Y(\mathcal O_Y/\mathcal K_z)\right)
        \ \middle|\
        \operatorname{IsIso}(\theta)\ \wedge\
        r_z\circ\theta=z\right\}.
  \]
  By (CSub3) of \Cref{def:nagata-closed-subscheme-class}, (SIC0) is the
  existence statement $\mathscr S_z\ne\varnothing$.  Define the chosen
  witness
  \[
    \sigma_z:=\operatorname{ch}(\mathscr S_z)\in\mathscr S_z,
    \qquad
    \sigma_z\colon Z\xrightarrow{\ \sim\ }
      \operatorname{Spec}_Y(\mathcal O_Y/\mathcal K_z),
    \qquad
    r_z\circ\sigma_z=z.
  \]
  If $\theta\in\mathscr S_z$, then
  \[
    r_z\circ\theta=z=r_z\circ\sigma_z
      \Longrightarrow\theta=\sigma_z,
  \]
  where the implication is cancellation against the monomorphism $r_z$.
  For every open $V\subseteq Y$, substitution of $e_{z,h}\colon g=z\circ h$
  into \Cref{def:nagata-direct-image-structure-map} is the component equality
  \[
    (\varepsilon_g)_V
      =(\varepsilon_{z\circ h})_V
      =(\varepsilon_h)_{z^{-1}(V)}\circ(\varepsilon_z)_V.
  \]
  Consequently, for $s\in\mathcal K_z(V)$,
  \[
    (\varepsilon_g)_V(s)
      =(\varepsilon_h)_{z^{-1}(V)}
        \bigl((\varepsilon_z)_V(s)\bigr)
      =(\varepsilon_h)_{z^{-1}(V)}(0)=0,
  \]
  The displayed vanishing for every $V$ and every
  $s\in\mathcal K_z(V)$ is the sheaf inclusion
  \[
    \mathcal K_z\subseteq\mathcal K_g.
  \]
  Define
  \[
    q_{g,z}\colon\mathcal O_Y/\mathcal K_z
      \twoheadrightarrow\mathcal O_Y/\mathcal K_g,
    \qquad
    q_{g,z}\circ\pi_z=\pi_g.
  \]
  Let
  \[
    \operatorname{Spec}_Y(q_{g,z})\colon
      \operatorname{SIm}_Y(g)\longrightarrow
      \operatorname{Spec}_Y(\mathcal O_Y/\mathcal K_z)
  \]
  denote the morphism induced contravariantly by $q_{g,z}$.  Define
  \[
    c_{g,z}:=\sigma_z^{-1}\circ\operatorname{Spec}_Y(q_{g,z})
      \colon\operatorname{SIm}_Y(g)\longrightarrow Z.
  \]
  Define the equality witness $e_{g,z}^{\mathrm{cmp}}$ by
  \[
    e_{g,z}^{\mathrm{cmp}}\colon z\circ c_{g,z}=i_g,
    \qquad
    z\circ c_{g,z}
      =(r_z\circ\sigma_z)\circ\sigma_z^{-1}
        \circ\operatorname{Spec}_Y(q_{g,z})
      =r_z\circ\operatorname{Spec}_Y(q_{g,z})
      =i_g.
  \]

  \paragraph{Formalization.}
  The comparison $c_{g,z}\colon\operatorname{SIm}_Y(g)\to Z$ is
  \texttt{\detokenize{IsClosedImmersion.lift z (Scheme.Hom.imageι g) H}}, where $H$ is a proof
  of $\mathcal K_z\le\mathcal K_{i_g}$,
  and $e_{g,z}^{\mathrm{cmp}}\colon z\circ c_{g,z}=i_g$ is
  \texttt{\detokenize{IsClosedImmersion.lift_fac}}.  There is no declaration named
  \texttt{\detokenize{Scheme.Hom.schemeTheoreticImageComparison}}; the blueprint symbol is an
  alias for that application of \texttt{\detokenize{lift}}.

  The hypothesis $H$ is the displayed inclusion $\mathcal K_z\subseteq\mathcal K_g$, supplied
  by the project lemma
  \texttt{\detokenize{AlgebraicGeometry.Scheme.Hom.ker_le_ker_imageι_comp}}, which combines
  Mathlib's \texttt{\detokenize{Scheme.Hom.le_ker_comp}} (the kernel of $z$ is contained in
  the kernel of $z\circ h$) with
  \texttt{\detokenize{Scheme.IdealSheafData.ker_subschemeι}} (the kernel of $i_g$ is
  $\mathcal K_g$).

  The Lean route \textbf{avoids the chosen isomorphism $\sigma_z$ entirely}.  Where the proof
  above replaces $Z$ by $\operatorname{Spec}_Y(\mathcal O_Y/\mathcal K_z)$ using a choice from
  the nonempty set $\mathscr S_z$ and then descends $q_{g,z}$,
  \texttt{\detokenize{lift}} takes the ideal-sheaf inequality directly and produces the
  morphism into $Z$ itself.  The class equality (SIC0) is then not needed as an input; the
  statement corresponding to it is
  \texttt{\detokenize{IsClosedImmersion.isIso_lift}}, which says that the lift is an
  isomorphism when the two ideal sheaves are equal.
\end{definition}

\begin{lemma}[Universal property of the scheme-theoretic image]
  \label{lemma:nagata-scheme-image-universal}
  \uses{def:nagata-scheme-image-factor,def:nagata-scheme-image-comparison}
  \lean{AlgebraicGeometry.Scheme.Hom.toImage_imageι,
    AlgebraicGeometry.IsClosedImmersion.lift_fac,
    AlgebraicGeometry.Scheme.Hom.toImage_lift,
    AlgebraicGeometry.IsClosedImmersion.isClosedImmersion_lift,
    AlgebraicGeometry.IsClosedImmersion.eq_lift}
  \leanok

  With the notation of
  \Cref{def:nagata-scheme-image-factor,def:nagata-scheme-image-comparison},
  \[
    i_g\circ g^{\mathrm{im}}=g,
    \qquad z\circ c_{g,z}=i_g,
    \qquad c_{g,z}\circ g^{\mathrm{im}}=h.
  \]
  The morphism $c_{g,z}$ is a closed immersion, and
  \[
    \{c\colon\operatorname{SIm}_Y(g)\to Z\mid z\circ c=i_g\}
      =\{c_{g,z}\}.
  \]

  \paragraph{Formalization.}
  All five assertions are formalized, though not as a single declaration --- the aggregate is
  blueprint-local.  Write $c_{g,z}$ for the lift of
  \Cref{def:nagata-scheme-image-comparison}.

  The equality $i_g\circ g^{\mathrm{im}}=g$ is
  \texttt{\detokenize{Scheme.Hom.toImage_imageι}} and $z\circ c_{g,z}=i_g$ is
  \texttt{\detokenize{IsClosedImmersion.lift_fac}}; both are Mathlib.  The third equality
  $c_{g,z}\circ g^{\mathrm{im}}=h$ is the project lemma
  \texttt{\detokenize{Scheme.Hom.toImage_lift}}, proved exactly as above by cancelling the
  monomorphism $z$.

  That $c_{g,z}$ is a closed immersion is the project instance
  \texttt{\detokenize{IsClosedImmersion.isClosedImmersion_lift}}.  Its Lean proof is not the
  affine argument given above: since $c_{g,z}$ composed with $z$ is the closed immersion
  $i_g$, and $z$ is a closed immersion,
  \texttt{\detokenize{IsClosedImmersion.of_comp_isClosedImmersion}} applies directly, so no
  affine chart and no surjectivity criterion is needed.

  The uniqueness clause is the project lemma
  \texttt{\detokenize{IsClosedImmersion.eq_lift}}, stated in the equivalent pointwise form
  ``any $c$ with $z\circ c=i_g$ equals the lift'', proved by
  \texttt{\detokenize{cancel_mono}}.

  The three project declarations are in
  \texttt{\detokenize{WeilConjectures/Mathlib/AlgebraicGeometry/Morphisms/ClosedImmersion.lean}}
  and are upstreamable as written: they complete the API of Mathlib's
  \texttt{\detokenize{IsClosedImmersion.lift}}, which previously supplied only
  \texttt{\detokenize{lift_fac}} and \texttt{\detokenize{isIso_lift}}.
\end{lemma}

\begin{proof}
  By the equality witnesses $e_g^{\mathrm{im}}$ and
  $e_{g,z}^{\mathrm{cmp}}$ of
  \Cref{def:nagata-scheme-image-factor,def:nagata-scheme-image-comparison},
  \[
    i_g\circ g^{\mathrm{im}}=g,
    \qquad
    z\circ c_{g,z}=i_g.
  \]
  By \Cref{def:nagata-scheme-image-comparison},
  \[
    \mathcal K_z\subseteq\mathcal K_g,
    \qquad
    q_{g,z}\colon\mathcal O_Y/\mathcal K_z
      \twoheadrightarrow\mathcal O_Y/\mathcal K_g.
  \]
  On an affine open $V=\operatorname{Spec}(A)\subseteq Y$, put
  \[
    I_z(V):=\mathcal K_z(V)\subseteq A,
    \qquad
    I_g(V):=\mathcal K_g(V)\subseteq A.
  \]
  The restriction of $\operatorname{Spec}_Y(q_{g,z})$ to $V$ is
  \[
    \operatorname{Spec}(A/I_g(V))\longrightarrow
      \operatorname{Spec}(A/I_z(V)),
    \qquad
    A/I_z(V)\twoheadrightarrow A/I_g(V).
  \]
  The affine closed-immersion criterion applied to the displayed
  surjection asserts that $\operatorname{Spec}_Y(q_{g,z})$ is a closed
  immersion.  Since
  $\sigma_z^{-1}$ is an isomorphism, the identity
  \[
    c_{g,z}=\sigma_z^{-1}\circ\operatorname{Spec}_Y(q_{g,z})
  \]
  shows that $c_{g,z}$ is a closed immersion.

  The equality $c_{g,z}\circ g^{\mathrm{im}}=h$ is verified by the chain
  \[
    z\circ c_{g,z}\circ g^{\mathrm{im}}
      =i_g\circ g^{\mathrm{im}}
      =g
      =z\circ h.
  \]
  Since a closed immersion is a monomorphism, cancellation of $z$ in this
  equality is
  \[
    z\circ(c_{g,z}\circ g^{\mathrm{im}})=z\circ h
      \quad\Longrightarrow\quad
    c_{g,z}\circ g^{\mathrm{im}}=h.
  \]
  If $c\colon\operatorname{SIm}_Y(g)\to Z$ satisfies $z\circ c=i_g$, then
  \[
    z\circ c=i_g=z\circ c_{g,z}
      \quad\Longrightarrow\quad c=c_{g,z},
  \]
  by the monomorphism property of $z$.  The asserted uniqueness is the set
  equality
  \[
    \{c\colon\operatorname{SIm}_Y(g)\to Z\mid z\circ c=i_g\}
      =\{c_{g,z}\}.
  \]
\end{proof}

\begin{lemma}[Cartesian square defined by the scheme-image factor]
  \label{lemma:nagata-scheme-image-factor-cartesian}
  \uses{def:nagata-scheme-image-factor}
  \lean{}

  Let $g\colon T\to Y$ be as in
  \Cref{def:nagata-scheme-theoretic-image}.  Define
  \[
    P_g:=T\times_{Y,g,i_g}\operatorname{SIm}_Y(g)
  \]
  with projections
  $r_T\colon P_g\to T$ and
  $r_{\mathrm{im}}\colon P_g\to\operatorname{SIm}_Y(g)$.  Define
  \[
    \theta_g:=(\operatorname{id}_T,g^{\mathrm{im}})
      \colon T\longrightarrow P_g.
  \]
  Then $\theta_g$ is an isomorphism with inverse $r_T$, and
  \[
    r_T\circ\theta_g=\operatorname{id}_T,
    \qquad
    r_{\mathrm{im}}\circ\theta_g=g^{\mathrm{im}}.
  \]
  Equivalently, the square
  \[
  \begin{tikzcd}
    T \arrow[r,"\operatorname{id}_T"] \arrow[d,"g^{\mathrm{im}}"'] &
    T \arrow[d,"g"]\\
    \operatorname{SIm}_Y(g) \arrow[r,"i_g"'] & Y
  \end{tikzcd}
  \]
  is cartesian, and $g^{\mathrm{im}}$ is the base change of $g$ along
  $i_g$.

  \paragraph{Mathlib relation.}
  This square follows from pullback extensionality and
  \texttt{\detokenize{Scheme.Hom.toImage_imageι}}.  It is not the different
  Mathlib theorem
  \texttt{\detokenize{IsImmersion.isPullback_toImage_liftCoborder}}, whose
  lower-left object is a coborder range; the proposed declaration name here
  is therefore not marked exact.
\end{lemma}

\begin{proof}
  The equality $e_g^{\mathrm{im}}$ of
  \Cref{def:nagata-scheme-image-factor} is
  \[
    g\circ\operatorname{id}_T
      =g
      =i_g\circ g^{\mathrm{im}},
  \]
  so the fibre-product universal property defines $\theta_g$ with
  \[
    r_T\circ\theta_g=\operatorname{id}_T,
    \qquad
    r_{\mathrm{im}}\circ\theta_g=g^{\mathrm{im}}.
  \]
  The defining equality of $P_g$ and $e_g^{\mathrm{im}}$ are the chain
  \[
    i_g\circ r_{\mathrm{im}}
      =g\circ r_T
      =(i_g\circ g^{\mathrm{im}})\circ r_T
      =i_g\circ(g^{\mathrm{im}}\circ r_T).
  \]
  Since the closed immersion $i_g$ is a monomorphism,
  \[
    r_{\mathrm{im}}=g^{\mathrm{im}}\circ r_T.
  \]
  Substitution of this equality into both projections is
  \[
    r_T\circ(\theta_g\circ r_T)
      =(r_T\circ\theta_g)\circ r_T=r_T,
    \qquad
    r_{\mathrm{im}}\circ(\theta_g\circ r_T)
      =(r_{\mathrm{im}}\circ\theta_g)\circ r_T
      =g^{\mathrm{im}}\circ r_T=r_{\mathrm{im}}.
  \]
  The uniqueness clause in the universal property of $P_g$ is the equality
  \[
    \theta_g\circ r_T=\operatorname{id}_{P_g},
    \qquad
    r_T\circ\theta_g=\operatorname{id}_T.
  \]
  These equalities define the asserted inverse and cartesian identification.
\end{proof}

\begin{lemma}[Restriction of a scheme-theoretic image to an open target]
  \label{lemma:nagata-scheme-image-restrict-open}
  \uses{def:nagata-closed-subscheme-class,def:nagata-scheme-image-factor,
    def:nagata-presented-relative-spec,
    def:nagata-scheme-theoretic-image,
    lemma:nagata-morphism-property-calculus}
  \lean{}

  Let \(g\colon T\to Y\) satisfy \(\operatorname{IsQCQS}(g)\), as in
  \Cref{def:nagata-scheme-theoretic-image}, and let
  \(w\colon W\to Y\) satisfy \(\operatorname{IsOpenImmersion}(w)\).  Define
  \[
    T_W:=T\times_{Y,g,w}W,
    \qquad
    \pi_T\colon T_W\to T,
    \qquad
    g_W\colon T_W\to W,
    \qquad
    g\circ\pi_T=w\circ g_W.
  \]
  The morphism $g_W$ is the base change of $g$ along $w$; hence (PC2) of
  \Cref{lemma:nagata-morphism-property-calculus} is the predicate chain
  \[
    \operatorname{IsQCQS}(g)
      \Longrightarrow\operatorname{IsQCQS}(g_W).                 \tag{SIRO0}
  \]
  Define
  \[
    P_{g,w}:=\operatorname{SIm}_Y(g)\times_{Y,i_g,w}W
  \]
  with projections
  $\rho_g\colon P_{g,w}\to\operatorname{SIm}_Y(g)$ and
  $\rho_W\colon P_{g,w}\to W$.  By
  \Cref{def:nagata-scheme-image-factor}, define
  \[
    g_W^{\mathrm{im}}\colon T_W\to\operatorname{SIm}_W(g_W),
    \qquad
    i_{g_W}\circ g_W^{\mathrm{im}}=g_W.                           \tag{SIRO1}
  \]
  There is a unique isomorphism
  \[
    \delta_{g,w}\colon\operatorname{SIm}_W(g_W)
      \xrightarrow{\ \sim\ }P_{g,w}
  \]
  satisfying, for
  $\widetilde w_g:=\rho_g\circ\delta_{g,w}$,
  \[
    \rho_W\circ\delta_{g,w}=i_{g_W},
    \qquad
    i_g\circ\widetilde w_g=w\circ i_{g_W},
    \qquad
    \widetilde w_g\circ g_W^{\mathrm{im}}
      =g^{\mathrm{im}}\circ\pi_T.
  \]
  Equivalently,
  \[
    [\,\operatorname{SIm}_W(g_W)
        \xhookrightarrow{i_{g_W}}W\,]
      =[\,P_{g,w}\xhookrightarrow{\rho_W}W\,]
      \quad\text{in }\operatorname{ClosedSub}(W).
  \]

  \paragraph{Mathlib relation.}
  The proof uses exact ideal-data comap, subscheme base-change, and scheme
  gluing interfaces, but no single current Mathlib declaration with the
  proposed \texttt{\detokenize{schemeTheoreticImage_restrict}} name was
  found.  This remains a project lemma over those APIs.
\end{lemma}

\begin{proof}
  Let
  \[
    \omega_w\colon w^{-1}\mathcal O_Y\xrightarrow{\ \sim\ }\mathcal O_W
  \]
  be the structure isomorphism of the open immersion $w$.  Define
  \[
    \beta_{g,w}\colon w^{-1}(g_*\mathcal O_T)
      \xrightarrow{\ \sim\ }(g_W)_*\mathcal O_{T_W}
  \]
  on the stalk at $x\in W$ by
  \[
    (\beta_{g,w})_x[U,s]
      :=\left[|w|^{-1}(U),
        s|_{|\pi_T|\bigl(|g_W|^{-1}(|w|^{-1}(U))\bigr)}\right],
    \quad
    U\subseteq |Y|\text{ open},\quad |w|(x)\in U,
    \quad s\in\mathcal O_T(|g|^{-1}(U)).
  \]
  For each open $V\subseteq|W|$, define the open subschemes
  \[
  \begin{gathered}
    T_{W,V}:=g_W^{-1}(V)\subseteq T_W,\qquad
    |T_{W,V}|=|g_W|^{-1}(V),\\
    T_{Y,V}:=g^{-1}(|w|(V))\subseteq T,\qquad
    |T_{Y,V}|=|g|^{-1}(|w|(V)),\\
    \pi_{T,V}:=\pi_T|_{T_{W,V}}\colon
      T_{W,V}\xrightarrow{\ \sim\ }T_{Y,V}.
  \end{gathered}
  \]
  Denote the induced ring isomorphism by
  \[
    \gamma_V:=
      \varepsilon_{\pi_{T,V},\,|T_{Y,V}|}\colon
      \mathcal O_{T_{Y,V}}(|T_{Y,V}|)\xrightarrow{\ \sim\ }
      \mathcal O_{T_{W,V}}(|T_{W,V}|),
  \]
  where $\varepsilon_{\pi_{T,V},\,|T_{Y,V}|}$ is the component of
  the direct-image structure map of
  \Cref{def:nagata-direct-image-structure-map}.  Define the inverse stalk
  map by
  \[
    (\beta_{g,w}^{-1})_x[V,t]
      :=[|w|(V),\gamma_V^{-1}(t)].
  \]
  The two component calculations are
  \[
    \beta_{g,w}^{-1}\circ\beta_{g,w}
      =\operatorname{id}_{w^{-1}g_*\mathcal O_T},
    \qquad
    \beta_{g,w}\circ\beta_{g,w}^{-1}
      =\operatorname{id}_{(g_W)_*\mathcal O_{T_W}}.
  \]
  Substitution of $g\circ\pi_T=w\circ g_W$ into
  \Cref{def:nagata-direct-image-structure-map} is the commutative square
  \[
  \begin{tikzcd}
    w^{-1}\mathcal O_Y
      \arrow[r,"w^{-1}\varepsilon_g"] \arrow[d,"\omega_w"',"\sim"] &
    w^{-1}g_*\mathcal O_T \arrow[d,"\beta_{g,w}","\sim"']\\
    \mathcal O_W \arrow[r,"\varepsilon_{g_W}"'] &
    (g_W)_*\mathcal O_{T_W}.
  \end{tikzcd}
  \]
  Applying $\ker(-)$ to this square and conjugating by the isomorphisms
  $\omega_w$ and $\beta_{g,w}$ is the ideal equality
  \[
    \mathcal K_{g_W}
      =\omega_w(w^{-1}\mathcal K_g)
      =w^{-1}\mathcal K_g\cdot\mathcal O_W.
  \]
  Define the pullback $\mathcal O_W$-algebra and the quotient-algebra map
  \[
  \begin{aligned}
    w^*(\mathcal O_Y/\mathcal K_g)
      &:={w^{-1}(\mathcal O_Y/\mathcal K_g)}
          \otimes_{w^{-1}\mathcal O_Y}\mathcal O_W,\\
    \overline\omega_{g,w}\colon
      w^*(\mathcal O_Y/\mathcal K_g)
      &\longrightarrow\mathcal O_W/\mathcal K_{g_W},\\
    \overline\omega_{g,w}((\overline a)\otimes b)
      &:=\overline{\omega_w(a)b}.
  \end{aligned}
  \]
  Here $a$ is a local section of $w^{-1}\mathcal O_Y$, $b$ is a local
  section of $\mathcal O_W$, and bars denote the indicated quotient classes.
  If $a-a'\in w^{-1}\mathcal K_g$ and
  $d\in w^{-1}\mathcal O_Y$, then
  \[
  \begin{aligned}
    \omega_w(a-a')b&\in\mathcal K_{g_W},\\
    \overline\omega_{g,w}((\overline a\,d)\otimes b)
      &=\overline{\omega_w(a)\omega_w(d)b}
       =\overline\omega_{g,w}
          (\overline a\otimes\omega_w(d)b).
  \end{aligned}
  \]
  The first formula is independence of the representative $a$; the second is
  the tensor-balancing relation.  Define
  \[
    \overline\omega_{g,w}^{-1}(\overline c)
      :=\overline{\omega_w^{-1}(c)}\otimes1.
  \]
  If $c-c'\in\mathcal K_{g_W}$, then
  \[
    \omega_w^{-1}(c-c')\in w^{-1}\mathcal K_g,
    \qquad
    \overline{\omega_w^{-1}(c)}
      =\overline{\omega_w^{-1}(c')},
  \]
  which is independence of the representative $c$.
  The two formulas are the identities
  \[
  \begin{aligned}
    (\overline\omega_{g,w}^{-1}\circ
       \overline\omega_{g,w})((\overline a)\otimes b)
      &=\overline{\omega_w^{-1}(\omega_w(a)b)}\otimes1
       =(\overline a)\otimes b,\\
    (\overline\omega_{g,w}\circ
       \overline\omega_{g,w}^{-1})(\overline c)
      &=\overline{\omega_w(\omega_w^{-1}(c))}
       =\overline c.
  \end{aligned}
  \]
  Thus $\overline\omega_{g,w}$ is an explicitly specified
  $\mathcal O_W$-algebra isomorphism.  The pair
  $(\mathcal K_{g_W},\overline\omega_{g,w}^{-1})$ is therefore the quotient
  presentation of $w^*(\mathcal O_Y/\mathcal K_g)$ used below, in the sense
  of \Cref{def:nagata-presented-relative-spec}.  Put
  \[
    Z_{g,w}:=
      \operatorname{Spec}_W\bigl(w^*(\mathcal O_Y/\mathcal K_g)\bigr),
  \]
  and let $z_W\colon Z_{g,w}\to W$ be its structure morphism.  Define the
  adjunction map and its associated relative-spectrum morphism by
  \[
  \begin{aligned}
    \eta_{g,w}\colon\mathcal O_Y/\mathcal K_g
      &\longrightarrow w_*w^*(\mathcal O_Y/\mathcal K_g),&
    \eta_{g,w}(\overline a)&:=\overline a\otimes1,\\
    z_g&:=\operatorname{Spec}(\eta_{g,w})\colon
      Z_{g,w}\longrightarrow\operatorname{SIm}_Y(g).
  \end{aligned}
  \]
  Here $z_g$ is precisely the cross-base morphism (MLPSP1a), with
  $f=w$ and $\mathcal A=\mathcal O_Y/\mathcal K_g$; no arbitrary
  relative-spectrum adjunction is being used.
  Define
  \[
    \operatorname{bc}_{g,w}\colon P_{g,w}
      \xrightarrow{\ \sim\ }
      Z_{g,w}
  \]
  to be the relative-spectrum base-change isomorphism characterized by
  \[
    z_W\circ\operatorname{bc}_{g,w}=\rho_W,
    \qquad
    z_g\circ\operatorname{bc}_{g,w}=\rho_g.
  \]
  By
  \Cref{def:nagata-scheme-theoretic-image}, define
  \[
  \begin{aligned}
    \delta_{g,w}
      &:=(\operatorname{bc}_{g,w})^{-1}\circ
        \operatorname{Spec}_W(\overline\omega_{g,w})\\
      &\colon
        \operatorname{SIm}_W(g_W)
        =\operatorname{Spec}_W(\mathcal O_W/\mathcal K_{g_W})
        \xrightarrow{\ \sim\ }P_{g,w}.
  \end{aligned}
  \]
  Functoriality of relative spectrum over $W$ is the structural equality
  \[
    \rho_W\circ\delta_{g,w}=i_{g_W}.
  \]
  The fibre-product equality for $P_{g,w}$ is the chain
  \[
    i_g\circ\widetilde w_g
      =i_g\circ\rho_g\circ\delta_{g,w}
      =w\circ\rho_W\circ\delta_{g,w}
      =w\circ i_{g_W}.
  \]
  Finally,
  \[
    i_g\circ\widetilde w_g\circ g_W^{\mathrm{im}}
      =w\circ i_{g_W}\circ g_W^{\mathrm{im}}
      =w\circ g_W
      =g\circ\pi_T
      =i_g\circ g^{\mathrm{im}}\circ\pi_T.
  \]
  Since $i_g$ is a monomorphism, cancellation in the preceding equality is
  \[
    \widetilde w_g\circ g_W^{\mathrm{im}}
      =g^{\mathrm{im}}\circ\pi_T.
  \]
  If $\delta'\colon\operatorname{SIm}_W(g_W)\to P_{g,w}$ is a second
  morphism with $\rho_W\circ\delta'=i_{g_W}$, then
  \[
    \rho_W\circ\delta'
      =i_{g_W}
      =\rho_W\circ\delta_{g,w}.
  \]
  The morphism $\rho_W$ is the base change of the closed immersion $i_g$,
  hence is a monomorphism; cancellation is the equality
  \[
    \delta'=\delta_{g,w}.
  \]
\end{proof}

\begin{lemma}[{\cite[\href{https://stacks.math.columbia.edu/tag/01RG}{Tag 01RG}]{stacks-project}}]
  \label{lemma:nagata-immersion-open-in-scheme-image}
  \uses{def:nagata-choice-notation,def:nagata-closed-subscheme-class,
    def:nagata-scheme-image-factor,
    lemma:nagata-immersion-qcqs-criteria,
    lemma:nagata-morphism-property-calculus}
  \lean{AlgebraicGeometry.IsImmersion.instIsOpenImmersionToImageOfQuasiCompact,
    AlgebraicGeometry.Scheme.instIsDominantToImageOfQuasiCompact}

  Let $g\colon T\to Y$ satisfy
  \[
    \operatorname{IsImmersion}(g),
    \qquad
    \operatorname{IsQuasiCompact}(g).
  \]
  By \Cref{lemma:nagata-immersion-qcqs-criteria},
  \[
    \operatorname{IsImmersion}(g)\wedge
      \operatorname{IsQuasiCompact}(g)
        \Longrightarrow\operatorname{IsQCQS}(g).                 \tag{IO-1}
  \]
  The morphism
  $g^{\mathrm{im}}\colon T\to\operatorname{SIm}_Y(g)$ of
  \Cref{def:nagata-scheme-image-factor} is an open immersion,
  and its image is topologically and scheme-theoretically dense.

  \paragraph{Formalization.}
  Two of the three conclusions are exact Mathlib instances, and the Lean proofs need neither
  the Stacks factorization $g=u\circ k$ nor the pullback $T_U$ constructed below.

  That $g^{\mathrm{im}}$ is an open immersion is
  \texttt{\detokenize{AlgebraicGeometry.IsImmersion.instIsOpenImmersionToImageOfQuasiCompact}}
  in \texttt{\detokenize{Mathlib/AlgebraicGeometry/Morphisms/Immersion.lean}}, which fires by
  instance search from the two hypotheses $\operatorname{IsImmersion}(g)$ and
  $\operatorname{IsQuasiCompact}(g)$.  That its image is topologically dense is
  \texttt{\detokenize{Scheme.instIsDominantToImageOfQuasiCompact}}, which gives
  $\operatorname{IsDominant}(g^{\mathrm{im}})$ from quasi-compactness of $g$ alone.  A third
  instance, \texttt{\detokenize{IsImmersion.instToImage}}, records the weaker statement that
  $g^{\mathrm{im}}$ is an immersion without any quasi-compactness hypothesis.

  \textbf{The scheme-theoretic density clause is not formalized, which is why this lemma is
  not marked as formalized.}  It asks for
  \texttt{\detokenize{IsSchemeTheoreticallyDominant g.toImage}}, equivalently
  $\ker(\varepsilon_{g^{\mathrm{im}}})=0$, i.e.\ vanishing of
  \texttt{\detokenize{Scheme.Hom.ker g.toImage}}.  Neither statement is in Mathlib and
  neither is found by instance search or by \texttt{\detokenize{exact?}}; there is no
  \texttt{\detokenize{Scheme.Hom.ker_toImage}}.  This single missing lemma is also what
  blocks the dense-refinement construction elsewhere in this chapter, so it should be proved
  once and cited from both places.
\end{lemma}

\begin{proof}
  By \cite[Tag 01RG]{stacks-project}, choose a scheme $U$, an open
  immersion $u\colon U\hookrightarrow Y$, a closed immersion
  $k\colon T\hookrightarrow U$, and an equality witness
  \[
    e_{u,k}\colon g=u\circ k.
  \]
  Put
  \[
    \overline T:=\operatorname{SIm}_Y(g),
    \qquad
    T_U:=T\times_{Y,g,u}U
  \]
  and denote the projections by
  $\pi_T\colon T_U\to T$ and $g_U\colon T_U\to U$.  Define
  \[
    \tau_{u,k}:=(\operatorname{id}_T,k)\colon T\longrightarrow T_U.
  \]
  Since $k$ is a closed immersion and $g_U$ is the base change of $g$ along
  $u$, \Cref{lemma:nagata-immersion-qcqs-criteria} and (PC2) of
  \Cref{lemma:nagata-morphism-property-calculus} are the predicate chain
  \[
  \begin{aligned}
    \operatorname{IsClosedImmersion}(k)
      &\Longrightarrow\operatorname{IsQCQS}(k),\\
    \operatorname{IsQCQS}(g)
      &\Longrightarrow\operatorname{IsQCQS}(g_U).
  \end{aligned}                                                     \tag{IO-2}
  \]
  By \Cref{def:nagata-scheme-image-factor}, define
  \[
    g_U^{\mathrm{im}}\colon T_U\to\operatorname{SIm}_U(g_U),
    \qquad i_{g_U}\circ g_U^{\mathrm{im}}=g_U.                    \tag{IO-3}
  \]
  Its projection identities are
  \[
    \pi_T\circ\tau_{u,k}=\operatorname{id}_T,
    \qquad
    g_U\circ\tau_{u,k}=k.
  \]
  Substitution of $e_{u,k}\colon g=u\circ k$ into the fibre-product identity
  is the chain
  \[
    u\circ g_U
      =g\circ\pi_T
      =u\circ k\circ\pi_T.
  \]
  Since $u$ is a monomorphism, cancellation is the equality
  \[
    g_U=k\circ\pi_T.
  \]
  The two projection identities are
  \[
    \pi_T\circ(\tau_{u,k}\circ\pi_T)=\pi_T,
    \qquad
    g_U\circ(\tau_{u,k}\circ\pi_T)
      =k\circ\pi_T=g_U,
  \]
  and the uniqueness clause for $T_U=T\times_YU$ is
  \[
    \tau_{u,k}\circ\pi_T=\operatorname{id}_{T_U}.
  \]
  The two equations
  \[
    \pi_T\circ\tau_{u,k}=\operatorname{id}_T,
    \qquad
    \tau_{u,k}\circ\pi_T=\operatorname{id}_{T_U}
  \]
  state that $\tau_{u,k}$ is an isomorphism with inverse $\pi_T$.

  From $k=g_U\circ\tau_{u,k}$ and
  \Cref{def:nagata-direct-image-structure-map},
  \[
    \varepsilon_k
      =\bigl((g_U)_*\varepsilon_{\tau_{u,k}}\bigr)
        \circ\varepsilon_{g_U},
    \qquad
    (g_U)_*\varepsilon_{\tau_{u,k}}\text{ is an isomorphism},
  \]
  Because postcomposition with an isomorphism does not change a kernel, the
  preceding factorization is the equality
  \[
    \mathcal K_k=\ker(\varepsilon_k)
      =\ker(\varepsilon_{g_U})=\mathcal K_{g_U}.
  \]
  The ideal equality is the class equality
  \[
    [\,\operatorname{SIm}_U(k)\xhookrightarrow{i_k}U\,]
      =[\,\operatorname{SIm}_U(g_U)\xhookrightarrow{i_{g_U}}U\,]
      \quad\text{in }\operatorname{ClosedSub}(U).                \tag{IO0}
  \]
  Define
  \[
    \mathscr L_{u,k}:=
      \left\{\lambda\in
        \operatorname{Hom}(\operatorname{SIm}_U(k),
          \operatorname{SIm}_U(g_U))
        \ \middle|\
        \operatorname{IsIso}(\lambda)\ \wedge\
        i_{g_U}\circ\lambda=i_k\right\}.
  \]
  By (CSub3) of \Cref{def:nagata-closed-subscheme-class}, (IO0) is
  $\mathscr L_{u,k}\ne\varnothing$.  Define the chosen witness by
  \[
    \lambda_{u,k}:=\operatorname{ch}(\mathscr L_{u,k})
      \in\mathscr L_{u,k},
    \qquad
    \lambda_{u,k}\colon\operatorname{SIm}_U(k)
      \xrightarrow{\ \sim\ }\operatorname{SIm}_U(g_U),
    \qquad
    i_{g_U}\circ\lambda_{u,k}=i_k.
  \]
  Since $k$ is a closed immersion, the isomorphism $\sigma_k$ of
  \Cref{def:nagata-scheme-image-comparison}, applied to
  $g=k$, $h=\operatorname{id}_T$, and $z=k$, satisfies
  \[
    \sigma_k\colon T\xrightarrow{\ \sim\ }\operatorname{SIm}_U(k),
    \qquad
    i_k\circ\sigma_k=k=i_k\circ k^{\mathrm{im}}.
  \]
  Since $i_k$ is a monomorphism, cancellation is
  \[
    \sigma_k=k^{\mathrm{im}}.
  \]
  Since $\sigma_k$ is an isomorphism, this equality states that
  $k^{\mathrm{im}}$ is an isomorphism.
  Cancellation of the monomorphism $i_{g_U}$ in
  \[
    i_{g_U}\circ\lambda_{u,k}\circ k^{\mathrm{im}}
      =i_k\circ k^{\mathrm{im}}
      =k
      =g_U\circ\tau_{u,k}
      =i_{g_U}\circ g_U^{\mathrm{im}}\circ\tau_{u,k}
  \]
  is the equality
  \[
    \lambda_{u,k}\circ k^{\mathrm{im}}
      =g_U^{\mathrm{im}}\circ\tau_{u,k}.
  \]

  Let $\delta_{g,u}$ and $\widetilde u_g$ be the maps of
  \Cref{lemma:nagata-scheme-image-restrict-open}, and define
  \[
    \Omega_{u,k}:=
      \delta_{g,u}\circ\lambda_{u,k}\circ k^{\mathrm{im}}
      \colon T\xrightarrow{\ \sim\ }\overline T\times_YU.
  \]
  Define the two projections by
  \[
    \rho_{\overline T}:=\rho_g\colon
      \overline T\times_YU\to\overline T,
    \qquad
    \rho_U:=\rho_W\colon\overline T\times_YU\to U.
  \]
  Then
  \[
    \rho_{\overline T}\circ\Omega_{u,k}
      =\widetilde u_g\circ g_U^{\mathrm{im}}\circ\tau_{u,k}
      =g^{\mathrm{im}}\circ\pi_T\circ\tau_{u,k}
      =g^{\mathrm{im}},
  \]
  and
  \[
    \rho_U\circ\Omega_{u,k}
      =i_{g_U}\circ\lambda_{u,k}\circ k^{\mathrm{im}}
      =k.
  \]
  The equality
  \[
    g^{\mathrm{im}}=\rho_{\overline T}\circ\Omega_{u,k}
  \]
  identifies $g^{\mathrm{im}}$, through the isomorphism $\Omega_{u,k}$,
  with the base change $\rho_{\overline T}$ of the open immersion $u$.
  Stability of open immersions under base change and composition with an
  isomorphism asserts
  \[
    \operatorname{IsOpenImmersion}(g^{\mathrm{im}}).
  \]
  By \Cref{lemma:nagata-scheme-image-factor-cartesian},
  $g^{\mathrm{im}}$ is the base change of $g$ along $i_g$; stability of
  quasi-compact morphisms under base change is
  \[
    \operatorname{IsQuasiCompact}(g^{\mathrm{im}}).
  \]

  By \Cref{def:nagata-scheme-image-factor},
  \[
    \overline\varepsilon_g
      =(i_g)_*\varepsilon_{g^{\mathrm{im}}}\circ\chi_g,
    \qquad
    \ker(\overline\varepsilon_g)=0.
  \]
  Since $\chi_g$ is an isomorphism and $(i_g)_*$ is faithful on sheaves,
  \[
    (i_g)_*\ker(\varepsilon_{g^{\mathrm{im}}})
      =\ker\bigl((i_g)_*\varepsilon_{g^{\mathrm{im}}}\bigr)
      =\ker(\overline\varepsilon_g\circ\chi_g^{-1})
      =0,
    \qquad
    \ker(\varepsilon_{g^{\mathrm{im}}})=0.
  \]
  By \Cref{def:nagata-direct-image-structure-map}, the last equality is
  \[
    \operatorname{SchematicallyDense}(g^{\mathrm{im}}).
  \]

  Put
  \[
    V:=\overline T\setminus
      \overline{g^{\mathrm{im}}(T)}
  \]
  with its open subscheme structure.  Then
  $(g^{\mathrm{im}})^{-1}(V)=\varnothing$, so the component over $V$ is
  \[
    (\varepsilon_{g^{\mathrm{im}}})_V\colon
      \mathcal O_{\overline T}(V)\longrightarrow
      \mathcal O_T(\varnothing)=0.
  \]
  Its kernel is zero, whence
  \[
    1_{\mathcal O_{\overline T}(V)}=0.
  \]
  For $x\in V$, the stalk homomorphism sends the preceding equality to
  \[
    1_{\mathcal O_{\overline T,x}}=0.
  \]
  A local ring attached to a point has a prime maximal ideal, so this
  equality forces $V=\varnothing$.  By the definition of $V$, this is
  \[
    \overline{g^{\mathrm{im}}(T)}=\overline T.
  \]
\end{proof}

\begin{definition}[Compactification constructed from an immersion]
  \label{def:nagata-immersion-compactification}
  \uses{def:nagata-compactification,
    def:nagata-scheme-image-factor,
    lemma:nagata-immersion-qcqs-criteria}
  \lean{}

  Let $g\colon X\to Y$ and $p_Y\colon Y\to S$ satisfy
  $f:=p_Y\circ g$,
  \[
    \operatorname{IsImmersion}(g),\qquad
    \operatorname{IsQuasiCompact}(g),\qquad
    \operatorname{IsProper}(p_Y).
  \]
  By \Cref{lemma:nagata-immersion-qcqs-criteria},
  \[
    \operatorname{IsImmersion}(g)\wedge
    \operatorname{IsQuasiCompact}(g)
      \Longrightarrow\operatorname{IsQCQS}(g).                  \tag{IC-1}
  \]
  With $i_g$ and $g^{\mathrm{im}}$ as in
  \Cref{def:nagata-scheme-theoretic-image,def:nagata-scheme-image-factor},
  define
  \[
    q_g:=p_Y\circ i_g\colon\operatorname{SIm}_Y(g)\to S
  \]
  and define the equality witness
  \[
    e_{\mathsf C(g)}\colon q_g\circ g^{\mathrm{im}}=f
  \]
  by the chain
  \[
    q_g\circ g^{\mathrm{im}}
      =(p_Y\circ i_g)\circ g^{\mathrm{im}}
      =p_Y\circ(i_g\circ g^{\mathrm{im}})
      =p_Y\circ g=f,
  \]
  where $i_g\circ g^{\mathrm{im}}=g$ is
  $e_g^{\mathrm{im}}$ of \Cref{def:nagata-scheme-image-factor}.  Define
  \[
    \mathsf C(g):=
      \bigl(\operatorname{SIm}_Y(g),g^{\mathrm{im}},q_g,
        e_{\mathsf C(g)}\bigr).
  \]
\end{definition}

\begin{lemma}[The immersion construction is a compactification]
  \label{lemma:nagata-immersion-into-proper}
  \uses{def:nagata-immersion-compactification}
  \lean{}

  The factorization $\mathsf C(g)$ of
  \Cref{def:nagata-immersion-compactification} is a compactification of
  $X\to S$.
\end{lemma}

\begin{proof}
  \uses{lemma:nagata-immersion-open-in-scheme-image,
    lemma:nagata-scheme-image-factor-cartesian}
  By \Cref{lemma:nagata-immersion-open-in-scheme-image},
  \[
    \operatorname{IsOpenImmersion}(g^{\mathrm{im}}).
  \]
  Let $P_g$, $r_X$, $r_{\mathrm{im}}$, and
  $\theta_g\colon X\xrightarrow{\sim}P_g$ be as in
  \Cref{lemma:nagata-scheme-image-factor-cartesian}.  Then
  \[
    g^{\mathrm{im}}=r_{\mathrm{im}}\circ\theta_g,
    \qquad
    P_g=X\times_Y\operatorname{SIm}_Y(g),
    \qquad
    \operatorname{IsQuasiCompact}(g)
      \Longrightarrow
    \operatorname{IsQuasiCompact}(r_{\mathrm{im}})
      \Longrightarrow
    \operatorname{IsQuasiCompact}(g^{\mathrm{im}}),
  \]
  where the first implication is stability of quasi-compact morphisms under
  base change and the second is stability under composition with an
  isomorphism.  By \Cref{def:nagata-scheme-theoretic-image},
  $i_g$ is a closed immersion; the standard implications and composition
  rule for proper morphisms are
  \[
    \operatorname{IsClosedImmersion}(i_g)
      \Longrightarrow\operatorname{IsProper}(i_g),
    \qquad
    \operatorname{IsProper}(p_Y)\wedge\operatorname{IsProper}(i_g)
      \Longrightarrow
    \operatorname{IsProper}(p_Y\circ i_g)
      =\operatorname{IsProper}(q_g).
  \]
  Finally, \Cref{def:nagata-immersion-compactification} defines
  \[
    e_{\mathsf C(g)}\colon
    q_g\circ g^{\mathrm{im}}=f.
  \]
  The preceding three property statements and the equality
  $q_g\circ g^{\mathrm{im}}=f$ satisfy
  \Cref{def:nagata-compactification}; explicitly,
  \[
    \mathsf C(g)\in\operatorname{Compactification}(f).
  \]
\end{proof}

\begin{definition}[Base-change construction for a compactification]
  \label{def:nagata-compactification-base-change}
  \uses{def:nagata-compactification}
  \lean{}

  Let $C=(\overline X,j,p,e)$ be a compactification of
  $f\colon X\to S$ as in \Cref{def:nagata-compactification}, and let
  $a\colon S'\to S$.  Define
  \[
    X':=X\times_{S,f,a}S',
    \qquad
    \overline X':=\overline X\times_{S,p,a}S'.
  \]
  Denote their projections by
  \[
    x_X\colon X'\to X,
    \quad x_{S'}\colon X'\to S',
    \quad
    \overline x_{\overline X}\colon\overline X'\to\overline X,
    \quad \overline x_{S'}\colon\overline X'\to S',
  \]
  with defining equalities
  \[
    f\circ x_X=a\circ x_{S'},
    \qquad
    p\circ\overline x_{\overline X}=a\circ\overline x_{S'}.
  \]
  Substitution of $e\colon p\circ j=f$ into the first projection is
  \[
    p\circ j\circ x_X
      =f\circ x_X
      =a\circ x_{S'},
  \]
  so the fibre-product universal property defines the unique morphism
  \[
    j':=(j\circ x_X,x_{S'})\colon X'\longrightarrow\overline X',
    \qquad
    \overline x_{\overline X}\circ j'=j\circ x_X,
    \qquad
    \overline x_{S'}\circ j'=x_{S'}.
  \]
  Define
  \[
    p':=\overline x_{S'}\colon\overline X'\to S',
    \qquad
    f':=x_{S'}\colon X'\to S'.
  \]
  Define the equality witness $e_{S'}$ by
  \[
    e_{S'}\colon p'\circ j'=f',
    \qquad
    p'\circ j'
      =\overline x_{S'}\circ j'
      =x_{S'}=f'.
  \]
  Finally, define
  \[
    C_{S'}:=(\overline X',j',p',e_{S'}).
  \]
\end{definition}

\begin{lemma}[Base change of a compactification]
  \label{lemma:nagata-compactification-base-change}
  \uses{def:nagata-compactification-base-change}
  \lean{}

  The factorization $C_{S'}$ of
  \Cref{def:nagata-compactification-base-change} is a compactification of
  $X\times_SS'\to S'$.
\end{lemma}

\begin{proof}
  By \Cref{def:nagata-compactification},
  \[
    \operatorname{IsQuasiCompactOpenImmersion}(j),
    \qquad
    \operatorname{IsProper}(p).
  \]
  By \Cref{def:nagata-compactification-base-change}, the two cartesian
  squares and their projection equalities are
  \[
  \begin{tikzcd}
    X' \arrow[r,"j'"] \arrow[d,"x_X"'] &
      \overline X' \arrow[d,"\overline x_{\overline X}"] &
    \overline X' \arrow[r,"p'"] \arrow[d,"\overline x_{\overline X}"'] &
      S' \arrow[d,"a"]\\
    X \arrow[r,"j"'] & \overline X &
    \overline X \arrow[r,"p"'] & S,
  \end{tikzcd}
  \]
  \[
    \overline x_{\overline X}\circ j'=j\circ x_X,
    \qquad
    \overline x_{S'}\circ j'=x_{S'},
    \qquad
    a\circ p'=p\circ\overline x_{\overline X}.
  \]
  Stability of quasi-compact open immersions and proper morphisms under base
  change is the pair of assertions
  \[
    \operatorname{IsQuasiCompactOpenImmersion}(j'),
    \qquad
    \operatorname{IsProper}(p').
  \]
  The specified equality witness is
  \[
    e_{S'}\colon p'\circ j'
      =\overline x_{S'}\circ j'
      =x_{S'}=f'.
  \]
  These two property assertions and $e_{S'}$ satisfy
  \Cref{def:nagata-compactification}; explicitly,
  \[
    C_{S'}\in\operatorname{Compactification}(f').
  \]
\end{proof}

\begin{definition}[Dense-refinement construction]
  \label{def:nagata-dense-refinement}
  \uses{def:nagata-compactification,
    def:nagata-scheme-image-factor,
    lemma:nagata-immersion-qcqs-criteria}
  \lean{}

  Let $C=(\overline X,j,p,e)$ be a compactification of $f\colon X\to S$.
  By \Cref{def:nagata-compactification,lemma:nagata-immersion-qcqs-criteria},
  \[
    \operatorname{IsQuasiCompactOpenImmersion}(j)
      \Longrightarrow\operatorname{IsQCQS}(j).                  \tag{DR-1}
  \]
  By
  \Cref{def:nagata-scheme-theoretic-image,def:nagata-scheme-image-factor},
  define
  \[
  \begin{gathered}
    i_j\colon\operatorname{SIm}_{\overline X}(j)\hookrightarrow\overline X,
    \qquad
    j^{\mathrm{im}}\colon X\to\operatorname{SIm}_{\overline X}(j),\\
    i_j\circ j^{\mathrm{im}}=j.
  \end{gathered}                                                     \tag{DR-2}
  \]
  Set
  \[
    p_{\mathrm{im}}:=p\circ i_j
      \colon\operatorname{SIm}_{\overline X}(j)\to S.
  \]
  Define the equality witness $e_{\mathrm{im}}$ by
  \[
    e_{\mathrm{im}}\colon p_{\mathrm{im}}\circ j^{\mathrm{im}}=f,
    \qquad
    p_{\mathrm{im}}\circ j^{\mathrm{im}}
      =(p\circ i_j)\circ j^{\mathrm{im}}
      =p\circ(i_j\circ j^{\mathrm{im}})
      =p\circ j=f,
  \]
  using $e_j^{\mathrm{im}}$ from
  \Cref{def:nagata-scheme-image-factor} and $e\colon p\circ j=f$ from
  \Cref{def:nagata-compactification}.  Define
  \[
    \mathsf{Dense}(C):=
      \left(\operatorname{SIm}_{\overline X}(j),
        j^{\mathrm{im}},p_{\mathrm{im}},e_{\mathrm{im}}\right).
  \]
\end{definition}

\begin{proposition}[Properties of the dense-refinement construction]
  \label{prop:nagata-dense-refinement}
  \uses{def:nagata-dense-compactification,def:nagata-dense-refinement}
  \lean{}

  The quadruple $\mathsf{Dense}(C)$ of
  \Cref{def:nagata-dense-refinement} is a scheme-theoretically dense
  compactification of the same morphism as $C$.
\end{proposition}

\begin{proof}
  \uses{lemma:nagata-immersion-open-in-scheme-image,
    lemma:nagata-scheme-image-factor-cartesian}
  By \Cref{lemma:nagata-immersion-open-in-scheme-image},
  \[
    \operatorname{IsOpenImmersion}(j^{\mathrm{im}}),
    \qquad
    \ker(\varepsilon_{j^{\mathrm{im}}})=0.
  \]
  Let
  \[
    P_j:=X\times_{\overline X,j,i_j}
      \operatorname{SIm}_{\overline X}(j)
  \]
  and let $r_{\mathrm{im}}\colon P_j\to
  \operatorname{SIm}_{\overline X}(j)$ and
  $\theta_j\colon X\xrightarrow{\sim}P_j$ be the maps of
  \Cref{lemma:nagata-scheme-image-factor-cartesian}.  Then
  \[
    j^{\mathrm{im}}=r_{\mathrm{im}}\circ\theta_j,
    \qquad
    \operatorname{IsQuasiCompact}(j)
      \Longrightarrow\operatorname{IsQuasiCompact}(r_{\mathrm{im}})
      \Longrightarrow\operatorname{IsQuasiCompact}(j^{\mathrm{im}}).
  \]
  The two implications are, respectively, base-change stability and
  composition stability for quasi-compact morphisms.  By
  \Cref{def:nagata-scheme-theoretic-image}, $i_j$ is a closed immersion, so
  \[
    \operatorname{IsClosedImmersion}(i_j)
      \Longrightarrow\operatorname{IsProper}(i_j),
    \qquad
    \operatorname{IsProper}(p)\wedge\operatorname{IsProper}(i_j)
      \Longrightarrow
    \operatorname{IsProper}(p\circ i_j)
      =\operatorname{IsProper}(p_{\mathrm{im}}).
  \]
  By \Cref{def:nagata-dense-refinement}, the equality witness is
  \[
    e_{\mathrm{im}}\colon
      p_{\mathrm{im}}\circ j^{\mathrm{im}}=f.
  \]
  Substitution of the displayed properties into
  \Cref{def:nagata-compactification,def:nagata-dense-compactification} is
  \[
    \mathsf{Dense}(C)\in\operatorname{Compactification}(f),
    \qquad
    \operatorname{Dense}(\mathsf{Dense}(C))
      \Longleftrightarrow
      \ker(\varepsilon_{j^{\mathrm{im}}})=0,
  \]
  by \Cref{def:nagata-compactification,def:nagata-dense-compactification}.
\end{proof}

\begin{definition}[Closed-subscheme compactification construction]
  \label{def:nagata-closed-subscheme-compactification}
  \uses{def:nagata-compactification,
    def:nagata-scheme-theoretic-image,
    def:nagata-scheme-image-factor,
    lemma:nagata-morphism-property-calculus}
  \lean{}

  Let $f\colon X\to S$ and $f'\colon X'\to S$ be morphisms.  Let
  $i\colon X\hookrightarrow X'$ be a closed immersion, let
  $e_i\colon f'\circ i=f$ be a specified equality, and let
  \[
    C'=(\overline X',j',p',e')
      \in\operatorname{Compactification}(f')
  \]
  be as in \Cref{def:nagata-compactification}, where the fourth component is
  the equality
  \[
    e'\colon p'\circ j'=f'.
  \]
  Define first
  \[
    g:=j'\circ i\colon X\to\overline X'.
  \]
  By \Cref{def:nagata-compactification,lemma:nagata-morphism-property-calculus},
  the required domain predicate
  for the scheme-theoretic image is the displayed chain
  \[
  \begin{aligned}
    \operatorname{IsClosedImmersion}(i)
      &\Longrightarrow
        \bigl(
          \operatorname{IsQuasiCompact}(i)\wedge
          \operatorname{QuasiSeparated}(i)
        \bigr)
       \Longrightarrow\operatorname{IsQCQS}(i),\\
    \operatorname{IsQuasiCompactOpenImmersion}(j')
      &\Longrightarrow
        \bigl(
          \operatorname{IsQuasiCompact}(j')\wedge
          \operatorname{QuasiSeparated}(j')
        \bigr)
       \Longrightarrow\operatorname{IsQCQS}(j'),\\
    \operatorname{IsQCQS}(i)\wedge\operatorname{IsQCQS}(j')
      &\Longrightarrow
        \operatorname{IsQCQS}(j'\circ i)
       =\operatorname{IsQCQS}(g).
  \end{aligned}
  \]
  Only after this predicate has been established, define
  \[
    \overline X:=\operatorname{SIm}_{\overline X'}(g),
    \qquad
    i_g\colon\overline X\hookrightarrow\overline X',
    \qquad
    g^{\mathrm{im}}\colon X\to\overline X,
  \]
  where \Cref{def:nagata-scheme-image-factor} gives
  \[
    i_g\circ g^{\mathrm{im}}=g.                                  \tag{CSC0}
  \]
  Define
  \[
    q:=p'\circ i_g\colon\overline X\to S.
  \]
  Define the equality witness $e_{i,C'}$ by
  \[
    e_{i,C'}\colon q\circ g^{\mathrm{im}}=f,
    \qquad
    q\circ g^{\mathrm{im}}
      =(p'\circ i_g)\circ g^{\mathrm{im}}
      =p'\circ(i_g\circ g^{\mathrm{im}})
      =p'\circ g
      =p'\circ j'\circ i
      =f'\circ i=f.
  \]
  Here the third equality is $e_g^{\mathrm{im}}$ from
  \Cref{def:nagata-scheme-image-factor}, the fifth is
  $e'\colon p'\circ j'=f'$, and the sixth is $e_i$.  Define
  \[
    \mathsf{Closed}(i,C'):=
      (\overline X,g^{\mathrm{im}},q,e_{i,C'}).
  \]
\end{definition}

\begin{proposition}[The closed-subscheme construction is a compactification]
  \label{prop:nagata-closed-subscheme-compactifiable}
  \uses{def:nagata-closed-subscheme-compactification}
  \lean{}

  The factorization $\mathsf{Closed}(i,C')$ of
  \Cref{def:nagata-closed-subscheme-compactification} is a compactification
  of $X\to S$.
\end{proposition}

\begin{proof}
  \uses{lemma:nagata-immersion-into-proper}
  By \Cref{def:nagata-closed-subscheme-compactification},
  \[
    g=j'\circ i,
    \qquad
    q=p'\circ i_g,
    \qquad
    e_{i,C'}\colon q\circ g^{\mathrm{im}}=f.
  \]
  The morphism $g$ is a quasi-compact immersion and $p'$ is proper, while
  the equality required by
  \Cref{def:nagata-immersion-compactification} is
  \[
    p'\circ g
      =p'\circ j'\circ i
      =f'\circ i=f.
  \]
  Substitution of $(g,p_Y)=(g,p')$ into
  \Cref{def:nagata-immersion-compactification} is the tuple equality
  \[
    \mathsf C(g)
      =\bigl(\operatorname{SIm}_{\overline X'}(g),
        g^{\mathrm{im}},p'\circ i_g,e_{\mathsf C(g)}\bigr)
      =\bigl(\overline X,g^{\mathrm{im}},q,e_{i,C'}\bigr)
      =\mathsf{Closed}(i,C'),
  \]
  because $e_{\mathsf C(g)}$ and $e_{i,C'}$ are the same displayed
  composition chain.  By
  \Cref{lemma:nagata-immersion-into-proper},
  \[
    \mathsf{Closed}(i,C')=\mathsf C(g)
      \in\operatorname{Compactification}(f).
  \]
\end{proof}

\section{Separated gluing}

\begin{lemma}[The diagonal on a mixed open]
  \label{lemma:nagata-diagonal-mixed-open}
  \lean{}

  Let $f\colon X\to S$, let $u\colon U\hookrightarrow X$ and
  $v\colon V\hookrightarrow X$ be open immersions with
  $u(U)\cup v(V)=|X|$, and put $W:=U\times_XV$ with projections
  $p_U\colon W\to U$ and $p_V\colon W\to V$.  Let
  \[
    \kappa_{UV}:=(p_U,p_V)\colon W\longrightarrow U\times_SV.
  \]
  If
  \[
    P_{UV}:=(U\times_SV)\times_{X\times_SX,\,(u,v)}X
  \]
  is the pullback of $\Delta_{X/S}\colon X\to X\times_SX$, then the morphism
  \[
    \phi_{UV}:=(\kappa_{UV},u\circ p_U)\colon W\longrightarrow P_{UV}
  \]
  is an isomorphism and identifies the base change of $\Delta_{X/S}$ with
  \[
    \kappa_{UV}\colon W\longrightarrow U\times_SV.
  \]
  Interchanging $U$ and $V$ defines
  $\phi_{VU}\colon V\times_XU\xrightarrow{\sim}P_{VU}$ with the same
  property over $V\times_SU$.
\end{lemma}

\begin{proof}
  Denote the projections $P_{UV}\to U\times_SV$ and $P_{UV}\to X$ by
  $r_{UV}$ and $r_X$.  The defining equality of $P_{UV}$ is
  \[
    (u,v)\circ r_{UV}=\Delta_{X/S}\circ r_X.
  \]
  Writing $r_{UV}=(r_U,r_V)$, this equality is equivalent to
  \[
    u\circ r_U=r_X=v\circ r_V.
  \]
  The equalities $u\circ r_U=r_X=v\circ r_V$ satisfy the universal property
  of $W=U\times_XV$ and define
  \[
    \psi_{UV}:=(r_U,r_V)\colon P_{UV}\longrightarrow W.
  \]
  The definitions of $\psi_{UV}$ and $\phi_{UV}$ are the projection
  equalities
  \[
    p_U\circ\psi_{UV}=r_U,
    \quad p_V\circ\psi_{UV}=r_V,
    \quad r_{UV}\circ\phi_{UV}=\kappa_{UV},
    \quad r_X\circ\phi_{UV}=u\circ p_U.
  \]
  The uniqueness clauses of the two fibre products are the identities
  \[
    \psi_{UV}\circ\phi_{UV}=\operatorname{id}_W,
    \qquad
    \phi_{UV}\circ\psi_{UV}=\operatorname{id}_{P_{UV}}.
  \]
  These two identities state that
  $\phi_{UV}^{-1}=\psi_{UV}$.  The equality
  $r_{UV}\circ\phi_{UV}=\kappa_{UV}$ is the claimed identification of the
  base-changed diagonal.
  The formulas with $U$ and $V$ interchanged prove the second assertion.
\end{proof}

\begin{lemma}[{\cite[\href{https://stacks.math.columbia.edu/tag/0F3U}{Tag 0F3U}]{stacks-project}}]
  \label{lemma:nagata-separated-two-open-cover}
  \lean{}

  Retain $f,u,v,W,p_U,p_V$ from
  \Cref{lemma:nagata-diagonal-mixed-open}.  Assume that
  $f\circ u\colon U\to S$ and $f\circ v\colon V\to S$ are separated and that
  \[
    \kappa_{UV}=(p_U,p_V)\colon W\longrightarrow U\times_SV
  \]
  is a closed immersion.  Then $X\to S$ is separated.
\end{lemma}

\begin{proof}
  \uses{lemma:nagata-diagonal-mixed-open}
  Put $Q:=X\times_SX$.  The four morphisms
  \[
    (u,u)\colon U\times_SU\to Q,\quad
    (u,v)\colon U\times_SV\to Q,\quad
    (v,u)\colon V\times_SU\to Q,\quad
    (v,v)\colon V\times_SV\to Q
  \]
  are open immersions whose images cover $Q$.  The four base changes of
  $\Delta_{X/S}$ along this cover are, by
  \Cref{lemma:nagata-diagonal-mixed-open},
  \[
    \Delta_{U/S},\quad \kappa_{UV},\quad
    \kappa_{VU},\quad \Delta_{V/S}.
  \]
  The first and fourth maps are closed immersions by the assumed
  separatedness, and the second is a closed immersion by hypothesis.  Put
  $W_{VU}:=V\times_XU$, and denote its projections by
  $q_V\colon W_{VU}\to V$ and $q_U\colon W_{VU}\to U$.  Define the
  factor-swap isomorphisms
  \[
    \tau_X:=(q_U,q_V)\colon W_{VU}\xrightarrow{\sim}W,
    \qquad
    \tau_S:=(\operatorname{pr}_U,\operatorname{pr}_V)\colon
      V\times_SU\xrightarrow{\sim}U\times_SV.
  \]
  Define their inverses explicitly by
  \[
    \tau_X^{-1}:=(p_V,p_U)\colon W\to W_{VU},
    \qquad
    \tau_S^{-1}:=(\operatorname{pr}_V,\operatorname{pr}_U)\colon
      U\times_SV\to V\times_SU.
  \]
  For
  \[
    \kappa_{VU}:=(q_V,q_U)\colon W_{VU}\longrightarrow V\times_SU
  \]
  the four projection equalities are
  \[
    \tau_S\circ\kappa_{VU}=\kappa_{UV}\circ\tau_X,
    \qquad
    \kappa_{VU}=\tau_S^{-1}\circ\kappa_{UV}\circ\tau_X.
  \]
  The second equality expresses $\kappa_{VU}$ as a composite of the closed
  immersion $\kappa_{UV}$ with isomorphisms.  Stability of closed immersions
  under composition with isomorphisms is
  \[
    \operatorname{IsClosedImmersion}(\kappa_{VU}).
  \]
  Target-locality of closed immersions on the displayed cover is
  \[
    \operatorname{IsClosedImmersion}(\Delta_{X/S}).
  \]
  The diagonal criterion for separatedness identifies this assertion with
  \[
    \operatorname{IsSeparated}(f).
  \]
\end{proof}

\section{Construction of blow-ups from Rees algebras}
\label{sec:nagata-blowup-foundations}

This section expands the blow-up package down to the algebraic and gluing
lemmas needed by Lean.  The later Nagata argument only uses the public
geometric interfaces, but those interfaces should not be marked complete until
the constructions in this section are available.

\subsection{The graded Rees algebra}

Until the localization datum (LOC0), fix the explicitly typed Rees datum
\[
  \mathsf D_{\mathrm{Rees}}:=(A,I),\qquad
  A\text{ a commutative ring},\qquad I\triangleleft A,\qquad
  t:=\operatorname{monomial}(1,1_A)\in A[X],\qquad A[t]:=A[X].      \tag{REES0}
\]
Every occurrence of $A$ or $I$ before (LOC0) denotes the corresponding
projection of $\mathsf D_{\mathrm{Rees}}$, and every occurrence of $t$ denotes
the element specified in (REES0).

\begin{definition}[Rees algebra]
  \label{def:nagata-rees-algebra}
  \lean{reesAlgebra}
  \leanok

  For the datum (REES0), define the Rees algebra by
  \[
    \mathcal R_A(I)=\bigoplus_{n\geq 0}I^nt^n\subseteq A[t].
  \]
  Its underlying ring is the subring of $A[t]$ consisting of polynomials whose
  coefficient of $t^n$ belongs to $I^n$.

  \paragraph{Mathlib relation.}
  This is exactly the lowercase declaration
  \texttt{\detokenize{reesAlgebra I}} in
  \texttt{\detokenize{Mathlib.RingTheory.ReesAlgebra}}.  Its defining
  membership theorem is \texttt{\detokenize{mem_reesAlgebra_iff}}.  The
  proposed uppercase name \texttt{\detokenize{ReesAlgebra}} is not the
  current Mathlib name.
\end{definition}

\begin{definition}[Homogeneous piece of the Rees algebra]
  \label{def:nagata-rees-piece}
  \uses{def:nagata-rees-algebra}
  \lean{AlgebraicGeometry.reesGrading}
  \leanok

  Let $\iota\colon\mathcal R_A(I)\hookrightarrow A[t]$ be the inclusion of
  \Cref{def:nagata-rees-algebra}.  For $n\in\mathbb N$, define
  \[
    \mathcal R_A(I)_n:=
    \{F\in\mathcal R_A(I)\mid
      \operatorname{coeff}_k(\iota(F))=0\text{ for every }k\ne n\}.
  \]
\end{definition}

\begin{definition}[Coefficients, homogeneous components, and support]
  \label{def:nagata-rees-components}
  \uses{def:nagata-rees-algebra,def:nagata-rees-piece}
  \lean{AlgebraicGeometry.reesMonomial, AlgebraicGeometry.reesX}
  \leanok

  Let $\iota\colon\mathcal R_A(I)\hookrightarrow A[t]$ be the inclusion of
  \Cref{def:nagata-rees-piece}.  For $F\in\mathcal R_A(I)$ and
  $n\in\mathbb N$, define
  \[
  \begin{aligned}
    F_n&:=\operatorname{coeff}_n(\iota(F))\in I^n,\\
    \operatorname{hcomp}_n(F)&:=F_nt^n\in\mathcal R_A(I)_n,\\
    \operatorname{supp}(F)
      &:=\{n\in\mathbb N\mid F_n\ne0\}
       =\{n\in\mathbb N\mid\operatorname{hcomp}_n(F)\ne0\}.
  \end{aligned}
  \]
  The first membership is the coefficient condition in
  \Cref{def:nagata-rees-algebra}; the second membership is the defining
  equality in \Cref{def:nagata-rees-piece}.

  \paragraph{Mathlib relation.}
  Coefficient extraction and polynomial support are
  \texttt{\detokenize{Polynomial.coeff}} and
  \texttt{\detokenize{Polynomial.support}} after coercion from the exact
  subalgebra \texttt{\detokenize{reesAlgebra I}}.  The homogeneous-component
  notation is a blueprint wrapper around those declarations.
\end{definition}

\begin{lemma}[Membership in one Rees piece]
  \label{lemma:nagata-rees-piece-membership}
  \uses{def:nagata-rees-piece}
  \lean{AlgebraicGeometry.mem_reesGrading_iff_exists}
  \leanok

  For $n\in\mathbb N$, an element $F\in\mathcal R_A(I)$ lies in
  $\mathcal R_A(I)_n$ if and only if
  there is $a\in I^n$ such that the underlying polynomial of $F$ is
  $at^n$.

  \paragraph{Mathlib relation.}
  Mathlib does not package the blueprint's homogeneous-piece submodule, but
  this statement is a direct adapter over
  \texttt{\detokenize{mem_reesAlgebra_iff}} and
  \texttt{\detokenize{reesAlgebra.monomial_mem}}.
\end{lemma}

\begin{proof}
  \leanok
  Let $F\in\mathcal R_A(I)$, and use $F_k$ as defined in
  \Cref{def:nagata-rees-components}.  If
  $F\in\mathcal R_A(I)_n$, coefficient extensionality is the equality
  \[
    \iota(F)=\operatorname{monomial}(n,F_n)=F_nt^n.
  \]
  The coefficient condition in \Cref{def:nagata-rees-algebra} is
  \[
    F_n\in I^n.
  \]
  Conversely, if $a\in I^n$, then
  \[
    \operatorname{coeff}_k(at^n)
      =\begin{cases}a,&k=n,\\0,&k\ne n,
        \end{cases}
  \]
  so $at^n\in\mathcal R_A(I)$ by
  \Cref{def:nagata-rees-algebra} and
  $at^n\in\mathcal R_A(I)_n$ by \Cref{def:nagata-rees-piece}.
\end{proof}

\begin{lemma}[Multiplication of Rees pieces]
  \label{lemma:nagata-rees-piece-mul}
  \lean{AlgebraicGeometry.reesMonomial_mul,
    AlgebraicGeometry.reesGrading.instGradedMonoid}
  \leanok

  For $m,n\in\mathbb N$, multiplication maps
  \[
    \mathcal R_A(I)_m\cdot\mathcal R_A(I)_n
      \subseteq\mathcal R_A(I)_{m+n}.
  \]
\end{lemma}

\begin{proof}
  \uses{lemma:nagata-rees-piece-membership}
  \leanok
  Let $F\in\mathcal R_A(I)_m$ and
  $G\in\mathcal R_A(I)_n$.  By
  \Cref{lemma:nagata-rees-piece-membership}, choose
  $a\in I^m$ and $b\in I^n$ with $F=at^m$ and $G=bt^n$.  Then
  \[
    FG=(at^m)(bt^n)=(ab)t^{m+n},
    \qquad ab\in I^mI^n=I^{m+n}.
  \]
  By \Cref{lemma:nagata-rees-piece-membership},
  $FG\in\mathcal R_A(I)_{m+n}$.
\end{proof}

\begin{lemma}[Finite support of the homogeneous decomposition]
  \label{lemma:nagata-rees-component-finite-support}
  \uses{def:nagata-rees-components}
  \lean{Polynomial.support}
  \leanok

  For $F\in\mathcal R_A(I)$, only finitely many homogeneous components
  $F_nt^n$ are nonzero.

  \paragraph{Mathlib relation.}
  This finiteness is built into
  \texttt{\detokenize{Polynomial.support}}, whose value is a finite set.
\end{lemma}

\begin{proof}
  \leanok
  By \Cref{def:nagata-rees-components},
  \[
    \operatorname{supp}(F)
      =\operatorname{supp}_{A[t]}(\iota(F)).
  \]
  On underlying sets, define the coefficient bijection
  \[
    \operatorname{Coeff}\colon |A[t]|\xrightarrow{\ \sim\ }
      A^{(\mathbb N)}
      :=\left\{(a_n)_{n\in\mathbb N}\in A^{\mathbb N}
        \ \middle|\
        \#\{n\in\mathbb N\mid a_n\ne0\}<\infty\right\},
    \qquad
    \operatorname{Coeff}(H)
      :=(\operatorname{coeff}_n(H))_{n\in\mathbb N}.
  \]
  the defining membership condition for $\iota(F)\in A[t]$ is
  \[
    \#\operatorname{supp}_{A[t]}(\iota(F))<\infty.
  \]
  Hence
  \[
    \#\operatorname{supp}(F)<\infty,
    \qquad
    n\notin\operatorname{supp}(F)
      \Longrightarrow F_n=0
      \Longrightarrow\operatorname{hcomp}_n(F)=F_nt^n=0.
  \]
\end{proof}

\begin{lemma}[Reconstruction from homogeneous components]
  \label{lemma:nagata-rees-sum-components}
  \lean{AlgebraicGeometry.reesGrading_isInternal}
  \leanok

  Every $F\in\mathcal R_A(I)$ is the finite sum of the monomials
  $F_nt^n$ regarded as elements of $\mathcal R_A(I)_n$.
\end{lemma}

\begin{proof}
  \uses{lemma:nagata-rees-component-finite-support,
    lemma:nagata-rees-piece-membership}
  \leanok
  Let $F_n$ and $\operatorname{supp}(F)$ be as in
  \Cref{def:nagata-rees-components}, and put
  $E:=\operatorname{supp}(F)$.  By
  \Cref{lemma:nagata-rees-piece-membership},
  $F_nt^n\in\mathcal R_A(I)_n$.  For every $k\in\mathbb N$,
  \[
    \operatorname{coeff}_k\!\left(\sum_{n\in E}F_nt^n\right)
      =\sum_{n\in E}\operatorname{coeff}_k(F_nt^n)
      =F_k
      =\operatorname{coeff}_k(\iota(F)).
  \]
  Polynomial extensionality is
  \[
    \iota(F)=\iota\!\left(\sum_{n\in E}F_nt^n\right).
  \]
  Injectivity of $\iota$ is the equality
  $F=\sum_{n\in E}F_nt^n$.
\end{proof}

\begin{lemma}[Uniqueness of the homogeneous decomposition]
  \label{lemma:nagata-rees-components-unique}
  \lean{AlgebraicGeometry.reesGrading_isInternal}
  \leanok

  If $E\subseteq\mathbb N$ is finite and
  $H_n\in\mathcal R_A(I)_n$ for $n\in E$, then
  \[
    \sum_{n\in E}H_n=0
      \Longrightarrow (\forall n\in E)\ H_n=0.
  \]
\end{lemma}

\begin{proof}
  \uses{lemma:nagata-rees-piece-membership}
  \leanok
  Let $E\subseteq\mathbb N$ be finite, let
  $H_n\in\mathcal R_A(I)_n$ for $n\in E$, and assume
  $\sum_{n\in E}H_n=0$.  Fix $k\in E$.  Applying
  $\operatorname{coeff}_k\circ\iota$ is the chain
  \[
    0
      =\operatorname{coeff}_k\!\left(\iota\!\left(
          \sum_{n\in E}H_n\right)\right)
      =\sum_{n\in E}\operatorname{coeff}_k(\iota(H_n))
      =\operatorname{coeff}_k(\iota(H_k)).
  \]
  By \Cref{lemma:nagata-rees-piece-membership},
  $\iota(H_k)=\operatorname{coeff}_k(\iota(H_k))t^k=0$; injectivity of
  $\iota$ is injective, so
  \[
    H_k=0.
  \]
\end{proof}

\begin{proposition}[Graded-ring structure on the Rees algebra]
  \label{prop:nagata-rees-graded}
  \lean{AlgebraicGeometry.reesGrading.instGradedRing}
  \leanok

  The family $n\mapsto\mathcal R_A(I)_n$ defines an
  $\mathbb N$-graded-ring structure on $\mathcal R_A(I)$.
\end{proposition}

\begin{proof}
  \uses{lemma:nagata-rees-piece-mul,lemma:nagata-rees-sum-components,
    lemma:nagata-rees-components-unique}
  \leanok
  For all $m,n\in\mathbb N$, graded multiplication is the inclusion
  \[
    \mathcal R_A(I)_m\mathcal R_A(I)_n
      \subseteq\mathcal R_A(I)_{m+n}
  \]
  of \Cref{lemma:nagata-rees-piece-mul}.  For every
  $F\in\mathcal R_A(I)$, \Cref{lemma:nagata-rees-sum-components} supplies a
  finite equality $F=\sum_nH_n$ with
  $H_n\in\mathcal R_A(I)_n$.  If $F=\sum_nK_n$ is a second such expression,
  then
  \[
    0=F-F=\sum_n(H_n-K_n),
  \]
  and \Cref{lemma:nagata-rees-components-unique} asserts
  \[
    H_n=K_n\qquad(n\in\mathbb N).
  \]
  Existence and uniqueness of these components state that the addition map
  \[
    \bigoplus_{n\ge0}\mathcal R_A(I)_n\longrightarrow\mathcal R_A(I),
    \qquad (H_n)_n\longmapsto\sum_nH_n
  \]
  is bijective and multiplication has the required degrees.
\end{proof}

\begin{lemma}[The degree-zero piece]
  \label{lemma:nagata-rees-degree-zero}
  \lean{AlgebraicGeometry.reesGradingZero}
  \leanok

  Let $\iota\colon\mathcal R_A(I)\hookrightarrow A[t]$ be the inclusion in
  \Cref{def:nagata-rees-algebra}.  Define
  \[
  \begin{gathered}
    c_0\colon\mathcal R_A(I)_0\to A,\qquad
      c_0(F):=\operatorname{coeff}_0(\iota(F)),\\
    u_0\colon A\to\mathcal R_A(I)_0,\qquad u_0(a):=at^0,\\
    c_0\circ u_0=\operatorname{id}_A,\qquad
    u_0\circ c_0=\operatorname{id}_{\mathcal R_A(I)_0},\qquad
    (c_0,u_0)\colon\mathcal R_A(I)_0\xrightarrow{\sim}A.
  \end{gathered}                                                    \tag{RZ0}
  \]
\end{lemma}

\begin{proof}
  \uses{lemma:nagata-rees-piece-membership}
  \leanok
  Define
  \[
    c_0\colon\mathcal R_A(I)_0\to A,
      \quad F\mapsto\operatorname{coeff}_0(\iota(F)),
    \qquad
    u_0\colon A\to\mathcal R_A(I)_0,
      \quad a\mapsto at^0.
  \]
  Since $I^0=A$, \Cref{lemma:nagata-rees-piece-membership} makes both maps
  well typed.  For $a\in A$ and $F\in\mathcal R_A(I)_0$,
  \[
    c_0(u_0(a))=a,
    \qquad
    u_0(c_0(F))
      =\operatorname{coeff}_0(\iota(F))t^0=F.
  \]
  For $F,G\in\mathcal R_A(I)_0$ and $a,b\in A$, the coefficient identities
  in $A[t]$ and the monomial multiplication identity are the explicit chains
  \[
  \begin{aligned}
    c_0(F+G)
      &=\operatorname{coeff}_0(\iota(F)+\iota(G))
       =\operatorname{coeff}_0(\iota(F))
          +\operatorname{coeff}_0(\iota(G))
       =c_0(F)+c_0(G),\\
    c_0(FG)
      &=\operatorname{coeff}_0(\iota(F)\iota(G))
       =\sum_{i+j=0}\operatorname{coeff}_i(\iota(F))
          \operatorname{coeff}_j(\iota(G))
       =c_0(F)c_0(G),\\
    c_0(1_{\mathcal R_A(I)_0})
      &=\operatorname{coeff}_0(1_{A[t]})=1_A,
  \end{aligned}
  \]
  and
  \[
  \begin{aligned}
    u_0(a+b)&=(a+b)t^0=at^0+bt^0=u_0(a)+u_0(b),\\
    u_0(ab)&=(ab)t^0=(at^0)(bt^0)=u_0(a)u_0(b),\\
    u_0(1_A)&=1_At^0=1_{\mathcal R_A(I)_0}.
  \end{aligned}
  \]
  These six equalities are the addition, multiplication, and unit axioms for
  $c_0$ and $u_0$; together with
  \[
    c_0\circ u_0=\operatorname{id}_A,
    \qquad
    u_0\circ c_0=\operatorname{id}_{\mathcal R_A(I)_0},
  \]
  they define the asserted ring equivalence.
\end{proof}

\begin{lemma}[Degree-one element associated to an ideal element]
  \label{lemma:nagata-rees-degree-one-element}
  \uses{def:nagata-rees-components}
  \lean{AlgebraicGeometry.reesX, AlgebraicGeometry.reesX_mem}
  \leanok

  Let $\iota\colon\mathcal R_A(I)\hookrightarrow A[t]$ be the inclusion of
  \Cref{def:nagata-rees-components}.  For $a\in I$, define
  \[
    \tau_I(a):=at=\operatorname{monomial}(1,a)\in\mathcal R_A(I)_1.
  \]
  Then
  \[
    \iota(\tau_I(a))=at,
    \qquad
    \operatorname{coeff}_k(\iota(\tau_I(a)))
      =\begin{cases}a,&k=1,\\0,&k\ne1.
        \end{cases}
  \]

  \paragraph{Mathlib relation.}
  The element is \texttt{\detokenize{Polynomial.monomial}}; its membership
  in the Rees algebra is
  \texttt{\detokenize{reesAlgebra.monomial_mem}}, and the coefficient
  formula is \texttt{\detokenize{Polynomial.coeff_monomial}}.
\end{lemma}

\begin{proof}
  \uses{lemma:nagata-rees-piece-membership}
  \leanok
  For $a\in I=I^1$, the definition in the statement and
  \Cref{lemma:nagata-rees-piece-membership} are the equalities
  \[
    \tau_I(a)\in\mathcal R_A(I)_1,
    \qquad
    \iota(\tau_I(a))=at,
    \qquad
    \operatorname{coeff}_k(\iota(\tau_I(a)))
      =\begin{cases}a,&k=1,\\0,&k\ne1.
        \end{cases}
  \]
\end{proof}

\begin{proposition}[Finite generation over degree zero]
  \label{prop:nagata-rees-finite-generation}
  \lean{}

  Let $r\in\mathbb N$ and $(f_i)_{0\leq i\leq r}\in I^{r+1}$.  If
  $I=(f_0,\ldots,f_r)$, then $\mathcal R_A(I)$ is generated as an algebra
  over its degree-zero piece by $f_0t,\ldots,f_rt$.  In particular it is of
  finite type over $\mathcal R_A(I)_0$.

  \paragraph{Mathlib relation.}
  The finite-generation endpoint is the exact theorem
  \texttt{\detokenize{reesAlgebra.fg}}; generation by degree-one monomials
  is \texttt{\detokenize{adjoin_monomial_eq_reesAlgebra}}.  The explicit
  degree-zero transport in this proposition is blueprint-local.
\end{proposition}

\begin{proof}
  \uses{lemma:nagata-rees-degree-zero,lemma:nagata-rees-degree-one-element}
  Let $B\subseteq\mathcal R_A(I)$ be the $A$-subalgebra generated by
  $f_0t,\ldots,f_rt$.  For $n\ge0$,
  \[
    I^n=(f_{i_1}\cdots f_{i_n}\mid
       0\le i_1,\ldots,i_n\le r).
  \]
  By the displayed equality of ideals, every $a\in I^n$ has a finite
  expression
  $a=\sum_\nu c_\nu f_{i_{\nu,1}}\cdots f_{i_{\nu,n}}$, and
  \[
    at^n
      =\sum_\nu c_\nu
        (f_{i_{\nu,1}}t)\cdots(f_{i_{\nu,n}}t)\in B.
  \]
  By \Cref{lemma:nagata-rees-piece-membership}, every homogeneous element has
  this form, and by \Cref{lemma:nagata-rees-sum-components} every Rees
  element is a finite sum of homogeneous elements.  The two inclusions are
  \[
    B\subseteq\mathcal R_A(I),\qquad
    \mathcal R_A(I)\subseteq B,
  \]
  whence $B=\mathcal R_A(I)$.  Under
  $\mathcal R_A(I)_0\xrightarrow{\sim}A$ from
  \Cref{lemma:nagata-rees-degree-zero}, this equality is precisely finite
  generation over the degree-zero piece by $f_0t,\ldots,f_rt$.
\end{proof}

\subsection{Affine blow-ups and their standard charts}

\begin{definition}[Affine blow-up]
  \label{def:nagata-affine-blowup}
  \uses{def:nagata-rees-algebra,def:nagata-rees-piece,
    lemma:nagata-rees-degree-zero}
  \lean{AlgebraicGeometry.AffineBlowup, AlgebraicGeometry.AffineBlowup.toSpec}
  \leanok

  For an ideal $I\subseteq A$, define
  \[
    \operatorname{Bl}_I(\operatorname{Spec}A)
      =\operatorname{Proj}(\mathcal R_A(I)),
  \]
  the absolute homogeneous \(\operatorname{Proj}\) of the graded Rees
  algebra of \Cref{def:nagata-rees-piece}.  Let
  $u_0\colon A\xrightarrow{\sim}\mathcal R_A(I)_0$ be the inverse of
  the equivalence in \Cref{lemma:nagata-rees-degree-zero}.  Define
  \[
    b_{A,I}\colon\operatorname{Proj}(\mathcal R_A(I))\longrightarrow
      \operatorname{Spec}(A)
  \]
  to be the composite of the absolute structure morphism
  $\operatorname{Proj}(\mathcal R_A(I))\to
   \operatorname{Spec}(\mathcal R_A(I)_0)$ with
  $\operatorname{Spec}(u_0)$.

  \paragraph{Mathlib relation.}
  The ambient constructions are
  \texttt{\detokenize{AlgebraicGeometry.Proj}} and its structure map
  \texttt{\detokenize{Proj.toSpecZero}}, specialized to the graded
  Rees algebra.  No relative-\texttt{Proj} construction is needed for the
  affine blow-up.
\end{definition}

\begin{definition}[Affine blow-up chart algebra]
  \label{def:nagata-blowup-chart-ring}
  \lean{}

  Let $a\in I$, let $S_a:=\{a^n\mid n\in\mathbb N\}$, put
  $A_a:=S_a^{-1}A$, and write $\lambda_a\colon A\to A_a$ for the localization
  map.  Define
  \[
    A[I/a]:=A\bigl[\lambda_a(x)\lambda_a(a)^{-1}\mid x\in I\bigr]
      \subseteq A_a.
  \]
\end{definition}

\begin{lemma}[Normal form in a chart algebra]
  \label{lemma:nagata-chart-ring-normal-form}
  \uses{def:nagata-blowup-chart-ring}
  \lean{}

  For $a\in I$, every element of $A[I/a]$ is represented in $A_a$ by $x/a^n$ for some
  $n\geq 0$ and $x\in I^n$.
\end{lemma}

\begin{proof}
  Put
  \[
    N_a:=\{x/a^n\in A_a\mid n\in\mathbb N,\ x\in I^n\}.
  \]
  Since $I^0=A$, one has $\lambda_a(A)\subseteq N_a$.  If
  $x\in I^m$ and $y\in I^n$, then
  \[
    \frac{x}{a^m}+\frac{y}{a^n}
      =\frac{a^nx+a^my}{a^{m+n}},
    \qquad a^nx+a^my\in I^{m+n},
  \]
  and
  \[
    \frac{x}{a^m}\frac{y}{a^n}
      =\frac{xy}{a^{m+n}},
    \qquad xy\in I^mI^n=I^{m+n}.
  \]
  These addition and multiplication identities state that $N_a$ is an
  $A$-subalgebra.  Since $x/a\in N_a$ for $x\in I$, the defining minimality
  of $A[I/a]$ is
  \[
    A[I/a]\subseteq N_a.
  \]
  Conversely, if $x\in I^n$, choose
  $x=\sum_\nu x_{\nu,1}\cdots x_{\nu,n}$ with $x_{\nu,j}\in I$.  Then
  \[
    \frac{x}{a^n}
      =\sum_\nu\prod_{j=1}^n\frac{x_{\nu,j}}a\in A[I/a].
  \]
  This calculation is the reverse inclusion
  $N_a\subseteq A[I/a]$.  Combining it with the preceding inclusion is
  \[
    A[I/a]=N_a.
  \]
\end{proof}

\begin{lemma}[Equality criterion for chart fractions]
  \label{lemma:nagata-chart-ring-fraction-equality}
  \lean{}

  Let $a\in I$ and $m,n\in\mathbb N$.  If $x\in I^m$ and $y\in I^n$, then
  \[
    \frac{x}{a^m}=\frac{y}{a^n}\text{ in }A[I/a]
    \quad\Longleftrightarrow\quad
    \exists N\in\mathbb N,\quad
      a^N(a^nx-a^my)=0\text{ in }A.
  \]
\end{lemma}

\begin{proof}
  \uses{lemma:nagata-chart-ring-normal-form}
  By \Cref{def:nagata-blowup-chart-ring}, $A[I/a]\subseteq A_a$, so equality
  in $A[I/a]$ is equality in $S_a^{-1}A$.  The localization equivalence
  relation is
  \[
    \frac{x}{a^m}=\frac{y}{a^n}
    \quad\Longleftrightarrow\quad
    \exists s\in S_a,\quad s(a^nx-a^my)=0.
  \]
  Substituting $s=a^N$ proves the displayed criterion.
\end{proof}

\begin{definition}[Homogeneous-localization chart map]
  \label{def:nagata-homogeneous-away-to-chart}
  \uses{def:nagata-rees-algebra,
    def:nagata-blowup-chart-ring}
  \lean{}

  For $a\in I$, put $R:=\mathcal R_A(I)$ and define
  \[
    R_{(at)}:=\bigl(R[(at)^{-1}]\bigr)_0.
  \]
  Define $\theta_a\colon R_{(at)}\to A[I/a]$ on a homogeneous fraction by
  \[
    \theta_a\!\left(\frac{xt^n}{(at)^n}\right)=\frac{x}{a^n},
    \qquad x\in I^n.
  \]
\end{definition}

\begin{lemma}[Well-definedness of the homogeneous-localization chart map]
  \label{lemma:nagata-homogeneous-away-to-chart}
  \uses{def:nagata-homogeneous-away-to-chart}
  \lean{}

  For $a\in I$, the assignment $\theta_a$ is independent of the chosen homogeneous
  representative and defines a ring homomorphism from the degree-zero
  homogeneous localization to $A[I/a]$.
\end{lemma}

\begin{proof}
  \uses{lemma:nagata-chart-ring-fraction-equality}
  Put $R:=\mathcal R_A(I)$.  Suppose
  $xt^m/(at)^m=yt^n/(at)^n$ in $R[(at)^{-1}]$, with
  $x\in I^m$ and $y\in I^n$.  The localization relation supplies
  $N\in\mathbb N$ such that
  \[
    (at)^N\bigl((at)^nxt^m-(at)^myt^n\bigr)=0.
  \]
  Applying $\operatorname{coeff}_{N+m+n}$ is the equality
  \[
    a^N(a^nx-a^my)=0.
  \]
  By \Cref{lemma:nagata-chart-ring-fraction-equality},
  $x/a^m=y/a^n$ in $A[I/a]$.  This equality for any two representatives is
  precisely the well-definedness of $\theta_a$.  For
  $x\in I^m$ and $y\in I^n$,
  \[
    \theta_a\!\left(\frac{xt^m}{(at)^m}
      \frac{yt^n}{(at)^n}\right)
      =\frac{xy}{a^{m+n}}
      =\frac{x}{a^m}\frac{y}{a^n},
  \]
  and
  \[
    \theta_a\!\left(\frac{xt^m}{(at)^m}
      +\frac{yt^n}{(at)^n}\right)
      =\frac{a^nx+a^my}{a^{m+n}}
      =\frac{x}{a^m}+\frac{y}{a^n}.
  \]
  Also $\theta_a(1)=1$, so $\theta_a$ is a ring homomorphism.
\end{proof}

\begin{lemma}[Surjectivity of the chart map]
  \label{lemma:nagata-homogeneous-away-to-chart-surjective}
  \lean{}

  For $a\in I$, the map $\theta_a$ of
  \Cref{def:nagata-homogeneous-away-to-chart} is
  surjective.
\end{lemma}

\begin{proof}
  \uses{lemma:nagata-homogeneous-away-to-chart,lemma:nagata-chart-ring-normal-form}
  Let $z\in A[I/a]$.  By
  \Cref{lemma:nagata-chart-ring-normal-form}, choose $n\in\mathbb N$ and
  $x\in I^n$ with $z=x/a^n$.  Then
  \[
    \theta_a\!\left(\frac{xt^n}{(at)^n}\right)=\frac{x}{a^n}=z.
  \]
\end{proof}

\begin{lemma}[Injectivity of the chart map]
  \label{lemma:nagata-homogeneous-away-to-chart-injective}
  \lean{}

  For $a\in I$, the map $\theta_a$ of
  \Cref{def:nagata-homogeneous-away-to-chart} is
  injective.
\end{lemma}

\begin{proof}
  \uses{lemma:nagata-homogeneous-away-to-chart,
    lemma:nagata-chart-ring-fraction-equality}
  Put $R:=\mathcal R_A(I)$ as in
  \Cref{def:nagata-homogeneous-away-to-chart}.  Let
  $z=xt^n/(at)^n\in R_{(at)}$ with $x\in I^n$, and suppose
  $\theta_a(z)=x/a^n=0$.  By
  \Cref{lemma:nagata-chart-ring-fraction-equality}, there is
  $N\in\mathbb N$ such that $a^Nx=0$.  Multiplication in the Rees algebra is
  \[
    (at)^N(xt^n)=a^Nxt^{N+n}=0,
  \]
  and the localization equivalence relation is the equality
  \[
    z=0\quad\text{in }R_{(at)}.
  \]
\end{proof}

\begin{definition}[Canonical affine blow-up chart morphism]
  \label{def:nagata-affine-blowup-chart-map}
  \uses{def:nagata-affine-blowup,
    def:nagata-homogeneous-away-to-chart}
  \lean{}

  For $a\in I$, let $R_{(at)}$ and $\theta_a$ be as in
  \Cref{def:nagata-homogeneous-away-to-chart}.  Define
  \[
    \eta_a\colon A[I/a]\longrightarrow R_{(at)},
    \qquad \eta_a(x/a^n):=xt^n/(at)^n.
  \]
  This is well defined by
  \Cref{lemma:nagata-chart-ring-fraction-equality} and the homogeneous
  localization equivalence relation.  If $x\in I^m$ and $y\in I^n$, then
  \[
  \begin{aligned}
    \eta_a\left(\frac{x}{a^m}+\frac{y}{a^n}\right)
      &=\eta_a\left(\frac{a^nx+a^my}{a^{m+n}}\right)\\
      &=\frac{(a^nx+a^my)t^{m+n}}{(at)^{m+n}}\\
      &=\frac{xt^m}{(at)^m}+\frac{yt^n}{(at)^n},
  \end{aligned}
  \]
  \[
    \eta_a\left(\frac{x}{a^m}\frac{y}{a^n}\right)
      =\frac{xyt^{m+n}}{(at)^{m+n}}
      =\frac{xt^m}{(at)^m}\frac{yt^n}{(at)^n},
    \qquad \eta_a(1)=1.
  \]
  These addition, multiplication, and unit identities state that
  $\eta_a$ is a ring homomorphism.  Under
  $D_+(at)=\operatorname{Spec}(R_{(at)})$, define
  \[
    \chi_a\colon D_+(at)\longrightarrow\operatorname{Spec}(A[I/a])
  \]
  by $\chi_a:=\operatorname{Spec}(\eta_a)$.
\end{definition}

\begin{proposition}[Affine chart identification]
  \label{prop:nagata-affine-blowup-chart}
  \uses{def:nagata-affine-blowup-chart-map}
  \lean{}

  For $a\in I$, the morphism $\chi_a$ of
  \Cref{def:nagata-affine-blowup-chart-map} is an
  isomorphism over
  $\operatorname{Spec}(A)$.
\end{proposition}

\begin{proof}
  \uses{lemma:nagata-homogeneous-away-to-chart-surjective,
    lemma:nagata-homogeneous-away-to-chart-injective}
  For $x\in I^n$, the two definitions are the identities
  \[
    (\theta_a\circ\eta_a)(x/a^n)=x/a^n,
    \qquad
    (\eta_a\circ\theta_a)(xt^n/(at)^n)=xt^n/(at)^n.
  \]
  By \Cref{lemma:nagata-chart-ring-normal-form}, the first equality holds on
  every element of $A[I/a]$; by the definition of $R_{(at)}$, the second
  holds on every homogeneous degree-zero fraction.  Consequently
  \[
    \theta_a\circ\eta_a=\operatorname{id}_{A[I/a]},
    \qquad
    \eta_a\circ\theta_a=\operatorname{id}_{R_{(at)}}.
  \]
  These identities state that $\eta_a$ and $\theta_a$ are inverse ring
  isomorphisms, so $\chi_a=\operatorname{Spec}(\eta_a)$ is an isomorphism.
  If $r\in A$, then
  \[
    \theta_a(rt^0/(at)^0)=r/1.
  \]
  Both structure maps are the ring map $A\to R_{(at)}$ defined by
  $r\mapsto rt^0/(at)^0$, which states that $\chi_a$ is over
  $\operatorname{Spec}(A)$.
\end{proof}

\begin{lemma}[Degree-one generators cover Proj]
  \label{lemma:nagata-rees-generators-cover-proj}
  \lean{}

  Let $r\in\mathbb N$ and $(f_i)_{0\leq i\leq r}\in I^{r+1}$.  If
  $I=(f_0,\ldots,f_r)$, then the standard opens
  $D_+(f_it)$ cover $\operatorname{Proj}(\mathcal R_A(I))$.
\end{lemma}

\begin{proof}
  \uses{prop:nagata-rees-finite-generation}
  Put $R:=\mathcal R_A(I)$ and
  $R_+:=\bigoplus_{n>0}R_n$.  By
  \Cref{prop:nagata-rees-finite-generation},
  \[
    R_+=(f_0t,\ldots,f_rt)R.
  \]
  For a homogeneous prime $\mathfrak p$ representing a point of
  $\operatorname{Proj}(R)$,
  \[
    R_+\nsubseteq\mathfrak p
      \Longrightarrow
      \exists i\in\{0,\ldots,r\},\quad f_it\notin\mathfrak p
      \Longrightarrow
      \mathfrak p\in D_+(f_it).
  \]
  Since every point belongs to at least one displayed standard open, the
  asserted cover is the equality
  \[
    \operatorname{Proj}(R)=\bigcup_{i=0}^rD_+(f_it).
  \]
\end{proof}

\begin{proposition}[Finite affine chart cover]
  \label{prop:nagata-affine-blowup-chart-cover}
  \lean{}

  Let $r\in\mathbb N$ and $(f_i)_{0\leq i\leq r}\in I^{r+1}$.  If
  $I=(f_0,\ldots,f_r)$, the affine opens
  \[
    \operatorname{Spec}(A[I/f_i])
  \]
  form a finite open cover of
  $\operatorname{Bl}_I(\operatorname{Spec}A)$.
\end{proposition}

\begin{proof}
  \uses{prop:nagata-affine-blowup-chart,lemma:nagata-rees-generators-cover-proj}
  By \Cref{lemma:nagata-rees-generators-cover-proj},
  \[
    \operatorname{Bl}_I(\operatorname{Spec}A)
      =\bigcup_{i=0}^rD_+(f_it).
  \]
  For each $i$, \Cref{prop:nagata-affine-blowup-chart} supplies
  \[
    \chi_{f_i}\colon D_+(f_it)
      \xrightarrow{\sim}\operatorname{Spec}(A[I/f_i])
  \]
  over $\operatorname{Spec}(A)$.  Transporting the displayed finite open
  cover through these isomorphisms proves the assertion.
\end{proof}

\begin{lemma}[Blow-up of the unit ideal]
  \label{lemma:nagata-affine-blowup-top}
  \uses{def:nagata-affine-blowup}
  \lean{AlgebraicGeometry.AffineBlowup.isIso_toSpec_of_eq_top}
  \leanok

  If $I=A$, then the structure morphism
  \[
    b_{A,I}\colon
      \operatorname{Bl}_{I}(\operatorname{Spec}A)
      \longrightarrow\operatorname{Spec}(A)
  \]
  is an isomorphism.  We write
  $\upsilon_A:=b_{A,A}$ for this isomorphism.
\end{lemma}

\begin{proof}
  \uses{lemma:nagata-rees-degree-zero,
    lemma:nagata-rees-degree-one-element}
  \leanok
  By \Cref{def:nagata-rees-algebra},
  \[
    \mathcal R_A(A)=\bigoplus_{n\ge0}A^nt^n=A[t].
  \]
  Every homogeneous element $at^n$ of positive degree factors as
  $(at^{n-1})\cdot(1\cdot t)$, so the irrelevant ideal is contained in the
  principal ideal $(t)$ and $\operatorname{Proj}(A[t])=D_+(t)$.  Hence the
  affine chart at the degree-one element $t=1\cdot t$ of
  \Cref{lemma:nagata-rees-degree-one-element} is an isomorphism onto
  $\operatorname{Bl}_{A}(\operatorname{Spec}A)$.  Composed with the
  structure morphism it is the morphism of spectra induced by the canonical
  ring homomorphism
  \[
    A\longrightarrow \mathcal R_A(A)_{(t)},\qquad
    a\longmapsto \frac{at^0}{1},
  \]
  which is bijective: injectivity because $t$ is a non-zerodivisor in
  $A[t]$, surjectivity because every degree-zero fraction $at^n/t^n$ equals
  $at^0/1$.  Two-out-of-three for isomorphisms proves the claim.
\end{proof}

\begin{lemma}[Blow-up of the zero ideal]
  \label{lemma:nagata-affine-blowup-bot}
  \uses{def:nagata-affine-blowup}
  \lean{}

  The defining equality is
  \[
    \operatorname{Bl}_{0}(\operatorname{Spec}A)
      =\operatorname{Proj}(\mathcal R_A(0))=\varnothing.           \tag{R0I0}
  \]
\end{lemma}

\begin{proof}
  If $I=0$, then for $n>0$,
  \[
    I^n=0,
    \qquad \mathcal R_A(0)_n=0
  \]
  by \Cref{lemma:nagata-rees-piece-membership}.  Define the irrelevant ideal
  and compute it by
  \[
    \mathcal R_A(0)_+:=\bigoplus_{n>0}\mathcal R_A(0)_n=0.
  \]
  Therefore
  \[
    |\operatorname{Proj}(\mathcal R_A(0))|
      =\{\mathfrak p\mid \mathfrak p\text{ homogeneous prime and }
        0\nsubseteq\mathfrak p\}=\varnothing.
  \]
  A scheme with empty underlying space is the empty scheme.
\end{proof}

\subsection{Localization compatibility and coherence}

Throughout this subsection, fix the explicitly typed datum
\[
  \mathsf D_{\mathrm{loc}}:=(A,I,s),\qquad
  A\text{ a commutative ring},\qquad I\triangleleft A,\qquad s\in A.
                                                                    \tag{LOC0}
\]
Write $\lambda_s\colon A\to A_s$ for a localization of $A$ away from $s$.

\paragraph{Roadmap note.}  Compared to an earlier draft of this blueprint,
the transports below are organized around the \emph{coefficientwise
functoriality} of the graded Rees algebra rather than around
tensor-product comparison isomorphisms: every ring homomorphism
$\varphi\colon A\to B$ with $I.\varphi$-extension $J$ induces a graded
homomorphism $\mathcal R(\varphi)\colon\mathcal R_A(I)\to\mathcal R_B(J)$,
and the identity, composition, and localization laws are stated for these
maps.  The tensor-product isomorphisms of the earlier draft become the
single chart-wise localization statement
\Cref{prop:nagata-rees-localization}, and the explicit cocycle bookkeeping
becomes the composition law \Cref{lemma:nagata-affine-blowup-restrict-cocycle}.
This matches the Lean implementation in
\texttt{WeilConjectures.Blowup.Functoriality} and
\texttt{WeilConjectures.Blowup.Localization}.

\begin{lemma}[Powers of the extended ideal]
  \label{lemma:nagata-localization-ideal-pow}
  \lean{Ideal.map_pow}
  \leanok

  Let $\varphi\colon A\to B$ be a ring homomorphism and let
  $I B:=\varphi(I)B$ denote the extended ideal.  For every $n\in\mathbb N$,
  \[
    (I^n)B=(IB)^n.
  \]
  In particular $(I^n)A_s=(IA_s)^n$.

  \paragraph{Mathlib relation.}
  Extension is \texttt{\detokenize{Ideal.map}}; the equality is exactly
  \texttt{\detokenize{Ideal.map_pow}}.
\end{lemma}

\begin{proof}
  \leanok
  Induction on $n$, using $(JK)B=(JB)(KB)$ for ideals $J,K$ and
  $AB=B$.
\end{proof}

\begin{lemma}[Coefficientwise transport of monomials]
  \label{lemma:nagata-localization-rees-piece}
  \uses{def:nagata-rees-components}
  \lean{AlgebraicGeometry.reesHom_reesMonomial}
  \leanok

  Let $\varphi\colon A\to B$ be a ring homomorphism and put
  $J:=IB$.  For $n\in\mathbb N$ and $a\in I^n$, one has
  $\varphi(a)\in J^n$, and the assignment
  \[
    a\mathbf t^n\longmapsto\varphi(a)\mathbf t^n
  \]
  sends the degree-$n$ piece $\mathcal R_A(I)_n$ into
  $\mathcal R_B(J)_n$.
\end{lemma}

\begin{proof}
  \uses{lemma:nagata-localization-ideal-pow,
    lemma:nagata-rees-piece-membership}
  \leanok
  By \Cref{lemma:nagata-localization-ideal-pow},
  $\varphi(I^n)\subseteq(I^n)B=J^n$, and
  \Cref{lemma:nagata-rees-piece-membership} identifies the degree-$n$
  pieces with monomials.
\end{proof}

\begin{definition}[Coefficientwise Rees homomorphism]
  \label{def:nagata-rees-localization-map}
  \uses{def:nagata-rees-algebra,def:nagata-rees-piece}
  \lean{AlgebraicGeometry.reesHom}
  \leanok

  Let $\varphi\colon A\to B$ be a ring homomorphism and put $J:=IB$.
  Define the graded ring homomorphism
  \[
    \mathcal R(\varphi)\colon\mathcal R_A(I)\longrightarrow\mathcal R_B(J),
    \qquad
    \textstyle\sum_n a_n\mathbf t^n\longmapsto\sum_n\varphi(a_n)\mathbf t^n,
  \]
  i.e.\ coefficientwise application of $\varphi$ to the underlying
  polynomial; it is well defined by
  \Cref{lemma:nagata-localization-rees-piece} and degree preserving.  For
  the datum (LOC0) put $\rho_s:=\mathcal R(\lambda_s)\colon
  \mathcal R_A(I)\to\mathcal R_{A_s}(IA_s)$.
\end{definition}

\begin{lemma}[Degree-one transport]
  \label{lemma:nagata-rees-localization-degree-one}
  \uses{def:nagata-rees-localization-map,
    lemma:nagata-rees-degree-one-element}
  \lean{AlgebraicGeometry.reesHom_reesX}
  \leanok

  For $a\in I$, the map $\mathcal R(\varphi)$ of
  \Cref{def:nagata-rees-localization-map} satisfies
  $\mathcal R(\varphi)(a\mathbf t)=\varphi(a)\mathbf t$.
\end{lemma}

\begin{proof}
  \leanok
  Immediate from the coefficientwise definition.
\end{proof}

\begin{lemma}[Identity law]
  \label{lemma:nagata-rees-localization-one}
  \uses{def:nagata-rees-localization-map}
  \lean{AlgebraicGeometry.reesHom_id}
  \leanok

  For $\varphi=\operatorname{id}_A$ (so $J=I$), one has
  $\mathcal R(\operatorname{id}_A)=\operatorname{id}_{\mathcal R_A(I)}$.
\end{lemma}

\begin{proof}
  \leanok
  Coefficientwise application of the identity is the identity.
\end{proof}

\begin{lemma}[Composition law]
  \label{lemma:nagata-rees-localization-mul}
  \uses{def:nagata-rees-localization-map}
  \lean{AlgebraicGeometry.reesHom_comp}
  \leanok

  Let $\varphi\colon A\to B$ and $\psi\colon B\to C$ be ring homomorphisms
  and put $J:=IB$, $K:=JC=IC$.  Then
  \[
    \mathcal R(\psi)\circ\mathcal R(\varphi)=\mathcal R(\psi\circ\varphi)
    \colon\mathcal R_A(I)\longrightarrow\mathcal R_C(K).
  \]
  In particular, for iterated localizations
  $A\to A_s\to(A_s)_{t_s}$ the composite of the coefficientwise maps is
  the coefficientwise map of the composite localization; no separate
  cocycle datum is required.
\end{lemma}

\begin{proof}
  \leanok
  Coefficientwise, $\psi\circ\varphi$ applied to each coefficient.
\end{proof}

\begin{lemma}[Congruence]
  \label{lemma:nagata-rees-localization-natural}
  \uses{def:nagata-rees-localization-map}
  \lean{AlgebraicGeometry.reesHom_congr}
  \leanok

  The homomorphism $\mathcal R(\varphi)$ depends only on $\varphi$: if
  $\varphi=\varphi'$ as ring homomorphisms $A\to B$, then
  $\mathcal R(\varphi)=\mathcal R(\varphi')$.  Together with
  \Cref{lemma:nagata-rees-localization-mul} this subsumes the naturality
  of the localization comparisons under ring isomorphisms.
\end{lemma}

\begin{proof}
  \leanok
  Immediate.
\end{proof}

\begin{definition}[Induced morphism of affine blow-ups]
  \label{def:nagata-affine-blowup-restriction-map}
  \uses{def:nagata-affine-blowup,def:nagata-rees-localization-map}
  \lean{AlgebraicGeometry.blowupHom}
  \leanok

  Let $\varphi\colon A\to B$ with $J:=IB$.  Contravariance of
  $\operatorname{Proj}$ applied to $\mathcal R(\varphi)$ defines
  \[
    \operatorname{Bl}(\varphi):=\operatorname{Proj}(\mathcal R(\varphi))
    \colon\operatorname{Bl}_J(\operatorname{Spec}B)
    \longrightarrow\operatorname{Bl}_I(\operatorname{Spec}A).
  \]
  The required irrelevant-ideal condition holds because every homogeneous
  element of $\mathcal R_B(J)$ of positive degree $n$ has coefficient in
  $J^n=(I^n)B$, hence lies in the ideal generated by the image of the
  irrelevant ideal of $\mathcal R_A(I)$.  For the datum (LOC0) put
  $r_s:=\operatorname{Bl}(\lambda_s)$.
\end{definition}

\begin{lemma}[Structure-map compatibility]
  \label{lemma:nagata-rees-localization-degree-zero}
  \uses{def:nagata-affine-blowup-restriction-map,
    lemma:nagata-rees-degree-zero}
  \lean{AlgebraicGeometry.blowupHom_toSpec}
  \leanok

  On degree zero, $\mathcal R(\varphi)$ is $\varphi$ itself under the
  identifications $u_0$ of \Cref{lemma:nagata-rees-degree-zero}.
  Consequently the square
  \[
  \begin{aligned}
    \operatorname{Bl}(\varphi)\circ b_{A,I}^{\phantom{1}}
    \colon\quad
    \operatorname{Bl}_J(\operatorname{Spec}B)
    &\longrightarrow\operatorname{Spec}(A)
  \end{aligned}
  \qquad
    b_{A,I}\circ\operatorname{Bl}(\varphi)
      =\operatorname{Spec}(\varphi)\circ b_{B,J}
  \]
  commutes.
\end{lemma}

\begin{proof}
  \uses{lemma:nagata-rees-degree-zero}
  \leanok
  Both composites agree chart-wise on the standard affine charts, where
  they are induced by the commuting square of degree-zero ring maps.
\end{proof}

\begin{lemma}[Identity and cocycle laws for the induced morphisms]
  \label{lemma:nagata-affine-blowup-restrict-cocycle}
  \uses{def:nagata-affine-blowup-restriction-map}
  \lean{AlgebraicGeometry.blowupHom_id, AlgebraicGeometry.blowupHom_comp}
  \leanok

  One has
  \[
    \operatorname{Bl}(\operatorname{id}_A)
      =\operatorname{id}_{\operatorname{Bl}_I(\operatorname{Spec}A)},
    \qquad
    \operatorname{Bl}(\psi\circ\varphi)
      =\operatorname{Bl}(\varphi)\circ\operatorname{Bl}(\psi)
  \]
  for composable $\varphi\colon A\to B$, $\psi\colon B\to C$ as in
  \Cref{lemma:nagata-rees-localization-mul}.  These two laws replace the
  explicit gluing cocycle of the earlier draft.
\end{lemma}

\begin{proof}
  \uses{lemma:nagata-rees-localization-one,
    lemma:nagata-rees-localization-mul}
  \leanok
  Apply functoriality of $\operatorname{Proj}$ to
  \Cref{lemma:nagata-rees-localization-one,lemma:nagata-rees-localization-mul}.
\end{proof}

\begin{proposition}[Chart-wise localization]
  \label{prop:nagata-rees-localization}
  \uses{def:nagata-rees-localization-map,def:nagata-affine-blowup,
    lemma:nagata-rees-piece-membership}
  \lean{AlgebraicGeometry.AffineBlowup.isLocalization_away_map}
  \leanok

  Let $J:=IA_s$ and let $a\in I$.  Write
  $\mathcal R_A(I)_{(a\mathbf t)}$ for the degree-zero homogeneous
  localization at the degree-one element $a\mathbf t$ (the coordinate ring
  of the standard chart $D_+(a\mathbf t)$ of
  $\operatorname{Bl}_I(\operatorname{Spec}A)$), and let
  \[
    \mathcal R(\lambda_s)_{(a\mathbf t)}\colon
    \mathcal R_A(I)_{(a\mathbf t)}
    \longrightarrow
    \mathcal R_{A_s}(J)_{(\lambda_s(a)\mathbf t)}
  \]
  be the map induced by $\rho_s$ on chart rings.  This map exhibits the
  target as the localization of the source away from the degree-zero
  fraction $s\mathbf t^0/1$.
\end{proposition}

\begin{proof}
  \uses{lemma:nagata-localization-ideal-pow,
    lemma:nagata-rees-piece-membership,
    lemma:nagata-rees-localization-degree-zero}
  \leanok
  Write elements of both chart rings in the normal form
  $x\mathbf t^n/(a\mathbf t)^n$ with $x\in I^n$, respectively
  $y\mathbf t^n/(\lambda_s(a)\mathbf t)^n$ with $y\in J^n$
  (\Cref{lemma:nagata-rees-piece-membership}).  The image of
  $s\mathbf t^0/1$ is the unit $\lambda_s(s)\mathbf t^0/1$.  For
  surjectivity up to powers of $s$, expand $y\in J^n=(I^n)A_s$
  (\Cref{lemma:nagata-localization-ideal-pow}) as
  $\lambda_s(s)^m y=\lambda_s(c)$ with $c\in I^n$; then
  $y\mathbf t^n/(\lambda_s(a)\mathbf t)^n$ times the image of
  $(s\mathbf t^0/1)^m$ is the image of
  $c\mathbf t^n/(a\mathbf t)^n$.  For the kernel condition, an equality of
  two images is, after clearing denominators, an equality of monomials in
  $\mathcal R_{A_s}(J)$; extracting the relevant polynomial coefficient
  and using the defining property of the localization $A_s$ produces
  $s^m\bigl(a^{k+n_2}c_1-a^{k+n_1}c_2\bigr)=0$ in $A$ for some
  $m,k$, which is precisely the statement that the two normal forms agree
  after multiplication by $(s\mathbf t^0/1)^m$.
\end{proof}

\begin{proposition}[Restriction of an affine blow-up]
  \label{prop:nagata-affine-blowup-restrict}
  \uses{def:nagata-affine-blowup-restriction-map}
  \lean{AlgebraicGeometry.isPullback_blowupHom_of_isLocalization,
    AlgebraicGeometry.isOpenImmersion_blowupHom_of_isLocalization}
  \leanok

  For the datum (LOC0) with $J:=IA_s$, the commuting square
  \[
    \begin{array}{ccc}
      \operatorname{Bl}_J(\operatorname{Spec}A_s)
        &\xrightarrow{\ r_s\ }
        &\operatorname{Bl}_I(\operatorname{Spec}A)\\
      \downarrow b_{A_s,J}& &\downarrow b_{A,I}\\
      \operatorname{Spec}(A_s)
        &\xrightarrow{\operatorname{Spec}(\lambda_s)}
        &\operatorname{Spec}(A)
    \end{array}
  \]
  is cartesian.  In particular $r_s$ is an open immersion with image
  $b_{A,I}^{-1}(D(s))$, i.e.\ the blow-up commutes with restriction to a
  basic open.
\end{proposition}

\begin{proof}
  \uses{prop:nagata-rees-localization,
    lemma:nagata-rees-localization-degree-zero,
    lemma:nagata-affine-blowup-restrict-cocycle}
  \leanok
  The standard charts $D_+(a\mathbf t)$, $a\in I$, cover
  $\operatorname{Bl}_I(\operatorname{Spec}A)$, and by
  \Cref{prop:nagata-rees-localization} the restriction of $r_s$ over each
  chart is the morphism of spectra of a localization away from one
  element, hence an open immersion with image the basic open of
  $s\mathbf t^0/1$.  Openness of immersions and the computation of the
  ranges are Zariski-local on the target, so $r_s$ is an open immersion
  whose range is the preimage of the range of
  $\operatorname{Spec}(\lambda_s)$ under $b_{A,I}$.  A commuting square of
  schemes whose vertical maps are open immersions with this preimage
  compatibility is cartesian.
\end{proof}

\subsection{Global blow-up by gluing}

\paragraph{Roadmap note.}  Compared to an earlier draft of this blueprint,
the global blow-up is glued over the \emph{canonical locally directed
cover by all affine opens} of $X$ rather than over a chosen affine open
cover with pairwise overlap isomorphisms.  Because the index category of
this cover is a thin preorder, no cocycle conditions arise: the gluing
input is a functor of local blow-ups together with equifibered structure
maps, and the overlap content of the earlier draft is carried by the
composition law \Cref{lemma:nagata-affine-blowup-restrict-cocycle} and by
\Cref{lemma:nagata-global-blowup-glue-data-valid}.  This matches the Lean
implementation in \texttt{WeilConjectures.Blowup.Gluing}, which uses the
Mathlib relative-gluing machinery
(\texttt{\detokenize{AlgebraicGeometry.Scheme.Cover.RelativeGluingData}},
\texttt{\detokenize{AlgebraicGeometry.Scheme.directedAffineCover}},
Stacks 01LH).  The comparison items after
\Cref{prop:nagata-global-blowup-affine-restrict} retain the notation of
the earlier draft and are not part of the current formalization.

Throughout this subsection let $X$ be a scheme and let
$\mathcal I\subseteq\mathcal O_X$ be a quasi-coherent ideal, encoded as
the datum of an ideal $I_U:=\mathcal I(U)\subseteq\Gamma(U,\mathcal O_X)$
for every affine open $U\subseteq X$, compatible with restriction to
basic affine opens.  For affine opens $U\subseteq V$ the compatibility
extends to $I_V\cdot\Gamma(U,\mathcal O_X)=I_U$.

\paragraph{Mathlib relation.}
The encoding is exactly
\texttt{\detokenize{AlgebraicGeometry.Scheme.IdealSheafData}}; the
restriction compatibility for arbitrary inclusions of affine opens is
\texttt{\detokenize{AlgebraicGeometry.Scheme.IdealSheafData.map_ideal}}.

\begin{definition}[Local blow-up models and their transition maps]
  \label{def:nagata-global-blowup-local-model}
  \uses{def:nagata-affine-blowup,def:nagata-affine-blowup-restriction-map}
  \lean{AlgebraicGeometry.Scheme.IdealSheafData.blowupFunctor}
  \leanok

  For an affine open $U\subseteq X$ with coordinate ring
  $A_U:=\Gamma(U,\mathcal O_X)$, define the local blow-up model
  \[
    \mathsf B_U(\mathcal I):=\operatorname{Bl}_{I_U}(\operatorname{Spec}A_U),
    \qquad
    b_U:=b_{A_U,I_U}\colon\mathsf B_U(\mathcal I)\longrightarrow
      \operatorname{Spec}(A_U).
  \]
  For an inclusion $U\subseteq V$ of affine opens, the restriction
  homomorphism
  $\operatorname{res}^V_U\colon A_V\to A_U$ satisfies
  $I_V A_U=I_U$, so \Cref{def:nagata-affine-blowup-restriction-map}
  supplies the transition morphism
  \[
    \operatorname{Bl}(\operatorname{res}^V_U)\colon
      \mathsf B_U(\mathcal I)\longrightarrow\mathsf B_V(\mathcal I).
  \]
  By \Cref{lemma:nagata-affine-blowup-restrict-cocycle} these data form a
  functor from the poset of affine opens of $X$ to schemes.
\end{definition}

\begin{definition}[Structure maps of the local models]
  \label{def:nagata-global-blowup-overlap-basic-map}
  \uses{def:nagata-global-blowup-local-model,
    lemma:nagata-rees-localization-degree-zero}
  \lean{AlgebraicGeometry.Scheme.IdealSheafData.blowupNatTrans}
  \leanok

  For every affine open $U\subseteq X$, composing the structure morphism
  $b_U$ with the canonical isomorphism
  $\operatorname{Spec}(A_U)\cong U$ defines
  \[
    \beta_U\colon\mathsf B_U(\mathcal I)\longrightarrow U.
  \]
  For every inclusion $U\subseteq V$ of affine opens, the square
  \[
    \begin{array}{ccc}
      \mathsf B_U(\mathcal I)
        &\xrightarrow{\operatorname{Bl}(\operatorname{res}^V_U)}
        &\mathsf B_V(\mathcal I)\\
      \downarrow\beta_U& &\downarrow\beta_V\\
      U&\hookrightarrow&V
    \end{array}
  \]
  commutes by \Cref{lemma:nagata-rees-localization-degree-zero}; i.e.\ the
  $\beta_U$ form a natural transformation from the functor of
  \Cref{def:nagata-global-blowup-local-model} to the inclusion diagram of
  the affine opens of $X$.
\end{definition}

\begin{lemma}[Basic-open inclusions give cartesian squares]
  \label{lemma:nagata-global-blowup-overlap-basic}
  \uses{def:nagata-global-blowup-overlap-basic-map}
  \lean{AlgebraicGeometry.Scheme.IdealSheafData.isPullback_blowupNatTrans_of_basicOpen}
  \leanok

  Let $U\subseteq V$ be affine opens of $X$ such that $U$ is a basic open
  of $V$, say $U=D_V(f)$ with $f\in A_V$.  Then the square of
  \Cref{def:nagata-global-blowup-overlap-basic-map} is cartesian.
\end{lemma}

\begin{proof}
  \uses{prop:nagata-affine-blowup-restrict}
  \leanok
  In this case $A_U$ is a localization of $A_V$ away from $f$ and
  $I_U=I_VA_U$, so the claim is
  \Cref{prop:nagata-affine-blowup-restrict} transported along the
  canonical isomorphisms $\operatorname{Spec}(A_U)\cong U$ and
  $\operatorname{Spec}(A_V)\cong V$.
\end{proof}

\begin{lemma}[General inclusions give open immersions with the expected range]
  \label{lemma:nagata-global-blowup-overlap-independent}
  \uses{def:nagata-global-blowup-overlap-basic-map}
  \lean{AlgebraicGeometry.Scheme.IdealSheafData.isOpenImmersion_blowupFunctor_map,
    AlgebraicGeometry.Scheme.IdealSheafData.opensRange_blowupFunctor_map}
  \leanok

  Let $U\subseteq V$ be arbitrary affine opens of $X$.  Then
  $\operatorname{Bl}(\operatorname{res}^V_U)$ is an open immersion whose
  range is $\beta_V^{-1}(U)$.
\end{lemma}

\begin{proof}
  \uses{lemma:nagata-global-blowup-overlap-basic,
    lemma:nagata-affine-blowup-restrict-cocycle}
  \leanok
  Every point of $U$ lies in an affine open $W$ that is basic in $V$ and
  contained in $U$; such a $W$ is then automatically basic in $U$ as
  well.  Over each such $W$, both inclusions $W\subseteq U$ and
  $W\subseteq V$ are covered by \Cref{lemma:nagata-global-blowup-overlap-basic},
  and the composition law
  \Cref{lemma:nagata-affine-blowup-restrict-cocycle} identifies the
  restriction of $\operatorname{Bl}(\operatorname{res}^V_U)$ over the
  corresponding part of $\mathsf B_V(\mathcal I)$ with
  $\operatorname{Bl}(\operatorname{res}^U_W)$ composed with an
  isomorphism.  Being an open immersion is Zariski-local on the target,
  and the ranges are computed on the same cover.
\end{proof}

\begin{lemma}[Identity and symmetry are automatic]
  \label{lemma:nagata-global-blowup-overlap-id-symm}
  \uses{def:nagata-global-blowup-local-model}
  \lean{AlgebraicGeometry.Scheme.IdealSheafData.blowupFunctor}
  \leanok

  The transition maps of \Cref{def:nagata-global-blowup-local-model}
  satisfy
  $\operatorname{Bl}(\operatorname{res}^U_U)=\operatorname{id}$ and, for
  $U\subseteq V\subseteq W$,
  $\operatorname{Bl}(\operatorname{res}^W_V)\circ
   \operatorname{Bl}(\operatorname{res}^V_U)
   =\operatorname{Bl}(\operatorname{res}^W_U)$;
  i.e.\ the local models form a functor on the poset of affine opens.
  There are no symmetry data to check: the index category is a poset, so
  distinct affine opens are related only through common refinements.
\end{lemma}

\begin{proof}
  \uses{lemma:nagata-affine-blowup-restrict-cocycle}
  \leanok
  Functoriality of $\operatorname{Bl}(-)$
  (\Cref{lemma:nagata-affine-blowup-restrict-cocycle}) applied to the
  presheaf identities
  $\operatorname{res}^U_U=\operatorname{id}$ and
  $\operatorname{res}^V_U\circ\operatorname{res}^W_V
    =\operatorname{res}^W_U$.
\end{proof}

\begin{lemma}[No cocycle condition]
  \label{lemma:nagata-global-blowup-overlap-cocycle}
  \uses{lemma:nagata-global-blowup-overlap-id-symm}
  \lean{AlgebraicGeometry.Scheme.IdealSheafData.blowupFunctor}
  \leanok

  In the locally directed gluing of
  \Cref{def:nagata-global-blowup-glue-data} the triple-overlap cocycle
  condition of the earlier draft is vacuous: the index category of the
  cover by all affine opens is a thin preorder, so any two parallel
  composites of transition maps agree by
  \Cref{lemma:nagata-global-blowup-overlap-id-symm}.
\end{lemma}

\begin{proof}
  \leanok
  Immediate from thinness of the index category and functoriality.
\end{proof}

\begin{definition}[Relative gluing datum]
  \label{def:nagata-global-blowup-glue-data}
  \uses{def:nagata-global-blowup-local-model,
    def:nagata-global-blowup-overlap-basic-map}
  \lean{AlgebraicGeometry.Scheme.IdealSheafData.blowupGluingData}
  \leanok

  Define $\mathsf G(X,\mathcal I)$ to be the relative gluing datum over
  the locally directed cover of $X$ by all of its affine opens whose
  functor is the functor of local blow-up models
  (\Cref{def:nagata-global-blowup-local-model}) and whose structure maps
  are the $\beta_U$ of
  \Cref{def:nagata-global-blowup-overlap-basic-map}.
\end{definition}

\begin{lemma}[The gluing datum is equifibered]
  \label{lemma:nagata-global-blowup-glue-data-valid}
  \uses{def:nagata-global-blowup-glue-data}
  \lean{AlgebraicGeometry.Scheme.IdealSheafData.blowupNatTrans_equifibered}
  \leanok

  For every inclusion $U\subseteq V$ of affine opens of $X$, the
  naturality square of \Cref{def:nagata-global-blowup-overlap-basic-map}
  is cartesian.  Hence $\mathsf G(X,\mathcal I)$ is a valid relative
  gluing datum in the sense of Stacks 01LH.
\end{lemma}

\begin{proof}
  \uses{lemma:nagata-global-blowup-overlap-basic,
    lemma:nagata-global-blowup-overlap-independent}
  \leanok
  A commuting square whose two vertical maps are open immersions and in
  which the range of the top map is the preimage of the range of the
  bottom map is cartesian.  The top map
  $\operatorname{Bl}(\operatorname{res}^V_U)$ is an open immersion with
  range $\beta_V^{-1}(U)$ by
  \Cref{lemma:nagata-global-blowup-overlap-independent}, and the bottom
  map $U\hookrightarrow V$ is an open immersion.
\end{proof}

\begin{definition}[Global blow-up]
  \label{def:nagata-global-blowup}
  \uses{def:nagata-global-blowup-glue-data,
    lemma:nagata-global-blowup-glue-data-valid}
  \lean{AlgebraicGeometry.Scheme.IdealSheafData.blowup,
    AlgebraicGeometry.Scheme.IdealSheafData.blowupι}
  \leanok

  Let $\mathsf G(X,\mathcal I)$ be the gluing datum of
  \Cref{def:nagata-global-blowup-glue-data}.  Define
  $\operatorname{Bl}_{\mathcal I}(X)$ to be its glued scheme (the
  colimit of the functor of local models), and denote the chart
  embeddings by
  \[
    e_U\colon\mathsf B_U(\mathcal I)
      \hookrightarrow\operatorname{Bl}_{\mathcal I}(X).
  \]
  They are open immersions, their images cover
  $\operatorname{Bl}_{\mathcal I}(X)$, and
  $e_V\circ\operatorname{Bl}(\operatorname{res}^V_U)=e_U$ for
  $U\subseteq V$.
\end{definition}

\begin{lemma}[Compatibility of the local structure maps]
  \label{lemma:nagata-global-blowup-local-maps-agree}
  \uses{def:nagata-global-blowup-overlap-basic-map}
  \lean{AlgebraicGeometry.Scheme.IdealSheafData.blowupι_blowupMap}
  \leanok

  With $v_U\colon U\hookrightarrow X$ the inclusions, one has for every
  inclusion $U\subseteq V$ of affine opens
  \[
    v_V\circ\beta_V\circ\operatorname{Bl}(\operatorname{res}^V_U)
      =v_U\circ\beta_U.
  \]
\end{lemma}

\begin{proof}
  \uses{lemma:nagata-rees-localization-degree-zero}
  \leanok
  Naturality of $\beta$
  (\Cref{def:nagata-global-blowup-overlap-basic-map}) composed with
  $v_V\circ(U\hookrightarrow V)=v_U$.
\end{proof}

\begin{definition}[Global blow-up morphism]
  \label{def:nagata-global-blowup-map}
  \uses{def:nagata-global-blowup,
    def:nagata-global-blowup-local-model,
    lemma:nagata-global-blowup-local-maps-agree}
  \lean{AlgebraicGeometry.Scheme.IdealSheafData.blowupMap}
  \leanok

  By \Cref{lemma:nagata-global-blowup-local-maps-agree} the composites
  $v_U\circ\beta_U\colon\mathsf B_U(\mathcal I)\to X$ are compatible
  with the transition maps, so they glue to a unique morphism
  \[
    b_{X,\mathcal I}\colon
      \operatorname{Bl}_{\mathcal I}(X)\longrightarrow X,
    \qquad
    b_{X,\mathcal I}\circ e_U=v_U\circ\beta_U\quad(\forall U).
  \]
  Uniqueness holds because the $e_U(\mathsf B_U(\mathcal I))$ cover
  $\operatorname{Bl}_{\mathcal I}(X)$.
\end{definition}

\begin{definition}[Affine-restriction comparison for the global blow-up]
  \label{def:nagata-global-blowup-affine-restriction-map}
  \uses{def:nagata-global-blowup-map,
    def:nagata-global-blowup-local-model}
  \lean{AlgebraicGeometry.Scheme.IdealSheafData.blowupMap_preimage}
  \leanok

  For every affine open $U\subseteq X$, the preimage
  $b_{X,\mathcal I}^{-1}(U)$ equals the image of the chart embedding
  $e_U$; the restriction of $b_{X,\mathcal I}$ over $U$ is thereby
  identified with $v_U\circ\beta_U$ on
  $\mathsf B_U(\mathcal I)$.
\end{definition}

\begin{proposition}[Affine restriction of the global blow-up]
  \label{prop:nagata-global-blowup-affine-restrict}
  \uses{def:nagata-global-blowup-affine-restriction-map}
  \lean{AlgebraicGeometry.Scheme.IdealSheafData.isPullback_blowupι}
  \leanok

  For every affine open $U\subseteq X$, the square
  \[
    \begin{array}{ccc}
      \mathsf B_U(\mathcal I)
        &\xrightarrow{\ e_U\ }
        &\operatorname{Bl}_{\mathcal I}(X)\\
      \downarrow\beta_U& &\downarrow b_{X,\mathcal I}\\
      U&\xhookrightarrow{\ v_U\ }&X
    \end{array}
  \]
  is cartesian; i.e.\
  $b_{X,\mathcal I}^{-1}(U)\cong
    \operatorname{Bl}_{I_U}(\operatorname{Spec}A_U)$
  compatibly with the two structure maps.
\end{proposition}

\begin{proof}
  \uses{lemma:nagata-global-blowup-glue-data-valid}
  \leanok
  This is the defining property of the relative gluing of an equifibered
  datum (Stacks 01LH): over each member of the cover, the glued object
  restricts to the given local object.  It is supplied directly by the
  Mathlib relative-gluing machinery applied to
  $\mathsf G(X,\mathcal I)$, which is admissible by
  \Cref{lemma:nagata-global-blowup-glue-data-valid}.
\end{proof}

\begin{definition}[Comparison between two affine-cover constructions]
  \label{def:nagata-global-blowup-cover-comparison}
  \uses{def:nagata-choice-notation,
    def:nagata-global-blowup-glue-data,def:nagata-global-blowup,
    def:nagata-global-blowup-affine-restriction-map,
    lemma:nagata-affine-blowup-restrict-cocycle,
    lemma:nagata-morphism-property-calculus}
  \lean{}

  Let
  \(\mathcal V=(v_i\colon V_i=\operatorname{Spec}(A_i)\to X)_{i\in I}\)
  and
  \(\mathcal W=(w_j\colon W_j=\operatorname{Spec}(B_j)\to X)_{j\in J}\)
  satisfy
  \[
    \operatorname{IsOpenImmersion}(v_i),\quad
    \operatorname{IsOpenImmersion}(w_j),\quad
    |X|=\bigcup_{i\in I}\operatorname{Im}(|v_i|)
       =\bigcup_{j\in J}\operatorname{Im}(|w_j|).                  \tag{CCR0}
  \]
  By \Cref{def:nagata-global-blowup-glue-data,def:nagata-global-blowup},
  define
  \[
    B_{\mathcal V}:=
      \operatorname{Bl}_{\mathcal I}(X;\mathcal V),\qquad
    B_{\mathcal W}:=
      \operatorname{Bl}_{\mathcal I}(X;\mathcal W),\\
    b_{\mathcal V}:=b_{X,\mathcal I}^{\mathcal V}
      \colon B_{\mathcal V}\to X,\qquad
    b_{\mathcal W}:=b_{X,\mathcal I}^{\mathcal W}
      \colon B_{\mathcal W}\to X.                               \tag{CCR0a}
  \]
  where
  $\operatorname{Bl}_{\mathcal I}(X;\mathcal V)$ and
  $\operatorname{Bl}_{\mathcal I}(X;\mathcal W)$ denote exactly the glued
  schemes of \Cref{def:nagata-global-blowup} instantiated with
  $\mathsf G(X,\mathcal I,(V_i)_{i\in I})$ and
  $\mathsf G(X,\mathcal I,(W_j)_{j\in J})$, respectively; the symbols
  $b_{X,\mathcal I}^{\mathcal V}$ and
  $b_{X,\mathcal I}^{\mathcal W}$ denote the corresponding morphisms of
  \Cref{def:nagata-global-blowup-map}.
  Define
  \[
    \mathscr R(\mathcal V,\mathcal W):=
    \left\{\bigl(K,a,b,(R_k,u_k,u_k^{\mathcal V},
        u_k^{\mathcal W})_{k\in K}\bigr)\ \middle|\
    \begin{array}{l}
      K\text{ is a set},\quad a\colon K\to I,\quad b\colon K\to J,\\
      a_k:=a(k),\quad b_k:=b(k),\\
      Q_k:=\operatorname{Spec}(R_k),\quad
      u_k\colon Q_k\to X,\quad
      u_k^{\mathcal V}\colon Q_k\to V_{a_k},\quad
      u_k^{\mathcal W}\colon Q_k\to W_{b_k},\\
      \operatorname{IsOpenImmersion}(u_k),\quad
      \operatorname{IsOpenImmersion}(u_k^{\mathcal V}),\quad
      \operatorname{IsOpenImmersion}(u_k^{\mathcal W}),\\
      v_{a_k}\circ u_k^{\mathcal V}=u_k
        =w_{b_k}\circ u_k^{\mathcal W},\\
      |X|=\displaystyle\bigcup_{k\in K}\operatorname{Im}(|u_k|)
    \end{array}\right\}.                                           \tag{CCR1}
  \]
  The affine-basis axiom applied to each
  \(V_i\times_XW_j\) and the cover equations (CCR0) are the following
  nonemptiness predicate:
  \[
    \mathscr R(\mathcal V,\mathcal W)\neq\varnothing.              \tag{CCR2}
  \]
  By \Cref{def:nagata-choice-notation}, define the common refinement
  \[
    \mathsf R_{\mathcal V,\mathcal W}
      :=\operatorname{ch}(\mathscr R(\mathcal V,\mathcal W))
      =\bigl(K,a,b,(R_k,u_k,u_k^{\mathcal V},
        u_k^{\mathcal W})_{k\in K}\bigr),\qquad
    a_k:=a(k),\quad b_k:=b(k),\quad
    Q_k:=\operatorname{Spec}(R_k),\quad
    I_{Q_k}:=\Gamma(Q_k,\mathcal I|_{Q_k}).                        \tag{CCR3}
  \]
  Let
  \[
    B_{\mathcal V}|_{Q_k}:=
      B_{\mathcal V}\times_{X,(b_{\mathcal V},u_k)}Q_k,
    \qquad
    B_{\mathcal W}|_{Q_k}:=
      B_{\mathcal W}\times_{X,(b_{\mathcal W},u_k)}Q_k,
    \qquad
    q_{\mathcal V,k}\colon
      B_{\mathcal V}|_{Q_k}\xrightarrow{\sim}
      \operatorname{Bl}_{I_{Q_k}}(Q_k),
    \qquad
    q_{\mathcal W,k}\colon
      B_{\mathcal W}|_{Q_k}\xrightarrow{\sim}
      \operatorname{Bl}_{I_{Q_k}}(Q_k)
  \]
  be the isomorphisms of
  \Cref{prop:nagata-global-blowup-affine-restrict}.  Define
  \[
    \gamma_k:=q_{\mathcal W,k}^{-1}\circ q_{\mathcal V,k}
      \colon B_{\mathcal V}|_{Q_k}\xrightarrow{\sim}
        B_{\mathcal W}|_{Q_k}.
  \]
  Let
  $\iota_{\mathcal W,k}:=\operatorname{pr}_{B_{\mathcal W}}\colon
  B_{\mathcal W}|_{Q_k}\hookrightarrow B_{\mathcal W}$ be the open
  immersion and put
  \[
    \widehat\gamma_k:=
      \iota_{\mathcal W,k}\circ
      q_{\mathcal W,k}^{-1}\circ q_{\mathcal V,k}\colon
      B_{\mathcal V}|_{Q_k}\longrightarrow B_{\mathcal W}.
  \]
  For $k,\ell\in K$, define
  \[
    Q_{k\ell}:=Q_k\times_XQ_\ell,\qquad
    I_{Q_{k\ell}}:=\Gamma(Q_{k\ell},\mathcal I|_{Q_{k\ell}}),\qquad
    B_{k\ell}:=\operatorname{Bl}_{I_{Q_{k\ell}}}(Q_{k\ell}),
    \qquad
    B_{\mathcal V}|_{Q_{k\ell}}:=B_{\mathcal V}\times_XQ_{k\ell},
    \qquad
    B_{\mathcal W}|_{Q_{k\ell}}:=B_{\mathcal W}\times_XQ_{k\ell},\\
    \iota_{\mathcal W,k\ell}:=\operatorname{pr}_{B_{\mathcal W}}\colon
      B_{\mathcal W}|_{Q_{k\ell}}\hookrightarrow B_{\mathcal W}.
                                                                    \tag{CCR4}
  \]
  The restriction-cocycle equalities of
  \Cref{lemma:nagata-affine-blowup-restrict-cocycle}, instantiated for
  the two covers, are
  \[
  \begin{gathered}
    q_{\mathcal V,k}|_{Q_{k\ell}}
      =q_{\mathcal V,k\ell}
      =q_{\mathcal V,\ell}|_{Q_{k\ell}}
      \colon B_{\mathcal V}|_{Q_{k\ell}}\xrightarrow{\sim}B_{k\ell},\\
    q_{\mathcal W,k}|_{Q_{k\ell}}
      =q_{\mathcal W,k\ell}
      =q_{\mathcal W,\ell}|_{Q_{k\ell}}
      \colon B_{\mathcal W}|_{Q_{k\ell}}\xrightarrow{\sim}B_{k\ell}.
  \end{gathered}                                                    \tag{CCR5}
  \]
  Substitution of (CCR5) into the definitions of
  $\widehat\gamma_k$ and $\widehat\gamma_\ell$ is the chain
  \[
  \begin{aligned}
    \widehat\gamma_k|_{B_{\mathcal V}|_{Q_{k\ell}}}
      &=\iota_{\mathcal W,k\ell}\circ
        (q_{\mathcal W,k}|_{Q_{k\ell}})^{-1}\circ
        q_{\mathcal V,k}|_{Q_{k\ell}}\\
      &=\iota_{\mathcal W,k\ell}\circ
        q_{\mathcal W,k\ell}^{-1}\circ q_{\mathcal V,k\ell}\\
      &=\iota_{\mathcal W,k\ell}\circ
        (q_{\mathcal W,\ell}|_{Q_{k\ell}})^{-1}\circ
        q_{\mathcal V,\ell}|_{Q_{k\ell}}\\
      &=\widehat\gamma_\ell|_{B_{\mathcal V}|_{Q_{k\ell}}}.
  \end{aligned}                                                     \tag{CCR6}
  \]
  Define
  \[
    \mathscr C_{\mathcal V,\mathcal W}:=
      \left\{h\in\operatorname{Hom}(B_{\mathcal V},B_{\mathcal W})
        \ \middle|\
        h|_{B_{\mathcal V}|_{Q_k}}=\widehat\gamma_k
        \text{ for every }k\right\}.
  \]
  In the notation (PC11) of
  \Cref{lemma:nagata-morphism-property-calculus}, (CCR6) is the membership
  \[
    \boldsymbol\gamma:=(\widehat\gamma_k)_{k\in K}
      \in\mathscr G_{B_{\mathcal W}}
        \bigl(B_{\mathcal V},(B_{\mathcal V}|_{Q_k})_{k\in K}\bigr).
                                                                    \tag{CCR7}
  \]
  The inverse of the Hom-equalizer bijection (PC12) defines
  \[
    \gamma_{\mathcal V,\mathcal W}:=
      \operatorname{res}_{B_{\mathcal V},B_{\mathcal W}}^{-1}
        (\boldsymbol\gamma)
      \colon B_{\mathcal V}\longrightarrow B_{\mathcal W},
    \qquad
    \gamma_{\mathcal V,\mathcal W}|_{B_{\mathcal V}|_{Q_k}}
      =\widehat\gamma_k.                                          \tag{CCR8}
  \]
  For every $h\colon B_{\mathcal V}\to B_{\mathcal W}$, (PC12) and
  (CCR8) are the following equivalence chain:
  \[
  \begin{aligned}
    h\in\mathscr C_{\mathcal V,\mathcal W}
      &\Longleftrightarrow
        \operatorname{res}_{B_{\mathcal V},B_{\mathcal W}}(h)
          =\boldsymbol\gamma\\
      &\Longleftrightarrow
        h=\operatorname{res}_{B_{\mathcal V},B_{\mathcal W}}^{-1}
          (\boldsymbol\gamma)
       \Longleftrightarrow h=\gamma_{\mathcal V,\mathcal W}.
  \end{aligned}
  \]
  Hence
  \[
    \mathscr C_{\mathcal V,\mathcal W}
      =\{\gamma_{\mathcal V,\mathcal W}\}.                       \tag{CCR9}
  \]
\end{definition}

\begin{proposition}[Independence of the affine cover]
  \label{prop:nagata-global-blowup-cover-independent}
  \uses{def:nagata-global-blowup-cover-comparison}
  \lean{}

  The morphism $\gamma_{\mathcal V,\mathcal W}$ of
  \Cref{def:nagata-global-blowup-cover-comparison} is the unique isomorphism
  over $X$ restricting
  to the identity comparison on every common affine refinement.
\end{proposition}

\begin{proof}
  \uses{prop:nagata-global-blowup-affine-restrict,
    lemma:nagata-global-blowup-overlap-cocycle,
    lemma:nagata-morphism-property-calculus}
  On $Q_k$ define
  \[
  \begin{gathered}
    \delta_k:=q_{\mathcal V,k}^{-1}\circ q_{\mathcal W,k}
      \colon B_{\mathcal W}|_{Q_k}\to B_{\mathcal V}|_{Q_k},\\
    \iota_{\mathcal V,k}:=\operatorname{pr}_{B_{\mathcal V}}\colon
      B_{\mathcal V}|_{Q_k}\hookrightarrow B_{\mathcal V},\qquad
    \widehat\delta_k:=\iota_{\mathcal V,k}\circ\delta_k
      \colon B_{\mathcal W}|_{Q_k}\to B_{\mathcal V}.
  \end{gathered}
  \]
  By \Cref{lemma:nagata-global-blowup-overlap-cocycle},
  the overlap equations are
  \[
    \widehat\delta_k|_{B_{\mathcal W}|_{Q_{k\ell}}}
      =\widehat\delta_\ell|_{B_{\mathcal W}|_{Q_{k\ell}}}.
  \]
  In the notation (PC11), they are the membership statement
  \[
    \mathbf\delta:=(\widehat\delta_k)_k
      \in\mathscr G_{B_{\mathcal W}}
        \bigl(B_{\mathcal V},(B_{\mathcal W}|_{Q_k})_k\bigr).
  \]
  Define, using the inverse of (PC12),
  \[
  \begin{gathered}
    \delta:=\operatorname{res}_{B_{\mathcal W},B_{\mathcal V}}^{-1}
      (\mathbf\delta)\colon B_{\mathcal W}\to B_{\mathcal V},\\
    \delta|_{B_{\mathcal W}|_{Q_k}}=\widehat\delta_k
      \qquad(k).
  \end{gathered}
  \]
  On each $Q_k$,
  \[
    (\delta\circ\gamma_{\mathcal V,\mathcal W})
      |_{B_{\mathcal V}|_{Q_k}}
      =q_{\mathcal V,k}^{-1}q_{\mathcal W,k}
        q_{\mathcal W,k}^{-1}q_{\mathcal V,k}
      =\operatorname{id}_{B_{\mathcal V}|_{Q_k}},
  \]
  \[
    (\gamma_{\mathcal V,\mathcal W}\circ\delta)
      |_{B_{\mathcal W}|_{Q_k}}
      =q_{\mathcal W,k}^{-1}q_{\mathcal V,k}
        q_{\mathcal V,k}^{-1}q_{\mathcal W,k}
      =\operatorname{id}_{B_{\mathcal W}|_{Q_k}}.
  \]
  Injectivity of the two restriction maps in (PC12) is the pair of global
  equalities
  \[
    \delta\circ\gamma_{\mathcal V,\mathcal W}
      =\operatorname{id}_{B_{\mathcal V}},\qquad
    \gamma_{\mathcal V,\mathcal W}\circ\delta
      =\operatorname{id}_{B_{\mathcal W}},\qquad
    \operatorname{IsIso}(\gamma_{\mathcal V,\mathcal W}).        \tag{CCR10}
  \]
  Define
  \[
    r_{\mathcal V,k}:=\operatorname{pr}_{Q_k}\colon
      B_{\mathcal V}|_{Q_k}\to Q_k,\qquad
    r_{\mathcal W,k}:=\operatorname{pr}_{Q_k}\colon
      B_{\mathcal W}|_{Q_k}\to Q_k.
  \]
  The two fibre-product equalities are
  \[
    b_{\mathcal V}\circ\iota_{\mathcal V,k}
      =u_k\circ r_{\mathcal V,k},\qquad
    b_{\mathcal W}\circ\iota_{\mathcal W,k}
      =u_k\circ r_{\mathcal W,k}.                                \tag{CCR10a}
  \]
  The assertion that $q_{\mathcal V,k}$ and $q_{\mathcal W,k}$ are over
  $Q_k$ is
  \[
    r_{\mathcal W,k}\circ q_{\mathcal W,k}^{-1}\circ
      q_{\mathcal V,k}=r_{\mathcal V,k}.
  \]
  Consequently, for every $k$,
  \[
  \begin{aligned}
    b_{\mathcal W}\circ\gamma_{\mathcal V,\mathcal W}circ
      \iota_{\mathcal V,k}
      &=b_{\mathcal W}\circ\iota_{\mathcal W,k}\circ
        q_{\mathcal W,k}^{-1}\circ q_{\mathcal V,k}\\
      &=u_k\circ r_{\mathcal W,k}\circ
        q_{\mathcal W,k}^{-1}\circ q_{\mathcal V,k}\\
      &=u_k\circ r_{\mathcal V,k}
       =b_{\mathcal V}\circ\iota_{\mathcal V,k}.
  \end{aligned}
  \]
  Injectivity of (PC12) is
  \[
    b_{\mathcal W}\circ\gamma_{\mathcal V,\mathcal W}
      =b_{\mathcal V}.                                           \tag{CCR11}
  \]
  If $\gamma'\colon B_{\mathcal V}\to B_{\mathcal W}$ has the restrictions
  in (CCR8), then
  \[
    \operatorname{res}_{B_{\mathcal V},B_{\mathcal W}}(\gamma')
      =\boldsymbol\gamma
      =\operatorname{res}_{B_{\mathcal V},B_{\mathcal W}}
        (\gamma_{\mathcal V,\mathcal W})
      \Longrightarrow
      \gamma'=\gamma_{\mathcal V,\mathcal W},                   \tag{CCR12}
  \]
  where the implication is injectivity of (PC12).
\end{proof}

\begin{definition}[Comparison of the glued blow-up with relative Proj]
  \label{def:nagata-global-blowup-relative-proj-comparison}
  \uses{def:nagata-qcoh-graded-algebra-interface,
    def:nagata-relative-proj,lemma:nagata-relative-proj-interface,
    def:nagata-global-blowup-glue-data,
    def:nagata-global-blowup,def:nagata-global-blowup-map,
    def:nagata-global-blowup-overlap-basic-map}

  Let $X$ be a scheme, let
  $\mathcal I\subseteq\mathcal O_X$ be a finite-type quasi-coherent ideal,
  and fix the affine open cover
  $(v_i\colon V_i=\operatorname{Spec}(A_i)\hookrightarrow X)_{i\in I}$
  used in \Cref{def:nagata-global-blowup-glue-data}.  With
  \[
    I_i:=\Gamma(V_i,\mathcal I|_{V_i}),\qquad
    B_i:=\mathsf B_{V_i}(\mathcal I),
  \]
  define the input tuple
  \[
    \mathsf D^{\mathrm{RP}}_{X,\mathcal I}:=
      \bigl(X,\mathcal I,I,(v_i,V_i,A_i,I_i,B_i,e_i)_{i\in I},
        \operatorname{Bl}_{\mathcal I}(X)\bigr),                 \tag{RP-1}
  \]
  where
  $e_i\colon B_i\hookrightarrow\operatorname{Bl}_{\mathcal I}(X)$ is the
  open immersion of \Cref{def:nagata-global-blowup}.  All symbols in
  (RP0)--(RP10) below are functions of
  $\mathsf D^{\mathrm{RP}}_{X,\mathcal I}$.

  Define
  \[
    \mathscr R_{X,\mathcal I}:=\bigoplus_{n\ge0}\mathcal I^n,\qquad
    \mathbb P_{X,\mathcal I}:=
      \operatorname{Proj}_X(\mathscr R_{X,\mathcal I}),\qquad
    p_{X,\mathcal I}\colon\mathbb P_{X,\mathcal I}\to X.       \tag{RP0}
  \]
  Retain the affine cover
  $v_i\colon V_i=\operatorname{Spec}(A_i)\hookrightarrow X$, the ideals
  $I_i:=\Gamma(V_i,\mathcal I|_{V_i})$, the affine models
  $B_i=\operatorname{Bl}_{I_i}(V_i)$, and the chart embeddings
  $e_i\colon B_i\hookrightarrow\operatorname{Bl}_{\mathcal I}(X)$ from
  \Cref{def:nagata-global-blowup}.  For every $i$, define the coefficientwise
  graded-algebra isomorphism
  \[
  \begin{gathered}
    \eta_i:=\bigoplus_{n\ge0}\eta_{i,n}\colon
      \widetilde{\mathcal R_{A_i}(I_i)}
      \xrightarrow{\ \sim\ }v_i^*\mathscr R_{X,\mathcal I},\\
    \eta_{i,n}\colon\widetilde{I_i^n}
      \xrightarrow{\ \sim\ }\mathcal I^n|_{V_i},\qquad
    \eta_{i,n}(\widetilde a):=a|_{V_i}\quad(a\in I_i^n),
  \end{gathered}                                               \tag{RP1}
  \]
  where the equalities
  $\widetilde{I_i^n}=\mathcal I^n|_{V_i}$ are the affine
  quasi-coherence identifications.  Put
  \[
    P_i:=V_i\times_X\mathbb P_{X,\mathcal I}
      =p_{X,\mathcal I}^{-1}(V_i),\qquad
    j_i\colon P_i\hookrightarrow\mathbb P_{X,\mathcal I}.      \tag{RP2}
  \]
  Relative-Proj base change
  \cite[\href{https://stacks.math.columbia.edu/tag/01O3}{Tag 01O3}]
  {stacks-project} and (RP1) define
  \[
  \begin{gathered}
    \operatorname{bc}_i\colon
      \operatorname{Proj}_{V_i}(v_i^*\mathscr R_{X,\mathcal I})
      \xrightarrow{\ \sim\ }P_i,\\
    d_i:=
      \operatorname{bc}_i\circ
      \operatorname{Proj}_{V_i}(\eta_i^{-1})\colon
      B_i\xrightarrow{\ \sim\ }P_i,\qquad
    \widehat d_i:=j_i\circ d_i\colon B_i\to\mathbb P_{X,\mathcal I}.
  \end{gathered}                                               \tag{RP3}
  \]
  For a basic affine $W=\operatorname{Spec}(A_W)\subseteq V_i\cap V_j$,
  let $I_W:=\Gamma(W,\mathcal I|_W)$ and define
  $\widehat d_W\colon\operatorname{Bl}_{I_W}(W)\to
  \mathbb P_{X,\mathcal I}$ by the formula (RP3) with $V_i$ replaced by
  $W$.  The coefficientwise definitions (RP1), (RP3), and the restriction
  isomorphisms $r_{\nu,W}$ of
  \Cref{def:nagata-global-blowup-overlap-basic-map} are the equalities
  \[
    \widehat d_\nu|_{b_{V_\nu}^{-1}(W)}
      =\widehat d_W\circ r_{\nu,W}\qquad(\nu=i,j).              \tag{RP4}
  \]
  Since $\omega_{ij,W}=r_{j,W}^{-1}\circ r_{i,W}$, substitution in (RP4)
  is the chain
  \[
  \begin{aligned}
    \widehat d_j\circ\omega_{ij,W}
      &=\widehat d_W\circ r_{j,W}\circ
        r_{j,W}^{-1}\circ r_{i,W}\\
      &=\widehat d_W\circ r_{i,W}
       =\widehat d_i
       \quad\text{on }b_{V_i}^{-1}(W).
  \end{aligned}                                                \tag{RP5}
  \]
  Source locality on the basic cover of $B_{ij}$ is
  \[
    \widehat d_j\circ\omega_{ij}
      =\widehat d_i|_{B_{ij}}.                                 \tag{RP6}
  \]
  Define
  \[
    \mathscr D_{X,\mathcal I}:=
      \left\{d\in
        \operatorname{Hom}(\operatorname{Bl}_{\mathcal I}(X),
          \mathbb P_{X,\mathcal I})
        \ \middle|\ d\circ e_i=\widehat d_i\text{ for every }i\right\}.
  \]
  The Hom equalizer for the cover $(e_i(B_i))_i$ and (RP6) is the singleton
  equality
  \[
    \mathscr D_{X,\mathcal I}=\{\delta_{X,\mathcal I}\},\qquad
    \delta_{X,\mathcal I}\circ e_i=j_i\circ d_i.               \tag{RP7}
  \]
  Define
  \[
    \widehat e_i:=e_i\circ d_i^{-1}\colon
      P_i\longrightarrow\operatorname{Bl}_{\mathcal I}(X)
  \]
  and
  \[
    \mathscr E_{X,\mathcal I}:=
      \left\{\epsilon\in
        \operatorname{Hom}(\mathbb P_{X,\mathcal I},
          \operatorname{Bl}_{\mathcal I}(X))
        \ \middle|\ \epsilon\circ j_i=\widehat e_i
          \text{ for every }i\right\}.
  \]
  Taking inverses in (RP6) is the overlap equality for
  $(\widehat e_i)_i$; the corresponding Hom equalizer is
  \[
    \mathscr E_{X,\mathcal I}=\{\epsilon_{X,\mathcal I}\},\qquad
    \epsilon_{X,\mathcal I}\circ j_i=e_i\circ d_i^{-1}.        \tag{RP8}
  \]
  Substitution of (RP7)--(RP8) is, for every $i$, the pair of chains
  \[
  \begin{aligned}
    (\epsilon_{X,\mathcal I}\circ\delta_{X,\mathcal I})\circ e_i
      &=e_i\circ d_i^{-1}\circ d_i=e_i,\\
    (\delta_{X,\mathcal I}\circ\epsilon_{X,\mathcal I})\circ j_i
      &=j_i\circ d_i\circ d_i^{-1}=j_i.
  \end{aligned}
  \]
  Source locality on the two displayed covers is
  \[
    \epsilon_{X,\mathcal I}\circ\delta_{X,\mathcal I}
      =\operatorname{id}_{\operatorname{Bl}_{\mathcal I}(X)},\qquad
    \delta_{X,\mathcal I}\circ\epsilon_{X,\mathcal I}
      =\operatorname{id}_{\mathbb P_{X,\mathcal I}}.           \tag{RP9}
  \]
  Thus $\delta_{X,\mathcal I}$ is an isomorphism with
  $\delta_{X,\mathcal I}^{-1}:=\epsilon_{X,\mathcal I}$.  Finally, (RP3)
  is over $V_i$, so
  \[
  \begin{aligned}
    (p_{X,\mathcal I}\circ\delta_{X,\mathcal I})\circ e_i
      &=p_{X,\mathcal I}\circ j_i\circ d_i\\
      &=v_i\circ b_{V_i}
       =b_{X,\mathcal I}\circ e_i.
  \end{aligned}
  \]
  Source locality is the structure-map equality
  \[
    p_{X,\mathcal I}\circ\delta_{X,\mathcal I}
      =b_{X,\mathcal I}.                                      \tag{RP10}
  \]

  \paragraph{Mathlib relation.}
  The local pieces use exact \texttt{\detokenize{AlgebraicGeometry.Proj}}
  and \texttt{\detokenize{reesAlgebra}} declarations, and the global maps
  use \texttt{\detokenize{Scheme.GlueData}} together with
  \texttt{\detokenize{Scheme.Cover.glueMorphisms}}.  Mathlib has no
  blow-up-specific declaration matching the proposed comparison; this
  definition verifies that the two blueprint-local glued constructions agree.
\end{definition}

\subsection{Complement, properness, and principalization}

\begin{lemma}[The center ideal is the unit ideal off its support]
  \label{lemma:nagata-ideal-unit-off-support}
  \lean{AlgebraicGeometry.Scheme.IdealSheafData.ideal_eq_top_of_disjoint_support}
  \leanok

  Let
  \[
    Z:=V(\mathcal I)
      :=\{x\in X\mid\mathcal I_x\subseteq\mathfrak m_x\},
    \qquad U:=X\setminus Z,
  \]
  where $\mathfrak m_x\subseteq\mathcal O_{X,x}$ is the maximal ideal.
  For every affine open $W\subseteq X$ disjoint from $Z$, one has
  $I_W=A_W$, the unit ideal.

  \paragraph{Mathlib relation.}
  The closed set $Z$ is
  \texttt{\detokenize{AlgebraicGeometry.Scheme.IdealSheafData.support}};
  the statement mirrors
  \texttt{\detokenize{AlgebraicGeometry.Scheme.IdealSheafData.support_eq_bot_iff}}
  localized at one affine open.
\end{lemma}

\begin{proof}
  \leanok
  Inside the affine open $W=\operatorname{Spec}(A_W)$, the intersection
  $Z\cap W$ is the image of the vanishing locus
  $V(I_W)\subseteq\operatorname{Spec}(A_W)$.  Disjointness makes this
  vanishing locus empty, and an ideal with empty vanishing locus is the
  unit ideal, since otherwise it would be contained in a maximal ideal.
\end{proof}

\begin{lemma}[Complement isomorphism over an affine open with unit ideal]
  \label{lemma:nagata-affine-blowup-complement}
  \uses{def:nagata-global-blowup-map}
  \lean{AlgebraicGeometry.Scheme.IdealSheafData.isIso_blowupMap_restrict_of_ideal_eq_top}
  \leanok

  Let $W\subseteq X$ be an affine open with $I_W=A_W$.  Then the
  restriction
  \[
    b_{X,\mathcal I}|_{b_{X,\mathcal I}^{-1}(W)}\colon
      b_{X,\mathcal I}^{-1}(W)\longrightarrow W
  \]
  is an isomorphism.
\end{lemma}

\begin{proof}
  \uses{prop:nagata-global-blowup-affine-restrict,
    lemma:nagata-affine-blowup-top}
  \leanok
  By \Cref{prop:nagata-global-blowup-affine-restrict}, the restriction of
  $b_{X,\mathcal I}$ over $W$ is identified with the structure morphism
  $\beta_W$ of the local model
  $\operatorname{Bl}_{I_W}(\operatorname{Spec}A_W)$.  Since $I_W=A_W$ is
  the unit ideal, $\beta_W$ is an isomorphism by
  \Cref{lemma:nagata-affine-blowup-top}.
\end{proof}

\begin{theorem}[The blow-up is an isomorphism away from the center]
  \label{thm:nagata-global-blowup-complement}
  \lean{AlgebraicGeometry.Scheme.IdealSheafData.isIso_blowupMap_restrict,
    AlgebraicGeometry.Scheme.IdealSheafData.isIso_blowupMap_restrict_support_compl}
  \leanok

  Let $U:=X\setminus V(\mathcal I)$ and put
  $B_U:=b_{X,\mathcal I}^{-1}(U)$.  The restriction
  \[
    b_{X,\mathcal I}|_{B_U}\colon B_U\longrightarrow U
  \]
  is an isomorphism.  More generally, the restriction of
  $b_{X,\mathcal I}$ over any open $U'\subseteq X$ disjoint from
  $V(\mathcal I)$ is an isomorphism.
\end{theorem}

\begin{proof}
  \uses{lemma:nagata-ideal-unit-off-support,
    lemma:nagata-affine-blowup-complement,
    prop:nagata-global-blowup-affine-restrict}
  \leanok
  Cover the open $U'$ by affine opens $W\subseteq U'$ of $X$.  Each such
  $W$ is disjoint from $V(\mathcal I)$, so $I_W=A_W$ by
  \Cref{lemma:nagata-ideal-unit-off-support}, and the restriction of
  $b_{X,\mathcal I}$ over $W$ is an isomorphism by
  \Cref{lemma:nagata-affine-blowup-complement}.  Being an isomorphism is
  Zariski-local on the target, so
  $b_{X,\mathcal I}|_{b_{X,\mathcal I}^{-1}(U')}$ is an isomorphism.
  Taking $U'=U$ proves the displayed statement.
\end{proof}

\begin{definition}[Homogeneous polynomial map to the Rees algebra]
  \label{def:nagata-rees-polynomial-map}
  \uses{def:nagata-rees-algebra}
  \lean{AlgebraicGeometry.reesPresentation,
    AlgebraicGeometry.reesGradedPresentation}
  \leanok

  Let $A$ be a commutative ring, let $r\geq0$, and let
  $I=(f_0,\ldots,f_r)\triangleleft A$.  Put
  $P:=A[T_0,\ldots,T_r]$ with
  $\deg(T_i)=1$.  Define the $A$-algebra map
  \[
    \rho_I\colon P\longrightarrow\mathcal R_A(I)
  \]
  by
  \[
    \rho_I|_A(a)=at^0,\qquad \rho_I(T_i)=f_it.
  \]

  \paragraph{Mathlib relation.}
  The formalization indexes the generators by an arbitrary type $\sigma$ rather
  than by $\{0,\ldots,r\}$, so that $P$ is
  \texttt{\detokenize{MvPolynomial}} $\sigma$ $A$ with its total-degree grading
  \texttt{\detokenize{MvPolynomial.homogeneousSubmodule}}; taking
  $\sigma=\operatorname{Fin}(r+1)$ recovers the statement above.  The map is
  \texttt{\detokenize{MvPolynomial.aeval}} applied to the degree-one elements
  $f_it$, and \texttt{\detokenize{reesGradedPresentation}} is the same map
  packaged as a homomorphism of graded rings.
\end{definition}

\begin{lemma}[Properties of the homogeneous polynomial map]
  \label{lemma:nagata-rees-polynomial-surjection}
  \uses{def:nagata-rees-polynomial-map}
  \lean{AlgebraicGeometry.reesPresentation_mem_reesGrading,
    AlgebraicGeometry.reesPresentation_surjective,
    AlgebraicGeometry.reesGradedPresentation_surjective}
  \leanok

  Let $A,r,I=(f_0,\ldots,f_r)$ be the datum of
  \Cref{def:nagata-rees-polynomial-map}.  The map $\rho_I$ of that
  definition is graded and surjective.
\end{lemma}

\begin{proof}
  \uses{lemma:nagata-rees-degree-one-element,
    prop:nagata-rees-finite-generation}
  \leanok
  By \Cref{lemma:nagata-rees-degree-one-element},
  $\rho_I(T_i)=f_it\in\mathcal R_A(I)_1$, while
  $\rho_I(a)=at^0\in\mathcal R_A(I)_0$ for $a\in A$.  For a homogeneous
  monomial of degree $n$,
  \[
    \rho_I(T_{i_1}\cdots T_{i_n})
      =(f_{i_1}\cdots f_{i_n})t^n\in\mathcal R_A(I)_n.
  \]
  The displayed degree calculation is
  \[
    \rho_I(P_n)\subseteq\mathcal R_A(I)_n\qquad(n\geq0),
  \]
  so $\rho_I$ is graded.  Since
  \[
    I^n=(f_{i_1}\cdots f_{i_n}\mid
      0\le i_1,\ldots,i_n\le r),
  \]
  every $at^n\in\mathcal R_A(I)_n$ is in $\rho_I(P_n)$.  By
  \Cref{lemma:nagata-rees-sum-components},
  \[
    \mathcal R_A(I)=\sum_{n\ge0}\mathcal R_A(I)_n
      \subseteq\operatorname{im}(\rho_I),
  \]
  proving surjectivity.
\end{proof}

\begin{definition}[Projective-space map of an affine blow-up]
  \label{def:nagata-affine-blowup-projective-embedding}
  \uses{def:nagata-affine-blowup,def:nagata-rees-polynomial-map,
    def:nagata-relative-projective-space,
    lemma:nagata-relative-proj-interface}
  \lean{AlgebraicGeometry.ProjectiveSpace,
    AlgebraicGeometry.ProjectiveSpace.toSpec,
    AlgebraicGeometry.AffineBlowup.toProjectiveSpace}
  \leanok

  Put
  \[
    \mathbf P_A^r:=\operatorname{Proj}(A[T_0,\ldots,T_r]).
  \]
  Define
  \[
    e_I:=\operatorname{Proj}_{\operatorname{Spec}(A)}(\rho_I)
      \colon\operatorname{Bl}_I(\operatorname{Spec}A)
      \longrightarrow\mathbf P^r_A
  \]
  for the graded homomorphism $\rho_I$ of
  \Cref{def:nagata-rees-polynomial-map}; relative Proj is contravariant.

  \paragraph{Mathlib relation.}
  The source and target are instances of
  \texttt{\detokenize{AlgebraicGeometry.Proj}}.  Contravariance and the
  quotient closed immersion are supplied in this chapter by
  \Cref{lemma:nagata-relative-proj-interface}; no exact
  \texttt{\detokenize{affineBlowupToProjectiveSpace}} declaration exists.

  \formalizationnote{%
    The formalization stays in the absolute setting, where both sides are
    literally \texttt{\detokenize{AlgebraicGeometry.Proj}} and no relative
    \(\operatorname{Proj}\) is needed: the target \(\mathbf P^r_A\) is
    \texttt{\detokenize{AlgebraicGeometry.ProjectiveSpace}}, defined as
    \(\operatorname{Proj}\) of \(A[T]\) with its total-degree grading, and
    \(e_I\) is
    \texttt{\detokenize{AlgebraicGeometry.AffineBlowup.toProjectiveSpace}},
    obtained by applying \texttt{\detokenize{AlgebraicGeometry.Proj.map}} to
    \(\rho_I\).  The variables are indexed by an arbitrary type
    \(\sigma\) rather than by \(\{0,\ldots,r\}\); the finitely generated
    case of the blueprint is recovered by taking \(\sigma\) to be a finite
    generating set, as in
    \texttt{\detokenize{AffineBlowup.exists_isClosedImmersion_comp_projectiveSpaceToSpec}}.
    Mathlib does not make the total-degree grading of \(A[T]\) a global
    instance, so it is enabled locally in that file.}
\end{definition}

\begin{lemma}[Closed immersion into relative projective space]
  \label{lemma:nagata-affine-blowup-projective-embedding}
  \uses{def:nagata-affine-blowup-projective-embedding}
  \lean{AlgebraicGeometry.AffineBlowup.isClosedImmersion_toProjectiveSpace,
    AlgebraicGeometry.AffineBlowup.toProjectiveSpace_toSpec}
  \leanok

  Let $A,r,I=(f_0,\ldots,f_r)$ be the datum of
  \Cref{def:nagata-affine-blowup-projective-embedding}.
  The morphism $e_I$ of
  \Cref{def:nagata-affine-blowup-projective-embedding} is a closed immersion
  over
  $\operatorname{Spec}(A)$.
\end{lemma}

\begin{proof}
  \uses{lemma:nagata-rees-polynomial-surjection}
  \leanok
  By \Cref{lemma:nagata-rees-polynomial-surjection}, $\rho_I$ is a surjective
  graded $A$-algebra map.  The relative-Proj quotient theorem applied to
  this named map is the closed immersion
  \[
    \operatorname{Proj}(\rho_I)\colon
      \operatorname{Proj}(\mathcal R_A(I))
      \hookrightarrow\operatorname{Proj}(A[T_0,\ldots,T_r])
      =\mathbf P_A^r.
  \]
  This is $e_I$ by
  \Cref{def:nagata-affine-blowup-projective-embedding}.  In degree zero,
  \[
    \rho_I|_A=\operatorname{id}_A
  \]
  under \Cref{lemma:nagata-rees-degree-zero}.  Functoriality on degree zero
  is the equality
  \[
    (\mathbf P_A^r\to\operatorname{Spec}A)\circ e_I=b_{A,I}.
  \]

  \formalizationnote{%
    The general input used here --- that \(\operatorname{Proj}\) of a
    surjective graded ring homomorphism is a closed immersion, which is the
    absolute case of (MLRP3) of
    \Cref{lemma:nagata-relative-proj-interface} --- is not in Mathlib.  It is
    proved in this project as
    \texttt{\detokenize{AlgebraicGeometry.Proj.isClosedImmersion_map}}, by
    Zariski-locality of \texttt{\detokenize{IsClosedImmersion}} on the target
    together with the identification of the restriction of
    \texttt{\detokenize{Proj.map}} over a standard affine open \(D_+(s)\)
    with \(\operatorname{Spec}\) of the induced map on the degree-zero parts
    of the localizations, which is surjective because a surjective graded ring
    homomorphism is surjective in each degree.

    The two conclusions are formalized separately: the closed immersion is
    \texttt{\detokenize{isClosedImmersion_toProjectiveSpace}}, and the
    factorization over \(\operatorname{Spec}(A)\) is
    \texttt{\detokenize{toProjectiveSpace_toSpec}}.  The latter is the formal
    content of ``\(\rho_I|_A=\operatorname{id}_A\)'': it identifies \(A\)
    with the degree-zero part of \(A[T]\) via \(a\mapsto C(a)\) and with
    the degree-zero part of \(\mathcal R_A(I)\) by
    \Cref{lemma:nagata-rees-degree-zero}, and checks that \(\rho_I\)
    intertwines the two.}
\end{proof}


\begin{proposition}[Projectivity of the affine blow-up]
  \label{prop:nagata-affine-blowup-projective}
  % Not formalized: this Mathlib has no morphism class `IsProjective' for
  % schemes, so ``is projective'' cannot be stated.  See the formalization
  % note after the proof.
  \lean{}

  Let $A$ be a commutative ring and let $I\triangleleft A$ be finitely
  generated.  The affine blow-up morphism
  \[
    \operatorname{Bl}_I(\operatorname{Spec}(A))
      \longrightarrow\operatorname{Spec}(A)
  \]
  is projective.
\end{proposition}

\begin{proof}
  \uses{lemma:nagata-affine-blowup-projective-embedding,
    def:nagata-morphism-predicate-notation}
  Choose $f_0,\ldots,f_r\in I$ with $I=(f_0,\ldots,f_r)$.  By
  \Cref{lemma:nagata-affine-blowup-projective-embedding}, $e_I$ is a closed
  immersion and
  \[
    b_{A,I}
      =\pi_A\circ e_I,\qquad
    \operatorname{Bl}_I(\operatorname{Spec}A)
      \xrightarrow{e_I}\mathbf P_A^r
      \xrightarrow{\pi_A}\operatorname{Spec}A.
  \]
  Define
  \[
    \mathcal E_A:=\mathcal O_{\operatorname{Spec}A}^{\oplus(r+1)},
    \qquad
    \mathbf P_{\operatorname{Spec}A}(\mathcal E_A)=\mathbf P_A^r,
    \qquad
    \operatorname{pr}_{\operatorname{Spec}A,\mathcal E_A}=\pi_A.
  \]
  The witness tuple required by (P0) is
  \[
    \mathcal E_A\in\operatorname{QCoh}(\mathcal O_{\operatorname{Spec}A}),
    \qquad
    \mathcal E_A\text{ is of finite type},
    \qquad
    \operatorname{IsClosedImmersion}(e_I),
    \qquad
    \operatorname{pr}_{\operatorname{Spec}A,\mathcal E_A}\circ e_I
      =\pi_A\circ e_I=b_{A,I}.
  \]
  Substitution of $(\mathcal E,\iota)=(\mathcal E_A,e_I)$ in (P0) is
  \[
    \operatorname{IsProjective}(b_{A,I}).
  \]

  \formalizationnote{%
    This proposition is not formalized, because this Mathlib has no morphism
    class \texttt{\detokenize{IsProjective}} for schemes --- there is no
    \texttt{\detokenize{Mathlib/AlgebraicGeometry/Morphisms/Projective.lean}}
    --- so its conclusion cannot be stated.  Everything the proof supplies
    before the final substitution \emph{is} formalized, as
    \texttt{\detokenize{AffineBlowup.exists_isClosedImmersion_comp_projectiveSpaceToSpec}}:
    for finitely generated \(I\) there is a finite index set \(\sigma\) and
    a morphism \(e_I\colon\operatorname{Bl}_I(\operatorname{Spec}A)\to
    \mathbf P^\sigma_A\) which is a closed immersion and satisfies
    \(\pi_A\circ e_I=b_{A,I}\).  Only the wrapping of that data into the
    predicate (P0) is missing.

    Nothing downstream is blocked: \Cref{lemma:nagata-affine-blowup-proper} is
    formalized without going through projectivity, as its own formalization
    note records.}
\end{proof}

\begin{lemma}[Properness of the affine blow-up]
  \label{lemma:nagata-affine-blowup-proper}
  \lean{AlgebraicGeometry.AffineBlowup.isProper_toSpec}
  \leanok

  If \(I\) is finitely generated, the affine blow-up morphism is proper.

  \paragraph{Mathlib relation.}
  The terminal properness input for the ambient absolute projective space is
  the exact instance
  \texttt{\detokenize{Proj.instIsProperToSpecZeroOfFiniteTypeSubtypeMemOfNatNat}}.
  The blow-up specialization remains a project lemma.
\end{lemma}

\begin{proof}
  \uses{prop:nagata-affine-blowup-projective}
  \leanok
  By \Cref{prop:nagata-affine-blowup-projective}, $b_{A,I}$ factors as
  \[
    \operatorname{Bl}_I(\operatorname{Spec}A)
      \xrightarrow{e_I}\mathbf P_A^r
      \xrightarrow{\pi_A}\operatorname{Spec}A.
  \]
  The theorem that projective morphisms are proper makes this composite
  proper.  The formalized proof takes the shorter route described in the note
  below, so the dependency on
  \Cref{prop:nagata-affine-blowup-projective} recorded above is that of this
  written argument only.
\end{proof}

\formalizationnote{%
  The Lean proof of \Cref{lemma:nagata-affine-blowup-proper} does not use the
  projective embedding, and therefore does not depend on
  \Cref{prop:nagata-affine-blowup-projective},
  \Cref{lemma:nagata-affine-blowup-projective-embedding} or
  \Cref{lemma:nagata-rees-polynomial-surjection}.  Instead it applies the
  Mathlib instance named in the ``Mathlib relation'' paragraph above directly
  to $\mathcal R_A(I)$: writing
  $b_{A,I}
    =\bigl(\operatorname{Proj}\mathcal R_A(I)
      \to\operatorname{Spec}\mathcal R_A(I)_0\bigr)
     \circ\operatorname{Spec}(\varepsilon)^{-1}$,
  where $\varepsilon\colon A\xrightarrow{\ \sim\ }\mathcal R_A(I)_0$ is the
  isomorphism of \Cref{lemma:nagata-rees-degree-zero}, properness of the first
  factor holds as soon as $\mathcal R_A(I)$ is of finite type over
  $\mathcal R_A(I)_0$, which follows from
  \Cref{prop:nagata-rees-finite-generation} transported along $\varepsilon$,
  and the second factor is an isomorphism.  The projective embedding is still
  wanted, for the strictly stronger
  \Cref{prop:nagata-affine-blowup-projective}.}

\begin{lemma}[Local properness of the global blow-up]
  \label{lemma:nagata-global-blowup-proper-local}
  \lean{AlgebraicGeometry.Scheme.IdealSheafData.isProper_blowupMap_restrict}
  \leanok

  If $\mathcal I$ is of finite type, the restriction of
  $b_{X,\mathcal I}$ over every affine open of $X$ is proper.
\end{lemma}

\begin{proof}
  \uses{prop:nagata-global-blowup-affine-restrict,lemma:nagata-affine-blowup-proper}
  \leanok
  Let $V=\operatorname{Spec}(A)\subseteq X$ be affine and put
  $I_V:=\Gamma(V,\mathcal I|_V)$.  Since $\mathcal I$ is of finite type,
  choose $f_0,\ldots,f_r\in I_V$ with $I_V=(f_0,\ldots,f_r)$.  By
  \Cref{prop:nagata-global-blowup-affine-restrict},
  \[
    q_V\colon b_{X,\mathcal I}^{-1}(V)
      \xrightarrow{\sim}\operatorname{Bl}_{I_V}(V),
    \qquad
    b_{X,\mathcal I}|_{b^{-1}(V)}
      =b_{A,I_V}\circ q_V.
  \]
  The morphism $b_{A,I_V}$ is proper by
  \Cref{lemma:nagata-affine-blowup-proper}.  Composition with the isomorphism
  $q_V$ is proper, so the displayed restriction is proper.
\end{proof}

\begin{theorem}[Projectivity of the global blow-up]
  \label{thm:nagata-global-blowup-projective}
  \lean{}

  If \(\mathcal I\) is of finite type, then
  \(b_{X,\mathcal I}\colon\operatorname{Bl}_{\mathcal I}(X)\to X\) is
  projective.
\end{theorem}

\begin{proof}
  \uses{def:nagata-global-blowup-relative-proj-comparison,
    lemma:nagata-morphism-property-calculus}
  Use the symbols of
  \Cref{def:nagata-global-blowup-relative-proj-comparison} and define
  \[
    \mu:=\bigoplus_{n\geq0}\mu_n\colon
      \operatorname{Sym}_{\mathcal O_X}^{\bullet}(\mathcal I)
        \longrightarrow
      \mathscr R_{X,\mathcal I}=\bigoplus_{n\geq0}\mathcal I^n,
    \qquad
    \mu_n(s_1\cdots s_n):=s_1\cdots s_n.                       \tag{GP1}
  \]
  By the definition of the sheaf power $\mathcal I^n$,
  \[
    \operatorname{im}(\mu_n)=\mathcal I^n
      =(\mathscr R_{X,\mathcal I})_n\qquad(n\geq0).             \tag{GP2}
  \]
  Thus $\mu$ is surjective, and (GP2) is exactly the assertion that
  $\mathscr R_{X,\mathcal I}$ is generated by
  $(\mathscr R_{X,\mathcal I})_1=\mathcal I$ over
  $(\mathscr R_{X,\mathcal I})_0=\mathcal O_X$.  Since $\mathcal I$ is of
  finite type,
  \cite[\href{https://stacks.math.columbia.edu/tag/07ZS}{Tag 07ZS}]
  {stacks-project}, applied with
  $\mathcal A=\mathscr R_{X,\mathcal I}$, states
  \[
    \operatorname{IsProjective}(p_{X,\mathcal I}).              \tag{GP3}
  \]
  Equations (RP9)--(RP10) of
  \Cref{def:nagata-global-blowup-relative-proj-comparison} are
  \[
    \operatorname{IsIso}(\delta_{X,\mathcal I}),
    \qquad
    b_{X,\mathcal I}=p_{X,\mathcal I}\circ\delta_{X,\mathcal I}
      =p_{X,\mathcal I}^{\delta_{X,\mathcal I},\operatorname{id}_X}.
                                                                    \tag{GP4}
  \]
  Substitution of (GP3)--(GP4) into (PC8) of
  \Cref{lemma:nagata-morphism-property-calculus} is the equivalence
  \[
    \operatorname{IsProjective}(b_{X,\mathcal I})
      \Longleftrightarrow
    \operatorname{IsProjective}(p_{X,\mathcal I}),
  \]
  whose right-hand side is (GP3).
\end{proof}

\begin{theorem}[Properness of the global blow-up]
  \label{thm:nagata-global-blowup-proper}
  \lean{AlgebraicGeometry.Scheme.IdealSheafData.isProper_blowupMap}
  \leanok

  If \(\mathcal I\) is of finite type, then
  \(b_{X,\mathcal I}\) is proper.
\end{theorem}

\begin{proof}
  \uses{thm:nagata-global-blowup-projective}
  \leanok
  By \Cref{thm:nagata-global-blowup-projective},
  \[
    b_{X,\mathcal I}\colon\operatorname{Bl}_{\mathcal I}(X)\to X
  \]
  is projective.  The theorem that projective morphisms are proper applied to
  this displayed morphism proves that it is proper.
\end{proof}

\formalizationnote{%
  The Lean proof of \Cref{thm:nagata-global-blowup-proper} does not go through
  projectivity, and therefore depends on neither
  \Cref{thm:nagata-global-blowup-projective} nor
  \Cref{def:nagata-global-blowup-relative-proj-comparison}.  Instead it uses
  that properness is Zariski-local on the target, which reduces the theorem to
  \Cref{lemma:nagata-global-blowup-proper-local} over the affine opens of $X$;
  that lemma is in turn \Cref{prop:nagata-global-blowup-affine-restrict}
  combined with \Cref{lemma:nagata-affine-blowup-proper}, exactly as its proof
  above describes.  \Cref{thm:nagata-global-blowup-projective} is still wanted,
  as the strictly stronger statement; the \texttt{\textbackslash uses} list of
  the proof printed above records the informal route, not the formalized one.

  The finite type condition on $\mathcal I$ is spelled out in Lean as
  ``$\Gamma(V,\mathcal I|_V)$ is a finitely generated ideal for every affine
  open $V\subseteq X$'', since this Mathlib has no predicate for a finite type
  quasi-coherent ideal sheaf.  On a locally Noetherian scheme the hypothesis is
  automatic, and the resulting unconditional statement is registered as an
  instance.}

\begin{lemma}[Principalization on one affine chart]
  \label{lemma:nagata-chart-extended-ideal-principal}
  \uses{def:nagata-blowup-chart-ring}
  \lean{AlgebraicGeometry.AffineBlowup.map_awayAlgebraMap_reesX_eq_span_singleton,
    AlgebraicGeometry.AffineBlowup.map_awayAlgebraMap_eq_span_singleton}
  \leanok

  Let $A$ be a commutative ring, let $I\triangleleft A$, and let $a\in I$.
  In $A[I/a]$, the extended ideal satisfies
  \[
    I\,A[I/a]=(a).
  \]
\end{lemma}

\begin{proof}
  \leanok
  Let $B:=A[I/a]$.  For every $x\in I$,
  \[
    x=a(x/a)\in aB,
  \]
  so $IB\subseteq aB$.  Since $a\in I$, one has $aB\subseteq IB$.
  The two inclusions combine to the equality
  \[
    IA[I/a]=IB=aB=(a).
  \]
\end{proof}

\begin{lemma}[Regularity of the chart generator]
  \label{lemma:nagata-chart-generator-regular}
  \uses{def:nagata-blowup-chart-ring}
  \lean{AlgebraicGeometry.AffineBlowup.isRegular_awayAlgebraMap_reesX,
    AlgebraicGeometry.AffineBlowup.isRegular_awayAlgebraMap}
  \leanok

  Let $A$ be a commutative ring, let $I\triangleleft A$, and let $a\in I$.
  Multiplication by $a$ on $A[I/a]$ is injective.
\end{lemma}

\begin{proof}
  \leanok
  The formalization argues directly in $\mathcal R_A(I)$ rather than through
  $\jmath$: if $z=F/(at)^n$ with $F$ homogeneous and $az=0$, then
  $(at)^k\,aF=0$ in $\mathcal R_A(I)$ for some $k$, hence
  $(at)^{k+1}F=t\,(at)^k\,aF=0$ in $A[t]$ and therefore $z=0$.  The argument of
  the blueprint is the following.

  By \Cref{def:nagata-blowup-chart-ring}, the inclusion
  $\jmath\colon A[I/a]\hookrightarrow A_a$ is injective.  If
  $z\in A[I/a]$ satisfies $az=0$, then
  \[
    \lambda_a(a)\jmath(z)=\jmath(az)=0.
  \]
  Since $\lambda_a(a)\in A_a^\times$, multiplication by
  $\lambda_a(a)^{-1}$ is the equality $\jmath(z)=0$.  Injectivity of
  $\jmath$ is then the equality
  \[
    z=0.
  \]
\end{proof}

\begin{lemma}[Transition between local exceptional generators]
  \label{lemma:nagata-chart-generators-associated}
  \uses{def:nagata-affine-blowup-chart-map}
  \lean{AlgebraicGeometry.AffineBlowup.exists_units_awayAlgebraMap_eq_mul}
  \leanok

  Let $A$ be a commutative ring, let $I\triangleleft A$, and let $a,b\in I$.
  Put
  \[
  \begin{gathered}
    R:=\mathcal R_A(I),\qquad
    R_a:=R[(at)^{-1}],\qquad R_b:=R[(bt)^{-1}],\\
    R_{ab}:=R[(at)^{-1},(bt)^{-1}],\qquad B_{ab}:=(R_{ab})_0.
  \end{gathered}                                                     \tag{LCE0}
  \]
  Define the graded localization maps and their degree-zero restrictions by
  \[
  \begin{gathered}
    \iota_{a,ab}\colon R_a\longrightarrow R_{ab},\qquad
      \iota_{a,ab}\!\left(\frac z{(at)^m}\right)
        :=\frac z{(at)^m},\\
    \iota_{b,ab}\colon R_b\longrightarrow R_{ab},\qquad
      \iota_{b,ab}\!\left(\frac z{(bt)^m}\right)
        :=\frac z{(bt)^m},\\
    (\iota_{a,ab})_0\colon R_{(at)}=(R_a)_0\longrightarrow B_{ab},
    \qquad
    (\iota_{b,ab})_0\colon R_{(bt)}=(R_b)_0\longrightarrow B_{ab}.
  \end{gathered}                                                     \tag{LCE1}
  \]
  With $\eta_a$ and $\eta_b$ defined by
  \Cref{def:nagata-affine-blowup-chart-map}, define
  \[
  \begin{gathered}
    \lambda_{a,ab}:=(\iota_{a,ab})_0\circ\eta_a
      \colon A[I/a]\longrightarrow B_{ab},\\
    \lambda_{b,ab}:=(\iota_{b,ab})_0\circ\eta_b
      \colon A[I/b]\longrightarrow B_{ab},\\
    \lambda_{a,ab}\!\left(\frac{x}{a^n}\right)
      =\frac{xt^n}{(at)^n}\quad(x\in I^n),\qquad
    \lambda_{b,ab}\!\left(\frac{y}{b^n}\right)
      =\frac{yt^n}{(bt)^n}\quad(y\in I^n).
  \end{gathered}                                                     \tag{LCE2}
  \]
  Thus
  \[
    \operatorname{Spec}(B_{ab})
      =\operatorname{Spec}(A[I/a])\cap
        \operatorname{Spec}(A[I/b])
  \]
  inside $\operatorname{Bl}_I(\operatorname{Spec}A)$.  In $B_{ab}$, define
  \[
    u_{ba}:=\lambda_{a,ab}(b/a),\qquad
    u_{ab}:=\lambda_{b,ab}(a/b).
  \]
  The images of $a$ and $b$ satisfy
  \[
    b=a\,u_{ba},\qquad a=b\,u_{ab},\qquad
    u_{ab}u_{ba}=u_{ba}u_{ab}=1,
  \]
  in the explicitly defined ring $B_{ab}$.

  \paragraph{Formalization.}
  The Lean statement is the last display: it produces a unit $u_{ba}$ of
  $B_{ab}=\texttt{\detokenize{HomogeneousLocalization.Away (reesGrading I) (reesX I a ha * reesX I b hb)}}$
  with $b=a\,u_{ba}$.  The identification of $\operatorname{Spec}(B_{ab})$ with
  the overlap of the two charts is the one carried by Mathlib's
  \texttt{\detokenize{Proj}} API and is not restated there.
\end{lemma}

\begin{proof}
  \uses{lemma:nagata-chart-extended-ideal-principal,
    lemma:nagata-chart-generator-regular}
  \leanok
  By \Cref{lemma:nagata-chart-extended-ideal-principal}, restriction to
  $B_{ab}$ is the ideal equality
  \[
    (a)B_{ab}=IB_{ab}=(b)B_{ab}.
  \]
  In the homogeneous localizations defining the overlap,
  \[
    u_{ba}:=\frac{b}{a},\qquad u_{ab}:=\frac{a}{b}
  \]
  belong to $B_{ab}$, and
  \[
    b=au_{ba},\qquad a=bu_{ab},\qquad
    u_{ab}u_{ba}=\frac ab\frac ba=1.
  \]
  The equality $u_{ab}u_{ba}=1$ states that
  \[
    u_{ba}\in B_{ab}^\times,
    \qquad u_{ba}^{-1}=u_{ab}.
  \]
\end{proof}

\begin{definition}[Invertible ideal and effective Cartier center]
  \label{def:nagata-invertible-ideal-effective-cartier}
  \uses{def:nagata-ideal-sheaf-data-adapter,
    def:nagata-quotient-relative-spec}
  \lean{AlgebraicGeometry.Scheme.IdealSheafData.IsInvertible,
    AlgebraicGeometry.IsEffectiveCartierCenter}
  \leanok

  Let $X$ be a scheme and let
  $\mathcal J\subseteq\mathcal O_X$ be a quasi-coherent ideal.  Define
  \[
  \begin{aligned}
    \operatorname{InvertibleIdeal}_X(\mathcal J)
      :\Longleftrightarrow{}&
      \text{there is an affine open cover }X=\bigcup_aU_a,\quad
      U_a=\operatorname{Spec}(A_a),\\
      &\text{and elements }s_a\in A_a\text{ such that}\quad
      \mathcal J|_{U_a}=\widetilde{(s_a)},\qquad
      s_a\in\operatorname{nonZeroDivisors}(A_a).
  \end{aligned}                                                   \tag{MLCI0}
  \]
  Thus the word ``invertible'' in this chapter includes the displayed
  regularity of the inclusion into $\mathcal O_X$.

  If $z\colon Z\hookrightarrow X$ is a closed immersion, define
  \(\operatorname{EffectiveCartierCenter}_X(z)\) to mean that there are a
  quasi-coherent ideal $\mathcal J$, an isomorphism
  \[
    \alpha\colon
      Z\xrightarrow{\sim}
      \operatorname{Spec}_X(\mathcal O_X/\mathcal J),
    \qquad
    i_{\mathsf D_X(\mathcal J)}\circ\alpha=z,                    \tag{MLCI1}
  \]
  and a proof of
  $\operatorname{InvertibleIdeal}_X(\mathcal J)$.

  \paragraph{Mathlib relation.}
  On an affine chart, regularity in (MLCI0) is the exact predicate
  \texttt{\detokenize{IsRegular}} (equivalently, membership in
  \texttt{\detokenize{nonZeroDivisors}}).  Mathlib also has
  \texttt{\detokenize{Module.Invertible}} and the sheaf predicate
  \texttt{\detokenize{SheafOfModules.IsLocallyFree}}.  The assembled
  scheme-level predicates (MLCI0)--(MLCI1), and their connection with the
  quotient subscheme, are blueprint-local adapters rather than existing
  Mathlib declarations.

  \paragraph{Formalization.}
  Both predicates are formalized.  (MLCI0) is
  \texttt{\detokenize{AlgebraicGeometry.Scheme.IdealSheafData.IsInvertible}}.  It is
  phrased pointwise -- every point has an affine open neighbourhood on which the ideal
  is generated by a single non-zerodivisor -- which is equivalent to the displayed
  cover formulation and is local by construction.  The affine-local shape
  ``generated by a single non-zerodivisor'' is
  \texttt{\detokenize{Ideal.IsRegularPrincipal}}.  (MLCI1) is
  \texttt{\detokenize{AlgebraicGeometry.IsEffectiveCartierCenter}}; the
  quotient/relative-$\operatorname{Spec}$ adapter of
  \Cref{def:nagata-quotient-relative-spec} that it needs is Mathlib's
  \texttt{\detokenize{Scheme.IdealSheafData.subscheme}}, so
  $\operatorname{Spec}_X(\mathcal O_X/\mathcal J)$ is available directly and no
  further adapter is required.  The predicate is stated for an arbitrary morphism $z$
  rather than assuming $z$ is a closed immersion, since (MLCI1) forces that; see
  \texttt{\detokenize{IsEffectiveCartierCenter.isClosedImmersion}}.
\end{definition}

\begin{lemma}[Interface for invertible ideals and Cartier centers]
  \label{lemma:nagata-invertible-ideal-effective-cartier}
  \uses{def:nagata-invertible-ideal-effective-cartier,
    lemma:nagata-ideal-sheaf-data-adapter,
    lemma:nagata-quotient-relative-spec}
  \lean{AlgebraicGeometry.Scheme.IdealSheafData.IsInvertible.of_iSup_opensRange_eq_top,
    AlgebraicGeometry.Scheme.IdealSheafData.IsInvertible.comap_of_isOpenImmersion,
    AlgebraicGeometry.Scheme.IdealSheafData.IsInvertible.of_isAffine,
    AlgebraicGeometry.Scheme.IdealSheafData.IsInvertible.comap_of_flat,
    AlgebraicGeometry.isEffectiveCartierCenter_iff_isInvertible_ker}

  Predicate (MLCI0) is independent of the chosen affine cover and is local
  on $X$.  It is equivalent to saying that the associated ideal-sheaf datum
  is affine-locally principal with a generator in the non-zero-divisors, and
  it implies that the underlying $\mathcal O_X$-module is locally free of
  rank one.  It is preserved by restriction to an open subscheme,
  isomorphism, and flat pullback.  Under the correspondence (MLRS0), (MLCI1)
  is equivalent to the usual affine-local assertion that the defining ideal
  of $Z$ is generated by one non-zero-divisor.

  \paragraph{Formalization.}
  Every assertion except local freeness of rank one is formalized.  Locality and cover
  independence are \texttt{\detokenize{IsInvertible.of_iSup_opensRange_eq_top}} (the
  predicate may be checked after pulling back along any family of open immersions whose
  ranges cover $X$) together with \texttt{\detokenize{IsInvertible.of_isAffine}} (on an
  affine scheme it suffices to check the global sections).  Flat pullback is
  \texttt{\detokenize{IsInvertible.comap_of_flat}}: on an affine chart the extended
  ideal is computed by \texttt{\detokenize{IdealSheafData.ideal_comap_appLE}}, the
  induced ring map is flat, and a flat ring map carries non-zerodivisors to
  non-zerodivisors because a flat module is torsion-free.  Stability under restriction
  to an open subscheme and under isomorphism,
  \texttt{\detokenize{IsInvertible.comap_of_isOpenImmersion}}, is the special case of
  that for an open immersion, which is flat.  The final assertion is
  \texttt{\detokenize{isEffectiveCartierCenter_iff_isInvertible_ker}}: a closed
  immersion $z$ is an effective Cartier center if and only if
  $\operatorname{InvertibleIdeal}_X(z.\mathrm{ker})$, and unfolding (MLCI0) that is
  exactly the affine-local assertion that the defining ideal of $Z$ is generated by one
  non-zero-divisor.  Local freeness of rank one is not formalized: it needs the
  quasi-coherent $\mathcal O_X$-module attached to an ideal-sheaf datum, which Mathlib
  does not provide (there is no bridge from
  \texttt{\detokenize{Scheme.IdealSheafData}} to
  \texttt{\detokenize{SheafOfModules}}), and it is not used anywhere downstream.
\end{lemma}

\begin{proof}
  Refine two affine covers by principal affine opens.  Localization preserves
  principal ideals and non-zero-divisors, which proves cover independence
  and locality.  Multiplication by a regular generator identifies the ideal
  with a free module of rank one.  Restriction and isomorphism preserve this
  calculation; for a flat pullback, tensoring the injective multiplication
  map by the target ring remains injective.  The final assertion is the
  affine description (MLRS2) of the quotient subscheme.
\end{proof}

\begin{theorem}[Principalization of the center ideal]
  \label{thm:nagata-blowup-principalizes-center}
  \uses{def:nagata-invertible-ideal-effective-cartier,
    lemma:nagata-invertible-ideal-effective-cartier}
  \lean{AlgebraicGeometry.Scheme.IdealSheafData.isInvertible_comap_blowupMap,
    AlgebraicGeometry.Scheme.IdealSheafData.isEffectiveCartierCenter_subschemeι_comap_blowupMap}
  \leanok

  Assume that $\mathcal I\subseteq\mathcal O_X$ is a finite-type
  quasi-coherent ideal.  On $\operatorname{Bl}_{\mathcal I}(X)$, the extended ideal
  $\mathcal I\mathcal O_{\operatorname{Bl}_{\mathcal I}(X)}$ is locally
  principal and generated locally by a non-zero-divisor; equivalently, it is
  an
  invertible ideal defining an effective Cartier divisor.
  More explicitly, if
  \[
    \mathcal J:=
      \mathcal I\mathcal O_{\operatorname{Bl}_{\mathcal I}(X)},
    \qquad J:=\mathsf D_{\operatorname{Bl}_{\mathcal I}(X)}(\mathcal J),
  \]
  then
  \[
    \operatorname{InvertibleIdeal}_{\operatorname{Bl}_{\mathcal I}(X)}
      (\mathcal J),
    \qquad
    \operatorname{EffectiveCartierCenter}_{\operatorname{Bl}_{\mathcal I}(X)}
      (i_J).
  \]

  \paragraph{Mathlib relation.}
  The conclusion is expressed through (MLCI0)--(MLCI1).  Mathlib supplies
  the affine regular-element and module interfaces cited there, but it does
  not currently contain this packaged blow-up principalization theorem.

  \paragraph{Formalization.}
  Both halves are formalized.  The (MLCI0) half,
  $\operatorname{InvertibleIdeal}_{\operatorname{Bl}_{\mathcal I}(X)}(\mathcal J)$ for
  $\mathcal J=\mathcal I\mathcal O_{\operatorname{Bl}_{\mathcal I}(X)}$, is
  \texttt{\detokenize{isInvertible_comap_blowupMap}}, where the extended ideal is
  Mathlib's \texttt{\detokenize{Scheme.IdealSheafData.comap}} along the blow-up map.
  The (MLCI1) half is
  \texttt{\detokenize{isEffectiveCartierCenter_subschemeι_comap_blowupMap}}, an
  immediate consequence of the first through
  \texttt{\detokenize{IsInvertible.isEffectiveCartierCenter}}; the closed immersion
  $i_J$ is Mathlib's \texttt{\detokenize{Scheme.IdealSheafData.subschemeι}} for
  $\mathcal J$.  No finiteness assumption on $\mathcal I$ is needed for either half:
  finite type only serves to keep the chart cover finite, which is irrelevant for a
  local statement.
\end{theorem}

\begin{proof}
  \uses{prop:nagata-affine-blowup-chart-cover,
    lemma:nagata-chart-extended-ideal-principal,
    lemma:nagata-chart-generator-regular}
  \leanok
  Choose an affine cover
  \[
    X=\bigcup_iV_i,\qquad V_i=\operatorname{Spec}(A_i).
  \]
  By the definition of a finite-type quasi-coherent ideal, choose symbols
  \[
    r_i\in\mathbf Z_{\ge0},\qquad
    a_{i,0},\ldots,a_{i,r_i}\in A_i
  \]
  satisfying
  \[
    I_i:=\Gamma(V_i,\mathcal I|_{V_i})
      =(a_{i,0},\ldots,a_{i,r_i}).
  \]
  Define the affine-restriction isomorphism and its standard chart opens by
  \[
  \begin{gathered}
    q_i\colon b_{X,\mathcal I}^{-1}(V_i)
      \xrightarrow{\ \sim\ }\operatorname{Bl}_{I_i}(V_i)
      \quad\text{as in \Cref{prop:nagata-global-blowup-affine-restrict}},\\
    D_{i,\nu}:=D_+(a_{i,\nu}\mathbf u)
      \subseteq\operatorname{Bl}_{I_i}(V_i),\qquad
    C_{i,\nu}:=q_i^{-1}(D_{i,\nu})
      \subseteq\operatorname{Bl}_{\mathcal I}(X),\\
    \kappa_{i,\nu}:=
      \chi_{a_{i,\nu}}\circ q_i|_{C_{i,\nu}}\colon
      C_{i,\nu}\xrightarrow{\ \sim\ }
      \operatorname{Spec}(A_i[I_i/a_{i,\nu}]),
  \end{gathered}
  \]
  where $\chi_{a_{i,\nu}}$ is
  \Cref{prop:nagata-affine-blowup-chart}.  By
  \Cref{prop:nagata-affine-blowup-chart-cover}, the typed affine opens
  \[
    (C_{i,\nu})_{i,\ 0\leq\nu\leq r_i}
  \]
  cover $\operatorname{Bl}_{\mathcal I}(X)$.  Put
  $\mathcal L:=\mathcal I\mathcal O_{\operatorname{Bl}_{\mathcal I}(X)}$.
  By \Cref{def:nagata-direct-image-structure-map}, define the affine
  global-section ring map and the section
  \[
  \begin{gathered}
    \kappa_{i,\nu}^{\#}
      :=\varepsilon_{\kappa_{i,\nu},
        |\operatorname{Spec}(A_i[I_i/a_{i,\nu}])|}\colon
      A_i[I_i/a_{i,\nu}]\longrightarrow
        \Gamma(C_{i,\nu},\mathcal O_{C_{i,\nu}}),\\
    \overline a_{i,\nu}:=
      \kappa_{i,\nu}^{\#}(a_{i,\nu})
      \in\Gamma(C_{i,\nu},\mathcal O_{C_{i,\nu}}).
  \end{gathered}
  \]
  On $C_{i,\nu}$,
  \[
    \mathcal L|_{C_{i,\nu}}=(\overline a_{i,\nu})
  \]
  by \Cref{lemma:nagata-chart-extended-ideal-principal}, and multiplication
  by $\overline a_{i,\nu}$ is injective by
  \Cref{lemma:nagata-chart-generator-regular}.  These two assertions define
  the module isomorphism
  \[
  \begin{aligned}
    m_{i,\nu}\colon\mathcal O_{C_{i,\nu}}
      &\xrightarrow{\ \sim\ }\mathcal L|_{C_{i,\nu}},&
    m_{i,\nu}(r)&:=r\overline a_{i,\nu},\\
    \operatorname{Ann}_{\mathcal O_{C_{i,\nu}}}
      (\overline a_{i,\nu})&=0,&
    \mathcal L|_{C_{i,\nu}}&=
      \overline a_{i,\nu}\mathcal O_{C_{i,\nu}}.
  \end{aligned}                                                \tag{PR0}
  \]
  The cover $(C_{i,\nu})_{i,\nu}$ and the isomorphisms $m_{i,\nu}$ in
  (PR0) are the local-triviality condition
  \[
    \mathcal L\ \text{is an invertible ideal}.
  \]
  The two equalities in (PR0) are, on every member of the same cover, the
  defining conditions for the ideal of an effective Cartier divisor:
  \[
    \mathcal L|_{C_{i,\nu}}
      =(\overline a_{i,\nu}),\qquad
    \operatorname{Ann}(\overline a_{i,\nu})=0.
  \]
\end{proof}

\subsection{Universal property, flat base change, and products of ideals}

\begin{definition}[Graded map induced by an invertible quotient]
  \label{def:nagata-blowup-universal-graded-map}
  \uses{def:nagata-qcoh-graded-algebra-interface,
    def:nagata-invertible-ideal-effective-cartier,
    lemma:nagata-qcoh-graded-algebra-interface}
  \lean{}

  Let $\mathcal I\subseteq\mathcal O_X$ be a finite-type quasi-coherent ideal,
  let $f\colon Y\to X$, and let
  \[
    \mu_f\colon f^*\mathcal I\longrightarrow\mathcal O_Y
  \]
  be the pullback of the inclusion
  $\mathcal I\hookrightarrow\mathcal O_X$.  Put
  \[
    \mathcal L_f:=\operatorname{im}(\mu_f)\subseteq\mathcal O_Y
  \]
  and assume that $\mathcal L_f$ is invertible.  Denote the image
  factorization by
  \[
    f^*\mathcal I\xrightarrow{q_f}\mathcal L_f
      \xrightarrow{\iota_f}\mathcal O_Y,
    \qquad \mu_f=\iota_f\circ q_f,
  \]
  where $q_f$ is surjective.  Define
  \[
    \mathcal R_X(\mathcal I)
      =\bigoplus_{n\ge0}\mathcal I^n\mathbf u^n
      \quad\text{as in (MLQC3)},
    \qquad
    \mathcal S_Y(\mathcal L_f):=
      \bigoplus_{n\ge0}\mathcal L_f^{\otimes n}.
  \]
  For $n\ge0$, define the multiplication epimorphism and its pullback by
  \[
    m_{\mathcal I,n}\colon\mathcal I^{\otimes n}
      \twoheadrightarrow\mathcal I^n,
    \qquad
    m^f_{\mathcal I,n}:=f^*(m_{\mathcal I,n})\colon
      (f^*\mathcal I)^{\otimes n}\twoheadrightarrow f^*(\mathcal I^n).
  \]
  Since
  \[
    q_f^{\otimes n}\bigl(\ker(m^f_{\mathcal I,n})\bigr)=0,
  \]
  define the unique map
  \[
    \alpha_{f,n}\colon f^*(\mathcal I^n)
      \longrightarrow\mathcal L_f^{\otimes n},
    \qquad
    \alpha_{f,n}\circ m^f_{\mathcal I,n}=q_f^{\otimes n}.
  \]
  Define
  \[
    \alpha_f:=\bigoplus_{n\ge0}\alpha_{f,n}\colon
      f^*\mathcal R_X(\mathcal I)
      \longrightarrow\mathcal S_Y(\mathcal L_f).
  \]
\end{definition}

\begin{lemma}[The quotient construction is a graded algebra map]
  \label{lemma:nagata-blowup-universal-graded-map}
  \uses{def:nagata-blowup-universal-graded-map}
  \lean{}

  The degreewise map $\alpha_f$ respects the unit and multiplication and is a
  graded algebra morphism.
\end{lemma}

\begin{proof}
  In degree zero,
  \[
    \alpha_{f,0}\colon f^*\mathcal O_X=\mathcal O_Y
      \longrightarrow\mathcal L_f^{\otimes0}=\mathcal O_Y
      =\operatorname{id}_{\mathcal O_Y}.
  \]
  For $m,n\ge0$, let
  $\mu^{\mathcal I}_{m,n}\colon\mathcal I^m\otimes\mathcal I^n
  \to\mathcal I^{m+n}$ and
  $\mu^{\mathcal L}_{m,n}\colon\mathcal L_f^{\otimes m}\otimes
  \mathcal L_f^{\otimes n}\to\mathcal L_f^{\otimes(m+n)}$ be multiplication.
  The definition of $\alpha_{f,n}$ by tensor powers is the commutative
  equality
  \[
    \alpha_{f,m+n}\circ f^*(\mu^{\mathcal I}_{m,n})
      =\mu^{\mathcal L}_{m,n}\circ
        (\alpha_{f,m}\otimes\alpha_{f,n}).
  \]
  The degree-zero equality is preservation of the unit, and the displayed
  equality is preservation of multiplication.  Since
  $\alpha_f(f^*\mathcal I^n)\subseteq\mathcal L_f^{\otimes n}$ for every
  $n$, $\alpha_f$ is a graded $\mathcal O_Y$-algebra morphism.
\end{proof}

\begin{definition}[Universal lift construction]
  \label{def:nagata-blowup-universal-lift}
  \uses{def:nagata-global-blowup-relative-proj-comparison,
    def:nagata-blowup-universal-graded-map,
    def:nagata-relative-proj,
    lemma:nagata-relative-proj-interface}
  \lean{}

  Retain the finite-type ideal $\mathcal I$, the morphism $f$, and the
  invertible image quotient $q_f$ of
  \Cref{def:nagata-blowup-universal-graded-map}.  Put
  \[
    B_{\mathcal I}:=\operatorname{Bl}_{\mathcal I}(X),\qquad
    b_{\mathcal I}:=b_{X,\mathcal I}\colon B_{\mathcal I}\to X,
  \]
  and define the fibre product and its projections by
  \[
  \begin{gathered}
    P_f:=Y\times_{X,(f,b_{\mathcal I})}B_{\mathcal I},\\
    p_{1,f}\colon P_f\to Y,\qquad
    p_{2,f}\colon P_f\to B_{\mathcal I},\qquad
    f\circ p_{1,f}=b_{\mathcal I}\circ p_{2,f}.
  \end{gathered}
  \tag{UL0}
  \]
  Retain $\mathbb P_{X,\mathcal I}$,
  $p_{X,\mathcal I}\colon\mathbb P_{X,\mathcal I}\to X$, and
  $\delta_{X,\mathcal I}\colon B_{\mathcal I}\xrightarrow{\sim}
  \mathbb P_{X,\mathcal I}$ from
  \Cref{def:nagata-global-blowup-relative-proj-comparison}.  The two Rees
  symbols have the equality
  \[
    \mathscr R_{X,\mathcal I}
      =\bigoplus_{n\ge0}\mathcal I^n
      =\mathcal R_X(\mathcal I).                              \tag{UL0a}
  \]
  Define
  \[
  \begin{gathered}
    Q_f:=Y\times_{X,(f,p_{X,\mathcal I})}\mathbb P_{X,\mathcal I},\\
    q_{1,f}\colon Q_f\to Y,\qquad
    q_{2,f}\colon Q_f\to\mathbb P_{X,\mathcal I},\qquad
    f\circ q_{1,f}=p_{X,\mathcal I}\circ q_{2,f}.
  \end{gathered}                                               \tag{UL0b}
  \]
  Relative-Proj base change
  \cite[\href{https://stacks.math.columbia.edu/tag/01O3}{Tag 01O3}]
  {stacks-project} is the first isomorphism below.  Equation (RP10) defines
  the second one by the fibre-product equations displayed:
  \[
  \begin{gathered}
    \operatorname{bc}_f\colon
      \operatorname{Proj}_Y\!\bigl(f^*\mathcal R_X(\mathcal I)\bigr)
      \xrightarrow{\ \sim\ }Q_f,\\
    \Delta_f:=
      \operatorname{id}_Y\times_X\delta_{X,\mathcal I}^{-1}
      \colon Q_f\xrightarrow{\ \sim\ }P_f,\\
    p_{1,f}\circ\Delta_f=q_{1,f},\qquad
    p_{2,f}\circ\Delta_f=
      \delta_{X,\mathcal I}^{-1}\circ q_{2,f},\\
    r_f:=\Delta_f\circ\operatorname{bc}_f\colon
      \operatorname{Proj}_Y\!\bigl(f^*\mathcal R_X(\mathcal I)\bigr)
      \xrightarrow{\ \sim\ }P_f.
  \end{gathered}                                               \tag{UL1}
  \]
  Since $\alpha_{f,1}=q_f$ is an epimorphism, relative-Proj functoriality
  \cite[\href{https://stacks.math.columbia.edu/tag/01O4}{Tag 01O4}]
  {stacks-project} defines
  \[
  \begin{gathered}
    \operatorname{Proj}_Y(\alpha_f)\colon
      \operatorname{Proj}_Y(\mathcal S_Y(\mathcal L_f))
      \longrightarrow
      \operatorname{Proj}_Y(f^*\mathcal R_X(\mathcal I)),\\
    \pi_f:=r_f\circ\operatorname{Proj}_Y(\alpha_f)\colon
      \operatorname{Proj}_Y(\mathcal S_Y(\mathcal L_f))
      \longrightarrow P_f.
  \end{gathered}
  \tag{UL2}
  \]
  Define the two relative-Proj structure morphisms
  \[
  \begin{gathered}
    s_f^{\mathcal R}\colon
      \operatorname{Proj}_Y\!\bigl(f^*\mathcal R_X(\mathcal I)\bigr)
        \longrightarrow Y,\\
    s_f\colon\operatorname{Proj}_Y(\mathcal S_Y(\mathcal L_f))
      \longrightarrow Y.
  \end{gathered}
  \]
  Relative-Proj base change in
  \cite[\href{https://stacks.math.columbia.edu/tag/01O3}{Tag 01O3}]
  {stacks-project} and relative-Proj functoriality in
  \cite[\href{https://stacks.math.columbia.edu/tag/01O4}{Tag 01O4}]
  {stacks-project} are the following component equalities:
  \[
    q_{1,f}\circ\operatorname{bc}_f=s_f^{\mathcal R},
    \qquad
    s_f^{\mathcal R}\circ\operatorname{Proj}_Y(\alpha_f)=s_f.
  \]
  Substitution of these equalities and (UL1)--(UL2) is the chain
  \[
  \begin{aligned}
    p_{1,f}\circ\pi_f
      &=p_{1,f}\circ\Delta_f\circ\operatorname{bc}_f
          \circ\operatorname{Proj}_Y(\alpha_f)\\
      &=q_{1,f}\circ\operatorname{bc}_f
          \circ\operatorname{Proj}_Y(\alpha_f)\\
      &=s_f^{\mathcal R}\circ\operatorname{Proj}_Y(\alpha_f)
       =s_f.
  \end{aligned}                                                  \tag{UL3}
  \]
  Let
  \[
    \tau_{\mathcal L_f}\colon
      Y\xrightarrow{\sim}
      \operatorname{Proj}_Y(\mathcal S_Y(\mathcal L_f))
  \]
  be the relative-Proj isomorphism characterized locally as follows: if
  $\mathcal L_f|_W=\mathcal O_W\ell$ on an open $W\subseteq Y$, define
  \[
    \theta_{W,\ell}\colon\mathcal O_W[T]\xrightarrow{\sim}
      \mathcal S_Y(\mathcal L_f)|_W,
    \qquad \deg(T)=1,\qquad T\longmapsto\ell,
  \]
  and require
  \[
    \tau_{\mathcal L_f}|_W
      :=\operatorname{Proj}_W(\theta_{W,\ell}^{-1})
      \colon W=\operatorname{Proj}_W(\mathcal O_W[T])
      \xrightarrow{\sim}
      \operatorname{Proj}_W(\mathcal S_Y(\mathcal L_f)|_W).
  \]
  In particular,
  \[
    s_f\circ\tau_{\mathcal L_f}=\operatorname{id}_Y.          \tag{UL4}
  \]
  Define
  \[
    \widetilde f^{(\mathcal I,q_f)}
      :=p_{2,f}\circ\pi_f\circ\tau_{\mathcal L_f}
      \colon Y\longrightarrow B_{\mathcal I},\qquad
    \widetilde f:=\widetilde f^{(\mathcal I,q_f)}.            \tag{UL5}
  \]
\end{definition}

\begin{definition}[Lift set realizing an invertible quotient]
  \label{def:nagata-blowup-lift-set}
  \uses{def:nagata-blowup-universal-lift,
    thm:nagata-blowup-principalizes-center}

  Retain the finite-type ideal $\mathcal I$ and the quotient $q_f$ of
  \Cref{def:nagata-blowup-universal-graded-map}.  Put
  \[
    B_{\mathcal I}:=\operatorname{Bl}_{\mathcal I}(X),
    \qquad b:=b_{X,\mathcal I}\colon B_{\mathcal I}\to X,
    \qquad
    \mu_b\colon b^*\mathcal I\longrightarrow\mathcal O_{B_{\mathcal I}},
    \qquad
    \mathcal L_b:=\mathcal I\mathcal O_{B_{\mathcal I}}
      :=\operatorname{im}(\mu_b),
  \]
  where $\mu_b$ is the multiplication map
  $r\otimes a\mapsto r\,b^\sharp(a)$ on local sections.
  Define the universal image quotient
  \[
    q_b\colon b^*\mathcal I\twoheadrightarrow\mathcal L_b
  \]
  to be the epimorphism in this image factorization.  Its target is
  invertible by \Cref{thm:nagata-blowup-principalizes-center}.
  For $h\colon Y\to B_{\mathcal I}$ and an equality
  $e_h\colon b\circ h=f$, let
  \[
  \begin{aligned}
    a_{h,b,\mathcal I}&\colon h^*b^*\mathcal I
      \xrightarrow{\sim}(b\circ h)^*\mathcal I,
      &a_{h,b,\mathcal I}(r\otimes(s\otimes x))&:=rs\otimes x,\\
    e_{h,\mathcal I}^*&\colon(b\circ h)^*\mathcal I
      \xrightarrow{\sim}f^*\mathcal I
  \end{aligned}
  \]
  be respectively the pullback associator and the equality-transport
  isomorphism defined by substituting $b\circ h=f$ in the source functor.
  Define
  \[
    \vartheta_{h,e_h}\colon h^*b^*\mathcal I
      \xrightarrow{\sim}f^*\mathcal I,
    \qquad
    \vartheta_{h,e_h}:=e_{h,\mathcal I}^*\circ a_{h,b,\mathcal I}.
  \]
  Define
  \[
  \begin{aligned}
    \operatorname{Realizes}_{q_f}(h,e_h)
      \quad\Longleftrightarrow\quad
    &\exists\bigl(\alpha_h\colon
        h^*\mathcal L_b\xrightarrow{\sim}\mathcal L_f\bigr),\quad
      \alpha_h\circ h^*(q_b)
        =q_f\circ\vartheta_{h,e_h},\\
    \operatorname{LiftPair}(f,q_f)
      :=\{(h,e_h)\mid\;&h\colon Y\to B_{\mathcal I},\ 
        e_h\colon b\circ h=f,\
        \operatorname{Realizes}_{q_f}(h,e_h)\},\\
    \operatorname{Lift}(f,q_f)
      :=\{h\colon Y\to B_{\mathcal I}\mid\;&
        \exists(e_h\colon b\circ h=f),\ 
        (h,e_h)\in\operatorname{LiftPair}(f,q_f)\}.
  \end{aligned}
  \]
\end{definition}

\begin{lemma}[Transport of an image quotient along an equality]
  \label{lemma:nagata-image-quotient-transport}
  \lean{}

  Let $\mathcal K\subseteq\mathcal O_X$ be an ideal, let
  $g_0,g_1\colon Y\to X$, and let $e\colon g_0=g_1$.  For $i\in\{0,1\}$,
  define
  \[
    \mu_i\colon g_i^*\mathcal K\to\mathcal O_Y,
    \qquad
    g_i^*\mathcal K\xrightarrow{q_i}\mathcal L_i
      \xrightarrow{\iota_i}\mathcal O_Y,
    \qquad
    \mu_i=\iota_i\circ q_i                              \tag{IQ0}
  \]
  to be the pullback of $\mathcal K\hookrightarrow\mathcal O_X$ and its
  epimorphism--monomorphism image factorization.  Define the object-valued
  family and its transport operation by equality induction:
  \[
  \begin{gathered}
    F_{\mathcal K}\colon\operatorname{Hom}(Y,X)
      \longrightarrow\operatorname{Ob}(\operatorname{QCoh}(Y)),
    \qquad F_{\mathcal K}(g):=g^*\mathcal K,\\
    \operatorname{tr}_{F_{\mathcal K}}
      (\operatorname{refl}_g)
      :=\operatorname{id}_{g^*\mathcal K}.
  \end{gathered}                                                    \tag{IQ0a}
  \]
  Thus, with no implicit transport notation, put
  \[
  \begin{gathered}
    e_{\mathcal K}^*
      :=\operatorname{tr}_{F_{\mathcal K}}(e)
      \colon g_0^*\mathcal K\xrightarrow{\sim}g_1^*\mathcal K,\\
    (e_{\mathcal K}^*)^{-1}
      :=\operatorname{tr}_{F_{\mathcal K}}(e^{-1})
      \colon g_1^*\mathcal K\xrightarrow{\sim}g_0^*\mathcal K,\\
    (e_{\mathcal K}^*)^{-1}\circ e_{\mathcal K}^*
      =\operatorname{id}_{g_0^*\mathcal K},
    \qquad
    e_{\mathcal K}^*\circ(e_{\mathcal K}^*)^{-1}
      =\operatorname{id}_{g_1^*\mathcal K}.
  \end{gathered}                                                    \tag{IQ0b}
  \]
  The last two equalities are the instances of equality induction at
  $e=\operatorname{refl}_{g_0}$.  There is a unique isomorphism
  \[
    \operatorname{iqt}_{\mathcal K}(e)\colon
      \mathcal L_0\xrightarrow{\sim}\mathcal L_1             \tag{IQ1}
  \]
  satisfying
  \[
    \operatorname{iqt}_{\mathcal K}(e)\circ q_0
      =q_1\circ e_{\mathcal K}^*,
    \qquad
    \iota_1\circ\operatorname{iqt}_{\mathcal K}(e)=\iota_0. \tag{IQ2}
  \]
\end{lemma}

\begin{proof}
  Substitution of $e\colon g_0=g_1$ in the two pullback maps is
  \[
    \mu_1\circ e_{\mathcal K}^*=\mu_0.                        \tag{IQ3}
  \]
  By the universal property of the two image factorizations in (IQ0), (IQ3)
  defines unique morphisms
  \[
  \begin{aligned}
    \alpha&\colon\mathcal L_0\to\mathcal L_1,
      &\alpha\circ q_0&=q_1\circ e_{\mathcal K}^*,
      &\iota_1\circ\alpha&=\iota_0,\\
    \beta&\colon\mathcal L_1\to\mathcal L_0,
      &\beta\circ q_1&=q_0\circ(e_{\mathcal K}^*)^{-1},
      &\iota_0\circ\beta&=\iota_1.
  \end{aligned}                                                  \tag{IQ4}
  \]
  Since $q_0$ and $q_1$ are epimorphisms, the chains
  \[
  \begin{aligned}
    (\beta\circ\alpha)\circ q_0
      &=\beta\circ q_1\circ e_{\mathcal K}^*
       =q_0\circ(e_{\mathcal K}^*)^{-1}\circ e_{\mathcal K}^*=q_0,\\
    (\alpha\circ\beta)\circ q_1
      &=\alpha\circ q_0\circ(e_{\mathcal K}^*)^{-1}
       =q_1\circ e_{\mathcal K}^*\circ(e_{\mathcal K}^*)^{-1}=q_1
  \end{aligned}
  \]
  are the inverse equalities
  \[
    \beta\circ\alpha=\operatorname{id}_{\mathcal L_0},
    \qquad
    \alpha\circ\beta=\operatorname{id}_{\mathcal L_1}.
  \]
  Define $\operatorname{iqt}_{\mathcal K}(e):=\alpha$.  If $\alpha'$ also
  satisfies (IQ2), then
  $\alpha'\circ q_0=q_1\circ e_{\mathcal K}^*=\alpha\circ q_0$; epimorphism
  cancellation for $q_0$ is $\alpha'=\alpha$.
\end{proof}

\begin{lemma}[The constructed lift lies over $X$]
  \label{lemma:nagata-blowup-universal-exists}
  \uses{def:nagata-blowup-universal-lift,def:nagata-blowup-lift-set}
  \lean{}

  The morphism $\widetilde f$ of
  \Cref{def:nagata-blowup-universal-lift} is an $X$-morphism:
  $b_{X,\mathcal I}\circ\widetilde f=f$.  If
  $e_{\widetilde f}$ denotes this equality, then
  \[
    (\widetilde f,e_{\widetilde f})
      \in\operatorname{LiftPair}(f,q_f),
    \qquad
    \widetilde f\in\operatorname{Lift}(f,q_f).
  \]
\end{lemma}

\begin{proof}
  \uses{lemma:nagata-blowup-universal-graded-map,
    lemma:nagata-global-blowup-glue-data-valid,
    lemma:nagata-global-blowup-local-maps-agree}
  Retain $P_f,p_{1,f},p_{2,f},\pi_f,s_f$ from
  \Cref{def:nagata-blowup-universal-lift}.  Equations (UL3) and (UL4) are
  \[
    p_{1,f}\circ\pi_f\circ\tau_{\mathcal L_f}
      =s_f\circ\tau_{\mathcal L_f}
      =\operatorname{id}_Y,                                  \tag{UL6}
  \]
  and (UL0) is
  \[
    b_{X,\mathcal I}\circ p_{2,f}=f\circ p_{1,f}.             \tag{UL7}
  \]
  Substitution of (UL5)--(UL7) is the chain
  \[
  \begin{aligned}
    b_{X,\mathcal I}\circ\widetilde f
      &=b_{X,\mathcal I}\circ p_{2,f}\circ\pi_f
        \circ\tau_{\mathcal L_f}\\
      &=f\circ p_{1,f}\circ\pi_f\circ\tau_{\mathcal L_f}\\
      &=f.
  \end{aligned}
  \]
  Denote this equality by $e_{\widetilde f}$.  Choose indexed affine covers
  \[
  \begin{gathered}
    X=\bigcup_{\lambda\in\Lambda}V_\lambda,\qquad
    V_\lambda=\operatorname{Spec}(A_\lambda),\\
    Y=\bigcup_{(\lambda,\rho)\in\Xi}W_{\lambda,\rho},\qquad
    W_{\lambda,\rho}=\operatorname{Spec}(C_{\lambda,\rho})
      \subseteq f^{-1}(V_\lambda),\\
    \mathcal L_f|_{W_{\lambda,\rho}}
      =\mathcal O_{W_{\lambda,\rho}}\ell_{\lambda,\rho}.
  \end{gathered}                                               \tag{UL7a}
  \]
  Finite type is the existence, for every $\lambda\in\Lambda$, of symbols
  \[
    r_\lambda\in\mathbf Z_{\ge0},\qquad
    a_{\lambda,0},\ldots,a_{\lambda,r_\lambda}\in A_\lambda,
    \qquad
    \Gamma(V_\lambda,\mathcal I|_{V_\lambda})
      =(a_{\lambda,0},\ldots,a_{\lambda,r_\lambda}).           \tag{UL7b}
  \]
  Fix $(\lambda,\rho)\in\Xi$ and abbreviate
  \[
  \begin{gathered}
    V:=V_\lambda=\operatorname{Spec}(A),\qquad
    W:=W_{\lambda,\rho}=\operatorname{Spec}(C),\qquad
    \ell:=\ell_{\lambda,\rho},\\
    r:=r_\lambda,\qquad a_i:=a_{\lambda,i},\qquad
    I:=\Gamma(V,\mathcal I|_V)=(a_0,\ldots,a_r).
  \end{gathered}                                               \tag{UL7c}
  \]
  Define the affine-restriction comparison and its charts by
  \[
  \begin{gathered}
    B_V^X:=b_{X,\mathcal I}^{-1}(V),\qquad
    b_V:=b_{X,\mathcal I}|_{B_V^X}\colon B_V^X\to V,\\
    q_V\colon B_V^X\xrightarrow{\ \sim\ }
      \operatorname{Bl}_I(V)
      \quad\text{as in
      \Cref{prop:nagata-global-blowup-affine-restrict}},\\
    D_i^A:=D_+(a_i\mathbf u)\subseteq\operatorname{Bl}_I(V),
    \qquad
    D_i^X:=q_V^{-1}(D_i^A)\subseteq B_V^X,\\
    \chi_i:=\chi_{a_i}\circ q_V|_{D_i^X}\colon
      D_i^X\xrightarrow{\ \sim\ }\operatorname{Spec}(A[I/a_i]),
  \end{gathered}                                               \tag{UL8}
  \]
  where $\chi_{a_i}$ is the isomorphism of
  \Cref{prop:nagata-affine-blowup-chart}.
  For $x\in I$, define $c_x\in C$ by
  \[
    q_f(f^*x)|_W=c_x\ell.
  \]
  Put $c_i:=c_{a_i}$ for $0\le i\le r$.
  Define the chosen trivialization by
  \[
    t_\ell\colon\mathcal O_W\xrightarrow{\ \sim\ }\mathcal L_f|_W,
    \qquad t_\ell(r):=r\ell.
  \]
  From $I=(a_0,\ldots,a_r)$ and the epimorphism
  $q_f|_W\colon f^*\mathcal I|_W\twoheadrightarrow\mathcal L_f|_W$, the
  ideal-image computation is
  \[
  \begin{aligned}
    \mathcal L_f|_W
      &=\operatorname{im}(q_f|_W)\\
      &=q_f\!\left(\sum_{i=0}^r
          \mathcal O_W(f|_W)^*a_i\right)\\
      &=\sum_{i=0}^r\mathcal O_W(c_i\ell)
       =t_\ell\!\left((c_0,\ldots,c_r)\mathcal O_W\right).
  \end{aligned}                                                \tag{UL7d}
  \]
  Applying $t_\ell^{-1}$ to (UL7d) is the equality
  \[
    (c_0,\ldots,c_r)\mathcal O_W=\mathcal O_W,
    \qquad (c_0,\ldots,c_r)C=C.                               \tag{UL7e}
  \]
  The affine distinguished-open identity applied to (UL7e) is
  \[
    W_i:=D_W(c_i)=\operatorname{Spec}(C_{c_i}),\qquad
    W=\operatorname{Spec}(C)
      =\bigcup_{i=0}^rD_W(c_i)=\bigcup_{i=0}^rW_i.
  \]
  For $x\in I$, define the universal and pulled-back local sections by
  \[
  \begin{gathered}
    \overline x^{\,b}:=
      q_b|_{B_V^X}(b_V^*x)
      \in\Gamma(B_V^X,\mathcal L_b|_{B_V^X}),\\
    \overline x^{\,f}:=
      q_f|_W((f|_W)^*x)=c_x\ell
      \in\Gamma(W,\mathcal L_f|_W),\\
    \overline a_i^b:=\overline{a_i}^{\,b},\qquad
    \overline a_i^f:=\overline{a_i}^{\,f}=c_i\ell.
  \end{gathered}                                               \tag{UL9}
  \]
  Relative-Proj functoriality in (UL2), restricted through (UL8), is the
  pair of formulas
  \[
    \widetilde f(W_i)\subseteq D_i^X,\qquad
    \widetilde f^*(\overline x^{\,b})|_{W_i}
      =\frac{c_x}{c_i}\,
        \widetilde f^*(\overline a_i^b)|_{W_i}
      \quad(x\in I).                                      \tag{UP2a}
      \label{eq:nagata-universal-lift-chart}
  \]
  Both $\widetilde f^*(\overline a_i^b)|_{W_i}$ and
  $\overline a_i^f|_{W_i}=c_i\ell$ generate their respective invertible
  modules.  Define the local isomorphism by the explicit formula
  \[
  \begin{aligned}
    \alpha_{\widetilde f,W_i}\colon
      \widetilde f^*\mathcal L_b|_{W_i}
      &\xrightarrow{\ \sim\ }\mathcal L_f|_{W_i},\\
    \alpha_{\widetilde f,W_i}
      \bigl(r\,\widetilde f^*(\overline a_i^b)\bigr)
      &:=r\,\overline a_i^f=r\,c_i\ell
      \qquad(r\in\mathcal O_Y(W_i)).
  \end{aligned}
  \]
  For every $x\in I$, substitution of (UP2a) into the defining formula for
  $\alpha_{\widetilde f,W_i}$ is the equality chain
  \[
  \begin{aligned}
    \alpha_{\widetilde f,W_i}
      \bigl(\widetilde f^*(\overline x^{\,b})\bigr)
      &=\alpha_{\widetilde f,W_i}
        \left(\frac{c_x}{c_i}\widetilde f^*(\overline a_i^b)\right)\\
      &=\frac{c_x}{c_i}c_i\ell
       =c_x\ell=q_f(f^*x).
  \end{aligned}                                             \tag{UP2b}
  \]
  Equality (UP2b) on the generators $x\in I$ is the sheaf-morphism equality
  \[
    \alpha_{\widetilde f,W_i}\circ
      \widetilde f^*(q_b)|_{W_i}
      =q_f|_{W_i}\circ
        \vartheta_{\widetilde f,e_{\widetilde f}}|_{W_i}.     \tag{UP2}
  \]
  Restore the indices suppressed in (UL7c) and put
  \[
    \Upsilon:=
      \{\,\xi=(\lambda,\rho,i)\mid
        (\lambda,\rho)\in\Xi,\ 0\le i\le r_\lambda\,\},\qquad
    U_\xi:=W_{\lambda,\rho,i}.
  \]
  Equations (UL7a) and $W_{\lambda,\rho}=\bigcup_iW_{\lambda,\rho,i}$
  are the cover equality
  \[
    Y=\bigcup_{\xi\in\Upsilon}U_\xi.                           \tag{UL10}
  \]
  Write $\alpha_\xi$ for the map
  $\alpha_{\widetilde f,W_{\lambda,\rho,i}}$ defined above.  For
  $\xi,\xi'\in\Upsilon$, set
  \[
    U_{\xi\xi'}:=U_\xi\cap U_{\xi'},\qquad
    Q_{\xi\xi'}:=
      \widetilde f^*(q_b)|_{U_{\xi\xi'}}.
  \]
  Since pullback and restriction preserve epimorphisms,
  $Q_{\xi\xi'}$ is an epimorphism.  The two instances of (UP2) form the
  equality chain
  \[
  \begin{aligned}
    (\alpha_\xi|_{U_{\xi\xi'}})\circ Q_{\xi\xi'}
      &=(q_f\circ\vartheta_{\widetilde f,e_{\widetilde f}})
        |_{U_{\xi\xi'}}\\
      &=(\alpha_{\xi'}|_{U_{\xi\xi'}})\circ Q_{\xi\xi'}.
  \end{aligned}
  \]
  Epimorphism cancellation is the overlap equality
  \[
    \alpha_\xi|_{U_{\xi\xi'}}
      =\alpha_{\xi'}|_{U_{\xi\xi'}}.                           \tag{UL11}
  \]
  Define the Hom fibre
  \[
    \mathscr A_f:=
      \left\{\alpha\in
        \operatorname{Hom}_{\mathcal O_Y}
          (\widetilde f^*\mathcal L_b,\mathcal L_f)
        \ \middle|\
        \alpha|_{U_\xi}=\alpha_\xi\ \text{for all }\xi\in\Upsilon
      \right\}.
  \]
  The sheaf-Hom equalizer for the cover (UL10) and equations (UL11) is the
  singleton equality
  \[
    \mathscr A_f=\{\alpha_{\widetilde f}\},\qquad
    \alpha_{\widetilde f}|_{U_\xi}=\alpha_\xi.                 \tag{UL12}
  \]
  Each $\alpha_\xi$ is an isomorphism.  Applying the same equalizer to
  $(\alpha_\xi^{-1})_\xi$ defines
  $\beta_{\widetilde f}\colon\mathcal L_f\to
  \widetilde f^*\mathcal L_b$ and the local equations
  \[
    (\beta_{\widetilde f}\circ\alpha_{\widetilde f})|_{U_\xi}
      =\operatorname{id}|_{U_\xi},\qquad
    (\alpha_{\widetilde f}\circ\beta_{\widetilde f})|_{U_\xi}
      =\operatorname{id}|_{U_\xi}.
  \]
  Morphism locality on (UL10) is
  \[
    \alpha_{\widetilde f}\colon
      \widetilde f^*\mathcal L_b\xrightarrow{\ \sim\ }\mathcal L_f,\qquad
    \beta_{\widetilde f}\circ\alpha_{\widetilde f}
      =\operatorname{id},\qquad
    \alpha_{\widetilde f}\circ\beta_{\widetilde f}
      =\operatorname{id}.                                    \tag{UL13}
  \]
  Equation (UP2) then holds globally, which is
  exactly
  \[
    \operatorname{Realizes}_{q_f}
      (\widetilde f,e_{\widetilde f}).
  \]
  Substitution in the two set definitions in
  \Cref{def:nagata-blowup-lift-set} is
  \[
    (\widetilde f,e_{\widetilde f})
      \in\operatorname{LiftPair}(f,q_f),
    \qquad
    \widetilde f\in\operatorname{Lift}(f,q_f).
  \]
\end{proof}

\begin{lemma}[Uniqueness of the universal lift]
  \label{lemma:nagata-blowup-universal-unique}
  \uses{def:nagata-blowup-universal-lift,def:nagata-blowup-lift-set}
  \lean{}

  Assume that $\mathcal I$ is of finite type.  Let
  $(h,e_h)\in\operatorname{LiftPair}(f,q_f)$ as defined in
  \Cref{def:nagata-blowup-lift-set}.  Then
  \[
    h=\widetilde f.
  \]
\end{lemma}

\begin{proof}
  By \Cref{def:nagata-blowup-lift-set}, choose
  \[
    e_h\colon b_{X,\mathcal I}\circ h=f,
    \qquad
    \alpha_h\colon h^*\mathcal L_b\xrightarrow{\sim}\mathcal L_f
  \]
  satisfying
  \[
    \alpha_h\circ h^*(q_b)=q_f\circ\vartheta_{h,e_h}.          \tag{UP3}
  \]
  Choose indexed affine covers and trivializations
  \[
  \begin{gathered}
    X=\bigcup_{\lambda\in\Lambda}V_\lambda,\qquad
    V_\lambda=\operatorname{Spec}(A_\lambda),\\
    Y=\bigcup_{(\lambda,\rho)\in\Xi}W_{\lambda,\rho},\qquad
    W_{\lambda,\rho}=\operatorname{Spec}(C_{\lambda,\rho})
      \subseteq f^{-1}(V_\lambda),\\
    \mathcal L_f|_{W_{\lambda,\rho}}
      =\mathcal O_{W_{\lambda,\rho}}\ell_{\lambda,\rho}.
  \end{gathered}                                               \tag{UQ0}
  \]
  For every $\lambda$, choose the finite presentation
  \[
    r_\lambda\in\mathbf Z_{\ge0},\qquad
    \Gamma(V_\lambda,\mathcal I|_{V_\lambda})
      =(a_{\lambda,0},\ldots,a_{\lambda,r_\lambda}).           \tag{UQ1}
  \]
  Fix $(\lambda,\rho)\in\Xi$ and abbreviate
  \[
  \begin{gathered}
    V:=V_\lambda=\operatorname{Spec}(A),\qquad
    W:=W_{\lambda,\rho}=\operatorname{Spec}(C),\qquad
    \ell:=\ell_{\lambda,\rho},\\
    r:=r_\lambda,\qquad a_i:=a_{\lambda,i},\qquad
    I:=\Gamma(V,\mathcal I|_V)=(a_0,\ldots,a_r).
  \end{gathered}                                               \tag{UQ2}
  \]
  Define
  \[
  \begin{gathered}
    B_V^X:=b_{X,\mathcal I}^{-1}(V),\qquad
    b_V:=b_{X,\mathcal I}|_{B_V^X}\colon B_V^X\to V,\\
    q_V\colon B_V^X\xrightarrow{\ \sim\ }\operatorname{Bl}_I(V)
      \quad\text{as in
      \Cref{prop:nagata-global-blowup-affine-restrict}},\\
    D_i^A:=D_+(a_i\mathbf u)\subseteq\operatorname{Bl}_I(V),
    \qquad D_i^X:=q_V^{-1}(D_i^A),\\
    \chi_i:=\chi_{a_i}\circ q_V|_{D_i^X}\colon
      D_i^X\xrightarrow{\ \sim\ }\operatorname{Spec}(A[I/a_i]).
  \end{gathered}                                               \tag{UQ3}
  \]
  For every $x\in I$, define $c_x\in C$ by
  \[
    q_f(f^*x)|_W=c_x\ell.
  \]
  Put $c_i:=c_{a_i}$ and define
  \[
    t_\ell\colon\mathcal O_W\xrightarrow{\ \sim\ }\mathcal L_f|_W,
    \qquad t_\ell(r):=r\ell.
  \]
  The equalities $I=(a_0,\ldots,a_r)$ and
  $\operatorname{im}(q_f|_W)=\mathcal L_f|_W$ are the chain
  \[
  \begin{aligned}
    \mathcal L_f|_W
      &=q_f\!\left(\sum_{i=0}^r
          \mathcal O_W(f|_W)^*a_i\right)\\
      &=\sum_{i=0}^r\mathcal O_W(c_i\ell)
       =t_\ell\!\left((c_0,\ldots,c_r)\mathcal O_W\right).
  \end{aligned}                                                \tag{UQ3a}
  \]
  Applying $t_\ell^{-1}$ and then $\operatorname{Spec}$ is the pair of
  equalities
  \[
    (c_0,\ldots,c_r)C=C,\qquad
    W=\operatorname{Spec}(C)
      =\bigcup_{i=0}^rD_W(c_i).                               \tag{UQ3b}
  \]
  Define
  \[
    \overline x^{\,b}:=
      q_b|_{B_V^X}(b_V^*x)
      \in\Gamma(B_V^X,\mathcal L_b|_{B_V^X}),\qquad
    \overline a_i^b:=\overline{a_i}^{\,b}.                    \tag{UQ4}
  \]
  Define
  \[
    W_i:=D_W(c_i)=\operatorname{Spec}(C_{c_i}).
  \]
  Substitution of this definition into (UQ3b) is the equality
  \[
    W=\bigcup_{i=0}^rD_W(c_i)=\bigcup_{i=0}^rW_i.                \tag{UQ3c}
  \]
  Evaluating (UP3) on $(f|_W)^*x$ is the equality
  \[
    \alpha_h\bigl(h^*(\overline x^{\,b})\bigr)|_W
      =c_x\ell\qquad(x\in I).                                 \tag{UQ5}
  \]
  For $y\in W_i$, (UQ5) at $x=a_i$ and
  $c_i|_{W_i}\in C_{c_i}^{\times}$ are
  \[
    \alpha_{h,y}\bigl((h^*\overline a_i^b)_y\bigr)
      =c_i(y)\ell_y,\qquad c_i(y)\in\mathcal O_{Y,y}^{\times}.
  \]
  The standard-chart membership criterion in
  \Cref{prop:nagata-affine-blowup-chart} is therefore the set inclusion
  \[
    h(W_i)\subseteq D_i^X.                                    \tag{UQ6}
  \]
  Equation \eqref{eq:nagata-universal-lift-chart} is
  \[
    \widetilde f(W_i)\subseteq D_i^X.                         \tag{UQ7}
  \]
  Define the
  corestrictions and chart morphisms
  \[
  \begin{gathered}
    h_i\colon W_i\to D_i^X,\qquad
    \widetilde f_i\colon W_i\to D_i^X,\\
    \psi_{h,i}:=\chi_i\circ h_i,\qquad
    \psi_{\widetilde f,i}:=\chi_i\circ\widetilde f_i
      \colon W_i\longrightarrow\operatorname{Spec}(A[I/a_i]).
  \end{gathered}
  \]
  Define the affine corestriction
  \[
    f_W:=f|_W\colon W\longrightarrow V.
  \]
  By \Cref{def:nagata-direct-image-structure-map}, their affine
  global-section ring maps have the source and targets
  \[
  \begin{gathered}
    f_W^{\#}:=\varepsilon_{f_W,|V|}\colon A\longrightarrow C,\\
    \psi_{h,i}^{\#}:=
      \varepsilon_{\psi_{h,i},
        |\operatorname{Spec}(A[I/a_i])|}\colon
      A[I/a_i]\longrightarrow C_{c_i},\\
    \psi_{\widetilde f,i}^{\#}:=
      \varepsilon_{\psi_{\widetilde f,i},
        |\operatorname{Spec}(A[I/a_i])|}\colon
      A[I/a_i]\longrightarrow C_{c_i}.
  \end{gathered}
  \]
  For $a\in A$, the two structural equalities
  $b\circ h=f=b\circ\widetilde f$ are
  \[
    \psi_{h,i}^{\#}(a)
      =f_W^{\#}(a)/1
      =\psi_{\widetilde f,i}^{\#}(a).                         \tag{UQ8}
  \]
  For $x\in I$, the chart identity
  $\overline x^{\,b}=(x/a_i)\overline a_i^b$ together with (UQ5) is
  \[
  \begin{aligned}
    c_x\ell
      &=\alpha_h(h^*\overline x^{\,b})
       =\psi_{h,i}^{\#}(x/a_i)\,
          \alpha_h(h^*\overline a_i^b)\\
      &=\psi_{h,i}^{\#}(x/a_i)c_i\ell.
  \end{aligned}
  \]
  Since $\ell$ is a basis and $c_i\in C_{c_i}^{\times}$, coefficient
  cancellation is
  \[
    \psi_{h,i}^{\#}(x/a_i)
      =c_x/c_i.                                               \tag{UQ9}
  \]
  The second equality of
  \eqref{eq:nagata-universal-lift-chart} is, under $\chi_i$,
  \[
    \psi_{\widetilde f,i}^{\#}(x/a_i)=c_x/c_i.                \tag{UQ10}
  \]
  By \Cref{def:nagata-blowup-chart-ring},
  \[
    A[I/a_i]
      =A[x/a_i\mid x\in I]\subseteq A_{a_i}.
  \]
  Equations (UQ8)--(UQ10) on the displayed generating set are the
  ring-homomorphism equality
  \[
    \psi_{h,i}^{\#}=\psi_{\widetilde f,i}^{\#}.
  \]
  Applying the contravariant functor $\operatorname{Spec}$ is the equality
  \[
    h|_{W_i}=\widetilde f|_{W_i}.
  \]
  Restoring indices, define
  \[
    \Upsilon:=
      \{(\lambda,\rho,i)\mid(\lambda,\rho)\in\Xi,\
        0\le i\le r_\lambda\},\qquad
    U_{\lambda,\rho,i}:=W_{\lambda,\rho,i}.
  \]
  Equations (UQ0) and
  $W_{\lambda,\rho}=\bigcup_iW_{\lambda,\rho,i}$ are
  \[
    Y=\bigcup_{(\lambda,\rho,i)\in\Upsilon}U_{\lambda,\rho,i}.
  \]
  The preceding equality holds on every $U_{\lambda,\rho,i}$; source
  locality for this displayed cover is
  \[
    h=\widetilde f.
  \]
\end{proof}

\begin{theorem}[{\cite[\href{https://stacks.math.columbia.edu/tag/0806}{Tag 0806}]{stacks-project}}]
  \label{thm:nagata-blowup-universal-property}
  \uses{def:nagata-blowup-universal-lift,def:nagata-blowup-lift-set}
  \lean{}

  Assume that $\mathcal I$ is of finite type.  With
  $\operatorname{Lift}(f,q_f)$ defined in
  \Cref{def:nagata-blowup-lift-set}, one has
  \[
    \operatorname{Lift}(f,q_f)=\{\widetilde f\}.
  \]
\end{theorem}

\begin{proof}
  \uses{lemma:nagata-blowup-universal-exists,
    lemma:nagata-blowup-universal-unique}
  \Cref{lemma:nagata-blowup-universal-exists} is the membership
  \[
    \widetilde f\in\operatorname{Lift}(f,q_f).
  \]
  If $h\in\operatorname{Lift}(f,q_f)$, its defining membership supplies an
  equality $e_h\colon b\circ h=f$ with
  $(h,e_h)\in\operatorname{LiftPair}(f,q_f)$.  Therefore
  \Cref{lemma:nagata-blowup-universal-unique} is the equality
  $h=\widetilde f$.  The two set inclusions are
  \[
    \{\widetilde f\}\subseteq\operatorname{Lift}(f,q_f),
    \qquad
    \operatorname{Lift}(f,q_f)\subseteq\{\widetilde f\},
  \]
  and antisymmetry of inclusion is
  $\operatorname{Lift}(f,q_f)=\{\widetilde f\}$.
\end{proof}

\begin{corollary}[Blow-up of an effective Cartier center]
  \label{cor:nagata-blowup-cartier-iso}
  \lean{AlgebraicGeometry.Scheme.IdealSheafData.isIso_blowupMap_of_isInvertible}
  \leanok

  If $\mathcal I$ is already invertible on $X$, then
  $b_{X,\mathcal I}$ is an isomorphism.

  \paragraph{Formalization.}
  The hypothesis is $\operatorname{InvertibleIdeal}_X(\mathcal I)$, predicate (MLCI0) of
  \Cref{def:nagata-invertible-ideal-effective-cartier}, i.e.\
  \texttt{\detokenize{AlgebraicGeometry.Scheme.IdealSheafData.IsInvertible}}, and the
  conclusion is \texttt{\detokenize{IsIso}} for the blow-up map
  \texttt{\detokenize{Scheme.IdealSheafData.blowupMap}}.  No finiteness assumption on
  $\mathcal I$ is needed.

  \textbf{The Lean proof does not follow the proof below, and
  \Cref{thm:nagata-blowup-universal-property} is not formalized.}  The universal property is
  phrased through relative $\operatorname{Proj}$, which Mathlib does not have, and this
  repository's blow-up is glued from affine blow-ups rather than constructed as a relative
  $\operatorname{Proj}$, so the argument below is not available.  Instead the Lean proof
  reduces to the affine charts: being an isomorphism is Zariski-local on the target, and on
  an affine open where $\mathcal I$ is generated by a single non-zerodivisor $a$ the Rees
  algebra is the polynomial ring $A[aX]$ --- regularity of $a$ is exactly what makes the
  powers $a^nX^n$ free over $A$ --- so the single chart at $a\cdot X$ already covers, and the
  structure morphism of the affine blow-up is $\operatorname{Spec}$ of a ring isomorphism.
  The affine statement is
  \texttt{\detokenize{AlgebraicGeometry.AffineBlowup.isIso_toSpec_of_isRegularPrincipal}}, of
  which \Cref{lemma:nagata-affine-blowup-top} is the special case $\mathcal I=(1)$.  The
  regularity hypothesis is indispensable rather than an artefact: for $A=k[x]/(x^2)$ and
  $a=x$ the ideal $(a)$ is principal but not invertible, and its blow-up is not
  $\operatorname{Spec}(A)$.
\end{corollary}

\begin{proof}
  \uses{thm:nagata-blowup-principalizes-center,thm:nagata-blowup-universal-property}
  Put $B:=\operatorname{Bl}_{\mathcal I}(X)$ and $b:=b_{X,\mathcal I}$.
  Since $\mathcal I$ is invertible, define
  $q_{\operatorname{id}}:=\operatorname{id}_{\mathcal I}$.  The singleton
  assertion of \Cref{thm:nagata-blowup-universal-property} for
  $(\operatorname{id}_X,q_{\operatorname{id}})$ defines
  $s\colon X\to B$ with
  \[
    b\circ s=\operatorname{id}_X.
  \]
  By \Cref{thm:nagata-blowup-principalizes-center},
  $\mathcal I\mathcal O_B$ is invertible.  Let
  $q_b\colon b^*\mathcal I\twoheadrightarrow\mathcal I\mathcal O_B$ be
  \Cref{def:nagata-blowup-lift-set}.  The pullback identities are
  \[
    \operatorname{id}_B,\ s\circ b\in\operatorname{Lift}(b,q_b).
  \]
  The singleton equality in
  \Cref{thm:nagata-blowup-universal-property}, applied to $(b,q_b)$, is
  \[
    s\circ b=\operatorname{id}_B.
  \]
  Together with $b\circ s=\operatorname{id}_X$, this is the pair of inverse
  identities defining $s=b^{-1}$.
\end{proof}

\begin{lemma}[Flat pullback commutes with ideal powers]
  \label{lemma:nagata-flat-pullback-ideal-pow}
  \lean{}

  Let $X$ and $Y$ be schemes, let
  $\mathcal I\subseteq\mathcal O_X$ be a quasi-coherent ideal, and let
  $g\colon Y\to X$ be flat.  For $n\ge0$, let
  \[
    \mu_{g,n}\colon g^*(\mathcal I^n)\longrightarrow\mathcal O_Y
  \]
  be the pullback of $\mathcal I^n\hookrightarrow\mathcal O_X$, and put
  $\mathcal I_Y:=\operatorname{im}(\mu_{g,1})$.  Then
  \[
    \operatorname{im}(\mu_{g,n})=\mathcal I_Y^n
  \]
  as ideal subsheaves of $\mathcal O_Y$ for every $n$.
\end{lemma}

\begin{proof}
  The assertion is local on $X$ and $Y$.  On affine opens with
  $X=\operatorname{Spec}(A)$, $Y=\operatorname{Spec}(B)$, and ring map
  $\varphi\colon A\to B$, flatness of $g$ is flatness of $B$ over $A$.
  If $I\subseteq A$ represents $\mathcal I$, tensoring the injection
  $I^n\hookrightarrow A$ with $B$ is the injection
  \[
    B\otimes_AI^n\hookrightarrow B\otimes_AA=B.
  \]
  Its image is the extended ideal $(I^n)B$.  Extension preserves products, so
  \[
    \operatorname{im}(\mu_{g,n})
      =(I^n)B=(IB)^n=\mathcal I_Y^n.
  \]
  For every localization $B\to B_t$, both sides localize to
  $B_t\otimes_AI^n=(IB_t)^n$.  Compatibility with every restriction map is
  the gluing equality of the asserted ideal sheaves.
\end{proof}

\begin{definition}[Flat pullback map for Rees algebras]
  \label{def:nagata-flat-pullback-rees-map}
  \uses{def:nagata-qcoh-graded-algebra-interface,
    def:nagata-rees-algebra,
    lemma:nagata-qcoh-graded-algebra-interface}
  \lean{}

  Let $X$ and $Y$ be schemes, let
  $\mathcal I\subseteq\mathcal O_X$ be a quasi-coherent ideal, and let
  $g\colon Y\to X$ be flat.  Retain $\mathcal I_Y$ and $\mu_{g,n}$ from
  \Cref{lemma:nagata-flat-pullback-ideal-pow}.  Let
  \[
    \beta_{g,n}\colon g^*(\mathcal I^n)\longrightarrow\mathcal I_Y^n
  \]
  be the corestriction of $\mu_{g,n}$ to its image, which equals
  $\mathcal I_Y^n$ by
  \Cref{lemma:nagata-flat-pullback-ideal-pow}.  Define
  \[
    \beta_g:=\bigoplus_{n\ge0}\beta_{g,n}\colon
      g^*\mathcal R_X(\mathcal I)\longrightarrow
      \mathcal R_Y(\mathcal I_Y).
  \]
\end{definition}

\begin{lemma}[Flat pullback of the Rees algebra]
  \label{lemma:nagata-flat-pullback-rees}
  \uses{def:nagata-flat-pullback-rees-map}
  \lean{}

  Let $X,Y,\mathcal I,g$ be the data of
  \Cref{def:nagata-flat-pullback-rees-map}.  The map $\beta_g$ is a graded
  algebra equivalence.
\end{lemma}

\begin{proof}
  \uses{lemma:nagata-flat-pullback-ideal-pow}
  By \Cref{lemma:nagata-flat-pullback-ideal-pow}, each
  $\beta_{g,n}$ is an isomorphism
  \[
    g^*(\mathcal I^n)\xrightarrow{\sim}\mathcal I_Y^n.
  \]
  For $m,n\ge0$, naturality of multiplication of ideal powers is the
  equality
  \[
    \beta_{g,m+n}\circ g^*(\mu^{\mathcal I}_{m,n})
      =\mu^{\mathcal I_Y}_{m,n}\circ
        (\beta_{g,m}\otimes\beta_{g,n}).
  \]
  Also $\beta_{g,0}=\operatorname{id}_{\mathcal O_Y}$.  These unit and
  multiplication equalities state that $\beta_g$ is a graded algebra
  morphism.  Define its inverse degreewise by
  \[
    \beta_g^{-1}:=\bigoplus_{n\ge0}\beta_{g,n}^{-1}.
  \]
  Then
  \[
    \beta_g^{-1}\circ\beta_g
      =\bigoplus_n(\beta_{g,n}^{-1}\beta_{g,n})
      =\operatorname{id}_{g^*\mathcal R_X(\mathcal I)},
    \qquad
    \beta_g\circ\beta_g^{-1}
      =\operatorname{id}_{\mathcal R_Y(\mathcal I_Y)}.
  \]
\end{proof}

\begin{definition}[Flat base-change comparison for blow-ups]
  \label{def:nagata-blowup-flat-base-change-map}
  \uses{def:nagata-global-blowup,
    def:nagata-flat-pullback-rees-map,
    def:nagata-relative-proj,
    lemma:nagata-relative-proj-interface}
  \lean{}

  Let $X$ and $Y$ be schemes, let
  $\mathcal I\subseteq\mathcal O_X$ be a finite-type quasi-coherent ideal,
  and let $g\colon Y\to X$ be flat.  Retain
  $\mathcal I_Y$ and $\beta_g$ from
  \Cref{def:nagata-flat-pullback-rees-map}.  
  Let $A,r,I=(f_0,\ldots,f_r)$ and $\rho_I$ be the datum and morphism of
  \Cref{def:nagata-rees-polynomial-map}.  Put
  \[
    P_g:=Y\times_X\operatorname{Bl}_{\mathcal I}(X).
  \]
  Define
  \[
    \operatorname{bc}_g:=\operatorname{Proj}_Y(\beta_g)\colon
    \operatorname{Bl}_{\mathcal I_Y}(Y)
      \longrightarrow
    P_g
  \]
  for the graded equivalence $\beta_g$ from
  \Cref{lemma:nagata-flat-pullback-rees} and contravariance of relative
  Proj.
\end{definition}

\begin{theorem}[Flat base change of blow-ups]
  \label{thm:nagata-blowup-flat-base-change}
  \uses{def:nagata-blowup-flat-base-change-map}
  \lean{}

  Let $X,Y,\mathcal I,g$ be the data of
  \Cref{def:nagata-blowup-flat-base-change-map}.  The morphism
  $\operatorname{bc}_g$ of
  \Cref{def:nagata-blowup-flat-base-change-map} is an isomorphism over
  $Y$.
\end{theorem}

\begin{proof}
  \uses{lemma:nagata-flat-pullback-rees,
    lemma:nagata-global-blowup-overlap-cocycle}
  By \Cref{lemma:nagata-flat-pullback-rees},
  \[
    \beta_g\colon g^*\mathcal R_X(\mathcal I)
      \xrightarrow{\sim}\mathcal R_Y(\mathcal I_Y)
  \]
  is a graded equivalence whose degree-zero component is
  $\operatorname{id}_{\mathcal O_Y}$.  Relative Proj sends $\beta_g$ to
  \[
    \operatorname{Proj}_Y(\mathcal R_Y(\mathcal I_Y))
      \xrightarrow{\sim}
    \operatorname{Proj}_Y(g^*\mathcal R_X(\mathcal I))
      =P_g.
  \]
  This morphism is $\operatorname{bc}_g$ by
  \Cref{def:nagata-blowup-flat-base-change-map}.  Functoriality on degree zero
  is the equality
  \[
    \operatorname{pr}_Y\circ\operatorname{bc}_g
      =b_{Y,\mathcal I_Y},
  \]
  so the isomorphism is over $Y$.
\end{proof}

\begin{lemma}[Product of invertible ideals is invertible]
  \label{lemma:nagata-invertible-ideal-mul}
  \lean{}

  Let $\mathcal I,\mathcal J\subseteq\mathcal O_X$ be invertible ideal
  sheaves.  Define the multiplication morphism and its image by
  \[
    \mu_{\mathcal I,\mathcal J}\colon
      \mathcal I\otimes_{\mathcal O_X}\mathcal J
        \longrightarrow\mathcal O_X,
    \qquad s\otimes t\longmapsto st,
    \qquad
    \mathcal I\mathcal J:=
      \operatorname{im}(\mu_{\mathcal I,\mathcal J}).
  \]
  Then $\mathcal I\mathcal J$ is an invertible ideal sheaf.
\end{lemma}

\begin{proof}
  Choose an open cover $(U_\lambda)_\lambda$ such that
  \[
    \mathcal I|_{U_\lambda}=(a_\lambda),\qquad
    \mathcal J|_{U_\lambda}=(b_\lambda),
  \]
  with multiplication by $a_\lambda$ and $b_\lambda$ injective.  Then
  \[
    (\mathcal I\mathcal J)|_{U_\lambda}=(a_\lambda b_\lambda).
  \]
  If $(a_\lambda b_\lambda)c=0$, then
  $a_\lambda(b_\lambda c)=0$, so $b_\lambda c=0$, and then $c=0$.
  The implication
  $(a_\lambda b_\lambda)c=0\Rightarrow c=0$ states that
  $a_\lambda b_\lambda$ is a non-zero-divisor, and
  \[
    \mathcal O_{U_\lambda}\xrightarrow{\cdot a_\lambda b_\lambda}
      (\mathcal I\mathcal J)|_{U_\lambda}
  \]
  is an isomorphism.  These local trivializations satisfy the definition of
  an invertible ideal, so
  \[
    \mathcal I\mathcal J\text{ is invertible}.
  \]
\end{proof}

\begin{lemma}[{\cite[\href{https://stacks.math.columbia.edu/tag/07ZV}{Tag 07ZV}]{stacks-project}}]
  \label{lemma:nagata-invertible-product-factor}
  \lean{}

  If the product $\mathcal I\mathcal J$ defines an effective Cartier divisor,
  then each of $\mathcal I$ and $\mathcal J$ defines an effective Cartier
  divisor.
\end{lemma}

\begin{proof}
  Fix $x\in X$ and put
  $R:=\mathcal O_{X,x}$, $I:=\mathcal I_x$, and $J:=\mathcal J_x$.
  The hypothesis supplies $f\in R$ with
  \[
    IJ=(f),\qquad
    (f r=0\Longrightarrow r=0).
  \]
  By the local algebra statement of \cite[Tag 07ZV]{stacks-project}, there
  exist $a,b\in R$ and $u\in R^\times$ such that
  \[
    I=(a),\qquad J=(b),\qquad ab=uf,
  \]
  and $a,b$ are non-zero-divisors.  After shrinking to an open neighborhood
  $U_x$ of $x$, these elements lift to sections satisfying
  \[
    \mathcal I|_{U_x}=(a),\qquad
    \mathcal J|_{U_x}=(b),
  \]
  with $a,b$ non-zero-divisors on $U_x$.  The multiplication maps
  $\mathcal O_{U_x}\xrightarrow{\cdot a}\mathcal I|_{U_x}$ and
  $\mathcal O_{U_x}\xrightarrow{\cdot b}\mathcal J|_{U_x}$ are isomorphisms.
  Since the neighborhoods $U_x$ cover $X$, both ideal sheaves are invertible.
\end{proof}

\begin{definition}[Comparison maps for the product-ideal blow-up]
  \label{def:nagata-blowup-product-comparison}
  \uses{def:nagata-blowup-universal-lift,def:nagata-global-blowup,
    lemma:nagata-blowup-universal-exists}
  \lean{}

  Let $\mathcal I,\mathcal J\subseteq\mathcal O_X$ be finite-type ideals.
  Define
  \[
    B_{\mathcal I}:=\operatorname{Bl}_{\mathcal I}(X),\qquad
    b_{\mathcal I}:=b_{X,\mathcal I},\qquad
    \mathcal J_{\mathcal I}:=
      \mathcal J\mathcal O_{B_{\mathcal I}},
  \]
  and
  \[
    X_{\mathcal I,\mathcal J}:=
      \operatorname{Bl}_{\mathcal J_{\mathcal I}}(B_{\mathcal I}),\qquad
    b_{\mathcal J|\mathcal I}:=
      b_{B_{\mathcal I},\mathcal J_{\mathcal I}}\colon
      X_{\mathcal I,\mathcal J}\to B_{\mathcal I},\qquad
    q_{\mathcal I,\mathcal J}:=
      b_{\mathcal I}\circ b_{\mathcal J|\mathcal I}.
  \]
  Let
  \[
    \iota_{\mathcal I}\colon\mathcal I\hookrightarrow\mathcal O_X,\qquad
    \iota_{\mathcal J}\colon\mathcal J\hookrightarrow\mathcal O_X,
    \qquad
    \varepsilon_q\colon
      q_{\mathcal I,\mathcal J}^*\mathcal O_X
      \xrightarrow{\ \sim\ }\mathcal O_{X_{\mathcal I,\mathcal J}}
  \]
  be the two inclusions and the pullback structure isomorphism.  Define
  \[
  \begin{gathered}
    \lambda_{\mathcal I}:=
      \varepsilon_q\circ q_{\mathcal I,\mathcal J}^*
        (\iota_{\mathcal I})\colon
      q_{\mathcal I,\mathcal J}^*\mathcal I
        \longrightarrow\mathcal O_{X_{\mathcal I,\mathcal J}},\\
    \lambda_{\mathcal J}:=
      \varepsilon_q\circ q_{\mathcal I,\mathcal J}^*
        (\iota_{\mathcal J})\colon
      q_{\mathcal I,\mathcal J}^*\mathcal J
        \longrightarrow\mathcal O_{X_{\mathcal I,\mathcal J}},\\
    \mathcal I':=\operatorname{im}(\lambda_{\mathcal I})
      =\mathcal I\mathcal O_{X_{\mathcal I,\mathcal J}},\qquad
    \mathcal J':=\operatorname{im}(\lambda_{\mathcal J})
      =\mathcal J\mathcal O_{X_{\mathcal I,\mathcal J}}.
  \end{gathered}                                               \tag{PCa}
  \]
  By \Cref{thm:nagata-blowup-principalizes-center}, both $\mathcal I'$ and
  $\mathcal J'$ are invertible; by
  \Cref{lemma:nagata-invertible-ideal-mul}, so is
  $\mathcal I'\mathcal J'=(\mathcal I\mathcal J)
  \mathcal O_{X_{\mathcal I,\mathcal J}}$.  Let
  $\iota_{\mathcal I'}\colon\mathcal I'\hookrightarrow
  \mathcal O_{X_{\mathcal I,\mathcal J}}$ and
  $\iota_{\mathcal J'}\colon\mathcal J'\hookrightarrow
  \mathcal O_{X_{\mathcal I,\mathcal J}}$ be the inclusions.  Define
  \[
  \begin{gathered}
    q_{\mathcal I'}\colon
      q_{\mathcal I,\mathcal J}^*\mathcal I\twoheadrightarrow\mathcal I',
    \qquad
    q_{\mathcal J'}\colon
      q_{\mathcal I,\mathcal J}^*\mathcal J\twoheadrightarrow\mathcal J',\\
    \mu_{\mathcal I,\mathcal J}\colon
      \mathcal I\otimes\mathcal J\twoheadrightarrow\mathcal I\mathcal J,
    \qquad
    \mu_{\mathcal I',\mathcal J'}\colon
      \mathcal I'\otimes\mathcal J'\twoheadrightarrow\mathcal I'\mathcal J'
  \end{gathered}
  \]
  by the equations
  \[
  \begin{gathered}
    \iota_{\mathcal I'}\circ q_{\mathcal I'}=\lambda_{\mathcal I},
    \qquad
    \iota_{\mathcal J'}\circ q_{\mathcal J'}=\lambda_{\mathcal J},\\
    \mu_{\mathcal I,\mathcal J}(x\otimes y):=xy
      \quad(x\in\mathcal I(U),\ y\in\mathcal J(U)),\\
    \mu_{\mathcal I',\mathcal J'}(x'\otimes y'):=x'y'
      \quad(x'\in\mathcal I'(U'),\ y'\in\mathcal J'(U')),
  \end{gathered}                                               \tag{PCb}
  \]
  for open $U\subseteq X$ and
  $U'\subseteq X_{\mathcal I,\mathcal J}$.
  Define
  $q^{\mathrm{prod}}_{\mathcal I,\mathcal J}$ uniquely by
  \[
  \begin{gathered}
    q^{\mathrm{prod}}_{\mathcal I,\mathcal J}\colon
      q_{\mathcal I,\mathcal J}^*(\mathcal I\mathcal J)
        \twoheadrightarrow\mathcal I'\mathcal J',\\
    q^{\mathrm{prod}}_{\mathcal I,\mathcal J}\circ
      q_{\mathcal I,\mathcal J}^*(\mu_{\mathcal I,\mathcal J})
      =\mu_{\mathcal I',\mathcal J'}\circ
        (q_{\mathcal I'}\otimes q_{\mathcal J'}).
  \end{gathered}
  \]
  The universal-lift construction of
  \Cref{def:nagata-blowup-universal-lift}, applied to
  $(q_{\mathcal I,\mathcal J},q^{\mathrm{prod}}_{\mathcal I,\mathcal J})$,
  is the typed definition
  \[
    \Phi_{\mathcal I,\mathcal J}
      :=\widetilde{q_{\mathcal I,\mathcal J}}^{
        (\mathcal I\mathcal J,q^{\mathrm{prod}}_{\mathcal I,\mathcal J})}
        \colon
      X_{\mathcal I,\mathcal J}
      \longrightarrow\operatorname{Bl}_{\mathcal I\mathcal J}(X)
  \]
  and \Cref{lemma:nagata-blowup-universal-exists} defines the equality
  \[
    e_\Phi:=e_{\widetilde{q_{\mathcal I,\mathcal J}}^{
      (\mathcal I\mathcal J,q^{\mathrm{prod}}_{\mathcal I,\mathcal J})}}
      \colon
      b_{X,\mathcal I\mathcal J}\circ\Phi_{\mathcal I,\mathcal J}
        =q_{\mathcal I,\mathcal J}.
  \]
  These symbols satisfy
  \[
    (\Phi_{\mathcal I,\mathcal J},e_{\Phi})
      \in\operatorname{LiftPair}
        (q_{\mathcal I,\mathcal J},
          q^{\mathrm{prod}}_{\mathcal I,\mathcal J}).
  \]

  Put $B_{\mathcal I\mathcal J}:=
  \operatorname{Bl}_{\mathcal I\mathcal J}(X)$ and
  $b_{\mathcal I\mathcal J}:=b_{X,\mathcal I\mathcal J}$.  By
  \Cref{thm:nagata-blowup-principalizes-center} and
  \Cref{lemma:nagata-invertible-product-factor}, the ideals
  \[
    \mathcal I'':=\mathcal I\mathcal O_{B_{\mathcal I\mathcal J}},
    \qquad
    \mathcal J'':=\mathcal J\mathcal O_{B_{\mathcal I\mathcal J}}
  \]
  are invertible.  Define
  \[
    \varepsilon_b\colon
      b_{\mathcal I\mathcal J}^*\mathcal O_X
        \xrightarrow{\ \sim\ }\mathcal O_{B_{\mathcal I\mathcal J}},
    \qquad
    \lambda_{\mathcal I}'':=
      \varepsilon_b\circ b_{\mathcal I\mathcal J}^*
        (\iota_{\mathcal I})\colon
      b_{\mathcal I\mathcal J}^*\mathcal I
        \longrightarrow\mathcal O_{B_{\mathcal I\mathcal J}}.
  \]
  Define the image quotient and its specified universal lift for
  $\mathcal I$ by
  \[
  \begin{gathered}
    q_{\mathcal I}''\colon
      b_{\mathcal I\mathcal J}^*\mathcal I
        \twoheadrightarrow\mathcal I'',\\
    \mathcal I''=\operatorname{im}(\lambda_{\mathcal I}''),\qquad
    \iota_{\mathcal I''}\circ q_{\mathcal I}''
      =\lambda_{\mathcal I}'',\\
    u_{\mathcal I}:=
      \widetilde{b_{\mathcal I\mathcal J}}^{
        (\mathcal I,q_{\mathcal I}'')}
      \colon B_{\mathcal I\mathcal J}\longrightarrow B_{\mathcal I},
  \end{gathered}
  \]
  where
  $\iota_{\mathcal I''}\colon\mathcal I''\hookrightarrow
  \mathcal O_{B_{\mathcal I\mathcal J}}$ is the inclusion and
  the last line is exactly the formula of
  \Cref{def:nagata-blowup-universal-lift} with
  $(f,\mathcal I,q_f)=(b_{\mathcal I\mathcal J},\mathcal I,q_{\mathcal I}'')$.
  By \Cref{lemma:nagata-blowup-universal-exists}, define
  \[
  \begin{gathered}
    e_{\mathcal I}\colon
      b_{\mathcal I}\circ u_{\mathcal I}
        =b_{\mathcal I\mathcal J},\\
    (u_{\mathcal I},e_{\mathcal I})
      \in\operatorname{LiftPair}
        (b_{\mathcal I\mathcal J},q_{\mathcal I}'').
  \end{gathered}                                             \tag{BPC0}
  \]
  Let
  \[
    \iota_{\mathcal J_{\mathcal I}}\colon
      \mathcal J_{\mathcal I}\hookrightarrow\mathcal O_{B_{\mathcal I}},
    \qquad
    \varepsilon_{u_{\mathcal I}}\colon
      u_{\mathcal I}^*\mathcal O_{B_{\mathcal I}}
      \xrightarrow{\ \sim\ }\mathcal O_{B_{\mathcal I\mathcal J}}.
  \]
  Define the second-stage pullback map, its image quotient, and its specified
  universal lift by
  \[
  \begin{gathered}
    \mu_{\mathcal J|\mathcal I}'':=
      \varepsilon_{u_{\mathcal I}}\circ
        u_{\mathcal I}^*(\iota_{\mathcal J_{\mathcal I}})\colon
      u_{\mathcal I}^*\mathcal J_{\mathcal I}
        \longrightarrow\mathcal O_{B_{\mathcal I\mathcal J}},\\
    \mathcal J''=
      \operatorname{im}(\mu_{\mathcal J|\mathcal I}''),
    \qquad
    q_{\mathcal J|\mathcal I}''\colon
      u_{\mathcal I}^*\mathcal J_{\mathcal I}
        \twoheadrightarrow\mathcal J'',\\
    u_{\mathcal J|\mathcal I}:=
      \widetilde{u_{\mathcal I}}^{
        (\mathcal J_{\mathcal I},q_{\mathcal J|\mathcal I}'')}
      \colon B_{\mathcal I\mathcal J}\longrightarrow
        X_{\mathcal I,\mathcal J},
  \end{gathered}
  \]
  If
  $\iota_{\mathcal J''}\colon\mathcal J''\hookrightarrow
  \mathcal O_{B_{\mathcal I\mathcal J}}$ is the inclusion, then the image
  factorization is the equality
  \[
    \mu_{\mathcal J|\mathcal I}''
      =\iota_{\mathcal J''}\circ q_{\mathcal J|\mathcal I}''. \tag{PCc}
  \]
  The last line of the preceding display is
  \Cref{def:nagata-blowup-universal-lift} with
  $(f,\mathcal I,q_f)=(u_{\mathcal I},\mathcal J_{\mathcal I},
  q_{\mathcal J|\mathcal I}'')$.  Define
  \[
  \begin{gathered}
    e_{\mathcal J|\mathcal I}\colon
      b_{\mathcal J|\mathcal I}\circ u_{\mathcal J|\mathcal I}
        =u_{\mathcal I},\\
    (u_{\mathcal J|\mathcal I},e_{\mathcal J|\mathcal I})
      \in\operatorname{LiftPair}
        (u_{\mathcal I},q_{\mathcal J|\mathcal I}'').
  \end{gathered}                                             \tag{BPC1}
  \]
  Define
  \[
    \Psi_{\mathcal I,\mathcal J}:=
      u_{\mathcal J|\mathcal I}\colon
      B_{\mathcal I\mathcal J}\longrightarrow X_{\mathcal I,\mathcal J}.
  \]
\end{definition}

\begin{theorem}[{\cite[\href{https://stacks.math.columbia.edu/tag/080A}{Tag 080A}]{stacks-project}}]
  \label{thm:nagata-blowup-product-ideal}
  \uses{def:nagata-blowup-product-comparison}
  \lean{}

  The maps $\Phi_{\mathcal I,\mathcal J}$ and
  $\Psi_{\mathcal I,\mathcal J}$ satisfy
  \[
    \Phi_{\mathcal I,\mathcal J}\circ\Psi_{\mathcal I,\mathcal J}
      =\operatorname{id}_{B_{\mathcal I\mathcal J}},
    \qquad
    \Psi_{\mathcal I,\mathcal J}\circ\Phi_{\mathcal I,\mathcal J}
      =\operatorname{id}_{X_{\mathcal I,\mathcal J}}.
  \]
\end{theorem}

\begin{proof}
  \uses{thm:nagata-blowup-universal-property,
    lemma:nagata-image-quotient-transport,
    lemma:nagata-invertible-ideal-mul,
    lemma:nagata-invertible-product-factor}
  The construction in
  \Cref{def:nagata-blowup-product-comparison} and
  \Cref{lemma:nagata-blowup-universal-exists} consists of the equalities
  \[
    b_{\mathcal I\mathcal J}\circ\Phi_{\mathcal I,\mathcal J}
      =q_{\mathcal I,\mathcal J},
    \qquad
    q_{\mathcal I,\mathcal J}\circ\Psi_{\mathcal I,\mathcal J}
      =b_{\mathcal I\mathcal J}.
  \]
  Define the universal quotient
  \[
    q_{\mathcal I\mathcal J}^{\mathrm{univ}}\colon
    b_{\mathcal I\mathcal J}^*(\mathcal I\mathcal J)
      \twoheadrightarrow
      (\mathcal I\mathcal J)\mathcal O_{B_{\mathcal I\mathcal J}}.
  \]
  Define
  \[
  \begin{gathered}
    H_0^{\mathrm{prod}}
      :=\Phi_{\mathcal I,\mathcal J}\circ\Psi_{\mathcal I,\mathcal J},
    \qquad
    H_1^{\mathrm{prod}}:=\operatorname{id}_{B_{\mathcal I\mathcal J}},\\
    e_0^{\mathrm{prod}}\colon
      b_{\mathcal I\mathcal J}\circ H_0^{\mathrm{prod}}
      =q_{\mathcal I,\mathcal J}\circ\Psi_{\mathcal I,\mathcal J}
      =b_{\mathcal I\mathcal J},\\
    e_1^{\mathrm{prod}}\colon
      b_{\mathcal I\mathcal J}\circ H_1^{\mathrm{prod}}
      =b_{\mathcal I\mathcal J}.
  \end{gathered}                                                   \tag{PI0}
  \]
  By \Cref{lemma:nagata-image-quotient-transport}, define, for $i=0,1$,
  \[
    \alpha_i^{\mathrm{prod}}
      :=\operatorname{iqt}_{\mathcal I\mathcal J}
          (e_i^{\mathrm{prod}})\colon
      (H_i^{\mathrm{prod}})^*
        ((\mathcal I\mathcal J)\mathcal O_{B_{\mathcal I\mathcal J}})
      \xrightarrow{\sim}
        (\mathcal I\mathcal J)\mathcal O_{B_{\mathcal I\mathcal J}}.
  \]
  Equation (IQ2), with the associator of
  \Cref{def:nagata-blowup-lift-set}, is the pair of quotient equalities
  \[
    \alpha_i^{\mathrm{prod}}\circ
      (H_i^{\mathrm{prod}})^*(q_{\mathcal I\mathcal J}^{\mathrm{univ}})
      =q_{\mathcal I\mathcal J}^{\mathrm{univ}}\circ
        \vartheta_{H_i^{\mathrm{prod}},e_i^{\mathrm{prod}}}
      \qquad(i=0,1).                                               \tag{PI1}
  \]
  Substitution of (PI0)--(PI1) in
  \Cref{def:nagata-blowup-lift-set} is
  \[
  \begin{gathered}
    (H_i^{\mathrm{prod}},e_i^{\mathrm{prod}})
      \in\operatorname{LiftPair}
        (b_{\mathcal I\mathcal J},q_{\mathcal I\mathcal J}^{\mathrm{univ}}),\\
    H_i^{\mathrm{prod}}
      \in\operatorname{Lift}
        (b_{\mathcal I\mathcal J},q_{\mathcal I\mathcal J}^{\mathrm{univ}})
      \qquad(i=0,1).
  \end{gathered}                                                   \tag{PI2}
  \]
  The singleton equality in
  \Cref{thm:nagata-blowup-universal-property} is
  \[
    \Phi_{\mathcal I,\mathcal J}\circ\Psi_{\mathcal I,\mathcal J}
      =\operatorname{id}_{B_{\mathcal I\mathcal J}}.
  \]

  Put
  \[
    H:=\Psi_{\mathcal I,\mathcal J}\circ
      \Phi_{\mathcal I,\mathcal J}.
  \]
  The two structural equalities at the start of the proof are the chain
  \[
  \begin{aligned}
    q_{\mathcal I,\mathcal J}\circ H
      &=q_{\mathcal I,\mathcal J}\circ
          \Psi_{\mathcal I,\mathcal J}\circ
          \Phi_{\mathcal I,\mathcal J}\\
      &=b_{\mathcal I\mathcal J}\circ
          \Phi_{\mathcal I,\mathcal J}
       =q_{\mathcal I,\mathcal J}.
  \end{aligned}                                                    \tag{PI3}
  \]
  Define
  \[
  \begin{gathered}
    q_{\mathcal I}^{\mathrm{univ}}\colon
      b_{\mathcal I}^*\mathcal I
        \twoheadrightarrow\mathcal I\mathcal O_{B_{\mathcal I}},\\
    q_{\mathcal I}'\colon
      q_{\mathcal I,\mathcal J}^*\mathcal I
        \twoheadrightarrow\mathcal I',\\
    K_0:=b_{\mathcal J|\mathcal I}\circ H,
    \qquad K_1:=b_{\mathcal J|\mathcal I},\\
    e_0^{\mathcal I}\colon
      b_{\mathcal I}\circ K_0
       =q_{\mathcal I,\mathcal J}\circ H
       =q_{\mathcal I,\mathcal J},\qquad
    e_1^{\mathcal I}\colon
      b_{\mathcal I}\circ K_1=q_{\mathcal I,\mathcal J}.
  \end{gathered}                                                    \tag{PI4}
  \]
  For $i=0,1$, define
  \[
    \alpha_i^{\mathcal I}
      :=\operatorname{iqt}_{\mathcal I}(e_i^{\mathcal I})\colon
      K_i^*(\mathcal I\mathcal O_{B_{\mathcal I}})
        \xrightarrow{\sim}\mathcal I'.
  \]
  Equation (IQ2) and (PI4) are
  \[
    \alpha_i^{\mathcal I}\circ K_i^*(q_{\mathcal I}^{\mathrm{univ}})
      =q_{\mathcal I}'\circ\vartheta_{K_i,e_i^{\mathcal I}},
    \qquad
    K_i\in\operatorname{Lift}
      (q_{\mathcal I,\mathcal J},q_{\mathcal I}')\quad(i=0,1).   \tag{PI5}
  \]
  The singleton equality for the $\mathcal I$-blow-up is
  \[
    b_{\mathcal J|\mathcal I}\circ H=b_{\mathcal J|\mathcal I}.
  \]
  Define the universal second-stage quotient
  \[
    q_{\mathcal J|\mathcal I}\colon
      b_{\mathcal J|\mathcal I}^*\mathcal J_{\mathcal I}
      \twoheadrightarrow\mathcal J'.
  \]
  Define
  \[
  \begin{gathered}
    H_0^{\mathcal J}:=H,\qquad
    H_1^{\mathcal J}:=\operatorname{id}_{X_{\mathcal I,\mathcal J}},\\
    e_0^{\mathcal J}\colon
      b_{\mathcal J|\mathcal I}\circ H_0^{\mathcal J}
        =b_{\mathcal J|\mathcal I},\qquad
    e_1^{\mathcal J}\colon
      b_{\mathcal J|\mathcal I}\circ H_1^{\mathcal J}
        =b_{\mathcal J|\mathcal I},\\
    \alpha_i^{\mathcal J}
      :=\operatorname{iqt}_{\mathcal J_{\mathcal I}}(e_i^{\mathcal J})
        \colon(H_i^{\mathcal J})^*\mathcal J'
          \xrightarrow{\sim}\mathcal J'\quad(i=0,1).
  \end{gathered}                                                    \tag{PI6}
  \]
  Equation (IQ2) is
  \[
    \alpha_i^{\mathcal J}\circ
      (H_i^{\mathcal J})^*(q_{\mathcal J|\mathcal I})
      =q_{\mathcal J|\mathcal I}\circ
        \vartheta_{H_i^{\mathcal J},e_i^{\mathcal J}},
    \qquad
    H_i^{\mathcal J}\in
      \operatorname{Lift}(b_{\mathcal J|\mathcal I},
        q_{\mathcal J|\mathcal I})\quad(i=0,1).                  \tag{PI7}
  \]
  The singleton equality of
  \Cref{thm:nagata-blowup-universal-property} is
  \[
    \Psi_{\mathcal I,\mathcal J}\circ\Phi_{\mathcal I,\mathcal J}
      =\operatorname{id}_{X_{\mathcal I,\mathcal J}}.
  \]
\end{proof}

\begin{definition}[Domination maps from the product blow-up]
  \label{def:nagata-blowup-product-domination-maps}
  \uses{def:nagata-blowup-product-comparison}
  \lean{}

  With the notation of
  \Cref{def:nagata-blowup-product-comparison}, define
  \[
    d_{\mathcal I}:=
      b_{\mathcal J|\mathcal I}\circ\Psi_{\mathcal I,\mathcal J}\colon
      B_{\mathcal I\mathcal J}\longrightarrow B_{\mathcal I}.
  \]
  Interchanging $\mathcal I$ and $\mathcal J$ defines
  \[
    d_{\mathcal J}:=
      b_{\mathcal I|\mathcal J}\circ\Psi_{\mathcal J,\mathcal I}\colon
      B_{\mathcal I\mathcal J}\longrightarrow B_{\mathcal J}.
  \]
\end{definition}

\begin{corollary}[Common refinement by the product ideal]
  \label{cor:nagata-blowup-product-common-refinement}
  \uses{def:nagata-blowup-product-domination-maps}
  \lean{}

  The morphisms $d_{\mathcal I}$ and $d_{\mathcal J}$ of
  \Cref{def:nagata-blowup-product-domination-maps} are over
  $X$.  If both ideals restrict to the unit ideal on an open $U$, the product
  blow-up and both maps restrict to the identity over $U$.
\end{corollary}

\begin{proof}
  \uses{thm:nagata-blowup-product-ideal,
    thm:nagata-global-blowup-complement}
  By \Cref{def:nagata-blowup-product-comparison} and
  \Cref{lemma:nagata-blowup-universal-exists},
  \[
    b_{\mathcal I}\circ d_{\mathcal I}
      =b_{\mathcal I}\circ b_{\mathcal J|\mathcal I}
        \circ\Psi_{\mathcal I,\mathcal J}
      =b_{\mathcal I\mathcal J},
  \]
  and, after interchanging the two ideals,
  \[
    b_{\mathcal J}\circ d_{\mathcal J}=b_{\mathcal I\mathcal J}.
  \]
  These two equalities state that both maps are over $X$.

  Assume $\mathcal I|_U=\mathcal O_U$ and
  $\mathcal J|_U=\mathcal O_U$.  Then
  $(\mathcal I\mathcal J)|_U=\mathcal O_U$, and
  \Cref{thm:nagata-global-blowup-complement} supplies isomorphisms
  \[
    B_{\mathcal I}|_U\xrightarrow{\sim}U,\qquad
    B_{\mathcal J}|_U\xrightarrow{\sim}U,\qquad
    B_{\mathcal I\mathcal J}|_U\xrightarrow{\sim}U.
  \]
  Under these isomorphisms, $d_{\mathcal I}|_U$ is a $U$-morphism
  $U\to U$; the equation
  $b_{\mathcal I}\circ d_{\mathcal I}=b_{\mathcal I\mathcal J}$ becomes
  \[
    \operatorname{id}_U\circ d_{\mathcal I}|_U=\operatorname{id}_U,
  \]
  so $d_{\mathcal I}|_U=\operatorname{id}_U$.  Interchanging
  $\mathcal I$ and $\mathcal J$ in the same chain is the equality
  \[
    d_{\mathcal J}|_U=\operatorname{id}_U.
  \]
\end{proof}


\subsection{An arbitrary composition of two blow-ups}

The product-ideal theorem treats a second center pulled back from the original
scheme.  Stability of admissible blow-ups under composition needs the more
general statement in which the second center is an arbitrary finite-presentation
closed subscheme of the first blow-up.

\begin{definition}[Finite-degree data attached to a closed center in Proj]
  \label{def:nagata-proj-closed-center-finite-degree}
  \uses{def:nagata-choice-notation,
    def:nagata-closed-subscheme-class,
    def:nagata-relative-proj,
    lemma:nagata-relative-proj-interface}
  \lean{}

  Let $X$ be quasi-compact and quasi-separated, let
  $\mathcal I\subseteq\mathcal O_X$ be a finite-type quasi-coherent ideal,
  and define
  \[
    \mathcal A:=\mathcal R_X(\mathcal I)
      =\bigoplus_{n\ge0}\mathcal I^n,\qquad
    X':=\operatorname{Proj}_X(\mathcal A),\qquad
    b:=b_{X,\mathcal I}\colon X'\to X,\qquad
    Z:=V(\mathcal I).
  \]
  Let $i'\colon Z'\hookrightarrow X'$ be a closed immersion of finite
  presentation and put
  \[
    U(Z'):=X\setminus\bigl(|Z|\cup |b|(|Z'|)\bigr).
  \]
  For every quasi-coherent $\mathcal O_X$-module $\mathcal M$, define
  \[
    \operatorname{Supp}_X(\mathcal M)
      :=\{x\in|X|\mid\mathcal M_x\neq0\},
    \qquad
    \operatorname{Supp}(\mathcal M):=
      \operatorname{Supp}_X(\mathcal M).                           \tag{FD0}
  \]
  For $d>0$ and
  $\mathcal F\subseteq\mathcal I^d$, define
  \[
    q_{d,\mathcal F}\colon\mathcal A\twoheadrightarrow
      \mathcal A/\mathcal F\mathcal A,
    \qquad
    \kappa_{d,\mathcal F}:=
      \operatorname{Proj}_X(q_{d,\mathcal F})\colon
      \operatorname{Proj}_X(\mathcal A/\mathcal F\mathcal A)
        \hookrightarrow X'.                                      \tag{FD1}
  \]
  Define $\mathcal P(Z')$ to be the set of pairs $(d,\mathcal F)$ satisfying
  \[
  \begin{aligned}
    &d\in\mathbb N,\quad d>0,\\
    &\mathcal F\subseteq\mathcal I^d
      \text{ is a finite-type quasi-coherent }\mathcal O_X\text{-submodule},\\
    &[\,Z'\xhookrightarrow{i'}X'\,]
      =[\,\operatorname{Proj}_X
        (\mathcal A/\mathcal F\mathcal A)
        \xhookrightarrow{\kappa_{d,\mathcal F}}X'\,]
      \quad\text{in }\operatorname{ClosedSub}(X'),\\
    &\operatorname{Supp}(\mathcal I^d/\mathcal F)\cap U(Z')=\varnothing.
  \end{aligned}
  \]
  By \cite[Tag 0803]{stacks-project}, $\mathcal P(Z')\ne\varnothing$.  By
  \Cref{def:nagata-choice-notation}, define
  \[
    \mathsf{FD}(Z'):=
      \operatorname{ch}(\mathcal P(Z'))
      =\bigl(d(Z'),\mathcal F(Z')\bigr).
  \]
\end{definition}

\begin{lemma}[Properties of the finite-degree center data]
  \label{lemma:nagata-proj-closed-center-finite-degree}
  \uses{def:nagata-proj-closed-center-finite-degree}
  \lean{}

  Let $X,\mathcal I,\mathcal A,X',b,Z,i',Z'$ be the datum of
  \Cref{def:nagata-proj-closed-center-finite-degree}.
  The integer $d(Z')$ and module $\mathcal F(Z')$ of
  \Cref{def:nagata-proj-closed-center-finite-degree} satisfy
  $d(Z')>0$, $\mathcal F(Z')$ is a finite
  type submodule of $\mathcal I^{d(Z')}$, and
  \[
    [\,Z'\xhookrightarrow{i'}X'\,]
      =[\,\operatorname{Proj}_X
        (\mathcal A/\mathcal F(Z')\mathcal A)
        \xhookrightarrow{\kappa_{d(Z'),\mathcal F(Z')}}X'\,]
      \quad\text{in }\operatorname{ClosedSub}(X').
  \]
  Moreover
  \[
    \operatorname{Supp}
      (\mathcal I^{d(Z')}/\mathcal F(Z'))
      \subseteq |Z|\cup |b|(|Z'|).
  \]
\end{lemma}

\begin{proof}
  First, $b$ is proper by
  \Cref{thm:nagata-global-blowup-proper}.  Since $i'$ is a closed immersion,
  the composite $b\circ i'\colon Z'\to X$ is proper.  Universal closedness
  of this proper morphism is the assertion that $|b|(|Z'|)$ is closed.  The
  complement formula is
  \[
    U(Z')=X\setminus\bigl(|Z|\cup |b|(|Z'|)\bigr),
  \]
  so $U(Z')$ is open.  The open $X\setminus Z$ is quasi-compact
  because $\mathcal I$ is of finite type, and under the isomorphism
  \[
    b^{-1}(X\setminus Z)\xrightarrow{\sim}X\setminus Z
  \]
  of \Cref{thm:nagata-global-blowup-complement}, the open
  $U(Z')$ corresponds to
  $b^{-1}(X\setminus Z)\setminus Z'$.  Since $i'$ is of finite presentation,
  that complement is quasi-compact.  These facts verify the hypotheses of
  \cite[Tag 0803]{stacks-project}, which supplies
  $\mathcal P(Z')\ne\varnothing$.

  By the defining property of the choice function,
  \[
    (d(Z'),\mathcal F(Z'))\in\mathcal P(Z').
  \]
  Expanding membership in the set defined in
  \Cref{def:nagata-proj-closed-center-finite-degree} is the list
  \[
    d(Z')>0,\qquad
    \mathcal F(Z')\subseteq\mathcal I^{d(Z')}
      \text{ finite type},
  \]
  \[
    [\,Z'\xhookrightarrow{i'}X'\,]
      =[\,\operatorname{Proj}_X
        \bigl(\mathcal A/\mathcal F(Z')\mathcal A\bigr)
        \xhookrightarrow{\kappa_{d(Z'),\mathcal F(Z')}}X'\,]
      \quad\text{in }\operatorname{ClosedSub}(X'),
  \]
  and
  \[
    \operatorname{Supp}
      (\mathcal I^{d(Z')}/\mathcal F(Z'))
      \cap U(Z')=\varnothing.
  \]
  Since $X\setminus U(Z')=|Z|\cup |b|(|Z'|)$, the last equality is equivalent
  to the asserted support inclusion.
\end{proof}

\begin{definition}[Center ideal representing a second blow-up]
  \label{def:nagata-composition-center-ideal}
  \uses{def:nagata-proj-closed-center-finite-degree}
  \lean{}

  With $\mathsf{FD}(Z')=(d,\mathcal F)$ as in
  \Cref{def:nagata-proj-closed-center-finite-degree}, let
  \[
    \jmath_{\mathcal F}\colon\mathcal F\hookrightarrow\mathcal I^d,
    \qquad
    \jmath_{\mathcal I^d}\colon\mathcal I^d\hookrightarrow\mathcal O_X,
    \qquad
    \iota_{\mathcal F}:=
      \jmath_{\mathcal I^d}\circ\jmath_{\mathcal F}\colon
      \mathcal F\hookrightarrow\mathcal O_X
  \]
  be the two inclusions and their explicitly displayed composite.  Define the
  quasi-coherent ideals
  \[
    \mathcal J(Z'):=\operatorname{im}(\iota_{\mathcal F})
      =\mathcal F\subseteq\mathcal O_X,
    \qquad
    \mathcal K(Z'):=\mathcal I\mathcal J(Z'),
    \qquad
    Y(Z'):=V(\mathcal K(Z'))\hookrightarrow X.
  \]
\end{definition}

\begin{lemma}[Support of the constructed composition center]
  \label{lemma:nagata-composition-center-ideal}
  \uses{def:nagata-composition-center-ideal}
  \lean{}

  The ideal $\mathcal K(Z')$ is finite type and
  \[
    |Y(Z')|=|V(\mathcal I)|\cup |b|(|Z'|).
  \]
\end{lemma}

\begin{proof}
  \uses{def:nagata-closed-subscheme-class,
    lemma:nagata-proj-closed-center-finite-degree,
    thm:nagata-global-blowup-complement}
  Write $d:=d(Z')$, $\mathcal F:=\mathcal F(Z')$,
  $\mathcal J:=\mathcal J(Z')$, and $U:=U(Z')$.  By
  \Cref{lemma:nagata-proj-closed-center-finite-degree},
  \[
    \operatorname{Supp}(\mathcal I^d/\mathcal J)\cap U=\varnothing.
  \]
  Since $U\subseteq X\setminus V(\mathcal I)$,
  \Cref{lemma:nagata-ideal-unit-off-support} is the equality
  $\mathcal I|_U=\mathcal O_U$, so
  $\mathcal I^d|_U=\mathcal O_U$.  The support equality is the chain
  \[
    \operatorname{Supp}(\mathcal I^d/\mathcal J)\cap U=\varnothing
      \Longrightarrow
    (\mathcal I^d/\mathcal J)|_U=0
      \Longrightarrow
    \mathcal J|_U=\mathcal I^d|_U.
  \]
  Substitution of $\mathcal I^d|_U=\mathcal O_U$ is
  \[
    \mathcal J|_U=\mathcal I^d|_U=\mathcal O_U,
    \qquad
    V(\mathcal J)\cap U=\varnothing,
  \]
  Taking complements in the equality
  $V(\mathcal J)\cap U=\varnothing$ is the inclusion
  \[
    |V(\mathcal J)|\subseteq X\setminus U
      =|V(\mathcal I)|\cup |b|(|Z'|).
  \]
  Since $\mathcal J\subseteq\mathcal I^d$,
  \[
    |V(\mathcal I)|=|V(\mathcal I^d)|
      \subseteq|V(\mathcal J)|.
  \]

  Put $X_0:=X\setminus V(\mathcal I)$ and
  $X'_0:=b^{-1}(X_0)$.  By
  \Cref{thm:nagata-global-blowup-complement},
  $\varphi:=b|_{X'_0}\colon X'_0\xrightarrow{\sim}X_0$.  Define
  \[
  \begin{gathered}
    Z'_0:=Z'\times_{X'}X'_0,
    \qquad i'_0:=\operatorname{pr}_{X'_0}\colon Z'_0\hookrightarrow X'_0,\\
    P_{\mathcal F,0}:=
      \operatorname{Proj}_{X_0}\!\left(
        \bigoplus_{n\ge0}\mathcal O_{X_0}t^n
          \big/ (\mathcal J|_{X_0})t^d\right),\\
    \kappa_{\mathcal F,0}\colon P_{\mathcal F,0}\hookrightarrow X'_0
      \quad\text{the base change of }
      \kappa_{d,\mathcal F}\colon
        \operatorname{Proj}_X(\mathcal A/\mathcal F\mathcal A)
          \hookrightarrow X',\\
    V'_{\mathcal J,0}:=
      X'_0\times_{X_0,\varphi}V(\mathcal J|_{X_0}),
    \qquad
    \nu_{\mathcal J,0}:=\operatorname{pr}_{X'_0}\colon
      V'_{\mathcal J,0}\hookrightarrow X'_0.
  \end{gathered}                                                   \tag{CCI0}
  \]
  Restricting the class equality of
  \Cref{lemma:nagata-proj-closed-center-finite-degree} to $X'_0$, and using
  $\mathcal I^d|_{X_0}=\mathcal O_{X_0}$, gives the typed chain
  \[
  \begin{aligned}
    [\,Z'_0\xhookrightarrow{i'_0}X'_0\,]
      &=[\,P_{\mathcal F,0}
        \xhookrightarrow{\kappa_{\mathcal F,0}}X'_0\,]\\
      &=[\,V'_{\mathcal J,0}
        \xhookrightarrow{\nu_{\mathcal J,0}}X'_0\,]
      \quad\text{in }\operatorname{ClosedSub}(X'_0).
  \end{aligned}
  \]
  Equality of the embedded supports in the first and last classes is
  \[
    |\varphi|\!\left(|i'_0|(|Z'_0|)\right)
      =|V(\mathcal J)|\cap|X_0|.
  \]
  On underlying sets this is
  \[
    |b|(|Z'|)\setminus|V(\mathcal I)|
      =|V(\mathcal J)|\setminus|V(\mathcal I)|.
  \]
  Set-theoretic union of the equality off $V(\mathcal I)$ with the two
  preceding inclusions is the equality
  \[
    |V(\mathcal J)|
      =|V(\mathcal I)|\cup |b|(|Z'|).
  \]
  For ideal sheaves, $V(\mathcal I\mathcal J)=
  V(\mathcal I)\cup V(\mathcal J)$; consequently
  \[
    |Y(Z')|=|V(\mathcal K(Z'))|
      =|V(\mathcal I)|\cup |b|(|Z'|).
  \]
  Finally, choose local finite presentations
  \[
    \mathcal I|_V=(a_1,\ldots,a_r),
    \qquad
    \mathcal J|_V=(c_1,\ldots,c_s)
  \]
  on each member $V$ of a finite affine cover.  Then
  \[
    \mathcal K(Z')|_V
      =\mathcal I\mathcal J|_V
      =(a_ic_j\mid1\leq i\leq r,\ 1\leq j\leq s),
  \]
  which is the finite-type assertion.
\end{proof}


\begin{lemma}[Comparison of the two second-center ideals]
  \label{lemma:nagata-composition-extended-ideal-equality}
  \lean{}

  Retain $b\colon X'\to X$, $i'\colon Z'\hookrightarrow X'$,
  $d:=d(Z')$, and $\mathcal J:=\mathcal J(Z')$ from
  \Cref{def:nagata-proj-closed-center-finite-degree} and
  \Cref{def:nagata-composition-center-ideal}.  Define
  \[
    \iota_{\mathcal I}\colon\mathcal I\hookrightarrow\mathcal O_X,
    \qquad
    m_b\colon b^*\mathcal O_X\to\mathcal O_{X'},\quad
      r\otimes a\longmapsto r\,b^\sharp(a),
    \qquad
    \mu_b:=m_b\circ b^*(\iota_{\mathcal I})\colon
      b^*\mathcal I\longrightarrow\mathcal O_{X'},
    \qquad
    \mathcal I_E:=\mathcal I\mathcal O_{X'}
      :=\operatorname{im}(\mu_b),
    \qquad
    \varepsilon_{i'}:=(i')^\flat\colon
      \mathcal O_{X'}\to i'_*\mathcal O_{Z'},
    \qquad
    \mathcal I_{Z'}:=\ker(\varepsilon_{i'}).
  \]
  Then, as ideal sheaves on $X'$,
  \[
    \mathcal J\mathcal O_{X'}
      =\mathcal I_E^d\mathcal I_{Z'}.
  \]
\end{lemma}

\begin{proof}
  \uses{lemma:nagata-composition-center-ideal,
    lemma:nagata-chart-extended-ideal-principal}
  Choose an affine open $V=\operatorname{Spec}(A)\subseteq X$, put
  $I:=\Gamma(V,\mathcal I)$ and
  $M:=\Gamma(V,\mathcal F(Z'))\subseteq I^d$, and choose $a\in I$.
  Let
  \[
    C_a:=\operatorname{Spec}(B_a),\qquad B_a:=A[I/a]
  \]
  be the chart of \Cref{prop:nagata-affine-blowup-chart}.  The Proj equality
  in \Cref{lemma:nagata-proj-closed-center-finite-degree} says
  \[
    \Gamma(C_a,\mathcal I_{Z'})
      =\left(\frac m{a^d}\ \middle|\ m\in M\right)B_a.
  \]
  By \Cref{lemma:nagata-chart-extended-ideal-principal},
  \[
    \Gamma(C_a,\mathcal I_E)=(a)B_a,
    \qquad
    \Gamma(C_a,\mathcal I_E^d)=(a^d)B_a.
  \]
  The ideal $\mathcal J\mathcal O_{X'}$ restricts to the extension
  $MB_a$, and for every $m\in M$ the equality in $B_a\subseteq A_a$ is
  \[
    m=a^d\left(\frac m{a^d}\right).
  \]
  Substitution of $m=a^d(m/a^d)$ for every generator $m\in M$ is
  \[
  \begin{aligned}
    \Gamma(C_a,\mathcal J\mathcal O_{X'})
      &= (m\mid m\in M)B_a\\
      &= \left(a^d\frac m{a^d}\ \middle|\ m\in M\right)B_a\\
      &= (a^d)B_a
        \left(\frac m{a^d}\ \middle|\ m\in M\right)B_a\\
      &=\Gamma(C_a,\mathcal I_E^d\mathcal I_{Z'}).
  \end{aligned}
  \]
  The charts $C_a$ cover $X'$ by
  \Cref{prop:nagata-affine-blowup-chart-cover}; equality of ideal sheaves is
  local, so
  $\mathcal J\mathcal O_{X'}=\mathcal I_E^d\mathcal I_{Z'}$.
\end{proof}

\begin{definition}[Comparison for a composite of arbitrary blow-ups]
  \label{def:nagata-composition-of-blowups-comparison}
  \uses{def:nagata-composition-center-ideal,
    def:nagata-blowup-product-comparison,
    def:nagata-global-blowup}
  \lean{}

  Retain the notation of
  \Cref{def:nagata-proj-closed-center-finite-degree},
  \Cref{def:nagata-composition-center-ideal}, and
  \Cref{lemma:nagata-composition-extended-ideal-equality}.  Put
  \[
    X'':=\operatorname{Bl}_{\mathcal I_{Z'}}(X'),\qquad
    b':=b_{X',\mathcal I_{Z'}}\colon X''\to X',
  \]
  \[
    C:=\operatorname{Bl}_{\mathcal J\mathcal O_{X'}}(X'),\qquad
    c:=b_{X',\mathcal J\mathcal O_{X'}}\colon C\to X',
    \qquad
    B:=\operatorname{Bl}_{\mathcal K(Z')}(X),\qquad
    k:=b_{X,\mathcal K(Z')}\colon B\to X.
  \]
  Since $\mathcal I_E^d$ is invertible by
  \Cref{thm:nagata-blowup-principalizes-center}, the ideal
  \[
    (b')^*(\mathcal J\mathcal O_{X'})
      = (b')^*\mathcal I_E^d\cdot
        \mathcal I_{Z'}\mathcal O_{X''}
  \]
  is invertible.  Let
  \[
  \begin{gathered}
    \iota_{\mathcal J,X'}\colon
      \mathcal J\mathcal O_{X'}\hookrightarrow\mathcal O_{X'},\qquad
    \varepsilon_{b'}\colon(b')^*\mathcal O_{X'}
      \xrightarrow{\ \sim\ }\mathcal O_{X''},\\
    \lambda_{\mathcal J,X''}:=
      \varepsilon_{b'}\circ(b')^*(\iota_{\mathcal J,X'})
      \colon(b')^*(\mathcal J\mathcal O_{X'})
        \longrightarrow\mathcal O_{X''}.
  \end{gathered}                                               \tag{CBa}
  \]
  Define its image quotient, the specified universal lift,
  and its factorization equality by
  \[
  \begin{gathered}
    q_{\mathcal J,X''}\colon
      (b')^*(\mathcal J\mathcal O_{X'})
        \twoheadrightarrow
        (\mathcal J\mathcal O_{X'})\mathcal O_{X''},\\
    (\mathcal J\mathcal O_{X'})\mathcal O_{X''}
      =\operatorname{im}(\lambda_{\mathcal J,X''}),\\
    \iota_{\mathcal J,X''}\circ q_{\mathcal J,X''}
      =\lambda_{\mathcal J,X''},\\
    \kappa:=
      \widetilde{b'}^{
        (\mathcal J\mathcal O_{X'},q_{\mathcal J,X''})}
      \colon X''\longrightarrow C,\\
    e_\kappa\colon c\circ\kappa=b',\qquad
    (\kappa,e_\kappa)\in
      \operatorname{LiftPair}(b',q_{\mathcal J,X''}).
  \end{gathered}                                             \tag{CB0}
  \]
  where
  $\iota_{\mathcal J,X''}\colon
  (\mathcal J\mathcal O_{X'})\mathcal O_{X''}
  \hookrightarrow\mathcal O_{X''}$ is the inclusion.  The formula for
  $\kappa$ is
  \Cref{def:nagata-blowup-universal-lift} with
  $(f,\mathcal I,q_f)=(b',\mathcal J\mathcal O_{X'},
  q_{\mathcal J,X''})$, and the last line is
  \Cref{lemma:nagata-blowup-universal-exists}.  On $C$, the product
  $\mathcal J\mathcal O_C$ is invertible; by
  \Cref{lemma:nagata-invertible-product-factor} and
  $\mathcal J\mathcal O_{X'}=\mathcal I_E^d\mathcal I_{Z'}$, the ideal
  $\mathcal I_{Z'}\mathcal O_C$ is invertible.  Let
  \[
  \begin{gathered}
    \iota_{Z'}\colon\mathcal I_{Z'}\hookrightarrow\mathcal O_{X'},\qquad
    \varepsilon_c\colon c^*\mathcal O_{X'}
      \xrightarrow{\ \sim\ }\mathcal O_C,\\
    \lambda_{Z',C}:=\varepsilon_c\circ c^*(\iota_{Z'})
      \colon c^*\mathcal I_{Z'}\longrightarrow\mathcal O_C.
  \end{gathered}                                               \tag{CBb}
  \]
  Define
  \[
  \begin{gathered}
    q_{Z',C}\colon c^*\mathcal I_{Z'}
      \twoheadrightarrow\mathcal I_{Z'}\mathcal O_C,\\
    \mathcal I_{Z'}\mathcal O_C
      =\operatorname{im}(\lambda_{Z',C}),\qquad
    \iota_{Z',C}\circ q_{Z',C}=\lambda_{Z',C},\\
    \lambda:=
      \widetilde c^{(\mathcal I_{Z'},q_{Z',C})}
      \colon C\longrightarrow X'',\\
    e_\lambda\colon b'\circ\lambda=c,\qquad
    (\lambda,e_\lambda)\in\operatorname{LiftPair}(c,q_{Z',C}).
  \end{gathered}                                             \tag{CB1}
  \]
  where
  $\iota_{Z',C}\colon\mathcal I_{Z'}\mathcal O_C
  \hookrightarrow\mathcal O_C$ is the inclusion.  The symbol $\lambda$ is
  the formula of
  \Cref{def:nagata-blowup-universal-lift} with
  $(f,\mathcal I,q_f)=(c,\mathcal I_{Z'},q_{Z',C})$, and
  $e_\lambda$ and the membership are
  \Cref{lemma:nagata-blowup-universal-exists}.

  By \Cref{def:nagata-blowup-product-comparison}, applied to
  $\mathcal I$ and $\mathcal J$, define
  \[
    \Phi_{\mathcal I,\mathcal J}\colon C\to B,\qquad
    \Psi_{\mathcal I,\mathcal J}\colon B\to C.
  \]
  Define
  \[
    \Omega_{Z,Z'}:=
      \Phi_{\mathcal I,\mathcal J}\circ\kappa\colon X''\longrightarrow B,
    \qquad
    \Lambda_{Z,Z'}:=
      \lambda\circ\Psi_{\mathcal I,\mathcal J}\colon B\longrightarrow X''.
  \]
\end{definition}

\begin{theorem}[{\cite[\href{https://stacks.math.columbia.edu/tag/080B}{Tag 080B}]{stacks-project}}]
  \label{thm:nagata-composition-of-blowups}
  \uses{def:nagata-composition-of-blowups-comparison}
  \lean{}

  The morphisms $\Omega_{Z,Z'}$ and $\Lambda_{Z,Z'}$ of
  \Cref{def:nagata-composition-of-blowups-comparison} satisfy
  \[
    \Lambda_{Z,Z'}\circ\Omega_{Z,Z'}=\operatorname{id}_{X''},
    \qquad
    \Omega_{Z,Z'}\circ\Lambda_{Z,Z'}=\operatorname{id}_{B}.
  \]
  Their equations over $X$ and their isomorphism predicate are
  \[
    b_{X,\mathcal K(Z')}\circ\Omega_{Z,Z'}=b\circ b',
    \qquad
    (b\circ b')\circ\Lambda_{Z,Z'}=b_{X,\mathcal K(Z')},
    \qquad
    \operatorname{IsIso}(\Omega_{Z,Z'}).
  \]
  The identified finite-presentation center and isomorphism are
  \[
    Y(Z'):=V(\mathcal K(Z'))\hookrightarrow X,
    \qquad
    (X'',b\circ b')\xrightarrow[\Omega_{Z,Z'}]{\sim}
      (\operatorname{Bl}_{\mathcal K(Z')}(X),b_{X,\mathcal K(Z')}).
  \]
  Its support satisfies
  \[
    |Y(Z')|=|Z|\cup |b|(|Z'|).
  \]
\end{theorem}

\begin{proof}
  \uses{lemma:nagata-composition-center-ideal,
    lemma:nagata-composition-extended-ideal-equality,
    thm:nagata-blowup-product-ideal,
    cor:nagata-blowup-cartier-iso}
  Use the notation $X'',C,B,\kappa,\lambda,\Phi,\Psi$ of
  \Cref{def:nagata-composition-of-blowups-comparison}.  By
  \Cref{lemma:nagata-composition-extended-ideal-equality},
  \[
    \mathcal J\mathcal O_{X'}
      =\mathcal I_E^d\mathcal I_{Z'}.
  \]
  Define
  \[
    q_{Z'}\colon (b')^*\mathcal I_{Z'}
      \twoheadrightarrow\mathcal I_{Z'}\mathcal O_{X''}.
  \]
  The comparison definitions are the membership pair
  \[
    \lambda\circ\kappa,\ \operatorname{id}_{X''}
      \in\operatorname{Lift}(b',q_{Z'}).
  \]
  The singleton equality in
  \Cref{thm:nagata-blowup-universal-property} is
  \[
    \lambda\circ\kappa=\operatorname{id}_{X''}.
  \]
  Define
  \[
    q_{\mathcal J,C}\colon c^*(\mathcal J\mathcal O_{X'})
      \twoheadrightarrow\mathcal J\mathcal O_C.
  \]
  The comparison definitions are the membership pair
  \[
    \kappa\circ\lambda,\ \operatorname{id}_C
      \in\operatorname{Lift}(c,q_{\mathcal J,C}).
  \]
  The singleton equality in
  \Cref{thm:nagata-blowup-universal-property} is
  \[
    \kappa\circ\lambda=\operatorname{id}_C.
  \]
  By \Cref{thm:nagata-blowup-product-ideal},
  \[
    \Psi\circ\Phi=\operatorname{id}_C,\qquad
    \Phi\circ\Psi=\operatorname{id}_B.
  \]
  Substitution of the definitions of $\Omega_{Z,Z'}$ and
  $\Lambda_{Z,Z'}$ is the pair of chains
  \[
  \begin{aligned}
    \Lambda_{Z,Z'}\circ\Omega_{Z,Z'}
      &=\lambda\circ\Psi\circ\Phi\circ\kappa\\
      &=\lambda\circ\operatorname{id}_C\circ\kappa\\
      &=\operatorname{id}_{X''},
  \end{aligned}
  \]
  \[
  \begin{aligned}
    \Omega_{Z,Z'}\circ\Lambda_{Z,Z'}
      &=\Phi\circ\kappa\circ\lambda\circ\Psi\\
      &=\Phi\circ\operatorname{id}_C\circ\Psi\\
      &=\operatorname{id}_B.
  \end{aligned}
  \]
  The four factorization equalities are
  \[
    c\circ\kappa=b',\qquad b'\circ\lambda=c,\qquad
    k\circ\Phi=b\circ c,\qquad b\circ c\circ\Psi=k.
  \]
  Substitution of the four factorization equalities is
  \[
    k\circ\Omega_{Z,Z'}
      =k\circ\Phi\circ\kappa
      =b\circ c\circ\kappa
      =b\circ b',
  \]
  and
  \[
    b\circ b'\circ\Lambda_{Z,Z'}
      =b\circ b'\circ\lambda\circ\Psi
      =b\circ c\circ\Psi=k.
  \]
  The two displayed structural equalities are the definition that the
  inverse isomorphisms are over $X$.  By
  \Cref{lemma:nagata-composition-center-ideal},
  \[
    |Y(Z')|=|V(\mathcal I)|\cup |b|(|Z'|)=|Z|\cup |b|(|Z'|).
  \]
  The same lemma states that $\mathcal K(Z')$ is finite type.  The
  finite-presentation criterion for a closed immersion defined by a
  finite-type quasi-coherent ideal is
  \[
    Y(Z')=V(\mathcal K(Z'))\hookrightarrow X
      \quad\text{of finite presentation}.
  \]
\end{proof}


\section{Admissible blow-ups and strict transforms}

\begin{definition}[$U$-admissible blow-up, {\cite[\href{https://stacks.math.columbia.edu/tag/080K}{Tag 080K}]{stacks-project}}]
  \label{def:nagata-admissible-blowup}
  \lean{}

  Let $j\colon U\hookrightarrow X$ be a quasi-compact open immersion.  A
  \emph{$U$-admissible blow-up} is a tuple
  \[
    (\mathcal I,X',\pi,\alpha_U)
  \]
  with the following explicitly typed components:
  \[
  \begin{gathered}
    \mathcal I\subseteq\mathcal O_X
       \text{ a finite-type quasi-coherent ideal},\qquad
    \mathcal I\mathcal O_U=\mathcal O_U,\\
    X'=\operatorname{Bl}_{\mathcal I}(X),\qquad
    \pi=b_{X,\mathcal I}\colon X'\longrightarrow X
       \quad\text{as in \Cref{def:nagata-global-blowup-map}},\\
    \alpha_U\colon X'\times_XU\xrightarrow{\sim}U,\qquad
    j\circ\alpha_U
      =\pi\circ\operatorname{pr}_{X'}.
  \end{gathered}
  \]
  Here $\alpha_U$ is the isomorphism of
  \Cref{thm:nagata-global-blowup-complement}.  Its inverse determines the
  explicitly named lift
  \[
    \widetilde j
      :=\operatorname{pr}_{X'}\circ\alpha_U^{-1}
      \colon U\longrightarrow X',
    \qquad
    \pi\circ\widetilde j=j.
  \]
\end{definition}

\begin{definition}[Strict transform]
  \label{def:nagata-strict-transform}
  \uses{def:nagata-scheme-image-factor,
    def:nagata-scheme-theoretic-image}
  \lean{}

  Let $f\colon Z\to X$, let $\pi\colon X'\to X$, let
  $j\colon U\hookrightarrow X$ be a quasi-compact open immersion, and fix
  an isomorphism
  \[
    \alpha_U\colon X'\times_XU\xrightarrow{\sim}U,\qquad
    j\circ\alpha_U=\pi\circ\operatorname{pr}_{X'}.
  \]
  For a blow-up defined by a finite-type quasi-coherent ideal
  $\mathcal I\subseteq\mathcal O_X$,
  $\pi=b_{X,\mathcal I}\colon\operatorname{Bl}_{\mathcal I}(X)\to X$, this
  means $U=X\setminus V(\mathcal I)$ and $\alpha_U$ is
  \Cref{thm:nagata-global-blowup-complement}.  Define
  \[
  \begin{gathered}
    Z_U:=Z\times_XU,\qquad
    k_U\colon Z_U\longrightarrow Z,\qquad
    f_U\colon Z_U\longrightarrow U,\\
    \lambda_{f,\pi}
      :=\bigl(k_U,
        \operatorname{pr}_{X'}\circ\alpha_U^{-1}\circ f_U\bigr)
      \colon Z_U\longrightarrow Z\times_XX'.
  \end{gathered}
  \]
  The equality required for the displayed fibre-product morphism is
  \[
    f\circ k_U
      =j\circ f_U
      =\pi\circ\operatorname{pr}_{X'}\circ\alpha_U^{-1}\circ f_U.
  \]
  Put $P:=Z\times_XX'$ with projection
  $p_Z\colon P\to Z$, and put
  \[
    P_U:=P\times_ZZ_U,
    \qquad
    q_P\colon P_U\to P,
    \qquad
    q_U\colon P_U\to Z_U.
  \]
  Define
  \[
    \theta_U:=(\lambda_{f,\pi},\operatorname{id}_{Z_U})
      \colon Z_U\longrightarrow P_U.
  \]
  The projection identities are
  \[
  \begin{gathered}
    q_U\circ\theta_U=\operatorname{id}_{Z_U},
    \qquad
    p_Z\circ q_P=k_U\circ q_U,\\
    \operatorname{pr}_{X'}\circ q_P
      =\operatorname{pr}_{X'}\circ\alpha_U^{-1}\circ f_U\circ q_U.
  \end{gathered}
  \]
  The fibre-product compatibility
  \[
    \pi\circ\operatorname{pr}_{X'}\circ q_P
      =f\circ p_Z\circ q_P
      =f\circ k_U\circ q_U
      =j\circ f_U\circ q_U
  \]
  defines
  \[
    \chi_U:=
      (\operatorname{pr}_{X'}\circ q_P,f_U\circ q_U)
      \colon P_U\longrightarrow X'\times_XU.
  \]
  The structural equality for $\alpha_U$ and the preceding chain compose to
  \[
  \begin{aligned}
    j\circ\alpha_U\circ\chi_U
      &=\pi\circ\operatorname{pr}_{X'}\circ\chi_U\\
      &=\pi\circ\operatorname{pr}_{X'}\circ q_P\\
      &=f\circ p_Z\circ q_P\\
      &=f\circ k_U\circ q_U\\
      &=j\circ f_U\circ q_U.
  \end{aligned}
  \]
  The open immersion $j$ satisfies
  $\operatorname{Mono}(j)$.  Applying left cancellation by $j$ to
  the equality displayed immediately above and substituting the inverse
  identities for $\alpha_U$ are the equalities
  \[
    \alpha_U\circ\chi_U=f_U\circ q_U,
    \qquad
    \chi_U=\alpha_U^{-1}\circ f_U\circ q_U,
  \]
  and application of $\operatorname{pr}_{X'}$ to the second equality is
  \[
    \operatorname{pr}_{X'}\circ q_P
      =\operatorname{pr}_{X'}\circ\alpha_U^{-1}\circ f_U\circ q_U.
  \]
  Together with $p_Z\circ q_P=k_U\circ q_U$, fibre-product uniqueness for
  $P=Z\times_XX'$ is the equality
  \[
    q_P=\lambda_{f,\pi}\circ q_U.
  \]
  The two components of $\theta_U\circ q_U$ are
  \[
  \begin{aligned}
    q_P\circ\theta_U\circ q_U
      &=\lambda_{f,\pi}\circ q_U=q_P,\\
    q_U\circ\theta_U\circ q_U&=q_U.
  \end{aligned}
  \]
  Fibre-product uniqueness for $P_U=P\times_ZZ_U$, together with the first
  projection identity for $\theta_U$, is the pair of inverse equalities
  \[
    \theta_U\circ q_U=\operatorname{id}_{P_U},
    \qquad
    q_U\circ\theta_U=\operatorname{id}_{Z_U},
    \qquad
    \operatorname{IsIso}(\theta_U),
    \qquad
    \theta_U^{-1}=q_U.
  \]
  The defining projection equation $q_P\circ\theta_U=\lambda_{f,\pi}$ is
  \[
    \lambda_{f,\pi}=q_P\circ\theta_U.
  \]
  The morphism $k_U$ is the base change of the quasi-compact open immersion
  $j$, and $q_P$ is the base change of $k_U$.  Stability under base change
  and composition records
  \[
    \lambda_{f,\pi}\colon Z_U\to P
      \text{ is a quasi-compact open immersion},
    \qquad
    \operatorname{IsQCQS}(\lambda_{f,\pi}).
  \]
  The \emph{strict transform} and its closed immersion are, by
  \Cref{def:nagata-scheme-theoretic-image},
  \[
    \operatorname{ST}_{\pi}(Z)
      :=\operatorname{SIm}(\lambda_{f,\pi}),
    \qquad
    \iota_{f,\pi}:=i_{\lambda_{f,\pi}}\colon
      \operatorname{ST}_{\pi}(Z)\hookrightarrow Z\times_XX'.
  \]
  By \Cref{def:nagata-scheme-image-factor}, define
  \[
    \lambda_{f,\pi}^{\mathrm{im}}\colon
      Z_U\longrightarrow\operatorname{ST}_{\pi}(Z),
    \qquad
    \iota_{f,\pi}\circ\lambda_{f,\pi}^{\mathrm{im}}
      =\lambda_{f,\pi}.                                           \tag{ST0f}
  \]
  Its projection to $X'$ is the explicitly defined morphism
  \[
    f^{\mathrm{st}}_{\pi}
      :=\operatorname{pr}_{X'}\circ\iota_{f,\pi}
      \colon\operatorname{ST}_{\pi}(Z)\longrightarrow X'.
  \]

  \paragraph{Mathlib relation.}
  The image object and its factorization are represented by the exact
  declarations \texttt{\detokenize{Scheme.Hom.image}},
  \texttt{\detokenize{Scheme.Hom.imageι}}, and
  \texttt{\detokenize{Scheme.Hom.toImage}}.  Mathlib does not currently
  package the data depending on $(U,\alpha_U)$ as a strict-transform
  construction; this definition specializes its scheme-image API.
\end{definition}

\begin{lemma}[Strict transform as an image in the modified space]
  \label{lemma:nagata-strict-transform-scheme-image-comparison}
  \uses{def:nagata-strict-transform,
    lemma:nagata-scheme-image-universal}

  In \Cref{def:nagata-strict-transform}, assume in addition that
  $f\colon Z\hookrightarrow X$ is a closed immersion and define
  \[
    q_{f,\pi}:=
      \operatorname{pr}_{X'}\circ\lambda_{f,\pi}
      \colon Z_U\longrightarrow X'.                             \tag{ST0g}
  \]
  Then $f^{\mathrm{st}}_\pi$ is a closed immersion and there is a unique
  isomorphism over $X'$
  \[
    c_{f,\pi}\colon
      \operatorname{ST}_{\pi}(Z)
        \xrightarrow{\sim}
      \operatorname{SIm}_{X'}(q_{f,\pi})                         \tag{ST0h}
  \]
  satisfying
  \[
  \begin{gathered}
    i_{q_{f,\pi}}\circ c_{f,\pi}=f^{\mathrm{st}}_\pi,\\
    c_{f,\pi}\circ\lambda_{f,\pi}^{\mathrm{im}}
      =q_{f,\pi}^{\mathrm{im}}.
  \end{gathered}                                                   \tag{ST0i}
  \]

  \paragraph{Mathlib relation.}
  This is a blueprint-local comparison between two applications of
  \texttt{\detokenize{Scheme.Hom.image}}.  Its factor maps are instances of
  \texttt{\detokenize{Scheme.Hom.toImage}}; no separate Mathlib strict-
  transform declaration is being asserted.
\end{lemma}

\begin{proof}
  Since $f$ is a closed immersion, its base change
  $\operatorname{pr}_{X'}\colon Z\times_XX'\hookrightarrow X'$ is a closed
  immersion.  Thus a closed subscheme of $Z\times_XX'$ containing
  $\lambda_{f,\pi}(Z_U)$ is carried to a closed subscheme of $X'$ containing
  $q_{f,\pi}(Z_U)$.  Conversely, pullback along this closed immersion sends
  every closed subscheme containing $q_{f,\pi}(Z_U)$ to one containing
  $\lambda_{f,\pi}(Z_U)$.  Apply the two scheme-image universal properties.
  Their comparison maps satisfy (ST0i), and uniqueness makes their two
  composites identities.  Equation (ST0h) follows, and (ST0i) identifies
  $f^{\mathrm{st}}_\pi$ with a scheme-image inclusion.
\end{proof}

\subsection{Power torsion and strict-transform charts}

The strict transform is the correction term for the failure of blow-ups to
commute with non-flat base change.  The following algebraic lemmas make that
correction explicit before it is sheafified and glued.

\begin{definition}[Power-torsion submodule]
  \label{def:nagata-power-torsion}
  \lean{}

  Let $R$ be a commutative ring, let $M$ be an $R$-module, and let $a\in R$.
  Put $S_a=\{a^n\mid n\in\mathbf N\}$ and define
  \[
    \Gamma_{(a)}(M)
      :=\{m\in M\mid \exists n\in\mathbf N,\ a^nm=0\}.
  \]

  \paragraph{Mathlib relation.}
  With $S_a=\texttt{\detokenize{Submonoid.powers}}\,a$, this is exactly
  \texttt{\detokenize{Submodule.torsion' R M S_a}}.
\end{definition}

\begin{lemma}[Power torsion is a submodule]
  \label{lemma:nagata-power-torsion-submodule}
  \uses{def:nagata-power-torsion}
  \lean{}

  The subset $\Gamma_{(a)}(M)$ contains zero and is closed under addition and
  scalar multiplication.

  \paragraph{Mathlib relation.}
  These closure properties are already encoded by the submodule-valued
  declaration \texttt{\detokenize{Submodule.torsion'}}.
\end{lemma}

\begin{proof}
  By \Cref{def:nagata-power-torsion},
  \[
    a^0\cdot0=0,
    \qquad 0\in\Gamma_{(a)}(M).
  \]
  If $m_i\in\Gamma_{(a)}(M)$, choose $n_i\in\mathbf N$ with
  $a^{n_i}m_i=0$.  Distributivity and the two annihilation equalities form the
  calculation
  \[
  \begin{aligned}
    a^{n_1+n_2}(m_1+m_2)
      &=a^{n_2}(a^{n_1}m_1)+a^{n_1}(a^{n_2}m_2)\\
      &=a^{n_2}0+a^{n_1}0=0,
  \end{aligned}
  \]
  By \Cref{def:nagata-power-torsion}, the last equality is the membership
  \[
    m_1+m_2\in\Gamma_{(a)}(M).
  \]
  For $r\in R$,
  \[
    a^{n_1}(rm_1)=r(a^{n_1}m_1)=r0=0,
  \]
  and \Cref{def:nagata-power-torsion} rewrites it as
  \[
    rm_1\in\Gamma_{(a)}(M).
  \]
  The three membership statements are the zero, addition, and scalar
  submodule axioms.
\end{proof}

\begin{lemma}[Membership in power torsion]
  \label{lemma:nagata-mem-power-torsion}
  \uses{def:nagata-power-torsion}
  \lean{}

  One has
  \[
    m\in\Gamma_{(a)}(M)
      \Longleftrightarrow
    \exists n,\ a^nm=0.
  \]
  The right-hand side should be the simp-normal form for membership.

  \paragraph{Mathlib relation.}
  After unwinding \texttt{\detokenize{Submonoid.powers}}, this is
  \texttt{\detokenize{Submodule.mem_torsion'_iff}}.
\end{lemma}

\begin{proof}
  By the set-builder equality in \Cref{def:nagata-power-torsion}, the
  membership predicate is
  \[
  \begin{aligned}
    m\in\Gamma_{(a)}(M)
      &\Longleftrightarrow
        m\in\{x\in M\mid\exists n\in\mathbf N,\ a^nx=0\}\\
      &\Longleftrightarrow
        \exists n\in\mathbf N,\ a^nm=0.
  \end{aligned}
  \]
\end{proof}

\begin{lemma}[Power torsion as a localization kernel]
  \label{lemma:nagata-power-torsion-kernel-localization}
  \uses{def:nagata-power-torsion}
  \lean{}

  Put $M_a:=S_a^{-1}M$ and let
  $\ell_{M,a}\colon M\to M_a$, $m\mapsto m/1$, be the localization map.
  \[
    \Gamma_{(a)}(M)=\ker(\ell_{M,a}).
  \]

  \paragraph{Mathlib relation.}
  This follows by extensionality from
  \texttt{\detokenize{LocalizedModule.mem_ker_mkLinearMap_iff}}, or from
  \texttt{\detokenize{IsLocalizedModule.mem_ker_iff}} for an arbitrary
  localization model, with
  $S=\texttt{\detokenize{Submonoid.powers}}\,a$.
\end{lemma}

\begin{proof}
  For $m\in M$, the zero criterion for the localization $S_a^{-1}M$ is
  \[
  \begin{aligned}
    m\in\ker(\ell_{M,a})
      &\Longleftrightarrow \ell_{M,a}(m)=0\\
      &\Longleftrightarrow \exists s\in S_a,\ sm=0\\
      &\Longleftrightarrow \exists n\in\mathbf N,\ a^nm=0\\
      &\Longleftrightarrow m\in\Gamma_{(a)}(M),
  \end{aligned}
  \]
  where the last equivalence is
  \Cref{lemma:nagata-mem-power-torsion}.  Submodule extensionality applied to
  the quantified equivalence
  \[
    \forall m\in M,\quad
    m\in\ker(\ell_{M,a})
      \Longleftrightarrow m\in\Gamma_{(a)}(M)
  \]
  is the equality
  \[
    \ker(\ell_{M,a})=\Gamma_{(a)}(M).
  \]
\end{proof}

\begin{lemma}[Localization of power torsion]
  \label{lemma:nagata-power-torsion-localization}
  \lean{}

  Let $s\in R$, put $S_s=\{s^n\mid n\in\mathbf N\}$, and write
  $M_s=S_s^{-1}M$.  For a submodule $N\subseteq M$, write
  $N_s:=S_s^{-1}N\subseteq M_s$ for the image under the localization map.
  \[
    \bigl(\Gamma_{(a)}(M)\bigr)_s
      =\Gamma_{(a/1)}(M_s).
  \]
\end{lemma}

\begin{proof}
  \uses{lemma:nagata-power-torsion-kernel-localization}
  Let
  \[
    \ell_{M,a}\colon M\to M_a,
    \qquad
    (\ell_{M,a})_s\colon M_s\to(M_a)_s
  \]
  be the map of \Cref{lemma:nagata-power-torsion-kernel-localization} and its
  localization at $S_s$.  Exactness of $S_s^{-1}(-)$ is the kernel identity
  \[
    S_s^{-1}\ker(\ell_{M,a})
      =\ker\bigl(S_s^{-1}\ell_{M,a}\bigr)
      =\ker\bigl((\ell_{M,a})_s\bigr).
  \]
  Substitution of this identity is the formal chain
  \[
  \begin{aligned}
    \bigl(\Gamma_{(a)}(M)\bigr)_s
      &=S_s^{-1}\ker(\ell_{M,a})
        &&\text{by \Cref{lemma:nagata-power-torsion-kernel-localization}}\\
      &=\ker\bigl((\ell_{M,a})_s\bigr).
  \end{aligned}
  \]
  Let
  \[
    \tau_{a,s}\colon(M_a)_s\xrightarrow{\sim}(M_s)_{a/1},
    \qquad
    \tau_{a,s}\!\left(\frac{m/a^r}{s^t}\right)
       =\frac{m/s^t}{(a/1)^r},
  \]
  be the iterated-localization isomorphism.  For every $m/s^t\in M_s$, its
  defining fraction formula satisfies
  \[
  \begin{aligned}
    \bigl(\tau_{a,s}\circ(\ell_{M,a})_s\bigr)(m/s^t)
      &=\tau_{a,s}\bigl((m/1)/s^t\bigr)\\
      &=(m/s^t)/1
       =\ell_{M_s,a/1}(m/s^t).
  \end{aligned}
  \]
  Map extensionality applied to this elementwise equality is
  \[
    \tau_{a,s}\circ(\ell_{M,a})_s
      =\ell_{M_s,a/1}\colon M_s\longrightarrow(M_s)_{a/1}.
  \]
  The isomorphism $\tau_{a,s}$ satisfies
  $\ker(\tau_{a,s})=0$.  Substitution of
  $\tau_{a,s}\circ(\ell_{M,a})_s=\ell_{M_s,a/1}$ and
  \Cref{lemma:nagata-power-torsion-kernel-localization} is the kernel chain
  \[
  \begin{aligned}
    \ker\bigl((\ell_{M,a})_s\bigr)
      &=\ker\bigl(\tau_{a,s}\circ(\ell_{M,a})_s\bigr)
        &&(\ker(\tau_{a,s})=0)\\
      &=\ker(\ell_{M_s,a/1})\\
      &=\Gamma_{(a/1)}(M_s)
        &&\text{by \Cref{lemma:nagata-power-torsion-kernel-localization}}.
  \end{aligned}
  \]
  Together with
  $\ker((\ell_{M,a})_s)=(\Gamma_{(a)}(M))_s$, this chain is the asserted
  equality
  \[
    (\Gamma_{(a)}(M))_s=\Gamma_{(a/1)}(M_s).
  \]
\end{proof}

\begin{lemma}[The power-torsion quotient has no exceptional torsion]
  \label{lemma:nagata-power-torsion-quotient-torsionfree}
  \uses{def:nagata-power-torsion}
  \lean{}

  Multiplication by $a$ on $M/\Gamma_{(a)}(M)$ is injective.  More generally,
  no nonzero class in the quotient is killed by a power of $a$.
\end{lemma}

\begin{proof}
  Let
  $q\colon M\to\overline M:=M/\Gamma_{(a)}(M)$ be the quotient map.  Its
  defining identities are
  \[
    q(rm)=r q(m),\qquad
    \ker(q)=\Gamma_{(a)}(M).
  \]
  For $m\in M$ satisfying $a^nq(m)=0$, these identities form the formal
  equivalence chain
  \[
  \begin{aligned}
    a^nq(m)=0
      &\Longleftrightarrow q(a^nm)=0\\
      &\Longleftrightarrow a^nm\in\Gamma_{(a)}(M)\\
      &\Longleftrightarrow
        \exists r\in\mathbf N,\ a^r(a^nm)=0
        &&\text{by \Cref{lemma:nagata-mem-power-torsion}}\\
      &\Longleftrightarrow a^{r+n}m=0\text{ for some }r\\
      &\Longleftrightarrow m\in\Gamma_{(a)}(M)
        &&\text{by \Cref{lemma:nagata-mem-power-torsion}}\\
      &\Longleftrightarrow q(m)=0.
  \end{aligned}
  \]
  The first and last terms of the equivalence chain state
  \[
    \forall n\in\mathbf N,\ \forall\overline m\in\overline M,\qquad
    a^n\overline m=0\Longrightarrow\overline m=0.
  \]
  At $n=1$ this is the kernel equation
  \[
    \ker(\overline m\mapsto a\overline m)=0.
  \]
\end{proof}

\begin{definition}[Affine base-change chart ring]
  \label{def:nagata-strict-transform-basechange-chart-ring}
  \uses{def:nagata-blowup-chart-ring}
  \lean{}

  Fix the standing algebra datum
  \[
    \mathsf D_{\mathrm{alg}}
      :=(A,B,I,a,\varphi),\qquad
    A,B\text{ commutative rings},\qquad
    I\triangleleft A,\qquad a\in I,\qquad
    \varphi\colon A\to B\text{ a ring homomorphism}.                \tag{STAlg0}
  \]
  Define
  \[
    C_a(A,B,I):=B\otimes_{A,\varphi}A[I/a],
  \]
  together with
  \[
    \eta_B(y):=y\otimes1,\qquad
    \eta_a(q):=1\otimes q,\qquad
    a_C:=\eta_B(\varphi(a))=\eta_a(a).
  \]
  The notation $\Gamma_{(a)}(C_a(A,B,I))$ means
  $\Gamma_{(a_C)}(C_a(A,B,I))$.
\end{definition}

\begin{lemma}[Ordinary base-change chart]
  \label{lemma:nagata-strict-transform-basechange-chart}
  \uses{def:nagata-strict-transform-basechange-chart-ring}
  \lean{}

  Fix the expanded standing datum
  \[
    \mathsf D_{\mathrm{alg}}=(A,B,I,a,\varphi),\qquad
    A,B\text{ commutative rings},\quad
    I\triangleleft A,\quad a\in I,\quad
    \varphi\colon A\to B
  \]
  of (STAlg0) in
  \Cref{def:nagata-strict-transform-basechange-chart-ring}.
  Let $\chi_a\colon D_+(at)\xrightarrow{\sim}
  \operatorname{Spec}(A[I/a])$ be
  \Cref{prop:nagata-affine-blowup-chart}.  Define $\beta_a$ by
  \[
    \beta_a\colon\operatorname{Spec}(C_a(A,B,I))
      \longrightarrow
      \operatorname{Spec}(B)\times_{\operatorname{Spec}A}D_+(at),
    \qquad
    \operatorname{pr}_B\circ\beta_a=\operatorname{Spec}(\eta_B),
  \]
  \[
    \operatorname{pr}_{D_+(at)}\circ\beta_a
      =\chi_a^{-1}\circ\operatorname{Spec}(\eta_a).
  \]
  Equivalently, $\operatorname{IsIso}(\beta_a)$.
\end{lemma}

\begin{proof}
  \uses{prop:nagata-affine-blowup-chart}
  Let
  \[
    \omega_a\colon
      \operatorname{Spec}(B\otimes_AA[I/a])
      \xrightarrow{\sim}
      \operatorname{Spec}(B)\times_{\operatorname{Spec}A}
        \operatorname{Spec}(A[I/a])
  \]
  be the affine fibre-product isomorphism defined by
  \[
    \operatorname{pr}_B\circ\omega_a=\operatorname{Spec}(\eta_B),
    \qquad
    \operatorname{pr}_{A[I/a]}\circ\omega_a
      =\operatorname{Spec}(\eta_a).
  \]
  By \Cref{prop:nagata-affine-blowup-chart},
  $\chi_a\colon D_+(at)\xrightarrow{\sim}\operatorname{Spec}(A[I/a])$.
  Define
  \[
    \gamma_a
      :=(\operatorname{id}_{\operatorname{Spec}B}\times\chi_a^{-1})
        \circ\omega_a
      \colon\operatorname{Spec}(C_a(A,B,I))
        \longrightarrow
        \operatorname{Spec}(B)\times_{\operatorname{Spec}A}D_+(at).
  \]
  Its projection equations are
  \[
  \begin{aligned}
    \operatorname{pr}_B\circ\gamma_a
      &=\operatorname{Spec}(\eta_B),\\
    \operatorname{pr}_{D_+(at)}\circ\gamma_a
      &=\chi_a^{-1}\circ\operatorname{Spec}(\eta_a).
  \end{aligned}
  \]
  The defining equations for $\beta_a$ and the displayed equations for
  $\gamma_a$ are
  \[
  \begin{gathered}
    \operatorname{pr}_B\circ\beta_a
      =\operatorname{Spec}(\eta_B)
      =\operatorname{pr}_B\circ\gamma_a,\\
    \operatorname{pr}_{D_+(at)}\circ\beta_a
      =\chi_a^{-1}\circ\operatorname{Spec}(\eta_a)
      =\operatorname{pr}_{D_+(at)}\circ\gamma_a.
  \end{gathered}
  \]
  The uniqueness clause in the fibre-product universal property is the
  implication
  \[
    \left(
      \operatorname{pr}_B\circ\beta_a
        =\operatorname{pr}_B\circ\gamma_a
      \ \wedge\ 
      \operatorname{pr}_{D_+(at)}\circ\beta_a
        =\operatorname{pr}_{D_+(at)}\circ\gamma_a
    \right)
    \Longrightarrow
    \beta_a=\gamma_a.
  \]
  The inverse equations for $\omega_a$ and
  $\operatorname{id}_{\operatorname{Spec}B}\times\chi_a^{-1}$ are
  \[
  \begin{gathered}
    \operatorname{IsIso}(\omega_a),\qquad
    \operatorname{IsIso}
      (\operatorname{id}_{\operatorname{Spec}B}\times\chi_a^{-1}),\\
    \operatorname{IsIso}(\gamma_a),\qquad
    \beta_a=\gamma_a,
    \qquad \operatorname{IsIso}(\beta_a)
      \quad\text{by \Cref{lemma:nagata-morphism-property-calculus}},
  \end{gathered}
  \]
\end{proof}

\begin{definition}[Affine strict-transform chart ring]
  \label{def:nagata-strict-transform-chart-ring}
  \uses{def:nagata-power-torsion,
    def:nagata-strict-transform-basechange-chart-ring}
  \lean{}

  Fix the expanded standing datum
  \[
    \mathsf D_{\mathrm{alg}}=(A,B,I,a,\varphi),\qquad
    A,B\text{ commutative rings},\quad
    I\triangleleft A,\quad a\in I,\quad
    \varphi\colon A\to B
  \]
  of (STAlg0) in
  \Cref{def:nagata-strict-transform-basechange-chart-ring}.
  Define
  \[
    C_a^{\mathrm{st}}(A,B,I)
      :=C_a(A,B,I)/\Gamma_{(a_C)}(C_a(A,B,I)),
    \qquad
    q_a\colon C_a(A,B,I)\twoheadrightarrow C_a^{\mathrm{st}}(A,B,I).
  \]
\end{definition}

\begin{lemma}[Closedness of the strict-transform chart]
  \label{lemma:nagata-strict-transform-chart-closed}
  \uses{def:nagata-strict-transform-chart-ring}
  \lean{}

  Fix the expanded standing datum
  \[
    \mathsf D_{\mathrm{alg}}=(A,B,I,a,\varphi),\qquad
    A,B\text{ commutative rings},\quad
    I\triangleleft A,\quad a\in I,\quad
    \varphi\colon A\to B
  \]
  of (STAlg0) in
  \Cref{def:nagata-strict-transform-basechange-chart-ring}.
  The morphism
  \[
    \kappa_a:=\operatorname{Spec}(q_a)\colon
      \operatorname{Spec}(C_a^{\mathrm{st}}(A,B,I))
      \hookrightarrow\operatorname{Spec}(C_a(A,B,I))
  \]
  is a closed immersion.
\end{lemma}

\begin{proof}
  \uses{lemma:nagata-power-torsion-submodule}
  Put $C=C_a(A,B,I)$ and $T=\Gamma_{(a_C)}(C)$.  By
  \Cref{lemma:nagata-power-torsion-submodule}, $T$ is a submodule of the
  regular $C$-module $C$.  Its scalar-closure axiom is
  \[
    cT\subseteq T\qquad(c\in C),
  \]
  and the ideal predicate for a submodule of the regular module is
  \[
    T\triangleleft C.
  \]
  The quotient identities are
  \[
    C_a^{\mathrm{st}}=C/T,\qquad\ker(q_a)=T,
  \]
  $q_a$ is surjective.  The affine closed-immersion criterion identifies
  $\operatorname{Spec}(q_a)=\kappa_a$ as a closed immersion.
\end{proof}

\begin{definition}[Map to the blow-up chart of the extended ideal]
  \label{def:nagata-strict-chart-to-extended-blowup-chart}
  \uses{def:nagata-strict-transform-basechange-chart-ring,
    def:nagata-blowup-chart-ring}
  \lean{}

  Fix the expanded standing datum
  \[
    \mathsf D_{\mathrm{alg}}=(A,B,I,a,\varphi),\qquad
    A,B\text{ commutative rings},\quad
    I\triangleleft A,\quad a\in I,\quad
    \varphi\colon A\to B
  \]
  of (STAlg0) in
  \Cref{def:nagata-strict-transform-basechange-chart-ring}.
  Put
  \[
    J:=IB=\varphi(I)B,\qquad b:=\varphi(a).
  \]
  Define
  \[
    \varphi_{I,a}\colon A[I/a]\to B[J/b],
    \qquad \varphi_{I,a}(x/a):=\varphi(x)/b,
  \]
  and define
  \[
    \sigma_a\colon C_a(A,B,I)\longrightarrow B[J/b]
  \]
  by
  \[
    \sigma_a(y\otimes q):=y\varphi_{I,a}(q).
  \]
\end{definition}

\begin{lemma}[Well-definedness of the extended-ideal chart map]
  \label{lemma:nagata-strict-chart-to-extended-blowup-chart}
  \uses{def:nagata-strict-chart-to-extended-blowup-chart}
  \lean{}

  Fix the expanded standing datum
  \[
    \mathsf D_{\mathrm{alg}}=(A,B,I,a,\varphi),\qquad
    A,B\text{ commutative rings},\quad
    I\triangleleft A,\quad a\in I,\quad
    \varphi\colon A\to B
  \]
  of (STAlg0) in
  \Cref{def:nagata-strict-transform-basechange-chart-ring}.
  The assignment defining $\sigma_a$ respects addition, multiplication, the
  localization relations in $A[I/a]$, and the two $A$-algebra structures.
  Equivalently,
  \[
    \sigma_a\in
      \operatorname{Hom}_{B\text{-alg}}
        \bigl(C_a(A,B,I),B[J/b]\bigr).
  \]
\end{lemma}

\begin{proof}
  Define the localization homomorphism
  \[
    \widetilde\varphi_a\colon A_a\longrightarrow B_b,
    \qquad
    \widetilde\varphi_a(r/a^n):=\varphi(r)/b^n.
  \]
  For every $x\in I$,
  \[
    \widetilde\varphi_a(x/a)=\varphi(x)/b\in B[J/b].
  \]
  The generator identities
  \[
    A[I/a]=A[\,x/a\mid x\in I\,]\subseteq A_a,
    \qquad
    B[J/b]=B[\,\varphi(x)/b\mid x\in I\,]\subseteq B_b,
  \]
  and
  $\widetilde\varphi_a(x/a)=\varphi(x)/b\in B[J/b]$ are the inclusion
  \[
    \widetilde\varphi_a(A[I/a])\subseteq B[J/b].
  \]
  Define the restricted homomorphism
  \[
    \varphi_{I,a}
      :=\widetilde\varphi_a|_{A[I/a]}
      \colon A[I/a]\longrightarrow B[J/b],
    \qquad
    \varphi_{I,a}(x/a)=\varphi(x)/b.
  \]
  Let $\iota_B\colon B\to B[J/b]$ and
  $\iota_A\colon A\to A[I/a]$ be the inclusion homomorphisms.  For every
  $r\in A$ one has
  \[
  \begin{aligned}
    (\iota_B\circ\varphi)(r)
      &=\varphi(r)
       =\varphi(r)/b^0\\
      &=\widetilde\varphi_a(r/a^0)
       =(\varphi_{I,a}\circ\iota_A)(r),
  \end{aligned}
  \]
  and the image equality is
  \[
    \iota_B\circ\varphi=\varphi_{I,a}\circ\iota_A.
  \]
  The tensor-product universal property applied to this equality defines the
  unique homomorphism
  \[
    \sigma_a\colon B\otimes_AA[I/a]\longrightarrow B[J/b]
  \]
  satisfying
  \[
    \sigma_a\circ\eta_B=\iota_B,
    \qquad
    \sigma_a\circ\eta_a=\varphi_{I,a}.
  \]
  For $y\in B$ and $q\in A[I/a]$, the two defining equations expand as
  \[
  \begin{aligned}
    \sigma_a(y\otimes q)
      &=\sigma_a((y\otimes1)(1\otimes q))\\
      &=\sigma_a(y\otimes1)\sigma_a(1\otimes q)\\
      &=y\varphi_{I,a}(q),
  \end{aligned}
  \]
  which is the formula in
  \Cref{def:nagata-strict-chart-to-extended-blowup-chart}.
\end{proof}

\begin{lemma}[Surjectivity onto the extended-ideal chart]
  \label{lemma:nagata-strict-chart-map-surjective}
  \uses{def:nagata-strict-chart-to-extended-blowup-chart}
  \lean{}

  Fix the expanded standing datum
  \[
    \mathsf D_{\mathrm{alg}}=(A,B,I,a,\varphi),\qquad
    A,B\text{ commutative rings},\quad
    I\triangleleft A,\quad a\in I,\quad
    \varphi\colon A\to B
  \]
  of (STAlg0) in
  \Cref{def:nagata-strict-transform-basechange-chart-ring}.
  The constructed homomorphism $\sigma_a$ is surjective.
\end{lemma}

\begin{proof}
  \uses{lemma:nagata-strict-chart-to-extended-blowup-chart,
    lemma:nagata-chart-ring-normal-form}
  By \Cref{lemma:nagata-chart-ring-normal-form}, an element of the target has
  the form $z/b^n$ with $z\in J^n$.  Write
  \[
    z=\sum_k y_k\prod_{r=1}^n\varphi(x_{kr}),
    \qquad y_k\in B,\quad x_{kr}\in I.
  \]
  Then
  \[
    \frac z{b^n}
      =\sigma_a\left(
        \sum_k y_k\otimes
          \frac{\prod_{r=1}^nx_{kr}}{a^n}\right).
  \]
  As $z/b^n$ ranges over
  \Cref{lemma:nagata-chart-ring-normal-form}, the preceding preimage formula
  is the pair of inclusions and equality
  \[
    B[J/b]\subseteq\operatorname{im}(\sigma_a)
      \subseteq B[J/b],
    \qquad
    \operatorname{im}(\sigma_a)=B[J/b].
  \]
\end{proof}

\begin{lemma}[Kernel of the map to the extended-ideal chart]
  \label{lemma:nagata-strict-chart-map-kernel}
  \uses{def:nagata-power-torsion,def:nagata-strict-chart-to-extended-blowup-chart}
  \lean{}

  Fix the expanded standing datum
  \[
    \mathsf D_{\mathrm{alg}}=(A,B,I,a,\varphi),\qquad
    A,B\text{ commutative rings},\quad
    I\triangleleft A,\quad a\in I,\quad
    \varphi\colon A\to B
  \]
  of (STAlg0) in
  \Cref{def:nagata-strict-transform-basechange-chart-ring}.
  The kernel of $\sigma_a$ is exactly
  $\Gamma_{(a)}(C_a(A,B,I))$.
\end{lemma}

\begin{proof}
  \uses{lemma:nagata-strict-chart-map-surjective,
    lemma:nagata-power-torsion-kernel-localization}
  Put $C=C_a(A,B,I)$ and let
  $\ell_{C,a}\colon C\to C_{a_C}$.  The formulas
  \[
    \delta_a\left(\frac{y\otimes q}{a_C^n}\right)
      :=\frac{y\varphi_{I,a}(q)}{b^n},
    \qquad
    \delta_a^{-1}(y/b^n):=(y\otimes1)/a_C^n
  \]
  define homomorphisms
  $\delta_a\colon C_{a_C}\to B_b$ and
  $\delta_a^{-1}\colon B_b\to C_{a_C}$.  The ring $C_{a_C}$ is generated
  by the images of $y\otimes1$ ($y\in B$), $1\otimes(x/a)$ ($x\in I$),
  and $a_C^{-1}$.  On these generators,
  \[
  \begin{aligned}
    \delta_a\!\left(\delta_a^{-1}(y/b^n)\right)
      &=\delta_a((y\otimes1)/a_C^n)=y/b^n,\\
    \delta_a^{-1}\!\left(\delta_a((y\otimes1)/1)\right)
      &=\delta_a^{-1}(y/1)=(y\otimes1)/1,\\
    \delta_a^{-1}\!\left(
      \delta_a((1\otimes(x/a))/1)\right)
      &=\delta_a^{-1}(\varphi(x)/b)
       =(\varphi(x)\otimes1)/a_C\\
      &=(1\otimes x)/(1\otimes a)
       =(1\otimes(x/a))/1,\\
    \delta_a^{-1}\!\left(\delta_a(a_C^{-1})\right)
      &=\delta_a^{-1}(b^{-1})=a_C^{-1}.
  \end{aligned}
  \]
  Ring-homomorphism extensionality applied to the preceding identities is
  \[
    \delta_a\circ\delta_a^{-1}=\operatorname{id}_{B_b},
    \qquad
    \delta_a^{-1}\circ\delta_a=\operatorname{id}_{C_{a_C}},
    \qquad
    \operatorname{IsIso}(\delta_a),
    \qquad
    \delta_a^{-1}\text{ is the displayed inverse}.
  \]
  Let $\jmath_a\colon B[J/b]\hookrightarrow B_b$ be the inclusion from
  \Cref{def:nagata-blowup-chart-ring}.  The defining formulas for
  $\sigma_a$, $\delta_a$, and $\ell_{C,a}$ are the equality
  \[
    \jmath_a\circ\sigma_a=\delta_a\circ\ell_{C,a}.
  \]
  The injectivity predicates are
  $\ker(\jmath_a)=0$ and $\ker(\delta_a)=0$.  The kernel calculation is
  \[
  \begin{aligned}
    \ker(\sigma_a)
      &=\ker(\jmath_a\circ\sigma_a)
        &&(\ker(\jmath_a)=0)\\
      &=\ker(\delta_a\circ\ell_{C,a})\\
      &=\ker(\ell_{C,a})
        &&(\ker(\delta_a)=0)\\
      &=\Gamma_{(a_C)}(C)
        &&\text{by \Cref{lemma:nagata-power-torsion-kernel-localization}}.
  \end{aligned}
  \]
\end{proof}

\begin{definition}[Strict-transform chart comparison]
  \label{def:nagata-strict-transform-chart-comparison}
  \uses{def:nagata-strict-transform-chart-ring,
    def:nagata-strict-chart-to-extended-blowup-chart}
  \lean{}

  Fix the expanded standing datum
  \[
    \mathsf D_{\mathrm{alg}}=(A,B,I,a,\varphi),\qquad
    A,B\text{ commutative rings},\quad
    I\triangleleft A,\quad a\in I,\quad
    \varphi\colon A\to B
  \]
  of (STAlg0) in
  \Cref{def:nagata-strict-transform-basechange-chart-ring}.
  Let $q_a$ be \Cref{def:nagata-strict-transform-chart-ring}.  Define
  \[
    \overline\sigma_a\colon
      C_a^{\mathrm{st}}(A,B,I)\longrightarrow B[IB/b]
  \]
  by the equation
  \[
    \overline\sigma_a\circ q_a=\sigma_a.
  \]
\end{definition}

\begin{proposition}[Affine strict-transform chart formula]
  \label{prop:nagata-strict-transform-chart-iso}
  \uses{def:nagata-strict-transform-chart-comparison}
  \lean{}

  Fix the expanded standing datum
  \[
    \mathsf D_{\mathrm{alg}}=(A,B,I,a,\varphi),\qquad
    A,B\text{ commutative rings},\quad
    I\triangleleft A,\quad a\in I,\quad
    \varphi\colon A\to B
  \]
  of (STAlg0) in
  \Cref{def:nagata-strict-transform-basechange-chart-ring}.
  The constructed comparison $\overline\sigma_a$ is a ring equivalence.
\end{proposition}

\begin{proof}
  \uses{lemma:nagata-strict-chart-map-surjective,
    lemma:nagata-strict-chart-map-kernel}
  By the two cited lemmas,
  \[
    \operatorname{im}(\sigma_a)=B[IB/b],\qquad
    \ker(\sigma_a)=\Gamma_{(a_C)}(C_a(A,B,I)).
  \]
  The first isomorphism theorem instantiated at $\sigma_a$ is the
  isomorphism
  \[
    C_a(A,B,I)/\ker(\sigma_a)
      \xrightarrow[\ [c]\mapsto\sigma_a(c)\ ]{\sim}B[IB/b].
  \]
  Substituting
  $C_a^{\mathrm{st}}=C_a/\Gamma_{(a_C)}(C_a)$ shows that this map is exactly
  $\overline\sigma_a$.
\end{proof}

\begin{definition}[Localized strict-transform chart comparison]
  \label{def:nagata-strict-transform-chart-comparison-localized}
  \uses{def:nagata-strict-transform-chart-comparison}
  \lean{}

  Fix the expanded standing datum
  \[
    \mathsf D_{\mathrm{alg}}=(A,B,I,a,\varphi),\qquad
    A,B\text{ commutative rings},\quad
    I\triangleleft A,\quad a\in I,\quad
    \varphi\colon A\to B
  \]
  of (STAlg0) in
  \Cref{def:nagata-strict-transform-basechange-chart-ring}.
  Let $s\in A$.  Define the localized ring homomorphism and the localized
  standing datum by
  \[
  \begin{gathered}
    A_s:=S_s^{-1}A,\qquad
    B_s:=S_{\varphi(s)}^{-1}B,\qquad
    \varphi_s\colon A_s\to B_s,\qquad
    \varphi_s(r/s^n):=\varphi(r)/\varphi(s)^n,\\
    I_s:=IA_s\triangleleft A_s,\qquad
    a_s:=a/1\in I_s,\qquad
    b:=\varphi(a),\qquad
    b_s:=\varphi_s(a_s)=b/1,\\
    \mathsf D_{\mathrm{alg},s}
      :=(A_s,B_s,I_s,a_s,\varphi_s).
  \end{gathered}                                                     \tag{STAlg1}
  \]
  Define
  \[
  \begin{aligned}
    \mu_{a,s}\colon
      C_a^{\mathrm{st}}(A,B,I)_{\varphi(s)}
        &\xrightarrow{\sim}C_{a_s}^{\mathrm{st}}(A_s,B_s,I_s),\\
    \nu_{a,s}\colon
      B[IB/b]_{\varphi(s)}
        &\xrightarrow{\sim}B_s[I_sB_s/b_s]
  \end{aligned}
  \]
  by
  \[
    \mu_{a,s}\left(
      \frac{[y\otimes(x/a)]}{\varphi(s)^n}\right)
      :=\left[\frac y{\varphi(s)^n}\otimes\frac{x/1}{a_s}\right],
    \qquad
    \nu_{a,s}\left(\frac{z/b^r}{\varphi(s)^n}\right)
      :=\frac{z/\varphi(s)^n}{b_s^r}.
  \]
  The first map uses
  \Cref{lemma:nagata-power-torsion-localization}.  Define
  \[
    (\overline\sigma_a)_s
      :=\nu_{a,s}\circ
        S_{\varphi(s)}^{-1}(\overline\sigma_a)\circ\mu_{a,s}^{-1}.
  \]
\end{definition}

\begin{lemma}[Strict-transform chart formula localizes]
  \label{lemma:nagata-strict-transform-chart-localization}
  \uses{def:nagata-strict-transform-chart-comparison-localized}
  \lean{}

  Fix the expanded standing datum
  \[
    \mathsf D_{\mathrm{alg}}=(A,B,I,a,\varphi),\qquad
    A,B\text{ commutative rings},\quad
    I\triangleleft A,\quad a\in I,\quad
    \varphi\colon A\to B
  \]
  of (STAlg0) in
  \Cref{def:nagata-strict-transform-basechange-chart-ring}, and fix $s\in A$.
  Fix also the expanded localized datum
  \[
  \begin{gathered}
    \mathsf D_{\mathrm{alg},s}
      =(A_s,B_s,I_s,a_s,\varphi_s),\\
    A_s=S_s^{-1}A,\qquad
    B_s=S_{\varphi(s)}^{-1}B,\qquad
    \varphi_s(r/s^n)=\varphi(r)/\varphi(s)^n,\\
    I_s=IA_s\triangleleft A_s,\qquad
    a_s=a/1\in I_s,\qquad
    b_s=\varphi_s(a_s)=\varphi(a)/1
  \end{gathered}
  \]
  of (STAlg1) in
  \Cref{def:nagata-strict-transform-chart-comparison-localized}.
  If
  \[
    \overline\sigma_{a_s}\colon
      C_{a_s}^{\mathrm{st}}(A_s,B_s,I_s)
      \to B_s[I_sB_s/b_s]
  \]
  denotes \Cref{def:nagata-strict-transform-chart-comparison} for the
  localized data.  The asserted localization compatibility is
  \[
    (\overline\sigma_a)_s=\overline\sigma_{a_s}.
  \]
\end{lemma}

\begin{proof}
  \uses{lemma:nagata-power-torsion-localization,
    prop:nagata-strict-transform-chart-iso}
  For $y\in B$, $x\in I$, and $n\in\mathbf N$,
  \[
  \begin{aligned}
    (\overline\sigma_a)_s\left(
      \left[\frac y{\varphi(s)^n}\otimes\frac{x/1}{a_s}\right]\right)
      &=\nu_{a,s}\left(
        \frac{\overline\sigma_a([y\otimes(x/a)])}{\varphi(s)^n}\right)\\
      &=\frac{(y/\varphi(s)^n)(\varphi(x)/1)}{b_s}\\
      &=\overline\sigma_{a_s}\left(
        \left[\frac y{\varphi(s)^n}\otimes\frac{x/1}{a_s}\right]\right).
  \end{aligned}
  \]
  Ring-map extensionality applied to the displayed equality on the generating
  classes is
  \[
    (\overline\sigma_a)_s=\overline\sigma_{a_s}.
  \]
\end{proof}

\begin{lemma}[Cocycle for strict-transform charts]
  \label{lemma:nagata-strict-transform-chart-cocycle}
  \uses{def:nagata-strict-transform-chart-comparison}
  \lean{}

  Let $A$ and $B$ be commutative rings, let $I\triangleleft A$, and let
  \(\varphi\colon A\to B\) be a ring homomorphism.
  Let $r,s\in I$, put $b_r:=\varphi(r)$ and $b_s:=\varphi(s)$, and define
  \[
  \begin{gathered}
    u_{r,s}:=q_r\bigl(1\otimes(s/r)\bigr)
      \in C_r^{\mathrm{st}}(A,B,I),\qquad
    v_{r,s}:=b_s/b_r\in B[IB/b_r],\\
    C^{\mathrm{st}}_{r;s}:=
      C_r^{\mathrm{st}}(A,B,I)_{u_{r,s}},\qquad
    E_{r;s}:=B[IB/b_r]_{v_{r,s}}.
  \end{gathered}
  \]
  Localizing $\overline\sigma_r$ defines
  \[
    \overline\sigma_{r;s}\colon
      C^{\mathrm{st}}_{r;s}\xrightarrow{\sim}E_{r;s},
    \qquad
    \overline\sigma_{r;s}(c/u_{r,s}^n)
      =\overline\sigma_r(c)/v_{r,s}^n.
  \]
  Define the extended-blow-up overlap isomorphism
  \[
    \omega^B_{r,s}\colon E_{r;s}\xrightarrow{\sim}E_{s;r}
  \]
  by the equations, for $y\in B$ and $x\in I$,
  \[
  \begin{aligned}
    \omega^B_{r,s}(y)&=y,\\
    \omega^B_{r,s}\!\left(\frac{\varphi(x)}{b_r}\right)
      &=\frac{\varphi(x)/b_s}{b_r/b_s},\\
    \omega^B_{r,s}(v_{r,s}^{-1})&=b_r/b_s.
  \end{aligned}
  \]
  Define the strict-chart ring isomorphism
  \[
    \rho^{\mathrm{alg}}_{r,s}:=
      \overline\sigma_{s;r}^{-1}\circ\omega^B_{r,s}
        \circ\overline\sigma_{r;s}
      \colon C^{\mathrm{st}}_{r;s}\xrightarrow{\sim}
        C^{\mathrm{st}}_{s;r}.
  \]
  Put
  \[
    Q_r:=\operatorname{Spec}(C_r^{\mathrm{st}}(A,B,I)),\qquad
    Q_{r;s}:=D(u_{r,s})\subseteq Q_r,
  \]
  and define the scheme overlap isomorphism
  \[
    \rho_{r,s}:=
      \operatorname{Spec}\bigl((\rho^{\mathrm{alg}}_{r,s})^{-1}\bigr)
      \colon Q_{r;s}\xrightarrow{\sim}Q_{s;r}.
  \]
  For $r,s,t\in I$, put
  \[
  \begin{gathered}
    C^{\mathrm{st}}_{r;s,t}:=
      C_r^{\mathrm{st}}(A,B,I)_{u_{r,s}u_{r,t}},\qquad
    E_{r;s,t}:=B[IB/b_r]_{v_{r,s}v_{r,t}},\\
    \overline\sigma_{r;s,t}\colon
      C^{\mathrm{st}}_{r;s,t}\xrightarrow{\sim}E_{r;s,t},\qquad
    Q_{r;s,t}:=\operatorname{Spec}(C^{\mathrm{st}}_{r;s,t})
      =Q_{r;s}\cap Q_{r;t}.
  \end{gathered}
  \]
  Let $\omega^B_{r,s;t}\colon E_{r;s,t}\to E_{s;r,t}$ and
  $\rho^{\mathrm{alg}}_{r,s;t}\colon
    C^{\mathrm{st}}_{r;s,t}\to C^{\mathrm{st}}_{s;r,t}$ denote the
  localizations of $\omega^B_{r,s}$ and $\rho^{\mathrm{alg}}_{r,s}$.
  On $Q_{r;s,t}$ one has
  \[
    \rho_{s,t}|_{\rho_{r,s}(Q_{r;s,t})}\circ
      \rho_{r,s}|_{Q_{r;s,t}}
      =\rho_{r,t}|_{Q_{r;s,t}}.
  \]
\end{lemma}

\begin{proof}
  \uses{prop:nagata-strict-transform-chart-iso}
  The source and target have the presentations
  \[
  \begin{aligned}
    E_{r;s}
      &=B[\,\varphi(x)/b_r\mid x\in I\,]
          [(b_s/b_r)^{-1}],\\
    E_{s;r}
      &=B[\,\varphi(x)/b_s\mid x\in I\,]
          [(b_r/b_s)^{-1}].
  \end{aligned}
  \]
  The relations respected by the three displayed formulas are
  \[
    b_r(\varphi(x)/b_r)=\varphi(x),\qquad
    (b_s/b_r)(b_r/b_s)=1.                                  \tag{STCoc1}
  \]
  The localization universal property applied to (STCoc1) defines
  \[
    \omega^B_{r,s}\in\operatorname{Hom}(E_{r;s},E_{s;r}).
  \]
  On the
  generators $y$, $\varphi(x)/b_r$, and $v_{r,s}^{-1}$, direct substitution
  is the pair of identities
  \[
    \omega^B_{s,r}\circ\omega^B_{r,s}=\operatorname{id}_{E_{r;s}},
    \qquad
    \omega^B_{r,s}\circ\omega^B_{s,r}=\operatorname{id}_{E_{s;r}},
    \qquad
    \operatorname{IsIso}(\omega^B_{r,s}),
    \qquad
    (\omega^B_{r,s})^{-1}=\omega^B_{s,r}.                    \tag{STCoc2}
  \]
  On the localization of
  the $r$-chart where both $v_{r,s}$ and $v_{r,t}$ are invertible, the same
  generator calculation is the identity
  \[
    \omega^B_{s,t;r}\circ\omega^B_{r,s;t}
      =\omega^B_{r,t;s}.                                          \tag{*}
  \]
  By \Cref{prop:nagata-strict-transform-chart-iso}, each
  $\overline\sigma_{r;s}$ is an isomorphism.  The definition of
  $\rho^{\mathrm{alg}}$ and equation $(*)$ expand the composite as the chain
  \[
  \begin{aligned}
    \rho^{\mathrm{alg}}_{s,t;r}\circ\rho^{\mathrm{alg}}_{r,s;t}
      &=\overline\sigma_{t;r,s}^{-1}\circ\omega^B_{s,t;r}
          \circ\overline\sigma_{s;r,t}
          \circ\overline\sigma_{s;r,t}^{-1}
          \circ\omega^B_{r,s;t}\circ\overline\sigma_{r;s,t}\\
      &=\overline\sigma_{t;r,s}^{-1}\circ
          (\omega^B_{s,t;r}\circ\omega^B_{r,s;t})
          \circ\overline\sigma_{r;s,t}\\
      &=\overline\sigma_{t;r,s}^{-1}\circ\omega^B_{r,t;s}
          \circ\overline\sigma_{r;s,t}\\
      &=\rho^{\mathrm{alg}}_{r,t;s}.
  \end{aligned}
  \]
  Taking inverses in the last equality is the identity
  \[
    (\rho^{\mathrm{alg}}_{r,s;t})^{-1}\circ
      (\rho^{\mathrm{alg}}_{s,t;r})^{-1}
      =(\rho^{\mathrm{alg}}_{r,t;s})^{-1}.
  \]
  Applying the contravariant functor $\operatorname{Spec}$ to the preceding
  ring-homomorphism identity is the scheme-morphism chain
  \[
  \begin{aligned}
    \rho_{s,t}|_{\rho_{r,s}(Q_{r;s,t})}\circ
      \rho_{r,s}|_{Q_{r;s,t}}
      &=\operatorname{Spec}\!\left(
          (\rho^{\mathrm{alg}}_{r,s;t})^{-1}\circ
          (\rho^{\mathrm{alg}}_{s,t;r})^{-1}\right)\\
      &=\operatorname{Spec}\!\left(
          (\rho^{\mathrm{alg}}_{r,t;s})^{-1}\right)
       =\rho_{r,t}|_{Q_{r;s,t}}.
  \end{aligned}
  \]
\end{proof}

\begin{definition}[Strict transform of a quasi-coherent module]
  \label{def:nagata-module-strict-transform}
  \uses{def:nagata-power-torsion}
  \lean{}

  Let $f\colon Y\to X$, let
  $\pi=b_{X,\mathcal I}\colon X'=\operatorname{Bl}_{\mathcal I}(X)\to X$,
  and put
  \[
    Y':=Y\times_XX',\qquad
    p\colon Y'\to Y,\qquad q\colon Y'\to X'.
  \]
  Define the exceptional ideal on $Y'$ by
  \[
    \mathcal E_Y
      :=q^{-1}(\mathcal I\mathcal O_{X'})\mathcal O_{Y'}
      =p^{-1}(\mathcal I\mathcal O_Y)\mathcal O_{Y'}.
  \]
  For an $\mathcal O_{Y'}$-module $\mathcal G$, define the
  $\mathcal E_Y$-power-torsion subsheaf $\Gamma_{\mathcal E_Y}(\mathcal G)$
  by the stalk equalities
  \[
    \Gamma_{\mathcal E_Y}(\mathcal G)_y
      :=\{m\in\mathcal G_y\mid
           \exists n\in\mathbf N,\ (\mathcal E_{Y,y})^n m=0\}
      \subseteq\mathcal G_y
      \qquad(y\in Y').
  \]
  For a quasi-coherent $\mathcal O_Y$-module $\mathcal F$, define its module
  strict transform to be
  \[
    \operatorname{ST}^{\mathrm{mod}}_\pi(\mathcal F)
      :=p^*\mathcal F\big/
         \Gamma_{\mathcal E_Y}(p^*\mathcal F).
  \]
\end{definition}

\begin{lemma}[Affine formula for a module strict transform]
  \label{lemma:nagata-module-strict-transform-chart}
  \uses{def:nagata-module-strict-transform,def:nagata-power-torsion}
  \lean{}

  Suppose $X=\operatorname{Spec}(A)$, $Y=\operatorname{Spec}(B)$,
  $\mathcal F=\widetilde M$, and $a\in I$.  With
  $C=C_a(A,B,I)$ and $a_C$ as in
  \Cref{def:nagata-strict-transform-basechange-chart-ring}, the restriction
  of $\operatorname{ST}^{\mathrm{mod}}_\pi(\mathcal F)$ to
  $\operatorname{Spec}(C)$ is
  \[
    \widetilde{\bigl(M\otimes_AA[I/a]\bigr)\big/
      \Gamma_{(a_C)}\bigl(M\otimes_AA[I/a]\bigr)}.
  \]
\end{lemma}

\begin{proof}
  Let $N:=M\otimes_AA[I/a]$.  Under the affine pullback formula and the chart
  identification of
  \Cref{lemma:nagata-strict-transform-basechange-chart}, one has
  \[
    (p^*\widetilde M)|_{\operatorname{Spec}(C)}=\widetilde N.
  \]
  By \Cref{lemma:nagata-chart-extended-ideal-principal}, the restriction of
  $\mathcal E_Y$ to this chart is the principal ideal sheaf
  \[
    \mathcal E_Y|_{\operatorname{Spec}(C)}=(a_C)^{\sim}.
  \]
  Substitution of the principal-ideal equality into
  \Cref{def:nagata-module-strict-transform} is, for every
  $\mathfrak p\subset C$, the stalk chain
  \[
  \begin{aligned}
    \Gamma_{\mathcal E_Y}(\widetilde N)_{\mathfrak p}
      &=\{m/1\in N_{\mathfrak p}\mid
           \exists n,\ a_C^n(m/1)=0\}\\
      &=\Gamma_{(a_C/1)}(N_{\mathfrak p})\\
      &=\Gamma_{(a_C)}(N)_{\mathfrak p}
        &&\text{by \Cref{lemma:nagata-power-torsion-localization}}.
  \end{aligned}
  \]
  Stalkwise extensionality for sheaves, applied to
  \[
    \forall\mathfrak p\in\operatorname{Spec}(C),\quad
    \Gamma_{\mathcal E_Y}(\widetilde N)_{\mathfrak p}
      =\widetilde{\Gamma_{(a_C)}(N)}_{\mathfrak p},
  \]
  is the equality
  \[
    \Gamma_{\mathcal E_Y}(\widetilde N)
      =\widetilde{\Gamma_{(a_C)}(N)}.
  \]
  The quotient equality in
  \Cref{def:nagata-module-strict-transform} is
  \[
  \begin{aligned}
    \operatorname{ST}^{\mathrm{mod}}_\pi(\mathcal F)
      |_{\operatorname{Spec}(C)}
      &=\widetilde N\big/
        \Gamma_{\mathcal E_Y}(\widetilde N)\\
      &=\widetilde N\big/
        \widetilde{\Gamma_{(a_C)}(N)}\\
      &=\widetilde{N/\Gamma_{(a_C)}(N)}.
  \end{aligned}
  \]
  which is the asserted formula.
\end{proof}

\begin{lemma}[Quasi-coherence of the module strict transform]
  \label{lemma:nagata-module-strict-transform-qc}
  \uses{def:nagata-module-strict-transform}
  \lean{}

  The module strict transform is quasi-coherent.
\end{lemma}

\begin{proof}
  \uses{lemma:nagata-module-strict-transform-chart,
    lemma:nagata-power-torsion-localization}
  On every affine chart $W_a=\operatorname{Spec}(C_a(A,B,I))$,
  \Cref{lemma:nagata-module-strict-transform-chart} identifies the
  restriction with $\widetilde{N_a/T_a}$, where
  \[
    N_a=M\otimes_AA[I/a],
    \qquad T_a=\Gamma_{(a_C)}(N_a).
  \]
  If $D(s)\subseteq W_a$, the exact localization functor and
  \Cref{lemma:nagata-power-torsion-localization} occur in the chain
  \[
  \begin{aligned}
    (N_a/T_a)_s
      &= (N_a)_s/(T_a)_s
         &&\text{by exactness of }S_s^{-1}(-)\\
      &= (N_a)_s/\Gamma_{(a_C/1)}((N_a)_s)
         &&\text{by \Cref{lemma:nagata-power-torsion-localization}}.
  \end{aligned}
  \]
  The affine quasi-coherence criterion applied to this displayed equality is
  the implication
  \[
    \left[
      \forall D(s)\subseteq W_a,\quad
      (N_a/T_a)_s=(N_a)_s/\Gamma_{(a_C/1)}((N_a)_s)
    \right]
    \Longrightarrow
    \operatorname{IsQuasiCoherent}
      \bigl(\operatorname{ST}^{\mathrm{mod}}_\pi(\mathcal F)\bigr).
  \]
\end{proof}

\begin{definition}[Closed chart model of the global strict transform]
  \label{def:nagata-strict-transform-closed-chart-model}
  \uses{def:nagata-strict-transform-chart-ring,
    def:nagata-direct-image-structure-map,
    lemma:nagata-strict-transform-basechange-chart,
    lemma:nagata-strict-transform-chart-cocycle}
  \lean{}

  Let $f\colon Y\to X$, let
  $\mathcal I\subseteq\mathcal O_X$ be a finite-type quasi-coherent ideal,
  and let
  $\pi=b_{X,\mathcal I}\colon X'=\operatorname{Bl}_{\mathcal I}(X)\to X$.
  Let $\mathscr A(f,\mathcal I)$ be the set of tuples
  \[
    \bigl(R,(L_r)_{r\in R},
          (V_r=\operatorname{Spec}(A_r))_{r\in R},\,
          (W_{r\ell}=\operatorname{Spec}(B_{r\ell}))_
            {r\in R,\ell\in L_r},\,
          (a_{r0},\ldots,a_{rm_r})_{r\in R}\bigr)
  \]
  satisfying
  \[
    X=\bigcup_{r\in R}V_r,\qquad
    f^{-1}(V_r)=\bigcup_{\ell\in L_r}W_{r\ell},\qquad
    \Gamma(V_r,\mathcal I)=(a_{r0},\ldots,a_{rm_r})
      \quad(r\in R).
  \]
  Here every $V_r\hookrightarrow X$ is an affine open immersion and every
  $W_{r\ell}\hookrightarrow Y$ is an affine open immersion whose image is
  contained in $f^{-1}(V_r)$.
  The affine-basis theorem and the affine-local finite-generation criterion
  for the finite-type ideal $\mathcal I$ assert the existence statement
  \[
    \exists\bigl(R,(L_r)_{r\in R},(V_r)_{r\in R},
      (W_{r\ell})_{r\in R,\ell\in L_r},
      (a_{r0},\ldots,a_{rm_r})_{r\in R}\bigr)
      \in\mathscr A(f,\mathcal I).
  \]
  Equivalently,
  $\mathscr A(f,\mathcal I)\ne\varnothing$.  By
  \Cref{def:nagata-choice-notation}, define the displayed triple by
  \[
    \bigl(R,(L_r)_{r\in R},(V_r)_{r\in R},
      (W_{r\ell})_{r\in R,\ell\in L_r},
      (a_{r0},\ldots,a_{rm_r})_{r\in R}\bigr)
      :=\operatorname{ch}(\mathscr A(f,\mathcal I)).
  \]
  Let
  \[
    v_r\colon V_r\hookrightarrow X,\qquad
    w_{r\ell}\colon W_{r\ell}\hookrightarrow Y,\qquad
    f_{r\ell}\colon W_{r\ell}\longrightarrow V_r,
    \qquad
    v_r\circ f_{r\ell}=f\circ w_{r\ell}.
  \]
  The morphism $f_{r\ell}$ is the unique factor through the open immersion
  $v_r$.  Define
  \[
    \varphi_{r\ell}:=f_{r\ell}^{\#}\colon A_r\to B_{r\ell},
    \qquad I_r:=\Gamma(V_r,\mathcal I).
  \]
  Define the chart-index set
  \[
    \Lambda:=\coprod_{r\in R}\coprod_{\ell\in L_r}
      \{0,\ldots,m_r\}.
  \]
  For each $\alpha=(r,\ell,k)\in\Lambda$ define
  \[
  \begin{gathered}
    C_\alpha=C_{r\ell k}:=C_{a_{rk}}(A_r,B_{r\ell},I_r),\qquad
    C^{\mathrm{st}}_\alpha=C^{\mathrm{st}}_{r\ell k}:=
      C_{a_{rk}}^{\mathrm{st}}(A_r,B_{r\ell},I_r),\\
    P_\alpha=P_{r\ell k}:=\operatorname{Spec}(C_{r\ell k}),\qquad
    Q_\alpha=Q_{r\ell k}:=\operatorname{Spec}(C^{\mathrm{st}}_{r\ell k}),\\
    \eta^B_\alpha\colon B_{r\ell}\to C_\alpha,\quad
      y\mapsto y\otimes1,\qquad
    \eta^A_\alpha\colon A_r[I_r/a_{rk}]\to C_\alpha,\quad
      q\mapsto1\otimes q,\\
    \kappa_\alpha=\kappa_{r\ell k}:=\operatorname{Spec}(q_{a_{rk}})
      \colon Q_\alpha\hookrightarrow P_\alpha.
  \end{gathered}
  \]
  Define
  \[
  \begin{gathered}
    B_{V_r}^X:=b_{X,\mathcal I}^{-1}(V_r),\qquad
    \jmath_{V_r}\colon B_{V_r}^X\hookrightarrow X',\\
    q_{V_r}\colon B_{V_r}^X\xrightarrow{\sim}
      \operatorname{Bl}_{I_r}(V_r)
      \quad\text{as in
        \Cref{def:nagata-global-blowup-affine-restriction-map}},\\
    \delta_{rk}\colon D_+(a_{rk}t)
      \hookrightarrow\operatorname{Bl}_{I_r}(V_r),\\
    d_{rk}:=\jmath_{V_r}\circ q_{V_r}^{-1}\circ\delta_{rk}
      \colon D_+(a_{rk}t)\hookrightarrow X'.
  \end{gathered}                                                   \tag{ST0a}
  \]
  Let
  $\beta_{r\ell k}$ be the isomorphism of
  \Cref{lemma:nagata-strict-transform-basechange-chart} for
  $(A_r,B_{r\ell},I_r,a_{rk})$.  Put
  \[
    H_{r\ell k}:=
      W_{r\ell}\times_{V_r}D_+(a_{rk}t),
    \qquad
    h_W\colon H_{r\ell k}\to W_{r\ell},
    \qquad
    h_D\colon H_{r\ell k}\to D_+(a_{rk}t),
  \]
  where $D_+(a_{rk}t)\to V_r$ is the restriction of $\pi$.  Define the
  fibre-product morphism
  \[
    \zeta_{r\ell k}\colon H_{r\ell k}\longrightarrow
      P:=Y\times_XX'
  \]
  by
  \[
    \operatorname{pr}_Y\circ\zeta_{r\ell k}
      =w_{r\ell}\circ h_W,
    \qquad
    \operatorname{pr}_{X'}\circ\zeta_{r\ell k}
      =d_{rk}\circ h_D.
  \]
  Its defining compatibility is the chain
  \[
    f\circ w_{r\ell}\circ h_W
      =v_r\circ f_{r\ell}\circ h_W
      =\pi\circ d_{rk}\circ h_D.
  \]
  Regard $\beta_{r\ell k}$ as the specified isomorphism
  $\beta_{r\ell k}\colon P_\alpha\xrightarrow{\sim}H_{r\ell k}$.  Define
  the ambient chart immersion
  \[
    o_\alpha=o_{r\ell k}:=
      \zeta_{r\ell k}\circ\beta_{r\ell k}
      \colon P_\alpha\hookrightarrow P:=Y\times_XX'.
  \]
  Stability of open immersions under base change and composition is the
  predicate chain
  \[
  \begin{aligned}
    \operatorname{IsOpenImmersion}(w_{r\ell})
      &\wedge\operatorname{IsOpenImmersion}(d_{rk})\\
      &\Longrightarrow
        \operatorname{IsOpenImmersion}(\zeta_{r\ell k})
        \wedge\operatorname{IsOpenImmersion}(o_\alpha).
  \end{aligned}
  \]
  For indices $\alpha,\beta$, define the ambient overlap and its projections
  by
  \[
    P_{\alpha\beta}:=P_\alpha\times_PP_\beta,\qquad
    p^\alpha_{\alpha\beta}\colon P_{\alpha\beta}\to P_\alpha,\qquad
    p^\beta_{\alpha\beta}\colon P_{\alpha\beta}\to P_\beta,
  \]
  and define the two restrictions of the closed charts by
  \[
  \begin{aligned}
    Q_{\alpha\beta}^{\alpha}
      &:=Q_\alpha\times_{P_\alpha}P_{\alpha\beta},&
    \kappa_{\alpha\beta}^{\alpha}
      &\colon Q_{\alpha\beta}^{\alpha}\hookrightarrow P_{\alpha\beta},\\
    Q_{\alpha\beta}^{\beta}
      &:=Q_\beta\times_{P_\beta}P_{\alpha\beta},&
    \kappa_{\alpha\beta}^{\beta}
      &\colon Q_{\alpha\beta}^{\beta}\hookrightarrow P_{\alpha\beta}.
  \end{aligned}
  \]
  On an affine principal cover of $P_{\alpha\beta}$,
  \Cref{lemma:nagata-strict-transform-chart-cocycle} identifies the two
  quotient ideals; explicitly,
  \[
    \ker\bigl(\mathcal O_{P_{\alpha\beta}}\to
      (\kappa_{\alpha\beta}^{\alpha})_*
        \mathcal O_{Q_{\alpha\beta}^{\alpha}}\bigr)
    =
    \ker\bigl(\mathcal O_{P_{\alpha\beta}}\to
      (\kappa_{\alpha\beta}^{\beta})_*
        \mathcal O_{Q_{\alpha\beta}^{\beta}}\bigr).               \tag{ST1}
  \]
  Define the unique overlap isomorphism
  \[
    \rho_{\alpha\beta}\colon
      Q_{\alpha\beta}^{\alpha}\xrightarrow{\sim}
      Q_{\alpha\beta}^{\beta},
    \qquad
    \kappa_{\alpha\beta}^{\beta}\circ\rho_{\alpha\beta}
      =\kappa_{\alpha\beta}^{\alpha}.                            \tag{ST2}
  \]
  For $\alpha=\beta$, equation (ST2) is
  \[
    \kappa_{\alpha\alpha}^{\alpha}\circ\rho_{\alpha\alpha}
      =\kappa_{\alpha\alpha}^{\alpha}
      =\kappa_{\alpha\alpha}^{\alpha}\circ
        \operatorname{id}_{Q_\alpha},
  \]
  Monicity of $\kappa_{\alpha\alpha}^{\alpha}$ is the implication
  \[
    \kappa_{\alpha\alpha}^{\alpha}\circ\rho_{\alpha\alpha}
      =\kappa_{\alpha\alpha}^{\alpha}\circ
        \operatorname{id}_{Q_\alpha}
    \quad\Longrightarrow\quad
    \rho_{\alpha\alpha}=\operatorname{id}_{Q_\alpha}.
  \]
  Equations (ST2) for
  $(\alpha,\beta)$ and $(\beta,\alpha)$ form the two chains
  \[
  \begin{aligned}
    \kappa_{\alpha\beta}^{\alpha}\circ
      (\rho_{\beta\alpha}\circ\rho_{\alpha\beta})
      &=\kappa_{\alpha\beta}^{\beta}\circ\rho_{\alpha\beta}
       =\kappa_{\alpha\beta}^{\alpha},\\
    \kappa_{\alpha\beta}^{\beta}\circ
      (\rho_{\alpha\beta}\circ\rho_{\beta\alpha})
      &=\kappa_{\alpha\beta}^{\alpha}\circ\rho_{\beta\alpha}
       =\kappa_{\alpha\beta}^{\beta}.
  \end{aligned}
  \]
  Cancellation against the two monomorphisms is the pair of equalities
  \[
    \rho_{\alpha\alpha}=\operatorname{id}_{Q_\alpha},
    \qquad
    \rho_{\beta\alpha}=\rho_{\alpha\beta}^{-1}.                  \tag{ST3}
  \]
  Define
  \[
    P_{\alpha\beta\gamma}:=
      P_\alpha\times_PP_\beta\times_PP_\gamma,
    \qquad
    Q_{\alpha\beta\gamma}^{\delta}:=
      Q_\delta\times_{P_\delta}P_{\alpha\beta\gamma},\qquad
    \kappa_{\alpha\beta\gamma}^{\delta}\colon
      Q_{\alpha\beta\gamma}^{\delta}
        \hookrightarrow P_{\alpha\beta\gamma}
      \quad(\delta\in\{\alpha,\beta,\gamma\}),
  \]
  and let
  \[
    \rho_{\alpha\beta\mid\gamma}\colon
      Q_{\alpha\beta\gamma}^{\alpha}\xrightarrow{\sim}
      Q_{\alpha\beta\gamma}^{\beta}
  \]
  be the base change of $\rho_{\alpha\beta}$.  After base change, (ST2)
  supplies the chain
  \[
  \begin{aligned}
    \kappa_{\alpha\beta\gamma}^{\gamma}\circ
      (\rho_{\beta\gamma\mid\alpha}\circ
        \rho_{\alpha\beta\mid\gamma})
      &=\kappa_{\alpha\beta\gamma}^{\beta}\circ
        \rho_{\alpha\beta\mid\gamma}\\
      &=\kappa_{\alpha\beta\gamma}^{\alpha}
       =\kappa_{\alpha\beta\gamma}^{\gamma}\circ
        \rho_{\alpha\gamma\mid\beta},
  \end{aligned}
  \]
  where $\kappa_{\alpha\beta\gamma}^{\delta}$ is the displayed base change
  of $\kappa_\delta$.  Monicity of
  $\kappa_{\alpha\beta\gamma}^{\gamma}$ is the typed cocycle equality
  \[
    \rho_{\beta\gamma\mid\alpha}\circ
      \rho_{\alpha\beta\mid\gamma}
      =\rho_{\alpha\gamma\mid\beta}
      \colon Q_{\alpha\beta\gamma}^{\alpha}\xrightarrow{\sim}
        Q_{\alpha\beta\gamma}^{\gamma}.                          \tag{ST4}
  \]
  Define the two overlap inclusions and record their base-change equations:
  \[
  \begin{gathered}
    j_{\alpha\beta}^{\alpha}:=\operatorname{pr}_{Q_\alpha}
      \colon Q_{\alpha\beta}^{\alpha}\hookrightarrow Q_\alpha,
    \qquad
    j_{\alpha\beta}^{\beta}:=\operatorname{pr}_{Q_\beta}
      \colon Q_{\alpha\beta}^{\beta}\hookrightarrow Q_\beta,\\
    p_{\alpha\beta}^{\alpha}\circ
      \kappa_{\alpha\beta}^{\alpha}
      =\kappa_\alpha\circ j_{\alpha\beta}^{\alpha},
    \qquad
    p_{\alpha\beta}^{\beta}\circ
      \kappa_{\alpha\beta}^{\beta}
      =\kappa_\beta\circ j_{\alpha\beta}^{\beta}.
  \end{gathered}                                                   \tag{ST4a}
  \]
  For every scheme $T$, define the set of transition-compatible chart maps
  \[
    \mathscr E_T^{\mathrm{st}}:=
      \left\{
        (h_\alpha)_{\alpha\in\Lambda}
          \in\prod_{\alpha\in\Lambda}\operatorname{Hom}(Q_\alpha,T)
        \ \middle|\ 
        h_\alpha\circ j_{\alpha\beta}^{\alpha}
          =h_\beta\circ j_{\alpha\beta}^{\beta}
             \circ\rho_{\alpha\beta}
        \quad(\alpha,\beta\in\Lambda)
      \right\}.                                                   \tag{ST4b}
  \]
  Define the set of gluing realizations
  \[
  \begin{aligned}
    \mathscr G_{\mathrm{st}}:=
    \bigl\{(G,(e_\alpha)_{\alpha\in\Lambda})\ \bigm|\ {}&
      G\text{ is a scheme},\quad
      e_\alpha\colon Q_\alpha\to G,\quad
      \operatorname{IsOpenImmersion}(e_\alpha),\\
    &|G|=\bigcup_{\alpha\in\Lambda}|e_\alpha|(|Q_\alpha|),\\
    &e_\alpha\circ j_{\alpha\beta}^{\alpha}
       =e_\beta\circ j_{\alpha\beta}^{\beta}
          \circ\rho_{\alpha\beta},\\
    &\operatorname{IsIso}\!\left(
       (j_{\alpha\beta}^{\alpha},
          j_{\alpha\beta}^{\beta}\circ\rho_{\alpha\beta})
       \colon Q_{\alpha\beta}^{\alpha}
          \longrightarrow Q_\alpha\times_GQ_\beta\right)
       \quad(\alpha,\beta\in\Lambda),\\
    &\forall T,\quad
      \left[
      \begin{array}{rcl}
        \operatorname{Hom}(G,T)&\longrightarrow&\mathscr E_T^{\mathrm{st}}\\
        h&\longmapsto&(h\circ e_\alpha)_{\alpha\in\Lambda}
      \end{array}
      \right]\text{ is bijective}
      \bigr\}.
  \end{aligned}                                                     \tag{ST4c}
  \]
  Equations (ST3) and (ST4) are respectively the inverse and cocycle hypotheses
  of the scheme-gluing theorem
  \cite[\href{https://stacks.math.columbia.edu/tag/01JA}{Tag 01JA}]
  {stacks-project}; its existence and universal-property assertions are the
  symbolic statement
  \[
    \mathscr G_{\mathrm{st}}\neq\varnothing.                      \tag{ST4d}
  \]
  By \Cref{def:nagata-choice-notation}, define
  \[
    \bigl(Y^{\mathrm{st}}_{\mathrm{chart}},
      (e_\alpha^{\mathrm{st}})_{\alpha\in\Lambda}\bigr)
      :=\operatorname{ch}(\mathscr G_{\mathrm{st}})
      \in\mathscr G_{\mathrm{st}}.                               \tag{ST4e}
  \]
  Unpacking this membership defines the open immersions
  $e_\alpha^{\mathrm{st}}\colon
    Q_\alpha\hookrightarrow Y^{\mathrm{st}}_{\mathrm{chart}}$ and is the
  conjunction
  \[
  \begin{gathered}
    |Y^{\mathrm{st}}_{\mathrm{chart}}|
       =\bigcup_{\alpha\in\Lambda}
          |e_\alpha^{\mathrm{st}}|(|Q_\alpha|),\\
    e_\alpha^{\mathrm{st}}\circ j_{\alpha\beta}^{\alpha}
      =e_\beta^{\mathrm{st}}\circ j_{\alpha\beta}^{\beta}
        \circ\rho_{\alpha\beta},\\
    \Gamma_T^{\mathrm{st}}\colon
      \operatorname{Hom}(Y^{\mathrm{st}}_{\mathrm{chart}},T)
        \xrightarrow{\sim}\mathscr E_T^{\mathrm{st}},
    \qquad
    \Gamma_T^{\mathrm{st}}(h)
      :=(h\circ e_\alpha^{\mathrm{st}})_{\alpha\in\Lambda}.
  \end{gathered}                                                   \tag{ST4f}
  \]
  The transition compatibility of the family
  $(o_\alpha\circ\kappa_\alpha)_{\alpha\in\Lambda}$ is the chain
  \[
  \begin{aligned}
    o_\alpha\circ\kappa_\alpha\circ j_{\alpha\beta}^{\alpha}
      &=o_\alpha\circ p_{\alpha\beta}^{\alpha}\circ
          \kappa_{\alpha\beta}^{\alpha}\\
      &=o_\beta\circ p_{\alpha\beta}^{\beta}\circ
          \kappa_{\alpha\beta}^{\alpha}\\
      &=o_\beta\circ p_{\alpha\beta}^{\beta}\circ
          \kappa_{\alpha\beta}^{\beta}\circ\rho_{\alpha\beta}\\
      &=o_\beta\circ\kappa_\beta\circ
          j_{\alpha\beta}^{\beta}\circ\rho_{\alpha\beta}.
  \end{aligned}
  \]
  By the set-builder defining
  $\mathscr E_{Y\times_XX'}^{\mathrm{st}}$, the preceding transition chain is
  the membership
  \[
    (o_\alpha\circ\kappa_\alpha)_{\alpha\in\Lambda}
      \in\mathscr E_{Y\times_XX'}^{\mathrm{st}}.
  \]
  Define the ambient morphism by the inverse of the bijection in (ST4f):
  \[
    \kappa_{\mathrm{chart}}
      :=(\Gamma_{Y\times_XX'}^{\mathrm{st}})^{-1}
          \bigl((o_\alpha\circ\kappa_\alpha)_{\alpha\in\Lambda}\bigr)
      \colon Y^{\mathrm{st}}_{\mathrm{chart}}\longrightarrow Y\times_XX'.
  \]
  Applying $\Gamma_{Y\times_XX'}^{\mathrm{st}}$ to this definition is the
  componentwise equality
  \[
    \kappa_{\mathrm{chart}}\circ e_\alpha^{\mathrm{st}}
      =o_\alpha\circ\kappa_\alpha
      \qquad(\alpha\in\Lambda).                                  \tag{ST5}
  \]
\end{definition}

\begin{proposition}[Global strict transform is a closed subscheme]
  \label{prop:nagata-strict-transform-closed-fiber-product}
  \uses{def:nagata-closed-subscheme-class,
    def:nagata-strict-transform-closed-chart-model}
  \lean{}

  Put $P:=Y\times_XX'$.  The constructed map
  $\kappa_{\mathrm{chart}}\colon
    Y^{\mathrm{st}}_{\mathrm{chart}}\hookrightarrow P$
  is a closed immersion, and
  \[
    [\,Y^{\mathrm{st}}_{\mathrm{chart}}
      \xhookrightarrow{\kappa_{\mathrm{chart}}}P\,]
      =[\,\operatorname{ST}_\pi(Y)
      \xhookrightarrow{\iota_{f,\pi}}P\,]
      \quad\text{in }\operatorname{ClosedSub}(P).                \tag{ST6a}
  \]
\end{proposition}

\begin{proof}
  \uses{def:nagata-choice-notation,def:nagata-closed-subscheme-class,
    lemma:nagata-immersion-qcqs-criteria,
    lemma:nagata-morphism-property-calculus,
    lemma:nagata-strict-transform-chart-closed,
    lemma:nagata-strict-transform-chart-cocycle,
    lemma:nagata-scheme-image-restrict-open}
  Put
  \[
    U:=X\setminus V(\mathcal I),
    \qquad j_U\colon U\hookrightarrow X,
    \qquad Y_U:=Y\times_XU,
    \qquad P:=Y\times_XX'.
  \]
  On an affine open $V=\operatorname{Spec}(A)\subseteq X$, finite type of
  $\mathcal I$ means
  \[
    \mathcal I|_V=\widetilde{(a_1,\ldots,a_n)},
    \qquad
    U\cap V=D(a_1,\ldots,a_n)=\bigcup_{i=1}^{n}D(a_i).
  \]
  Each $D(a_i)$ is affine.  The affine-local definition of a quasi-compact
  morphism is the predicate chain
  \[
    \forall V=\operatorname{Spec}(A)\subseteq X,\quad
      j_U^{-1}(V)=\bigcup_{i=1}^{n}D(a_i)\text{ is quasi-compact}
    \quad\Longrightarrow\quad
      \operatorname{IsQuasiCompact}(j_U).                       \tag{ST5c}
  \]
  Let
  \[
    \alpha_U\colon X'\times_XU\xrightarrow{\sim}U,
    \qquad
    \lambda_{f,\pi}\colon Y_U\to P
  \]
  be the morphisms of
  \Cref{thm:nagata-global-blowup-complement,def:nagata-strict-transform}.
  For $\alpha=(r,\ell,k)$, define
  \[
  \begin{gathered}
    a^C_\alpha:=
      \eta^B_\alpha(\varphi_{r\ell}(a_{rk}))
      =\eta^A_\alpha(a_{rk})\in C_\alpha,\\
    C_\alpha[(a^C_\alpha)^{-1}]
      :=\{(a^C_\alpha)^n\mid n\in\mathbf N\}^{-1}C_\alpha,\qquad
    \ell_\alpha\colon C_\alpha\to C_\alpha[(a^C_\alpha)^{-1}],
      \quad c\mapsto c/1,\\
    D_\alpha:=\operatorname{Spec}(C_\alpha[(a^C_\alpha)^{-1}]),\qquad
    d^{\mathrm{loc}}_\alpha:=\operatorname{Spec}(\ell_\alpha)
      \colon D_\alpha\hookrightarrow P_\alpha.
  \end{gathered}                                                  \tag{ST6}
  \]
  Put
  \[
    y_\alpha:=\operatorname{pr}_Y\circ o_\alpha\circ
      d^{\mathrm{loc}}_\alpha\colon D_\alpha\to Y,
    \qquad
    x_\alpha:=f\circ y_\alpha\colon D_\alpha\to X.
  \]
  Invertibility of $a^C_\alpha$ is the factorization of $x_\alpha$ through
  $D_X(a_{rk})\subseteq U$.  Define that factor by
  \[
    u_\alpha\colon D_\alpha\to U,
    \qquad
    j_U\circ u_\alpha=x_\alpha,
  \]
  and define the source transport
  \[
    s_\alpha:=(y_\alpha,u_\alpha)
      \colon D_\alpha\longrightarrow Y_U,
    \qquad
    \operatorname{pr}_Y\circ s_\alpha=y_\alpha,
    \quad
    \operatorname{pr}_U\circ s_\alpha=u_\alpha.               \tag{ST7}
  \]
  The two component calculations
  \[
  \begin{aligned}
    \operatorname{pr}_Y\circ\lambda_{f,\pi}\circ s_\alpha
      &=y_\alpha
       =\operatorname{pr}_Y\circ o_\alpha\circ d^{\mathrm{loc}}_\alpha,\\
    \operatorname{pr}_{X'}\circ\lambda_{f,\pi}\circ s_\alpha
      &=\operatorname{pr}_{X'}\circ\alpha_U^{-1}\circ u_\alpha
       =\operatorname{pr}_{X'}\circ o_\alpha\circ
          d^{\mathrm{loc}}_\alpha
  \end{aligned}
  \]
  and fibre-product uniqueness are the typed equality
  \[
    \lambda_{f,\pi}\circ s_\alpha
      =o_\alpha\circ d^{\mathrm{loc}}_\alpha
      \colon D_\alpha\to P.                                    \tag{ST8}
  \]
  Define the actual source of the restriction to $P_\alpha$ by
  \[
  \begin{gathered}
    R_\alpha:=Y_U\times_{P,\lambda_{f,\pi},o_\alpha}P_\alpha,
    \qquad
    r_U^\alpha\colon R_\alpha\to Y_U,
    \qquad
    r_P^\alpha\colon R_\alpha\to P_\alpha,\\
    \lambda_\alpha:=r_P^\alpha\colon R_\alpha\to P_\alpha,
    \qquad
    \tau_\alpha:=(s_\alpha,d^{\mathrm{loc}}_\alpha)
      \colon D_\alpha\longrightarrow R_\alpha.
  \end{gathered}                                                  \tag{ST9}
  \]
  Equation (ST8) defines $\tau_\alpha$.  The affine-chart identity in
  \Cref{prop:nagata-affine-blowup-chart,lemma:nagata-affine-blowup-complement}
  defines
  \[
    v_\alpha\colon R_\alpha\longrightarrow D_X(a_{rk})
  \]
  by the equations
  \[
  \begin{gathered}
    \iota_{a_{rk}}^U\circ v_\alpha
      =\operatorname{pr}_U\circ r_U^\alpha,\\
    (r_P^\alpha)^*(a^C_\alpha)
      =v_\alpha^*(a_{rk}|_{D_X(a_{rk})}).
  \end{gathered}                                                   \tag{ST9a}
  \]
  The section $a_{rk}|_{D_X(a_{rk})}$ has inverse
  $(a_{rk}|_{D_X(a_{rk})})^{-1}$, and pullback preserves the two inverse
  equations.  Substitution in (ST9a) is
  \[
  \begin{gathered}
    v_\alpha^*(a_{rk}|_{D_X(a_{rk})})
      \cdot v_\alpha^*((a_{rk}|_{D_X(a_{rk})})^{-1})=1,\\
    v_\alpha^*((a_{rk}|_{D_X(a_{rk})})^{-1})
      \cdot v_\alpha^*(a_{rk}|_{D_X(a_{rk})})=1,\\
    (r_P^\alpha)^*(a^C_\alpha)
      \in\Gamma(R_\alpha,\mathcal O_{R_\alpha})^\times.
  \end{gathered}
  \]
  The universal property of the
  distinguished open $d^{\mathrm{loc}}_\alpha\colon D_\alpha\hookrightarrow
  P_\alpha$ defines the unique factor
  \[
    r_P^\alpha=d^{\mathrm{loc}}_\alpha\circ\bar r_P^\alpha,
    \qquad
    \bar r_P^\alpha\colon R_\alpha\to D_\alpha.
  \]
  The projection calculations are
  \[
  \begin{aligned}
    d^{\mathrm{loc}}_\alpha\circ\bar r_P^\alpha\circ\tau_\alpha
      &=r_P^\alpha\circ\tau_\alpha
       =d^{\mathrm{loc}}_\alpha,\\
    r_P^\alpha\circ\tau_\alpha\circ\bar r_P^\alpha
      &=d^{\mathrm{loc}}_\alpha\circ\bar r_P^\alpha=r_P^\alpha,\\
    \lambda_{f,\pi}\circ r_U^\alpha\circ\tau_\alpha\circ
      \bar r_P^\alpha
      &=o_\alpha\circ r_P^\alpha\circ\tau_\alpha\circ
        \bar r_P^\alpha\\
      &=o_\alpha\circ r_P^\alpha
       =\lambda_{f,\pi}\circ r_U^\alpha.
  \end{aligned}
  \]
  After cancellation against the monomorphisms
  $d^{\mathrm{loc}}_\alpha$ and $\lambda_{f,\pi}$, the two component
  equalities in the preceding display satisfy the fibre-product uniqueness
  criterion; the resulting equalities are
  \[
    \bar r_P^\alpha\circ\tau_\alpha=\operatorname{id}_{D_\alpha},
    \qquad
    \tau_\alpha\circ\bar r_P^\alpha=\operatorname{id}_{R_\alpha}.
  \]
  The pair displays $\tau_\alpha$ as an isomorphism.  Its correctly typed
  restriction identity is
  \[
    \lambda_\alpha\circ\tau_\alpha=d^{\mathrm{loc}}_\alpha.   \tag{ST10}
  \]
  The base-change identity defining $R_\alpha$, the inverse equations for
  $\tau_\alpha$, (PC2)--(PC3) of
  \Cref{lemma:nagata-morphism-property-calculus},
  \Cref{def:nagata-strict-transform}, and
  \Cref{lemma:nagata-immersion-qcqs-criteria} form the predicate chain
  \[
  \begin{aligned}
    \operatorname{IsQCQS}(\lambda_{f,\pi})
      &\Longrightarrow\operatorname{IsQCQS}(\lambda_\alpha),\\
    \operatorname{IsIso}(\tau_\alpha)
      &\Longrightarrow\operatorname{IsQCQS}(\tau_\alpha),\\
    \operatorname{IsQCQS}(\lambda_\alpha)\wedge
      \operatorname{IsQCQS}(\tau_\alpha)
      &\Longrightarrow
        \operatorname{IsQCQS}(\lambda_\alpha\circ\tau_\alpha),\\
    (\lambda_\alpha\circ\tau_\alpha=d^{\mathrm{loc}}_\alpha)
      &\Longrightarrow\operatorname{IsQCQS}(d^{\mathrm{loc}}_\alpha).
  \end{aligned}                                                     \tag{ST10a}
  \]
  By \Cref{lemma:nagata-power-torsion-kernel-localization},
  \[
    \ker(\ell_\alpha)
      =\Gamma_{(a^C_\alpha)}(C_\alpha).
  \]
  For every $\alpha=(r,\ell,k)$, define the base-changed closed immersions
  \[
  \begin{gathered}
    \kappa_{\mathrm{chart},\alpha}\colon
      Y^{\mathrm{st}}_{\mathrm{chart}}\times_PP_\alpha
        \hookrightarrow P_\alpha,\\
    \iota_{f,\pi,\alpha}\colon
      \operatorname{ST}_\pi(Y)\times_PP_\alpha
        \hookrightarrow P_\alpha,\\
    i_{\lambda,\alpha}:=\operatorname{pr}_{P_\alpha}\colon
      \operatorname{SIm}_P(\lambda_{f,\pi})\times_PP_\alpha
        \hookrightarrow P_\alpha.
  \end{gathered}
  \]
  The first is the base change of $\kappa_{\mathrm{chart}}$ and the second
  is the base change of $\iota_{f,\pi}$.  Substitution of the restriction
  identity in \Cref{lemma:nagata-scheme-image-restrict-open}, source
  invariance under $\tau_\alpha$, (ST10), and the displayed kernel equality
  is the class chain
  \[
  \begin{aligned}
    [\,\operatorname{ST}_\pi(Y)\times_PP_\alpha
      \xhookrightarrow{\iota_{f,\pi,\alpha}}P_\alpha\,]
      &=[\,\operatorname{SIm}_P(\lambda_{f,\pi})\times_PP_\alpha
        \xhookrightarrow{i_{\lambda,\alpha}}P_\alpha\,]\\
      &=[\,\operatorname{SIm}_{P_\alpha}(\lambda_\alpha)
        \xhookrightarrow{i_{\lambda_\alpha}}P_\alpha\,]\\
      &=[\,\operatorname{SIm}_{P_\alpha}
          (\lambda_\alpha\circ\tau_\alpha)
        \xhookrightarrow{i_{\lambda_\alpha\circ\tau_\alpha}}P_\alpha\,]\\
      &=[\,\operatorname{SIm}_{P_\alpha}(d^{\mathrm{loc}}_\alpha)
        \xhookrightarrow{i_{d^{\mathrm{loc}}_\alpha}}P_\alpha\,]\\
      &=[\,Q_\alpha\xhookrightarrow{\kappa_\alpha}P_\alpha\,]
      \quad\text{in }\operatorname{ClosedSub}(P_\alpha),
  \end{aligned}                                                    \tag{ST11a}
  \]
  where the last equality is the explicit ideal chain
  \[
    \ker(\ell_\alpha)
      =\Gamma_{(a^C_\alpha)}(C_\alpha),
    \qquad
    Q_\alpha
      =\operatorname{Spec}\!\left(
        C_\alpha/\Gamma_{(a^C_\alpha)}(C_\alpha)\right).
  \]
  Equations (ST2)--(ST5) and
  \Cref{lemma:nagata-strict-transform-chart-cocycle} are the second class
  equality
  \[
    [\,Y^{\mathrm{st}}_{\mathrm{chart}}\times_PP_\alpha
      \xhookrightarrow{\kappa_{\mathrm{chart},\alpha}}P_\alpha\,]
      =[\,Q_\alpha\xhookrightarrow{\kappa_\alpha}P_\alpha\,]
      \quad\text{in }\operatorname{ClosedSub}(P_\alpha).       \tag{ST11b}
  \]
  Define the two witness sets
  \[
  \begin{aligned}
    \mathscr X^{\mathrm{chart}}_\alpha
      &:=\left\{\chi\in
        \operatorname{Hom}(Y^{\mathrm{st}}_{\mathrm{chart}}
          \times_PP_\alpha,Q_\alpha)
        \ \middle|\ 
        \operatorname{IsIso}(\chi)\ \wedge\
        \kappa_\alpha\circ\chi=\kappa_{\mathrm{chart},\alpha}\right\},\\
    \mathscr X^{\mathrm{im}}_\alpha
      &:=\left\{\chi\in
        \operatorname{Hom}(\operatorname{ST}_\pi(Y)
          \times_PP_\alpha,Q_\alpha)
        \ \middle|\ 
        \operatorname{IsIso}(\chi)\ \wedge\
        \kappa_\alpha\circ\chi=\iota_{f,\pi,\alpha}\right\}.
  \end{aligned}
  \]
  By (CSub3) of \Cref{def:nagata-closed-subscheme-class}, (ST11a)--(ST11b)
  are the two nonemptiness statements
  \[
    \mathscr X^{\mathrm{chart}}_\alpha\ne\varnothing,
    \qquad
    \mathscr X^{\mathrm{im}}_\alpha\ne\varnothing.
  \]
  Define the chosen witnesses and expand their membership as
  \[
  \begin{gathered}
    \chi^{\mathrm{chart}}_\alpha
      :=\operatorname{ch}(\mathscr X^{\mathrm{chart}}_\alpha)
      \in\mathscr X^{\mathrm{chart}}_\alpha,\\
    \chi^{\mathrm{chart}}_\alpha\colon
      Y^{\mathrm{st}}_{\mathrm{chart}}\times_PP_\alpha
        \xrightarrow{\sim}Q_\alpha,
    \qquad
    \kappa_\alpha\circ\chi^{\mathrm{chart}}_\alpha
      =\kappa_{\mathrm{chart},\alpha},\\
    \chi^{\mathrm{im}}_\alpha
      :=\operatorname{ch}(\mathscr X^{\mathrm{im}}_\alpha)
      \in\mathscr X^{\mathrm{im}}_\alpha,\\
    \chi^{\mathrm{im}}_\alpha\colon
      \operatorname{ST}_\pi(Y)\times_PP_\alpha
        \xrightarrow{\sim}Q_\alpha,
    \qquad
    \kappa_\alpha\circ\chi^{\mathrm{im}}_\alpha
      =\iota_{f,\pi,\alpha}.
  \end{gathered}                                                  \tag{ST11}
  \]
  Closed-subscheme extensionality on the open cover
  $P=\bigcup_\alpha P_\alpha$ is the equality
  \[
    [\,Y^{\mathrm{st}}_{\mathrm{chart}}
        \xhookrightarrow{\kappa_{\mathrm{chart}}}P\,]
      =[\,\operatorname{ST}_\pi(Y)
        \xhookrightarrow{\iota_{f,\pi}}P\,]
      \quad\text{in }\operatorname{ClosedSub}(P).
  \]
  In particular, $\kappa_{\mathrm{chart}}$ is a closed immersion.
\end{proof}

\begin{definition}[Comparison from the strict transform to the extended-ideal blow-up]
  \label{def:nagata-strict-transform-blowup-comparison}
  \uses{def:nagata-strict-transform-closed-chart-model,
    def:nagata-strict-transform-chart-comparison,
    def:nagata-global-blowup,
    prop:nagata-strict-transform-closed-fiber-product,
    prop:nagata-global-blowup-affine-restrict,
    lemma:nagata-strict-transform-chart-localization,
    lemma:nagata-strict-transform-chart-cocycle,
    lemma:nagata-morphism-property-calculus}
  \lean{}

  Put $\mathcal J:=\mathcal I\mathcal O_Y$.  For
  $\alpha=(r,\ell,k)\in\Lambda$ in the chart notation of
  \Cref{def:nagata-strict-transform-closed-chart-model}, define
  \[
  \begin{gathered}
    \widehat Q_\alpha=\widehat Q_{r\ell k}:=
      \operatorname{Spec}
       \bigl(B_{r\ell}[I_rB_{r\ell}/\varphi(a_{rk})]\bigr),\\
    \Theta_\alpha=\Theta_{r\ell k}:=
      \operatorname{Spec}(\overline\sigma_{a_{rk}}^{-1})
      \colon Q_{r\ell k}\longrightarrow\widehat Q_{r\ell k}.
  \end{gathered}                                                   \tag{ST12}
  \]
  Let
  \[
  \begin{gathered}
    \jmath_\alpha^{\mathrm{st}}\colon
      Q_\alpha\hookrightarrow\operatorname{ST}_\pi(Y),\qquad
    \widehat\jmath_\alpha\colon
      \widehat Q_\alpha\hookrightarrow
        \operatorname{Bl}_{\mathcal J}(Y),\\
    \operatorname{ST}_\pi(Y)
      =\bigcup_{\alpha\in\Lambda}\jmath_\alpha^{\mathrm{st}}(Q_\alpha),
    \qquad
    \operatorname{Bl}_{\mathcal J}(Y)
      =\bigcup_{\alpha\in\Lambda}\widehat\jmath_\alpha(\widehat Q_\alpha)
  \end{gathered}                                                   \tag{ST13}
  \]
  be the chart open immersions supplied respectively by
  \Cref{prop:nagata-strict-transform-closed-fiber-product} and
  \Cref{prop:nagata-global-blowup-affine-restrict}.  For
  $\alpha,\beta\in\Lambda$, define
  \[
    R_{\alpha\beta}:=Q_\alpha\times_{\operatorname{ST}_\pi(Y)}Q_\beta,
    \qquad
    u_{\alpha\beta}^\alpha\colon R_{\alpha\beta}\to Q_\alpha,
    \qquad
    u_{\alpha\beta}^\beta\colon R_{\alpha\beta}\to Q_\beta.
  \]
  The localization equality of
  \Cref{lemma:nagata-strict-transform-chart-localization} and the generator
  cocycle equality of \Cref{lemma:nagata-strict-transform-chart-cocycle},
  on a principal affine cover of $R_{\alpha\beta}$, are the equality
  \[
    \widehat\jmath_\alpha\circ\Theta_\alpha\circ
      u_{\alpha\beta}^\alpha
      =\widehat\jmath_\beta\circ\Theta_\beta\circ
        u_{\alpha\beta}^\beta.                                  \tag{ST13a}
  \]
  By (PC11), the family of equalities (ST13a) is the membership statement
  \[
    \mathbf\Theta:=
      (\widehat\jmath_\alpha\circ\Theta_\alpha)_{\alpha\in\Lambda}
      \in\mathscr G_{\operatorname{ST}_\pi(Y)}
        \bigl(\operatorname{Bl}_{\mathcal J}(Y),(Q_\alpha)_\alpha\bigr).
                                                                    \tag{ST14}
  \]
  Define, using the inverse of (PC12),
  \[
  \begin{gathered}
    \Theta_{f,\mathcal I}:=
      \operatorname{res}_{\operatorname{ST}_\pi(Y),
        \operatorname{Bl}_{\mathcal J}(Y)}^{-1}(\mathbf\Theta)
      \colon\operatorname{ST}_\pi(Y)\longrightarrow
        \operatorname{Bl}_{\mathcal J}(Y),\\
    \Theta_{f,\mathcal I}\circ\jmath_\alpha^{\mathrm{st}}
      =\widehat\jmath_\alpha\circ\Theta_\alpha
      \qquad(\alpha\in\Lambda).                                
  \end{gathered} \tag{ST15}
  \]
\end{definition}

\begin{theorem}[{\cite[\href{https://stacks.math.columbia.edu/tag/080E}{Tag 080E}]{stacks-project}}]
  \label{thm:nagata-strict-transform-is-blowup}
  \uses{def:nagata-strict-transform-blowup-comparison}
  \lean{}

  The constructed morphism $\Theta_{f,\mathcal I}$ is an isomorphism over
  $Y$ and identifies the strict transform with the blow-up of $Y$ in the
  extended ideal $\mathcal I\mathcal O_Y$.
\end{theorem}

\begin{proof}
  \uses{prop:nagata-strict-transform-chart-iso,
    prop:nagata-strict-transform-closed-fiber-product,
    prop:nagata-global-blowup-affine-restrict,
    lemma:nagata-strict-transform-chart-localization,
    lemma:nagata-strict-transform-chart-cocycle,
    lemma:nagata-morphism-property-calculus}
  For each $(r,\ell,k)$,
  \Cref{prop:nagata-strict-transform-chart-iso} supplies inverse ring maps
  \[
    \overline\sigma_{a_{rk}}\colon
      C^{\mathrm{st}}_{r\ell k}
        \xrightarrow{\sim}
      B_{r\ell}[I_rB_{r\ell}/\varphi(a_{rk})],
    \qquad
    \overline\sigma_{a_{rk}}^{-1}.
  \]
  Retain $\widehat Q_{r\ell k}$, $\Theta_{r\ell k}$,
  $\jmath_{r\ell k}^{\mathrm{st}}$, and
  $\widehat\jmath_{r\ell k}$ from (ST12)--(ST15), and define
  \[
    \Xi_{r\ell k}
      :=\operatorname{Spec}(\overline\sigma_{a_{rk}})
      \colon\widehat Q_{r\ell k}\longrightarrow Q_{r\ell k}
  \]
  and record its inverse identities
  \[
    \Xi_{r\ell k}\circ\Theta_{r\ell k}
      =\operatorname{id}_{Q_{r\ell k}},
    \qquad
    \Theta_{r\ell k}\circ\Xi_{r\ell k}
      =\operatorname{id}_{\widehat Q_{r\ell k}}.
  \]
  Define
  \[
    \widehat\Xi_{r\ell k}:=
      \jmath_{r\ell k}^{\mathrm{st}}\circ\Xi_{r\ell k}
      \colon\widehat Q_{r\ell k}\to\operatorname{ST}_\pi(Y).
  \]
  For $\alpha=(r,\ell,k)$, set
  \[
    \widehat\Xi_\alpha:=\widehat\Xi_{r\ell k}.
  \]
  For $\alpha,\beta\in\Lambda$, define
  \[
    \widehat R_{\alpha\beta}:=
      \widehat Q_\alpha\times_{\operatorname{Bl}_{\mathcal J}(Y)}
        \widehat Q_\beta,
    \qquad
    \widehat u_{\alpha\beta}^\alpha\colon
      \widehat R_{\alpha\beta}\to\widehat Q_\alpha,
    \qquad
    \widehat u_{\alpha\beta}^\beta\colon
      \widehat R_{\alpha\beta}\to\widehat Q_\beta.
  \]
  The localization and cocycle equalities used in (ST13a), applied to the
  inverse ring maps, are
  \[
    \widehat\Xi_\alpha\circ\widehat u_{\alpha\beta}^\alpha
      =\widehat\Xi_\beta\circ\widehat u_{\alpha\beta}^\beta.
                                                                    \tag{ST16}
  \]
  By (PC11), (ST16) is the membership equality
  \[
    \mathbf\Xi:=(\widehat\Xi_{r\ell k})_{r,\ell,k}
      \in\mathscr G_{\operatorname{Bl}_{\mathcal J}(Y)}
        \bigl(\operatorname{ST}_\pi(Y),
          (\widehat Q_{r\ell k})_{r,\ell,k}\bigr)
  \]
  in the notation (PC11).  Define, using the inverse of (PC12),
  \[
  \begin{gathered}
    \Xi_{f,\mathcal I}:=
      \operatorname{res}_{\operatorname{Bl}_{\mathcal J}(Y),
        \operatorname{ST}_\pi(Y)}^{-1}(\mathbf\Xi)
      \colon\operatorname{Bl}_{\mathcal J}(Y)\to
        \operatorname{ST}_\pi(Y),\\
    \Xi_{f,\mathcal I}\circ\widehat\jmath_{r\ell k}
      =\widehat\Xi_{r\ell k}
      =\jmath_{r\ell k}^{\mathrm{st}}\circ\Xi_{r\ell k}.
  \end{gathered}
  \]
  Equations (ST15) and the preceding component equality are the two chains
  \[
  \begin{aligned}
    (\Xi_{f,\mathcal I}\circ\Theta_{f,\mathcal I})
      \circ\jmath_{r\ell k}^{\mathrm{st}}
      &=\Xi_{f,\mathcal I}\circ\widehat\jmath_{r\ell k}
          \circ\Theta_{r\ell k}\\
      &=\jmath_{r\ell k}^{\mathrm{st}}\circ\Xi_{r\ell k}
          \circ\Theta_{r\ell k}
       =\jmath_{r\ell k}^{\mathrm{st}},\\
    (\Theta_{f,\mathcal I}\circ\Xi_{f,\mathcal I})
      \circ\widehat\jmath_{r\ell k}
      &=\Theta_{f,\mathcal I}\circ\jmath_{r\ell k}^{\mathrm{st}}
          \circ\Xi_{r\ell k}\\
      &=\widehat\jmath_{r\ell k}\circ\Theta_{r\ell k}
          \circ\Xi_{r\ell k}
       =\widehat\jmath_{r\ell k}.
  \end{aligned}
  \]
  Morphism extensionality on these open covers is the pair of equalities
  \[
    \Xi_{f,\mathcal I}\circ\Theta_{f,\mathcal I}
      =\operatorname{id}_{\operatorname{ST}_\pi(Y)},
    \qquad
    \Theta_{f,\mathcal I}\circ\Xi_{f,\mathcal I}
      =\operatorname{id}_{\operatorname{Bl}_{\mathcal J}(Y)}.
  \]
  On coordinate rings both structural maps to $Y$ restrict to
  $B_{r\ell}\to C^{\mathrm{st}}_{r\ell k}$, $y\mapsto[y\otimes1]$.
  The chartwise structural equation is
  \[
    \overline\sigma_{a_{rk}}([y\otimes1])=y.
  \]
  Morphism extensionality on the chart cover is the equality
  \[
    b_{Y,\mathcal J}\circ\Theta_{f,\mathcal I}
      = f^{\mathrm{st}}_\pi,
  \]
\end{proof}


\begin{lemma}[Scheme-theoretic closure characterization]
  \label{lemma:nagata-strict-transform-closure-characterization}
  \lean{}

  Retain the notation
  \[
    P=Y\times_XX',\qquad Y_U=Y\times_XU,\qquad
    \lambda_{f,\pi}\colon Y_U\to P,\qquad
    \iota_{f,\pi}\colon\operatorname{ST}_\pi(Y)\hookrightarrow P
  \]
  from \Cref{def:nagata-strict-transform}.  For every closed immersion
  $h\colon W\hookrightarrow P$ and every morphism $r\colon Y_U\to W$ with
  \[
    h\circ r=\lambda_{f,\pi},
  \]
  there is a unique closed immersion
  \[
    c_{h,r}\colon\operatorname{ST}_\pi(Y)\hookrightarrow W
  \]
  satisfying
  \[
    h\circ c_{h,r}=\iota_{f,\pi},
    \qquad
    c_{h,r}\circ\lambda_{f,\pi}^{\mathrm{im}}=r,
  \]
  where $\lambda_{f,\pi}^{\mathrm{im}}$ and its factorization equality are
  (ST0f) of \Cref{def:nagata-strict-transform}.  This universal property is the
  meaning of ``the scheme-theoretic closure of $Y_U$ in $P$''.
\end{lemma}

\begin{proof}
  \uses{def:nagata-closed-subscheme-class,
    lemma:nagata-power-torsion-quotient-torsionfree,
    prop:nagata-strict-transform-closed-fiber-product}
  The existence and uniqueness of $c_{h,r}$ are the universal property
  \Cref{lemma:nagata-scheme-image-universal} applied to
  \[
    \lambda_{f,\pi}=h\circ r
  \]
  and to the class equality
  \[
    [\,\operatorname{ST}_\pi(Y)
      \xhookrightarrow{\iota_{f,\pi}}P\,]
      =[\,\operatorname{SIm}_P(\lambda_{f,\pi})
      \xhookrightarrow{i_{\lambda_{f,\pi}}}P\,]
      \quad\text{in }\operatorname{ClosedSub}(P),
  \]
  together with the exact defining identities
  \[
    \operatorname{ST}_\pi(Y)
      :=\operatorname{SIm}_P(\lambda_{f,\pi}),
    \qquad
    \iota_{f,\pi}:=i_{\lambda_{f,\pi}}
  \]
  of \Cref{def:nagata-strict-transform}.  For an affine verification, let
  $P_a=\operatorname{Spec}(C)$ be a chart, let $a\in C$ be its distinguished
  exceptional element, and define
  \[
    C[a^{-1}]:=\{a^n\mid n\in\mathbf N\}^{-1}C,
    \qquad
    \ell_{C,a}\colon C\to C[a^{-1}],\quad c\mapsto c/1.
  \]
  Let $W\times_PP_a=\operatorname{Spec}(C/K)$.  The equation
  $h\circ r=\lambda_{f,\pi}$ restricts to the factorization of the typed
  localization morphism
  \[
    \operatorname{Spec}(\ell_{C,a})\colon
      \operatorname{Spec}(C[a^{-1}])\longrightarrow\operatorname{Spec}(C)
  \]
  through $\operatorname{Spec}(C/K)$.  The corresponding ring maps satisfy
  \[
  \begin{aligned}
    K&\subseteq\ker(\ell_{C,a})\\
     &=\Gamma_{(a)}(C)
       &&\text{by \Cref{lemma:nagata-power-torsion-kernel-localization}}.
  \end{aligned}
  \]
  Define
  \[
    C^{\mathrm{st}}_{(a)}:=C/\Gamma_{(a)}(C).
  \]
  The inclusion $K\subseteq\Gamma_{(a)}(C)$ defines the quotient map
  \[
    C/K\twoheadrightarrow C/\Gamma_{(a)}(C)
      =C^{\mathrm{st}}_{(a)}
  \]
  and the explicitly named closed immersion
  \[
    c_{h,r,a}:=
      \operatorname{Spec}
        \bigl(C/K\twoheadrightarrow C/\Gamma_{(a)}(C)\bigr)
      \colon
      \operatorname{Spec}(C^{\mathrm{st}}_{(a)})\hookrightarrow
      \operatorname{Spec}(C/K).
  \]
  Restriction of the two defining equalities for $c_{h,r}$ to $P_a$ and
  uniqueness in \Cref{lemma:nagata-scheme-image-universal} are the exact
  equality
  \[
    c_{h,r}|_{P_a}=c_{h,r,a}.                                    \tag{SC1}
  \]
  Finally, \Cref{lemma:nagata-power-torsion-quotient-torsionfree} is the
  saturation equality
  \[
    \Gamma_{(a)}(C/\Gamma_{(a)}(C))=0,
  \]
\end{proof}

\begin{lemma}[Flatness kills the strict-transform torsion]
  \label{lemma:nagata-flat-strict-transform-no-torsion}
  \uses{def:nagata-power-torsion}
  \lean{}

  If $B$ is flat over $A$, then the $a$-power-torsion in
  $B\otimes_AA[I/a]$ is zero.
\end{lemma}

\begin{proof}
  \uses{lemma:nagata-chart-generator-regular}
  Put $R=A[I/a]$ and $C=B\otimes_AR$.  By
  \Cref{lemma:nagata-chart-generator-regular}, the $A$-linear map
  \[
    m_a\colon R\to R,\qquad r\mapsto ar
  \]
  is injective.  Flatness of the $A$-module $B$ is the exactness identity
  \[
    \ker(\operatorname{id}_B\otimes m_a)
      =B\otimes_A\ker(m_a)=B\otimes_A0=0.
  \]
  Define the explicitly typed map
  \[
    m_{a,C}:=\operatorname{id}_B\otimes m_a\colon
      C=B\otimes_AR\longrightarrow B\otimes_AR=C.
  \]
  The preceding kernel equality is
  \[
    \ker(m_{a,C})=0,
    \qquad
    m_{a,C}(c)=0\Longrightarrow c=0\quad(c\in C).               \tag{FT0}
  \]
  For $y\otimes r\in C$,
  \[
    m_{a,C}(y\otimes r)
      =y\otimes ar
      =a_C(y\otimes r),
  \]
  Pure-tensor extensionality applied to the preceding equality is
  \[
    m_{a,C}=(c\mapsto a_Cc),\qquad
    a_Cc=0\Longrightarrow m_{a,C}(c)=0\Longrightarrow c=0       \tag{FT1}
  \]
  by (FT0).  For $n\geq1$, iteration of (FT1) is the chain
  \[
    a_C^nc=0
      \Longrightarrow a_C^{n-1}c=0
      \Longrightarrow\cdots\Longrightarrow a_Cc=0
      \Longrightarrow c=0
      \qquad(c\in C,n\in\mathbf N).
  \]
  By \Cref{def:nagata-power-torsion}, the iteration chain is the set equality
  \[
    \Gamma_{(a_C)}(C)
      =\{c\in C\mid\exists n\in\mathbf N,\ a_C^nc=0\}
      =\{0\}.
  \]
\end{proof}

\begin{corollary}[Flat strict transform equals ordinary base change]
  \label{cor:nagata-flat-strict-transform-base-change}
  \lean{}

  If $f\colon Y\to X$ is flat, if
  $\mathcal I\subseteq\mathcal O_X$ is a finite-type quasi-coherent ideal,
  and if
  $\pi=b_{X,\mathcal I}\colon X'=\operatorname{Bl}_{\mathcal I}(X)\to X$,
  then the closed immersion of \Cref{def:nagata-strict-transform}
  \[
    \iota_{f,\pi}\colon\operatorname{ST}_\pi(Y)\hookrightarrow Y\times_XX'
  \]
  is an isomorphism.
\end{corollary}

\begin{proof}
  \uses{lemma:nagata-flat-strict-transform-no-torsion,
    prop:nagata-strict-transform-closed-fiber-product}
  On every affine chart of
  \Cref{def:nagata-strict-transform-closed-chart-model}, put
  \[
    C_a^{\mathrm{bc}}:=C_a(A,B,I).
  \]
  Flatness of $f$ makes $B$ flat over $A$; the instantiated equality in
  \Cref{lemma:nagata-flat-strict-transform-no-torsion} is
  \[
    \Gamma_{(a_C)}(C_a^{\mathrm{bc}})=0.
  \]
  Substitution of the zero-torsion equality into
  \Cref{def:nagata-strict-transform-chart-ring} is
  \[
    C_a^{\mathrm{st}}(A,B,I)
      =C_a^{\mathrm{bc}}/0=C_a^{\mathrm{bc}},
  \]
  and the chart restriction of $\iota_{f,\pi}$ is
  \[
    \operatorname{Spec}(\operatorname{id}_{C_a^{\mathrm{bc}}})
      =\operatorname{id}_{\operatorname{Spec}(C_a^{\mathrm{bc}})}.
  \]
  The chart equations and their covering equation are
  \[
    Y\times_XX'=\bigcup_a\operatorname{Spec}(C_a^{\mathrm{bc}}),\qquad
    \iota_{f,\pi}|_{\operatorname{Spec}(C_a^{\mathrm{bc}})}
      =\operatorname{id}_{\operatorname{Spec}(C_a^{\mathrm{bc}})}.
  \]
  The local isomorphism criterion, applied to this cover and to
  \Cref{prop:nagata-strict-transform-closed-fiber-product}, is
  \[
    \operatorname{IsIso}(\iota_{f,\pi}).
  \]
\end{proof}

\begin{lemma}[Strict transform on the preserved open]
  \label{lemma:nagata-strict-transform-over-complement}
  \lean{}

  Let $f\colon Y\to X$, let
  $\mathcal I\subseteq\mathcal O_X$ be a finite-type quasi-coherent ideal,
  and let
  $\pi=b_{X,\mathcal I}\colon X'=\operatorname{Bl}_{\mathcal I}(X)\to X$.
  Put
  \[
  \begin{gathered}
    U:=X\setminus V(\mathcal I),\qquad
    j_U\colon U\hookrightarrow X,\qquad
    X'_U:=X'\times_XU,\qquad
    Y_U:=Y\times_XU,\\
    P:=Y\times_XX',\qquad
    P_U:=P\times_{X'}X'_U,\qquad
    \operatorname{ST}_\pi(Y)_U
      :=\operatorname{ST}_\pi(Y)\times_{X'}X'_U.
  \end{gathered}
  \]
  Let
  \[
  \begin{gathered}
    q_P\colon P_U\to P,\qquad q_{X'}\colon P_U\to X'_U,\\
    \iota_{f,\pi,U}\colon\operatorname{ST}_\pi(Y)_U\hookrightarrow P_U
  \end{gathered}
  \]
  be the projections and the base change of $\iota_{f,\pi}$.
  There is a unique isomorphism
  \[
    \alpha^{\mathrm{st}}_U\colon
      \operatorname{ST}_\pi(Y)_U\xrightarrow{\sim}Y_U
  \]
  satisfying
  \[
    \operatorname{pr}_Y\circ\alpha^{\mathrm{st}}_U
      =\operatorname{pr}_Y\circ q_P\circ\iota_{f,\pi,U},
    \qquad
    \operatorname{pr}_U\circ\alpha^{\mathrm{st}}_U
      =\alpha_U\circ q_{X'}\circ\iota_{f,\pi,U},
  \]
  where $\alpha_U\colon X'_U\xrightarrow{\sim}U$ is the isomorphism of
  \Cref{thm:nagata-global-blowup-complement}.
\end{lemma}

\begin{proof}
  \uses{def:nagata-choice-notation,def:nagata-closed-subscheme-class,
    lemma:nagata-immersion-qcqs-criteria,
    thm:nagata-global-blowup-complement,
    lemma:nagata-strict-transform-closure-characterization}
  By \Cref{thm:nagata-global-blowup-complement},
  $\alpha_U\colon X'_U\xrightarrow{\sim}U$.  Define
  \[
    a_U:=
      \bigl(\operatorname{pr}_Y\circ q_P,
        \alpha_U\circ q_{X'}\bigr)
      \colon P_U\xrightarrow{\sim}Y_U.                         \tag{STOC1}
  \]
  Its inverse is the fibre-product morphism determined by
  \[
  \begin{gathered}
    q_P\circ a_U^{-1}=\lambda_{f,\pi},\\
    q_{X'}\circ a_U^{-1}
      =\alpha_U^{-1}\circ\operatorname{pr}_U,
    \qquad
    a_U\circ a_U^{-1}=\operatorname{id}_{Y_U},
    \qquad
    a_U^{-1}\circ a_U=\operatorname{id}_{P_U}.
  \end{gathered}                                                    \tag{STOC2}
  \]
  Let
  \[
    c_U\colon
      \operatorname{ST}_\pi(Y)_U
        \xrightarrow{\sim}
      \operatorname{ST}_\pi(Y)\times_PP_U
  \]
  be the fibre-product associativity isomorphism.  Define
  \[
    \iota_{f,\pi,U}^{P}:=\operatorname{pr}_{P_U}\colon
      \operatorname{ST}_\pi(Y)\times_PP_U\hookrightarrow P_U.
  \]
  The inverse identities in (STOC2) are
  \[
    a_U^{-1}\circ a_U=\operatorname{id}_{P_U},\qquad
    a_U\circ a_U^{-1}=\operatorname{id}_{Y_U},\qquad
    \operatorname{IsIso}(a_U^{-1}).                              \tag{STOC2a}
  \]
  By \Cref{lemma:nagata-immersion-qcqs-criteria},
  \[
    \operatorname{IsIso}(a_U^{-1})
      \Longrightarrow\operatorname{IsQCQS}(a_U^{-1}).           \tag{STOC2b}
  \]
  Substitution of (STOC2) into
  \Cref{lemma:nagata-scheme-image-restrict-open} is the class chain
  \[
  \begin{aligned}
    [\,\operatorname{ST}_\pi(Y)\times_PP_U
      \xhookrightarrow{\iota_{f,\pi,U}^{P}}P_U\,]
      &=[\,\operatorname{SIm}_{P_U}(a_U^{-1})
        \xhookrightarrow{i_{a_U^{-1}}}P_U\,]
        &&\text{by \Cref{lemma:nagata-scheme-image-restrict-open}}\\
      &=[\,P_U\xhookrightarrow{\operatorname{id}_{P_U}}P_U\,]
        \quad\text{in }\operatorname{ClosedSub}(P_U)
        &&\text{by \Cref{def:nagata-scheme-theoretic-image} and (STOC2a)}.
  \end{aligned}                                                  \tag{STOC3}
  \]
  Define
  \[
    \mathscr B_U:=
      \left\{\beta\in
        \operatorname{Hom}(\operatorname{ST}_\pi(Y)\times_PP_U,P_U)
        \ \middle|\
        \operatorname{IsIso}(\beta)\ \wedge\
        \operatorname{id}_{P_U}\circ\beta
          =\iota_{f,\pi,U}^{P}\right\}.
  \]
  By (CSub3) of \Cref{def:nagata-closed-subscheme-class}, (STOC3) is
  $\mathscr B_U\ne\varnothing$.  Define
  \[
    \beta_U:=\operatorname{ch}(\mathscr B_U)\in\mathscr B_U,
    \qquad
    \operatorname{IsIso}(\beta_U),
    \qquad
    \beta_U=\iota_{f,\pi,U}^{P}.                                \tag{STOC3a}
  \]
  Compose the isomorphism in (STOC3a) with $c_U$ and define
  \[
    \iota_{f,\pi,U}:=
      \iota_{f,\pi,U}^{P}\circ c_U
      \colon\operatorname{ST}_\pi(Y)_U\xrightarrow{\sim}P_U.   \tag{STOC3b}
  \]
  Define
  \[
    \alpha_U^{\mathrm{st}}
      :=a_U\circ\iota_{f,\pi,U}
      \colon\operatorname{ST}_\pi(Y)_U
        \xrightarrow{\sim}Y_U.                                 \tag{STOC4}
  \]
  Substitution of (STOC1) into (STOC4) is the component pair
  \[
  \begin{aligned}
    \operatorname{pr}_Y\circ\alpha_U^{\mathrm{st}}
      &=\operatorname{pr}_Y\circ
        q_P\circ\iota_{f,\pi,U},\\
    \operatorname{pr}_U\circ\alpha_U^{\mathrm{st}}
      &=\alpha_U\circ q_{X'}\circ\iota_{f,\pi,U}.
  \end{aligned}
  \]
\end{proof}

\begin{lemma}[Strict transform under a dominating blow-up]
  \label{lemma:nagata-strict-transform-dominating}
  \lean{}

  Let $f\colon Y\to X$, let $\pi_1\colon X_1\to X$ and
  $\rho\colon X_2\to X_1$ be blow-ups, and put
  $\pi_2:=\pi_1\circ\rho$.  Let $j\colon U\hookrightarrow X$ be a
  quasi-compact open immersion, and set
  \[
    X_{i,U}:=X_i\times_XU,\qquad
    j_i^U:=\operatorname{pr}_{X_i}\colon X_{i,U}\hookrightarrow X_i,\qquad
    \rho_U:=(\rho\circ\operatorname{pr}_{X_2},\operatorname{pr}_U)
      \colon X_{2,U}\to X_{1,U},
  \]
  and fix compatible isomorphisms
  \[
    \alpha_i\colon X_{i,U}\xrightarrow{\sim}U,\qquad
    j\circ\alpha_i=\pi_i\circ\operatorname{pr}_{X_i},\qquad
    \alpha_1\circ\rho_U=\alpha_2.                                 \tag{D1}
  \]
  Put
  \[
    X_{2,1,U}:=X_2\times_{X_1}X_{1,U},
  \]
  and define
  \[
    c_U:=(\operatorname{pr}_{X_2},
      \operatorname{pr}_U\circ\operatorname{pr}_{X_{1,U}})
      \colon X_{2,1,U}\xrightarrow{\sim}X_{2,U},
    \qquad
    \alpha_\rho:=\rho_U\circ c_U
      \colon X_{2,1,U}\xrightarrow{\sim}X_{1,U}.                 \tag{D2}
  \]
  Its structure equation is
  \[
    j_1^U\circ\alpha_\rho
      =\rho\circ\operatorname{pr}_{X_2}\colon X_{2,1,U}\to X_1.
  \]
  The type and structure equation in (D2) identify
  $(\alpha_\rho,j_1^U\circ\alpha_\rho
    =\rho\circ\operatorname{pr}_{X_2})$
  as the preserved-open datum of
  \Cref{def:nagata-strict-transform}.  Define the ambient associativity
  isomorphism
  \[
    a_{Y,\pi_1,\rho}:=
      \bigl(
        \operatorname{pr}_Y\circ\operatorname{pr}_{Y\times_XX_1},
        \operatorname{pr}_{X_2}
      \bigr)
      \colon (Y\times_XX_1)\times_{X_1}X_2
        \xrightarrow{\sim}Y\times_XX_2.                           \tag{D3}
  \]
  Its inverse is
  \[
    a_{Y,\pi_1,\rho}^{-1}
      =\bigl(
          (\operatorname{pr}_Y,\rho\circ\operatorname{pr}_{X_2}),
          \operatorname{pr}_{X_2}
        \bigr).
  \]
  There is a unique isomorphism
  \[
    \delta_{\pi_1,\rho,Y}\colon
      \operatorname{ST}_{\rho}(\operatorname{ST}_{\pi_1}(Y))
      \xrightarrow{\sim}\operatorname{ST}_{\pi_2}(Y)             \tag{D4}
  \]
  whose composite with $\iota_{f,\pi_2}$ equals
  \[
    a_{Y,\pi_1,\rho}\circ
      (\iota_{f,\pi_1}\times_{X_1}\operatorname{id}_{X_2})\circ
      \iota_{f^{\mathrm{st}}_{\pi_1},\rho}.                       \tag{D5}
  \]
\end{lemma}

\begin{proof}
  \uses{def:nagata-choice-notation,def:nagata-closed-subscheme-class,
    def:nagata-strict-transform,
    lemma:nagata-immersion-qcqs-criteria,
    lemma:nagata-morphism-property-calculus,
    lemma:nagata-strict-transform-over-complement,
    lemma:nagata-strict-transform-closure-characterization}
  Put
  \[
    T_1:=\operatorname{ST}_{\pi_1}(Y),\qquad
    f_1:=f^{\mathrm{st}}_{\pi_1}\colon T_1\to X_1,\qquad
    T_{1,U}:=T_1\times_{X_1}X_{1,U},\qquad
    Y_U:=Y\times_XU.
  \]
  By \Cref{lemma:nagata-strict-transform-over-complement}, define
  \[
    \theta_U:=\alpha_U^{\mathrm{st}}\colon
      T_{1,U}\xrightarrow{\sim}Y_U.                               \tag{D6}
  \]
  Let
  \[
    \lambda_{f_1,\rho}\colon T_{1,U}\longrightarrow T_1\times_{X_1}X_2
  \]
  be the morphism of \Cref{def:nagata-strict-transform} formed with
  $\alpha_\rho$ from (D2), and define
  \[
    m:=a_{Y,\pi_1,\rho}\circ
      (\iota_{f,\pi_1}\times_{X_1}\operatorname{id}_{X_2})
      \colon T_1\times_{X_1}X_2\longrightarrow Y\times_XX_2.      \tag{D7}
  \]
  Let
  \[
    f_{1,U}:=\operatorname{pr}_{X_{1,U}}\colon T_{1,U}\to X_{1,U},
    \qquad
    f_U:=\operatorname{pr}_U\colon Y_U\to U.
  \]
  The second projection equation in
  \Cref{lemma:nagata-strict-transform-over-complement}, instantiated at
  $(f,\pi_1)$, is
  \[
    f_U\circ\theta_U=\alpha_1\circ f_{1,U}.                        \tag{D7a}
  \]
  Expanding the definitions of $\lambda_{f_1,\rho}$,
  $\lambda_{f,\pi_2}$, and $m$ is the first-projection chain
  \[
  \begin{aligned}
    \operatorname{pr}_Y\circ m\circ\lambda_{f_1,\rho}
      &=\operatorname{pr}_Y\circ\theta_U
       =\operatorname{pr}_Y\circ\lambda_{f,\pi_2}\circ\theta_U.
  \end{aligned}
  \]
  From (D1)--(D2) one has
  \[
    c_U\circ\alpha_\rho^{-1}
      =\alpha_2^{-1}\circ\alpha_1\colon X_{1,U}\to X_{2,U}.        \tag{D7b}
  \]
  Substitution of (D7a)--(D7b) is the second-projection chain
  \[
  \begin{aligned}
    \operatorname{pr}_{X_2}\circ m\circ\lambda_{f_1,\rho}
      &=\operatorname{pr}_{X_2}\circ c_U\circ\alpha_\rho^{-1}
          \circ f_{1,U}\\
      &=\operatorname{pr}_{X_2}\circ\alpha_2^{-1}\circ\alpha_1
          \circ f_{1,U} &&\text{by (D7b)}\\
      &=\operatorname{pr}_{X_2}\circ\alpha_2^{-1}\circ f_U\circ\theta_U
          &&\text{by (D7a)}\\
      &=\operatorname{pr}_{X_2}\circ\lambda_{f,\pi_2}\circ\theta_U.
  \end{aligned}
  \]
  The fibre-product comparison map
  \[
    \operatorname{Hom}(T_{1,U},Y\times_XX_2)
      \longrightarrow
      \operatorname{Hom}(T_{1,U},Y)\times
      \operatorname{Hom}(T_{1,U},X_2),
    \qquad h\longmapsto
      (\operatorname{pr}_Y\circ h,\operatorname{pr}_{X_2}\circ h),
  \]
  is injective.  Applied to the two preceding projection equalities, its
  conclusion is
  \[
    m\circ\lambda_{f_1,\rho}
      =\lambda_{f,\pi_2}\circ\theta_U.                            \tag{D8}
  \]
  Let
  \[
    E:=m\circ\iota_{f_1,\rho}\colon
      \operatorname{ST}_{\rho}(T_1)\hookrightarrow Y\times_XX_2. \tag{D9}
  \]
  The definition of $m$ and the stability of closed immersions under base
  change and composition are the predicate chain
  \[
  \begin{gathered}
    \operatorname{IsClosedImmersion}(\iota_{f_1,\rho}),\\
    \operatorname{IsClosedImmersion}
      (\iota_{f,\pi_1}\times_{X_1}\operatorname{id}_{X_2}),
    \qquad \operatorname{IsIso}(a_{Y,\pi_1,\rho}),\\
    m=(\iota_{f,\pi_1}\times_{X_1}\operatorname{id}_{X_2})
        \circ a_{Y,\pi_1,\rho},\\
    E=m\circ\iota_{f_1,\rho},
    \qquad \operatorname{IsClosedImmersion}(E)
      \quad\text{by \Cref{lemma:nagata-morphism-property-calculus}}.
  \end{gathered}                                                  \tag{D9a}
  \]
  By \Cref{lemma:nagata-immersion-qcqs-criteria}, (PC3) of
  \Cref{lemma:nagata-morphism-property-calculus}, and the QCQS predicates in
  \Cref{def:nagata-strict-transform}, the predicates required below are the
  chain
  \[
  \begin{aligned}
    \operatorname{IsClosedImmersion}
      (\iota_{f,\pi_1}\times_{X_1}\operatorname{id}_{X_2})
      &\Longrightarrow
        \operatorname{IsQCQS}
          (\iota_{f,\pi_1}\times_{X_1}\operatorname{id}_{X_2}),\\
    \operatorname{IsIso}(a_{Y,\pi_1,\rho})
      &\Longrightarrow\operatorname{IsQCQS}(a_{Y,\pi_1,\rho}),\\
    \operatorname{IsQCQS}
      (\iota_{f,\pi_1}\times_{X_1}\operatorname{id}_{X_2})\wedge
      \operatorname{IsQCQS}(a_{Y,\pi_1,\rho})\wedge
      \bigl(m=(\iota_{f,\pi_1}\times_{X_1}\operatorname{id}_{X_2})
        \circ a_{Y,\pi_1,\rho}\bigr)
      &\Longrightarrow\operatorname{IsQCQS}(m),\\
    \operatorname{IsQCQS}(m)\wedge
      \operatorname{IsQCQS}(\lambda_{f_1,\rho})
      &\Longrightarrow
        \operatorname{IsQCQS}(m\circ\lambda_{f_1,\rho}),\\
    \operatorname{IsIso}(\theta_U)
      &\Longrightarrow\operatorname{IsQCQS}(\theta_U),\\
    \operatorname{IsQCQS}(\lambda_{f,\pi_2})\wedge
      \operatorname{IsQCQS}(\theta_U)
      &\Longrightarrow
        \operatorname{IsQCQS}(\lambda_{f,\pi_2}\circ\theta_U).
  \end{aligned}                                                     \tag{D9b}
  \]
  Equations
  (D8)--(D9), invariance of the scheme-theoretic image under the source
  isomorphism $\theta_U$, and \Cref{def:nagata-strict-transform} are the
  class chain
  \[
  \begin{aligned}
    [\,\operatorname{ST}_{\rho}(T_1)
      \xhookrightarrow{E}Y\times_XX_2\,]
      &=[\,\operatorname{SIm}_{Y\times_XX_2}
          (m\circ\lambda_{f_1,\rho})
        \xhookrightarrow{i_{m\circ\lambda_{f_1,\rho}}}
          Y\times_XX_2\,]\\
      &=[\,\operatorname{SIm}_{Y\times_XX_2}
          (\lambda_{f,\pi_2}\circ\theta_U)
        \xhookrightarrow{i_{\lambda_{f,\pi_2}\circ\theta_U}}
          Y\times_XX_2\,]\\
      &=[\,\operatorname{SIm}_{Y\times_XX_2}(\lambda_{f,\pi_2})
        \xhookrightarrow{i_{\lambda_{f,\pi_2}}}Y\times_XX_2\,]\\
      &=[\,\operatorname{ST}_{\pi_2}(Y)
        \xhookrightarrow{\iota_{f,\pi_2}}Y\times_XX_2\,]
      \quad\text{in }\operatorname{ClosedSub}(Y\times_XX_2).
  \end{aligned}                                                     \tag{D10}
  \]
  Define
  \[
    \mathscr D_{\pi_1,\rho,Y}:=
      \left\{\delta\in
        \operatorname{Hom}(\operatorname{ST}_{\rho}(T_1),
          \operatorname{ST}_{\pi_2}(Y))
        \ \middle|\
        \operatorname{IsIso}(\delta)\ \wedge\
        \iota_{f,\pi_2}\circ\delta=E\right\}.
  \]
  By (CSub3) of \Cref{def:nagata-closed-subscheme-class}, equality of the
  first and last classes in (D10) is
  \[
    \mathscr D_{\pi_1,\rho,Y}\ne\varnothing.
  \]
  Define
  \[
    \delta_{\pi_1,\rho,Y}:=
      \operatorname{ch}(\mathscr D_{\pi_1,\rho,Y})
      \in\mathscr D_{\pi_1,\rho,Y}.
  \]
  Its membership expands to (D4)--(D5).  For
  $\delta\in\mathscr D_{\pi_1,\rho,Y}$, the two membership equations and
  monicity of $\iota_{f,\pi_2}$ form the cancellation chain
  \[
    \iota_{f,\pi_2}\circ\delta=E
      =\iota_{f,\pi_2}\circ\delta_{\pi_1,\rho,Y}
      \Longrightarrow\delta=\delta_{\pi_1,\rho,Y},
  \]
\end{proof}

\begin{lemma}[Disjoint strict transforms stay disjoint]
  \label{lemma:nagata-strict-transform-disjoint-dominating}
  \lean{}

  With the notation of
  \Cref{lemma:nagata-strict-transform-dominating}, let
  $f_i\colon Y_i\to X$ and put
  $S_i^1:=\operatorname{ST}_{\pi_1}(Y_i)$.  If
  \[
    S_1^1\times_{X_1}S_2^1=\varnothing,
  \]
  then, for $S_i^2:=\operatorname{ST}_{\pi_2}(Y_i)$,
  \[
    S_1^2\times_{X_2}S_2^2=\varnothing.
  \]
\end{lemma}

\begin{proof}
  \uses{lemma:nagata-strict-transform-dominating}
  For $i=1,2$, define
  \[
  \begin{gathered}
    f_i^1:=(f_i)_{\pi_1}^{\mathrm{st}}\colon S_i^1\to X_1,\\
    \delta_i:=\delta_{\pi_1,\rho,Y_i}\colon
      \operatorname{ST}_{\rho}(S_i^1)\xrightarrow{\sim}S_i^2
      \quad\text{by \Cref{lemma:nagata-strict-transform-dominating}},\\
    \iota_i^\rho:=\iota_{f_i^1,\rho}\colon
      \operatorname{ST}_{\rho}(S_i^1)
        \hookrightarrow S_i^1\times_{X_1}X_2,\\
    e_i:=\iota_i^\rho\circ\delta_i^{-1}\colon
      S_i^2\hookrightarrow S_i^1\times_{X_1}X_2.
  \end{gathered}
  \]
  Define their fibre product by
  \[
    e_{12}:=e_1\times_{X_2}e_2\colon
      S_1^2\times_{X_2}S_2^2
      \hookrightarrow
      (S_1^1\times_{X_1}X_2)
        \times_{X_2}(S_2^1\times_{X_1}X_2).
  \]
  The target has the explicit associativity chain
  \[
  \begin{aligned}
    (S_1^1\times_{X_1}X_2)
      \times_{X_2}(S_2^1\times_{X_1}X_2)
      &\cong(S_1^1\times_{X_1}S_2^1)\times_{X_1}X_2\\
      &=\varnothing\times_{X_1}X_2=\varnothing.
  \end{aligned}
  \]
  Under the displayed identification of the target with $\varnothing$, the
  underlying continuous map of $e_{12}$ is
  \[
    |e_{12}|\colon
      |S_1^2\times_{X_2}S_2^2|
        \longrightarrow|\varnothing|=\varnothing.
  \]
  \[
  \begin{aligned}
    |e_{12}|&\in
      \operatorname{Hom}_{\mathbf{Set}}
        (|S_1^2\times_{X_2}S_2^2|,\varnothing)\\
    &\Longrightarrow
      \operatorname{Hom}_{\mathbf{Set}}
        (|S_1^2\times_{X_2}S_2^2|,\varnothing)\neq\varnothing\\
    &\Longrightarrow
      |S_1^2\times_{X_2}S_2^2|=\varnothing\\
    &\Longrightarrow
      S_1^2\times_{X_2}S_2^2=\varnothing.
  \end{aligned}
  \]
\end{proof}

\begin{lemma}[Composition of admissible blow-ups, {\cite[\href{https://stacks.math.columbia.edu/tag/080L}{Tag 080L}]{stacks-project}}]
  \label{lemma:nagata-admissible-blowup-comp}
  \uses{def:nagata-admissible-blowup}
  \lean{}

  Let $X$ be quasi-compact and quasi-separated, let
  $j\colon U\hookrightarrow X$ be a quasi-compact open immersion, and let
  \[
    \pi_1\colon X_1=\operatorname{Bl}_{\mathcal I_1}(X)\to X
  \]
  be $U$-admissible.  Put $U_1=X_1\times_XU$ and let
  \[
    \pi_2\colon X_2=\operatorname{Bl}_{\mathcal I_2}(X_1)\to X_1
  \]
  be $U_1$-admissible.  Define
  $\pi_{21}:=\pi_1\circ\pi_2\colon X_2\to X$.  If
  \[
    \alpha_1\colon U_1\xrightarrow{\sim}U,
    \qquad
    \alpha_2\colon X_2\times_{X_1}U_1\xrightarrow{\sim}U_1
  \]
  are the isomorphisms included in the two data of
  \Cref{def:nagata-admissible-blowup}, then the isomorphism included in the
  composite preserved-open isomorphism on the given source $X_2$ is
  \[
    \alpha_{21}:=\alpha_1\circ\alpha_2\circ\gamma
      \colon X_2\times_XU\xrightarrow{\sim}U,
  \]
  where
  $\gamma\colon X_2\times_XU\xrightarrow{\sim}
    X_2\times_{X_1}U_1$ is the fibre-product associativity isomorphism.
  There exist a finite-type quasi-coherent ideal
  $\mathcal K\subseteq\mathcal O_X$ and an isomorphism over $X$
  \[
    \omega\colon X_2\xrightarrow{\sim}\operatorname{Bl}_{\mathcal K}(X),
    \qquad
    b_{X,\mathcal K}\circ\omega=\pi_{21}.                      \tag{AB1}
  \]
  Define
  \[
    \omega_U:=\omega\times_X\operatorname{id}_U
      \colon X_2\times_XU\xrightarrow{\sim}
        \operatorname{Bl}_{\mathcal K}(X)\times_XU,
    \qquad
    \alpha_{21}^{\mathcal K}:=
      \alpha_{21}\circ\omega_U^{-1}.                          \tag{AB2}
  \]
  Then the transported tuple
  \[
    \mathsf D_{21}:=
      (\mathcal K,\operatorname{Bl}_{\mathcal K}(X),
        b_{X,\mathcal K},\alpha_{21}^{\mathcal K})             \tag{AB3}
  \]
  is a $U$-admissible blow-up in the literal sense of
  \Cref{def:nagata-admissible-blowup}; the composite morphism
  $\pi_{21}$ represents $\mathsf D_{21}$ through the specified isomorphism
  $\omega$ in (AB1).
\end{lemma}

\begin{proof}
  \uses{thm:nagata-composition-of-blowups}
  Let $Z_1=V(\mathcal I_1)\subseteq X$ and
  $Z_2=V(\mathcal I_2)\subseteq X_1$.  By
  \Cref{thm:nagata-composition-of-blowups}, there is a finite-type
  quasi-coherent ideal $\mathcal K\subseteq\mathcal O_X$ and an isomorphism
  over $X$
  \[
    \omega\colon X_2\xrightarrow{\sim}\operatorname{Bl}_{\mathcal K}(X)
  \]
  satisfying
  \[
    b_{X,\mathcal K}\circ\omega=\pi_1\circ\pi_2,
    \qquad
    |V(\mathcal K)|=|Z_1|\cup\pi_1(|Z_2|).
  \]
  Admissibility of $\pi_1$ and $\pi_2$ means
  \[
    |Z_1|\cap U=\varnothing,
    \qquad
    |Z_2|\cap U_1=\varnothing.
  \]
  Since $\pi_1^{-1}(U)=U_1$,
  \[
  \begin{aligned}
    |V(\mathcal K)|\cap U
      &= (|Z_1|\cap U)\cup(\pi_1(|Z_2|)\cap U)\\
      &=\varnothing\cup
        \pi_1(|Z_2|\cap\pi_1^{-1}(U))\\
      &=\varnothing\cup\pi_1(|Z_2|\cap U_1)=\varnothing.
  \end{aligned}
  \]
  The equality
  $|V(\mathcal K)|\cap U=\varnothing$ is equivalent to
  \[
    \mathcal K\mathcal O_U=\mathcal O_U.                        \tag{AB4}
  \]
  Define $\omega_U$ and $\alpha_{21}^{\mathcal K}$ by (AB2).  Expanding the
  definitions of $\gamma$, $\alpha_1$, and $\alpha_2$ is the chain
  \[
    j\circ\alpha_{21}
      =\pi_1\circ\operatorname{pr}_{X_1}\circ\alpha_2\circ\gamma
      =\pi_1\circ\pi_2\circ\operatorname{pr}_{X_2},
                                                                    \tag{AB5}
  \]
  Base change of (AB1) is
  \[
    \omega\circ\operatorname{pr}_{X_2}\circ\omega_U^{-1}
      =\operatorname{pr}_{\operatorname{Bl}_{\mathcal K}(X)}.
  \]
  Substitution of (AB1), (AB2), and (AB5) is
  \[
  \begin{aligned}
    j\circ\alpha_{21}^{\mathcal K}
      &=j\circ\alpha_{21}\circ\omega_U^{-1}\\
      &=\pi_{21}\circ\operatorname{pr}_{X_2}\circ\omega_U^{-1}\\
      &=b_{X,\mathcal K}\circ\omega\circ
          \operatorname{pr}_{X_2}\circ\omega_U^{-1}\\
      &=b_{X,\mathcal K}\circ
          \operatorname{pr}_{\operatorname{Bl}_{\mathcal K}(X)}.
  \end{aligned}                                                   \tag{AB6}
  \]
  Equations (AB1), (AB4), and (AB6) are exactly the defining conditions of
  \Cref{def:nagata-admissible-blowup} for the transported tuple
  $\mathsf D_{21}$ in (AB3).
\end{proof}

\begin{lemma}[Properness of a blow-up]
  \label{lemma:nagata-blowup-proper}
  \lean{}

  If $\mathcal I\subseteq\mathcal O_X$ is a finite-type quasi-coherent ideal,
  then the morphism
  \[
    b_{X,\mathcal I}\colon\operatorname{Bl}_{\mathcal I}(X)\to X
  \]
  of \Cref{def:nagata-global-blowup-map} is proper.  For every datum
  $(\mathcal I,X',\pi,\alpha_U)$ of
  \Cref{def:nagata-admissible-blowup}, substitution of
  $\pi=b_{X,\mathcal I}$ is the predicate
  \[
    \operatorname{IsProper}(\pi).
  \]
\end{lemma}

\begin{proof}
  By \Cref{def:nagata-global-blowup-map}, the displayed morphism is exactly
  $b_{X,\mathcal I}$.  Applying
  \Cref{thm:nagata-global-blowup-proper} to the finite-type ideal
  $\mathcal I$ is the predicate
  \[
    \operatorname{IsProper}(b_{X,\mathcal I}).
  \]
  If $(\mathcal I,X',\pi,\alpha_U)$ is a datum of
  \Cref{def:nagata-admissible-blowup}, then its defining equality is
  \[
    \pi=b_{X,\mathcal I}.
  \]
  The two preceding formulas form the implication
  \[
    \pi=b_{X,\mathcal I}
      \ \wedge\ 
    \operatorname{IsProper}(b_{X,\mathcal I})
      \quad\Longrightarrow\quad
    \operatorname{IsProper}(\pi).
  \]
\end{proof}

\begin{lemma}[Strict transform and restriction]
  \label{lemma:nagata-strict-transform-restrict}
  \uses{def:nagata-strict-transform}
  \lean{}

  Retain $f\colon Y\to X$, the finite-type quasi-coherent ideal
  $\mathcal I$, the blow-up $\pi\colon X'\to X$, and
  $j_U\colon U=X\setminus V(\mathcal I)\hookrightarrow X$ from
  \Cref{lemma:nagata-strict-transform-over-complement}.  Let
  $\ell\colon W\hookrightarrow U$ be an open immersion and define
  \[
    j_W:=j_U\circ\ell\colon W\hookrightarrow X.
  \]
  In particular, $W\subseteq U$.  Put
  \[
  \begin{gathered}
    X'_W:=X'\times_XW,\qquad
    \alpha_W:=\alpha_U\times_U\operatorname{id}_W
      \colon X'_W\xrightarrow{\sim}W,\\
    Y_W:=Y\times_XW,\qquad
    P:=Y\times_XX',\qquad
    P_W:=P\times_{X'}X'_W,\\
    S_W:=\operatorname{ST}_\pi(Y)\times_{X'}X'_W,\qquad
    \iota_{f,\pi,W}\colon S_W\hookrightarrow P_W
      \quad\text{the base change of }\iota_{f,\pi}.
  \end{gathered}                                                   \tag{SR1}
  \]
  There is a unique isomorphism
  \[
    r_W\colon S_W\xrightarrow{\sim}Y_W
  \]
  such that
  \[
    \operatorname{pr}_Y\circ r_W
      =\operatorname{pr}_Y\circ\operatorname{pr}_P\circ
        \iota_{f,\pi,W},
    \qquad
    \operatorname{pr}_W\circ r_W
      =\alpha_W\circ\operatorname{pr}_{X'_W}\circ
        \iota_{f,\pi,W}.
  \]
\end{lemma}

\begin{proof}
  Let
  \[
  \begin{gathered}
    c_S\colon S_W\xrightarrow{\sim}
      \operatorname{ST}_\pi(Y)_U\times_UW,\\
    c_Y\colon Y_W\xrightarrow{\sim}Y_U\times_UW
  \end{gathered}                                                    \tag{SR2}
  \]
  be the two fibre-product associativity isomorphisms.  Base change of
  $\alpha_U^{\mathrm{st}}$ from
  \Cref{lemma:nagata-strict-transform-over-complement} along
  $\ell\colon W\to U$ is the isomorphism
  \[
    \alpha_U^{\mathrm{st}}\times_U\operatorname{id}_W\colon
      \operatorname{ST}_\pi(Y)_U\times_UW
        \xrightarrow{\sim}Y_U\times_UW.
  \]
  Define
  \[
    r_W:=c_Y^{-1}\circ
      (\alpha_U^{\mathrm{st}}\times_U\operatorname{id}_W)\circ c_S.
                                                                    \tag{SR3}
  \]
  Base change of the two component identities for
  $\alpha_U^{\mathrm{st}}$ is
  \[
  \begin{aligned}
    \operatorname{pr}_Y\circ r_W
      &=\operatorname{pr}_Y\circ\operatorname{pr}_P\circ
        \iota_{f,\pi,W},\\
    \operatorname{pr}_W\circ r_W
      &=\alpha_W\circ\operatorname{pr}_{X'_W}\circ
        \iota_{f,\pi,W}.
  \end{aligned}
  \]
  The formula (SR3) determines $r_W$ uniquely.
\end{proof}

\begin{lemma}[{\cite[\href{https://stacks.math.columbia.edu/tag/080E}{Tag 080E}]{stacks-project}}]
  \label{lemma:nagata-strict-transform-ambient-closed}
  \uses{def:nagata-closed-subscheme-class,def:nagata-strict-transform,
    lemma:nagata-immersion-qcqs-criteria,
    lemma:nagata-morphism-property-calculus,
    thm:nagata-global-blowup-complement}
  \lean{}

  Let $i\colon Z\hookrightarrow X$ be a closed immersion, let
  $\mathcal I\subseteq\mathcal O_X$ be a finite-type quasi-coherent ideal,
  let
  $\pi=b_{X,\mathcal I}\colon X'=\operatorname{Bl}_{\mathcal I}(X)\to X$,
  and put
  \[
    \mathcal I_Z:=i^{-1}\mathcal I\mathcal O_Z,\qquad
    U:=X\setminus V(\mathcal I),\qquad
    j_U\colon U\hookrightarrow X.                               \tag{AC1}
  \]
  There is a unique
  isomorphism over $Z$
  \[
    \vartheta_{i,\mathcal I}\colon
      \operatorname{ST}_\pi(Z)\xrightarrow{\sim}
      \operatorname{Bl}_{\mathcal I_Z}(Z)
  \]
  satisfying
  \[
    b_{Z,\mathcal I_Z}\circ\vartheta_{i,\mathcal I}
      =\operatorname{pr}_Z\circ\iota_{i,\pi}.
  \]
  Define the base change of $\pi$ along $j_U$ and its source by
  \[
    X'_U:=X'\times_XU,
    \qquad
    \pi_U:=\operatorname{pr}_U\colon X'_U\longrightarrow U,
    \qquad
    j_U\circ\pi_U=\pi\circ\operatorname{pr}_{X'}.
  \]
  
  \Cref{thm:nagata-global-blowup-complement}, applied to the objects in
  (AC1), is the predicate $\operatorname{IsIso}(\pi_U)$.  Regard the same
  morphism as the resulting isomorphism and define
  \[
    \alpha_U:=\pi_U\colon X'_U\xrightarrow{\sim}U,
    \qquad
    s_U:=\operatorname{pr}_{X'}\circ\alpha_U^{-1}\colon U\to X'.
  \]
  The base-change equation and the inverse equation
  $\alpha_U\circ\alpha_U^{-1}=\operatorname{id}_U$ form the chain
  \[
    \pi\circ s_U
      =\pi\circ\operatorname{pr}_{X'}\circ\alpha_U^{-1}
      =j_U\circ\alpha_U\circ\alpha_U^{-1}
      =j_U.                                                       \tag{AC2}
  \]
  If $k\colon T\hookrightarrow Z$ is closed, define
  \[
    t:=i\circ k\colon T\to X,
  \]
  and define
  \[
    T_U:=T\times_{X,t,j_U}U,\qquad
    p_T\colon T_U\to T,\qquad p_U\colon T_U\to U,
  \]
  \[
    \lambda_{T,\pi}:=
      (p_T,s_U\circ p_U)
      \colon T_U\to T\times_XX'.
  \]
  The equality defining this fibre-product morphism is
  \[
    t\circ p_T=j_U\circ p_U
      =\pi\circ s_U\circ p_U.                                   \tag{AC3}
  \]
  Define the direct ambient construction by
  \[
    \mathsf A_{\mathrm{dir}}(T/X,\pi):=
      [\,\operatorname{ST}_\pi(T)
        \xhookrightarrow{\iota_{t,\pi}}T\times_XX'\,]
      \in\operatorname{ClosedSub}(T\times_XX').
  \]
  Put
  \[
    \pi_Z:=b_{Z,\mathcal I_Z}\colon
      \operatorname{Bl}_{\mathcal I_Z}(Z)\to Z
  \]
  and define the transport of the ambient fibre product by
  \[
  \begin{aligned}
    e_{T,Z}\colon
      T\times_Z\operatorname{Bl}_{\mathcal I_Z}(Z)
        &\longrightarrow T\times_XX',\\
    \operatorname{pr}_T\circ e_{T,Z}
      &=\operatorname{pr}_T,\\
    \operatorname{pr}_{X'}\circ e_{T,Z}
      &=\operatorname{pr}_{X'}\circ\iota_{i,\pi}\circ
        \vartheta_{i,\mathcal I}^{-1}\circ
        \operatorname{pr}_{\operatorname{Bl}_{\mathcal I_Z}(Z)}.
  \end{aligned}                                                   \tag{AC4}
  \]
  Define the transported ambient construction by
  \[
    \mathsf A_Z(T):=
      [\,\operatorname{ST}_{\pi_Z}(T)
        \xhookrightarrow{e_{T,Z}\circ\iota_{k,\pi_Z}}
          T\times_XX'\,]
      \in\operatorname{ClosedSub}(T\times_XX').
  \]
  Both named ambient constructions satisfy the class equality
  \[
    \mathsf A_{\mathrm{dir}}(T/X,\pi)
      =[\,\operatorname{SIm}_{T\times_XX'}(\lambda_{T,\pi})
        \xhookrightarrow{i_{\lambda_{T,\pi}}}T\times_XX'\,]
      =\mathsf A_Z(T)
      \quad\text{in }\operatorname{ClosedSub}(T\times_XX').    \tag{AC0}
  \]
\end{lemma}

\begin{proof}
  Apply \Cref{thm:nagata-strict-transform-is-blowup} with $Y=Z$ and $f=i$.
  The comparison of
  \Cref{def:nagata-strict-transform-blowup-comparison} is an isomorphism
  \[
    \Theta_{i,\mathcal I}\colon
      \operatorname{ST}_\pi(Z)\xrightarrow{\sim}
      \operatorname{Bl}_{\mathcal I_Z}(Z)
  \]
  and satisfies
  \[
    b_{Z,\mathcal I_Z}\circ\Theta_{i,\mathcal I}
      =\operatorname{pr}_Z\circ\iota_{i,\pi}.
  \]
  Define
  \[
    \vartheta_{i,\mathcal I}:=\Theta_{i,\mathcal I}.             \tag{AC5}
  \]
  This is the asserted isomorphism and the preceding displayed equation is
  its required equality over $Z$.

  For the last assertion, put
  \[
    U_Z:=Z\times_XU,
    \qquad
    c_T\colon T_U\xrightarrow{\sim}T\times_ZU_Z
  \]
  with $c_T$ the fibre-product associativity isomorphism.  Let
  \[
    \lambda_{k,\pi_Z}\colon
      T\times_ZU_Z\longrightarrow
      T\times_Z\operatorname{Bl}_{\mathcal I_Z}(Z)
  \]
  be the morphism of \Cref{def:nagata-strict-transform} for
  $(k,\pi_Z)$.  The preserved-open component equations for
  $\vartheta_{i,\mathcal I}$ are
  \[
  \begin{aligned}
    \operatorname{pr}_T\circ e_{T,Z}\circ
      \lambda_{k,\pi_Z}\circ c_T
      &=p_T,\\
    \operatorname{pr}_{X'}\circ e_{T,Z}\circ
      \lambda_{k,\pi_Z}\circ c_T
      &=s_U\circ p_U.
  \end{aligned}                                                   \tag{AC6}
  \]
  The corresponding two components of $\lambda_{T,\pi}$ are, by its
  definition,
  \[
    \operatorname{pr}_T\circ\lambda_{T,\pi}=p_T,
    \qquad
    \operatorname{pr}_{X'}\circ\lambda_{T,\pi}=s_U\circ p_U.
  \]
  Fibre-product uniqueness applied to these four equalities is
  \[
    e_{T,Z}\circ\lambda_{k,\pi_Z}\circ c_T
      =\lambda_{T,\pi}\colon T_U\to T\times_XX'.                \tag{AC7}
  \]
  The direct construction is, by \Cref{def:nagata-strict-transform},
  \[
    \mathsf A_{\mathrm{dir}}(T/X,\pi)
      =[\,\operatorname{SIm}_{T\times_XX'}(\lambda_{T,\pi})
        \xhookrightarrow{i_{\lambda_{T,\pi}}}T\times_XX'\,]
      \quad\text{in }\operatorname{ClosedSub}(T\times_XX').    \tag{AC8}
  \]
  Since $e_{T,Z}$ is the base change of the closed immersion
  $\iota_{i,\pi}\circ\vartheta_{i,\mathcal I}^{-1}$, it is a closed
  immersion.  By \Cref{lemma:nagata-immersion-qcqs-criteria}, (PC3) of
  \Cref{lemma:nagata-morphism-property-calculus}, and the QCQS predicate in
  \Cref{def:nagata-strict-transform}, the required predicates are
  \[
  \begin{aligned}
    \operatorname{IsClosedImmersion}(e_{T,Z})
      &\Longrightarrow\operatorname{IsQCQS}(e_{T,Z}),\\
    \operatorname{IsQCQS}(e_{T,Z})\wedge
      \operatorname{IsQCQS}(\lambda_{k,\pi_Z})
      &\Longrightarrow
        \operatorname{IsQCQS}(e_{T,Z}\circ\lambda_{k,\pi_Z}),\\
    \operatorname{IsIso}(c_T)
      &\Longrightarrow\operatorname{IsQCQS}(c_T),\\
    \operatorname{IsQCQS}(e_{T,Z}\circ\lambda_{k,\pi_Z})\wedge
      \operatorname{IsQCQS}(c_T)
      &\Longrightarrow
        \operatorname{IsQCQS}
          (e_{T,Z}\circ\lambda_{k,\pi_Z}\circ c_T).
  \end{aligned}                                                     \tag{AC8a}
  \]
  The universal property of the scheme-theoretic image,
  source invariance under the isomorphism $c_T$, and (AC7) form the chain
  \[
  \begin{aligned}
    \mathsf A_Z(T)
      &=[\,\operatorname{SIm}_{T\times_XX'}
          (e_{T,Z}\circ\lambda_{k,\pi_Z})
        \xhookrightarrow{i_{e_{T,Z}\circ\lambda_{k,\pi_Z}}}
          T\times_XX'\,]\\
      &=[\,\operatorname{SIm}_{T\times_XX'}
          (e_{T,Z}\circ\lambda_{k,\pi_Z}\circ c_T)
        \xhookrightarrow{i_{e_{T,Z}\circ\lambda_{k,\pi_Z}\circ c_T}}
          T\times_XX'\,]\\
      &=[\,\operatorname{SIm}_{T\times_XX'}(\lambda_{T,\pi})
        \xhookrightarrow{i_{\lambda_{T,\pi}}}T\times_XX'\,]
      \quad\text{in }\operatorname{ClosedSub}(T\times_XX').
  \end{aligned}                                                   \tag{AC9}
  \]
  Equations (AC8)--(AC9) are exactly (AC0).
\end{proof}

\begin{definition}[Finite-type extension of an ideal from a quasi-compact open]
  \label{def:nagata-finite-type-ideal-extension}
  \lean{}

  Let $X$ be quasi-compact and quasi-separated, let
  $\iota\colon W\hookrightarrow X$ be a quasi-compact open immersion, and let
  $\mathcal J\subseteq\mathcal O_W$ be a finite-type quasi-coherent ideal.
  Define the set
  \[
    \mathscr E_{X,W}(\mathcal J)
      :=\left\{
        \mathcal I\subseteq\mathcal O_X\
        \middle|\
        \begin{array}{l}
          \mathcal I\text{ is a finite-type quasi-coherent ideal},\\
          \iota^{-1}\mathcal I\cdot\mathcal O_W=\mathcal J
        \end{array}\right\}.
  \]
  The extension theorem
  \cite[\href{https://stacks.math.columbia.edu/tag/01PF}{Tag 01PF}]
    {stacks-project} is the assertion
  \[
    \mathscr E_{X,W}(\mathcal J)\neq\varnothing.
  \]
  By \Cref{def:nagata-choice-notation}, define
  \[
    \operatorname{Ext}_{X,W}(\mathcal J)
      :=\operatorname{ch}(\mathscr E_{X,W}(\mathcal J))
      \subseteq\mathcal O_X.
  \]
\end{definition}

\begin{lemma}[Finite-type ideal extension,
  {\cite[\href{https://stacks.math.columbia.edu/tag/01PF}{Tag 01PF}]{stacks-project}}]
  \label{lemma:nagata-extend-finite-type-ideal}
  \uses{def:nagata-finite-type-ideal-extension}
  \lean{}

  The sheaf $\operatorname{Ext}_{X,W}(\mathcal J)$ is a finite type
  quasi-coherent ideal on $X$, and its restriction to $W$ is
  $\mathcal J$.
\end{lemma}

\begin{proof}
  The extension theorem
  {\cite[\href{https://stacks.math.columbia.edu/tag/01PF}{Tag 01PF}]
  {stacks-project}}, applied to the tuple
  $(X,\iota\colon W\hookrightarrow X,\mathcal J)$, supplies an ideal
  $\mathcal I_0\subseteq\mathcal O_X$ with
  \[
    \mathcal I_0\text{ finite-type and quasi-coherent},
    \qquad
    \iota^{-1}\mathcal I_0\mathcal O_W=\mathcal J.
  \]
  By the set-builder in
  \Cref{def:nagata-finite-type-ideal-extension},
  \[
    \mathcal I_0\in\mathscr E_{X,W}(\mathcal J),
    \qquad
    \mathscr E_{X,W}(\mathcal J)\ne\varnothing.
  \]
  The defining property of the selected minimum is
  \[
    \operatorname{Ext}_{X,W}(\mathcal J)
      \in\mathscr E_{X,W}(\mathcal J).
  \]
  By the set-builder in
  \Cref{def:nagata-finite-type-ideal-extension}, this membership is the
  equivalence
  \[
  \begin{aligned}
    \operatorname{Ext}_{X,W}(\mathcal J)
      \in\mathscr E_{X,W}(\mathcal J)
    \quad\Longleftrightarrow\quad
    &\operatorname{Ext}_{X,W}(\mathcal J)
       \text{ is finite-type and quasi-coherent},\\
    &\iota^{-1}\operatorname{Ext}_{X,W}(\mathcal J)
       \mathcal O_W=\mathcal J.
  \end{aligned}
  \]
  Its right-hand side is the pair
  \[
  \begin{gathered}
    \operatorname{Ext}_{X,W}(\mathcal J)
       \text{ is finite-type and quasi-coherent},\\
    \iota^{-1}\operatorname{Ext}_{X,W}(\mathcal J)
       \mathcal O_W=\mathcal J.
  \end{gathered}
  \]
\end{proof}

\begin{definition}[Extension construction for an admissible blow-up]
  \label{def:nagata-extended-admissible-blowup}
  \uses{def:nagata-admissible-blowup,
    def:nagata-finite-type-ideal-extension,
    def:nagata-global-blowup}
  \lean{}

  Let $X$ be quasi-compact and quasi-separated, let
  $\iota_W\colon W\hookrightarrow X$ and
  $\iota_U\colon U\hookrightarrow X$ be quasi-compact open immersions, and
  put $U_W:=W\times_XU$.  Let
  \[
    (\mathcal J,W',\pi_W,\alpha_{U_W})
  \]
  be a $U_W$-admissible blow-up datum on $W$ in the sense of
  \Cref{def:nagata-admissible-blowup}; expanding that definition is
  \[
    W'=\operatorname{Bl}_{\mathcal J}(W),\quad
    \pi_W=b_{W,\mathcal J},\quad
    \mathcal J\mathcal O_{U_W}=\mathcal O_{U_W}.
  \]
  Define $\mathcal J^+\subseteq\mathcal O_{W\cup U}$ by
  \[
    \mathcal J^+\mathcal O_W=\mathcal J,\qquad
    \mathcal J^+\mathcal O_U=\mathcal O_U.
  \]
  This gluing is defined because its overlap equation is
  \[
    \mathcal J\mathcal O_{U_W}
      =\mathcal O_{U_W}
      =\mathcal O_U\mathcal O_{U_W}.
  \]
  Define
  \[
    \mathcal I_{X}(\mathcal J)
      :=\operatorname{Ext}_{X,W\cup U}(\mathcal J^+),
    \qquad
    X^{\mathrm{ext}}
      :=\operatorname{Bl}_{\mathcal I_X(\mathcal J)}(X),
    \qquad
    \pi^{\mathrm{ext}}
      :=b_{X,\mathcal I_X(\mathcal J)}.
  \]
  Define the restricted-ideal equality
  \[
    \vartheta_W\colon
      \mathcal I_X(\mathcal J)\mathcal O_W=\mathcal J.             \tag{EA0}
  \]
  Since \(\iota_W\) is flat, let
  \[
    \operatorname{bc}_{\iota_W,\vartheta_W}\colon
      \operatorname{Bl}_{\mathcal J}(W)
        \xrightarrow{\sim}X^{\mathrm{ext}}\times_XW                \tag{EA1}
  \]
  be the isomorphism of
  \Cref{thm:nagata-blowup-flat-base-change}, transported through
  \(\vartheta_W\), and define
  \[
    \rho_W:=\operatorname{bc}_{\iota_W,\vartheta_W}^{-1}
      \colon X^{\mathrm{ext}}\times_XW
        \xrightarrow{\sim}\operatorname{Bl}_{\mathcal J}(W)=W'.    \tag{EA2}
  \]
  Composing the structural equation for (EA1) with \(\rho_W\) is the chain
  \[
    \pi_W\circ\rho_W
      =\operatorname{pr}_W\circ
        \operatorname{bc}_{\iota_W,\vartheta_W}\circ\rho_W
      =\operatorname{pr}_W.                                      \tag{EA3}
  \]
  Define
  \[
    \mathsf{ExtBl}_{X,W,U}(\mathcal J)
      :=(\mathcal I_X(\mathcal J),X^{\mathrm{ext}},
         \pi^{\mathrm{ext}},\alpha_U^{\mathrm{ext}}),
  \]
  where
  \[
    \alpha_U^{\mathrm{ext}}\colon
      X^{\mathrm{ext}}\times_XU\xrightarrow{\sim}U
  \]
  is the complement isomorphism of
  \Cref{thm:nagata-global-blowup-complement}.
\end{definition}

\begin{lemma}[Property of the extended admissible blow-up,
  {\cite[\href{https://stacks.math.columbia.edu/tag/080M}{Tag 080M}]{stacks-project}}]
  \label{lemma:nagata-extend-admissible-blowup}
  \uses{def:nagata-extended-admissible-blowup}
  \lean{}

  The morphism
  $\pi^{\mathrm{ext}}\colon X^{\mathrm{ext}}\to X$ in
  $\mathsf{ExtBl}_{X,W,U}(\mathcal J)$ is $U$-admissible, and
  \[
    \rho_W\colon X^{\mathrm{ext}}\times_XW
      \xrightarrow{\sim}\operatorname{Bl}_{\mathcal J}(W)
  \]
  satisfies $\pi_W\circ\rho_W=\operatorname{pr}_W$.
\end{lemma}

\begin{proof}
  \uses{lemma:nagata-extend-finite-type-ideal,
    thm:nagata-blowup-flat-base-change}
  By \Cref{lemma:nagata-extend-finite-type-ideal}, applied to
  $\mathcal J^+$ on $W\cup U$,
  \[
  \begin{gathered}
    \mathcal I_X(\mathcal J)
      \text{ is finite-type and quasi-coherent},\\
    \mathcal I_X(\mathcal J)\mathcal O_{W\cup U}=\mathcal J^+.
  \end{gathered}
  \]
  Restriction to the two opens is the pair of equality chains
  \[
  \begin{aligned}
    \mathcal I_X(\mathcal J)\mathcal O_W
      &=(\mathcal I_X(\mathcal J)\mathcal O_{W\cup U})\mathcal O_W
       =\mathcal J^+\mathcal O_W=\mathcal J,\\
    \mathcal I_X(\mathcal J)\mathcal O_U
      &=(\mathcal I_X(\mathcal J)\mathcal O_{W\cup U})\mathcal O_U
       =\mathcal J^+\mathcal O_U=\mathcal O_U.
  \end{aligned}
  \]
  By \Cref{def:nagata-admissible-blowup}, the preceding finite-type and
  restriction equalities, together with the defining equations of
  $X^{\mathrm{ext}}$, $\pi^{\mathrm{ext}}$, and
  $\alpha_U^{\mathrm{ext}}$, are the conjunction
  \[
  \begin{gathered}
    \mathcal I_X(\mathcal J)\text{ finite-type and quasi-coherent},\qquad
    \mathcal I_X(\mathcal J)\mathcal O_U=\mathcal O_U,\\
    X^{\mathrm{ext}}
      =\operatorname{Bl}_{\mathcal I_X(\mathcal J)}(X),\qquad
    \pi^{\mathrm{ext}}=b_{X,\mathcal I_X(\mathcal J)},\\
    \operatorname{IsIso}(\alpha_U^{\mathrm{ext}}),\qquad
    \iota_U\circ\alpha_U^{\mathrm{ext}}
      =\pi^{\mathrm{ext}}\circ\operatorname{pr}_{X^{\mathrm{ext}}}.
  \end{gathered}
  \]
  Equations (EA0)--(EA3) are the
  typed isomorphism and structural equation
  \[
    \rho_W\colon
      X^{\mathrm{ext}}\times_XW
        \xrightarrow{\sim}\operatorname{Bl}_{\mathcal J}(W)=W',
    \qquad
    \pi_W\circ\rho_W=\operatorname{pr}_W.
  \]
  The equality
  $\mathcal I_X(\mathcal J)\mathcal O_U=\mathcal O_U$ and
  \Cref{thm:nagata-global-blowup-complement} define
  $\alpha_U^{\mathrm{ext}}$.
\end{proof}

\begin{lemma}[{\cite[\href{https://stacks.math.columbia.edu/tag/080A}{Tag 080A}]{stacks-project}}]
  \label{lemma:nagata-product-ideal-common-blowup}
  \lean{}

  Let $\mathcal I_1,\mathcal I_2\subseteq\mathcal O_X$ be finite-type
  quasi-coherent ideals.  The morphisms
  \[
  \begin{aligned}
    d_1&:=d_{\mathcal I_1}\colon
      \operatorname{Bl}_{\mathcal I_1\mathcal I_2}(X)
        \longrightarrow\operatorname{Bl}_{\mathcal I_1}(X),\\
    d_2&:=d_{\mathcal I_2}\colon
      \operatorname{Bl}_{\mathcal I_1\mathcal I_2}(X)
        \longrightarrow\operatorname{Bl}_{\mathcal I_2}(X)
  \end{aligned}
  \]
  are the domination maps of
  \Cref{def:nagata-blowup-product-domination-maps} and satisfy
  \[
    b_{X,\mathcal I_i}\circ d_i=b_{X,\mathcal I_1\mathcal I_2}
      \qquad(i=1,2).
  \]
  If $\mathcal I_i\mathcal O_U=\mathcal O_U$ for $i=1,2$, all three
  blow-up morphisms are $U$-admissible and
  \[
    d_i\times_XU=\operatorname{id}_U
  \]
  under the complement isomorphisms.
\end{lemma}

\begin{proof}
  \uses{cor:nagata-blowup-product-common-refinement,
    thm:nagata-global-blowup-complement}
  By \Cref{cor:nagata-blowup-product-common-refinement}, the maps
  $d_{\mathcal I_1}$ and $d_{\mathcal I_2}$ satisfy
  \[
    b_{X,\mathcal I_i}\circ d_{\mathcal I_i}
      =b_{X,\mathcal I_1\mathcal I_2}.
  \]
  If $\mathcal I_i\mathcal O_U=\mathcal O_U$, then
  \[
    (\mathcal I_1\mathcal I_2)\mathcal O_U
      =(\mathcal I_1\mathcal O_U)(\mathcal I_2\mathcal O_U)
      =\mathcal O_U\mathcal O_U=\mathcal O_U.
  \]
  The three unit-ideal equations required in
  \Cref{def:nagata-admissible-blowup} are
  \[
    \mathcal I_1\mathcal O_U=\mathcal O_U,\qquad
    \mathcal I_2\mathcal O_U=\mathcal O_U,\qquad
    (\mathcal I_1\mathcal I_2)\mathcal O_U=\mathcal O_U.
  \]
  If
  $\alpha_{12},\alpha_i$ are their complement isomorphisms from
  \Cref{thm:nagata-global-blowup-complement}, then
  \[
    \alpha_i\circ(d_i\times_XU)
      =\alpha_{12}.
  \]
  After identifying both source and target with $U$, this equality reads
  $d_i\times_XU=\operatorname{id}_U$.
\end{proof}

\begin{definition}[Elimination-of-indeterminacy construction]
  \label{def:nagata-elimination-indeterminacy}
  \uses{def:nagata-admissible-blowup,def:nagata-strict-transform}
  \lean{}

  Let $q\colon Z\to X$ be separated and of finite type, let $X$ be qcqs, let
  $j\colon U\hookrightarrow X$ be quasi-compact, and fix an isomorphism
  \[
    \eta\colon Z\times_XU\xrightarrow{\sim}U,
    \qquad
    j\circ\eta=q\circ\operatorname{pr}_Z.
  \]
  For every finite-type quasi-coherent ideal
  $\mathcal I\subseteq\mathcal O_X$ with
  $\mathcal I\mathcal O_U=\mathcal O_U$, put
  \[
  \begin{gathered}
    X_{\mathcal I}:=\operatorname{Bl}_{\mathcal I}(X),\qquad
    \pi_{\mathcal I}:=b_{X,\mathcal I},\\
    Z_{\mathcal I}:=\operatorname{ST}_{\pi_{\mathcal I}}(Z),\qquad
    q_{\mathcal I}:=q^{\mathrm{st}}_{\pi_{\mathcal I}}
       \colon Z_{\mathcal I}\to X_{\mathcal I},\\
    X_{\mathcal I,U}:=X_{\mathcal I}\times_XU,\qquad
    Z_{\mathcal I,U}:=
      Z_{\mathcal I}\times_{X_{\mathcal I}}X_{\mathcal I,U},\\
    (q_{\mathcal I})_U\colon
      Z_{\mathcal I,U}\to X_{\mathcal I,U}.
  \end{gathered}
  \]
  Define
  \[
    r_{\mathcal I}:=
      \alpha_U^{\mathrm{st}}(q,\pi_{\mathcal I})\colon
      Z_{\mathcal I,U}\xrightarrow{\sim}Z\times_XU                \tag{EI0}
  \]
  by \Cref{lemma:nagata-strict-transform-over-complement}, and define
  \[
    \alpha_{\mathcal I}\colon
      X_{\mathcal I,U}\xrightarrow{\sim}U                         \tag{EI1}
  \]
  by \Cref{thm:nagata-global-blowup-complement}.
  Define the set of elimination ideals
  \[
    \mathscr F(q,U,\eta):=
    \left\{\mathcal I\
      \middle|\
      \begin{array}{l}
      \mathcal I\subseteq\mathcal O_X
         \text{ is finite-type and quasi-coherent},\\
      \mathcal I\mathcal O_U=\mathcal O_U,\\
      q_{\mathcal I}\text{ is an open immersion},\\
      \alpha_{\mathcal I}\circ(q_{\mathcal I})_U
        =\eta\circ r_{\mathcal I}
      \end{array}\right\},
  \]
  The theorem
  {\cite[\href{https://stacks.math.columbia.edu/tag/081S}{Tag 081S}]
  {stacks-project}} states that $\mathscr F(q,U,\eta)\ne\varnothing$.
  By \Cref{def:nagata-choice-notation}, define
  \[
    \mathcal F_U(q,\eta):=\operatorname{ch}(\mathscr F(q,U,\eta)).
  \]
  Define
  \[
    X^{\mathrm{el}}(q,U)
      :=\operatorname{Bl}_{\mathcal F_U(q,\eta)}(X),
    \qquad
    Z^{\mathrm{el}}(q,U)=
      \operatorname{ST}_{b_{X,\mathcal F_U(q,\eta)}}(Z),
  \]
  and
  \[
    \pi^{\mathrm{el}}:=b_{X,\mathcal F_U(q,\eta)}
      \colon X^{\mathrm{el}}(q,U)\to X,
    \qquad
    q^{\mathrm{el}}:=q^{\mathrm{st}}_{\pi^{\mathrm{el}}}
      \colon Z^{\mathrm{el}}(q,U)\to X^{\mathrm{el}}(q,U).
  \]
\end{definition}

\begin{lemma}[Elimination of indeterminacy,
  {\cite[\href{https://stacks.math.columbia.edu/tag/081S}{Tag 081S}]{stacks-project}}]
  \label{lemma:nagata-elimination-indeterminacy}
  \uses{def:nagata-elimination-indeterminacy}
  \lean{}

  The morphism $X^{\mathrm{el}}(q,U)\to X$ is $U$-admissible, and the
  constructed strict-transform projection
  $q^{\mathrm{el}}\colon Z^{\mathrm{el}}(q,U)\to X^{\mathrm{el}}(q,U)$ is
  an open immersion.  Its restriction over $U$ is the original chosen
  isomorphism $q^{-1}(U)\cong U$.
\end{lemma}

\begin{proof}
  The cited elimination theorem in
  \Cref{def:nagata-elimination-indeterminacy} supplies the pair of predicates
  \[
    \mathcal I_0\in\mathscr F(q,U,\eta),
    \qquad
    \mathscr F(q,U,\eta)\ne\varnothing.
  \]
  By the defining property of the selected minimum,
  \[
    \mathcal F_U(q,\eta)\in\mathscr F(q,U,\eta).
  \]
  The set-builder in \Cref{def:nagata-elimination-indeterminacy} identifies
  this membership with a four-term conjunction.  Substitution of
  \[
    \pi^{\mathrm{el}}=b_{X,\mathcal F_U(q,\eta)},\qquad
    q^{\mathrm{el}}=q^{\mathrm{st}}_{\pi^{\mathrm{el}}}
  \]
  into that conjunction is
  \[
  \begin{gathered}
    \mathcal F_U(q,\eta)\text{ is finite-type and quasi-coherent},\\
    \mathcal F_U(q,\eta)\mathcal O_U=\mathcal O_U,\\
    q^{\mathrm{el}}\colon Z^{\mathrm{el}}(q,U)
       \hookrightarrow X^{\mathrm{el}}(q,U)
       \text{ is an open immersion},\\
    \alpha_{\mathcal F_U(q,\eta)}\circ(q^{\mathrm{el}})_U
      =\eta\circ r_{\mathcal F_U(q,\eta)}.
  \end{gathered}
  \]
  The first two lines and
  \Cref{def:nagata-admissible-blowup} state exactly that
  $\pi^{\mathrm{el}}$ is $U$-admissible; the last two lines are the remaining
  assertions.
\end{proof}

\begin{lemma}[A proper dense open immersion is an isomorphism]
  \label{lemma:nagata-proper-open-dense-iso}
  \uses{def:nagata-closed-subscheme-class,
    lemma:nagata-immersion-qcqs-criteria}
  \lean{}

  Let $q\colon Z\to X$ be proper and an open immersion.  Let
  $d\colon D\hookrightarrow X$ satisfy
  $\operatorname{IsQuasiCompactOpenImmersion}(d)$.  By
  \Cref{lemma:nagata-immersion-qcqs-criteria},
  \[
    \operatorname{IsQuasiCompactOpenImmersion}(d)
      \Longrightarrow\operatorname{IsQCQS}(d).                   \tag{PDO0}
  \]
  Suppose that
  \[
    [\,\operatorname{SIm}_X(d)\xhookrightarrow{i_d}X\,]
      =[\,X\xhookrightarrow{\operatorname{id}_X}X\,]
      \quad\text{in }\operatorname{ClosedSub}(X)
  \]
  as in \Cref{def:nagata-scheme-theoretic-image}, and assume there is
  $e\colon D\to Z$ with $q\circ e=d$.  Then $q$ is an isomorphism.
\end{lemma}

\begin{proof}
  \uses{def:nagata-choice-notation,def:nagata-closed-subscheme-class,
    def:nagata-morphism-predicate-notation,
    lemma:nagata-scheme-image-universal}
  By the hypotheses and (P0)--(P0a),
  \[
    \operatorname{IsProper}(q),
    \qquad
    \operatorname{IsOpenImmersion}(q)
      \Longrightarrow\operatorname{IsImmersion}(q).
  \]
  The proper-immersion criterion
  \cite[\href{https://stacks.math.columbia.edu/tag/01IQ}{Tag 01IQ}]
  {stacks-project} is the implication
  \[
    \operatorname{IsProper}(q)\wedge\operatorname{IsImmersion}(q)
      \Longrightarrow\operatorname{IsClosedImmersion}(q).       \tag{PD1}
  \]
  Combining the two displayed implications is the predicate chain
  \[
    \operatorname{IsProper}(q)\wedge
      \operatorname{IsOpenImmersion}(q)
    \Longrightarrow
    \operatorname{IsProper}(q)\wedge
      \operatorname{IsImmersion}(q)
    \Longrightarrow
    \operatorname{IsClosedImmersion}(q).
  \]
  The factorization data are the equations
  \[
    d=q\circ e,\qquad q\colon Z\hookrightarrow X.
  \]
  Apply \Cref{lemma:nagata-scheme-image-universal} to this factorization.  It
  defines
  \[
    c_{d,Z}\colon\operatorname{SIm}(d)\longrightarrow Z
  \]
  satisfying
  \[
    q\circ c_{d,Z}
      =i_d\colon\operatorname{SIm}_X(d)\hookrightarrow X.
  \]
  Define
  \[
    \mathscr R_d:=
      \left\{\delta\in\operatorname{Hom}(\operatorname{SIm}_X(d),X)
        \ \middle|\
        \operatorname{IsIso}(\delta)\ \wedge\
        \operatorname{id}_X\circ\delta=i_d\right\}.
  \]
  By (CSub3) of \Cref{def:nagata-closed-subscheme-class}, the hypothesis is
  $\mathscr R_d\ne\varnothing$.  Define the chosen witness
  \[
    \delta_d:=\operatorname{ch}(\mathscr R_d)\in\mathscr R_d,
    \qquad
    \delta_d\colon\operatorname{SIm}_X(d)\xrightarrow{\sim}X,
    \qquad
    \delta_d=i_d.
  \]
  Put $r:=c_{d,Z}\circ\delta_d^{-1}\colon X\to Z$.  Expanding this
  definition and the preceding equalities is the chain
  \[
    q\circ r
      =q\circ c_{d,Z}\circ\delta_d^{-1}
      =i_d\circ\delta_d^{-1}
      =\delta_d\circ\delta_d^{-1}
      =\operatorname{id}_X.
  \]
  Since every immersion is a monomorphism,
  \[
    q\circ(r\circ q)=(q\circ r)\circ q=q
      =q\circ\operatorname{id}_Z
    \quad\Longrightarrow\quad
    r\circ q=\operatorname{id}_Z.
  \]
  The two inverse equations are
  \[
    q\circ r=\operatorname{id}_X,\qquad
    r\circ q=\operatorname{id}_Z,\qquad
    r=q^{-1},\qquad
    \operatorname{IsIso}(q).
  \]
\end{proof}

\section{Separating boundary pieces}

\begin{lemma}[Finite-presentation closed immersions and finite type ideals]
  \label{lemma:nagata-closed-fp-ideal-finite-type}
  \lean{}

  Let $i\colon Z\hookrightarrow X$ be a closed immersion.  By
  \Cref{def:nagata-direct-image-structure-map}, define
  \[
    \varepsilon_i:=i^\flat\colon
      \mathcal O_X\twoheadrightarrow i_*\mathcal O_Z,
    \qquad
    \mathcal I_i:=\ker(\varepsilon_i)
      \subseteq\mathcal O_X.
  \]
  Then
  \[
    i\text{ is of finite presentation}
      \Longleftrightarrow
    \mathcal I_i\text{ is a finite-type quasi-coherent ideal}.
  \]
\end{lemma}

\begin{proof}
  Let $V=\operatorname{Spec}(A)\subseteq X$ be affine and put
  \[
    I_V:=\Gamma(V,\mathcal I_i).
  \]
  The affine closed-immersion correspondence identifies
  \[
    Z\times_XV=\operatorname{Spec}(A/I_V),
    \qquad
    i|_V=\operatorname{Spec}(A\twoheadrightarrow A/I_V).
  \]
  For elements $f_1,\ldots,f_n\in A$, the following equivalence is the
  presentation criterion for a quotient algebra:
  \[
  \begin{aligned}
    I_V=(f_1,\ldots,f_n)
      &\Longleftrightarrow
        A/I_V\cong
        A[x_1,\ldots,x_n]/(x_1,\ldots,x_n,f_1,\ldots,f_n)\\
      &\Longleftrightarrow
        A/I_V\text{ is finitely presented as an }A\text{-algebra}.
  \end{aligned}
  \]
  The affine closed-immersion correspondence substitutes the preceding
  quotient-algebra equivalence as
  \[
    i|_V\text{ is of finite presentation}
      \Longleftrightarrow I_V\text{ is finitely generated}.
  \]
  Both predicates are local on the affine open cover $(V)$, and
  $\mathcal I_i|_V=\widetilde{I_V}$.  Their local-to-global equivalences form
  the chain
  \[
  \begin{aligned}
    i\text{ is of finite presentation}
      &\Longleftrightarrow
        \forall V,\ i|_V\text{ is of finite presentation}\\
      &\Longleftrightarrow
        \forall V,\ I_V\text{ is finitely generated}\\
      &\Longleftrightarrow
        \mathcal I_i\text{ is finite-type and quasi-coherent}.
  \end{aligned}
  \]
\end{proof}

\begin{definition}[Finite-type ideal attached to a quasi-compact open]
  \label{def:nagata-qc-open-complement-ideal}
  \lean{}

  Let $Y$ be qcqs and let $w\colon W\hookrightarrow Y$ be a quasi-compact
  open immersion.  Define
  \[
    \mathscr C_Y(W):=
      \left\{\mathcal I\subseteq\mathcal O_Y\
      \middle|\
      \begin{array}{l}
        \mathcal I\text{ is finite-type and quasi-coherent},\\
        V(\mathcal I)=Y\setminus W
      \end{array}\right\}.
  \]
  The finite-type ideal theorem
  {\cite[\href{https://stacks.math.columbia.edu/tag/01PH}{Tag 01PH}]
  {stacks-project}} states that $\mathscr C_Y(W)\ne\varnothing$.  By
  \Cref{def:nagata-choice-notation}, define
  \[
    \mathcal I_Y(W):=\operatorname{ch}(\mathscr C_Y(W))
      \subseteq\mathcal O_Y.
  \]
\end{definition}

\begin{lemma}[Vanishing set of the open-complement ideal]
  \label{lemma:nagata-qc-open-complement-ideal}
  \uses{def:nagata-qc-open-complement-ideal}
  \lean{}

  The ideal $\mathcal I_Y(W)$ is finite type and quasi-coherent, and
  \[
    V(\mathcal I_Y(W))=Y\setminus W.
  \]
\end{lemma}

\begin{proof}
  Apply
  {\cite[\href{https://stacks.math.columbia.edu/tag/01PH}{Tag 01PH}]
  {stacks-project}} to the closed subset $T:=Y\setminus W$, whose complement
  $W$ is quasi-compact.  It produces
  $\mathcal I_0\subseteq\mathcal O_Y$ satisfying
  \[
    \mathcal I_0\text{ finite-type and quasi-coherent},
    \qquad V(\mathcal I_0)=T=Y\setminus W.
  \]
  The two displayed properties and the set equality in
  \Cref{def:nagata-qc-open-complement-ideal} establish
  \[
    \mathcal I_0\in\mathscr C_Y(W),
    \qquad \mathscr C_Y(W)\ne\varnothing.
  \]
  By \Cref{def:nagata-qc-open-complement-ideal},
  \[
    \mathcal I_Y(W)=\operatorname{ch}(\mathscr C_Y(W))
      \in\mathscr C_Y(W).
  \]
  The set-builder in \Cref{def:nagata-qc-open-complement-ideal} is the
  equivalence
  \[
    \mathcal I_Y(W)\in\mathscr C_Y(W)
      \Longleftrightarrow
      \left(
        \mathcal I_Y(W)\text{ is finite-type and quasi-coherent}
        \ \wedge\ 
        V(\mathcal I_Y(W))=Y\setminus W
      \right).
  \]
  Its right-hand side is
  \[
    \mathcal I_Y(W)\text{ finite-type and quasi-coherent},
    \qquad V(\mathcal I_Y(W))=Y\setminus W.
  \]
\end{proof}

\begin{definition}[Closed-boundary separation blow-up]
  \label{def:nagata-separate-closed-subschemes}
  \uses{def:nagata-global-blowup,def:nagata-strict-transform}
  \lean{}

  Let $j\colon U\hookrightarrow X$ be a quasi-compact open immersion and let
  $i_r\colon Z_r=V(\mathcal I_r)\hookrightarrow X$ $(r=1,2)$ be
  finite-presentation closed immersions such that
  \[
    Z_1\times_XZ_2\times_XU=\varnothing.
  \]
  By \Cref{lemma:nagata-closed-fp-ideal-finite-type}, each $\mathcal I_r$ is
  finite-type.  Define
  \[
  \begin{gathered}
    \mathcal H_{12}:=\mathcal I_1+\mathcal I_2
      \subseteq\mathcal O_X,\qquad
    X^{\mathrm{sep}}_{12}:=
      \operatorname{Bl}_{\mathcal H_{12}}(X),\\
    \pi_{12}^{\mathrm{sep}}:=b_{X,\mathcal H_{12}}
      \colon X^{\mathrm{sep}}_{12}\to X,\\
    S_r:=\operatorname{ST}_{\pi_{12}^{\mathrm{sep}}}(Z_r),\\
    s_r:=\operatorname{pr}_{X^{\mathrm{sep}}_{12}}\circ
      \iota_{i_r,\pi_{12}^{\mathrm{sep}}}
      \colon S_r\hookrightarrow X^{\mathrm{sep}}_{12}
      \qquad(r=1,2).
  \end{gathered}
  \]
  The notation
  \[
    \mathsf{SepBl}_U(Z_1,Z_2)
      :=(\mathcal H_{12},X^{\mathrm{sep}}_{12},
          \pi_{12}^{\mathrm{sep}},\alpha_U)
  \]
  denotes this tuple, where $\alpha_U$ is the complement isomorphism of
  \Cref{thm:nagata-global-blowup-complement}.
\end{definition}

\begin{lemma}[Properties of the closed-boundary separation blow-up,
  {\cite[\href{https://stacks.math.columbia.edu/tag/0F3V}{Tag 0F3V}]{stacks-project}}]
  \label{lemma:nagata-separate-closed-subschemes}
  \uses{def:nagata-separate-closed-subschemes}
  \lean{}

  The morphism $\mathsf{SepBl}_U(Z_1,Z_2)\to X$ is $U$-admissible, and the
  strict transforms of $Z_1$ and $Z_2$ in it are disjoint.
\end{lemma}

\begin{proof}
  \uses{lemma:nagata-closed-fp-ideal-finite-type,
    lemma:nagata-qc-open-complement-ideal,
    lemma:nagata-strict-transform-ambient-closed}
  By \Cref{lemma:nagata-closed-fp-ideal-finite-type},
  $\mathcal I_1$ and $\mathcal I_2$ are finite-type, whence so is
  $\mathcal H_{12}=\mathcal I_1+\mathcal I_2$.  The zero-locus identities are
  \[
  \begin{aligned}
    V(\mathcal H_{12})
      &=V(\mathcal I_1+\mathcal I_2)\\
      &=V(\mathcal I_1)\cap V(\mathcal I_2)
       =Z_1\times_XZ_2.
  \end{aligned}
  \]
  Therefore
  \[
    V(\mathcal H_{12})\times_XU
      =Z_1\times_XZ_2\times_XU=\varnothing,
    \qquad
    \mathcal H_{12}\mathcal O_U=\mathcal O_U.
  \]
  This is exactly the admissibility condition of
  \Cref{def:nagata-admissible-blowup}.

  Work on an affine $V=\operatorname{Spec}(A)\subseteq X$ with
  \[
    I_1=(a_1,\ldots,a_m),\qquad
    I_2=(b_1,\ldots,b_n),\qquad H=I_1+I_2.
  \]
  The charts indexed by $a_r$ and $b_s$ cover
  $\operatorname{Bl}_H(V)$ by
  \Cref{prop:nagata-affine-blowup-chart-cover}.  On the $a_r$-chart, apply
  \Cref{def:nagata-strict-transform-chart-ring} to $B=A/I_1$.  Since the
  distinguished element is $a_r=0$ in $B$, every $c$ in the base-change
  chart ring satisfies
  \[
    0^1c=0,\qquad
    \Gamma_{(0)}(C_{a_r}(A,A/I_1,H))=C_{a_r}(A,A/I_1,H),
  \]
  and consequently
  \[
    C_{a_r}^{\mathrm{st}}(A,A/I_1,H)=0.
  \]
  Thus $S_1$ is empty on every $a_r$-chart.  Under the substitution
  \[
    (I_1,a_r,S_1)\longleftrightarrow(I_2,b_s,S_2),
  \]
  the preceding chart equality is
  \[
    C_{b_s}^{\mathrm{st}}(A,A/I_2,H)=0,
  \]
  so $S_2$ is empty on every $b_s$-chart.  On each chart of the displayed
  cover, at least one factor of $S_1\times_{X^{\mathrm{sep}}_{12}}S_2$ is
  empty.  Hence
  \[
    S_1\times_{X^{\mathrm{sep}}_{12}}S_2=\varnothing.
  \]
  The identification of these chartwise transforms with the transforms in
  the ambient blow-up is
  \Cref{lemma:nagata-strict-transform-ambient-closed}.
\end{proof}

\begin{definition}[Ideal whose support is a constructible closed subset]
  \label{def:nagata-constructible-closed-support}
  \uses{def:nagata-qc-open-complement-ideal}
  \lean{}

  Let $T\subseteq|U|$ be constructible and closed, with $U$ qcqs.  Define
  \[
    W_T:=U\setminus T,\qquad
    \mathcal J_U(T):=\mathcal I_U(W_T).
  \]
  Here $W_T$ is quasi-compact because it is open and constructible in the
  spectral space $|U|$.
\end{definition}

\begin{lemma}[Property of the constructible-support ideal,
  {\cite[\href{https://stacks.math.columbia.edu/tag/01PH}{Tag 01PH}]{stacks-project}}]
  \label{lemma:nagata-constructible-closed-support}
  \uses{def:nagata-constructible-closed-support}
  \lean{}

  The ideal $\mathcal J_U(T)$ is finite type and quasi-coherent, and its
  vanishing set is exactly $T$.
\end{lemma}

\begin{proof}
  \uses{lemma:nagata-qc-open-complement-ideal}
  By \Cref{def:nagata-constructible-closed-support},
  \[
    \mathcal J_U(T)=\mathcal I_U(W_T),
    \qquad W_T=U\setminus T.
  \]
  The instance of \Cref{lemma:nagata-qc-open-complement-ideal} at
  $(Y,W)=(U,W_T)$ is
  \[
    \mathcal J_U(T)\text{ finite-type and quasi-coherent},
    \qquad
    V(\mathcal J_U(T))=U\setminus W_T.
  \]
  Substitution of $W_T=U\setminus T$ is the equality chain
  \[
    V(\mathcal J_U(T))
      =U\setminus(U\setminus T)=T.
  \]
\end{proof}

\begin{definition}[Constructible-boundary separation blow-up]
  \label{def:nagata-separate-constructible-closed}
  \uses{def:nagata-constructible-closed-support,
    def:nagata-finite-type-ideal-extension,
    def:nagata-separate-closed-subschemes,
    lemma:nagata-morphism-property-calculus}
  \lean{}

  Let $j\colon U\hookrightarrow X$ be a quasi-compact open immersion and let
  $T_1,T_2\subseteq|U|$ be disjoint constructible subsets closed in $|U|$.
  Define
  \[
    \mathcal J_i:=\mathcal J_U(T_i),\qquad
    \mathcal I_i:=\operatorname{Ext}_{X,U}(\mathcal J_i),\qquad
    Z_i:=V(\mathcal I_i),\qquad
    z_i\colon Z_i\hookrightarrow X\qquad(i=1,2),
  \]
  and define the four components of the separation datum by
  \[
  \begin{gathered}
    (\mathcal H_C,X_C,\pi_C,\alpha_U)
      :=\mathsf{CSepBl}_U(T_1,T_2)
      :=\mathsf{SepBl}_U(Z_1,Z_2),\\
    \pi_C\colon X_C\to X,\qquad
    \alpha_U\colon X_C\times_XU\xrightarrow{\sim}U.
  \end{gathered}                                                    \tag{KCS0}
  \]
  Define every morphism entering the two scheme-theoretic images by
  \[
  \begin{gathered}
    \mathbf T_i:=V(\sqrt{\mathcal J_i}),\qquad
    \iota_i\colon\mathbf T_i\hookrightarrow U,\\
    p_C:=\operatorname{pr}_{X_C}\colon X_C\times_XU\to X_C,
    \qquad
    a_U^C:=p_C\circ\alpha_U^{-1}\colon U\to X_C,\\
    \tau_i:=a_U^C\circ\iota_i\colon\mathbf T_i\to X_C
      \qquad(i=1,2).
  \end{gathered}                                                    \tag{KCS1}
  \]
  The base-change and composition implications (PC2)--(PC3) of
  \Cref{lemma:nagata-morphism-property-calculus} are the following predicate
  chain:
  \[
  \begin{gathered}
    \operatorname{IsQuasiCompactOpenImmersion}(j)
      \Longrightarrow
    \operatorname{IsQuasiCompactOpenImmersion}(p_C),\\
    \operatorname{IsIso}(\alpha_U^{-1})
      \Longrightarrow
    \operatorname{IsQuasiCompactOpenImmersion}(\alpha_U^{-1}),\\
    \operatorname{IsQuasiCompactOpenImmersion}(p_C)\ \wedge\
      \operatorname{IsQuasiCompactOpenImmersion}(\alpha_U^{-1})
      \Longrightarrow
    \operatorname{IsQuasiCompactOpenImmersion}(a_U^C)
      \Longrightarrow\operatorname{IsQCQS}(a_U^C),\\
    \operatorname{IsClosedImmersion}(\iota_i)
      \Longrightarrow\operatorname{IsQCQS}(\iota_i),\\
    \operatorname{IsQCQS}(a_U^C)\ \wedge\
      \operatorname{IsQCQS}(\iota_i)\ \wedge\
      (\tau_i=a_U^C\circ\iota_i)
      \Longrightarrow\operatorname{IsQCQS}(\tau_i)
      \qquad(i=1,2).
  \end{gathered}                                                    \tag{KCS2}
  \]
  Only after (KCS2), define
  \[
  \begin{gathered}
    \overline{\mathbf T_i}^{\,X_C}
      :=\operatorname{SIm}_{X_C}(\tau_i),
    \qquad
    i_i^C:=i_{\tau_i}\colon
      \overline{\mathbf T_i}^{\,X_C}\hookrightarrow X_C.
  \end{gathered}
  \]
\end{definition}

\begin{lemma}[Properties of constructible-boundary separation]
  \label{lemma:nagata-separate-constructible-closed}
  \uses{def:nagata-separate-constructible-closed}
  \lean{}

  The morphism $\pi_C\colon X_C\to X$ is $U$-admissible, and
  \[
    \overline{\mathbf T_1}^{\,X_C}
      \times_{X_C}\overline{\mathbf T_2}^{\,X_C}
      =\varnothing.
  \]
\end{lemma}

\begin{proof}
  \uses{lemma:nagata-constructible-closed-support,
    lemma:nagata-extend-finite-type-ideal,
    lemma:nagata-separate-closed-subschemes}
  By \Cref{lemma:nagata-constructible-closed-support} and
  \Cref{lemma:nagata-extend-finite-type-ideal},
  \[
  \begin{aligned}
    \mathcal I_i\mathcal O_U&=\mathcal J_i,\\
    |Z_i\times_XU|
      &=V(\mathcal I_i\mathcal O_U)
       =V(\mathcal J_i)=T_i.
  \end{aligned}
  \]
  Since $T_1\cap T_2=\varnothing$,
  \[
    |Z_1\times_XZ_2\times_XU|
      =T_1\cap T_2=\varnothing,
  \]
  hence $Z_1\times_XZ_2\times_XU=\varnothing$.  Applying
  the conclusion of \Cref{lemma:nagata-separate-closed-subschemes} is
  \[
    \pi_C\text{ is }U\text{-admissible},
    \qquad
    S_1\times_{X_C}S_2=\varnothing,
  \]
  where $S_i:=\operatorname{ST}_{\pi_C}(Z_i)$.  Define its closed immersion
  into $X_C$ by
  \[
    s_i^C:=\operatorname{pr}_{X_C}\circ\iota_{z_i,\pi_C}
      \colon S_i\hookrightarrow X_C.
  \]
  The morphism $\tau_i$ factors through $S_i$: on the preserved open this is
  the factorization
  \[
    \mathbf T_i\hookrightarrow Z_i\times_XU
      \longrightarrow S_i\times_{X_C}(X_C\times_XU)=Z_i\times_XU
  \]
  from \Cref{lemma:nagata-strict-transform-over-complement}.  By
  \Cref{lemma:nagata-scheme-image-universal}, this factorization defines the
  typed closed immersion
  \[
    c_i^T\colon
      \overline{\mathbf T_i}^{\,X_C}\hookrightarrow S_i,
    \qquad
    s_i^C\circ c_i^T=i_i^C.
  \]
  Its fibre product is the closed immersion
  \[
    c_1^T\times_{X_C}c_2^T\colon
      \overline{\mathbf T_1}^{\,X_C}
        \times_{X_C}\overline{\mathbf T_2}^{\,X_C}
      \hookrightarrow S_1\times_{X_C}S_2=\varnothing,
  \]
  so the left-hand side is empty.
\end{proof}

\begin{definition}[Constructible separation data in a spectral space]
  \label{def:nagata-spectral-constructible-separation}
  \lean{}

  Let $T,T'\subseteq V$ be disjoint closed subsets of a spectral space $V$.
  Define
  \[
    \mathscr S_V(T,T'):=
      \left\{(W,W')\
      \middle|\
      \begin{array}{l}
        W,W'\subseteq V\text{ are quasi-compact open},\\
        T'\subseteq W\subseteq V\setminus T,\\
        V\setminus W\subseteq W'\subseteq V\setminus T'
      \end{array}\right\}.
  \]
  Define the quasi-compact-open basis
  \[
    \mathcal B_{\mathrm{qc}}(V)
      :=\{Q\subseteq V\mid Q\text{ is open and quasi-compact}\}.
  \]
  Since $T'\subseteq V\setminus T$ and $T'$ is quasi-compact, choose a finite
  set $A$ and $(Q_a)_{a\in A}\in\mathcal B_{\mathrm{qc}}(V)^A$ satisfying
  \[
    T'\subseteq W_0:=\bigcup_{a\in A}Q_a\subseteq V\setminus T.   \tag{CS0}
  \]
  The set $V\setminus W_0$ is closed in the quasi-compact space $V$, hence
  quasi-compact, and (CS0) is the containment
  $V\setminus W_0\subseteq V\setminus T'$.  Choose a finite set $B$ and
  $(Q'_b)_{b\in B}\in\mathcal B_{\mathrm{qc}}(V)^B$ satisfying
  \[
    V\setminus W_0\subseteq
      W'_0:=\bigcup_{b\in B}Q'_b\subseteq V\setminus T'.          \tag{CS1}
  \]
  Finite unions in (CS0)--(CS1) are quasi-compact open, and substitution in
  the defining set is the membership
  \[
    (W_0,W'_0)\in\mathscr S_V(T,T'),
    \qquad
    \mathscr S_V(T,T')\neq\varnothing.                            \tag{CS2}
  \]
  By \Cref{def:nagata-choice-notation}, define
  \[
    (W(T,T'),W'(T,T'))
      :=\operatorname{ch}(\mathscr S_V(T,T')),
  \]
  and
  \[
    T_c:=V\setminus W(T,T'),\qquad
    T'_c:=V\setminus W'(T,T').
  \]
\end{definition}

\begin{lemma}[Properties of the constructible separation data]
  \label{lemma:nagata-spectral-constructible-separation}
  \uses{def:nagata-spectral-constructible-separation}
  \lean{}

  The constructed subsets $T_c,T'_c$ are constructible and closed in $V$,
  and they satisfy
  \[
    T\subseteq T_c,\qquad T'\subseteq T'_c,\qquad
    T_c\cap T'_c=\varnothing.
  \]
\end{lemma}

\begin{proof}
  For each $x\in T'$, the open $V\setminus T$ contains $x$.  The
  quasi-compact-open basis supplies $B_x$ with
  \[
    x\in B_x\subseteq V\setminus T.
  \]
  Since $T'$ is closed in the quasi-compact space $V$, choose
  $x_1,\ldots,x_m$ such that
  \[
    T'\subseteq W_0:=\bigcup_{r=1}^mB_{x_r}\subseteq V\setminus T.
  \]
  The finite union $W_0$ is quasi-compact open.  For each
  $y\in V\setminus W_0$, choose a quasi-compact open
  $B'_y$ satisfying
  \[
    y\in B'_y\subseteq V\setminus T'.
  \]
  Since $V\setminus W_0$ is closed in the quasi-compact space $V$, choose
  $y_1,\ldots,y_n$ and put
  \[
    V\setminus W_0\subseteq
      W'_0:=\bigcup_{s=1}^nB'_{y_s}\subseteq V\setminus T'.
  \]
  Thus $(W_0,W'_0)\in\mathscr S_V(T,T')$, so the selected
  $(W,W')=(W(T,T'),W'(T,T'))$ belongs to this set.  The set-builder in
  \Cref{def:nagata-spectral-constructible-separation} is the equivalence
  \[
  \begin{aligned}
    (W,W')\in\mathscr S_V(T,T')
      \Longleftrightarrow{}&
      W,W'\text{ are quasi-compact open},\\
      &T'\subseteq W\subseteq V\setminus T,\\
      &V\setminus W\subseteq W'\subseteq V\setminus T'.
  \end{aligned}
  \]
  Its right-hand side contains the relations
  \[
    T'\subseteq W,\quad W\cap T=\varnothing,\quad
    V\setminus W\subseteq W',\quad W'\cap T'=\varnothing.
  \]
  Therefore
  \[
  \begin{aligned}
    T\subseteq V\setminus W&=T_c,\\
    T'\subseteq V\setminus W'&=T'_c,\\
    T_c\cap T'_c
      &=(V\setminus W)\cap(V\setminus W')
       \subseteq W'\cap(V\setminus W')=\varnothing.
  \end{aligned}
  \]
  Finally, $T_c$ and $T'_c$ are closed and are complements of quasi-compact
  opens; hence they are constructible by the definition of the constructible
  topology on a spectral space.
\end{proof}

\begin{definition}[Modification making a proper map an isomorphism near a closure]
  \label{def:nagata-isomorphism-near-closure}
  \uses{def:nagata-spectral-constructible-separation,
    def:nagata-separate-constructible-closed,
    def:nagata-elimination-indeterminacy,
    def:nagata-extended-admissible-blowup,
    lemma:nagata-strict-transform-restrict}
  \lean{}

  Let $q\colon Z\to X$ be proper, let $v\colon V\hookrightarrow X$ be a
  quasi-compact open immersion, and let $T\subseteq|V|$ be closed.  Put
  \[
    Z_V:=Z\times_{X,q,v}V,
    \qquad
    q_V:=\operatorname{pr}_V\colon Z_V\longrightarrow V.           \tag{INC-1}
  \]
  For every open subscheme $N\subseteq V$, define
  \[
    n_N\colon N\hookrightarrow V,\qquad
    Z_{V,N}:=Z_V\times_{V,q_V,n_N}N,\qquad
    q_{V,N}:=\operatorname{pr}_N\colon Z_{V,N}\to N.               \tag{INC-0}
  \]
  Define the isomorphism locus and its closed complement by
  \[
  \begin{aligned}
    V^{\mathrm{iso}}(q)
      &:=\bigcup\{N\subseteq V\mid N\text{ open and }
         \operatorname{IsIso}(q_{V,N})\},\\
    E(q,V)&:=V\setminus V^{\mathrm{iso}}(q).
  \end{aligned}
  \]
  Assume $T\subseteq V^{\mathrm{iso}}(q)$, so
  $T\cap E(q,V)=\varnothing$.  Apply
  \Cref{def:nagata-spectral-constructible-separation} and define
  \[
    (T_c,E_c):=
      (T_c(T,E(q,V)),T'_c(T,E(q,V))).
  \]
  Apply \Cref{def:nagata-separate-constructible-closed} and define its four
  components by
  \[
    (\mathcal H_1,X_1,\pi_1,\alpha_V^1)
      :=\mathsf{CSepBl}_V(T_c,E_c),\qquad
    \pi_1\colon X_1\to X,\quad
    \alpha_V^1\colon X_1\times_XV\xrightarrow{\sim}V.              \tag{INC0a}
  \]
  Under the substitution
  $(U,T_1,T_2,X_C,\alpha_U)=(V,T_c,E_c,X_1,\alpha_V^1)$ in
  (KCS1)--(KCS2) of \Cref{def:nagata-separate-constructible-closed}, define
  \[
  \begin{gathered}
    \mathbf T_c:=V(\sqrt{\mathcal J_V(T_c)}),\qquad
    \iota_{T_c}\colon\mathbf T_c\hookrightarrow V,\\
    \mathbf E_c:=V(\sqrt{\mathcal J_V(E_c)}),\qquad
    \iota_{E_c}\colon\mathbf E_c\hookrightarrow V,\\
    p_1:=\operatorname{pr}_{X_1}\colon X_1\times_XV\to X_1,
    \qquad
    a_V^1:=p_1\circ(\alpha_V^1)^{-1}\colon V\to X_1,\\
    \tau_T:=a_V^1\circ\iota_{T_c}\colon\mathbf T_c\to X_1,
    \qquad
    \tau_E:=a_V^1\circ\iota_{E_c}\colon\mathbf E_c\to X_1,\\
    \operatorname{IsQCQS}(\tau_T),\qquad
    \operatorname{IsQCQS}(\tau_E).
  \end{gathered}                                                    \tag{INC0b}
  \]
  Only after the final line of (INC0b), define
  \[
  \begin{gathered}
    B_T:=\operatorname{SIm}_{X_1}(\tau_T),\qquad
    i_{\tau_T}\colon B_T\hookrightarrow X_1,\\
    B_E:=\operatorname{SIm}_{X_1}(\tau_E),\qquad
    i_{\tau_E}\colon B_E\hookrightarrow X_1.
  \end{gathered}                                                    \tag{INC0c}
  \]
  Define
  \[
    X_0:=X_1\setminus|B_E|,\qquad h_0\colon X_0\hookrightarrow X_1.
  \]
  Put
  \[
  \begin{gathered}
    Z_1:=\operatorname{ST}_{\pi_1}(Z),\qquad
    q_1:=q^{\mathrm{st}}_{\pi_1}\colon Z_1\to X_1,\\
    Z_0:=Z_1\times_{X_1}X_0,\qquad
    q_0:=q_1\times_{X_1}X_0\colon Z_0\to X_0,\\
    U_0:=V\setminus E_c=X_0\times_{X_1}V,\\
    T_{\mathrm{red}}
      :=\operatorname{Spec}_{U_0}
        (\mathcal O_{U_0}/\mathcal I_{U_0,T}),\qquad
    i_T:=i_{U_0,T}\colon T_{\mathrm{red}}\hookrightarrow U_0
      \quad\text{as in
        \Cref{def:nagata-reduced-closed-subscheme}}.
  \end{gathered}
  \]
  Define the open immersion, the restricted source, and the restricted morphism
  by
  \[
  \begin{gathered}
    j_0:=\operatorname{pr}_{X_0}\colon
      U_0=X_0\times_{X_1}V\hookrightarrow X_0,\qquad
    v_0:=\operatorname{pr}_V\colon
      U_0=X_0\times_{X_1}V\hookrightarrow V,\\
    Z_{0,U_0}:=Z_0\times_{X_0}U_0,\\
    q_{0,U_0}:=\operatorname{pr}_{U_0}
      \colon Z_{0,U_0}\longrightarrow U_0,\qquad
    j_0\circ q_{0,U_0}=q_0\circ\operatorname{pr}_{Z_0},\\
    Z_{V,U_0}:=Z_V\times_VU_0,\qquad
    q_{V,U_0}:=\operatorname{pr}_{U_0}
      \colon Z_{V,U_0}\longrightarrow U_0,\qquad
    v_0\circ q_{V,U_0}=q_V\circ\operatorname{pr}_{Z_V}.
  \end{gathered}                                                    \tag{INC0}
  \]
  The defining union in $V^{\mathrm{iso}}(q)$ and the equality
  $U_0\subseteq V^{\mathrm{iso}}(q)$ are the predicate
  $\operatorname{IsIso}(q_{V,U_0})$.  Applying
  \Cref{lemma:nagata-strict-transform-restrict} first to $\pi_1$ over $V$
  and then to $U_0\hookrightarrow V$ is the predicate
  \[
    \operatorname{IsIso}(q_{0,U_0}).                              \tag{INC1}
  \]
  Regard this morphism as an isomorphism and define
  \[
    \eta_0:=q_{0,U_0}\colon Z_{0,U_0}\xrightarrow{\sim}U_0,
    \qquad
    j_0\circ\eta_0=q_0\circ\operatorname{pr}_{Z_0}.              \tag{INC2}
  \]
  Apply \Cref{def:nagata-elimination-indeterminacy} to the fully specified
  datum $(q_0,j_0,\eta_0)$ and define
  \[
  \begin{gathered}
    \mathcal I_0:=\mathcal F_{U_0}(q_0,\eta_0),\\
    X_0^{\mathrm{el}}:=X^{\mathrm{el}}(q_0,U_0)
      =\operatorname{Bl}_{\mathcal I_0}(X_0),\qquad
    \pi_0:=b_{X_0,\mathcal I_0}\colon X_0^{\mathrm{el}}\to X_0,\\
    Z_0^{\mathrm{el}}:=Z^{\mathrm{el}}(q_0,U_0)
      =\operatorname{ST}_{\pi_0}(Z_0),\qquad
    q_0^{\mathrm{el}}:=(q_0)^{\mathrm{st}}_{\pi_0}
      \colon Z_0^{\mathrm{el}}\to X_0^{\mathrm{el}},\\
    \alpha_0\colon X_0^{\mathrm{el}}\times_{X_0}U_0
      \xrightarrow{\sim}U_0
      \quad\text{as in \Cref{def:nagata-elimination-indeterminacy}}.
  \end{gathered}                                                    \tag{INC3}
  \]
  Thus
  \[
    \mathsf D_0:=(\mathcal I_0,X_0^{\mathrm{el}},\pi_0,\alpha_0)
  \]
  is the $U_0=X_0\times_{X_1}V$-admissible datum supplied to
  \Cref{def:nagata-extended-admissible-blowup}.  Define its extension from
  $X_0$ to $X_1$, with preserved open $V$, by
  \[
  \begin{gathered}
    (\mathcal I_2,X_2,\pi_2,\alpha_V^2)
      :=\mathsf{ExtBl}_{X_1,X_0,V}(\mathcal I_0),\\
    \pi_2=b_{X_1,\mathcal I_2}\colon X_2\to X_1,\qquad
    \rho_0:=\rho_{X_0}\colon
      X_2\times_{X_1}X_0\xrightarrow{\sim}X_0^{\mathrm{el}},\\
    \pi_0\circ\rho_0=\operatorname{pr}_{X_0}.
  \end{gathered}                                                    \tag{INC4}
  \]
  Finally define
  \[
    \mathsf{IsoBl}_{V,T}(q)
      :=(\pi:=\pi_1\circ\pi_2\colon X_2\to X),
    \qquad
    q_2:=q^{\mathrm{st}}_\pi\colon
      \operatorname{ST}_\pi(Z)\to X_2.
  \]
\end{definition}

\begin{lemma}[Property of the isomorphism-near-closure construction,
  {\cite[\href{https://stacks.math.columbia.edu/tag/0F3W}{Tag 0F3W}]{stacks-project}}]
  \label{lemma:nagata-isomorphism-near-closure}
  \uses{def:nagata-isomorphism-near-closure}
  \lean{}

  For the data of \Cref{def:nagata-isomorphism-near-closure}, including
  \(T\subseteq V^{\mathrm{iso}}(q)\), the morphism
  \[
    \pi=\pi_1\circ\pi_2\colon X_2\to X
  \]
  is \(V\)-admissible.  Define
  \[
  \begin{gathered}
    \alpha_V^\pi\colon X_2\times_XV\xrightarrow{\sim}V,\qquad
    a_V^\pi:=\operatorname{pr}_{X_2}\circ(\alpha_V^\pi)^{-1}
      \colon V\to X_2,\\
    T_{\mathrm{red}}^V
      :=\operatorname{Spec}_V(\mathcal O_V/\mathcal I_{V,T}),\qquad
    i_{V,T}\colon T_{\mathrm{red}}^V\hookrightarrow V
      \quad\text{as in
        \Cref{def:nagata-reduced-closed-subscheme}},\\
    t_{\pi,T}:=a_V^\pi\circ i_{V,T}
      \colon T_{\mathrm{red}}^V\to X_2.
  \end{gathered}                                                   \tag{IC0a}
  \]
  The complement square for the $V$-admissible morphism $\pi$ identifies
  $a_V^\pi$ with the base change of the quasi-compact open immersion
  $V\hookrightarrow X$.  Stability under base change and composition is the
  predicate chain
  \[
    \operatorname{IsQuasiCompactOpenImmersion}(a_V^\pi)
      \wedge\operatorname{IsClosedImmersion}(i_{V,T})
      \Longrightarrow
    \operatorname{IsQCQS}(a_V^\pi)
      \wedge\operatorname{IsQCQS}(i_{V,T})
      \Longrightarrow
    \operatorname{IsQCQS}(t_{\pi,T}).                             \tag{IC0b}
  \]
  By \Cref{def:nagata-scheme-theoretic-image,def:nagata-scheme-image-factor},
  define
  \[
  \begin{gathered}
    B_{\pi,T}:=\operatorname{SIm}_{X_2}(t_{\pi,T}),\qquad
    i_{\pi,T}:=i_{t_{\pi,T}}\colon B_{\pi,T}\hookrightarrow X_2.
  \end{gathered}                                                   \tag{IC0}
  \]
  There exist an open immersion \(n\colon N\hookrightarrow X_2\), a closed
  immersion \(c_{\pi,T}\colon B_{\pi,T}\hookrightarrow N\), and an
  isomorphism
  \[
    q_{2,N}\colon
      \operatorname{ST}_\pi(Z)\times_{X_2}N\xrightarrow{\sim}N
  \]
  satisfying
  \[
    n\circ c_{\pi,T}=i_{\pi,T},\qquad
    n\circ q_{2,N}
      =q_2\circ\operatorname{pr}_{\operatorname{ST}_\pi(Z)}.       \tag{IC1}
  \]
\end{lemma}

\begin{proof}
  \uses{def:nagata-choice-notation,def:nagata-closed-subscheme-class,
    lemma:nagata-spectral-constructible-separation,
    lemma:nagata-separate-constructible-closed,
    lemma:nagata-scheme-image-universal,
    lemma:nagata-scheme-image-restrict-open,
    lemma:nagata-elimination-indeterminacy,
    lemma:nagata-proper-open-dense-iso,
    lemma:nagata-extend-admissible-blowup,
    lemma:nagata-admissible-blowup-comp}
  From \Cref{lemma:nagata-spectral-constructible-separation},
  \[
    T\subseteq T_c,\qquad E(q,V)\subseteq E_c,\qquad
    T_c\cap E_c=\varnothing.
  \]
  By \Cref{lemma:nagata-separate-constructible-closed}, applied to the
  objects $B_T,B_E$ defined by (INC0c),
  \[
    B_T\times_{X_1}B_E=\varnothing.
  \]
  Define the closed immersion $b_T$ by
  \[
    b_T\colon B_T\hookrightarrow X_0,
    \qquad
    h_0\circ b_T=i_{\tau_T},
  \]
  where $i_{\tau_T}\colon B_T\hookrightarrow X_1$ is the closed immersion of
  \Cref{def:nagata-scheme-image-factor}.  By the equality
  $U_0=X_0\times_{X_1}V$ in
  \Cref{def:nagata-isomorphism-near-closure}, define
  \[
  \begin{gathered}
    u_0:=j_0=\operatorname{pr}_{X_0}\colon
      U_0=X_0\times_{X_1}V\hookrightarrow X_0,\\
    t_1:=h_0\circ u_0\circ i_T\colon T_{\mathrm{red}}\to X_1.
  \end{gathered}
  \]
  The inclusion $T\subseteq T_c$ and
  \Cref{def:nagata-separate-constructible-closed} define
  \[
    a_T\colon T_{\mathrm{red}}\hookrightarrow
      \mathbf T_c=V\!\left(\sqrt{\mathcal J_V(T_c)}\right),
    \qquad
    \tau_T\circ a_T=t_1.
  \]
  The final line of (INC0b), the closed-immersion implication, and (PC3) of
  \Cref{lemma:nagata-morphism-property-calculus} are the following predicate
  chain:
  \[
    \operatorname{IsClosedImmersion}(a_T)
      \Longrightarrow\operatorname{IsQCQS}(a_T),
    \qquad
    \operatorname{IsQCQS}(\tau_T)\ \wedge\
      \operatorname{IsQCQS}(a_T)\ \wedge\
      (t_1=\tau_T\circ a_T)
      \Longrightarrow\operatorname{IsQCQS}(t_1).                  \tag{IC0c}
  \]
  Only after (IC0c), define
  \[
  \begin{gathered}
    B_T^{(0)}:=\operatorname{SIm}_{X_1}(t_1),
    \qquad
    i_T^{(0)}:=i_{t_1}\colon B_T^{(0)}\hookrightarrow X_1.
  \end{gathered}
  \]
  By \Cref{def:nagata-scheme-image-factor}, define
  \[
    t_1^{\mathrm{im}}\colon T_{\mathrm{red}}\longrightarrow B_T^{(0)},
    \qquad
    i_T^{(0)}\circ t_1^{\mathrm{im}}=t_1.                         \tag{IC0d}
  \]
  Substitution of
  $\tau_T=i_{\tau_T}\circ\tau_T^{\mathrm{im}}$ from
  \Cref{def:nagata-scheme-image-factor} is the equality
  \[
    t_1=i_{\tau_T}\circ\tau_T^{\mathrm{im}}\circ a_T.
  \]
  Hence \Cref{lemma:nagata-scheme-image-universal} defines the closed
  immersions
  \[
  \begin{gathered}
    c_T^{(0)}:=c_{t_1,i_{\tau_T}}\colon B_T^{(0)}\hookrightarrow B_T,
    \qquad
    i_{\tau_T}\circ c_T^{(0)}=i_T^{(0)},\\
    b_T^{(0)}:=b_T\circ c_T^{(0)}\colon B_T^{(0)}\hookrightarrow X_0,
    \qquad
    h_0\circ b_T^{(0)}=i_T^{(0)}.
  \end{gathered}
  \]

  Because $E(q,V)\subseteq E_c$,
  \[
    U_0=V\setminus E_c
      \subseteq V\setminus E(q,V)=V^{\mathrm{iso}}(q).
  \]
  By (INC0)--(INC2), the restriction and the previously defined isomorphism
  are exactly
  \[
    Z_{0,U_0}=q_0^{-1}(U_0),
    \qquad
    \eta_0=q_{0,U_0}\colon Z_{0,U_0}\xrightarrow{\sim}U_0,
    \qquad
    \operatorname{IsIso}(\eta_0).
  \]
  By \Cref{lemma:nagata-elimination-indeterminacy},
  \[
    q_0^{\mathrm{el}}\colon Z_0^{\mathrm{el}}
      \hookrightarrow X_0^{\mathrm{el}}
  \]
  is an open immersion and $\pi_0$ is $U_0$-admissible.  The morphism $q_0$
  is proper.  Define
  \[
  \begin{gathered}
    \iota_0\colon Z_0\hookrightarrow Z\times_XX_0,\qquad
    p_0\colon Z\times_XX_0\to X_0,\qquad
    q_0=p_0\circ\iota_0,\\
    \iota_0^{\mathrm{el}}\colon Z_0^{\mathrm{el}}\hookrightarrow
      Z_0\times_{X_0}X_0^{\mathrm{el}},\qquad
    p_0^{\mathrm{el}}\colon
      Z_0\times_{X_0}X_0^{\mathrm{el}}\to X_0^{\mathrm{el}},\\
    q_0^{\mathrm{el}}
      =p_0^{\mathrm{el}}\circ\iota_0^{\mathrm{el}}.
  \end{gathered}                                                   \tag{IC1a}
  \]
  Here \(\iota_0,\iota_0^{\mathrm{el}}\) are the closed immersions of
  \Cref{def:nagata-strict-transform}, \(p_0\) is the base change of \(q\),
  and \(p_0^{\mathrm{el}}\) is the base change of \(q_0\).  The predicate
  chains are
  \[
  \begin{aligned}
    \operatorname{IsClosedImmersion}(\iota_0)
      \wedge\operatorname{IsProper}(p_0)
      &\Longrightarrow\operatorname{IsProper}(q_0),\\
    \operatorname{IsClosedImmersion}(\iota_0^{\mathrm{el}})
      \wedge\operatorname{IsProper}(p_0^{\mathrm{el}})
      &\Longrightarrow\operatorname{IsProper}(q_0^{\mathrm{el}}).
  \end{aligned}                                                    \tag{IC1b}
  \]
  The open-immersion assertion above, (IC1b), (P0a), and the
  proper-immersion criterion
  \cite[\href{https://stacks.math.columbia.edu/tag/01IQ}{Tag 01IQ}]
  {stacks-project} form the predicate chain
  \[
  \begin{aligned}
    &\operatorname{IsOpenImmersion}(q_0^{\mathrm{el}})
      \Longrightarrow\operatorname{IsImmersion}(q_0^{\mathrm{el}}),\\
    &\operatorname{IsProper}(q_0^{\mathrm{el}})
      \wedge\operatorname{IsImmersion}(q_0^{\mathrm{el}})
      \Longrightarrow\operatorname{IsClosedImmersion}(q_0^{\mathrm{el}}).
  \end{aligned}                                                   \tag{IC1c}
  \]
  The last implication in (IC1c) and the standard closed-immersion implication
  are the following implication:
  \[
    \operatorname{IsClosedImmersion}(q_0^{\mathrm{el}})
      \Longrightarrow\operatorname{IsQCQS}(q_0^{\mathrm{el}}).     \tag{IC1d}
  \]
  Only after (IC1d), apply
  \Cref{def:nagata-scheme-theoretic-image,def:nagata-scheme-image-factor}
  and define the scheme-theoretic image, closed immersion, and image factor by
  \[
  \begin{gathered}
    N_0:=\operatorname{SIm}_{X_0^{\mathrm{el}}}(q_0^{\mathrm{el}}),\qquad
    i_{N_0}:=i_{q_0^{\mathrm{el}}}\colon
      N_0\hookrightarrow X_0^{\mathrm{el}},\\
    (q_0^{\mathrm{el}})^{\mathrm{im}}\colon
      Z_0^{\mathrm{el}}\longrightarrow N_0,\qquad
    i_{N_0}\circ(q_0^{\mathrm{el}})^{\mathrm{im}}=q_0^{\mathrm{el}}.
  \end{gathered}
  \]
  Since \(q_0^{\mathrm{el}}\) is both an open and a closed immersion,
  \[
    (q_0^{\mathrm{el}})^{\mathrm{im}}\colon
      Z_0^{\mathrm{el}}\xrightarrow{\sim}N_0,\qquad
    i_{N_0}\text{ is both an open and a closed immersion}.        \tag{IC2}
  \]
  Let
  \[
    \alpha_0\colon X_0^{\mathrm{el}}\times_{X_0}U_0
      \xrightarrow{\sim}U_0
  \]
  be the complement isomorphism in
  \Cref{def:nagata-elimination-indeterminacy}, and define
  \[
    s_0:=\operatorname{pr}_{X_0^{\mathrm{el}}}\circ\alpha_0^{-1}
      \colon U_0\to X_0^{\mathrm{el}}.
  \]
  The two fibre-product projection equalities are
  \[
  \begin{aligned}
    \pi_0\circ s_0
      &=\pi_0\circ\operatorname{pr}_{X_0^{\mathrm{el}}}
          \circ\alpha_0^{-1}\\
      &=u_0\circ\operatorname{pr}_{U_0}\circ\alpha_0^{-1}
       =u_0.
  \end{aligned}
  \]
  The restriction equation in
  \Cref{lemma:nagata-elimination-indeterminacy} defines
  \[
    \bar s_0\colon U_0\to N_0,\qquad
    i_{N_0}\circ\bar s_0=s_0,
  \]
  where $i_{N_0}\colon N_0\hookrightarrow X_0^{\mathrm{el}}$ is the closed
  immersion.  Define
  \[
    \bar t_0:=\bar s_0\circ i_T\colon T_{\mathrm{red}}\to N_0,
    \qquad
    t_0:=s_0\circ i_T=i_{N_0}\circ\bar t_0.
  \]
  The complement square for the $U_0$-admissible morphism $\pi_0$, together
  with (PC2)--(PC3) of
  \Cref{lemma:nagata-morphism-property-calculus}, is the following predicate
  chain:
  \[
  \begin{gathered}
    \operatorname{IsQuasiCompactOpenImmersion}(j_0)
      \Longrightarrow
    \operatorname{IsQuasiCompactOpenImmersion}(s_0)
      \Longrightarrow\operatorname{IsQCQS}(s_0),\\
    \operatorname{IsClosedImmersion}(i_T)
      \Longrightarrow\operatorname{IsQCQS}(i_T),\\
    \operatorname{IsQCQS}(s_0)\ \wedge\
      \operatorname{IsQCQS}(i_T)\ \wedge\
      (t_0=s_0\circ i_T)
      \Longrightarrow\operatorname{IsQCQS}(t_0).
  \end{gathered}                                                    \tag{IC2a}
  \]
  Only after (IC2a), the universal property in
  \Cref{lemma:nagata-scheme-image-universal} defines
  \[
    c_{T,N_0}\colon\operatorname{SIm}(t_0)\hookrightarrow N_0,
    \qquad
    i_{N_0}\circ c_{T,N_0}=i_{t_0},
  \]
  where $i_{t_0}\colon\operatorname{SIm}(t_0)\hookrightarrow
  X_0^{\mathrm{el}}$ is the closed immersion of
  \Cref{def:nagata-scheme-theoretic-image}.
  Equation (IC2), the definition of $c_{T,N_0}$, and the displayed
  factorization are the predicate-and-image chain
  \[
  \begin{gathered}
    \operatorname{IsOpenImmersion}(i_{N_0}),\qquad
    \operatorname{IsClosedImmersion}(c_{T,N_0}),\qquad
    i_{N_0}\circ c_{T,N_0}=i_{t_0},\\
    |i_{t_0}|\bigl(|\operatorname{SIm}(t_0)|\bigr)
      =|i_{N_0}|\bigl(
        |c_{T,N_0}|\bigl(|\operatorname{SIm}(t_0)|\bigr)\bigr)
      \subseteq |i_{N_0}|(|N_0|).
  \end{gathered}
  \]

  By \Cref{lemma:nagata-extend-admissible-blowup},
  \[
    \pi_2\text{ is }V\text{-admissible},\qquad
    \rho_0\colon X_2\times_{X_1}X_0\xrightarrow{\sim}X_0^{\mathrm{el}}.
  \]
  By \Cref{lemma:nagata-admissible-blowup-comp},
  \[
    \pi=\pi_1\circ\pi_2\text{ is }V\text{-admissible}.
  \]
  Define
  \[
  \begin{gathered}
    P_0:=X_2\times_{X_1}X_0,
    \qquad
    r_2\colon P_0\to X_2,
    \qquad
    r_0\colon P_0\to X_0,\\
    \pi_2\circ r_2=h_0\circ r_0,
    \qquad
    \pi_0\circ\rho_0=r_0.
  \end{gathered}
  \]
  Here the second equality is the structural equality of
  \Cref{lemma:nagata-extend-admissible-blowup}.  Define the fibre product
  \[
  \begin{gathered}
    N:=P_0\times_{X_0^{\mathrm{el}},\rho_0,i_{N_0}}N_0
       =\rho_0^{-1}(N_0),\\
    n_0\colon N\hookrightarrow P_0,
    \qquad
    \rho_N\colon N\xrightarrow{\sim}N_0,
    \qquad
    \rho_0\circ n_0=i_{N_0}\circ\rho_N,\\
    n_2:=r_2\circ n_0\colon N\hookrightarrow X_2.
  \end{gathered}
  \]
  The morphism $r_2$ is the base change of the open immersion $h_0$ along
  $\pi_2$, and $n_0$ is the base change of the open immersion $i_{N_0}$
  along $\rho_0$.  Base-change stability and composition stability in
  \Cref{lemma:nagata-morphism-property-calculus} are the following predicate
  chain:
  \[
  \begin{gathered}
    \operatorname{IsOpenImmersion}(h_0)
      \Longrightarrow\operatorname{IsOpenImmersion}(r_2),\\
    \operatorname{IsOpenImmersion}(i_{N_0})
      \Longrightarrow\operatorname{IsOpenImmersion}(n_0),\\
    \operatorname{IsOpenImmersion}(r_2)\ \wedge\
      \operatorname{IsOpenImmersion}(n_0)
      \Longrightarrow
      \operatorname{IsOpenImmersion}(r_2\circ n_0)
      =\operatorname{IsOpenImmersion}(n_2).                     
  \end{gathered} \tag{IC3a}
  \]
  Define
  \[
  \begin{gathered}
    \widetilde t_2:=\rho_0^{-1}\circ t_0
      \colon T_{\mathrm{red}}\to P_0,
    \qquad
    t_2:=r_2\circ\widetilde t_2
      \colon T_{\mathrm{red}}\to X_2,\\
    \rho_0\circ\widetilde t_2=t_0.
  \end{gathered}
  \]
  Equation (IC2a), the isomorphism $\rho_0$ in (INC4), and (PC3) of
  \Cref{lemma:nagata-morphism-property-calculus} are the following predicate
  chain:
  \[
    \operatorname{IsIso}(\rho_0^{-1})
      \Longrightarrow\operatorname{IsQCQS}(\rho_0^{-1}),
    \qquad
    \operatorname{IsQCQS}(\rho_0^{-1})\ \wedge\
      \operatorname{IsQCQS}(t_0)\ \wedge\
      (\widetilde t_2=\rho_0^{-1}\circ t_0)
      \Longrightarrow\operatorname{IsQCQS}(\widetilde t_2).       \tag{IC3b}
  \]
  Successive substitution of the structural equalities above is the chain
  \[
  \begin{aligned}
    \pi_2\circ t_2
      &=\pi_2\circ r_2\circ\widetilde t_2
       =h_0\circ r_0\circ\widetilde t_2\\
      &=h_0\circ\pi_0\circ\rho_0\circ\widetilde t_2
       =h_0\circ\pi_0\circ t_0\\
      &=h_0\circ\pi_0\circ s_0\circ i_T
       =h_0\circ u_0\circ i_T=t_1.
  \end{aligned}
  \]
  Retain the open immersion \(v_0\colon U_0\hookrightarrow V\) from (INC0).  By
  \Cref{def:nagata-reduced-closed-subscheme}, define the unique
  isomorphism
  \[
    \chi_T\colon T_{\mathrm{red}}\xrightarrow{\sim}T_{\mathrm{red}}^V,
    \qquad
    i_{V,T}\circ\chi_T=v_0\circ i_T.                              \tag{IC3}
  \]
  The complement-isomorphism equations for
  \(\alpha_V^\pi\), applied to the equality
  \(\pi\circ t_2=v\circ v_0\circ i_T\), are
  \[
    t_2
      =a_V^\pi\circ v_0\circ i_T
      =a_V^\pi\circ i_{V,T}\circ\chi_T
      =t_{\pi,T}\circ\chi_T.                                     \tag{IC4}
  \]
  Equations (IC0b) and (IC4), together with (PC3) of
  \Cref{lemma:nagata-morphism-property-calculus}, are the following predicate
  chain:
  \[
    \operatorname{IsIso}(\chi_T)
      \Longrightarrow\operatorname{IsQCQS}(\chi_T),
    \qquad
    \operatorname{IsQCQS}(t_{\pi,T})\ \wedge\
      \operatorname{IsQCQS}(\chi_T)\ \wedge\
      (t_2=t_{\pi,T}\circ\chi_T)
      \Longrightarrow\operatorname{IsQCQS}(t_2).                 \tag{IC4a}
  \]
  Precomposition by the isomorphism \(\chi_T\) in (IC4) is the following
  class equality:
  \[
    [\,\operatorname{SIm}_{X_2}(t_2)\xhookrightarrow{i_{t_2}}X_2\,]
      =[\,B_{\pi,T}\xhookrightarrow{i_{\pi,T}}X_2\,]
      \quad\text{in }\operatorname{ClosedSub}(X_2).              \tag{IC5}
  \]
  Define
  \[
    \mathscr B_{2,T}:=
      \left\{\beta\in
        \operatorname{Hom}(\operatorname{SIm}_{X_2}(t_2),B_{\pi,T})
        \ \middle|\
        \operatorname{IsIso}(\beta)\ \wedge\
        i_{\pi,T}\circ\beta=i_{t_2}\right\}.
  \]
  By (CSub3) of \Cref{def:nagata-closed-subscheme-class}, (IC5) is
  $\mathscr B_{2,T}\ne\varnothing$.  Define
  \[
    \beta_{2,T}:=\operatorname{ch}(\mathscr B_{2,T})
      \in\mathscr B_{2,T},
    \qquad
    \beta_{2,T}\colon\operatorname{SIm}_{X_2}(t_2)
      \xrightarrow{\sim}B_{\pi,T},
    \qquad
    i_{\pi,T}\circ\beta_{2,T}=i_{t_2}.                         \tag{IC5a}
  \]
  Define
  \[
  \begin{gathered}
    B_T^{(2)}:=B_{\pi,T},
    \qquad
    i_T^{(2)}:=i_{\pi,T}\colon B_T^{(2)}\hookrightarrow X_2,\\
    t_2^{\mathrm{im}}:=\beta_{2,T}\circ(t_2)^{\mathrm{im}}
      \colon T_{\mathrm{red}}\to B_T^{(2)},
    \qquad
    i_T^{(2)}\circ t_2^{\mathrm{im}}
      =i_{\pi,T}\circ\beta_{2,T}\circ(t_2)^{\mathrm{im}}
      =i_{t_2}\circ(t_2)^{\mathrm{im}}=t_2.
  \end{gathered}
  \]
  Let
  \[
  \begin{gathered}
    P_T:=X_2\times_{X_1,i_T^{(0)}}B_T^{(0)},
    \qquad
    p_2\colon P_T\hookrightarrow X_2,
    \qquad
    p_T\colon P_T\to B_T^{(0)},\\
    \widehat t_2:=(t_2,t_1^{\mathrm{im}})
      \colon T_{\mathrm{red}}\to P_T,
    \qquad
    p_2\circ\widehat t_2=t_2.
  \end{gathered}
  \]
  The map $p_2$ is the base change of the closed immersion $i_T^{(0)}$.
  Thus \Cref{lemma:nagata-scheme-image-universal} defines
  \[
    c_{2,T}\colon B_T^{(2)}\hookrightarrow P_T,
    \qquad
    p_2\circ c_{2,T}=i_T^{(2)}.
  \]
  The equality $h_0\circ b_T^{(0)}=i_T^{(0)}$ defines the base-changed
  closed immersion
  \[
  \begin{gathered}
    e_T\colon P_T\hookrightarrow P_0,
    \qquad
    r_2\circ e_T=p_2,
    \qquad
    r_0\circ e_T=b_T^{(0)}\circ p_T,\\
    b_2:=e_T\circ c_{2,T}\colon B_T^{(2)}\hookrightarrow P_0,
    \qquad
    r_2\circ b_2=i_T^{(2)}.
  \end{gathered}
  \]
  Put
  \[
    Z_{2,N}:=\operatorname{ST}_\pi(Z)\times_{X_2}N,\qquad
    Z^{\mathrm{el}}_{0,N_0}:=
      Z_0^{\mathrm{el}}\times_{X_0^{\mathrm{el}}}N_0.
  \]
  The isomorphism of
  \Cref{lemma:nagata-strict-transform-dominating}, base-changed to $N$ and
    transported by $\rho_N=\rho_0|_N\colon N\xrightarrow{\sim}N_0$, is the typed
  isomorphism
  \[
    \delta_N\colon Z_{2,N}\xrightarrow{\sim}Z^{\mathrm{el}}_{0,N_0},
    \qquad
    q^{\mathrm{el}}_{0,N_0}\circ\delta_N
      =\rho_N\circ q_{2,N},
  \]
  where $q_{2,N}\colon Z_{2,N}\to N$ and
  $q^{\mathrm{el}}_{0,N_0}\colon Z^{\mathrm{el}}_{0,N_0}\to N_0$ are the
  base changes of $q_2$ and $q_0^{\mathrm{el}}$.  Since
  $q^{\mathrm{el}}_{0,N_0}$, $\delta_N$, and $\rho_N$ are isomorphisms,
  \[
    q_{2,N}=\rho_N^{-1}\circ
      q^{\mathrm{el}}_{0,N_0}\circ\delta_N
  \]
  is an isomorphism.
  Let
  \[
    C_0:=\operatorname{SIm}_{P_0}(\widetilde t_2),
    \qquad
    i_{\widetilde t_2}\colon C_0\hookrightarrow P_0.
  \]
  By \Cref{lemma:nagata-scheme-image-restrict-open}, applied to
  $r_2\colon P_0\hookrightarrow X_2$, the factorization
  $t_2=r_2\circ\widetilde t_2$ and the equality
  $r_2\circ b_2=i_T^{(2)}$ define the unique isomorphism
  \[
    \varepsilon_0\colon C_0\xrightarrow{\sim}B_T^{(2)},
    \qquad
    b_2\circ\varepsilon_0=i_{\widetilde t_2}.
  \]
  Since $t_0=\rho_0\circ\widetilde t_2$, two applications of
  \Cref{lemma:nagata-scheme-image-universal} define closed immersions
  \[
  \begin{gathered}
    d_0\colon\operatorname{SIm}(t_0)\hookrightarrow C_0,
    \qquad
    (\rho_0\circ i_{\widetilde t_2})\circ d_0=i_{t_0},\\
    e_0\colon C_0\hookrightarrow\operatorname{SIm}(t_0),
    \qquad
    i_{t_0}\circ e_0=\rho_0\circ i_{\widetilde t_2}.
  \end{gathered}
  \]
  Cancellation of the closed immersions $i_{t_0}$ and
  $\rho_0\circ i_{\widetilde t_2}$ in
  \[
  \begin{aligned}
    i_{t_0}\circ e_0\circ d_0
      &=\rho_0\circ i_{\widetilde t_2}\circ d_0=i_{t_0},\\
    \rho_0\circ i_{\widetilde t_2}\circ d_0\circ e_0
      &=i_{t_0}\circ e_0=\rho_0\circ i_{\widetilde t_2}
  \end{aligned}
  \]
  is the pair of equalities
  \[
    e_0\circ d_0=\operatorname{id}_{\operatorname{SIm}(t_0)},
    \qquad
    d_0\circ e_0=\operatorname{id}_{C_0}.
  \]
  Hence the comparison isomorphism and its comparison to $N_0$ are
  \[
  \begin{gathered}
    \lambda_0:=e_0\circ\varepsilon_0^{-1}
      \colon B_T^{(2)}\xrightarrow{\sim}\operatorname{SIm}(t_0),
    \qquad
    i_{t_0}\circ\lambda_0=\rho_0\circ b_2,\\
    c_{2,N_0}:=c_{T,N_0}\circ\lambda_0
      \colon B_T^{(2)}\hookrightarrow N_0,\\
    i_{N_0}\circ c_{2,N_0}
      =i_{N_0}\circ c_{T,N_0}\circ\lambda_0
      =i_{t_0}\circ\lambda_0
      =\rho_0\circ b_2.
  \end{gathered}
  \]
  The universal property of
  $N=P_0\times_{X_0^{\mathrm{el}}}N_0$ therefore defines the closed immersion
  \[
  \begin{gathered}
    c_{2,N}:=(b_2,c_{2,N_0})\colon B_T^{(2)}\hookrightarrow N,\\
    n_0\circ c_{2,N}=b_2,
    \qquad
    \rho_N\circ c_{2,N}=c_{2,N_0},\\
    n_2\circ c_{2,N}
      =r_2\circ n_0\circ c_{2,N}
      =r_2\circ b_2=i_T^{(2)}.
  \end{gathered}
  \]
  \[
  \begin{gathered}
    n:=n_2\colon N\hookrightarrow X_2,\qquad
    c_{\pi,T}:=c_{2,N}\colon B_{\pi,T}\hookrightarrow N,\\
    n\circ c_{\pi,T}
      =n_2\circ c_{2,N}=i_T^{(2)}=i_{\pi,T},\\
    q_{2,N}=\rho_N^{-1}\circ q^{\mathrm{el}}_{0,N_0}\circ\delta_N
      \colon\operatorname{ST}_\pi(Z)\times_{X_2}N\xrightarrow{\sim}N,\\
    n\circ q_{2,N}
      =q_2\circ\operatorname{pr}_{\operatorname{ST}_\pi(Z)}.
  \end{gathered}                                                   \tag{IC6}
  \]
\end{proof}


\section{Common modifications and graph extension}

\begin{lemma}[Scheme-theoretic image of a partial section]
  \label{lemma:nagata-scheme-image-partial-section}
  \lean{}

  Let $x\colon X\to S$, let $y\colon Y\to S$ be separated, let
  $j\colon U\hookrightarrow X$ be a quasi-compact open immersion, and let
  \[
    \gamma\colon U\longrightarrow X\times_SY,\qquad
    \operatorname{pr}_X\circ\gamma=j.
  \]
  Define
  \[
  \begin{gathered}
    s_\gamma:=(\operatorname{id}_U,\operatorname{pr}_Y\circ\gamma)
      \colon U\to U\times_SY,\\
    \jmath_{U,Y}:=j\times_S\operatorname{id}_Y
      \colon U\times_SY\hookrightarrow X\times_SY,\\
    \operatorname{pr}_U\circ s_\gamma=\operatorname{id}_U,\qquad
    \gamma=\jmath_{U,Y}\circ s_\gamma.
  \end{gathered}                                                    \tag{PS0}
  \]
  Base change of $y$ and $j$, the closed-section criterion for a separated
  morphism, and (PC2)--(PC3) of
  \Cref{lemma:nagata-morphism-property-calculus} are the following predicate
  chain:
  \[
  \begin{gathered}
    \operatorname{IsSeparated}(y)
      \Longrightarrow\operatorname{IsSeparated}(\operatorname{pr}_U),\\
    \operatorname{IsSeparated}(\operatorname{pr}_U)\ \wedge\
      (\operatorname{pr}_U\circ s_\gamma=\operatorname{id}_U)
      \Longrightarrow\operatorname{IsClosedImmersion}(s_\gamma)
      \Longrightarrow\operatorname{IsQCQS}(s_\gamma),\\
    \operatorname{IsQuasiCompactOpenImmersion}(j)
      \Longrightarrow
    \operatorname{IsQuasiCompactOpenImmersion}(\jmath_{U,Y})
      \Longrightarrow\operatorname{IsQCQS}(\jmath_{U,Y}),\\
    \operatorname{IsQCQS}(s_\gamma)\ \wedge\
      \operatorname{IsQCQS}(\jmath_{U,Y})\ \wedge\
      (\gamma=\jmath_{U,Y}\circ s_\gamma)
      \Longrightarrow\operatorname{IsQCQS}(\gamma).
  \end{gathered}                                                    \tag{PS1}
  \]
  Only after (PS1), apply
  \Cref{def:nagata-scheme-theoretic-image,def:nagata-scheme-image-factor}
  to $\gamma$ and define
  \[
  \begin{gathered}
    Z:=\operatorname{SIm}(\gamma),\qquad
    i_Z:=i_\gamma\colon Z\hookrightarrow X\times_SY,\\
    \gamma^{\mathrm{im}}\colon U\to Z,\qquad
    i_Z\circ\gamma^{\mathrm{im}}=\gamma,\\
    p:=\operatorname{pr}_X\circ i_Z\colon Z\to X,\\
    Z_U:=Z\times_{X,p,j}U,\qquad
    p_U:=\operatorname{pr}_U\colon Z_U\to U.
  \end{gathered}
  \]
  There is a unique isomorphism
  \[
    \zeta_\gamma\colon U\xrightarrow{\sim}Z_U
  \]
  satisfying
  \[
    \operatorname{pr}_Z\circ\zeta_\gamma=\gamma^{\mathrm{im}},
    \qquad
    p_U\circ\zeta_\gamma=\operatorname{id}_U.
  \]
\end{lemma}

\begin{proof}
  \uses{lemma:nagata-morphism-property-calculus,
    lemma:nagata-scheme-image-universal,
    lemma:nagata-scheme-image-restrict-open}
  Put $W_\gamma:=U\times_SY$.  Base-change stability of separated
  morphisms and the closed-section criterion are the following predicate
  chain:
  \[
  \begin{gathered}
    \operatorname{IsSeparated}(y)
      \Longrightarrow\operatorname{IsSeparated}
        (\operatorname{pr}_U\colon W_\gamma\to U),\\
    \operatorname{pr}_U\circ s_\gamma=\operatorname{id}_U
      \ \wedge\
      \operatorname{IsSeparated}(\operatorname{pr}_U)
      \Longrightarrow\operatorname{IsClosedImmersion}(s_\gamma). 
  \end{gathered} \tag{PS2}
  \]
  Define the restricted source of $\gamma$ along $\jmath_{U,Y}$ by
  \[
  \begin{gathered}
    T_\gamma:=
      U\times_{X\times_SY,(\gamma,\jmath_{U,Y})}W_\gamma,\\
    a_U\colon T_\gamma\to U,\qquad
    a_W\colon T_\gamma\to W_\gamma,\qquad
    \gamma\circ a_U=\jmath_{U,Y}\circ a_W.                    
  \end{gathered} \tag{PS3}
  \]
  Equation (PS0) defines the fibre-product morphism
  \[
    \theta_\gamma:=(\operatorname{id}_U,s_\gamma)
      \colon U\longrightarrow T_\gamma,\qquad
    a_U\circ\theta_\gamma=\operatorname{id}_U,\qquad
    a_W\circ\theta_\gamma=s_\gamma.                            \tag{PS4}
  \]
  Since the open immersion $\jmath_{U,Y}$ is a monomorphism, (PS0) and
  (PS3) are the following cancellation chain:
  \[
  \begin{aligned}
    \jmath_{U,Y}\circ a_W
      &=\gamma\circ a_U
       =\jmath_{U,Y}\circ s_\gamma\circ a_U\\
    &\Longrightarrow a_W=s_\gamma\circ a_U.
  \end{aligned}                                                  \tag{PS5}
  \]
  The two fibre-product components of $\theta_\gamma\circ a_U$ are therefore
  \[
    a_U\circ\theta_\gamma\circ a_U=a_U,
    \qquad
    a_W\circ\theta_\gamma\circ a_U
      =s_\gamma\circ a_U=a_W.
  \]
  Fibre-product uniqueness and (PS4) are the following inverse equalities:
  \[
    \theta_\gamma\circ a_U=\operatorname{id}_{T_\gamma},\qquad
    \theta_\gamma^{-1}=a_U.                                    \tag{PS6}
  \]

  Define the restricted morphism by
  \[
    g_\gamma:=a_W=s_\gamma\circ a_U\colon T_\gamma\to W_\gamma. \tag{PS6a}
  \]
  Equations (PS2), (PS6), and (PS6a), together with (PC3) of
  \Cref{lemma:nagata-morphism-property-calculus}, are the following predicate
  chain:
  \[
  \begin{gathered}
    \operatorname{IsIso}(a_U)
      \Longrightarrow\operatorname{IsClosedImmersion}(a_U),\\
    \operatorname{IsClosedImmersion}(s_\gamma)\ \wedge\
      \operatorname{IsClosedImmersion}(a_U)\ \wedge\
      (g_\gamma=s_\gamma\circ a_U)
      \Longrightarrow\operatorname{IsClosedImmersion}(g_\gamma)
      \Longrightarrow\operatorname{IsQCQS}(g_\gamma).
  \end{gathered}                                                    \tag{PS6b}
  \]
  Only after (PS6b), apply
  \Cref{def:nagata-scheme-theoretic-image,def:nagata-scheme-image-factor}
  to $g_\gamma$ and define all scheme-image symbols by
  \[
  \begin{gathered}
    M_\gamma:=\operatorname{SIm}_{W_\gamma}(g_\gamma),\qquad
    i_{g_\gamma}\colon M_\gamma\hookrightarrow W_\gamma,\qquad
    g_\gamma^{\mathrm{im}}\colon T_\gamma\to M_\gamma,\\
    i_{g_\gamma}\circ g_\gamma^{\mathrm{im}}=g_\gamma.       
  \end{gathered} \tag{PS7}
  \]
  Apply \Cref{lemma:nagata-scheme-image-universal} to the closed factorization
  \[
    g_\gamma=g_\gamma\circ\operatorname{id}_{T_\gamma}.
  \]
  With the notation of
  \Cref{def:nagata-scheme-image-comparison}, this defines
  \[
    \epsilon_\gamma:=c_{g_\gamma,g_\gamma}\colon
      M_\gamma\to T_\gamma,\qquad
    g_\gamma\circ\epsilon_\gamma=i_{g_\gamma},\qquad
    \epsilon_\gamma\circ g_\gamma^{\mathrm{im}}
      =\operatorname{id}_{T_\gamma}.                            \tag{PS8}
  \]
  Since the closed immersion $i_{g_\gamma}$ is a monomorphism,
  \[
  \begin{aligned}
    i_{g_\gamma}\circ g_\gamma^{\mathrm{im}}
        \circ\epsilon_\gamma
      &=g_\gamma\circ\epsilon_\gamma=i_{g_\gamma}\\
    &\Longrightarrow
      g_\gamma^{\mathrm{im}}\circ\epsilon_\gamma
        =\operatorname{id}_{M_\gamma}.
  \end{aligned}                                                  \tag{PS9}
  \]
  Thus $\epsilon_\gamma$ is an isomorphism with inverse
  $g_\gamma^{\mathrm{im}}$.

  Define the target-side fibre product occurring in
  \Cref{lemma:nagata-scheme-image-restrict-open} by
  \[
  \begin{gathered}
    P_\gamma:=Z\times_{X\times_SY,(i_Z,\jmath_{U,Y})}W_\gamma,\\
    \rho_Z\colon P_\gamma\to Z,\qquad
    \rho_W\colon P_\gamma\to W_\gamma,\qquad
    i_Z\circ\rho_Z=\jmath_{U,Y}\circ\rho_W.                  
  \end{gathered} \tag{PS10}
  \]
  Its instance at $(g,w)=(\gamma,\jmath_{U,Y})$ is the isomorphism
  \[
  \begin{gathered}
    \delta_\gamma:=\delta_{\gamma,\jmath_{U,Y}}
      \colon M_\gamma\xrightarrow{\sim}P_\gamma,\\
    \rho_W\circ\delta_\gamma=i_{g_\gamma},\qquad
    \rho_Z\circ\delta_\gamma\circ g_\gamma^{\mathrm{im}}
      =\gamma^{\mathrm{im}}\circ a_U.                           
  \end{gathered} \tag{PS11}
  \]

  The fibre-product equations for $Z\times_XU$ and $X\times_SY$ form the
  chain
  \[
  \begin{aligned}
    x\circ j\circ p_U
      &=x\circ p\circ\operatorname{pr}_Z
       =y\circ\operatorname{pr}_Y\circ i_Z\circ\operatorname{pr}_Z.
  \end{aligned}                                                  \tag{PS12a}
  \]
  Thus the first ordered pair below is a morphism to
  $W_\gamma=U\times_SY$.  Define
  \[
  \begin{gathered}
    w_\gamma^Z:=
      (p_U,\operatorname{pr}_Y\circ i_Z\circ\operatorname{pr}_Z)
      \colon Z_U\to W_\gamma,\\
    \beta_\gamma:=(\rho_Z,\operatorname{pr}_U\circ\rho_W)
      \colon P_\gamma\to Z_U,\\
    \beta_\gamma^{-1}:=(\operatorname{pr}_Z,w_\gamma^Z)
      \colon Z_U\to P_\gamma.                                  
  \end{gathered} \tag{PS12}
  \]
  The compatibility equations defining $\beta_\gamma$ and
  $\beta_\gamma^{-1}$ are
  \[
  \begin{aligned}
    p\circ\rho_Z
      &=\operatorname{pr}_X\circ i_Z\circ\rho_Z
       =\operatorname{pr}_X\circ\jmath_{U,Y}\circ\rho_W
       =j\circ\operatorname{pr}_U\circ\rho_W,\\
    \operatorname{pr}_X\circ i_Z\circ\operatorname{pr}_Z
      &=p\circ\operatorname{pr}_Z
       =j\circ p_U
       =\operatorname{pr}_X\circ\jmath_{U,Y}\circ w_\gamma^Z,\\
    \operatorname{pr}_Y\circ i_Z\circ\operatorname{pr}_Z
      &=\operatorname{pr}_Y\circ\jmath_{U,Y}\circ w_\gamma^Z.
  \end{aligned}                                                   \tag{PS13}
  \]
  Product-morphism extensionality applied to the last two lines is
  \[
    i_Z\circ\operatorname{pr}_Z
      =\jmath_{U,Y}\circ w_\gamma^Z.
  \]
  The projection equalities in (PS12)--(PS13) are the following inverse
  equalities:
  \[
    \beta_\gamma^{-1}\circ\beta_\gamma
      =\operatorname{id}_{P_\gamma},\qquad
    \beta_\gamma\circ\beta_\gamma^{-1}
      =\operatorname{id}_{Z_U},                                  \tag{PS14}
  \]
  by fibre-product uniqueness.

  Define the required morphism and its inverse by
  \[
  \begin{gathered}
    \zeta_\gamma:=\beta_\gamma\circ\delta_\gamma
      \circ g_\gamma^{\mathrm{im}}\circ\theta_\gamma
      \colon U\to Z_U,\\
    \zeta_\gamma^{-1}:=a_U\circ\epsilon_\gamma
      \circ\delta_\gamma^{-1}\circ\beta_\gamma^{-1}
      \colon Z_U\to U.                                         
  \end{gathered} \tag{PS15}
  \]
  Equations (PS6), (PS8), (PS9), and (PS14) are the following inverse
  equality chains:
  \[
  \begin{aligned}
    \zeta_\gamma^{-1}\circ\zeta_\gamma
      &=a_U\circ\epsilon_\gamma\circ g_\gamma^{\mathrm{im}}
          \circ\theta_\gamma
       =a_U\circ\theta_\gamma=\operatorname{id}_U,\\
    \zeta_\gamma\circ\zeta_\gamma^{-1}
      &=\beta_\gamma\circ\delta_\gamma
          \circ g_\gamma^{\mathrm{im}}\circ\epsilon_\gamma
          \circ\delta_\gamma^{-1}\circ\beta_\gamma^{-1}
       =\operatorname{id}_{Z_U}.
  \end{aligned}                                                  \tag{PS16}
  \]
  Finally, (PS4), (PS7), (PS11), and (PS12) are the following two required
  component equality chains:
  \[
  \begin{aligned}
    \operatorname{pr}_Z\circ\zeta_\gamma
      &=\rho_Z\circ\delta_\gamma\circ g_\gamma^{\mathrm{im}}
          \circ\theta_\gamma
       =\gamma^{\mathrm{im}}\circ a_U\circ\theta_\gamma
       =\gamma^{\mathrm{im}},\\
    p_U\circ\zeta_\gamma
      &=\operatorname{pr}_U\circ\rho_W\circ\delta_\gamma
          \circ g_\gamma^{\mathrm{im}}\circ\theta_\gamma\\
      &=\operatorname{pr}_U\circ i_{g_\gamma}
          \circ g_\gamma^{\mathrm{im}}\circ\theta_\gamma
       =\operatorname{pr}_U\circ a_W\circ\theta_\gamma
       =\operatorname{id}_U.
  \end{aligned}                                                  \tag{PS17}
  \]
  If $\zeta'\colon U\to Z_U$ has the two components in (PS17), then
  \[
    \operatorname{pr}_Z\circ\zeta'
      =\gamma^{\mathrm{im}}
      =\operatorname{pr}_Z\circ\zeta_\gamma,
    \qquad
    p_U\circ\zeta'
      =\operatorname{id}_U
      =p_U\circ\zeta_\gamma.
  \]
  Fibre-product uniqueness is the implication
  \[
    \bigl(
      \operatorname{pr}_Z\circ\zeta'
        =\operatorname{pr}_Z\circ\zeta_\gamma
      \ \wedge\
      p_U\circ\zeta'=p_U\circ\zeta_\gamma
    \bigr)
      \Longrightarrow \zeta'=\zeta_\gamma.
  \]
\end{proof}

\begin{lemma}[{\cite[\href{https://stacks.math.columbia.edu/tag/01W6}{Tag 01W6}]{stacks-project}}]
  \label{lemma:nagata-proper-to-separated-target}
  \lean{}

  Let
  \[
    f\colon T\to Y,\qquad g\colon Y\to B,\qquad h\colon T\to B,
    \qquad g\circ f=h.
  \]
  If $h$ is proper and $g$ is separated, then $f$ is proper.

  \paragraph{Mathlib relation.}
  This is exactly
  \texttt{\detokenize{AlgebraicGeometry.IsProper.of_comp}}; the equality
  $g\circ f=h$ only identifies the composite in its hypothesis.
\end{lemma}

\begin{proof}
  Define
  \[
    \Gamma_f:=(\operatorname{id}_T,f)\colon T\to T\times_BY,
    \qquad p_Y\colon T\times_BY\to Y.
  \]
  The fibre-product equation defining $\Gamma_f$ is
  \[
    h\circ\operatorname{id}_T=g\circ f.
  \]
  The graph $\Gamma_f$ is a base change of
  $\Delta_{Y/B}\colon Y\to Y\times_BY$.  Hence
  \[
    g\text{ separated}
      \Longrightarrow\Delta_{Y/B}\text{ closed}
      \Longrightarrow\Gamma_f\text{ closed}
      \Longrightarrow\Gamma_f\text{ proper}.
  \]
  The map $p_Y$ is the base change of the proper morphism $h$, so $p_Y$ is
  proper.  Since
  \[
    f=p_Y\circ\Gamma_f,
  \]
  stability of proper morphisms under composition is the implication
  \[
    \operatorname{IsProper}(\Gamma_f)
      \ \wedge\ 
    \operatorname{IsProper}(p_Y)
      \Longrightarrow
    \operatorname{IsProper}(p_Y\circ\Gamma_f)
      =\operatorname{IsProper}(f).
  \]
\end{proof}

\begin{lemma}[Finite type over a separated target]
  \label{lemma:nagata-finite-type-to-separated-target}
  \lean{}

  Let
  \[
    f\colon T\to Y,\qquad
    g\colon Y\to B,\qquad
    h\colon T\to B,\qquad
    g\circ f=h,                                                    \tag{FTST0}
  \]
  with the symbols typed as in
  \Cref{lemma:nagata-proper-to-separated-target}.  If \(h\) is of finite
  type and \(g\) is separated, then \(f\) is of finite type.
\end{lemma}

\begin{proof}
  Define
  \[
    \Gamma_f=(\operatorname{id}_T,f)\colon T\to T\times_BY,
    \qquad p_Y\colon T\times_BY\to Y.
  \]
  The graph calculation in
  \Cref{lemma:nagata-proper-to-separated-target}, instantiated at
  $(f,g,h)$, is
  \[
    \Gamma_f\text{ is a closed immersion},\qquad f=p_Y\circ\Gamma_f.
  \]
  A closed immersion is of finite type, and $p_Y$ is the base change of the
  finite-type morphism $h$.  Therefore
  \[
    \operatorname{FiniteType}(\Gamma_f)
      \wedge\operatorname{FiniteType}(p_Y)
      \Longrightarrow
    \operatorname{FiniteType}(p_Y\circ\Gamma_f)
      =\operatorname{FiniteType}(f).
  \]
\end{proof}

\begin{corollary}[Quasi-compactness of an open map to a separated target]
  \label{lemma:nagata-open-immersion-qc-over-separated-target}
  \lean{}

  In the triangle of
  \Cref{lemma:nagata-finite-type-to-separated-target}, assume in addition
  that $f\colon T\hookrightarrow Y$ is an open immersion.  Then $f$ is
  quasi-compact.
\end{corollary}

\begin{proof}
  \uses{lemma:nagata-finite-type-to-separated-target}
  By \Cref{lemma:nagata-finite-type-to-separated-target},
  \[
    \operatorname{FiniteType}(f).
  \]
  For an open immersion one has the standard equivalence
  \[
    \operatorname{FiniteType}(f)
      \Longleftrightarrow\operatorname{QuasiCompact}(f).
  \]
  The preceding predicate and equivalence form the implication
  \[
    \operatorname{FiniteType}(f)
      \Longrightarrow\operatorname{QuasiCompact}(f).
  \]
\end{proof}

\begin{definition}[Common-modification construction]
  \label{def:nagata-common-modification}
  \uses{def:nagata-scheme-theoretic-image,
    def:nagata-scheme-image-factor,
    def:nagata-choice-notation,
    def:nagata-admissible-blowup,
    def:nagata-elimination-indeterminacy,
    def:nagata-blowup-product-domination-maps,
    def:nagata-global-blowup,
    def:nagata-extended-admissible-blowup,
    lemma:nagata-scheme-image-partial-section,
    lemma:nagata-elimination-indeterminacy,
    thm:nagata-blowup-flat-base-change,
    thm:nagata-blowup-product-ideal,
    thm:nagata-strict-transform-is-blowup}
  \lean{}

  Let $s_i\colon X_i\to S$ and $u\colon U\to S$ be separated and of finite
  type, let $S$ be qcqs, and let
  $j_i\colon U\hookrightarrow X_i$ be quasi-compact open immersions satisfying
  $s_i\circ j_i=u$.  Define the following symbols in order.
  \begin{enumerate}
    \item Define
      \[
        \gamma_{12}:=(j_1,j_2)\colon U\to X_1\times_SX_2.          \tag{CM0}
      \]
      Define
      \[
      \begin{gathered}
        \sigma_{12}:=(\operatorname{id}_U,j_2)
          \colon U\to U\times_SX_2,\qquad
        \jmath_{12}:=j_1\times_S\operatorname{id}_{X_2}
          \colon U\times_SX_2\hookrightarrow X_1\times_SX_2,\\
        \operatorname{pr}_U\circ\sigma_{12}=\operatorname{id}_U,\qquad
        \gamma_{12}=\jmath_{12}\circ\sigma_{12}.
      \end{gathered}                                                \tag{CM0a}
      \]
      Since \(\operatorname{pr}_U\) is separated as the base change of
      \(s_2\), \(\sigma_{12}\) is a closed immersion; since \(j_1\) is
      quasi-compact, \(\jmath_{12}\) is a quasi-compact open immersion.
      Stability of qcqs morphisms under composition is the implication
      \[
        \operatorname{IsQCQS}(\sigma_{12})\wedge
        \operatorname{IsQCQS}(\jmath_{12})
          \Longrightarrow
        \operatorname{IsQCQS}(\gamma_{12}).                       \tag{CM0b}
      \]
      Define
      \[
        X_{12}:=\operatorname{SIm}_{X_1\times_SX_2}(\gamma_{12}),
      \]
      define the closed immersion of
      \Cref{def:nagata-scheme-theoretic-image} by
      \[
        c_{12}:=i_{\gamma_{12}}\colon
          X_{12}\hookrightarrow X_1\times_SX_2,
      \]
      By \Cref{def:nagata-scheme-image-factor}, define
      \[
        p_i:=\operatorname{pr}_i\circ c_{12}\colon X_{12}\to X_i,
        \qquad
        u_{12}:=\gamma_{12}^{\mathrm{im}}\colon U\to X_{12},
        \qquad
        c_{12}\circ u_{12}=\gamma_{12}.                           \tag{CM0c'}
      \]
      Stability of separatedness and finite type under base change and
      composition, applied to the closed immersion \(c_{12}\) and the two
      projections, is the predicate pair
      \[
        \operatorname{Separated}(p_i),\qquad
        \operatorname{FiniteType}(p_i).
      \]
    \item For \(i=1,2\), define the instance of
      \Cref{lemma:nagata-scheme-image-partial-section} for the projection
      \(p_i\) by
      \[
        \zeta_i:=(u_{12},\operatorname{id}_U)\colon
          U\xrightarrow{\sim}X_{12}\times_{X_i}U,\qquad
        \eta_i:=\zeta_i^{-1}\colon
          X_{12}\times_{X_i}U\xrightarrow{\sim}U.                 \tag{CM0d}
      \]
      Apply
      \Cref{def:nagata-elimination-indeterminacy} to $(p_i,U,\eta_i)$ and
      define
      \[
      \begin{gathered}
        \mathcal I_i:=\mathcal F_U(p_i,\eta_i),\qquad
        \pi_i\colon X_i^i:=
          \operatorname{Bl}_{\mathcal I_i}(X_i)\to X_i,\\
        A_i:=\operatorname{ST}_{\pi_i}(X_{12}),\qquad
        e_i:=p_i^{\mathrm{el}}\colon A_i\hookrightarrow X_i^i,\\
        \mathcal J_i:=p_i^{-1}\mathcal I_i\mathcal O_{X_{12}},\\
        q_i^A:=\operatorname{pr}_{X_{12}}\circ\iota_{p_i,\pi_i}
          \colon A_i\longrightarrow X_{12}.
      \end{gathered}
      \]
      The projection equation for
      \(\iota_{p_i,\pi_i}\colon
        A_i\hookrightarrow X_{12}\times_{X_i}X_i^i\) is
      \[
        p_i\circ q_i^A=\pi_i\circ e_i.
      \]
      Define the preserved open in \(X_i^i\) and its lift to \(A_i\) by
      \[
      \begin{gathered}
        \alpha_i^U\colon X_i^i\times_{X_i}U\xrightarrow{\sim}U,\qquad
        \overline j_i^U:=\operatorname{pr}_{X_i^i}
          \colon X_i^i\times_{X_i}U\hookrightarrow X_i^i,\\
        \widetilde j_i:=
          \overline j_i^U\circ(\alpha_i^U)^{-1}
          \colon U\hookrightarrow X_i^i,\\
        \sigma_i\colon U\hookrightarrow A_i,\qquad
        e_i\circ\sigma_i=\widetilde j_i,\qquad
        q_i^A\circ\sigma_i=u_{12}.
      \end{gathered}                                                \tag{CM0c}
      \]
      Here \(\alpha_i^U\) is the complement isomorphism from
      \Cref{def:nagata-admissible-blowup}, and \(\sigma_i\) is the inverse
      of the restriction isomorphism in
      \Cref{lemma:nagata-elimination-indeterminacy}.
      Define the actual overlap required by
      \Cref{def:nagata-extended-admissible-blowup} and its two projections by
      \[
      \begin{gathered}
        U_i^A:=A_i\times_{X_i^i,e_i,\widetilde j_i}U,
        \qquad
        \iota_i^A:=\operatorname{pr}_{A_i}\colon U_i^A\to A_i,
        \qquad
        \delta_i^A:=\operatorname{pr}_U\colon U_i^A\to U,\\
        \zeta_i^A:=(\sigma_i,\operatorname{id}_U)
          \colon U\longrightarrow U_i^A.
      \end{gathered}
      \]
      The fibre-product equation and (CM0c) form the chain
      \[
        e_i\circ\iota_i^A
          =\widetilde j_i\circ\delta_i^A
          =e_i\circ\sigma_i\circ\delta_i^A.
      \]
      Monicity of \(e_i\) is the equality
      \[
        \iota_i^A=\sigma_i\circ\delta_i^A.
      \]
      Together with \(\delta_i^A\circ\zeta_i^A=\operatorname{id}_U\), the two
      inverse equations are
      \[
      \begin{gathered}
        \delta_i^A\circ\zeta_i^A=\operatorname{id}_U,
        \qquad
        \iota_i^A=\sigma_i\circ\delta_i^A,\\
        \zeta_i^A\circ\delta_i^A
          =(\sigma_i\circ\delta_i^A,\delta_i^A)
          =(\iota_i^A,\delta_i^A)=\operatorname{id}_{U_i^A}.
      \end{gathered}                                               \tag{CM0f}
      \]
      Equations (CM0f) define
      \(\zeta_i^A\colon U\xrightarrow{\sim}U_i^A\) with inverse
      \(\delta_i^A\); the preserved-open morphism used below is
      \(\iota_i^A\colon U_i^A\hookrightarrow A_i\).
      Since \(\overline j_i^U\) is the base change of the quasi-compact open
      immersion \(j_i\), stability under base change and composition with an
      isomorphism is the predicate chain
      \[
      \begin{gathered}
        \operatorname{IsQuasiCompactOpenImmersion}(\overline j_i^U),\\
        \widetilde j_i=\overline j_i^U\circ(\alpha_i^U)^{-1}
          \quad\Longrightarrow\quad
        \operatorname{IsQuasiCompactOpenImmersion}(\widetilde j_i),\\
        \operatorname{IsQuasiCompactOpenImmersion}(\iota_i^A),
      \end{gathered}                                               \tag{CM0f'}
      \]
      where the last predicate is base-change stability for
      \(U_i^A=A_i\times_{X_i^i}U\).
    \item Put \(\widehat 1:=2\) and \(\widehat 2:=1\).  For \(i=1,2\),
      define
      \[
      \begin{gathered}
        B_i:=\operatorname{Bl}_{\mathcal J_i}(X_{12}),\qquad
        b_i^B:=b_{X_{12},\mathcal J_i}\colon B_i\to X_{12},\\
        \mathcal H_i:=\mathcal J_{\widehat i}\mathcal O_{B_i},
        \qquad
        \mathcal K_{12}:=\mathcal J_1\mathcal J_2,\\
        X'_{12}:=\operatorname{Bl}_{\mathcal K_{12}}(X_{12}),
        \qquad
        b_{12}:=b_{X_{12},\mathcal K_{12}}
          \colon X'_{12}\to X_{12}.
      \end{gathered}
      \]
      The strict-transform comparison and product-ideal comparison define
      the isomorphisms and domination map
      \[
      \begin{gathered}
        \Theta_i\colon A_i\xrightarrow{\sim}B_i,\qquad
        b_i^B\circ\Theta_i=q_i^A,\\
        \Psi_i:=\Psi_{\mathcal J_i,\mathcal J_{\widehat i}}
          \colon X'_{12}\xrightarrow{\sim}
            \operatorname{Bl}_{\mathcal H_i}(B_i),\\
        d_i:=b_{B_i,\mathcal H_i}\circ\Psi_i
          \colon X'_{12}\longrightarrow B_i,\qquad
        b_i^B\circ d_i=b_{12}.
      \end{gathered}                                               \tag{CM0g}
      \]
      Here \(\Theta_i\) is the isomorphism of
      \Cref{thm:nagata-strict-transform-is-blowup}, and
      \(\Psi_i\) and \(d_i\) are the instances of
      \Cref{thm:nagata-blowup-product-ideal,def:nagata-blowup-product-domination-maps}.

      Define the transported center on \(A_i\) by
      \[
        \mathcal L_i:=
          \Theta_i^{-1}\mathcal H_i\mathcal O_{A_i}
          =\Theta_i^*\mathcal H_i
          \subseteq\mathcal O_{A_i}.                              \tag{CM0h}
      \]
      The ideal \(\mathcal L_i\) is finite-type and quasi-coherent.  The
      base-map and ideal-restriction chains defined by
      (CM0c), (CM0f), and (CM0g) are
      \[
      \begin{aligned}
        b_i^B\circ\Theta_i\circ\iota_i^A
          &=q_i^A\circ\iota_i^A
           =q_i^A\circ\sigma_i\circ\delta_i^A
           =u_{12}\circ\delta_i^A,\\
        u_{12}^{-1}\mathcal J_{\widehat i}\mathcal O_U
          &=u_{12}^{-1}
              (p_{\widehat i}^{-1}\mathcal I_{\widehat i}
                \mathcal O_{X_{12}})\mathcal O_U\\
          &=(p_{\widehat i}\circ u_{12})^{-1}
              \mathcal I_{\widehat i}\mathcal O_U
           =j_{\widehat i}^{-1}\mathcal I_{\widehat i}\mathcal O_U
           =\mathcal O_U,\\
        \mathcal L_i\mathcal O_{U_i^A}
          &=(\Theta_i\circ\iota_i^A)^{-1}
              \mathcal H_i\mathcal O_{U_i^A}\\
          &=(u_{12}\circ\delta_i^A)^{-1}
              \mathcal J_{\widehat i}\mathcal O_{U_i^A}
           =(\delta_i^A)^{-1}\mathcal O_U\mathcal O_{U_i^A}
           =\mathcal O_{U_i^A}.
      \end{aligned}                                                \tag{CM0i}
      \]

      To define the transported blow-up realization, put
      \[
      \begin{gathered}
        F_i:=A_i\times_{B_i,\Theta_i,b_{B_i,\mathcal H_i}}
          \operatorname{Bl}_{\mathcal H_i}(B_i),\\
        f_i^A\colon F_i\to A_i,\qquad
        f_i^B\colon F_i\to\operatorname{Bl}_{\mathcal H_i}(B_i),\\
        \upsilon_i:=
          (\Theta_i^{-1}\circ b_{B_i,\mathcal H_i},
            \operatorname{id}_{\operatorname{Bl}_{\mathcal H_i}(B_i)})
          \colon\operatorname{Bl}_{\mathcal H_i}(B_i)
            \xrightarrow{\sim}F_i,
        \qquad \upsilon_i^{-1}=f_i^B.
      \end{gathered}
      \]
      Since \(\Theta_i\) is flat and
      \(\mathcal L_i=\Theta_i^*\mathcal H_i\),
      \Cref{thm:nagata-blowup-flat-base-change} defines
      \[
        \operatorname{bc}_{\Theta_i}\colon
          \operatorname{Bl}_{\mathcal L_i}(A_i)
            \xrightarrow{\sim}F_i.
      \]
      Define
      \[
      \begin{gathered}
        \tau_i:=\operatorname{bc}_{\Theta_i}^{-1}\circ\upsilon_i
          \colon\operatorname{Bl}_{\mathcal H_i}(B_i)
            \xrightarrow{\sim}\operatorname{Bl}_{\mathcal L_i}(A_i),\\
        \Omega_i:=\tau_i\circ\Psi_i
          \colon X'_{12}\xrightarrow{\sim}
            \operatorname{Bl}_{\mathcal L_i}(A_i),\\
        r_i:=\Theta_i^{-1}\circ d_i\colon X'_{12}\longrightarrow A_i.
      \end{gathered}
      \]
      The structural equality for flat base change and (CM0g) form the chain
      \[
      \begin{aligned}
        b_{A_i,\mathcal L_i}\circ\Omega_i
          &=b_{A_i,\mathcal L_i}\circ\tau_i\circ\Psi_i\\
          &=\Theta_i^{-1}\circ b_{B_i,\mathcal H_i}\circ\Psi_i\\
          &=\Theta_i^{-1}\circ d_i=r_i.
      \end{aligned}                                                \tag{CM0j}
      \]

      Define the complement and the factorization of the preserved open by
      \[
      \begin{gathered}
        V_i^{\mathcal L}:=A_i\setminus V(\mathcal L_i),\qquad
        j_i^{\mathcal L}\colon V_i^{\mathcal L}\hookrightarrow A_i,\\
        \nu_i^{\mathcal L}\colon U_i^A\longrightarrow V_i^{\mathcal L},
        \qquad
        j_i^{\mathcal L}\circ\nu_i^{\mathcal L}=\iota_i^A.
      \end{gathered}
      \]
      Equation (CM0i) is the factorization criterion defining
      \(\nu_i^{\mathcal L}\).  Let
      \[
        \alpha_i^V\colon
          \operatorname{Bl}_{\mathcal L_i}(A_i)
            \times_{A_i}V_i^{\mathcal L}
              \xrightarrow{\sim}V_i^{\mathcal L}
      \]
      be \Cref{thm:nagata-global-blowup-complement}, and define its base change
      along \(\nu_i^{\mathcal L}\) by
      \[
      \begin{gathered}
        \alpha_i^{\mathcal L}\colon
          \operatorname{Bl}_{\mathcal L_i}(A_i)
            \times_{A_i}U_i^A\xrightarrow{\sim}U_i^A,\\
        \iota_i^A\circ\alpha_i^{\mathcal L}
          =b_{A_i,\mathcal L_i}\circ
            \operatorname{pr}_{\operatorname{Bl}_{\mathcal L_i}(A_i)}.
      \end{gathered}
      \]
      Define the base change of \(\Omega_i\) and the transported complement
      isomorphism by
      \[
      \begin{gathered}
        \Omega_{i,U}:=
          (\Omega_i\circ\operatorname{pr}_{X'_{12}},
            \operatorname{pr}_{U_i^A})
          \colon X'_{12}\times_{A_i,r_i,\iota_i^A}U_i^A
          \xrightarrow{\sim}
          \operatorname{Bl}_{\mathcal L_i}(A_i)
            \times_{A_i,b_{A_i,\mathcal L_i},\iota_i^A}U_i^A,\\
        \alpha_i^r:=\alpha_i^{\mathcal L}\circ\Omega_{i,U}
          \colon X'_{12}\times_{A_i}U_i^A\xrightarrow{\sim}U_i^A.
      \end{gathered}
      \]
      Its structural equality is the chain
      \[
      \begin{aligned}
        \iota_i^A\circ\alpha_i^r
          &=\iota_i^A\circ\alpha_i^{\mathcal L}\circ\Omega_{i,U}\\
          &=b_{A_i,\mathcal L_i}\circ
              \operatorname{pr}_{\operatorname{Bl}_{\mathcal L_i}(A_i)}
                \circ\Omega_{i,U}\\
          &=b_{A_i,\mathcal L_i}\circ\Omega_i\circ
              \operatorname{pr}_{X'_{12}}\\
          &=r_i\circ\operatorname{pr}_{X'_{12}}.
      \end{aligned}                                                \tag{CM0k}
      \]
      Define the explicitly transported admissible-blow-up datum on \(A_i\) by
      \[
        D_i^r:=(\mathcal L_i,X'_{12},r_i,\alpha_i^r),              \tag{CM0l}
      \]
      with source realization \(\Omega_i\) and preserved open
      \(\iota_i^A\colon U_i^A\hookrightarrow A_i\).
      Define the structural morphism
      \[
        s_S\colon S\longrightarrow\operatorname{Spec}(\mathbf Z).
      \]
      The qcqs/finite-type stability implication and the quasi-compactness
      criterion for an open immersion are
      \[
      \begin{gathered}
        \operatorname{IsQCQS}(s_S)\ \wedge\
          \operatorname{FiniteType}(s_i)\ \wedge\
        (\mathcal I_i,\mathcal J_i,\mathcal L_i
          \text{ finite-type quasi-coherent ideals})\\
        \Longrightarrow
        \operatorname{QuasiCompact}(A_i)\ \wedge\
          \operatorname{QuasiSeparated}(X_i^i),\\
        \operatorname{QuasiCompact}(A_i)\ \wedge\
          \operatorname{QuasiSeparated}(X_i^i)\ \wedge\
          \operatorname{IsOpenImmersion}(e_i)
        \Longrightarrow
          \operatorname{IsQuasiCompactOpenImmersion}(e_i).
      \end{gathered}                                               \tag{CM0l'}
      \]

    \item Apply \Cref{def:nagata-extended-admissible-blowup} with the exact
      input assignment
      \[
      \begin{aligned}
        &(X,W,U,\iota_W,\iota_U,
          (\mathcal J,W',\pi_W,\alpha_{U_W}))\\
        &\qquad:=
          (X_i^i,A_i,U,e_i,\widetilde j_i,
            (\mathcal L_i,X'_{12},r_i,\alpha_i^r)),
        \qquad U_W=U_i^A.
      \end{aligned}                                                \tag{CM0m}
      \]
      Denote its output morphism and its untransported base-change
      isomorphism by
      \[
      \begin{gathered}
        \varepsilon_i\colon X'_i\longrightarrow X_i^i,\\
        \rho_i^{\mathcal L}\colon
          X'_i\times_{X_i^i}A_i
            \xrightarrow{\sim}\operatorname{Bl}_{\mathcal L_i}(A_i),
        \qquad
        b_{A_i,\mathcal L_i}\circ\rho_i^{\mathcal L}
          =\operatorname{pr}_{A_i}.
      \end{gathered}
      \]
      Define
      \[
      \begin{gathered}
        \rho_i:=\Omega_i^{-1}\circ\rho_i^{\mathcal L}
          \colon X'_i\times_{X_i^i}A_i\xrightarrow{\sim}X'_{12},\\
        \jmath_i^A:=\operatorname{pr}_{X'_i}
          \colon X'_i\times_{X_i^i}A_i\hookrightarrow X'_i,\\
        k_i:=\jmath_i^A\circ\rho_i^{-1}
          \colon X'_{12}\hookrightarrow X'_i.
      \end{gathered}                                               \tag{CM0n}
      \]
      Stability of open immersions under base change and composition, applied
      to \(\jmath_i^A\), \(\rho_i^{-1}\), and (CM0n), is the predicate
      \[
        \operatorname{IsOpenImmersion}(k_i).                      \tag{CM0o}
      \]

    \item For every scheme \(T\), define the set of compatible two-chart maps
      \[
        \mathscr E_T^{\mathsf{CM}}:=
          \left\{(h_1,h_2)\in
            \operatorname{Hom}(X'_1,T)\times\operatorname{Hom}(X'_2,T)
          \ \middle|\ h_1\circ k_1=h_2\circ k_2\right\}.       \tag{CM1a}
      \]
      Define the set of two-chart gluing realizations
      \[
      \begin{aligned}
        \mathscr G_{\mathsf{CM}}:=
        \bigl\{(G,e_1^G,e_2^G)\ \bigm|\ {}&
          G\text{ is a scheme},\quad
          e_i^G\colon X'_i\to G,\quad
          \operatorname{IsOpenImmersion}(e_i^G)\ (i=1,2),\\
        &|G|=|e_1^G|(|X'_1|)\cup |e_2^G|(|X'_2|),\\
        &e_1^G\circ k_1=e_2^G\circ k_2,\\
        &\operatorname{IsIso}\!\left(
          (k_1,k_2)\colon X'_{12}\longrightarrow
            X'_1\times_{G,e_1^G,e_2^G}X'_2\right),\\
        &\forall T,\quad
          \left[
          \begin{array}{rcl}
            \operatorname{Hom}(G,T)&\longrightarrow&
              \mathscr E_T^{\mathsf{CM}}\\
            h&\longmapsto&(h\circ e_1^G,h\circ e_2^G)
          \end{array}
          \right]\text{ is bijective}\bigr\}.
      \end{aligned}                                                \tag{CM1b}
      \]
      Equation (CM0o), together with the identity transition on
      \(X'_{12}\), satisfies the inverse and cocycle equations
      \[
        \operatorname{id}_{X'_{12}}^{-1}
          =\operatorname{id}_{X'_{12}},\qquad
        \operatorname{id}_{X'_{12}}\circ
          \operatorname{id}_{X'_{12}}
          =\operatorname{id}_{X'_{12}}.
      \]
      The existence assertion of the scheme-gluing theorem
      \cite[\href{https://stacks.math.columbia.edu/tag/01JA}{Tag 01JA}]
      {stacks-project} is the statement
      \[
        \mathscr G_{\mathsf{CM}}\ne\varnothing.                  \tag{CM1c}
      \]
      By \Cref{def:nagata-choice-notation}, define
      \[
        (\mathsf{CM}(X_1,X_2;U),\lambda_1,\lambda_2)
          :=\operatorname{ch}(\mathscr G_{\mathsf{CM}})
          \in\mathscr G_{\mathsf{CM}}.                           \tag{CM1d}
      \]
      Membership (CM1d) defines the open immersions
      \(\lambda_i\colon X'_i\hookrightarrow
        \mathsf{CM}(X_1,X_2;U)\) and is the following equality-and-bijection
      tuple:
      \[
      \begin{gathered}
        |\mathsf{CM}(X_1,X_2;U)|
          =\lambda_1(|X'_1|)\cup\lambda_2(|X'_2|),\qquad
        \lambda_1\circ k_1=\lambda_2\circ k_2,\\
        \Gamma_T^{\mathsf{CM}}\colon
          \operatorname{Hom}(\mathsf{CM}(X_1,X_2;U),T)
            \xrightarrow{\sim}\mathscr E_T^{\mathsf{CM}},\qquad
        \Gamma_T^{\mathsf{CM}}(h)
          :=(h\circ\lambda_1,h\circ\lambda_2).
      \end{gathered}                                               \tag{CM1e}
      \]
      Finally define
      \[
      \begin{gathered}
        M_{12}:=X'_1\times_{
          \mathsf{CM}(X_1,X_2;U),\lambda_1,\lambda_2}X'_2,
        \qquad
        m_i:=\operatorname{pr}_{X'_i}\colon M_{12}\to X'_i,\\
        \theta_{12}:=(k_1,k_2)
          \colon X'_{12}\xrightarrow{\sim}M_{12},
        \qquad
        m_i\circ\theta_{12}=k_i\quad(i=1,2).
      \end{gathered}                                               \tag{CM1}
      \]
  \end{enumerate}
  Finally define
  \[
    \phi_i:=\pi_i\circ\varepsilon_i\colon X'_i\to X_i.
  \]
  Define
  \[
    s_{12}:=s_1\circ p_1=s_2\circ p_2\colon X_{12}\to S           \tag{CM2a}
  \]
  by the fibre-product equality
  \[
    s_1\circ p_1
      =s_1\circ\operatorname{pr}_1\circ c_{12}
      =s_2\circ\operatorname{pr}_2\circ c_{12}
      =s_2\circ p_2.                                               \tag{CM2b}
  \]
  For \(i=1,2\), define
  \[
    \xi_i:=\operatorname{pr}_{A_i}\circ\rho_i^{-1}
      \colon X'_{12}\to A_i.
  \]
  The substitutions from (CM0j), (CM0n), the structural equation for
  \(\rho_i^{\mathcal L}\), and the two fibre-product projections are the
  equality chains
  \[
  \begin{aligned}
    \xi_i
      &=\operatorname{pr}_{A_i}\circ
          (\rho_i^{\mathcal L})^{-1}\circ\Omega_i
       =b_{A_i,\mathcal L_i}\circ\Omega_i=r_i,\\
    \varepsilon_i\circ k_i
      &=\varepsilon_i\circ\jmath_i^A\circ\rho_i^{-1}
       =e_i\circ\operatorname{pr}_{A_i}\circ\rho_i^{-1}
       =e_i\circ\xi_i,\\
    \pi_i\circ e_i\circ r_i
      &=p_i\circ q_i^A\circ\Theta_i^{-1}\circ d_i
       =p_i\circ b_i^B\circ d_i
       =p_i\circ b_{12}.
  \end{aligned}                                                     \tag{CM2c}
  \]
  Define the equality witness \(e_{i,12}\) to be the chain
  \[
  \begin{aligned}
    e_{i,12}\colon\quad
    \phi_i\circ k_i
      &=\pi_i\circ\varepsilon_i\circ k_i\\
      &=\pi_i\circ e_i\circ\xi_i\\
      &=\pi_i\circ e_i\circ r_i\\
      &=p_i\circ b_{12}.
  \end{aligned}                                                     \tag{CM2d}
  \]
  Define
  \[
    s_i':=s_i\circ\phi_i\colon X'_i\to S.                          \tag{CM2e}
  \]
  Equations (CM2a), (CM2d), and (CM2e) form the overlap chain
  \[
  \begin{aligned}
    s_1'\circ k_1
      &=s_1\circ\phi_1\circ k_1
       =s_1\circ p_1\circ b_{12}
       =s_{12}\circ b_{12}\\
      &=s_2\circ p_2\circ b_{12}
       =s_2\circ\phi_2\circ k_2
       =s_2'\circ k_2.
  \end{aligned}                                                     \tag{CM2f}
  \]
  By (CM1a), equation (CM2f) is precisely the membership and its image under
  the bijection (CM1e):
  \[
    (s_1',s_2')\in\mathscr E_S^{\mathsf{CM}},\qquad
    \Gamma_S^{\mathsf{CM}}\!\left(
      (\Gamma_S^{\mathsf{CM}})^{-1}(s_1',s_2')\right)
      =(s_1',s_2').                                                \tag{CM2g}
  \]
  Define
  \[
    s_{\mathsf{CM}}:=(\Gamma_S^{\mathsf{CM}})^{-1}(s_1',s_2')
      \colon\mathsf{CM}(X_1,X_2;U)\to S,\qquad
    s_{\mathsf{CM}}\circ\lambda_i=s_i'=s_i\circ\phi_i.             \tag{CM2h}
  \]
  Since $\mathcal K_{12}|_U=\mathcal O_U$, let
  \[
    \alpha_{12}\colon X'_{12}\times_{X_{12}}U\xrightarrow{\sim}U
  \]
  be \Cref{thm:nagata-global-blowup-complement}, and define
  \[
    a_{12}:=\operatorname{pr}_{X'_{12}}\circ\alpha_{12}^{-1}
      \colon U\to X'_{12},
    \qquad
    a_i':=k_i\circ a_{12}\colon U\to X'_i.                        \tag{CM2}
  \]
  Define the identified morphism from $U$ by
  \[
    j_{\mathsf{CM}}:=\lambda_1\circ a_1'
      =\lambda_2\circ a_2'\colon U\longrightarrow
        \mathsf{CM}(X_1,X_2;U),                                   \tag{CM3}
  \]
  The equality in (CM3) is the chain
  \[
    \lambda_1\circ a_1'
      =\lambda_1\circ k_1\circ a_{12}
      =\lambda_2\circ k_2\circ a_{12}
      =\lambda_2\circ a_2'.                                      \tag{CM3a}
  \]
  The complement-isomorphism equations are
  \[
    \phi_i\circ a_i'=j_i\qquad(i=1,2).                            \tag{CM3b}
  \]
  Equations (CM2h), (CM3), and (CM3b) form the structural chain
  \[
    s_{\mathsf{CM}}\circ j_{\mathsf{CM}}
      =s_{\mathsf{CM}}\circ\lambda_1\circ a_1'
      =s_1\circ\phi_1\circ a_1'
      =s_1\circ j_1=u.                                            \tag{CM3c}
  \]
\end{definition}

\begin{lemma}[Properties of the common modification,
  {\cite[\href{https://stacks.math.columbia.edu/tag/0F3X}{Tag 0F3X}]{stacks-project}}]
  \label{lemma:nagata-common-modification}
  \uses{def:nagata-common-modification}
  \lean{}

  Each $\phi_i\colon X'_i\to X_i$ is a $U$-admissible blow-up.  Each
  $\lambda_i$ is an open immersion,
  \[
    \lambda_1\circ k_1=\lambda_2\circ k_2,
  \]
  the intersection isomorphism $\theta_{12}$ satisfies (CM1), the two
  restrictions from $U$ are both $j_{\mathsf{CM}}$ of (CM3), and
  \(s_{\mathsf{CM}}\colon\mathsf{CM}(X_1,X_2;U)\to S\) is separated
  and of finite type.
\end{lemma}

\begin{proof}
  \uses{def:nagata-choice-notation,
    def:nagata-closed-subscheme-class,
    lemma:nagata-scheme-image-partial-section,
    lemma:nagata-elimination-indeterminacy,
    lemma:nagata-product-ideal-common-blowup,
    thm:nagata-blowup-product-ideal,
    def:nagata-admissible-blowup,
    lemma:nagata-extend-admissible-blowup,
    lemma:nagata-blowup-proper,
    lemma:nagata-morphism-property-calculus,
    lemma:nagata-proper-to-separated-target,
    lemma:nagata-separated-two-open-cover,
    lemma:nagata-admissible-blowup-comp}
  The instance of \Cref{lemma:nagata-scheme-image-partial-section} at
  \(\gamma_{12}\), with \(\zeta_i,\eta_i\) defined by (CM0d), is
  \[
    \eta_i:=\zeta_i^{-1}\colon
      X_{12}\times_{X_i}U\xrightarrow{\sim}U,\qquad
      \eta_i\circ\zeta_i=\operatorname{id}_U,\qquad
    \zeta_i\circ\eta_i
      =\operatorname{id}_{X_{12}\times_{X_i}U}.
  \]
  The conclusion of \Cref{lemma:nagata-elimination-indeterminacy}, applied to
  $(p_i,U,\eta_i)$, is
  \[
    \pi_i\text{ is }U\text{-admissible},
    \qquad e_i\colon A_i\hookrightarrow X_i^i\text{ is open}.
  \]
  Equations (CM0h)--(CM0l) are the component-by-component verification
  \[
  \begin{gathered}
    \mathcal L_i\subseteq\mathcal O_{A_i}
      \text{ finite-type and quasi-coherent},\qquad
    \mathcal L_i\mathcal O_{U_i^A}=\mathcal O_{U_i^A},\\
    \Omega_i\colon X'_{12}\xrightarrow{\sim}
      \operatorname{Bl}_{\mathcal L_i}(A_i),\qquad
    r_i=b_{A_i,\mathcal L_i}\circ\Omega_i,\\
    \alpha_i^r\colon X'_{12}\times_{A_i}U_i^A
      \xrightarrow{\sim}U_i^A,\qquad
    \iota_i^A\circ\alpha_i^r
      =r_i\circ\operatorname{pr}_{X'_{12}}.
  \end{gathered}
  \]
  This displayed component tuple is precisely
  \Cref{def:nagata-admissible-blowup} for
  \(D_i^r=(\mathcal L_i,X'_{12},r_i,\alpha_i^r)\), its specified realization
  \(\Omega_i\), and
  \(\iota_i^A\colon U_i^A\hookrightarrow A_i\).
  The exact input substitution (CM0m) in
  \Cref{lemma:nagata-extend-admissible-blowup} is
  \[
  \begin{gathered}
    \varepsilon_i\colon X'_i\to X_i^i
      \text{ is }U\text{-admissible},\\
    \rho_i^{\mathcal L}\colon X'_i\times_{X_i^i}A_i
      \xrightarrow{\sim}\operatorname{Bl}_{\mathcal L_i}(A_i),
    \qquad
    b_{A_i,\mathcal L_i}\circ\rho_i^{\mathcal L}
      =\operatorname{pr}_{A_i}.
  \end{gathered}
  \]
  The substitution in (CM0n) is the transported isomorphism
  \[
    \rho_i=\Omega_i^{-1}\circ\rho_i^{\mathcal L}
      \colon X'_i\times_{X_i^i}A_i\xrightarrow{\sim}X'_{12}.
  \]
  The instance of \Cref{lemma:nagata-admissible-blowup-comp} at
  $(\pi_i,\varepsilon_i)$ is
  \[
    \phi_i=\pi_i\circ\varepsilon_i
      \text{ is }U\text{-admissible}.
  \]
  Membership (CM1d) in the explicitly defined set (CM1b) is the predicate and
  equality tuple
  \[
  \begin{gathered}
    \operatorname{IsOpenImmersion}(\lambda_i)\quad(i=1,2),\\
    |\mathsf{CM}(X_1,X_2;U)|
      =\lambda_1(|X'_1|)\cup\lambda_2(|X'_2|),\qquad
    \lambda_1\circ k_1=\lambda_2\circ k_2,\\
    \operatorname{IsIso}(\theta_{12}),\qquad
    m_i\circ\theta_{12}=k_i\quad(i=1,2).
  \end{gathered}
  \]
  Substitution of (CM2) into (CM3) is the equality chain
  \[
    \lambda_1\circ a_1'=j_{\mathsf{CM}}
      =\lambda_2\circ a_2'.                                      \tag{CM4}
  \]

  Define
  \[
    \kappa:=(k_1,k_2)\colon X'_{12}\to X'_1\times_SX'_2,
  \]
  \[
    g:=\phi_1\times_S\phi_2\colon
      X'_1\times_SX'_2\to X_1\times_SX_2.
  \]
  Define
  \[
  \begin{gathered}
    P_{k_2}:=X'_{12}\times_{S,(s_1'\circ k_1,s_2')}X'_2,\\
    r_{12}\colon P_{k_2}\to X'_{12},\qquad
    r_2'\colon P_{k_2}\to X'_2,\qquad
    s_1'\circ k_1\circ r_{12}=s_2'\circ r_2',\\
    \Gamma_{k_2}:=(\operatorname{id}_{X'_{12}},k_2)
      \colon X'_{12}\to P_{k_2},\\
    \lambda_{k_1}:=k_1\times_S\operatorname{id}_{X'_2}
      \colon P_{k_2}\to X'_1\times_SX'_2.                       
  \end{gathered} \tag{CM4a}
  \]
  Equation (CM2f) is exactly the compatibility equation defining
  $\Gamma_{k_2}$, and the fibre-product component equalities are
  \[
  \begin{gathered}
    r_{12}\circ\Gamma_{k_2}=\operatorname{id}_{X'_{12}},
    \qquad r_2'\circ\Gamma_{k_2}=k_2,\\
    \operatorname{pr}_1\circ\lambda_{k_1}=k_1\circ r_{12},
    \qquad \operatorname{pr}_2\circ\lambda_{k_1}=r_2'.          
  \end{gathered} \tag{CM4b}
  \]
  The instance of \Cref{lemma:nagata-admissible-blowup-comp} at
  $(\pi_2,\varepsilon_2)$, together with
  \Cref{def:nagata-admissible-blowup,lemma:nagata-blowup-proper}, is the
  following predicate chain:
  \[
    \phi_2=\pi_2\circ\varepsilon_2
      \text{ is }U\text{-admissible}
      \Longrightarrow\operatorname{IsProper}(\phi_2).           \tag{CM4c0}
  \]
  Separatedness of proper morphisms and composition stability are the
  following predicate chain:
  \[
  \begin{gathered}
    \operatorname{IsProper}(\phi_2)
      \Longrightarrow\operatorname{IsSeparated}(\phi_2),\\
    \operatorname{IsSeparated}(\phi_2)\ \wedge\
      \operatorname{IsSeparated}(s_2)
      \Longrightarrow
      \operatorname{IsSeparated}(s_2\circ\phi_2)
      =\operatorname{IsSeparated}(s_2').                         
  \end{gathered} \tag{CM4c}
  \]
  The morphism $r_{12}$ is the base change of $s_2'$, and
  $\Gamma_{k_2}$ is its section.  The morphism $\lambda_{k_1}$ is the base
  change of $k_1$.  Thus
  \[
  \begin{gathered}
    \operatorname{IsSeparated}(s_2')
      \Longrightarrow\operatorname{IsSeparated}(r_{12}),\\
    r_{12}\circ\Gamma_{k_2}=\operatorname{id}_{X'_{12}}
      \ \wedge\ \operatorname{IsSeparated}(r_{12})
      \Longrightarrow\operatorname{IsClosedImmersion}(\Gamma_{k_2}),\\
    \operatorname{IsOpenImmersion}(k_1)
      \Longrightarrow\operatorname{IsOpenImmersion}(\lambda_{k_1}). 
  \end{gathered} \tag{CM4d}
  \]
  Equations (CM4b) and the definition of $\kappa$ are the following component
  equality chains:
  \[
  \begin{aligned}
    \operatorname{pr}_1\circ\lambda_{k_1}\circ\Gamma_{k_2}
      &=k_1,\\
    \operatorname{pr}_2\circ\lambda_{k_1}\circ\Gamma_{k_2}
      &=k_2,
  \end{aligned}
  \qquad
    \kappa=\lambda_{k_1}\circ\Gamma_{k_2}.                     \tag{CM4e}
  \]
  Finally, a closed immersion is an immersion, an open immersion is an
  immersion, and immersions are stable under composition; (CM4d)--(CM4e) are
  the following predicate chain:
  \[
    \operatorname{IsImmersion}(\Gamma_{k_2})\ \wedge\
      \operatorname{IsImmersion}(\lambda_{k_1})
      \Longrightarrow
      \operatorname{IsImmersion}(\lambda_{k_1}\circ\Gamma_{k_2})
      =\operatorname{IsImmersion}(\kappa).                       \tag{CM4f}
  \]
  For
  $i=1,2$, the definitions of $p_i$, $\phi_i$, $k_i$, $r_i$, and $d_i$ form
  the component chain
  \[
  \begin{aligned}
    \operatorname{pr}_i\circ g\circ\kappa
      &=\phi_i\circ k_i\\
      &=\pi_i\circ e_i\circ r_i\\
      &=\pi_i\circ e_i\circ\Theta_i^{-1}\circ d_i\\
      &=p_i\circ b_{12}\\
      &=\operatorname{pr}_i\circ c_{12}\circ b_{12}.
  \end{aligned}
  \]
  The morphism
  \[
    \operatorname{Hom}(X'_{12},X_1\times_SX_2)
      \longrightarrow
    \prod_{i=1}^2\operatorname{Hom}(X'_{12},X_i),\qquad
    h\longmapsto(\operatorname{pr}_1\circ h,
      \operatorname{pr}_2\circ h),
  \]
  is injective; applied to the two component chains, it is the equality
  \[
    g\circ\kappa=c_{12}\circ b_{12}.
  \]
  \Cref{lemma:nagata-blowup-proper}, stability under composition, and
  stability under products are the predicate chains
  \[
  \begin{aligned}
    \operatorname{IsProper}(b_{12})
      \wedge\operatorname{IsClosedImmersion}(c_{12})
      &\Longrightarrow
        \operatorname{IsProper}(c_{12}\circ b_{12}),\\
    \operatorname{Separated}(\phi_1)
      \wedge\operatorname{Separated}(\phi_2)
      &\Longrightarrow\operatorname{Separated}(g).
  \end{aligned}
  \]
  Applying
  \Cref{lemma:nagata-proper-to-separated-target} to the displayed equality
  is the predicate
  \[
    \kappa\text{ is proper}.
  \]
  The proper-immersion criterion is the implication
  \[
    \operatorname{IsProper}(\kappa)
      \wedge\operatorname{IsImmersion}(\kappa)
      \Longrightarrow\operatorname{IsClosedImmersion}(\kappa).
  \]
  Define the actual mixed-overlap morphism by
  \[
    \kappa_{\mathsf{CM}}:=(m_1,m_2)\colon
      M_{12}\longrightarrow X'_1\times_SX'_2.
  \]
  Substitution of the two equations in (CM1) is the chain
  \[
    \kappa_{\mathsf{CM}}\circ\theta_{12}
      =(m_1\circ\theta_{12},m_2\circ\theta_{12})
      =(k_1,k_2)=\kappa,
    \qquad
    \kappa_{\mathsf{CM}}=\kappa\circ\theta_{12}^{-1}.            \tag{CM5}
  \]
  Equation (CM5) and invariance under composition with an isomorphism are the
  predicate
  \[
    \operatorname{IsClosedImmersion}(\kappa)
      \wedge\operatorname{IsIso}(\theta_{12}^{-1})
      \Longrightarrow
        \operatorname{IsClosedImmersion}(\kappa_{\mathsf{CM}}).
  \]
  \Cref{lemma:nagata-separated-two-open-cover}, applied to
  \[
    \mathsf{CM}(X_1,X_2;U)=\lambda_1(X'_1)\cup\lambda_2(X'_2)
  \]
  with intersection $M_{12}$ and mixed-overlap morphism
  $\kappa_{\mathsf{CM}}$, is the predicate
  \[
    \operatorname{Separated}(s_{\mathsf{CM}}).
  \]
  Finite type of blow-ups and stability under composition are the predicate
  pair
  \[
    \operatorname{FiniteType}(\phi_i),\qquad
    \operatorname{FiniteType}(\phi_i)
      \wedge\operatorname{FiniteType}(s_i)
      \Longrightarrow\operatorname{FiniteType}(s_i\circ\phi_i).
  \]
  Equation (CM2h), the cover equation (CM1e), and source-locality of finite
  type form the implication
  \[
  \begin{gathered}
    s_{\mathsf{CM}}\circ\lambda_i=s_i\circ\phi_i,
    \qquad
    \operatorname{FiniteType}(s_i\circ\phi_i)\quad(i=1,2),\\
    |\mathsf{CM}(X_1,X_2;U)|
      =\lambda_1(|X'_1|)\cup\lambda_2(|X'_2|)
      \quad\Longrightarrow\quad
    \operatorname{FiniteType}(s_{\mathsf{CM}}).
  \end{gathered}
  \]
\end{proof}

\begin{lemma}[{\cite[\href{https://stacks.math.columbia.edu/tag/0F3Y}{Tag 0F3Y}]{stacks-project}}]
  \label{lemma:nagata-scheme-image-strict-transform}
  \uses{def:nagata-closed-subscheme-class,
    def:nagata-scheme-theoretic-image,def:nagata-scheme-image-factor,
    def:nagata-strict-transform,lemma:nagata-morphism-property-calculus}
  \lean{}

  Let $v\colon V\to P:=X\times_SY$ be quasi-compact and
  quasi-separated, and suppose
  \[
    \operatorname{pr}_X\circ v=j\circ u
  \]
  for $u\colon V\to U$ and an open immersion $j\colon U\hookrightarrow X$.
  The two hypotheses on $v$ and
  \Cref{def:nagata-morphism-predicate-notation} are the equivalence
  \[
    \operatorname{IsQuasiCompact}(v)\ \wedge\
      \operatorname{QuasiSeparated}(v)
      \Longleftrightarrow \operatorname{IsQCQS}(v).               \tag{SI-1}
  \]
  Only after (SI-1), define
  \[
    Z:=\operatorname{SIm}_P(v),\qquad
    i_Z:=i_v\colon Z\hookrightarrow P.                            \tag{SI-0}
  \]
  Let $(\mathcal I,X',\pi,\alpha_U)$ be a $U$-admissible blow-up and put
  \[
  \begin{gathered}
    P':=X'\times_SY,\qquad
    \pi_P:=\pi\times_S\operatorname{id}_Y\colon P'\to P,\\
    \widetilde u:=
      \operatorname{pr}_{X'}\circ\alpha_U^{-1}\circ u\colon V\to X',\\
    v':=(\widetilde u,\operatorname{pr}_Y\circ v)\colon V\to P',\\
    Z':=\operatorname{ST}_{\pi_P}(Z),\qquad
    \iota_{i_Z,\pi_P}\colon Z'\hookrightarrow Z\times_PP',\\
    i_{Z'}:=\operatorname{pr}_{P'}\circ\iota_{i_Z,\pi_P}
      \colon Z'\hookrightarrow P'.
  \end{gathered}
  \]
  Define
  \[
  \begin{gathered}
    G_{v'}:=V\times_{P,v,\pi_P}P',\qquad
    p_V\colon G_{v'}\to V,\qquad p_{P'}\colon G_{v'}\to P',\\
    \Gamma_{v'}:=(\operatorname{id}_V,v')\colon V\to G_{v'},\qquad
    p_V\circ\Gamma_{v'}=\operatorname{id}_V,\qquad
    v'=p_{P'}\circ\Gamma_{v'}.
  \end{gathered}                                                   \tag{SI0}
  \]
  By \Cref{lemma:nagata-blowup-proper} and stability under base change,
  \[
  \begin{gathered}
    \operatorname{IsProper}(\pi)
      \Longrightarrow \operatorname{IsSeparated}(\pi)
      \Longrightarrow \operatorname{IsSeparated}(\pi_P)
      \Longrightarrow \operatorname{IsSeparated}(p_V),\\
    p_V\circ\Gamma_{v'}=\operatorname{id}_V
      \ \wedge\ \operatorname{IsSeparated}(p_V)
      \Longrightarrow
      \operatorname{IsClosedImmersion}(\Gamma_{v'})
      \Longrightarrow\operatorname{IsQCQS}(\Gamma_{v'}),\\
    \operatorname{IsQCQS}(v)
      \Longrightarrow\operatorname{IsQCQS}(p_{P'}),\\
    v'=p_{P'}\circ\Gamma_{v'}
      \ \wedge\ \operatorname{IsQCQS}(\Gamma_{v'})
      \ \wedge\ \operatorname{IsQCQS}(p_{P'})
      \Longrightarrow\operatorname{IsQCQS}(v').
  \end{gathered}                                                   \tag{SI1}
  \]
  Define
  \[
    Z'':=\operatorname{SIm}_{P'}(v'),\qquad
    i_{Z''}:=i_{v'}\colon Z''\hookrightarrow P'.                  \tag{SI2}
  \]
  Then $v'$ is well-defined by
  \[
    \pi\circ\widetilde u=j\circ u
      =\operatorname{pr}_X\circ v,
  \]
  and
  \[
    [\,Z''\xhookrightarrow{i_{Z''}}P'\,]
      =[\,Z'\xhookrightarrow{i_{Z'}}P'\,]
      \quad\text{in }\operatorname{ClosedSub}(P').
  \]
\end{lemma}

\begin{proof}
  \uses{def:nagata-choice-notation,
    lemma:nagata-scheme-image-restrict-open}
  By \Cref{def:nagata-scheme-image-factor}, define
  \[
    v^{\mathrm{im}}\colon V\to Z,
    \qquad
    i_Z\circ v^{\mathrm{im}}=v.                                  \tag{SI3}
  \]
  Define the preserved-open objects and maps
  \[
  \begin{gathered}
    P'_U:=X'\times_XU\times_SY,
    \qquad
    \jmath'_U:=(\operatorname{pr}_{X'},\operatorname{pr}_Y)
      \colon P'_U\hookrightarrow P',\\
    Z'_U:=Z'\times_{P'}P'_U,
    \qquad
    \jmath'_{Z,U}:=\operatorname{pr}_{Z'}\colon Z'_U\hookrightarrow Z',\\
    i_{Z'_U}:=\operatorname{pr}_{P'_U}\colon
      Z'_U\hookrightarrow P'_U,
    \beta_U:=\alpha_U\times_S\operatorname{id}_Y
      \colon P'_U\xrightarrow{\sim}U\times_SY.
  \end{gathered}                                                   \tag{SI4}
  \]
  The component equations defining $v'$ are the equalities
  \[
  \begin{aligned}
    \pi_P\circ v'
      &=(\pi\circ\widetilde u,\operatorname{pr}_Y\circ v)\\
      &=(j\circ u,\operatorname{pr}_Y\circ v)=v.
  \end{aligned}                                                    \tag{SI5}
  \]
  Thus the fibre-product universal property defines the unique morphism
  \[
    v_U:=(u,\operatorname{pr}_Y\circ v)\colon V\to U\times_SY,
    \qquad
    v'_U\colon V\to P'_U,
    \qquad
    \jmath'_U\circ v'_U=v',
    \qquad
    \beta_U\circ v'_U=v_U.                                      \tag{SI6}
  \]
  Define the remaining projections and the preserved-open immersion by
  \[
  \begin{gathered}
    \rho_Z^P:=\operatorname{pr}_Z\colon Z\times_PP'\to Z,
    \qquad
    \rho_Z^U:=\operatorname{pr}_Z\colon
      Z\times_P(U\times_SY)\to Z,\\
    \rho_{U,Y}^Z:=\operatorname{pr}_{U\times_SY}\colon
      Z\times_P(U\times_SY)\to U\times_SY,
    \qquad
    \delta_U:=j\times_S\operatorname{id}_Y\colon U\times_SY\hookrightarrow P.
  \end{gathered}                                                    \tag{SI7a}
  \]
  Define the pullback of $v$ along $\delta_U$ and its factorization map by
  \[
  \begin{gathered}
    F_{v,U}:=V\times_{P,v,\delta_U}(U\times_SY),\qquad
    r_V\colon F_{v,U}\to V,\qquad
    r_{U,Y}\colon F_{v,U}\to U\times_SY,\\
    \rho_{v,U}:=(\operatorname{id}_V,v_U)\colon V\to F_{v,U},\qquad
    r_V\circ\rho_{v,U}=\operatorname{id}_V,\qquad
    r_{U,Y}\circ\rho_{v,U}=v_U.
  \end{gathered}                                                   \tag{SI7c}
  \]
  The pullback equality, $v=\delta_U\circ v_U$, and monicity of the open
  immersion $\delta_U$ are the cancellation chain
  \[
  \begin{aligned}
    \delta_U\circ v_U\circ r_V
      &=v\circ r_V
       =\delta_U\circ r_{U,Y}
      \Longrightarrow v_U\circ r_V=r_{U,Y},\\
    r_V\circ\rho_{v,U}\circ r_V&=r_V,\qquad
    r_{U,Y}\circ\rho_{v,U}\circ r_V=r_{U,Y},\\
    \rho_{v,U}\circ r_V&=\operatorname{id}_{F_{v,U}}
      \quad\text{by (PC10) of
        \Cref{lemma:nagata-morphism-property-calculus}},\\
    \operatorname{IsIso}(\rho_{v,U}),\qquad
    \rho_{v,U}^{-1}&=r_V.
  \end{aligned}                                                    \tag{SI7d}
  \]
  By (PC2)--(PC3) of
  \Cref{lemma:nagata-morphism-property-calculus}, (SI-1), (SI7c), and (SI7d)
  are the predicate chain
  \[
  \begin{aligned}
    \operatorname{IsQCQS}(v)
      &\Longrightarrow\operatorname{IsQCQS}(r_{U,Y}),\\
    \operatorname{IsIso}(\rho_{v,U})
      &\Longrightarrow\operatorname{IsQCQS}(\rho_{v,U}),\\
    \operatorname{IsQCQS}(r_{U,Y})\ \wedge\
      \operatorname{IsQCQS}(\rho_{v,U})\ \wedge\
      (v_U=r_{U,Y}\circ\rho_{v,U})
      &\Longrightarrow\operatorname{IsQCQS}(v_U).
  \end{aligned}                                                    \tag{SI7e}
  \]
  By \Cref{lemma:nagata-strict-transform-over-complement}, with
  $(X,Y,\pi,U)=(P,Z,\pi_P,U\times_SY)$, define the unique isomorphism
  \[
  \begin{gathered}
    \tau_{Z,U}\colon
      Z'_U\xrightarrow{\sim}Z\times_P(U\times_SY),\\
    \rho_Z^U\circ\tau_{Z,U}
      =\rho_Z^P\circ\iota_{i_Z,\pi_P}\circ\jmath'_{Z,U},
    \qquad
    \rho_{U,Y}^Z\circ\tau_{Z,U}=\beta_U\circ i_{Z'_U}.
  \end{gathered}                                                    \tag{SI7}
  \]
  Equations \textup{(SI3)}, \textup{(SI5)}, and \textup{(SI6)} are the
  following equality chain:
  \[
  \begin{aligned}
    i_Z\circ v^{\mathrm{im}}
      &=v
       =(j\circ u,\operatorname{pr}_Y\circ v)\\
      &=\delta_U\circ(u,\operatorname{pr}_Y\circ v)
       =\delta_U\circ\beta_U\circ v'_U.
  \end{aligned}                                                     \tag{SI7b}
  \]
  Hence the fibre-product universal property defines the unique map
  \[
  \begin{gathered}
    a_{v'}:=(v^{\mathrm{im}},\beta_U\circ v'_U)
      \colon V\to Z\times_P(U\times_SY),\\
    \rho_Z^U\circ a_{v'}=v^{\mathrm{im}},
    \qquad
    \rho_{U,Y}^Z\circ a_{v'}=\beta_U\circ v'_U,
    \qquad
    w_{v'}:=\tau_{Z,U}^{-1}\circ a_{v'}\colon V\to Z'_U.
  \end{gathered}                                                    \tag{SI8a}
  \]
  The second projection equality in \textup{(SI7)} is the following equality
  chain:
  \[
  \begin{aligned}
    \beta_U\circ i_{Z'_U}\circ w_{v'}
      &=\rho_{U,Y}^Z\circ\tau_{Z,U}\circ\tau_{Z,U}^{-1}\circ a_{v'}\\
      &=\rho_{U,Y}^Z\circ a_{v'}
       =\beta_U\circ v'_U,\\
    i_{Z'_U}\circ w_{v'}&=v'_U,\\
    \jmath'_U\circ i_{Z'_U}\circ w_{v'}
      &=\jmath'_U\circ v'_U=v'.
  \end{aligned}                                                     \tag{SI8}
  \]
  Put
  \[
    \overline v':=\jmath'_{Z,U}\circ w_{v'}\colon V\to Z'.       \tag{SI9}
  \]
  From \textup{(SI8)} and the pullback square in \textup{(SI4)},
  \[
    i_{Z'}\circ\overline v'
      =i_{Z'}\circ\jmath'_{Z,U}\circ w_{v'}
      =\jmath'_U\circ i_{Z'_U}\circ w_{v'}=v'.                    \tag{SI10}
  \]
  Applying \Cref{lemma:nagata-scheme-image-universal} to
  \textup{(SI10)} defines the unique closed immersion
  \[
    c\colon Z''\hookrightarrow Z',
    \qquad
    i_{Z''}=i_{Z'}\circ c.                                        \tag{SI11}
  \]
  Only after (SI7e), define the restricted closed embeddings and, by
  \Cref{def:nagata-scheme-theoretic-image}, their common scheme image in the
  ambient scheme $U\times_SY$ by
  \[
  \begin{gathered}
    Z''_U:=Z''\times_{P'}P'_U,
    \qquad i_{Z''_U}:=\operatorname{pr}_{P'_U}\colon Z''_U\hookrightarrow P'_U,\\
    c''_U:=\beta_U\circ i_{Z''_U}\colon Z''_U\hookrightarrow U\times_SY,\\
    c'_U:=\beta_U\circ i_{Z'_U}\colon Z'_U\hookrightarrow U\times_SY,\\
    M_U:=\operatorname{SIm}_{U\times_SY}(v_U),\qquad
    i_{v_U}:=i_{(v_U\colon V\to U\times_SY)}\colon
      M_U\hookrightarrow U\times_SY.
  \end{gathered}                                                   \tag{SI12a}
  \]
  The two instances of
  \Cref{lemma:nagata-scheme-image-restrict-open}, transported through
  $\beta_U$, are the class equalities
  \[
  \begin{aligned}
    [\,Z''_U\xhookrightarrow{c''_U}U\times_SY\,]
      &=[\,M_U\xhookrightarrow{i_{v_U}}U\times_SY\,],\\
    [\,Z'_U\xhookrightarrow{c'_U}U\times_SY\,]
      &=[\,M_U\xhookrightarrow{i_{v_U}}U\times_SY\,]
      \quad\text{in }\operatorname{ClosedSub}(U\times_SY).
  \end{aligned}                                                    \tag{SI12b}
  \]
  Define
  \[
  \begin{aligned}
    \mathscr T''_U
      &:=\left\{\theta\in\operatorname{Hom}(Z''_U,M_U)
        \ \middle|\
        \operatorname{IsIso}(\theta)\ \wedge\
        i_{v_U}\circ\theta=c''_U\right\},\\
    \mathscr T'_U
      &:=\left\{\theta\in\operatorname{Hom}(Z'_U,M_U)
        \ \middle|\
        \operatorname{IsIso}(\theta)\ \wedge\
        i_{v_U}\circ\theta=c'_U\right\}.
  \end{aligned}
  \]
  By (CSub3) of \Cref{def:nagata-closed-subscheme-class}, (SI12b) is the
  following pair of nonemptiness predicates:
  \[
    \mathscr T''_U\ne\varnothing,
    \qquad
    \mathscr T'_U\ne\varnothing.
  \]
  Define the chosen witnesses
  \[
  \begin{gathered}
    \theta''_U:=\operatorname{ch}(\mathscr T''_U)\in\mathscr T''_U,
    \qquad
    \theta'_U:=\operatorname{ch}(\mathscr T'_U)\in\mathscr T'_U,\\
    i_{v_U}\circ\theta''_U=c''_U,
    \qquad
    i_{v_U}\circ\theta'_U=c'_U.
  \end{gathered}                                                    \tag{SI12c}
  \]
  Define
  \[
    c_U:=c\times_{P'}\operatorname{id}_{P'_U}
      \colon Z''_U\longrightarrow Z'_U.
  \]
  The equality $i_{Z'_U}\circ c_U=i_{Z''_U}$ and (SI12a)--(SI12c) are the
  following equality chain:
  \[
  \begin{aligned}
    i_{v_U}\circ\theta'_U\circ c_U
      &=c'_U\circ c_U
       =\beta_U\circ i_{Z'_U}\circ c_U\\
      &=\beta_U\circ i_{Z''_U}
       =c''_U
       =i_{v_U}\circ\theta''_U,\\
    \theta'_U\circ c_U&=\theta''_U,\\
    c_U&=(\theta'_U)^{-1}\circ\theta''_U,
    \qquad \operatorname{IsIso}(c_U).
  \end{aligned}                                                     \tag{SI12}
  \]

  On a strict-transform chart, write
  \[
    Z'|_{\operatorname{Spec}(C)}
      =\operatorname{Spec}(R),\qquad
    R=C/\Gamma_{(a)}(C),
  \]
  as in \Cref{def:nagata-strict-transform-chart-ring}, and write
  \[
    Z''|_{\operatorname{Spec}(C)}=\operatorname{Spec}(R/K)
  \]
  for the uniquely determined ideal $K\triangleleft R$.  The last predicate
  in (SI12) says that
  $\operatorname{Spec}(R_a/K_a)\hookrightarrow\operatorname{Spec}(R_a)$ is
  an isomorphism.  The affine closed-immersion criterion is the equation
  \[
    K_a=0.
  \]
  For $k\in K$, the localization zero criterion is the implication chain
  \[
    k/1=0\text{ in }R_a
      \Longrightarrow \exists n\in\mathbf N,\ a^nk=0
      \Longrightarrow k\in\Gamma_{(a)}(R).
  \]
  By \Cref{lemma:nagata-power-torsion-quotient-torsionfree},
  \[
    \Gamma_{(a)}(R)
      =\Gamma_{(a)}(C/\Gamma_{(a)}(C))=0.
  \]
  For every $k\in K$, the preceding equalities form the chain
  \[
    k\in K
      \Longrightarrow k/1=0\text{ in }R_a
      \Longrightarrow k\in\Gamma_{(a)}(R)=0
      \Longrightarrow k=0,
    \qquad
    K=(0),
    \qquad
    \operatorname{Spec}(R/K)=\operatorname{Spec}(R).              \tag{SI13}
  \]
  Let $(\mathcal C_\lambda)_{\lambda\in\Lambda}$ denote the
  strict-transform chart cover of $Z'$.  The isomorphism-locality criterion
  applied to \textup{(SI13)} is the implication
  \[
    \left(\bigcup_{\lambda\in\Lambda}\mathcal C_\lambda=Z'
      \ \wedge\
      \bigwedge_{\lambda\in\Lambda}
        \operatorname{IsIso}(c|_{c^{-1}(\mathcal C_\lambda)})\right)
      \Longrightarrow\operatorname{IsIso}(c),
    \qquad
    \operatorname{IsIso}(c)\ \wedge\ i_{Z''}=i_{Z'}\circ c
      \Longrightarrow
      [\,Z''\xhookrightarrow{i_{Z''}}P'\,]
        =[\,Z'\xhookrightarrow{i_{Z'}}P'\,]
        \quad\text{in }\operatorname{ClosedSub}(P').             \tag{SI14}
  \]
\end{proof}

\begin{definition}[Graph-extension construction]
  \label{def:nagata-graph-extension}
  \uses{def:nagata-compactification,
    def:nagata-scheme-theoretic-image,
    def:nagata-elimination-indeterminacy,
    def:nagata-strict-transform,
    lemma:nagata-morphism-property-calculus}
  \lean{}

  Let $u\colon U\to S$ be separated and of finite type, let
  $\iota\colon V\hookrightarrow U$ be a quasi-compact open immersion, and
  fix a compactification
  \[
    j\colon V\hookrightarrow Y,\qquad
    p\colon Y\to S,\qquad
    p\circ j=u\circ\iota,
  \]
  with $p$ proper.  Define
  \[
  \begin{gathered}
    \gamma_V:=(j,\iota)\colon V\to Y\times_SU,\\
    s_{\gamma_V}:=(\operatorname{id}_V,\iota)
      \colon V\to V\times_SU,\\
    \jmath_{V,U}:=j\times_S\operatorname{id}_U
      \colon V\times_SU\hookrightarrow Y\times_SU,\\
    \operatorname{pr}_V\circ s_{\gamma_V}=\operatorname{id}_V,
    \qquad
    \gamma_V=\jmath_{V,U}\circ s_{\gamma_V}.
  \end{gathered}                                                   \tag{GX0}
  \]
  The separatedness of $u$, the quasi-compact-open predicate for $j$ in
  \Cref{def:nagata-compactification}, and the property calculus in
  \Cref{lemma:nagata-morphism-property-calculus} are the following predicate
  chain:
  \[
  \begin{gathered}
    \operatorname{IsSeparated}(u)
      \Longrightarrow\operatorname{IsSeparated}
        (\operatorname{pr}_V\colon V\times_SU\to V),\\
    \operatorname{pr}_V\circ s_{\gamma_V}=\operatorname{id}_V
      \ \wedge\ \operatorname{IsSeparated}(\operatorname{pr}_V)
      \Longrightarrow\operatorname{IsClosedImmersion}(s_{\gamma_V})
      \Longrightarrow\operatorname{IsQCQS}(s_{\gamma_V}),\\
    \operatorname{IsOpenImmersion}(j)\ \wedge\
      \operatorname{QuasiCompact}(j)
      \Longrightarrow
      \operatorname{IsOpenImmersion}(\jmath_{V,U})\ \wedge\
      \operatorname{QuasiCompact}(\jmath_{V,U})
      \Longrightarrow\operatorname{IsQCQS}(\jmath_{V,U}),\\
    \operatorname{IsQCQS}(s_{\gamma_V})\ \wedge\
      \operatorname{IsQCQS}(\jmath_{V,U})\ \wedge\
      \gamma_V=\jmath_{V,U}\circ s_{\gamma_V}
      \Longrightarrow\operatorname{IsQCQS}(\gamma_V).
  \end{gathered}                                                   \tag{GX1}
  \]
  Only after (GX1), define
  \[
  \begin{gathered}
    Z_V:=\operatorname{SIm}(\gamma_V),\qquad
    c_V:=i_{\gamma_V}\colon Z_V\hookrightarrow Y\times_SU,\\
    p_1:=\operatorname{pr}_Y\circ c_V\colon Z_V\to Y,\qquad
    p_2:=\operatorname{pr}_U\circ c_V\colon Z_V\to U.
  \end{gathered}
  \]
  Since $c_V$ is a closed immersion and
  $\operatorname{pr}_Y\colon Y\times_SU\to Y$ is the base change of $u$,
  \[
    \operatorname{Separated}(p_1),\qquad
    \operatorname{FiniteType}(p_1).
  \]
  Define
  \[
    \eta_V:=\zeta_{\gamma_V}^{-1}
      \colon Z_V\times_YV\xrightarrow{\sim}V,
    \qquad
    j\circ\eta_V=p_1\circ\operatorname{pr}_{Z_V},
  \]
  where $\zeta_{\gamma_V}$ is
  \Cref{lemma:nagata-scheme-image-partial-section}.  Apply
  \Cref{def:nagata-elimination-indeterminacy} to the datum
  $(p_1,j,\eta_V)$ and define
  \[
  \begin{gathered}
    \mathcal I_{\mathrm{gr}}:=\mathcal F_V(p_1,\eta_V),\qquad
    Y^{\mathrm{gr}}:=X^{\mathrm{el}}(p_1,V)
      =\operatorname{Bl}_{\mathcal I_{\mathrm{gr}}}(Y),\\
    \pi_{\mathrm{gr}}:=b_{Y,\mathcal I_{\mathrm{gr}}}
      \colon Y^{\mathrm{gr}}\to Y,\\
    V^{\mathrm{gr}}:=Z^{\mathrm{el}}(p_1,V)
      =\operatorname{ST}_{\pi_{\mathrm{gr}}}(Z_V),\\
    e_{\mathrm{gr}}:=p_1^{\mathrm{el}}\colon
      V^{\mathrm{gr}}\hookrightarrow Y^{\mathrm{gr}}.
  \end{gathered}
  \]
  Define the closed immersions
  \[
  \begin{gathered}
    \iota_{\mathrm{gr}}:=\iota_{p_1,\pi_{\mathrm{gr}}}
      \colon V^{\mathrm{gr}}\hookrightarrow Z_V\times_YY^{\mathrm{gr}},\\
    \widehat c_{\mathrm{gr}}:=
      c_V\times_Y\operatorname{id}_{Y^{\mathrm{gr}}}
      \colon Z_V\times_YY^{\mathrm{gr}}\hookrightarrow
        (Y\times_SU)\times_YY^{\mathrm{gr}}
        =Y^{\mathrm{gr}}\times_SU,\\
    \widetilde c_{\mathrm{gr}}:=
      \widehat c_{\mathrm{gr}}\circ\iota_{\mathrm{gr}}
      \colon V^{\mathrm{gr}}\hookrightarrow Y^{\mathrm{gr}}\times_SU.
  \end{gathered}
  \]
  Finally define
  \[
    \psi^{\mathrm{gr}}
      :=\operatorname{pr}_U\circ\widetilde c_{\mathrm{gr}}
      \colon V^{\mathrm{gr}}\to U.
  \]
\end{definition}

\begin{lemma}[Properties of the graph-extension construction,
  {\cite[\href{https://stacks.math.columbia.edu/tag/0F3Z}{Tag 0F3Z}]{stacks-project}}]
  \label{lemma:nagata-graph-extension}
  \uses{def:nagata-graph-extension}
  \lean{}

  The morphism $\pi_{\mathrm{gr}}$ is $V$-admissible,
  $e_{\mathrm{gr}}\colon V^{\mathrm{gr}}\hookrightarrow Y^{\mathrm{gr}}$ is
  open, and
  $\psi^{\mathrm{gr}}\colon V^{\mathrm{gr}}\to U$ is proper.  Its
  restriction along the open immersion
  $\widetilde j\colon V\hookrightarrow V^{\mathrm{gr}}$ from
  \Cref{lemma:nagata-strict-transform-over-complement} satisfies
  \[
    \psi^{\mathrm{gr}}\circ\widetilde j=\iota.
  \]
  Moreover, with
  $V^{\mathrm{gr}}_V:=V^{\mathrm{gr}}\times_UV$, define the typed
  fibre-product morphism
  \[
    \widetilde j_V:=(\widetilde j,\operatorname{id}_V)
      \colon V\xrightarrow{\sim}V^{\mathrm{gr}}_V,\qquad
    \operatorname{pr}_{V^{\mathrm{gr}}}\circ\widetilde j_V
        =\widetilde j,\qquad
    \operatorname{pr}_V\circ\widetilde j_V=\operatorname{id}_V.
  \]
\end{lemma}

\begin{proof}
  \uses{lemma:nagata-scheme-image-partial-section,
    lemma:nagata-scheme-image-strict-transform,
    lemma:nagata-scheme-image-universal,
    lemma:nagata-scheme-image-restrict-open,
    lemma:nagata-strict-transform-over-complement,
    lemma:nagata-elimination-indeterminacy,
    lemma:nagata-morphism-property-calculus,
    lemma:nagata-blowup-proper}
  The instance of \Cref{lemma:nagata-scheme-image-partial-section} at
  $\gamma_V=(j,\iota)$ is the typed isomorphism
  \[
    \eta_V\colon Z_V\times_YV\xrightarrow{\sim}V.
  \]
  The conclusion of \Cref{lemma:nagata-elimination-indeterminacy}, applied to
  $(p_1,V,\eta_V)$, is
  \[
    \pi_{\mathrm{gr}}\text{ is }V\text{-admissible},
    \qquad
    e_{\mathrm{gr}}\colon V^{\mathrm{gr}}\hookrightarrow
      Y^{\mathrm{gr}}\text{ is open}.
  \]

  Define the structure morphism, the preserved-open base changes, and their
  projections by
  \[
  \begin{gathered}
    p^{\mathrm{gr}}:=p\circ\pi_{\mathrm{gr}}\colon
      Y^{\mathrm{gr}}\to S,\\
    Y^{\mathrm{gr}}_V:=Y^{\mathrm{gr}}\times_{Y,(\pi_{\mathrm{gr}},j)}V,
    \qquad
    \alpha_V^{\mathrm{gr}}\colon Y^{\mathrm{gr}}_V\xrightarrow{\sim}V,\\
    \jmath_V^{\mathrm{gr}}:=\operatorname{pr}_{Y^{\mathrm{gr}}}
      \colon Y^{\mathrm{gr}}_V\hookrightarrow Y^{\mathrm{gr}},
    \qquad
    j\circ\alpha_V^{\mathrm{gr}}
      =\pi_{\mathrm{gr}}\circ\jmath_V^{\mathrm{gr}},\\
    V^{\mathrm{gr}}_{/V}:=
      V^{\mathrm{gr}}\times_{Y^{\mathrm{gr}}}Y^{\mathrm{gr}}_V,
    \qquad
    q_{\mathrm{gr}}\colon V^{\mathrm{gr}}_{/V}\to V^{\mathrm{gr}},\\
    q_Y\colon V^{\mathrm{gr}}_{/V}\to Y^{\mathrm{gr}}_V,\\
    r_{\mathrm{gr}}:=
      \alpha_V^{\mathrm{st}}(p_1,\pi_{\mathrm{gr}})
      \colon V^{\mathrm{gr}}_{/V}\xrightarrow{\sim}Z_V\times_YV.
  \end{gathered} \tag{GE0}
  \]
  Here $\alpha_V^{\mathrm{gr}}$ is the complement isomorphism of
  \Cref{def:nagata-admissible-blowup}, and $r_{\mathrm{gr}}$ is the displayed
  instance of \Cref{lemma:nagata-strict-transform-over-complement}.  Define
  \[
  \begin{gathered}
    \widetilde j_Y:=
      \jmath_V^{\mathrm{gr}}\circ(\alpha_V^{\mathrm{gr}})^{-1}
      \colon V\to Y^{\mathrm{gr}},
    \qquad
    \pi_{\mathrm{gr}}\circ\widetilde j_Y=j,\\
    \widetilde j:=q_{\mathrm{gr}}\circ r_{\mathrm{gr}}^{-1}
      \circ\eta_V^{-1}\colon V\to V^{\mathrm{gr}},\\
    \sigma_U^{\mathrm{gr}}:=(\widetilde j_Y,\iota)
      \colon V\to Y^{\mathrm{gr}}\times_SU.                       
  \end{gathered} \tag{GE0a}
  \]
  Base-change and composition stability in
  \Cref{lemma:nagata-morphism-property-calculus} are the following predicate
  chain:
  \[
  \begin{gathered}
    \operatorname{IsOpenImmersion}(j)
      \Longrightarrow
      \operatorname{IsOpenImmersion}(\jmath_V^{\mathrm{gr}})
      \Longrightarrow\operatorname{IsOpenImmersion}(q_{\mathrm{gr}}),\\
    \operatorname{IsIso}(r_{\mathrm{gr}}^{-1})\ \wedge\
      \operatorname{IsIso}(\eta_V^{-1})
      \Longrightarrow
      \operatorname{IsOpenImmersion}(r_{\mathrm{gr}}^{-1})\ \wedge\
      \operatorname{IsOpenImmersion}(\eta_V^{-1}),\\
    \operatorname{IsOpenImmersion}(q_{\mathrm{gr}})\ \wedge\
      \operatorname{IsOpenImmersion}(r_{\mathrm{gr}}^{-1})\ \wedge\
      \operatorname{IsOpenImmersion}(\eta_V^{-1})
      \Longrightarrow\operatorname{IsOpenImmersion}(\widetilde j).
  \end{gathered} \tag{GE0a1}
  \]
  The last ordered pair is defined because
  \[
  \begin{aligned}
    p^{\mathrm{gr}}\circ\widetilde j_Y
      &=p\circ\pi_{\mathrm{gr}}\circ\widetilde j_Y
       =p\circ j
       =u\circ\iota.
  \end{aligned}                                                  \tag{GE0b}
  \]
  The projection identities in
  \Cref{lemma:nagata-strict-transform-over-complement}, together with the two
  component identities of $\eta_V$, are the following explicit component
  chains:
  \[
  \begin{aligned}
    e_{\mathrm{gr}}\circ\widetilde j
      &=e_{\mathrm{gr}}\circ q_{\mathrm{gr}}
          \circ r_{\mathrm{gr}}^{-1}\circ\eta_V^{-1}
       =\widetilde j_Y,\\
    \operatorname{pr}_U\circ\widetilde c_{\mathrm{gr}}
        \circ\widetilde j
      &=p_2\circ\operatorname{pr}_{Z_V}
          \circ r_{\mathrm{gr}}\circ r_{\mathrm{gr}}^{-1}
          \circ\eta_V^{-1}
       =\iota\circ\eta_V\circ\eta_V^{-1}
       =\iota.
  \end{aligned}                                                  \tag{GE0c}
  \]
  Since
  $\operatorname{pr}_{Y^{\mathrm{gr}}}\circ\widetilde c_{\mathrm{gr}}
    =e_{\mathrm{gr}}$ by
  \Cref{def:nagata-graph-extension}, the two component chains in (GE0c) are
  the following component equalities:
  \[
    \widetilde c_{\mathrm{gr}}\circ\widetilde j
      =(\widetilde j_Y,\iota)=\sigma_U^{\mathrm{gr}},
    \qquad
    \psi^{\mathrm{gr}}\circ\widetilde j=\iota.                  \tag{GE0d}
  \]

  \[
    \kappa_U^{\mathrm{gr}}:=
      \widetilde j_Y\times_S\operatorname{id}_U
      \colon V\times_SU\to Y^{\mathrm{gr}}\times_SU,
    \qquad
    \sigma_U^{\mathrm{gr}}
      =\kappa_U^{\mathrm{gr}}\circ s_{\gamma_V}.                 \tag{GE0q0}
  \]
  The qcqs predicate (GX1), base-change stability, and composition stability
  in \Cref{lemma:nagata-morphism-property-calculus} form the chain
  \[
  \begin{gathered}
    \operatorname{IsOpenImmersion}(j)\ \wedge\
      \operatorname{QuasiCompact}(j)
      \Longrightarrow\operatorname{IsQCQS}(\jmath_V^{\mathrm{gr}}),\\
    \operatorname{IsIso}((\alpha_V^{\mathrm{gr}})^{-1})
      \Longrightarrow\operatorname{IsQCQS}
        ((\alpha_V^{\mathrm{gr}})^{-1}),\\
    \operatorname{IsQCQS}(\jmath_V^{\mathrm{gr}})\ \wedge\
      \operatorname{IsQCQS}((\alpha_V^{\mathrm{gr}})^{-1})\ \wedge\
      \widetilde j_Y
        =\jmath_V^{\mathrm{gr}}\circ(\alpha_V^{\mathrm{gr}})^{-1}
      \Longrightarrow\operatorname{IsQCQS}(\widetilde j_Y),\\
    \operatorname{IsQCQS}(\widetilde j_Y)
      \Longrightarrow\operatorname{IsQCQS}(\kappa_U^{\mathrm{gr}}),\\
    \operatorname{IsQCQS}(\kappa_U^{\mathrm{gr}})\ \wedge\
      \operatorname{IsQCQS}(s_{\gamma_V})\ \wedge\
      \sigma_U^{\mathrm{gr}}
        =\kappa_U^{\mathrm{gr}}\circ s_{\gamma_V}
      \Longrightarrow\operatorname{IsQCQS}(\sigma_U^{\mathrm{gr}}).
  \end{gathered}                                                   \tag{GE0q}
  \]
  Only after (GE0q), define the scheme image by
  \[
  \begin{gathered}
    M_U^{\mathrm{gr}}:=
      \operatorname{SIm}_{Y^{\mathrm{gr}}\times_SU}
        (\sigma_U^{\mathrm{gr}}),\\
    i_U^{\mathrm{gr}}:=i_{\sigma_U^{\mathrm{gr}}}\colon
      M_U^{\mathrm{gr}}\hookrightarrow Y^{\mathrm{gr}}\times_SU.
  \end{gathered} \tag{GE0e}
  \]
  In \Cref{lemma:nagata-scheme-image-strict-transform}, make the substitutions
  \[
  \begin{gathered}
    (X,Y,P,v,U,j,u,\pi)
      =(Y,U,Y\times_SU,\gamma_V,V,j,\operatorname{id}_V,
        \pi_{\mathrm{gr}}),\\
    v'=\sigma_U^{\mathrm{gr}},\qquad
    Z'=V^{\mathrm{gr}},\qquad i_{Z'}=\widetilde c_{\mathrm{gr}}.
  \end{gathered}
  \]
  Its conclusion is the class equality
  \[
    [\,V^{\mathrm{gr}}
      \xhookrightarrow{\ \widetilde c_{\mathrm{gr}}\ }
        Y^{\mathrm{gr}}\times_SU\,]
      =[\,M_U^{\mathrm{gr}}
      \xhookrightarrow{\ i_U^{\mathrm{gr}}\ }
        Y^{\mathrm{gr}}\times_SU\,]
      \quad\text{in }\operatorname{ClosedSub}(Y^{\mathrm{gr}}\times_SU).
                                                                    \tag{GE0f}
  \]

  By \Cref{lemma:nagata-blowup-proper} and stability under composition and
  base change,
  \[
  \begin{gathered}
    \operatorname{IsProper}(\pi_{\mathrm{gr}})
      \ \wedge\ \operatorname{IsProper}(p)
      \Longrightarrow\operatorname{IsProper}(p^{\mathrm{gr}}),\\
    \operatorname{pr}_U
      =p^{\mathrm{gr}}\times_S\operatorname{id}_U
      \colon Y^{\mathrm{gr}}\times_SU\to U,\\
    \operatorname{IsProper}(p^{\mathrm{gr}})
      \Longrightarrow\operatorname{IsProper}(\operatorname{pr}_U).
  \end{gathered}                                                    \tag{GE1}
  \]
  Since $\widetilde c_{\mathrm{gr}}$ is the closed immersion defined in
  \Cref{def:nagata-graph-extension},
  \[
    \operatorname{IsClosedImmersion}(\widetilde c_{\mathrm{gr}})
      \Longrightarrow\operatorname{IsProper}(\widetilde c_{\mathrm{gr}}),
    \qquad
    \psi^{\mathrm{gr}}
      =\operatorname{pr}_U\circ\widetilde c_{\mathrm{gr}}.        \tag{GE2}
  \]
  Stability under composition is the implication
  \[
    \operatorname{IsProper}(\widetilde c_{\mathrm{gr}})
      \ \wedge\ 
    \operatorname{IsProper}(\operatorname{pr}_U)
      \Longrightarrow
    \operatorname{IsProper}
      (\operatorname{pr}_U\circ\widetilde c_{\mathrm{gr}})
      =\operatorname{IsProper}(\psi^{\mathrm{gr}}).
  \]

  Define the open target morphism and the two morphisms into its source and
  target by
  \[
  \begin{gathered}
    w_V:=\operatorname{id}_{Y^{\mathrm{gr}}}\times_S\iota
      \colon Y^{\mathrm{gr}}\times_SV
        \hookrightarrow Y^{\mathrm{gr}}\times_SU,\\
    r_V:=\operatorname{pr}_V\colon Y^{\mathrm{gr}}\times_SV\to V,
    \qquad
    \sigma_V:=(\widetilde j_Y,\operatorname{id}_V)
      \colon V\to Y^{\mathrm{gr}}\times_SV,\\
    w_V\circ\sigma_V=\sigma_U^{\mathrm{gr}},\qquad
    r_V\circ\sigma_V=\operatorname{id}_V.                       
  \end{gathered} \tag{GE3}
  \]
  Equation (GE1), separatedness of proper morphisms, base-change stability,
  and the closed-section criterion are the following predicate chain:
  \[
  \begin{gathered}
    \operatorname{IsProper}(p^{\mathrm{gr}})
      \Longrightarrow\operatorname{IsSeparated}(p^{\mathrm{gr}}),\\
    \operatorname{IsSeparated}(p^{\mathrm{gr}})
      \Longrightarrow\operatorname{IsSeparated}(r_V),\\
    r_V\circ\sigma_V=\operatorname{id}_V\ \wedge\
      \operatorname{IsSeparated}(r_V)
      \Longrightarrow\operatorname{IsClosedImmersion}(\sigma_V)
      \Longrightarrow\operatorname{IsQCQS}(\sigma_V).
  \end{gathered}                                                   \tag{GE3q}
  \]
  The square
  \[
  \begin{array}{ccc}
    V & \xrightarrow{\ \sigma_V\ } & Y^{\mathrm{gr}}\times_SV\\
    {\scriptstyle\operatorname{id}_V}\big\downarrow
      &&\big\downarrow{\scriptstyle w_V}\\
    V & \xrightarrow{\ \sigma_U^{\mathrm{gr}}\ }
      & Y^{\mathrm{gr}}\times_SU
  \end{array}                                                     \tag{GE3a}
  \]
  is cartesian: for $a\colon T\to V$ and
  $b\colon T\to Y^{\mathrm{gr}}\times_SV$, the equality
  $\sigma_U^{\mathrm{gr}}\circ a=w_V\circ b$ has the component chain
  \[
  \begin{gathered}
    \iota\circ a=\iota\circ\operatorname{pr}_V\circ b
      \Longrightarrow a=\operatorname{pr}_V\circ b,\\
    \widetilde j_Y\circ a
      =\operatorname{pr}_{Y^{\mathrm{gr}}}\circ b,\\
    b=(\widetilde j_Y\circ a,a)=\sigma_V\circ a,                
  \end{gathered} \tag{GE3b}
  \]
  where the first implication cancels the open immersion $\iota$.

  Put
  \[
  \begin{gathered}
    V^{\mathrm{gr}}_V:=V^{\mathrm{gr}}\times_{U,(\psi^{\mathrm{gr}},\iota)}V,
    \qquad
    c_V^{\mathrm{gr}}:=
      (e_{\mathrm{gr}}\circ\operatorname{pr}_{V^{\mathrm{gr}}},
        \operatorname{pr}_V)
      \colon V^{\mathrm{gr}}_V\hookrightarrow
        Y^{\mathrm{gr}}\times_SV,\\
    M_V^{\mathrm{gr}}:=
      \operatorname{SIm}_{Y^{\mathrm{gr}}\times_SV}(\sigma_V),
    \qquad
    i_V^{\mathrm{gr}}:=i_{\sigma_V}\colon
      M_V^{\mathrm{gr}}\hookrightarrow Y^{\mathrm{gr}}\times_SV,\\
    \sigma_V^{\mathrm{im}}\colon V\to M_V^{\mathrm{gr}},\qquad
    i_V^{\mathrm{gr}}\circ\sigma_V^{\mathrm{im}}=\sigma_V.     
  \end{gathered} \tag{GE3c}
  \]
  The first closed immersion in (GE3c) is the base change of
  $\widetilde c_{\mathrm{gr}}$ along $w_V$.  Applying
  \Cref{lemma:nagata-scheme-image-restrict-open} to the cartesian square
  (GE3a), together with the substitution of (GE0f), is the following class
  equality:
  \[
    [\,V^{\mathrm{gr}}_V
      \xhookrightarrow{\ c_V^{\mathrm{gr}}\ }
        Y^{\mathrm{gr}}\times_SV\,]
      =[\,M_V^{\mathrm{gr}}
      \xhookrightarrow{\ i_V^{\mathrm{gr}}\ }
        Y^{\mathrm{gr}}\times_SV\,]
      \quad\text{in }\operatorname{ClosedSub}(Y^{\mathrm{gr}}\times_SV).
                                                                    \tag{GE3d}
  \]
  Define
  \[
    \mathscr T_V:=\left\{\theta\in
      \operatorname{Hom}(V^{\mathrm{gr}}_V,M_V^{\mathrm{gr}})
      \ \middle|\ 
      \operatorname{IsIso}(\theta)\ \wedge\
      i_V^{\mathrm{gr}}\circ\theta=c_V^{\mathrm{gr}}\right\}.
  \]
  By (CSub3) of \Cref{def:nagata-closed-subscheme-class}, (GE3d) is
  $\mathscr T_V\neq\varnothing$.  Define
  \[
    \theta_V:=\operatorname{ch}(\mathscr T_V)\in\mathscr T_V,
    \qquad
    \theta_V\colon V^{\mathrm{gr}}_V\xrightarrow{\sim}M_V^{\mathrm{gr}},
    \qquad
    i_V^{\mathrm{gr}}\circ\theta_V=c_V^{\mathrm{gr}}.          \tag{GE3e}
  \]
  If $\theta\in\mathscr T_V$, the monicity of the closed immersion
  $i_V^{\mathrm{gr}}$ is the cancellation implication
  \[
    \operatorname{IsClosedImmersion}(i_V^{\mathrm{gr}})
      \ \wedge\
    \bigl(i_V^{\mathrm{gr}}\circ\theta
      =c_V^{\mathrm{gr}}
      =i_V^{\mathrm{gr}}\circ\theta_V\bigr)
      \Longrightarrow \theta=\theta_V.
  \]

  Since $r_V$ is the base change of $p^{\mathrm{gr}}$, the section criterion
  is the predicate chain
  \[
  \begin{gathered}
    \operatorname{IsSeparated}(p^{\mathrm{gr}})
      \Longrightarrow\operatorname{IsSeparated}(r_V),\\
    r_V\circ\sigma_V=\operatorname{id}_V
      \ \wedge\ \operatorname{IsSeparated}(r_V)
      \Longrightarrow\operatorname{IsClosedImmersion}(\sigma_V).
  \end{gathered}                                                   \tag{GE3f}
  \]
  Apply \Cref{lemma:nagata-scheme-image-universal} to the closed
  factorization
  \[
    \sigma_V=\sigma_V\circ\operatorname{id}_V.
  \]
  With \Cref{def:nagata-scheme-image-comparison}, this defines
  \[
    \epsilon_V:=c_{\sigma_V,\sigma_V}\colon M_V^{\mathrm{gr}}\to V,
    \qquad
    \sigma_V\circ\epsilon_V=i_V^{\mathrm{gr}},\qquad
    \epsilon_V\circ\sigma_V^{\mathrm{im}}=\operatorname{id}_V. \tag{GE3g}
  \]
  Monomorphism cancellation for $i_V^{\mathrm{gr}}$ is the chain
  \[
  \begin{aligned}
    i_V^{\mathrm{gr}}\circ\sigma_V^{\mathrm{im}}\circ\epsilon_V
      &=\sigma_V\circ\epsilon_V=i_V^{\mathrm{gr}}\\
    &\Longrightarrow
      \sigma_V^{\mathrm{im}}\circ\epsilon_V
        =\operatorname{id}_{M_V^{\mathrm{gr}}}.
  \end{aligned}                                                  \tag{GE3h}
  \]
  Equations (GE3g)--(GE3h) are the inverse identities
  \[
    \epsilon_V\circ\sigma_V^{\mathrm{im}}=\operatorname{id}_V,
    \qquad
    \sigma_V^{\mathrm{im}}\circ\epsilon_V
      =\operatorname{id}_{M_V^{\mathrm{gr}}},
  \]
  equivalently
  $\operatorname{IsIso}(\epsilon_V)$ and
  $\epsilon_V^{-1}=\sigma_V^{\mathrm{im}}$ by
  \Cref{def:nagata-morphism-predicate-notation}.

  Equation (GE0d) defines the fibre-product morphism
  \[
    \widetilde j_V:=(\widetilde j,\operatorname{id}_V)
      \colon V\to V^{\mathrm{gr}}_V,\qquad
    \operatorname{pr}_{V^{\mathrm{gr}}}\circ\widetilde j_V
      =\widetilde j,\qquad
    \operatorname{pr}_V\circ\widetilde j_V=\operatorname{id}_V. \tag{GE4}
  \]
  Equations (GE0c), (GE3), and (GE4) are the following equality chain:
  \[
    c_V^{\mathrm{gr}}\circ\widetilde j_V
      =(e_{\mathrm{gr}}\circ\widetilde j,\operatorname{id}_V)
      =(\widetilde j_Y,\operatorname{id}_V)=\sigma_V.           \tag{GE4a}
  \]
  Equations (GE3c), (GE3e), and (GE4a), together with cancellation of the
  closed immersion $i_V^{\mathrm{gr}}$, are the following equality chain:
  \[
  \begin{aligned}
    i_V^{\mathrm{gr}}\circ\theta_V\circ\widetilde j_V
      &=c_V^{\mathrm{gr}}\circ\widetilde j_V
       =\sigma_V
       =i_V^{\mathrm{gr}}\circ\sigma_V^{\mathrm{im}}\\
    &\Longrightarrow
      \theta_V\circ\widetilde j_V=\sigma_V^{\mathrm{im}}.
  \end{aligned}                                                  \tag{GE4b}
  \]
  Define
  \[
    \vartheta_V:=\epsilon_V\circ\theta_V
      \colon V^{\mathrm{gr}}_V\xrightarrow{\sim}V.              \tag{GE4c}
  \]
  The inverse equations are now the explicit chains
  \[
  \begin{aligned}
    \vartheta_V\circ\widetilde j_V
      &=\epsilon_V\circ\sigma_V^{\mathrm{im}}
       =\operatorname{id}_V,\\
    \theta_V\circ\widetilde j_V\circ\vartheta_V
      &=\sigma_V^{\mathrm{im}}\circ\epsilon_V\circ\theta_V
       =\theta_V
      \Longrightarrow
      \widetilde j_V\circ\vartheta_V
       =\operatorname{id}_{V^{\mathrm{gr}}_V},
  \end{aligned}                                                  \tag{GE4d}
  \]
  where the last implication cancels the isomorphism $\theta_V$.  Thus
  $\widetilde j_V=\vartheta_V^{-1}$, and (GE0d) is the required equality
  $\psi^{\mathrm{gr}}\circ\widetilde j=\iota$.
\end{proof}


\section{The dense-source valuative criterion}

\begin{lemma}[Factorization through a closed subscheme over a valuation ring]
  \label{lemma:nagata-valuation-factor-closed}
  \uses{def:nagata-direct-image-structure-map}
  \lean{}

  Let $A$ be a valuation ring, let $K=\operatorname{Frac}(A)$, and define
  \[
    \nu_A:=\operatorname{Spec}(A\hookrightarrow K)
      \colon\operatorname{Spec}(K)\to\operatorname{Spec}(A).
  \]
  Let $i\colon Z\hookrightarrow Y$ be a closed
  immersion, and let $g\colon\operatorname{Spec}(A)\to Y$.  If there is a
  morphism $g_K\colon\operatorname{Spec}(K)\to Z$ and a specified equality
  \[
    e_K\colon g\circ\nu_A=i\circ g_K,
  \]
  then there is a unique morphism
  $\widetilde g\colon\operatorname{Spec}(A)\to Z$ satisfying
  $i\circ\widetilde g=g$.
\end{lemma}

\begin{proof}
  By \Cref{def:nagata-direct-image-structure-map}, define
  \[
  \begin{gathered}
    \varepsilon_i:=i^\flat\colon
      \mathcal O_Y\twoheadrightarrow i_*\mathcal O_Z,
    \qquad
    \mathcal I:=\ker(\varepsilon_i),
    \qquad
    \iota_{\mathcal I}\colon\mathcal I\hookrightarrow\mathcal O_Y,\\
    \varepsilon_g:=g^\flat\colon
      \mathcal O_Y\longrightarrow g_*\mathcal O_{\operatorname{Spec}(A)}.
  \end{gathered}                                                   \tag{VF0}
  \]
  For a morphism $u\colon T\to Y$, define the pullback structure map and
  the pullback ideal map by
  \[
  \begin{gathered}
    m_u\colon
      u^*\mathcal O_Y
      =u^{-1}\mathcal O_Y\otimes_{u^{-1}\mathcal O_Y}\mathcal O_T
        \longrightarrow\mathcal O_T,
    \qquad
    m_u(s\otimes a):=u^\sharp(s)a,\\
    \mu_u^{\mathcal I}:=
      m_u\circ u^*(\iota_{\mathcal I})
      \colon u^*\mathcal I\longrightarrow\mathcal O_T.
  \end{gathered}                                                   \tag{VF1}
  \]
  In particular, define the ideal pulled back by $g$ as the image
  factorization
  \[
    \mu_g:=\mu_g^{\mathcal I}\colon
      g^*\mathcal I\longrightarrow
        \mathcal O_{\operatorname{Spec}(A)},
    \qquad
    \mathcal J:=\operatorname{im}(\mu_g)
      \hookrightarrow\mathcal O_{\operatorname{Spec}(A)}.        \tag{VF2}
  \]
  The adjunction between $u^*$ and $u_*$ identifies the restriction of
  $\varepsilon_u$ to $\mathcal I$ with $\mu_u^{\mathcal I}$; for
  $u=i$ this is the equivalence
  \[
    \varepsilon_i\circ\iota_{\mathcal I}=0
      \Longleftrightarrow \mu_i^{\mathcal I}=0.                  \tag{VF3}
  \]

  Define the functorial pullback isomorphisms
  \[
  \begin{gathered}
    c_{\nu,g}\colon
      \nu_A^*g^*\mathcal I\xrightarrow{\sim}
        (g\circ\nu_A)^*\mathcal I,\\
    c_{g_K,i}\colon
      g_K^*i^*\mathcal I\xrightarrow{\sim}
        (i\circ g_K)^*\mathcal I,
  \end{gathered}                                                   \tag{VF4}
  \]
  by the equations
  \[
  \begin{aligned}
    \mu_{g\circ\nu_A}^{\mathcal I}\circ c_{\nu,g}
      &=\nu_A^*(\mu_g),\\
    \mu_{i\circ g_K}^{\mathcal I}\circ c_{g_K,i}
      &=g_K^*(\mu_i^{\mathcal I}).
  \end{aligned}                                                    \tag{VF5}
  \]
  The specified equality
  $e_K\colon g\circ\nu_A=i\circ g_K$ defines the transport isomorphism
  \[
    t_{e_K,\mathcal I}\colon
      (g\circ\nu_A)^*\mathcal I\xrightarrow{\sim}
        (i\circ g_K)^*\mathcal I,
    \qquad
    \mu_{i\circ g_K}^{\mathcal I}\circ t_{e_K,\mathcal I}
      =\mu_{g\circ\nu_A}^{\mathcal I}.                          \tag{VF6}
  \]
  Since $\nu_A$ is induced by the localization $A\hookrightarrow K$, it is
  flat.  Pullback by $\nu_A$ therefore preserves the image in (VF2), and
  (VF3)--(VF6) are the following formal equality chain:
  \[
  \begin{aligned}
    \nu_A^*\mathcal J
      &=\nu_A^*\operatorname{im}(\mu_g)\\
      &=\operatorname{im}(\nu_A^*\mu_g)\\
      &=\operatorname{im}
        (\mu_{g\circ\nu_A}^{\mathcal I}\circ c_{\nu,g})\\
      &=\operatorname{im}
        (\mu_{i\circ g_K}^{\mathcal I}\circ
          t_{e_K,\mathcal I}\circ c_{\nu,g})\\
      &=\operatorname{im}
        (g_K^*(\mu_i^{\mathcal I})\circ c_{g_K,i}^{-1}\circ
          t_{e_K,\mathcal I}\circ c_{\nu,g})\\
      &=\operatorname{im}(0)=0.
  \end{aligned}                                                    \tag{VF7}
  \]
  Put
  \[
    J_A:=\Gamma(\operatorname{Spec}(A),\mathcal J)\subseteq A.
  \]
  Quasi-coherence of $\mathcal J$ and the affine pullback formula transform
  (VF7) into
  \[
    K\otimes_AJ_A
      =\Gamma(\operatorname{Spec}(K),\nu_A^*\mathcal J)=0.        \tag{VF8}
  \]
  Define the three $A$-linear maps
  \[
  \begin{gathered}
    \eta_{J_A,K}\colon J_A\longrightarrow K\otimes_AJ_A,
    \qquad a\longmapsto1\otimes a,\\
    \kappa_{J_A,K}\colon K\otimes_AJ_A\longrightarrow K,
    \qquad k\otimes a\longmapsto k(a/1),\\
    \lambda_{J_A,K}:=\kappa_{J_A,K}\circ\eta_{J_A,K}
      \colon J_A\longrightarrow K,
    \qquad \lambda_{J_A,K}(a)=a/1.
  \end{gathered}                                                   \tag{VF8a}
  \]
  Equations (VF8)--(VF8a) and injectivity of $A\hookrightarrow K$ are the
  following implications:
  \[
  \begin{aligned}
    K\otimes_AJ_A=0
      &\Longrightarrow\operatorname{im}(\lambda_{J_A,K})=0,\\
    J_A\subseteq A\ \wedge\ \ker(A\hookrightarrow K)=0
      &\Longrightarrow\ker(\lambda_{J_A,K})=0,\\
    \operatorname{im}(\lambda_{J_A,K})=0\ \wedge\
      \ker(\lambda_{J_A,K})=0
      &\Longrightarrow J_A=0
       \Longrightarrow\mathcal J=\widetilde{J_A}=0.
  \end{aligned}                                                    \tag{VF9}
  \]

  By (VF2) and (VF9), $\mu_g=0$.  Define the quotient morphism
  \[
    q_{\mathcal I}\colon\mathcal O_Y\twoheadrightarrow
      \mathcal O_Y/\mathcal I.
  \]
  Adjunction and the quotient universal property are the following
  equivalence-and-implication chain:
  \[
    \mu_g=0
      \Longleftrightarrow
        \varepsilon_g\circ\iota_{\mathcal I}=0
      \Longrightarrow
        \exists!\,\overline\varepsilon_g\colon
          \mathcal O_Y/\mathcal I\to
            g_*\mathcal O_{\operatorname{Spec}(A)}:
          \overline\varepsilon_g\circ q_{\mathcal I}=\varepsilon_g,
                                                                    \tag{VF10}
  \]
  Define the closed model classified by this quotient and the comparison
  with $i$ by
  \[
  \begin{gathered}
    Y_{\mathcal I}:=\operatorname{Spec}_Y(\mathcal O_Y/\mathcal I),
    \qquad
    r_{\mathcal I}:=\operatorname{Spec}_Y(q_{\mathcal I})
      \colon Y_{\mathcal I}\hookrightarrow Y,\\
    \chi_i\colon\mathcal O_Y/\mathcal I
      \xrightarrow{\sim}i_*\mathcal O_Z,
    \qquad
    \chi_i\circ q_{\mathcal I}=\varepsilon_i,\\
    \sigma_i\colon Z\xrightarrow{\sim}Y_{\mathcal I},
    \qquad r_{\mathcal I}\circ\sigma_i=i.
  \end{gathered}                                                   \tag{VF11}
  \]
  Under the relative-spectrum adjunction, let
  \[
    \widehat g\colon\operatorname{Spec}(A)\to Y_{\mathcal I},
    \qquad r_{\mathcal I}\circ\widehat g=g                     \tag{VF12}
  \]
  be the morphism corresponding to $\overline\varepsilon_g$ in (VF10), and
  define
  \[
    \widetilde g:=\sigma_i^{-1}\circ\widehat g
      \colon\operatorname{Spec}(A)\to Z.
  \]
  Equations (VF11)--(VF12) are the following equality chain:
  \[
    i\circ\widetilde g
      =(r_{\mathcal I}\circ\sigma_i)\circ
        \sigma_i^{-1}\circ\widehat g
      =r_{\mathcal I}\circ\widehat g=g.                         \tag{VF13}
  \]
  If $\widetilde g'$ is a second such morphism, then
  \[
    i\circ\widetilde g'=g=i\circ\widetilde g;
  \]
  since a closed immersion is a monomorphism,
  $\widetilde g'=\widetilde g$.
\end{proof}

\begin{lemma}[Scheme-theoretic density survives flat pullback]
  \label{lemma:nagata-dense-flat-pullback}
  \uses{def:nagata-closed-subscheme-class,
    def:nagata-scheme-theoretic-image,
    lemma:nagata-morphism-property-calculus}
  \lean{}

  Let $h\colon U\to X$ satisfy $\operatorname{IsQCQS}(h)$ and
  \[
    [\,\operatorname{SIm}_X(h)\xhookrightarrow{i_h}X\,]
      =[\,X\xhookrightarrow{\operatorname{id}_X}X\,]
      \quad\text{in }\operatorname{ClosedSub}(X).                \tag{DF0}
  \]
  Let $q\colon P\to X$ be flat,
  put $U_P=U\times_XP$, and denote the second projection by
  $h_P\colon U_P\to P$.  Then
  \[
    \operatorname{IsQCQS}(h_P),
    \qquad
    [\,\operatorname{SIm}_P(h_P)\xhookrightarrow{i_{h_P}}P\,]
      =[\,P\xhookrightarrow{\operatorname{id}_P}P\,]
      \quad\text{in }\operatorname{ClosedSub}(P).
  \]

  \paragraph{Mathlib relation.}
  The density hypothesis is the scheme-theoretic dominance predicate
  \texttt{\detokenize{AlgebraicGeometry.IsSchemeTheoreticallyDominant}}.
  Its pullback interfaces are the exact declarations
  \texttt{\detokenize{IsSchemeTheoreticallyDominant.of_isPullback}},
  \texttt{\detokenize{IsSchemeTheoreticallyDominant.pullbackSnd}}, and
  \texttt{\detokenize{IsSchemeTheoreticallyDominant.pullbackFst}}.  The
  present statement adds the blueprint's scheme-image and
  $\operatorname{IsQCQS}$ adapters to those declarations.
\end{lemma}

\begin{proof}
  By base-change stability in
  \Cref{lemma:nagata-morphism-property-calculus},
  \[
    \operatorname{IsQCQS}(h)
      \Longrightarrow\operatorname{IsQCQS}(h_P).
  \]
  By \Cref{def:nagata-scheme-theoretic-image}, set
  \[
    \mathcal I_h
      =\ker\bigl(\mathcal O_X\to h_*\mathcal O_U\bigr),
    \qquad
    \mathcal I_{h_P}
      =\ker\bigl(\mathcal O_P\to(h_P)_*\mathcal O_{U_P}\bigr).
  \]
  Flat base change of scheme-theoretic images for a quasi-compact morphism
  (\cite[\href{https://stacks.math.columbia.edu/tag/081I}{Tag 081I}]
  {stacks-project}) is the equality of quasi-coherent ideals
  \[
    \mathcal I_{h_P}=q^*\mathcal I_h.
  \]
  By (CSub3) of \Cref{def:nagata-closed-subscheme-class}, (DF0) and the
  relative-spectrum description of $i_h$ are the following predicate chain:
  \[
  \begin{aligned}
    &(\mathrm{DF0})\\
    &\Longrightarrow
      \exists\delta\colon\operatorname{SIm}_X(h)\to X:\quad
        \operatorname{IsIso}(\delta)\ \wedge\
        \operatorname{id}_X\circ\delta=i_h\\
    &\Longrightarrow\operatorname{IsIso}(i_h)
     \Longrightarrow\mathcal I_h=0.
  \end{aligned}
  \]
  Substitution into flat base change is
  \[
    \mathcal I_{h_P}=q^*\mathcal I_h=q^*(0)=0.
  \]
  By \Cref{def:nagata-scheme-theoretic-image}, the last ideal equality is
  the class equality
  \[
    [\,\operatorname{SIm}_P(h_P)\xhookrightarrow{i_{h_P}}P\,]
      =[\,P\xhookrightarrow{\operatorname{id}_P}P\,]
      \quad\text{in }\operatorname{ClosedSub}(P).
  \]
\end{proof}

\begin{lemma}[Density of the preserved open after an admissible blow-up]
  \label{lemma:nagata-dense-admissible-blowup}
  \uses{def:nagata-admissible-blowup,
    def:nagata-closed-subscheme-class,
    def:nagata-scheme-theoretic-image,
    lemma:nagata-morphism-property-calculus}

  Let $j\colon U\hookrightarrow X$ be a quasi-compact,
  scheme-theoretically dense open immersion, let
  $\mathcal I\subseteq\mathcal O_X$ be a finite-type quasi-coherent ideal
  with $\mathcal I|_U=\mathcal O_U$, and let
  \[
    b\colon X'=\operatorname{Bl}_{\mathcal I}(X)\longrightarrow X.
  \]
  Put $U'=X'\times_XU$ and let $j'\colon U'\hookrightarrow X'$ be the first
  projection.  Base-change stability and
  \Cref{lemma:nagata-morphism-property-calculus} are the following predicate
  chain:
  \[
    \operatorname{IsQuasiCompactOpenImmersion}(j)
      \Longrightarrow
        \operatorname{IsQuasiCompactOpenImmersion}(j')
      \Longrightarrow
        \operatorname{IsQCQS}(j').
  \]
  By \Cref{def:nagata-scheme-theoretic-image}, define
  $\operatorname{SIm}_{X'}(j')$ using the last predicate.  Then
  \[
    [\,\operatorname{SIm}_{X'}(j')\xhookrightarrow{i_{j'}}X'\,]
      =[\,X'\xhookrightarrow{\operatorname{id}_{X'}}X'\,]
      \quad\text{in }\operatorname{ClosedSub}(X').              \tag{DAB0}
  \]
  Under the isomorphism
  $U'\xrightarrow{\sim}U$ of
  \Cref{thm:nagata-global-blowup-complement}, the preserved copy of $U$ is
  scheme-theoretically dense in $X'$.
\end{lemma}

\begin{proof}
  Cover $X$ by affine opens $X_\lambda=\operatorname{Spec}(A_\lambda)$ on
  which $\mathcal I$ is generated by
  $f_{\lambda,1},\ldots,f_{\lambda,n_\lambda}\in A_\lambda$.
  For $1\leq a\leq n_\lambda$, let
  \[
    X'_{\lambda,a}
      =\operatorname{Spec}
        \bigl(A_\lambda[\mathcal I(X_\lambda)/f_{\lambda,a}]\bigr)
      \subseteq X'
  \]
  be the affine chart of \Cref{def:nagata-blowup-chart-ring}.
  By its definition there is an injective ring map
  \[
    A_\lambda[\mathcal I(X_\lambda)/f_{\lambda,a}]
      \hookrightarrow (A_\lambda)_{f_{\lambda,a}}.
  \]
  Scheme-theoretic density of $j$ and
  \Cref{lemma:nagata-scheme-image-restrict-open} are the injections
  \[
    A_\lambda\hookrightarrow
      \Gamma(U\cap X_\lambda,\mathcal O_U),
    \qquad
    (A_\lambda)_{f_{\lambda,a}}\hookrightarrow
      \Gamma(U\cap D(f_{\lambda,a}),\mathcal O_U).
  \]
  Over $U$, the chart $X'_{\lambda,a}$ identifies with
  $U\cap D(f_{\lambda,a})$ by
  \Cref{thm:nagata-global-blowup-complement}.  Hence the restriction
  homomorphism on the chart is the chain of injections
  \[
    A_\lambda[\mathcal I(X_\lambda)/f_{\lambda,a}]
      \hookrightarrow(A_\lambda)_{f_{\lambda,a}}
      \hookrightarrow
      \Gamma(U\cap D(f_{\lambda,a}),\mathcal O_U)
      \xrightarrow{\sim}
      \Gamma(U'\cap X'_{\lambda,a},\mathcal O_{U'}).
  \]
  The composite is the restriction map from the chart ring, so its kernel is
  zero.
  The charts $X'_{\lambda,a}$ cover $X'$ by
  \Cref{prop:nagata-affine-blowup-chart-cover}; thus
  \[
    \ker\bigl(\mathcal O_{X'}\to j'_*\mathcal O_{U'}\bigr)=0.
  \]
  By \Cref{def:nagata-scheme-theoretic-image}, the preceding kernel
  equality is the class equality
  \[
    [\,\operatorname{SIm}_{X'}(j')\xhookrightarrow{i_{j'}}X'\,]
      =[\,\operatorname{Spec}_{X'}(\mathcal O_{X'}/0)
        \xhookrightarrow{\operatorname{Spec}_{X'}
          (\mathcal O_{X'}\twoheadrightarrow\mathcal O_{X'}/0)}X'\,]
      =[\,X'\xhookrightarrow{\operatorname{id}_{X'}}X'\,]
      \quad\text{in }\operatorname{ClosedSub}(X').
  \]
\end{proof}

\begin{definition}[Valuation probe through a schematically dense source]
  \label{def:nagata-dense-valuation-probe}
  \uses{def:nagata-choice-notation,def:nagata-closed-subscheme-class,
    def:nagata-scheme-theoretic-image}
  \lean{}

  Let $h\colon T\to P$ satisfy $\operatorname{IsQCQS}(h)$, assume
  \[
    [\,\operatorname{SIm}_P(h)\xhookrightarrow{i_h}P\,]
      =[\,P\xhookrightarrow{\operatorname{id}_P}P\,]
      \quad\text{in }\operatorname{ClosedSub}(P),                \tag{VP0}
  \]
  and let $z\in|P|$.  For a valuation ring $A'$,
  write
  \[
    K'=\operatorname{Frac}(A'),\qquad
    \nu_{A'}\colon\operatorname{Spec}(K')\to\operatorname{Spec}(A'),
    \qquad c_{A'}\in|\operatorname{Spec}(A')|
  \]
  for its fraction field, generic-point morphism, and closed point.  Define
  $\mathscr V(h,z)$ to be the set of tuples
  \[
    (A',K',\nu_{A'},\eta,\alpha)
  \]
  such that $A'$ is a valuation ring,
  \[
    \eta\colon\operatorname{Spec}(K')\to T,
    \qquad
    \alpha\colon\operatorname{Spec}(A')\to P,
  \]
  and
  \[
    h\circ\eta=\alpha\circ\nu_{A'},
    \qquad
    \alpha(c_{A'})=z.
  \]
  The valuation realization of a point of a scheme-theoretic image
  (\cite[\href{https://stacks.math.columbia.edu/tag/02JQ}{Tag 02JQ}]
  {stacks-project}) asserts
  $\mathscr V(h,z)\neq\varnothing$.  Using
  \Cref{def:nagata-choice-notation}, define
  \[
    \mathsf{VP}(h,z)=
      (A'_{h,z},K'_{h,z},\nu_{h,z},\eta_{h,z},\alpha_{h,z})
      :=\operatorname{ch}(\mathscr V(h,z)).
  \]
\end{definition}

\begin{lemma}[Properties of the valuation probe]
  \label{lemma:nagata-dense-valuation-probe}
  \uses{def:nagata-closed-subscheme-class,
    def:nagata-dense-valuation-probe}
  \lean{}

  Let $h\colon T\to P$ satisfy $\operatorname{IsQCQS}(h)$ and
  \[
    [\,\operatorname{SIm}_P(h)\xhookrightarrow{i_h}P\,]
      =[\,P\xhookrightarrow{\operatorname{id}_P}P\,]
      \quad\text{in }\operatorname{ClosedSub}(P),
  \]
  and let $z\in|P|$.  The tuple
  $\mathsf{VP}(h,z)$ satisfies
  \[
    K'_{h,z}=\operatorname{Frac}(A'_{h,z}),
    \qquad
    h\circ\eta_{h,z}=\alpha_{h,z}\circ\nu_{h,z},
    \qquad
    \alpha_{h,z}(c_{A'_{h,z}})=z.
  \]
\end{lemma}

\begin{proof}
  By \Cref{def:nagata-dense-valuation-probe},
  \[
    \mathsf{VP}(h,z)=\operatorname{ch}(\mathscr V(h,z))
      \in\mathscr V(h,z).
  \]
  Expanding the three defining predicates of membership in
  $\mathscr V(h,z)$ is exactly the displayed triple of equalities.
\end{proof}


\begin{lemma}[Composite of two valuation rings]
  \label{lemma:nagata-composite-valuation-ring}
  \lean{}

  Let $A'$ be a valuation ring, put
  $K'=\operatorname{Frac}(A')$, let $\mathfrak m'$ be its maximal ideal, and
  put $k'=A'/\mathfrak m'$.  Let $A''\subseteq k'$ be a valuation ring with
  $\operatorname{Frac}(A'')=k'$, and let
  $\rho\colon A'\twoheadrightarrow k'$ be the quotient map.  Define
  \[
    C=\rho^{-1}(A'')=\{a\in A'\mid\rho(a)\in A''\}.
  \]
  Then $C$ is a valuation ring and
  \[
    \operatorname{Frac}(C)=K',
    \qquad
    \mathfrak m'\subseteq C,
    \qquad
    C/\mathfrak m'\xrightarrow{\sim}A'',\quad
    c+\mathfrak m'\longmapsto\rho(c).
  \]
\end{lemma}

\begin{proof}
  The equality $\rho(\mathfrak m')=0\subseteq A''$ is
  $\mathfrak m'\subseteq C$, and the restriction of $\rho$ is the exact
  sequence
  \[
    0\longrightarrow\mathfrak m'\longrightarrow C
      \xrightarrow{\rho|_C}A''\longrightarrow0.
  \]
  Thus $C/\mathfrak m'\xrightarrow{\sim}A''$ by the first isomorphism theorem.

  Let $x\in K'^\times$.  Since $A'$ is a valuation ring,
  \[
    x\in A'\quad\text{or}\quad x^{-1}\in A'.
  \]
  If $x\in\mathfrak m'$, then $\rho(x)=0\in A''$ and $x\in C$.  If
  $x^{-1}\in\mathfrak m'$, the same equality applied to $x^{-1}$ is
  $x^{-1}\in C$.  In the remaining case $x,x^{-1}\in(A')^\times$, so
  \[
    \rho(x)\rho(x^{-1})=1\quad\text{in }k'.
  \]
  Since $A''$ is a valuation ring of $k'$, either
  $\rho(x)\in A''$ or $\rho(x)^{-1}=\rho(x^{-1})\in A''$.  By the definition
  $C=\rho^{-1}(A'')$, either $x\in C$ or $x^{-1}\in C$.  The valuation-ring
  criterion therefore shows that $C$ is a valuation ring.  Since
  $\mathfrak m'\subseteq C\subseteq A'$ and every $y\in A'$ satisfies
  \[
    y\in C\quad\text{or}\quad
    (y\notin C\Rightarrow y\in(A')^\times\text{ and }y^{-1}\in C),
  \]
  every element of $A'$ lies in $\operatorname{Frac}(C)$.  Hence
  \[
    K'=\operatorname{Frac}(A')
       \subseteq\operatorname{Frac}(C)
       \subseteq K',
  \]
  which proves $\operatorname{Frac}(C)=K'$.
\end{proof}

\begin{lemma}[Intersection with a dominating valuation ring]
  \label{lemma:nagata-valuation-intersection}
  \lean{}

  Let $A$ be a valuation ring, put $K=\operatorname{Frac}(A)$, let
  $\iota\colon K\hookrightarrow L$ be a field embedding, and let
  $B\subseteq L$ be a valuation ring satisfying
  \[
    \iota(A)\subseteq B,
    \qquad
    \mathfrak m_B\cap\iota(A)=\iota(\mathfrak m_A).
  \]
  Then
  \[
    \iota(K)\cap B=\iota(A)
  \]
  inside \(L\).
\end{lemma}

\begin{proof}
  The first domination equality is
  $\iota(A)\subseteq\iota(K)\cap B$.  Let $x\in K$ satisfy $\iota(x)\in B$.
  If $x\notin A$, the valuation-ring alternative is
  \[
    x^{-1}\in\mathfrak m_A.
  \]
  The second domination equality is
  $\iota(x^{-1})\in\mathfrak m_B$.  Since $\iota(x)\in B$,
  \[
    1=\iota(x)\iota(x^{-1})\in B\mathfrak m_B=\mathfrak m_B,
  \]
  a contradiction.  Therefore $x\in A$, so
  $\iota(K)\cap B\subseteq\iota(A)$.  Antisymmetry of the two inclusions is
  \[
    \iota(K)\cap B=\iota(A).
  \]
\end{proof}

\begin{lemma}[{\cite[\href{https://stacks.math.columbia.edu/tag/0894}{Tag 0894}]{stacks-project}}]
  \label{lemma:nagata-dense-source-valuative}
  \lean{}

  Let $f\colon X\to S$ be quasi-compact, let $h\colon U\to X$ satisfy
  $\operatorname{IsQCQS}(h)$, and assume
  $\overline{|h|(|U|)}=|X|$.  Suppose that for every valuation ring $A$,
  with $K=\operatorname{Frac}(A)$ and
  $\nu_A\colon\operatorname{Spec}(K)\to\operatorname{Spec}(A)$, every pair
  \[
    u_K\colon\operatorname{Spec}(K)\to U,
    \qquad
    s_A\colon\operatorname{Spec}(A)\to S
  \]
  satisfying
  \[
    f\circ h\circ u_K=s_A\circ\nu_A
  \]
  admits a unique morphism
  $x_A\colon\operatorname{Spec}(A)\to X$ satisfying
  \[
    x_A\circ\nu_A=h\circ u_K,
    \qquad
    f\circ x_A=s_A.
  \]
  Then $f$ is universally closed.  If $f$ is quasi-separated, then $f$ is
  separated.

  \paragraph{Mathlib relation.}
  The quantified square is represented by
  \texttt{\detokenize{AlgebraicGeometry.ValuativeCommSq}}.  Its existence
  and uniqueness predicates are the exact morphism properties
  \texttt{\detokenize{AlgebraicGeometry.ValuativeCriterion.Existence}} and
  \texttt{\detokenize{AlgebraicGeometry.ValuativeCriterion.Uniqueness}}.
  The lemma's dense-source reduction to those predicates is a
  blueprint-local theorem, not an existing declaration with the proposed
  name above.
\end{lemma}

\begin{proof}
  \uses{def:nagata-closed-subscheme-class,
    lemma:nagata-dense-flat-pullback,lemma:nagata-dense-valuation-probe,
    lemma:nagata-composite-valuation-ring,lemma:nagata-valuation-intersection,
    def:nagata-morphism-predicate-notation}
  Let $A$ be a valuation ring, put $K=\operatorname{Frac}(A)$, and fix a
  commutative square
  \[
    \begin{tikzcd}
      \operatorname{Spec}(K) \arrow[r,"x_K"] \arrow[d,"\nu_A"']
        & X \arrow[d,"f"]\\
      \operatorname{Spec}(A) \arrow[r,"s_A"'] & S
    \end{tikzcd}
    \qquad\text{with}\qquad
    f\circ x_K=s_A\circ\nu_A.
  \]
  Define the reduction immersions
  \[
    r_U\colon U_0:=U_{\mathrm{red}}\hookrightarrow U,\qquad
    r_X\colon X_0:=X_{\mathrm{red}}\hookrightarrow X,\qquad
    r_S\colon S_0:=S_{\mathrm{red}}\hookrightarrow S.
  \]
  Because $\operatorname{Spec}(A)$ and $\operatorname{Spec}(K)$ are reduced,
  the reduction factorization lemma
  (\cite[\href{https://stacks.math.columbia.edu/tag/0356}{Tag 0356}]
  {stacks-project}) defines unique morphisms
  \[
    x_{K,0}\colon\operatorname{Spec}(K)\to X_0,\qquad
    s_{A,0}\colon\operatorname{Spec}(A)\to S_0
  \]
  satisfying
  \[
    r_X\circ x_{K,0}=x_K,\qquad r_S\circ s_{A,0}=s_A.             \tag{0a}
  \]
  Define $h_0\colon U_0\to X_0$ and $f_0\colon X_0\to S_0$ by
  \[
    r_X\circ h_0=h\circ r_U,\qquad
    r_S\circ f_0=f\circ r_X.                                    \tag{0b}
  \]
  Define the two diagonal targets and their reduction comparison by
  \[
  \begin{gathered}
    P_0:=U_0\times_{X_0}U_0,
    \qquad P:=U\times_XU,\\
    r_P:=(r_U\circ\operatorname{pr}_1,
      r_U\circ\operatorname{pr}_2)\colon P_0\longrightarrow P,\\
    r_P\circ\Delta_{h_0}=\Delta_h\circ r_U.
  \end{gathered}                                                   \tag{0bq}
  \]
  By \Cref{def:nagata-reduced-closed-subscheme}, the maps
  $|r_U|$, $|r_X|$, and $|r_P|$ are homeomorphisms.  Substitution of (0bq)
  into the definitions in
  \Cref{def:nagata-morphism-predicate-notation} is the predicate chain
  \[
  \begin{aligned}
    \operatorname{IsQCQS}(h)
      &\Longleftrightarrow
        \bigl(
          \operatorname{IsQuasiCompact}(h)\ \wedge\
          \operatorname{IsQuasiCompact}(\Delta_h)
        \bigr)\\
      &\Longrightarrow
        \bigl(
          \operatorname{IsQuasiCompact}(h_0)\ \wedge\
          \operatorname{IsQuasiCompact}(\Delta_{h_0})
        \bigr)\\
      &\Longleftrightarrow \operatorname{IsQCQS}(h_0).
  \end{aligned}                                                    \tag{0bqcqs}
  \]
  The homeomorphisms $|r_U|$ and $|r_X|$, together with (0b), are the
  following closure equality:
  \[
    \overline{h_0(|U_0|)}=|X_0|.
  \]
  The reducedness predicates for $U_0$ and $X_0$, the scheme-image formula
  (\cite[\href{https://stacks.math.columbia.edu/tag/056B}{Tag 056B}]
  {stacks-project}), and the preceding support equality are the following
  class equality:
  \[
    [\,\operatorname{SIm}_{X_0}(h_0)\xhookrightarrow{i_{h_0}}X_0\,]
      =[\,X_0\xhookrightarrow{\operatorname{id}_{X_0}}X_0\,]
      \quad\text{in }\operatorname{ClosedSub}(X_0).              \tag{0b'}
  \]
  Equations (0a)--(0b), the original square, and monicity of $r_S$ form the
  chain
  \[
    r_S\circ f_0\circ x_{K,0}
      =f\circ x_K
      =s_A\circ\nu_A
      =r_S\circ s_{A,0}\circ\nu_A,
    \qquad
    f_0\circ x_{K,0}=s_{A,0}\circ\nu_A.                          \tag{0c}
  \]
  Every valuation-spectrum morphism into $U$, $X$, or $S$ factors uniquely
  through $r_U$, $r_X$, or $r_S$, respectively.  Under these bijections, the
  original existence-and-uniqueness predicate equals the predicate for
  $(f_0,h_0)$.  Make the explicit substitution
  \[
    (U,X,S,h,f,x_K,s_A)
      :=(U_0,X_0,S_0,h_0,f_0,x_{K,0},s_{A,0}).
  \]

  Let $a\in|\operatorname{Spec}(A)|$ be the closed point and choose an affine
  open immersion
  $v_1\colon S_1=\operatorname{Spec}(R)\hookrightarrow S$ containing
  $s_A(a)$.  The open
  subset $s_A^{-1}(S_1)$ contains $a$; every point of
  $\operatorname{Spec}(A)$ is a generization of $a$, so
  \[
    s_A^{-1}(S_1)=\operatorname{Spec}(A).
  \]
  Define the two base-change squares and their projections by
  \[
    X_1:=X\times_SS_1,\quad
    p_X\colon X_1\to X,\quad f_1\colon X_1\to S_1,
    \qquad
    U_1:=U\times_SS_1,\quad
    p_U\colon U_1\to U,\quad s_{U,1}\colon U_1\to S_1.
  \]
  Define $h_1\colon U_1\to X_1$ by
  \[
    p_X\circ h_1=h\circ p_U,\qquad
    f_1\circ h_1=s_{U,1}.                                       \tag{0d}
  \]
  The square defining $X_1$ and the stalkwise flatness predicate (P0) are the
  implications
  \[
    \operatorname{IsOpenImmersion}(v_1)
      \Longrightarrow\operatorname{IsOpenImmersion}(p_X)
      \Longrightarrow\operatorname{IsFlat}(p_X).
  \]
  By \Cref{lemma:nagata-dense-flat-pullback},
  \[
    [\,\operatorname{SIm}_{X_1}(h_1)\xhookrightarrow{i_{h_1}}X_1\,]
      =[\,X_1\xhookrightarrow{\operatorname{id}_{X_1}}X_1\,]
      \quad\text{in }\operatorname{ClosedSub}(X_1).             \tag{1a}
  \]
  The equality $s_A^{-1}(S_1)=\operatorname{Spec}(A)$ defines
  $s_{A,1}\colon\operatorname{Spec}(A)\to S_1$ by
  \[
    v_1\circ s_{A,1}=s_A.
  \]
  The fibre-product property defines
  $x_{K,1}\colon\operatorname{Spec}(K)\to X_1$ by
  \[
    p_X\circ x_{K,1}=x_K,\qquad
    f_1\circ x_{K,1}=s_{A,1}\circ\nu_A.                          \tag{0e}
  \]
  Base change identifies the original lifting predicate with the lifting
  predicate for $(f_1,h_1)$.  Make the explicit substitution
  \[
    (U,X,S,h,f,x_K,s_A)
      :=(U_1,X_1,S_1,h_1,f_1,x_{K,1},s_{A,1}),
    \qquad S=\operatorname{Spec}(R).
  \]

  Choose an affine open immersion
  $\xi_B\colon X_B=\operatorname{Spec}(B)\hookrightarrow X$ and the unique
  factor
  \[
    x_{K,B}\colon\operatorname{Spec}(K)\to X_B,
    \qquad \xi_B\circ x_{K,B}=x_K.
  \]
  By \Cref{def:nagata-direct-image-structure-map}~(DIM0), define
  \[
  \begin{gathered}
    r:=(f\circ\xi_B)^\#\colon R\to B,
    \qquad b:=x_{K,B}^\#\colon B\to K,\\
    a:=s_A^\#\colon R\to A,
    \qquad \iota_A:=\nu_A^\#\colon A\hookrightarrow K.
  \end{gathered}                                                   \tag{1-0}
  \]
  The affine-spectrum correspondence applied to (1-0) is the chain
  \[
  \begin{gathered}
    f\circ\xi_B=\operatorname{Spec}(r),\qquad
    x_{K,B}=\operatorname{Spec}(b),\qquad
    x_K=\xi_B\circ\operatorname{Spec}(b),\\
    s_A=\operatorname{Spec}(a),\qquad
    \nu_A=\operatorname{Spec}(\iota_A).
  \end{gathered}                                                   \tag{1-0a}
  \]
  Commutativity of the original square is the equality
  \[
    b\circ r=\iota_A\circ a.                                    \tag{1}
  \]
  Choose a polynomial $B$-algebra $P_B=B[t_\lambda]_{\lambda\in\Lambda}$ and
  a surjective $B$-algebra map
  \[
    \pi\colon P_B\twoheadrightarrow K,
    \qquad \pi|_B=b.
  \]
  Let $\jmath\colon B\hookrightarrow P_B$ be the polynomial-algebra
  structure map, and define
  \[
    Q:=\operatorname{Spec}(P_B),\qquad
    q:=\xi_B\circ\operatorname{Spec}(\jmath)\colon Q\to X.
  \]
  Define $z_K:=\operatorname{Spec}(\pi)\colon\operatorname{Spec}(K)\to Q$
  and $z:=|z_K|(*)\in|Q|$; then $\kappa(z)=K$.  The stalkwise definition of
  $\operatorname{IsFlat}$ in
  \Cref{def:nagata-morphism-predicate-notation} is the chain
  \[
  \begin{gathered}
    P_B=\bigoplus_{\mathbf n\in\mathbf N^{(\Lambda)}}B\,t^{\mathbf n}
      \Longrightarrow P_B\text{ is a free }B\text{-module}
      \Longrightarrow P_B\text{ is a flat }B\text{-module},\\
    \operatorname{IsFlat}(\operatorname{Spec}(\jmath)),\qquad
    \operatorname{IsOpenImmersion}(\xi_B)
      \Longrightarrow\operatorname{IsFlat}(\xi_B),\\
    q=\xi_B\circ\operatorname{Spec}(\jmath)
      \wedge
      \operatorname{IsFlat}(\operatorname{Spec}(\jmath))
      \wedge
      \operatorname{IsFlat}(\xi_B)
      \Longrightarrow\operatorname{IsFlat}(q).
  \end{gathered}                                                   \tag{1b}
  \]

  Put
  \[
    T=U\times_XQ,\qquad h_Q=\operatorname{pr}_Q\colon T\to Q.
  \]
  By \Cref{lemma:nagata-dense-flat-pullback},
  \[
    [\,\operatorname{SIm}_Q(h_Q)\xhookrightarrow{i_{h_Q}}Q\,]
      =[\,Q\xhookrightarrow{\operatorname{id}_Q}Q\,]
      \quad\text{in }\operatorname{ClosedSub}(Q).               \tag{2}
  \]
  Apply \Cref{def:nagata-dense-valuation-probe} to $(h_Q,z)$ and abbreviate
  \[
    \mathsf{VP}(h_Q,z)=(A',K',\nu',\eta,\alpha),
    \qquad
    \mathfrak m'=\mathfrak m_{A'},
    \qquad k'=A'/\mathfrak m'.
  \]
  By \Cref{def:nagata-direct-image-structure-map}~(DIM0), define
  \[
    \alpha^\#:=\varepsilon_{\alpha,|Q|}\colon P_B\to A',
  \]
  and let
  $\rho\colon A'\twoheadrightarrow k'$ be the quotient.  The equality
  $|\alpha|(c_{A'})=z$ in
  \Cref{lemma:nagata-dense-valuation-probe} is
  \[
    (\alpha^\#)^{-1}(\mathfrak m')=\ker(\pi).
  \]
  The quotient universal property for
  $K=P_B/\ker(\pi)$ defines the unique field embedding
  $\varphi\colon K\hookrightarrow k'$ satisfying
  \[
    \rho\circ\alpha^\#=\varphi\circ\pi.                       \tag{3}
  \]

  By
  \cite[\href{https://stacks.math.columbia.edu/tag/00IA}{Tag 00IA}]
  {stacks-project}, choose a valuation ring $A''\subseteq k'$ such that
  \[
    \operatorname{Frac}(A'')=k',\qquad
    \operatorname{im}(\varphi\circ\iota_A)
      =\varphi(\iota_A(A))\subseteq A'',\qquad
    \mathfrak m_{A''}\cap\varphi(\iota_A(A))
      =\varphi(\iota_A(\mathfrak m_A)).                         \tag{4}
  \]
  Define
  \[
    C=\rho^{-1}(A'')\subseteq A'.
  \]
  By \Cref{lemma:nagata-composite-valuation-ring},
  \[
    \operatorname{Frac}(C)=K',
    \qquad C/\mathfrak m'=A''.                                  \tag{5}
  \]
  For the polynomial-algebra structure map
  $\jmath\colon B\hookrightarrow P_B$ defined above, substitution of (1) and
  (3) is the chain
  \[
    \rho\circ\alpha^\#\circ\jmath\circ r
      =\varphi\circ\pi\circ\jmath\circ r
      =\varphi\circ b\circ r
      =\varphi\circ\iota_A\circ a.
                                                                    \tag{6}
  \]
  The inclusion $\varphi(\iota_A(A))\subseteq A''$ in (4), substituted into
  (6), is the successive pair of image inclusions
  \[
  \begin{aligned}
    \operatorname{im}(\rho\circ\alpha^\#\circ\jmath\circ r)
      &=\operatorname{im}(\varphi\circ\iota_A\circ a)
        \subseteq A'',\\
    \operatorname{im}(\alpha^\#\circ\jmath\circ r)
      &\subseteq\rho^{-1}(A'')=C.
  \end{aligned}
  \]
  Let $\iota_{C,A'}\colon C\hookrightarrow A'$ denote the inclusion.  Define
  $c_R\colon R\to C$ to be the unique ring homomorphism satisfying
  \[
    \iota_{C,A'}\circ c_R
      =\alpha^\#\circ\jmath\circ r,
  \]
  and define
  \[
    s_C:=\operatorname{Spec}(c_R)\colon\operatorname{Spec}(C)\to
      \operatorname{Spec}(R)=S.
  \]
  Let $\iota_{C,K'}\colon C\hookrightarrow K'$ be the fraction-field
  inclusion supplied by (5).  Define
  \[
    \lambda:=\operatorname{Spec}(\iota_{C,A'})
      \colon\operatorname{Spec}(A')\to\operatorname{Spec}(C),
    \qquad
    \nu_C:=\operatorname{Spec}(\iota_{C,K'})
      \colon\operatorname{Spec}(K')\to\operatorname{Spec}(C)
  \]
  and let
  \(\iota_{A',K'}\colon A'\hookrightarrow K'\) be the fraction-field
  inclusion defining \(\nu'\).  Contravariance of
  \(\operatorname{Spec}\) is the explicit chain
  \[
  \begin{aligned}
    \nu_C
      &=\operatorname{Spec}(\iota_{C,K'})\\
      &=\operatorname{Spec}
        (\iota_{A',K'}\circ\iota_{C,A'})\\
      &=\operatorname{Spec}(\iota_{C,A'})
          \circ\operatorname{Spec}(\iota_{A',K'})\\
      &=\lambda\circ\nu'.
  \end{aligned}                                                    \tag{DV0}
  \]

  Let
  \[
    u'=\operatorname{pr}_U\circ\eta
      \colon\operatorname{Spec}(K')\to U.
  \]
  The fibre-product equality and
  \Cref{lemma:nagata-dense-valuation-probe} are the chain
  \[
    h\circ u'
      =q\circ h_Q\circ\eta
      =q\circ\alpha\circ\nu'.                                 \tag{7}
  \]
  The definition of $c_R\colon R\to C$ is the ring-map equality
  \[
    (C\hookrightarrow A')\circ c_R
      =\alpha^\#\circ\jmath\circ r.
  \]
  Consequently
  \[
  \begin{aligned}
    f\circ h\circ u'
      &=f\circ q\circ\alpha\circ\nu'\\
      &=s_C\circ\lambda\circ\nu'\\
      &=s_C\circ\nu_C.
  \end{aligned}                                                   \tag{8}
  \]
  The hypothesis applied to the valuation ring $C$, its fraction field $K'$,
  and (8) defines a unique map
  \[
    x_C\colon\operatorname{Spec}(C)\to X,
    \qquad
    x_C\circ\nu_C=h\circ u',
    \qquad
    f\circ x_C=s_C.                                               \tag{9}
  \]

  Define
  \[
    y_1=q\circ\alpha,
    \qquad
    y_2=x_C\circ\lambda
    \colon\operatorname{Spec}(A')\to X.
  \]
  Equations (7)--(9) are the pair
  \[
    y_1\circ\nu'=h\circ u'=y_2\circ\nu',
    \qquad
    f\circ y_1=f\circ y_2.
  \]
  The uniqueness hypothesis applied to $A'$ is the equality
  \[
    q\circ\alpha=y_1=y_2=x_C\circ\lambda.                       \tag{10}
  \]

  Define
  \[
    \sigma:=\operatorname{Spec}(C\twoheadrightarrow A'')
      \colon\operatorname{Spec}(A'')\to\operatorname{Spec}(C),
    \qquad
    \tau:=\operatorname{Spec}(A''\hookrightarrow k')
      \colon\operatorname{Spec}(k')\to\operatorname{Spec}(A''),
  \]
  and define
  \[
    \varphi^*:=\operatorname{Spec}(\varphi)
      \colon\operatorname{Spec}(k')\to\operatorname{Spec}(K).
  \]
  Restricting (10) to $\operatorname{Spec}(k')$ and substituting (3) is the
  equality
  \[
    (x_C\circ\sigma)\circ\tau=x_K\circ\varphi^*.                \tag{11}
  \]
  Define the points and maps
  \[
  \begin{gathered}
    g'':=x_C\circ\sigma\colon\operatorname{Spec}(A'')\to X,
    \qquad c_{A''}:=\mathfrak m_{A''},
    \qquad \eta_{A''}:=(0),\\
    \ast_K:=(0)\in|\operatorname{Spec}(K)|,
    \qquad \ast_{k'}:=(0)\in|\operatorname{Spec}(k')|,\\
    x:=|g''|(c_{A''}),
    \qquad y_K:=|x_K|(\ast_K).
  \end{gathered}                                                   \tag{11a}
  \]
  By \Cref{def:nagata-morphism-predicate-notation}~(P-1), define
  \[
    \ell:=(g'')_{c_{A''}}^\sharp
      \colon\mathcal O_{X,x}\longrightarrow
      \mathcal O_{\operatorname{Spec}(A''),c_{A''}}=A''.
                                                                    \tag{11b}
  \]
  Since
  \[
    |\tau|(\ast_{k'})=\eta_{A''},
    \qquad |\varphi^*|(\ast_{k'})=\ast_K,
    \qquad c_{A''}\in\overline{\{\eta_{A''}\}},
  \]
  the equality (11) and continuity of $|g''|$ are the following equality and
  specialization chain:
  \[
  \begin{aligned}
    |g''|(\eta_{A''})
      &=|g''\circ\tau|(\ast_{k'})
       =|x_K\circ\varphi^*|(\ast_{k'})
       =|x_K|(\ast_K)=y_K,\\
    x=|g''|(c_{A''})
      &\in|g''|\bigl(\overline{\{\eta_{A''}\}}\bigr)
       \subseteq\overline{\{|g''|(\eta_{A''})\}}
       =\overline{\{y_K\}}.
  \end{aligned}                                                    \tag{11c}
  \]
  For any generization pair $p\in\overline{\{q\}}$ in a scheme $V$,
  define the generization homomorphism
  \[
    \gamma^V_{p,q}\colon\mathcal O_{V,p}\longrightarrow\mathcal O_{V,q},
    \qquad \gamma^V_{p,q}([W,s]_p):=[W,s]_q.                    \tag{11d}
  \]
  In particular, (11c) defines
  $\gamma^X_{x,y_K}\colon\mathcal O_{X,x}\to\mathcal O_{X,y_K}$.
  By (P-1), define
  \[
  \begin{gathered}
    (x_K)_{\ast_K}^\sharp\colon
      \mathcal O_{X,y_K}\longrightarrow
      \mathcal O_{\operatorname{Spec}(K),\ast_K}=K,\\
    \ell_{K,x}:=(x_K)_{\ast_K}^\sharp\circ\gamma^X_{x,y_K}
      \colon\mathcal O_{X,x}\longrightarrow K.
  \end{gathered}                                                   \tag{11e}
  \]
  Let $\iota_{A'',k'}\colon A''\hookrightarrow k'$ denote the fraction-field
  inclusion.  The definitions of $\tau$ and $\varphi^*$, (P-1), and (11d)
  are the identities
  \[
  \begin{gathered}
    \tau_{\ast_{k'}}^\sharp\circ
      \gamma^{\operatorname{Spec}(A'')}_{c_{A''},\eta_{A''}}
       =\iota_{A'',k'},
    \qquad (\varphi^*)_{\ast_{k'}}^\sharp=\varphi,\\
    \gamma^{\operatorname{Spec}(A'')}_{c_{A''},\eta_{A''}}
      \circ (g'')_{c_{A''}}^\sharp
      =(g'')_{\eta_{A''}}^\sharp\circ\gamma^X_{x,y_K}.
  \end{gathered}                                                   \tag{11f}
  \]
  Functoriality of the stalk maps in (P-1) applied to (11), together with
  substitutions (11b), (11e), and (11f), is the following equality chain:
  \[
  \begin{aligned}
    \iota_{A'',k'}\circ\ell
      &=\tau_{\ast_{k'}}^\sharp\circ
        \gamma^{\operatorname{Spec}(A'')}_{c_{A''},\eta_{A''}}
          \circ (g'')_{c_{A''}}^\sharp\\
      &=\tau_{\ast_{k'}}^\sharp\circ (g'')_{\eta_{A''}}^\sharp
          \circ\gamma^X_{x,y_K}\\
      &=(g''\circ\tau)_{\ast_{k'}}^\sharp
          \circ\gamma^X_{x,y_K}\\
      &=(x_K\circ\varphi^*)_{\ast_{k'}}^\sharp
          \circ\gamma^X_{x,y_K}\\
      &=(\varphi^*)_{\ast_{k'}}^\sharp\circ
          (x_K)_{\ast_K}^\sharp\circ\gamma^X_{x,y_K}\\
      &=\varphi\circ\ell_{K,x}.
  \end{aligned}                                                    \tag{11g}
  \]
  Consequently the image computation is
  \[
  \begin{aligned}
    \iota_{A'',k'}(\operatorname{im}(\ell))
      &=\operatorname{im}(\varphi\circ\ell_{K,x})
        \subseteq\varphi(K),\\
    \operatorname{im}(\ell)
      &\subseteq A''\cap\varphi(K),
  \end{aligned}
  \]
  The conclusion of \Cref{lemma:nagata-valuation-intersection} under the
  substitution $(\iota,B)=(\varphi,A'')$ in (4) is
  \[
    A''\cap\varphi(K)=\varphi(\iota_A(A)),
    \qquad
    \operatorname{im}(\ell)\subseteq\varphi(\iota_A(A)).       \tag{12}
  \]
  Define
  \[
  \begin{gathered}
    \psi_A\colon\varphi(\iota_A(A))\xrightarrow{\sim}A,
    \qquad \psi_A\bigl(\varphi(\iota_A(a_0))\bigr):=a_0,\\
    \overline\ell\colon\mathcal O_{X,x}\longrightarrow
      \varphi(\iota_A(A)),
    \qquad \overline\ell(t):=\iota_{A'',k'}(\ell(t)),\\
    \ell_A:=\psi_A\circ\overline\ell
      \colon\mathcal O_{X,x}\longrightarrow A.
  \end{gathered}                                                   \tag{12-0}
  \]
  Equations (4), (12), and (12-0), together with locality of the stalk map
  $\ell=(g'')_{c_{A''}}^\sharp$, are the following inverse-image chain:
  \[
  \begin{aligned}
    \ell_A^{-1}(\mathfrak m_A)
      &=\overline\ell^{-1}
          \bigl(\varphi(\iota_A(\mathfrak m_A))\bigr)\\
      &=\ell^{-1}
          \bigl(\mathfrak m_{A''}\cap\varphi(\iota_A(A))\bigr)\\
      &=\ell^{-1}(\mathfrak m_{A''})
       =\mathfrak m_x.
  \end{aligned}                                                    \tag{12-0a}
  \]
  Hence $\ell_A$ is local.  Choose an affine open immersion
  $j_x\colon X_x=\operatorname{Spec}(R_x)\hookrightarrow X$ with $x\in|X_x|$,
  and define
  \[
  \begin{gathered}
    \lambda_{R_x,x}\colon R_x\longrightarrow\mathcal O_{X,x},
    \qquad \lambda_{R_x,x}(s):=[X_x,s]_x,\\
    \lambda_x:=j_x\circ\operatorname{Spec}(\lambda_{R_x,x})
      \colon\operatorname{Spec}(\mathcal O_{X,x})\to X,\\
    x_A:=\lambda_x\circ\operatorname{Spec}(\ell_A)
      \colon\operatorname{Spec}(A)\to X.
  \end{gathered}                                                   \tag{12-0b}
  \]
  Locality in (12-0a) is the equality
  \[
    |x_A|(\mathfrak m_A)=x.
  \]
  Put $y:=|f|(x)\in|S|$.  By (P-1), define the structural local-ring map;
  define also the quotient map induced by $\rho$:
  \[
  \begin{gathered}
    \lambda_{R,y}\colon R=\mathcal O_S(|S|)\longrightarrow
      \mathcal O_{S,y},
    \qquad \lambda_{R,y}(r_0):=[S,r_0]_y,\\
    f_x^\sharp\colon\mathcal O_{S,y}\longrightarrow\mathcal O_{X,x},
    \qquad
    \theta_{f,x}:=f_x^\sharp\circ\lambda_{R,y}
      \colon R\longrightarrow\mathcal O_{X,x},\\
    \rho_C:=\rho|_C\colon C\twoheadrightarrow A''=C/\mathfrak m',
    \qquad
    \sigma=\operatorname{Spec}(\rho_C),\\
    \overline c_R:=\rho_C\circ c_R\colon R\to A''.
  \end{gathered}                                                   \tag{12a}
  \]
  Composing (9) with $\sigma$ and applying global sections over the affine
  target $S=\operatorname{Spec}(R)$ is the equality
  \[
  \begin{aligned}
    f\circ g''
      &=f\circ x_C\circ\sigma
       =s_C\circ\sigma,\\
    \ell\circ\theta_{f,x}
      &=\rho_C\circ c_R=\overline c_R.
  \end{aligned}                                                  \tag{12b}
  \]
  The definition of $\ell_A$ and the local-ring pullback of (11) are the
  following chain of ring homomorphisms:
  \[
  \begin{aligned}
    \varphi\circ\iota_A\circ\ell_A
      &=\iota_{A'',k'}\circ\ell
       =\varphi\circ\ell_{K,x}.
  \end{aligned}
  \]
  Injectivity of $\varphi$ and the preceding equality are the implication
  chain
  \[
    \bigl(
      \varphi\circ\iota_A\circ\ell_A
        =\varphi\circ\ell_{K,x}
      \ \wedge\ \operatorname{Injective}(\varphi)
    \bigr)
      \Longrightarrow
      \iota_A\circ\ell_A=\ell_{K,x}.                           \tag{12c}
  \]
  Define the set of $K$-points generizing $x$ by
  \[
    \mathscr G_x(K):=
      \left\{q\colon\operatorname{Spec}(K)\to X
        \ \middle|\ x\in
          \overline{\{|q|(\ast_K)\}}\right\}.                 \tag{12c0}
  \]
  For each $q\in\mathscr G_x(K)$, define
  \[
  \begin{gathered}
    y_q:=|q|(\ast_K),
    \qquad
    \gamma^X_{x,y_q}\colon\mathcal O_{X,x}\to\mathcal O_{X,y_q}
      \quad\text{by (11d)},\\
    q_{x,K}^\sharp:=q_{\ast_K}^\sharp\circ\gamma^X_{x,y_q}
      \colon\mathcal O_{X,x}\longrightarrow K,
  \end{gathered}                                                   \tag{12c0a}
  \]
  where $q_{\ast_K}^\sharp$ is the stalk homomorphism of (P-1).  Define
  \[
    \Phi_{x,K}\colon\mathscr G_x(K)\longrightarrow
      \operatorname{Hom}_{\mathrm{Ring}}(\mathcal O_{X,x},K),
    \qquad \Phi_{x,K}(q):=q_{x,K}^\sharp.                       \tag{12c0b}
  \]
  Since $x\in|X_x|$ and $x\in\overline{\{y_q\}}$, one has
  $y_q\in|X_x|$.  The open-immersion factorization therefore defines the
  unique morphism and its affine ring map
  \[
  \begin{gathered}
    q^{(x)}\colon\operatorname{Spec}(K)\to X_x,
    \qquad j_x\circ q^{(x)}=q,\\
    \beta_q:=(q^{(x)})^\#\colon R_x\to K.
  \end{gathered}                                                   \tag{12c0c}
  \]
  For $s\in R_x$, (11d), (12c0a), and (12c0c) are the following equality
  chain:
  \[
  \begin{aligned}
    (q_{x,K}^\sharp\circ\lambda_{R_x,x})(s)
      &=q_{\ast_K}^\sharp
          \bigl(\gamma^X_{x,y_q}([X_x,s]_x)\bigr)\\
      &=q_{\ast_K}^\sharp([X_x,s]_{y_q})
       =\beta_q(s).
  \end{aligned}                                                    \tag{12c0d}
  \]
  Applying $\operatorname{Spec}$ to (12c0d), together with substitutions
  (12-0b) and (12c0c), is the following factorization chain:
  \[
  \begin{aligned}
    \lambda_x\circ\operatorname{Spec}(q_{x,K}^\sharp)
      &=j_x\circ\operatorname{Spec}(\lambda_{R_x,x})
          \circ\operatorname{Spec}(q_{x,K}^\sharp)\\
      &=j_x\circ\operatorname{Spec}
          (q_{x,K}^\sharp\circ\lambda_{R_x,x})\\
      &=j_x\circ\operatorname{Spec}(\beta_q)
       =j_x\circ q^{(x)}=q.
  \end{aligned}                                                    \tag{12c0e}
  \]
  Hence, for $q_1,q_2\in\mathscr G_x(K)$,
  \[
  \begin{aligned}
    \Phi_{x,K}(q_1)=\Phi_{x,K}(q_2)
      &\Longrightarrow
        \operatorname{Spec}((q_1)_{x,K}^\sharp)
          =\operatorname{Spec}((q_2)_{x,K}^\sharp)\\
      &\Longrightarrow
        \lambda_x\circ\operatorname{Spec}((q_1)_{x,K}^\sharp)
          =\lambda_x\circ\operatorname{Spec}((q_2)_{x,K}^\sharp)\\
      &\overset{\text{(12c0e)}}{\Longrightarrow}q_1=q_2.
  \end{aligned}
  \]
  Thus
  \[
    \operatorname{Injective}(\Phi_{x,K}).                        \tag{12c1}
  \]
  The membership $\mathfrak m_A\in\overline{\{(0)\}}$, equations (11c)
  and (12-0b), and $|\nu_A|(\ast_K)=(0)$ are the following specialization
  and membership chain:
  \[
  \begin{aligned}
    x=|x_A|(\mathfrak m_A)
      &\in\overline{\{|x_A|((0))\}}
       =\overline{\{|x_A\circ\nu_A|(\ast_K)\}},\\
    x&\in\overline{\{|x_K|(\ast_K)\}},\\
    x_A\circ\nu_A&\in\mathscr G_x(K),
    \qquad x_K\in\mathscr G_x(K).
  \end{aligned}                                                    \tag{12c2}
  \]
  Contravariance of $\operatorname{Spec}$ and (12-0b) are the following
  equality chain:
  \[
  \begin{aligned}
    x_A\circ\nu_A
      &=\lambda_x\circ\operatorname{Spec}(\ell_A)
          \circ\operatorname{Spec}(\iota_A)\\
      &=\lambda_x\circ\operatorname{Spec}(\iota_A\circ\ell_A).
  \end{aligned}                                                    \tag{12c3}
  \]
  Put $c_A:=\mathfrak m_A$ and $\eta_A:=(0)$.  The construction (12-0b),
  the generization homomorphism (11d), and the stalk definition (P-1) are the
  following equality tuple:
  \[
  \begin{gathered}
    (x_A)_{c_A}^\sharp=\ell_A,\\
    (\nu_A)_{\ast_K}^\sharp\circ
      \gamma^{\operatorname{Spec}(A)}_{c_A,\eta_A}=\iota_A,\\
    \gamma^{\operatorname{Spec}(A)}_{c_A,\eta_A}
      \circ(x_A)_{c_A}^\sharp
      =(x_A)_{\eta_A}^\sharp\circ
        \gamma^X_{x,|x_A|(\eta_A)}.
  \end{gathered}                                                   \tag{12c4}
  \]
  The equality $|x_A\circ\nu_A|(\ast_K)=|x_A|(\eta_A)$, the composition
  law for the stalk maps in (P-1), and (12c4) are the following equality
  chain:
  \[
  \begin{aligned}
    (x_A\circ\nu_A)_{x,K}^\sharp
      &=(x_A\circ\nu_A)_{\ast_K}^\sharp\circ
        \gamma^X_{x,|x_A|(\eta_A)}\\
      &=(\nu_A)_{\ast_K}^\sharp\circ(x_A)_{\eta_A}^\sharp
          \circ\gamma^X_{x,|x_A|(\eta_A)}\\
      &=(\nu_A)_{\ast_K}^\sharp\circ
          \gamma^{\operatorname{Spec}(A)}_{c_A,\eta_A}
          \circ(x_A)_{c_A}^\sharp\\
      &=\iota_A\circ\ell_A,\\
    (x_K)_{x,K}^\sharp
      &=(x_K)_{\ast_K}^\sharp\circ\gamma^X_{x,y_K}
       =\ell_{K,x}.
  \end{aligned}                                                    \tag{12c4a}
  \]
  Substitution of (12c) and (12c4) into (12c0b) is the chain
  \[
  \begin{aligned}
    \Phi_{x,K}(x_A\circ\nu_A)
      &=(x_A\circ\nu_A)_{x,K}^\sharp
       =\iota_A\circ\ell_A\\
      &=\ell_{K,x}
       =(x_K)_{x,K}^\sharp
       =\Phi_{x,K}(x_K).
  \end{aligned}                                                    \tag{12c5}
  \]
  Equations (12c1) and (12c5) are the following implication chain:
  \[
    \Phi_{x,K}(x_A\circ\nu_A)=\Phi_{x,K}(x_K)
      \ \wedge\ \operatorname{Injective}(\Phi_{x,K})
      \Longrightarrow
      x_A\circ\nu_A=x_K.                                       \tag{12d}
  \]

  The identities
  \[
    \iota_{A'',k'}\circ\rho_C
      =\rho\circ\iota_{C,A'},\qquad
    \iota_{C,A'}\circ c_R
      =\alpha^\#\circ\jmath\circ r
  \]
  together with (6), (12a), and (12b) form the structural ring-map chain
  \[
  \begin{aligned}
    \varphi\circ\iota_A\circ\ell_A\circ\theta_{f,x}
      &=\iota_{A'',k'}\circ\ell\circ\theta_{f,x}\\
      &=\iota_{A'',k'}\circ\rho_C\circ c_R\\
      &=\rho\circ\iota_{C,A'}\circ c_R\\
      &=\rho\circ\alpha^\#\circ\jmath\circ r\\
      &=\varphi\circ\iota_A\circ a.
  \end{aligned}                                                  \tag{12e}
  \]
  Injectivity of $\varphi\circ\iota_A$ and (12e) are the following
  implication chain:
  \[
    \bigl(
      \varphi\circ\iota_A\circ\ell_A\circ\theta_{f,x}
        =\varphi\circ\iota_A\circ a
      \ \wedge\
      \operatorname{Injective}(\varphi\circ\iota_A)
    \bigr)
      \Longrightarrow
      \ell_A\circ\theta_{f,x}=a.                              \tag{12f}
  \]
  By \Cref{def:nagata-direct-image-structure-map}, define the affine-spectrum
  bijection
  \[
    \Gamma_{R,A}\colon
      \operatorname{Hom}(\operatorname{Spec}(A),\operatorname{Spec}(R))
      \xrightarrow{\sim}\operatorname{Hom}_{\mathrm{Ring}}(R,A),
    \qquad q\longmapsto q^\#.
  \]
  Applying $\Gamma_{R,A}$ to $f\circ x_A$ and $s_A$, and using (12f), is
  the chain
  \[
  \begin{aligned}
    \Gamma_{R,A}(f\circ x_A)
      &=(f\circ x_A)^\#
      &=\ell_A\circ\theta_{f,x}
       =a=s_A^\#
       =\Gamma_{R,A}(s_A)
      \Longrightarrow f\circ x_A=s_A.
  \end{aligned}                                                  \tag{12g}
  \]
  Combining (12d) and (12g) is the required pair
  \[
    x_A\circ\nu_A=x_K,
    \qquad
    f\circ x_A=s_A.                                              \tag{13}
  \]
  Thus $f$ satisfies the existence part of the valuative criterion.  Since
  $f$ is quasi-compact, the valuative criterion for universal closedness
  (\cite[\href{https://stacks.math.columbia.edu/tag/01KF}{Tag 01KF}]
  {stacks-project}), applied to the datum (13), is the predicate
  \[
    \operatorname{UniversallyClosed}(f).
  \]

  Assume now that $f$ is quasi-separated and put
  \[
    \Delta_f\colon X\to X\times_SX.
  \]
  By (P0), the quasi-separatedness hypothesis is the predicate equality
  \[
    \operatorname{QuasiSeparated}(f)
      \Longleftrightarrow\operatorname{IsQuasiCompact}(\Delta_f). \tag{13a}
  \]
  Let a valuation datum for
  $(\Delta_f,h)$ be given by
  \[
    u_K\colon\operatorname{Spec}(K)\to U,
    \qquad
    y_A=(y_{A,1},y_{A,2})\colon
      \operatorname{Spec}(A)\to X\times_SX,
  \]
  with
  \[
    y_{A,1}\circ\nu_A=h\circ u_K
      =y_{A,2}\circ\nu_A,
    \qquad
    f\circ y_{A,1}=f\circ y_{A,2}.                               \tag{14}
  \]
  Both $y_{A,1}$ and $y_{A,2}$ satisfy the lifting equations for the same
  datum $(A,K,u_K,f\circ y_{A,1})$.  The assumed uniqueness is the equality
  \[
    y_{A,1}=y_{A,2}.                                               \tag{15}
  \]
  By the universal property of the fibre product, (15) defines
  $\widetilde y_A\colon\operatorname{Spec}(A)\to X$ with
  \[
    \Delta_f\circ\widetilde y_A=y_A.
  \]
  Define the lift set for this exact datum by
  \[
    \mathscr L_{\Delta_f}(A,u_K,y_A):=
      \left\{x_A\in\operatorname{Hom}(\operatorname{Spec}A,X)
      \ \middle|\ 
      x_A\circ\nu_A=h\circ u_K\ \wedge\
      \Delta_f\circ x_A=y_A\right\}.                           \tag{15a}
  \]
  The first projection of
  $\Delta_f\circ x_A=y_A=(y_{A,1},y_{A,2})$ is $x_A=y_{A,1}$.
  Together with (14)--(15), this is the singleton equality
  \[
    \mathscr L_{\Delta_f}(A,u_K,y_A)=\{\widetilde y_A\}.         \tag{15b}
  \]
  The universal-closedness implication established in (0a)--(13), applied
  with the exact substitution
  \[
    (X,S,f,h)\longmapsto(X,X\times_SX,\Delta_f,h),
  \]
  and using (13a)--(15b), is
  \[
    \operatorname{UniversallyClosed}(\Delta_f).                 \tag{15c}
  \]
  The diagonal-immersion theorem
  \cite[\href{https://stacks.math.columbia.edu/tag/01KJ}{Tag 01KJ}]
  {stacks-project}, the definition (P0), and the closed-image criterion for
  immersions
  \cite[\href{https://stacks.math.columbia.edu/tag/01IQ}{Tag 01IQ}]
  {stacks-project} are the successive predicates
  \[
  \begin{aligned}
    &\operatorname{IsImmersion}(\Delta_f),\\
    &\operatorname{UniversallyClosed}(\Delta_f)
      \Longrightarrow
      |\Delta_f|(|X|)\text{ is closed in }|X\times_SX|,\\
    &\operatorname{IsImmersion}(\Delta_f)
      \wedge\bigl(|\Delta_f|(|X|)\text{ is closed in }|X\times_SX|\bigr)
      \Longrightarrow\operatorname{IsClosedImmersion}(\Delta_f),\\
    &\operatorname{IsClosedImmersion}(\Delta_f)
      \Longleftrightarrow\operatorname{IsSeparated}(f).
  \end{aligned}                                                   \tag{15d}
  \]
  Substitution of (15c) into (15d) proves
  $\operatorname{IsSeparated}(f)$.
\end{proof}
\section{The two-open compactification theorem}

Throughout this section fix a Noetherian scheme $S$, a separated finite-type
morphism
\[
  s\colon U\longrightarrow S,
\]
and open immersions $u_i\colon U_i\hookrightarrow U$ for $i\in\{1,2\}$ such
that
\[
  U=U_1\cup U_2.
\]
Put
\[
  U_{12}=U_1\times_UU_2,\qquad
  e_i\colon U_{12}\to U_i,\qquad
  e=u_1\circ e_1=u_2\circ e_2\colon U_{12}\to U,
\]
where $e_i$ are the fibre-product projections.  Assume
\[
  \overline{|e|(|U_{12}|)}=|U|.
\]
For $s_i=s\circ u_i$, let
\[
  \operatorname{Comp}(s_i)
    :=\operatorname{Compactification}(s_i)
    \qquad\text{by \Cref{def:nagata-compactification}}.
\]
Assume $\operatorname{Comp}(s_i)\neq\varnothing$ for $i=1,2$.

\begin{definition}[Chosen dense compactifications of the two opens]
  \label{def:nagata-two-open-dense-compactifications}
  \uses{def:nagata-compactifiable,def:nagata-dense-refinement}
  \lean{}

  Using \Cref{def:nagata-choice-notation}, define
  \[
    C_i^0
      :=\operatorname{ch}(\operatorname{Comp}(s_i))
      =\bigl(
          U_i\xrightarrow{j_i^0}X_i^0\xrightarrow{q_i^0}S,\,
          q_i^0\circ j_i^0=s_i
        \bigr).
  \]
  By \Cref{def:nagata-choice-notation}, the selected tuple satisfies
  \[
    C_i^0\in\operatorname{Comp}(s_i).                            \tag{TODC0}
  \]
  Define
  \[
    C_i=\mathsf{Dense}(C_i^0)
      =\bigl(
          U_i\xrightarrow{j_i}X_i\xrightarrow{q_i}S,\,
          q_i\circ j_i=s_i
        \bigr)
  \]
  by \Cref{def:nagata-dense-refinement}.
\end{definition}

\begin{lemma}[Properties of the chosen dense compactifications]
  \label{lemma:nagata-two-open-dense-compactifications}
  \uses{def:nagata-two-open-dense-compactifications}
  \lean{}

  For $i=1,2$,
  \[
    q_i\colon X_i\to S\text{ is proper},\qquad
    j_i\colon U_i\hookrightarrow X_i\text{ is a quasi-compact open immersion},
  \]
  and
  \[
    \ker\bigl(\mathcal O_{X_i}\to(j_i)_*\mathcal O_{U_i}\bigr)=0.
  \]
\end{lemma}

\begin{proof}
  \uses{prop:nagata-dense-refinement}
  For each $i$, the defining equality in
  \Cref{def:nagata-two-open-dense-compactifications} is
  $C_i=\mathsf{Dense}(C_i^0)$.  Applying
  \Cref{prop:nagata-dense-refinement} to $C_i^0$ is the pair
  \[
    C_i\in\operatorname{Comp}(s_i)
    \quad\text{and}\quad
    \ker\bigl(\mathcal O_{X_i}\to(j_i)_*\mathcal O_{U_i}\bigr)=0.
  \]
  Expanding $C_i\in\operatorname{Comp}(s_i)$ by
  \Cref{def:nagata-compactification} is the predicate tuple
  \[
    \operatorname{Proper}(q_i),\qquad
    \operatorname{IsOpenImmersion}(j_i),\qquad
    \operatorname{QuasiCompact}(j_i),\qquad
    q_i\circ j_i=s_i.
  \]
\end{proof}

\begin{definition}[Chosen graph extensions]
  \label{def:nagata-two-open-graph-extensions}
  \uses{def:nagata-two-open-dense-compactifications,
    def:nagata-graph-extension}
  \lean{}

  Before changing notation, define the input symbols
  \[
    X_i^{\mathrm{gr,in}}:=X_i,\qquad
    q_i^{\mathrm{gr,in}}:=q_i\colon X_i^{\mathrm{gr,in}}\to S,\qquad
    j_i^{\mathrm{gr,in}}:=j_i\colon U_i\hookrightarrow X_i^{\mathrm{gr,in}}.
  \]
  Apply \Cref{def:nagata-graph-extension} to
  \[
    U_i\xrightarrow{u_i}U,\qquad
    C_i=(U_i\xrightarrow{j_i^{\mathrm{gr,in}}}X_i^{\mathrm{gr,in}}
      \xrightarrow{q_i^{\mathrm{gr,in}}}S).
  \]
  Denote its output by
  \[
    b_i^{\mathrm{gr}}\colon X_i^{\mathrm{gr}}\to X_i^{\mathrm{gr,in}},\qquad
    \ell_i^{\mathrm{gr}}\colon V_i^{\mathrm{gr}}\hookrightarrow
      X_i^{\mathrm{gr}},\qquad
    \psi_i^{\mathrm{gr}}\colon V_i^{\mathrm{gr}}\to U,
  \]
  and denote the preserved open immersion by
  \[
    \widetilde j_i^{\mathrm{gr}}\colon
      U_i\hookrightarrow V_i^{\mathrm{gr}}.
  \]
  Define the ambient open immersion
  \[
    j_i^{\mathrm{gr}}
      =\ell_i^{\mathrm{gr}}\circ\widetilde j_i^{\mathrm{gr}}
      \colon U_i\hookrightarrow X_i^{\mathrm{gr}}.
  \]
  Define
  \[
    q_i^{\mathrm{gr}}=q_i^{\mathrm{gr,in}}\circ b_i^{\mathrm{gr}}
      \colon X_i^{\mathrm{gr}}\to S.
  \]
  Define the unadorned symbols to be the following abbreviations:
  \[
    (X_i,q_i,V_i,\ell_i,\psi_i,\widetilde j_i,j_i)
      :=
    (X_i^{\mathrm{gr}},q_i^{\mathrm{gr}},V_i^{\mathrm{gr}},
      \ell_i^{\mathrm{gr}},\psi_i^{\mathrm{gr}},
      \widetilde j_i^{\mathrm{gr}},j_i^{\mathrm{gr}}).
  \]
\end{definition}

\begin{lemma}[Properties of the chosen graph extensions]
  \label{lemma:nagata-two-open-graph-extensions}
  \uses{def:nagata-two-open-graph-extensions}
  \lean{}

  For $i=1,2$, $b_i^{\mathrm{gr}}$ is $U_i$-admissible,
  $\psi_i\colon V_i\to U$ is proper, and
  \[
    \psi_i\circ\widetilde j_i=u_i,\qquad
    \psi_i^{-1}(U_i)=\widetilde j_i(U_i),\qquad
    \psi_i|_{\widetilde j_i(U_i)}
      =u_i\circ\widetilde j_i^{-1}\colon
        \widetilde j_i(U_i)\xrightarrow{\sim}U_i
        \hookrightarrow U.
  \]
  Moreover $q_i\colon X_i\to S$ is proper and
  \[
    j_i=\ell_i\circ\widetilde j_i,
    \qquad q_i\circ j_i=s_i.
  \]
\end{lemma}

\begin{proof}
  \uses{lemma:nagata-graph-extension,lemma:nagata-blowup-proper}
  Applying \Cref{lemma:nagata-graph-extension} to the tuple in
  \Cref{def:nagata-two-open-graph-extensions} is the predicate pair
  \[
    b_i^{\mathrm{gr}}\text{ is }U_i\text{-admissible},\qquad
    \psi_i\text{ is proper},
  \]
  together with the three displayed restriction equalities.  By
  \Cref{lemma:nagata-blowup-proper}, $b_i^{\mathrm{gr}}$ is proper.  The
  defining equality in \Cref{def:nagata-two-open-graph-extensions} is
  \[
    q_i^{\mathrm{gr}}=q_i^{\mathrm{gr,in}}\circ b_i^{\mathrm{gr}}.
  \]
  Both factors on the right are proper.  Stability under composition makes
  $q_i^{\mathrm{gr}}$, denoted $q_i$ after the explicit substitution, proper.
  The preserved-open equations in
  \Cref{def:nagata-graph-extension} are
  \[
    b_i^{\mathrm{gr}}\circ j_i^{\mathrm{gr}}=j_i^{\mathrm{gr,in}},
    \qquad
    q_i^{\mathrm{gr,in}}\circ j_i^{\mathrm{gr,in}}=s_i.
  \]
  Therefore
  \[
  \begin{aligned}
    q_i^{\mathrm{gr}}\circ j_i^{\mathrm{gr}}
      &=q_i^{\mathrm{gr,in}}\circ b_i^{\mathrm{gr}}
          \circ j_i^{\mathrm{gr}}\\
      &=q_i^{\mathrm{gr,in}}\circ j_i^{\mathrm{gr,in}}
       =s_i.
  \end{aligned}
  \]
\end{proof}

\begin{definition}[The four boundary loci]
  \label{def:nagata-two-open-boundary-loci}
  \uses{def:nagata-two-open-graph-extensions,
    def:nagata-reduced-closed-subscheme}
  \lean{}

  For $\{i,j\}=\{1,2\}$, let
  \[
  \begin{gathered}
    F_i:=|U|\setminus|U_j|,\qquad
    \mathcal I_i:=\mathcal I_{U,F_i},\\
    Z_i:=(F_i)_{\mathrm{red}}^U
      =\operatorname{Spec}_U(\mathcal O_U/\mathcal I_i),
    \qquad
    z_i:=i_{U,F_i}\colon Z_i\hookrightarrow U.
  \end{gathered}
  \]
  The symbols $\mathcal I_{U,F_i}$ and $i_{U,F_i}$ are those of
  \Cref{def:nagata-reduced-closed-subscheme}.  For $k\in\{1,2\}$ define
  \[
    Z_{i,k}=V_i\times_UZ_k,\qquad
    a_{i,k}\colon Z_{i,k}\to V_i,\qquad
    \psi_{i,k}\colon Z_{i,k}\to Z_k
  \]
  by the cartesian square
  \[
    \begin{tikzcd}
      Z_{i,k} \arrow[r,"\psi_{i,k}"] \arrow[d,"a_{i,k}"']
        & Z_k \arrow[d,"z_k"]\\
      V_i \arrow[r,"\psi_i"'] & U.
    \end{tikzcd}
  \]
\end{definition}

\begin{lemma}[Closedness and disjointness of the four boundary loci]
  \label{lemma:nagata-boundary-loci}
  \lean{}

  For $\{i,j\}=\{1,2\}$,
  \[
    |Z_i|=|U|\setminus|U_j|\subseteq|U_i|,
    \qquad
    Z_{i,i}\times_{V_i}Z_{i,j}=\varnothing.
  \]
  Each $a_{i,k}\colon Z_{i,k}\to V_i$ is a closed immersion, and
  $\psi_i$ is an isomorphism over the open $U_i\subseteq U$, which contains
  $Z_i$.
\end{lemma}

\begin{proof}
  \uses{lemma:nagata-two-open-graph-extensions}
  Since $U=U_i\cup U_j$,
  \[
    |Z_i|=|U|\setminus|U_j|
      \subseteq |U_i|.
  \]
  The maps $a_{i,k}$ are base changes of the closed immersions
  $z_k\colon Z_k\hookrightarrow U$, so they are closed immersions.  Moreover
  \[
    Z_i\times_UZ_j
      =(U\setminus U_j)\times_U(U\setminus U_i)
      =\varnothing.
  \]
  Base change of this equality along $\psi_i$ is
  \[
    Z_{i,i}\times_{V_i}Z_{i,j}
      =V_i\times_U(Z_i\times_UZ_j)=\varnothing.
  \]
  Finally, \Cref{lemma:nagata-two-open-graph-extensions} states
  \[
    \psi_i^{-1}(U_i)=\widetilde j_i(U_i),\qquad
    \psi_i|_{\widetilde j_i(U_i)}\colon
      \widetilde j_i(U_i)\xrightarrow{\sim}U_i.
  \]
  Since $Z_i\subseteq U_i$, the last equality is the asserted
  neighbourhood statement.
\end{proof}

\begin{definition}[Prepared ambient schemes with separated boundary closures]
  \label{def:nagata-two-open-prepared-boundaries}
  \uses{def:nagata-two-open-boundary-loci,
    def:nagata-separate-constructible-closed}
  \lean{}

  For $\{i,j\}=\{1,2\}$, first define the input-stage symbols
  \[
  \begin{gathered}
    X_i^{\mathrm{prep,in}}:=X_i,\qquad
    q_i^{\mathrm{prep,in}}:=q_i,\qquad
    V_i^{\mathrm{prep,in}}:=V_i,\qquad
    \ell_i^{\mathrm{prep,in}}:=\ell_i,\\
    \psi_i^{\mathrm{prep,in}}:=\psi_i,\qquad
    \widetilde j_i^{\mathrm{prep,in}}:=\widetilde j_i,\qquad
    j_i^{\mathrm{prep,in}}:=j_i,\qquad
    Z_{i,k}^{\mathrm{prep,in}}:=Z_{i,k}.
  \end{gathered}
  \]
  Apply
  \Cref{def:nagata-separate-constructible-closed} to the disjoint closed
  subsets
  $|Z_{i,i}^{\mathrm{prep,in}}|,|Z_{i,j}^{\mathrm{prep,in}}|
  \subseteq|V_i^{\mathrm{prep,in}}|$.  Denote the output by
  \[
    b_i^{\mathrm{prep}}\colon
      X_i^{\mathrm{prep}}
      =\mathsf{CSepBl}_{V_i}(|Z_{i,i}|,|Z_{i,j}|)
      \longrightarrow X_i^{\mathrm{prep,in}}.
  \]
  Put
  \[
    V_i^{\mathrm{prep}}
      =X_i^{\mathrm{prep}}\times_{X_i^{\mathrm{prep,in}}}
        V_i^{\mathrm{prep,in}},
    \qquad
    \ell_i^{\mathrm{prep}}:=\operatorname{pr}_{X_i^{\mathrm{prep}}}
      \colon V_i^{\mathrm{prep}}\hookrightarrow X_i^{\mathrm{prep}}.
  \]
  and denote the selected preserved-open isomorphism by
  \[
    \beta_i\colon V_i^{\mathrm{prep}}\xrightarrow{\sim}
      V_i^{\mathrm{prep,in}},\qquad
    \ell_i^{\mathrm{prep,in}}\circ\beta_i
      =b_i^{\mathrm{prep}}|_{V_i^{\mathrm{prep}}}.
  \]
  Define
  \[
    \psi_i^{\mathrm{prep}}=\psi_i^{\mathrm{prep,in}}\circ\beta_i,\qquad
    \widetilde j_i^{\mathrm{prep}}
      =\beta_i^{-1}\circ\widetilde j_i^{\mathrm{prep,in}},
    \qquad
    j_i^{\mathrm{prep}}
      =\ell_i^{\mathrm{prep}}\circ\widetilde j_i^{\mathrm{prep}},
    \qquad
    Z_{i,k}^{\mathrm{prep}}
      =V_i^{\mathrm{prep}}\times_{V_i^{\mathrm{prep,in}}}
        Z_{i,k}^{\mathrm{prep,in}},\qquad
    q_i^{\mathrm{prep}}=q_i^{\mathrm{prep,in}}\circ b_i^{\mathrm{prep}}.
  \]
  Substitution of the preserved-open equality
  $b_i^{\mathrm{prep}}\circ\ell_i^{\mathrm{prep}}
    =\ell_i^{\mathrm{prep,in}}\circ\beta_i$ is the chain
  \[
  \begin{aligned}
    q_i^{\mathrm{prep}}\circ j_i^{\mathrm{prep}}
      &=q_i^{\mathrm{prep,in}}\circ b_i^{\mathrm{prep}}
          \circ\ell_i^{\mathrm{prep}}
          \circ\beta_i^{-1}\circ\widetilde j_i\\
      &=q_i^{\mathrm{prep,in}}\circ\ell_i^{\mathrm{prep,in}}
          \circ\widetilde j_i^{\mathrm{prep,in}}
       =q_i^{\mathrm{prep,in}}\circ j_i^{\mathrm{prep,in}}
       =s_i.
  \end{aligned}
  \]
  By \Cref{def:nagata-reduced-closed-subscheme}, define
  \[
  \begin{aligned}
    F_{i,k}^{\mathrm{prep}}
      &:=\overline{|Z_{i,k}^{\mathrm{prep}}|}^{\,|X_i^{\mathrm{prep}}|},\\
    \overline Z_{i,k}^{\mathrm{prep}}
      &:=(F_{i,k}^{\mathrm{prep}})_{\mathrm{red}}^{X_i^{\mathrm{prep}}}
        =\operatorname{Spec}_{X_i^{\mathrm{prep}}}
          \!\left(\mathcal O_{X_i^{\mathrm{prep}}}/
            \mathcal I_{X_i^{\mathrm{prep}},F_{i,k}^{\mathrm{prep}}}\right),\\
    \overline z_{i,k}^{\mathrm{prep}}
      &:=i_{X_i^{\mathrm{prep}},F_{i,k}^{\mathrm{prep}}}
        \colon\overline Z_{i,k}^{\mathrm{prep}}
          \hookrightarrow X_i^{\mathrm{prep}}.
  \end{aligned}
  \]
  Make the explicit substitution
  \[
    (X_i,q_i,V_i,\ell_i,\psi_i,\widetilde j_i,j_i,
      Z_{i,k},\overline Z_{i,k},\overline z_{i,k})
      :=
    (X_i^{\mathrm{prep}},q_i^{\mathrm{prep}},
      V_i^{\mathrm{prep}},\ell_i^{\mathrm{prep}},
      \psi_i^{\mathrm{prep}},\widetilde j_i^{\mathrm{prep}},
      j_i^{\mathrm{prep}},
      Z_{i,k}^{\mathrm{prep}},\overline Z_{i,k}^{\mathrm{prep}},
      \overline z_{i,k}^{\mathrm{prep}}).
  \]
\end{definition}

\begin{lemma}[Properties of the prepared boundary closures]
  \label{lemma:nagata-prepared-boundaries}
  \uses{def:nagata-two-open-prepared-boundaries}
  \lean{}

  For $\{i,j\}=\{1,2\}$,
  \[
    b_i^{\mathrm{prep}}\text{ is }V_i\text{-admissible},\qquad
    \overline Z_{i,i}\times_{X_i}\overline Z_{i,j}=\varnothing.
  \]
  Under $\beta_i$, the morphism on the preserved open is exactly
  \[
    \psi_i^{\mathrm{prep}}
      =\psi_i^{\mathrm{prep,in}}\circ\beta_i.
  \]
\end{lemma}

\begin{proof}
  \uses{lemma:nagata-boundary-loci,
    lemma:nagata-separate-constructible-closed}
  By \Cref{lemma:nagata-boundary-loci},
  $|Z_{i,i}|$ and $|Z_{i,j}|$ are disjoint constructible subsets closed in
  $|V_i|$.  Substitution into
  \Cref{lemma:nagata-separate-constructible-closed} is the pair
  \[
    b_i^{\mathrm{prep}}\text{ is }V_i\text{-admissible}
  \]
  and
  \[
    \overline{|\beta_i^{-1}(Z_{i,i})|}
      \cap
    \overline{|\beta_i^{-1}(Z_{i,j})|}
    =\varnothing.
  \]
  By \Cref{def:nagata-two-open-prepared-boundaries}, these two closures are
  $|\overline Z_{i,i}|$ and $|\overline Z_{i,j}|$, so their fibre product is
  empty.  The equality
  $\psi_i^{\mathrm{prep}}
    =\psi_i^{\mathrm{prep,in}}\circ\beta_i$ is the defining
  equality in that same definition.
\end{proof}

\begin{lemma}[Disjoint boundary closures persist]
  \label{lemma:nagata-boundary-disjoint-persists}
  \lean{}

  Let $b\colon X'\to X_i$ be a $V_i$-admissible blow-up.  For
  $k\in\{i,j\}$ let $\overline Z'_{i,k}\hookrightarrow X'$ be the strict
  transform of $\overline Z_{i,k}\hookrightarrow X_i$.  If
  $\overline Z_{i,i}\times_{X_i}\overline Z_{i,j}=\varnothing$, then
  \[
    \overline Z'_{i,i}\times_{X'}\overline Z'_{i,j}=\varnothing.
  \]
\end{lemma}

\begin{proof}
  \uses{lemma:nagata-strict-transform-restrict}
  By the closed-subscheme definition of strict transform,
  \[
    \overline Z'_{i,k}
      \hookrightarrow X'\times_{X_i}\overline Z_{i,k}
      \qquad(k=i,j).
  \]
  Therefore
  \[
  \begin{aligned}
    \overline Z'_{i,i}\times_{X'}\overline Z'_{i,j}
      &\hookrightarrow
        (X'\times_{X_i}\overline Z_{i,i})
          \times_{X'}
        (X'\times_{X_i}\overline Z_{i,j})\\
      &\cong
        X'\times_{X_i}
          (\overline Z_{i,i}\times_{X_i}\overline Z_{i,j})\\
      &=X'\times_{X_i}\varnothing
       =\varnothing.
  \end{aligned}
  \]
\end{proof}

\begin{definition}[The two-open comparison space]
  \label{def:nagata-two-open-comparison-space}
  \uses{def:nagata-two-open-prepared-boundaries,
    def:nagata-scheme-theoretic-image,
    def:nagata-scheme-image-factor,
    lemma:nagata-two-open-graph-extensions,
    lemma:nagata-prepared-boundaries,
    lemma:nagata-blowup-proper,
    lemma:nagata-morphism-property-calculus}
  \lean{}

  Define
  \[
    V_{12}=V_1\times_UV_2,\qquad
    r_i\colon V_{12}\to V_i,\qquad
    \psi_{12}=\psi_1\circ r_1=\psi_2\circ r_2.
  \]
  Define
  \[
    \delta_{12}:=(r_1,r_2)\colon V_{12}\to V_1\times_SV_2,
    \qquad
    \operatorname{pr}_i\circ\delta_{12}=r_i\quad(i=1,2).
  \]
  Since $s\colon U\to S$ is separated,
  $\delta_{12}$ is the base change of $\Delta_{U/S}$.  Thus
  \[
    \operatorname{IsSeparated}(s)
      \Longleftrightarrow
        \operatorname{IsClosedImmersion}(\Delta_{U/S})
      \Longrightarrow
        \operatorname{IsClosedImmersion}(\delta_{12})
      \Longrightarrow
        \operatorname{IsQCQS}(\delta_{12}).
  \]
  By \Cref{lemma:nagata-two-open-graph-extensions,def:nagata-two-open-prepared-boundaries,lemma:nagata-blowup-proper},
  the exact predicate chain is
  \[
  \begin{aligned}
    \operatorname{IsProper}(q_i^{\mathrm{prep,in}})\wedge
      \operatorname{IsProper}(b_i^{\mathrm{prep}})
      &\Longrightarrow
        \operatorname{IsProper}
          (q_i^{\mathrm{prep,in}}\circ b_i^{\mathrm{prep}})
       =\operatorname{IsProper}(q_i)\\
      &\Longrightarrow\operatorname{FiniteType}(q_i),\\
    \operatorname{Noetherian}(S)\wedge\operatorname{FiniteType}(q_i)
      &\Longrightarrow\operatorname{Noetherian}(X_i)
       &&\text{by \cite[Tag 01T6]{stacks-project}},\\
    \operatorname{Noetherian}(X_i)\wedge
      \operatorname{IsOpenImmersion}(\ell_i)
      &\Longrightarrow\operatorname{IsQuasiCompact}(\ell_i),\\
    \operatorname{IsOpenImmersion}(\ell_i)\wedge
      \operatorname{IsQuasiCompact}(\ell_i)
      &\Longrightarrow\operatorname{IsQCQS}(\ell_i)
      \qquad(i=1,2),\\
    \operatorname{IsQCQS}(\ell_1)\wedge
      \operatorname{IsQCQS}(\ell_2)
      &\Longrightarrow
        \operatorname{IsQCQS}(\ell_1\times_S\ell_2).
  \end{aligned}
  \]
  With
  $\ell_1\times_S\ell_2\colon V_1\times_SV_2\hookrightarrow X_1\times_SX_2$,
  define
  \[
    v_{12}
      =(\ell_1\times_S\ell_2)\circ\delta_{12}
      \colon V_{12}\to X_1\times_SX_2.
  \]
  Composition stability in
  \Cref{lemma:nagata-morphism-property-calculus} is the implication
  \[
    \operatorname{IsQCQS}(\delta_{12})\wedge
      \operatorname{IsQCQS}(\ell_1\times_S\ell_2)
      \Longrightarrow
      \operatorname{IsQCQS}
        ((\ell_1\times_S\ell_2)\circ\delta_{12})
      =\operatorname{IsQCQS}(v_{12}).
  \]
  With $i_{v_{12}}$ and $v_{12}^{\mathrm{im}}$ as in
  \Cref{def:nagata-scheme-theoretic-image,def:nagata-scheme-image-factor},
  define
  \[
    X_{12}:=\operatorname{SIm}(v_{12}),\qquad
    d_{12}:=v_{12}^{\mathrm{im}}\colon V_{12}\to X_{12},\qquad
    \iota_{12}:=i_{v_{12}}\colon X_{12}\hookrightarrow X_1\times_SX_2.
  \]
  The factorization equality of
  \Cref{def:nagata-scheme-image-factor} is
  \[
    \iota_{12}\circ d_{12}=v_{12}.
  \]
  For $i=1,2$, define
  \[
    p_i=\operatorname{pr}_i\circ\iota_{12}\colon X_{12}\to X_i.
  \]
\end{definition}

\begin{lemma}[Properness of the comparison projections]
  \label{lemma:nagata-comparison-projections-proper}
  \lean{}

  The morphisms $p_i\colon X_{12}\to X_i$ are proper for $i=1,2$.
\end{lemma}

\begin{proof}
  By \Cref{def:nagata-two-open-comparison-space},
  \[
    p_i=\operatorname{pr}_i\circ\iota_{12}.
  \]
  The map $\iota_{12}$ is a closed immersion, hence proper.  If
  $\{i,j\}=\{1,2\}$, the projection
  \[
    \operatorname{pr}_i\colon X_1\times_SX_2\to X_i
  \]
  is the base change of the proper morphism $q_j\colon X_j\to S$, so it is
  proper.  Stability under composition is the implication
  \[
    \operatorname{Proper}(\iota_{12})
      \wedge\operatorname{Proper}(\operatorname{pr}_i)
      \Longrightarrow\operatorname{Proper}(p_i).
  \]
\end{proof}

\begin{definition}[Distinguished comparison boundary]
  \label{def:nagata-two-open-comparison-boundary}
  \uses{def:nagata-closed-subscheme-class,
    def:nagata-two-open-boundary-loci,
    def:nagata-two-open-comparison-space}
  \lean{}

  Define the three fibre products and their closed immersions by
  \[
  \begin{gathered}
    Z_{12,2}:=V_{12}\times_{V_1}Z_{1,2},
    \qquad
    c_{12,2}^{(1)}:=\operatorname{pr}_{V_{12}}\colon
      Z_{12,2}\hookrightarrow V_{12},\\
    Z_{12,2}^{(2)}:=V_{12}\times_{V_2}Z_{2,2},
    \qquad
    c_{12,2}^{(2)}:=\operatorname{pr}_{V_{12}}\colon
      Z_{12,2}^{(2)}\hookrightarrow V_{12},\\
    Z_{12,2}^{(U)}:=V_{12}\times_UZ_2,
    \qquad
    c_{12,2}^{(U)}:=\operatorname{pr}_{V_{12}}\colon
      Z_{12,2}^{(U)}\hookrightarrow V_{12}.
  \end{gathered}
  \]
  Here the first fibre product uses $r_1$ and $a_{1,2}$, the second uses
  $r_2$ and $a_{2,2}$, and the third uses the structure morphisms to $U$.
  Fibre-product associativity is the following class chain:
  \[
    [\,Z_{12,2}\xhookrightarrow{c_{12,2}^{(1)}}V_{12}\,]
      =[\,Z_{12,2}^{(2)}\xhookrightarrow{c_{12,2}^{(2)}}V_{12}\,]
      =[\,Z_{12,2}^{(U)}\xhookrightarrow{c_{12,2}^{(U)}}V_{12}\,]
      \quad\text{in }\operatorname{ClosedSub}(V_{12}).          \tag{CB0a}
  \]
\end{definition}

\begin{lemma}[Restriction of the comparison image over the proper-extension open]
  \label{lemma:nagata-comparison-restrict-proper-open}
  \lean{}

  For $\{i,j\}=\{1,2\}$, put
  \[
    X_{12,i}=X_{12}\times_{X_i}V_i
  \]
  and let $a_i^{12}\colon V_{12}\to X_{12,i}$ be the unique map satisfying
  \[
    \operatorname{pr}_{X_{12}}\circ a_i^{12}=d_{12},\qquad
    \operatorname{pr}_{V_i}\circ a_i^{12}=r_i.
  \]
  Then $a_i^{12}$ is an isomorphism.  Under this isomorphism,
  $p_i|_{X_{12,i}}$ is the base change of
  $\psi_j\colon V_j\to U$ along $\psi_i\colon V_i\to U$.
\end{lemma}

\begin{proof}
  \uses{def:nagata-two-open-comparison-space,
    def:nagata-choice-notation,
    def:nagata-closed-subscheme-class,
    def:nagata-scheme-image-factor,
    lemma:nagata-scheme-image-restrict-open,
    lemma:nagata-proper-to-separated-target}
  Only after (IC3b), apply \Cref{def:nagata-scheme-theoretic-image} and set
  \[
    O_i=V_i\times_SX_j\subseteq X_i\times_SX_j,\qquad
    e_i^{12}=(r_i,\ell_j\circ r_j)\colon V_{12}\to O_i.
  \]
  The factorization
  \[
    V_{12}\xrightarrow{\delta_{12}}
      V_i\times_SV_j
      \xrightarrow{\mathrm{id}_{V_i}\times\ell_j}
      V_i\times_SX_j
  \]
  is the predicate
  \[
    \operatorname{IsImmersion}(e_i^{12}).                         \tag{16b}
  \]
  Define the structural projection and record
  \[
    s_{O_i}:=\operatorname{pr}_{V_i}\colon O_i\to V_i,
    \qquad s_{O_i}\circ e_i^{12}=r_i.
  \]
  Base-change stability for the two explicitly identified cartesian squares
  is the pair of implications
  \[
  \begin{aligned}
    \operatorname{IsProper}(\psi_j)
      &\Longrightarrow\operatorname{IsProper}(r_i),\\
    \operatorname{IsSeparated}(q_j)
      &\Longrightarrow\operatorname{IsSeparated}(s_{O_i}).
  \end{aligned}                                                   \tag{16c}
  \]
  Applying \Cref{lemma:nagata-proper-to-separated-target} with
  \[
    (f,g,h)=(e_i^{12},s_{O_i},r_i)
  \]
  and substituting $s_{O_i}\circ e_i^{12}=r_i$ and (16c) is the predicate
  \[
    \operatorname{IsProper}(e_i^{12}).                            \tag{16d}
  \]
  The proper-immersion criterion
  \cite[\href{https://stacks.math.columbia.edu/tag/01IQ}{Tag 01IQ}]
  {stacks-project}, applied to (16b) and (16d), is
  \[
    \operatorname{IsClosedImmersion}(e_i^{12}).                  \tag{16e}
  \]
  In particular,
  \[
    \operatorname{IsClosedImmersion}(e_i^{12})
      \Longrightarrow\operatorname{IsQCQS}(e_i^{12}).            \tag{16e0}
  \]
  Define the base-change closed immersion
  \[
  \begin{gathered}
    \kappa_{12,i}:=
      \bigl(\operatorname{pr}_{V_i},
        p_j\circ\operatorname{pr}_{X_{12}}\bigr)
      \colon X_{12,i}=X_{12}\times_{X_i}V_i\hookrightarrow O_i,\\
    \operatorname{pr}_{V_i}\circ\kappa_{12,i}
      =\operatorname{pr}_{V_i},\qquad
    \operatorname{pr}_{X_j}\circ\kappa_{12,i}
      =p_j\circ\operatorname{pr}_{X_{12}}.
  \end{gathered}                                                   \tag{16e1}
  \]
  Since $|v_{12}|^{-1}(|O_i|)=|V_{12}|$, the conclusion of
  \Cref{lemma:nagata-scheme-image-restrict-open}, applied after (16e0), is
  the class equality
  \[
    [\,X_{12,i}\xhookrightarrow{\kappa_{12,i}}O_i\,]
      =[\,\operatorname{SIm}_{O_i}(e_i^{12})
          \xhookrightarrow{i_{e_i^{12}}}O_i\,]
      \quad\text{in }\operatorname{ClosedSub}(O_i).              \tag{16}
  \]
  By \Cref{def:nagata-scheme-image-factor}, define the scheme-image factor
  and its factorization equality by
  \[
    \theta_i:=(e_i^{12})^{\mathrm{im}}
      \colon V_{12}\longrightarrow\operatorname{SIm}(e_i^{12}),
    \qquad
    i_{e_i^{12}}\circ\theta_i=e_i^{12}.                            \tag{16a}
  \]
  Equation (16e) and the scheme-image formula for a closed immersion are the
  predicate
  \[
    \operatorname{IsIso}(\theta_i).
  \]
  Define the comparison-isomorphism set
  \[
    \mathscr Z_i:=\left\{z\in
      \operatorname{Hom}(\operatorname{SIm}(e_i^{12}),X_{12,i})
      \ \middle|\ 
      \operatorname{IsIso}(z)\ \wedge\
      \kappa_{12,i}\circ z=i_{e_i^{12}}\right\}.
  \]
  By (CSub3) of \Cref{def:nagata-closed-subscheme-class}, equality (16) is
  the existence statement
  \[
    \exists z\colon\operatorname{SIm}_{O_i}(e_i^{12})
      \xrightarrow{\sim}X_{12,i},\qquad
      \kappa_{12,i}\circ z=i_{e_i^{12}},
    \qquad
    \mathscr Z_i\neq\varnothing.                                \tag{16f0}
  \]
  By \Cref{def:nagata-choice-notation}, define
  \[
    \zeta_i:=\operatorname{ch}(\mathscr Z_i)
      \in\mathscr Z_i.
  \]
  Its membership expands to
  \[
    \zeta_i\colon\operatorname{SIm}_{O_i}(e_i^{12})
      \xrightarrow{\sim}X_{12,i},\qquad
    \kappa_{12,i}\circ\zeta_i=i_{e_i^{12}}.                     \tag{16f}
  \]
  If $z\in\mathscr Z_i$, then
  \[
    \kappa_{12,i}\circ z=i_{e_i^{12}}
      =\kappa_{12,i}\circ\zeta_i
      \Longrightarrow z=\zeta_i,
  \]
  where the implication is cancellation against the monomorphism
  $\kappa_{12,i}$.  Hence
  \[
    \mathscr Z_i=\{\zeta_i\}.                                    \tag{16f1}
  \]
  Define the map in the statement by
  \[
    a_i^{12}=\zeta_i\circ\theta_i.
  \]
  The factorization equality in
  \Cref{def:nagata-two-open-comparison-space} is the pair of component
  equalities
  \[
  \begin{aligned}
    p_i\circ d_{12}
      &=\operatorname{pr}_i\circ\iota_{12}\circ d_{12}
       =\operatorname{pr}_i\circ v_{12}
       =\ell_i\circ r_i,\\
    p_j\circ d_{12}
      &=\operatorname{pr}_j\circ\iota_{12}\circ d_{12}
       =\operatorname{pr}_j\circ v_{12}
       =\ell_j\circ r_j.
  \end{aligned}                                                   \tag{16g}
  \]
  The first equality in (16g) defines the fibre-product morphism
  \[
    \overline a_i^{12}:=(d_{12},r_i)
      \colon V_{12}\longrightarrow
        X_{12}\times_{X_i}V_i=X_{12,i}.                         \tag{16h}
  \]
  The two projections of the morphisms
  $\kappa_{12,i}\circ\overline a_i^{12}$ and $e_i^{12}$ to the factors of
  $O_i=V_i\times_SX_j$ are, by (16g),
  \[
  \begin{aligned}
    \operatorname{pr}_{V_i}\circ\kappa_{12,i}\circ
      \overline a_i^{12}
      &=r_i
       =\operatorname{pr}_{V_i}\circ e_i^{12},\\
    \operatorname{pr}_{X_j}\circ\kappa_{12,i}\circ
      \overline a_i^{12}
      &=p_j\circ d_{12}
       =\ell_j\circ r_j
       =\operatorname{pr}_{X_j}\circ e_i^{12}.
  \end{aligned}
  \]
  Fibre-product uniqueness, (16a), and (16f) are the chain
  \[
  \begin{aligned}
    \kappa_{12,i}\circ\overline a_i^{12}
      &=e_i^{12},\\
    \kappa_{12,i}\circ a_i^{12}
      &=\kappa_{12,i}\circ\zeta_i\circ\theta_i
       =i_{e_i^{12}}\circ\theta_i
       =e_i^{12}
       =\kappa_{12,i}\circ\overline a_i^{12}.
  \end{aligned}
  \]
  Since $\kappa_{12,i}$ is a monomorphism and both $\zeta_i$ and $\theta_i$
  are isomorphisms, cancellation and stability under composition are
  \[
    a_i^{12}=\overline a_i^{12}=(d_{12},r_i),
    \qquad
    \operatorname{IsIso}(a_i^{12}).                             \tag{16i}
  \]
  The two projections of (16i), together with the defining equality of
  $V_{12}$, are the following component equalities:
  \[
    \operatorname{pr}_{X_{12}}\circ a_i^{12}=d_{12},\qquad
    \operatorname{pr}_{V_i}\circ a_i^{12}=r_i,\qquad
    V_{12}=V_i\times_UV_j.
  \]
  The last equality is the cartesian square identifying
  $p_i|_{X_{12,i}}$ with the base change of $\psi_j$ along $\psi_i$.
\end{proof}

\begin{definition}[Improved first comparison projection]
  \label{def:nagata-two-open-improved-comparison}
  \uses{def:nagata-two-open-comparison-space,
    def:nagata-two-open-comparison-boundary,
    def:nagata-isomorphism-near-closure,
    def:nagata-strict-transform,
    lemma:nagata-comparison-restrict-proper-open,
    lemma:nagata-scheme-image-strict-transform}
  \lean{}

  Define the input-stage symbols
  \[
  \begin{gathered}
    X_1^{\mathrm{cmp,in}}:=X_1,\qquad
    q_1^{\mathrm{cmp,in}}:=q_1,\qquad
    V_1^{\mathrm{cmp,in}}:=V_1,\qquad
    \ell_1^{\mathrm{cmp,in}}:=\ell_1,\\
    \psi_1^{\mathrm{cmp,in}}:=\psi_1,\qquad
    \widetilde j_1^{\mathrm{cmp,in}}:=\widetilde j_1,\qquad
    j_1^{\mathrm{cmp,in}}:=j_1,\\
    X_{12}^{\mathrm{cmp,in}}:=X_{12},\qquad
    d_{12}^{\mathrm{cmp,in}}:=d_{12},\qquad
    \iota_{12}^{\mathrm{cmp,in}}:=\iota_{12},\\
    p_i^{\mathrm{cmp,in}}:=p_i,\qquad
    Z_{1,k}^{\mathrm{cmp,in}}:=Z_{1,k},\qquad
    \overline Z_{1,k}^{\mathrm{cmp,in}}:=\overline Z_{1,k},\qquad
    \overline z_{1,k}^{\mathrm{cmp,in}}:=\overline z_{1,k}\colon
      \overline Z_{1,k}^{\mathrm{cmp,in}}
        \hookrightarrow X_1^{\mathrm{cmp,in}},\qquad
    Z_{12,2}^{\mathrm{cmp,in}}:=Z_{12,2}.
  \end{gathered}                                                   \tag{CMP0}
  \]
  For every open subscheme $N\subseteq U$, define
  \[
    \psi_{2,N}:=\psi_2|_{\psi_2^{-1}(N)}\colon
      \psi_2^{-1}(N)\longrightarrow N.
  \]
  Define
  \[
    \mathscr N_2
      =\left\{
          N\subseteq U\ \middle|\
          N\text{ is open and }\psi_{2,N}
          \text{ is an isomorphism}
        \right\},
    \qquad
    N_2=\bigcup_{N\in\mathscr N_2}N.
  \]
  Apply \Cref{def:nagata-isomorphism-near-closure} to
  \[
    p_1^{\mathrm{cmp,in}}\colon X_{12}^{\mathrm{cmp,in}}
      \to X_1^{\mathrm{cmp,in}},\qquad
    (p_1^{\mathrm{cmp,in}})^{-1}(V_1^{\mathrm{cmp,in}})=V_{12},\qquad
    |Z_{1,2}^{\mathrm{cmp,in}}|\subseteq|V_1^{\mathrm{cmp,in}}|,
  \]
  and denote the output by
  \[
    b_1^{\mathrm{cmp}}\colon
      X_1^{\mathrm{cmp}}
      =\mathsf{IsoBl}_{V_1^{\mathrm{cmp,in}},
        |Z_{1,2}^{\mathrm{cmp,in}}|}(p_1^{\mathrm{cmp,in}})
      \longrightarrow X_1^{\mathrm{cmp,in}}.                     \tag{C1}
  \]
  Define
  \[
    P^{\mathrm{cmp,in}}:=X_1^{\mathrm{cmp,in}}\times_SX_2,\qquad
    P^{\mathrm{cmp}}:=X_1^{\mathrm{cmp}}\times_SX_2,\qquad
    \pi_P^{\mathrm{cmp}}:=b_1^{\mathrm{cmp}}\times_S
      \operatorname{id}_{X_2}\colon P^{\mathrm{cmp}}\to P^{\mathrm{cmp,in}},
  \]
  and form the strict transform, in the sense of
  \Cref{def:nagata-strict-transform},
  \[
    X_{12}^{\mathrm{cmp}}
      :=\operatorname{ST}_{\pi_P^{\mathrm{cmp}}}
          (X_{12}^{\mathrm{cmp,in}}),
    \qquad
    \iota_{12}^{\mathrm{st,cmp}}\colon X_{12}^{\mathrm{cmp}}
      \hookrightarrow
      X_{12}^{\mathrm{cmp,in}}\times_{P^{\mathrm{cmp,in}}}P^{\mathrm{cmp}}.
                                                                    \tag{C2}
  \]
  Define its structure maps by
  \[
  \begin{aligned}
    b_{12}^{\mathrm{cmp}}
      &:=\operatorname{pr}_{X_{12}^{\mathrm{cmp,in}}}
          \circ\iota_{12}^{\mathrm{st,cmp}}
        \colon X_{12}^{\mathrm{cmp}}\to X_{12}^{\mathrm{cmp,in}},\\
    \iota_{12}^{\mathrm{cmp}}
      &:=\operatorname{pr}_{P^{\mathrm{cmp}}}
          \circ\iota_{12}^{\mathrm{st,cmp}}
        \colon X_{12}^{\mathrm{cmp}}\hookrightarrow P^{\mathrm{cmp}},\\
    p_1^{\mathrm{cmp}}
      &:=\operatorname{pr}_{X_1^{\mathrm{cmp}}}\circ
          \iota_{12}^{\mathrm{cmp}},\qquad
    p_2^{\mathrm{cmp}}
      :=\operatorname{pr}_{X_2}\circ\iota_{12}^{\mathrm{cmp}}.
  \end{aligned}                                                     \tag{C3}
  \]
  Let
  $V_1^{\mathrm{cmp}}:=(b_1^{\mathrm{cmp}})^{-1}
    (V_1^{\mathrm{cmp,in}})$, and denote the preserved-open isomorphism by
  \[
    \beta_1^{\mathrm{cmp}}\colon V_1^{\mathrm{cmp}}
      \xrightarrow{\sim}V_1^{\mathrm{cmp,in}}.                    \tag{C4}
  \]
  Define the pullback open immersion
  $\ell_1^{\mathrm{cmp}}\colon V_1^{\mathrm{cmp}}
    \hookrightarrow X_1^{\mathrm{cmp}}$ by
  \[
    b_1^{\mathrm{cmp}}\circ\ell_1^{\mathrm{cmp}}
      =\ell_1^{\mathrm{cmp,in}}\circ\beta_1^{\mathrm{cmp}}.        \tag{C4a}
  \]
  Define
  \[
    q_1^{\mathrm{cmp}}=q_1^{\mathrm{cmp,in}}\circ b_1^{\mathrm{cmp}},
  \]
  \[
    \psi_1^{\mathrm{cmp}}
      =\psi_1^{\mathrm{cmp,in}}\circ\beta_1^{\mathrm{cmp}},\qquad
    \widetilde j_1^{\mathrm{cmp}}
      =(\beta_1^{\mathrm{cmp}})^{-1}
        \circ\widetilde j_1^{\mathrm{cmp,in}},
    \qquad
    j_1^{\mathrm{cmp}}
      =\ell_1^{\mathrm{cmp}}\circ\widetilde j_1^{\mathrm{cmp}}.
  \]
  Substitution of the preserved-open equality (C4a) is the chain
  \[
  \begin{aligned}
    q_1^{\mathrm{cmp}}\circ j_1^{\mathrm{cmp}}
      &=q_1^{\mathrm{cmp,in}}\circ b_1^{\mathrm{cmp}}
          \circ\ell_1^{\mathrm{cmp}}
          \circ(\beta_1^{\mathrm{cmp}})^{-1}
          \circ\widetilde j_1^{\mathrm{cmp,in}}\\
      &=q_1^{\mathrm{cmp,in}}\circ\ell_1^{\mathrm{cmp,in}}
          \circ\widetilde j_1^{\mathrm{cmp,in}}
       =q_1^{\mathrm{cmp,in}}\circ j_1^{\mathrm{cmp,in}}
       =s_1.
  \end{aligned}
  \]
  Define
  \[
    V_{12}^{\mathrm{cmp}}
      =X_{12}^{\mathrm{cmp}}\times_{X_1^{\mathrm{cmp}}}V_1^{\mathrm{cmp}}.
  \]
  The restriction formula of
  \Cref{lemma:nagata-comparison-restrict-proper-open} and
  \Cref{lemma:nagata-scheme-image-strict-transform} define an isomorphism
  \[
    \beta_{12}^{\mathrm{cmp}}\colon
      V_{12}^{\mathrm{cmp}}\xrightarrow{\sim}V_{12},\qquad
    d_{12}^{\mathrm{cmp,in}}\circ\beta_{12}^{\mathrm{cmp}}
      =b_{12}^{\mathrm{cmp}}\circ
        \operatorname{pr}_{X_{12}^{\mathrm{cmp}}}.                \tag{C5}
  \]
  Define
  \[
    d_{12}^{\mathrm{cmp}}
      =\bigl(V_{12}\xrightarrow{(\beta_{12}^{\mathrm{cmp}})^{-1}}
        V_{12}^{\mathrm{cmp}}\hookrightarrow X_{12}^{\mathrm{cmp}}\bigr),
  \]
  The second equality in (C3) is the closed immersion
  \[
    \iota_{12}^{\mathrm{cmp}}\colon X_{12}^{\mathrm{cmp}}
      \hookrightarrow X_1^{\mathrm{cmp}}\times_SX_2.              \tag{C6}
  \]
  For $k=1,2$, define the pulled-back boundary and its strict-transform
  closure by
  \[
  \begin{aligned}
    Z_{1,k}^{\mathrm{cmp}}
      &:=V_1^{\mathrm{cmp}}\times_{V_1^{\mathrm{cmp,in}}}
          Z_{1,k}^{\mathrm{cmp,in}},\\
    \overline Z_{1,k}^{\mathrm{cmp}}
      &:=\operatorname{ST}_{b_1^{\mathrm{cmp}}}
          (\overline Z_{1,k}^{\mathrm{cmp,in}}),\\
    \overline z_{1,k}^{\mathrm{cmp}}
      &:=\operatorname{pr}_{X_1^{\mathrm{cmp}}}\circ
          \iota_{\overline z_{1,k}^{\mathrm{cmp,in}},b_1^{\mathrm{cmp}}}
        \colon\overline Z_{1,k}^{\mathrm{cmp}}\hookrightarrow
          X_1^{\mathrm{cmp}}.
  \end{aligned}                                                     \tag{C7}
  \]
  By (CMP0) and \Cref{def:nagata-two-open-prepared-boundaries},
  \[
    \overline z_{1,k}^{\mathrm{cmp,in}}
      =\overline z_{1,k}
      =i_{X_1^{\mathrm{cmp,in}},
          \overline{|Z_{1,k}^{\mathrm{cmp,in}}|}^{\,|X_1^{\mathrm{cmp,in}}|}}.
  \]
  Define
  \[
    Z_{12,2}^{\mathrm{cmp}}
      :=V_{12}^{\mathrm{cmp}}\times_{V_{12}}
          Z_{12,2}^{\mathrm{cmp,in}}.                              \tag{C8}
  \]
  Make the explicit substitution
  \[
    (X_1,q_1,V_1,\ell_1,\psi_1,\widetilde j_1,j_1,
      Z_{1,k},\overline Z_{1,k},X_{12},d_{12},\iota_{12},p_1,p_2,
      Z_{12,2})
      :=
    (X_1^{\mathrm{cmp}},q_1^{\mathrm{cmp}},V_1^{\mathrm{cmp}},
      \ell_1^{\mathrm{cmp}},\psi_1^{\mathrm{cmp}},
      \widetilde j_1^{\mathrm{cmp}},j_1^{\mathrm{cmp}},
      Z_{1,k}^{\mathrm{cmp}},\overline Z_{1,k}^{\mathrm{cmp}},
      X_{12}^{\mathrm{cmp}},d_{12}^{\mathrm{cmp}},
      \iota_{12}^{\mathrm{cmp}},p_1^{\mathrm{cmp}},p_2^{\mathrm{cmp}},
      Z_{12,2}^{\mathrm{cmp}}).                                   \tag{C9}
  \]
\end{definition}

\begin{lemma}[Properties of the improved first comparison projection]
  \label{lemma:nagata-comparison-isIso-near-boundary}
  \uses{def:nagata-two-open-improved-comparison}
  \lean{}

  The morphism
  $b_1^{\mathrm{cmp}}\colon X_1^{\mathrm{cmp}}\to X_1^{\mathrm{cmp,in}}$
  is $V_1^{\mathrm{cmp,in}}$-admissible.  There is an open immersion
  $\omega_1\colon\Omega_1\hookrightarrow X_1$ with
  \[
    \overline Z_{1,2}\subseteq\Omega_1,\qquad
    p_1^{-1}(\Omega_1)\xrightarrow[\ p_1\ ]{\sim}\Omega_1.
  \]
  The strict transforms satisfy
  \[
    \overline Z_{1,1}\times_{X_1}\overline Z_{1,2}=\varnothing.
  \]
\end{lemma}

\begin{proof}
  \uses{lemma:nagata-boundary-loci,
    lemma:nagata-comparison-projections-proper,
    lemma:nagata-comparison-restrict-proper-open,
    lemma:nagata-isomorphism-near-closure,
    lemma:nagata-scheme-image-strict-transform,
    lemma:nagata-boundary-disjoint-persists}
  By \Cref{lemma:nagata-two-open-graph-extensions},
  $U_2\in\mathscr N_2$.  By \Cref{lemma:nagata-boundary-loci},
  \[
    Z_2\subseteq U_2\subseteq N_2.                                \tag{17}
  \]
  Set
  \[
    O_1^{\mathrm{cmp,in}}
      :=(\psi_1^{\mathrm{cmp,in}})^{-1}(N_2)
      \subseteq V_1^{\mathrm{cmp,in}}.
  \]
  By
  \Cref{lemma:nagata-comparison-restrict-proper-open}, the restriction of
  $p_1^{\mathrm{cmp,in}}$ over $O_1^{\mathrm{cmp,in}}$ is the base change of
  $\psi_2^{-1}(N_2)\xrightarrow{\sim}N_2$; hence
  \[
    (p_1^{\mathrm{cmp,in}})^{-1}(O_1^{\mathrm{cmp,in}})
      \xrightarrow{\sim}O_1^{\mathrm{cmp,in}}.                    \tag{18}
  \]
  The equality
  $Z_{1,2}^{\mathrm{cmp,in}}=V_1^{\mathrm{cmp,in}}\times_UZ_2$
  and (17) are the inclusion
  \[
    Z_{1,2}^{\mathrm{cmp,in}}\subseteq O_1^{\mathrm{cmp,in}}.     \tag{19}
  \]
  The map $p_1^{\mathrm{cmp,in}}$ is proper by
  \Cref{lemma:nagata-comparison-projections-proper}.  Equations (18)--(19)
  verify the hypotheses of
  \Cref{lemma:nagata-isomorphism-near-closure}; under this substitution its
  conclusion is the $V_1^{\mathrm{cmp,in}}$-admissibility of
  $b_1^{\mathrm{cmp}}$ and an open immersion
  $\omega_1\colon\Omega_1\hookrightarrow X_1^{\mathrm{cmp}}$ for which
  \[
    \overline Z_{1,2}^{\mathrm{cmp}}\subseteq\Omega_1,
    \qquad
    (p_1^{\mathrm{cmp}})^{-1}(\Omega_1)
      \xrightarrow[ p_1^{\mathrm{cmp}}\ ]{\sim}\Omega_1.        \tag{19a}
  \]
  By \Cref{lemma:nagata-scheme-image-strict-transform}, the strict transform
  in that lemma is $p_1^{\mathrm{cmp}}$, which is denoted $p_1$ after the
  substitution (C9).  Finally,
  \Cref{lemma:nagata-prepared-boundaries} and
  \Cref{lemma:nagata-boundary-disjoint-persists} are the equality
  \[
    \overline Z_{1,1}^{\mathrm{cmp}}
      \times_{X_1^{\mathrm{cmp}}}
    \overline Z_{1,2}^{\mathrm{cmp}}=\varnothing.
  \]
  Substitution (C9) is the displayed equality in the statement.
\end{proof}

\begin{lemma}[Density over the prepared boundary]
  \label{lemma:nagata-comparison-boundary-density}
  \lean{}

  The image of $Z_{12,2}$ under $d_{12}\colon V_{12}\to X_{12}$ is dense in
  \[
    Y_{12,2}=p_1^{-1}(\overline Z_{1,2})\subseteq X_{12}.
  \]
\end{lemma}

\begin{proof}
  \uses{lemma:nagata-comparison-isIso-near-boundary}
  Restriction of the isomorphism over $\Omega_1$ is
  \[
    p_1|_{Y_{12,2}}\colon
      Y_{12,2}\xrightarrow{\sim}\overline Z_{1,2}.                 \tag{20}
  \]
  By \Cref{def:nagata-two-open-comparison-boundary},
  \[
    |p_1|\bigl(|d_{12}|(|Z_{12,2}|)\bigr)=|Z_{1,2}|.              \tag{21}
  \]
  The definition of closure is
  \[
    \overline{|Z_{1,2}|}^{\,|\overline Z_{1,2}|}
      =|\overline Z_{1,2}|.                                      \tag{22}
  \]
  Transport of (21)--(22) through the homeomorphism (20) is
  \[
    \overline{|d_{12}|(|Z_{12,2}|)}^{\,|Y_{12,2}|}=|Y_{12,2}|.
  \]
\end{proof}

\begin{lemma}[Key boundary disjointness]
  \label{lemma:nagata-key-boundary-disjointness}
  \lean{}

  Inside $X_1\times_SX_2$,
  \[
    X_{12}\times_{X_1\times_SX_2}
      (\overline Z_{1,2}\times_S\overline Z_{2,1})
    =\varnothing.
  \]
\end{lemma}

\begin{proof}
  \uses{lemma:nagata-prepared-boundaries,
    lemma:nagata-comparison-boundary-density}
  Put $Y=p_1^{-1}(\overline Z_{1,2})$.  By
  \Cref{lemma:nagata-comparison-boundary-density},
  \[
    \overline{|d_{12}|(|Z_{12,2}|)}^{\,|Y|}=|Y|.                  \tag{23}
  \]
  By \Cref{def:nagata-two-open-comparison-boundary},
  \[
    |p_2|\bigl(|d_{12}|(|Z_{12,2}|)\bigr)\subseteq|Z_{2,2}|.      \tag{24}
  \]
  Since $|\overline Z_{2,2}|$ is closed, continuity is the inclusion chain
  \[
    p_2(|Y|)
      \subseteq
      \overline{|p_2|(|d_{12}|(|Z_{12,2}|))}^{\,|X_2|}
      \subseteq|\overline Z_{2,2}|.                               \tag{25}
  \]
  By \Cref{lemma:nagata-prepared-boundaries},
  \[
    \overline Z_{2,2}\times_{X_2}\overline Z_{2,1}=\varnothing.   \tag{26}
  \]
  Therefore
  \[
  \begin{aligned}
    X_{12}\times_{X_1\times_SX_2}
      (\overline Z_{1,2}\times_S\overline Z_{2,1})
      &=
        p_1^{-1}(\overline Z_{1,2})
          \times_{X_{12}}
        p_2^{-1}(\overline Z_{2,1})\\
      &\subseteq
        p_2^{-1}(
          \overline Z_{2,2}\times_{X_2}\overline Z_{2,1})\\
      &=\varnothing.
  \end{aligned}
  \]
\end{proof}

\begin{definition}[Intermediate two-open schemes]
  \label{def:nagata-two-open-intermediate-schemes}
  \uses{def:nagata-two-open-improved-comparison,
    def:nagata-two-open-prepared-boundaries,
    def:nagata-choice-notation}
  \lean{}

  For $\{i,j\}=\{1,2\}$, define
  \[
    X_i^\circ=X_i\setminus|\overline Z_{i,j}|,\qquad
    k_i\colon X_i^\circ\hookrightarrow X_i.
  \]
  By \Cref{def:nagata-two-open-boundary-loci} and
  \Cref{lemma:nagata-two-open-graph-extensions},
  \[
  \begin{aligned}
    |\widetilde j_i(U_i)|\cap|Z_{i,j}|
      &=|\widetilde j_i(U_i)|
          \cap|\psi_i^{-1}(U\setminus U_i)|\\
      &=|\widetilde j_i(U_i)|
          \cap|\psi_i^{-1}(U_i)|
          \cap|\psi_i^{-1}(U\setminus U_i)|\\
      &=\varnothing.
  \end{aligned}                                                   \tag{26a}
  \]
  Since $|Z_{i,j}|$ is closed in the open subspace
  $|\ell_i|(|V_i|)\subseteq|X_i|$, the definition of
  $\overline Z_{i,j}$ in
  \Cref{def:nagata-two-open-prepared-boundaries} is the equality
  \[
    |\overline Z_{i,j}|\cap|\ell_i|(|V_i|)
      =|\ell_i|(|Z_{i,j}|).                                      \tag{26b}
  \]
  Substitution of (26a)--(26b) and
  $j_i=\ell_i\circ\widetilde j_i$ is the chain
  \[
  \begin{aligned}
    |j_i|(|U_i|)\cap|\overline Z_{i,j}|
      &=\ell_i(|\widetilde j_i(U_i)|)
          \cap|\overline Z_{i,j}|\\
      &=\ell_i\bigl(
          |\widetilde j_i(U_i)|\cap|Z_{i,j}|
        \bigr)
       =\varnothing.
  \end{aligned}                                                   \tag{26c}
  \]
  Define $j_i^\circ\colon U_i\hookrightarrow X_i^\circ$ by
  \[
    k_i\circ j_i^\circ=j_i.
  \]
  For every scheme $T$, define the overlap equalizer
  \[
    \mathscr E_{i,T}:=
      \left\{(h_U,h_X)\in
        \operatorname{Hom}(U,T)\times\operatorname{Hom}(X_i^\circ,T)
        \ \middle|\ 
        h_U\circ u_i=h_X\circ j_i^\circ
      \right\}.                                                    \tag{26d}
  \]
  Define the set of two-chart gluing realizations
  \[
  \begin{aligned}
    \mathscr G_i:=\bigl\{(G,a,c)\ \bigm|\ {}&
      G\text{ is a scheme},\quad
      a\colon U\to G,\quad c\colon X_i^\circ\to G,\\
    &\operatorname{IsOpenImmersion}(a)\wedge
      \operatorname{IsOpenImmersion}(c),\\
    &|G|=|a|(|U|)\cup |c|(|X_i^\circ|),\qquad
      a\circ u_i=c\circ j_i^\circ,\\
    &\operatorname{IsIso}\!\left(
       (u_i,j_i^\circ)\colon
       U_i\longrightarrow U\times_GX_i^\circ\right),\\
    &\forall T,\quad
      \left[
      \begin{array}{rcl}
        \operatorname{Hom}(G,T)&\longrightarrow&\mathscr E_{i,T}\\
        h&\longmapsto&(h\circ a,h\circ c)
      \end{array}
      \right]\text{ is bijective}
      \bigr\}.
  \end{aligned}                                                     \tag{26e}
  \]
  The scheme-gluing theorem
  \cite[\href{https://stacks.math.columbia.edu/tag/01JA}{Tag 01JA}]
  {stacks-project}, applied to the two open immersions $u_i$ and $j_i^\circ$
  with transition $\operatorname{id}_{U_i}$, is the assertion
  \[
    \mathscr G_i\neq\varnothing.                                  \tag{26f}
  \]
  By \Cref{def:nagata-choice-notation}, define
  \[
    (W_i,a_i,c_i):=\operatorname{ch}(\mathscr G_i)\in\mathscr G_i.
  \]
  Unpacking this membership defines the two open immersions and the Hom
  bijection
  \[
    a_i\colon U\hookrightarrow W_i,\qquad
    c_i\colon X_i^\circ\hookrightarrow W_i,\qquad
    a_i\circ u_i=c_i\circ j_i^\circ.                              \tag{27}
  \]
  \[
  \begin{gathered}
    \Gamma_{i,T}\colon\operatorname{Hom}(W_i,T)
      \xrightarrow{\sim}\mathscr E_{i,T},\qquad
    \Gamma_{i,T}(h):=(h\circ a_i,h\circ c_i).
  \end{gathered}                                                    \tag{27a}
  \]
  The structural equality $q_i\circ j_i=s\circ u_i$ and
  $k_i\circ j_i^\circ=j_i$ form the chain
  \[
    (q_i\circ k_i)\circ j_i^\circ
      =q_i\circ j_i
      =s\circ u_i,
  \]
  so $(s,q_i\circ k_i)\in\mathscr E_{i,S}$.  Define
  \[
    w_i:=\Gamma_{i,S}^{-1}(s,q_i\circ k_i)\colon W_i\to S.
  \]
  Applying $\Gamma_{i,S}$ to this definition is the pair of equations
  \[
    w_i\circ a_i=s,\qquad
    w_i\circ c_i=q_i\circ k_i.                                    \tag{28}
  \]
  Define the mixed-overlap immersion
  \[
    \kappa_i=(u_i,j_i^\circ)\colon U_i\to U\times_SX_i^\circ.     \tag{29}
  \]
\end{definition}

\begin{lemma}[Closed image of the proper extension graph]
  \label{lemma:nagata-proper-extension-graph-closed}
  \lean{}

  Define
  \[
    \gamma_i=(\psi_i,\ell_i)\colon V_i\to U\times_SX_i.
  \]
  Then $\gamma_i$ is proper and $|\gamma_i|(|V_i|)$ is closed.
\end{lemma}

\begin{proof}
  \uses{lemma:nagata-proper-to-separated-target}
  Let $\operatorname{pr}_U\colon U\times_SX_i\to U$.  Then
  \[
    \operatorname{pr}_U\circ\gamma_i=\psi_i.                       \tag{30}
  \]
  The map $\psi_i$ is proper by
  \Cref{lemma:nagata-two-open-graph-extensions}.  The projection
  $\operatorname{pr}_U$ is the base change of the proper, hence separated,
  morphism $q_i\colon X_i\to S$.  Applying
  \Cref{lemma:nagata-proper-to-separated-target} to
  $(\gamma_i,\operatorname{pr}_U,\psi_i)$ and (30) is
  \[
    \operatorname{Proper}(\gamma_i).
  \]
  The implication
  $\operatorname{Proper}(\gamma_i)\Rightarrow
    |\gamma_i|(|V_i|)\text{ is closed}$
  proves the second assertion.
\end{proof}

\begin{lemma}[Constructibility of the overlap image]
  \label{lemma:nagata-intermediate-overlap-constructible}
  \lean{}

  The subset $|\kappa_i|(|U_i|)\subseteq|U\times_SX_i^\circ|$ is constructible.
\end{lemma}

\begin{proof}
  The schemes $U_i$ and $U\times_SX_i^\circ$ are of finite type over the
  Noetherian scheme $S$, hence are Noetherian.  The target is separated over
  $S$.  By \Cref{lemma:nagata-finite-type-to-separated-target},
  $\kappa_i$ is of finite type, hence of finite presentation over its
  Noetherian target.  Chevalley's theorem
  (\cite[\href{https://stacks.math.columbia.edu/tag/054K}{Tag 054K}]
  {stacks-project}) asserts
  \[
    |\kappa_i|(|U_i|)\subseteq|U\times_SX_i^\circ|
    \quad\text{constructible}.
  \]
\end{proof}

\begin{lemma}[A specialization leaving $U_i$ meets the forbidden boundary]
  \label{lemma:nagata-intermediate-specialization-boundary}
  \lean{}

  Let $u_0\in|U_i|$ and $(u,x)\in|U\times_SX_i|$ satisfy
  \[
    (u_i(u_0),j_i(u_0))\leadsto(u,x).
  \]
  If $u\notin|U_i|$, then $x\in|\overline Z_{i,j}|$.
\end{lemma}

\begin{proof}
  \uses{lemma:nagata-proper-extension-graph-closed,lemma:nagata-boundary-loci}
  By \Cref{lemma:nagata-two-open-graph-extensions},
  \[
  \begin{aligned}
    \gamma_i(\widetilde j_i(u_0))
      &=(\psi_i(\widetilde j_i(u_0)),
          \ell_i(\widetilde j_i(u_0)))\\
      &=(u_i(u_0),j_i(u_0)).
  \end{aligned}
  \]
  Since $|\gamma_i|(|V_i|)$ is closed, there exists $v\in|V_i|$ satisfying
  \[
    \gamma_i(v)=(u,x),\qquad
    \psi_i(v)=u,\qquad \ell_i(v)=x.                               \tag{31}
  \]
  Since $U=U_i\cup U_j$,
  \[
    u\notin|U_i|\Longrightarrow u\in|U|\setminus|U_i|=|Z_j|.
  \]
  Thus
  \[
    v\in|\psi_i^{-1}(Z_j)|=|Z_{i,j}|,
    \qquad
    x=\ell_i(v)\in|Z_{i,j}|\subseteq|\overline Z_{i,j}|.
  \]
\end{proof}

\begin{lemma}[Specialization stability inside the boundary complement]
  \label{lemma:nagata-intermediate-overlap-specialization-closed}
  \lean{}

  The subset
  $|\kappa_i|(|U_i|)\subseteq|U\times_SX_i^\circ|$ is stable under
  specialization.
\end{lemma}

\begin{proof}
  \uses{lemma:nagata-intermediate-specialization-boundary}
  Let $|\kappa_i|(u_0)\leadsto(u,x)$ in $|U\times_SX_i^\circ|$.  If
  $u\notin|U_i|$, then
  \Cref{lemma:nagata-intermediate-specialization-boundary} is the membership
  \[
    x\in|\overline Z_{i,j}|\cap|X_i^\circ|=\varnothing.
  \]
  Therefore $u\in|U_i|$; define
  \[
    u^\sharp=|u_i|^{-1}(u)\in|U_i|.                               \tag{31a}
  \]
  Define the open immersion
  \[
    \eta_i=\operatorname{id}_U\times k_i\colon
      U\times_SX_i^\circ\hookrightarrow U\times_SX_i.
  \]
  The definitions in
  \Cref{def:nagata-two-open-intermediate-schemes} and
  \Cref{lemma:nagata-proper-extension-graph-closed} are the equality chain
  \[
  \begin{aligned}
    \eta_i\circ\kappa_i
      &=(u_i,k_i\circ j_i^\circ)
       =(u_i,j_i)\\
      &=(\psi_i\circ\widetilde j_i,
          \ell_i\circ\widetilde j_i)
       =\gamma_i\circ\widetilde j_i.
  \end{aligned}
  \]
  Functoriality of specialization under $|\eta_i|$ is
  \[
    |\gamma_i|(|\widetilde j_i|(u_0))
      \leadsto|\eta_i|(u,x)
      =(u,k_i(x))
      \quad\text{in }|U\times_SX_i|.
  \]
  By \Cref{lemma:nagata-proper-extension-graph-closed},
  $|\gamma_i|(|V_i|)$ is closed.  Hence there is $v\in|V_i|$ satisfying
  \[
    |\gamma_i|(v)=(u,|k_i|(x)),\qquad
    |\psi_i|(v)=u,\qquad
    |\ell_i|(v)=|k_i|(x).                                         \tag{31b}
  \]
  By (31a)--(31b) and
  \Cref{lemma:nagata-two-open-graph-extensions},
  \[
  \begin{aligned}
    v
      &=|\widetilde j_i|(u^\sharp),\\
    |k_i|(x)
      &=|\ell_i|(v)
       =|\ell_i|(|\widetilde j_i|(u^\sharp))
       =|j_i|(u^\sharp)
       =|k_i|(|j_i^\circ|(u^\sharp)),\\
    x&=|j_i^\circ|(u^\sharp),\\
    (u,x)
      &=(|u_i|(u^\sharp),|j_i^\circ|(u^\sharp))
       =|\kappa_i|(u^\sharp).
  \end{aligned}
  \]
\end{proof}

\begin{lemma}[Closedness of the intermediate overlap]
  \label{lemma:nagata-intermediate-overlap-closed}
  \lean{}

  The morphism
  \[
    \kappa_i=(u_i,j_i^\circ)\colon U_i\to U\times_SX_i^\circ
  \]
  is a closed immersion.
\end{lemma}

\begin{proof}
  \uses{lemma:nagata-intermediate-overlap-constructible,
    lemma:nagata-intermediate-overlap-specialization-closed}
  By \Cref{lemma:nagata-intermediate-overlap-constructible},
  $|\kappa_i|(|U_i|)$ is constructible.  By
  \Cref{lemma:nagata-intermediate-overlap-specialization-closed}, it is stable
  under specialization.  Since the target is Noetherian,
  \[
    |\kappa_i|(|U_i|)\text{ is closed}.
  \]
  The map $\kappa_i$ is the mixed-overlap immersion of the two open
  immersions in (27).  An immersion with closed image is a closed immersion,
  so $\kappa_i$ is a closed immersion.
\end{proof}

\begin{lemma}[Separatedness of the intermediate schemes]
  \label{lemma:nagata-intermediate-separated}
  \lean{}

  Each $w_i\colon W_i\to S$ is separated and of finite type.
\end{lemma}

\begin{proof}
  \uses{lemma:nagata-intermediate-overlap-closed,
    lemma:nagata-separated-two-open-cover}
  The open cover is
  \[
    W_i=a_i(U)\cup c_i(X_i^\circ),
  \]
  and (28) is the pair of restriction equalities
  \[
    w_i|_{a_i(U)}=s,\qquad
    w_i|_{c_i(X_i^\circ)}=q_i|_{X_i^\circ}.
  \]
  Both restrictions are separated and of finite type.  Their mixed-overlap
  map is $\kappa_i$, which is closed by
  \Cref{lemma:nagata-intermediate-overlap-closed}.  Hence
  \Cref{lemma:nagata-separated-two-open-cover} is the predicate
  $\operatorname{Separated}(w_i)$.  Since finite type is local on the source,
  the two displayed restrictions are the predicate
  $\operatorname{FiniteType}(w_i)$.
\end{proof}

\begin{lemma}[The pair of valuation extensions factors through $X_{12}$]
  \label{lemma:nagata-valuation-pair-factors-comparison}
  \lean{}

  Let $A$ be a valuation ring, $K=\operatorname{Frac}(A)$, and
  $\nu_A\colon\operatorname{Spec}(K)\to\operatorname{Spec}(A)$.  Let
  \[
    \gamma\colon\operatorname{Spec}(K)\to U_{12},\qquad
    s_A\colon\operatorname{Spec}(A)\to S
  \]
  satisfy $s\circ e\circ\gamma=s_A\circ\nu_A$.  Suppose
  $g_i\colon\operatorname{Spec}(A)\to X_i$ satisfy
  \[
    q_i\circ g_i=s_A,\qquad
    g_i\circ\nu_A=j_i\circ e_i\circ\gamma.
  \]
  Define
  \[
    g=(g_1,g_2)\colon\operatorname{Spec}(A)\to X_1\times_SX_2.
  \]
  Then there is a unique
  $\widetilde g\colon\operatorname{Spec}(A)\to X_{12}$ satisfying
  $\iota_{12}\circ\widetilde g=g$.
\end{lemma}

\begin{proof}
  \uses{lemma:nagata-valuation-factor-closed}
  The fibre-product universal property applied to
  $\widetilde j_i\circ e_i\circ\gamma$ defines
  \[
    \gamma_V\colon\operatorname{Spec}(K)\to V_1\times_UV_2=V_{12}.
  \]
  Its defining equations are
  \[
    r_i\circ\gamma_V=\widetilde j_i\circ e_i\circ\gamma
      \qquad(i=1,2).                                             \tag{31c}
  \]
  By the definitions of $v_{12}$ and $d_{12}$,
  \[
  \begin{aligned}
    g\circ\nu_A
      &=(j_1\circ e_1\circ\gamma,j_2\circ e_2\circ\gamma)\\
      &=(\ell_1\circ\widetilde j_1\circ e_1\circ\gamma,
          \ell_2\circ\widetilde j_2\circ e_2\circ\gamma)\\
      &=v_{12}\circ\gamma_V\\
      &=\iota_{12}\circ d_{12}\circ\gamma_V.
  \end{aligned}                                                    \tag{32}
  \]
  Apply \Cref{lemma:nagata-valuation-factor-closed} to
  $\iota_{12}$, $g$, and the generic factor in (32).  This defines the unique
  $\widetilde g$.
\end{proof}

\begin{lemma}[One valuation extension avoids the forbidden boundary]
  \label{lemma:nagata-one-valuation-avoids-boundary}
  \uses{def:nagata-two-open-comparison-space,
    def:nagata-two-open-intermediate-schemes}
  \lean{}

  Under the hypotheses of
  \Cref{lemma:nagata-valuation-pair-factors-comparison}, there is an
  $i\in\{1,2\}$, with $\{i,j\}=\{1,2\}$, and a unique
  \[
    g_i^\circ\colon\operatorname{Spec}(A)\to X_i^\circ
  \]
  satisfying $k_i\circ g_i^\circ=g_i$.  The morphism
  \[
    h_i=c_i\circ g_i^\circ\colon\operatorname{Spec}(A)\to W_i
  \]
  satisfies
  \[
    h_i\circ\nu_A=a_i\circ e\circ\gamma,\qquad
    w_i\circ h_i=s_A.
  \]
\end{lemma}

\begin{proof}
  \uses{lemma:nagata-key-boundary-disjointness,
    lemma:nagata-valuation-pair-factors-comparison}
  Let $a$ be the closed point of $\operatorname{Spec}(A)$.  Suppose no $g_i$
  factors through $X_i^\circ$.  The open $g_i^{-1}(X_i^\circ)$ contains the
  generic point.  An open of a valuation spectrum is the whole spectrum if
  and only if it contains the closed point; therefore
  \[
    g_i(a)\in|\overline Z_{i,j}|\quad(i=1,2).                     \tag{33}
  \]
  By \Cref{lemma:nagata-valuation-pair-factors-comparison},
  $g=\iota_{12}\circ\widetilde g$.  Thus (33) puts $\widetilde g(a)$ in
  \[
    \left|
      X_{12}\times_{X_1\times_SX_2}
        (\overline Z_{1,2}\times_S\overline Z_{2,1})
    \right|,
  \]
  contradicting \Cref{lemma:nagata-key-boundary-disjointness}.  Hence some
  $g_i$ factors uniquely through the open immersion $k_i$.

  Substitution of (27)--(29) is the chain
  \[
  \begin{aligned}
    h_i\circ\nu_A
      &=c_i\circ j_i^\circ\circ e_i\circ\gamma\\
      &=a_i\circ u_i\circ e_i\circ\gamma
       =a_i\circ e\circ\gamma,\\
    w_i\circ h_i
      &=q_i\circ k_i\circ g_i^\circ
       =q_i\circ g_i
       =s_A.
  \end{aligned}
  \]
\end{proof}

\begin{definition}[Final common modification for the two-open construction]
  \label{def:nagata-two-open-final-modification}
  \uses{def:nagata-two-open-intermediate-schemes,
    def:nagata-common-modification}
  \lean{}

  Apply \Cref{def:nagata-common-modification} to
  $a_i\colon U\hookrightarrow W_i$.  Denote the output by
  \[
    b_i'\colon W_i'\to W_i,\qquad
    o_i\colon W_i'\hookrightarrow W=\mathsf{CM}(W_1,W_2;U).
  \]
  Let $a_i'\colon U\hookrightarrow W_i'$ be the preserved-open maps and
  record their defining equations
  \[
    b_i'\circ a_i'=a_i\qquad(i=1,2).                             \tag{33a}
  \]
  Then define
  \[
    j_W=o_1\circ a_1'=o_2\circ a_2'\colon U\to W.                 \tag{34}
  \]
  Define $w\colon W\to S$ by
  \[
    w\circ o_i=w_i\circ b_i'\quad(i=1,2).                         \tag{35}
  \]
\end{definition}

\begin{lemma}[Structural properties of the final common modification]
  \label{lemma:nagata-two-open-final-modification}
  \uses{def:nagata-two-open-final-modification}
  \lean{}

  Each $b_i'$ is $U$-admissible, each $o_i$ is an open immersion, and
  \[
    W=o_1(W_1')\cup o_2(W_2'),\qquad
    j_W=o_i\circ a_i'.
  \]
  The morphism $w\colon W\to S$ is separated and of finite type,
  $j_W\colon U\hookrightarrow W$ is an open immersion, and
  \[
    w\circ j_W=s.
  \]
\end{lemma}

\begin{proof}
  \uses{lemma:nagata-intermediate-separated,
    lemma:nagata-common-modification}
  By \Cref{lemma:nagata-intermediate-separated}, both $W_i\to S$ are
  separated and of finite type.  Substitution into
  \Cref{lemma:nagata-common-modification} is the predicate-and-equality tuple
  \[
  \begin{gathered}
    b_i'\text{ is }U\text{-admissible},\qquad
    \operatorname{IsOpenImmersion}(o_i),\\
    W=o_1(W_1')\cup o_2(W_2'),\qquad
    \operatorname{Separated}(w),\qquad
    \operatorname{FiniteType}(w).
  \end{gathered}
  \]
  Equation (34) is $j_W=o_i\circ a_i'$, a composite of open immersions.
  Using (28), (34), and (35),
  \[
  \begin{aligned}
    w\circ j_W
      &=w\circ o_i\circ a_i'
       =w_i\circ b_i'\circ a_i'
       =w_i\circ a_i
       =s.
  \end{aligned}
  \]
\end{proof}

\begin{definition}[Two-open compactification candidate]
  \label{def:nagata-two-open-compactification-candidate}
  \uses{def:nagata-two-open-final-modification}
  \lean{}

  For arbitrary
  \[
    D_i\in\operatorname{Comp}(s_i)\qquad(i=1,2),
  \]
  repeat Definitions
  \ref{def:nagata-two-open-dense-compactifications}--
  \ref{def:nagata-two-open-final-modification} after the explicit
  substitution
  \[
    (C_1^0,C_2^0):=(D_1,D_2).
  \]
  Denote the three output symbols of
  \Cref{def:nagata-two-open-final-modification} by
  \[
    W(D_1,D_2),\qquad
    j_W(D_1,D_2)\colon U\to W(D_1,D_2),\qquad
    w(D_1,D_2)\colon W(D_1,D_2)\to S.
  \]
  Define
  \[
    \mathsf C_{12}(U/S;D_1,D_2)
      =
      \bigl(
        U\xrightarrow{j_W(D_1,D_2)}W(D_1,D_2)
          \xrightarrow{w(D_1,D_2)}S,\,
        w(D_1,D_2)\circ j_W(D_1,D_2)=s
      \bigr).
  \]
  For the selected $C_i^0$ of
  \Cref{def:nagata-two-open-dense-compactifications}, abbreviate this by
  \[
    \mathsf C_{12}(U/S)
      =\mathsf C_{12}(U/S;C_1^0,C_2^0),
  \]
  and abbreviate $W(C_1^0,C_2^0)$, $j_W(C_1^0,C_2^0)$, and
  $w(C_1^0,C_2^0)$ by $W$, $j_W$, and $w$.
\end{definition}

\begin{lemma}[Density after the final common modification]
  \label{lemma:nagata-final-density}
  \uses{def:nagata-two-open-compactification-candidate}
  \lean{}

  The morphism
  \[
    d=j_W\circ e\colon U_{12}\to W
  \]
  is quasi-compact and satisfies
  $\overline{|d|(|U_{12}|)}=|W|$.
\end{lemma}

\begin{proof}
  \uses{def:nagata-closed-subscheme-class,
    def:nagata-morphism-predicate-notation,
    def:nagata-scheme-theoretic-image,
    lemma:nagata-intermediate-separated,
    lemma:nagata-two-open-dense-compactifications,
    lemma:nagata-two-open-graph-extensions,
    lemma:nagata-two-open-final-modification,
    lemma:nagata-dense-admissible-blowup}
  The standing hypothesis $\operatorname{Noetherian}(S)$, the finite-type
  predicates in \Cref{lemma:nagata-intermediate-separated,lemma:nagata-two-open-final-modification},
  Tag 01T6, and the open-immersion
  predicates (27) and \Cref{lemma:nagata-two-open-final-modification} are the
  following predicate chain:
  \[
  \begin{aligned}
    \operatorname{Noetherian}(S)\ \wedge\
      \operatorname{FiniteType}(w_i)
      &\Longrightarrow\operatorname{Noetherian}(W_i)
        &&\text{by \cite[Tag 01T6]{stacks-project}},\\
    \operatorname{Noetherian}(S)\ \wedge\
      \operatorname{FiniteType}(w)
      &\Longrightarrow\operatorname{Noetherian}(W)
        &&\text{by \cite[Tag 01T6]{stacks-project}},\\
    \operatorname{Noetherian}(W_i)\ \wedge\
      \operatorname{IsOpenImmersion}(a_i)
      &\Longrightarrow\operatorname{IsQuasiCompact}(a_i),\\
    \operatorname{Noetherian}(W)\ \wedge\
      \operatorname{IsOpenImmersion}(j_W)
      &\Longrightarrow\operatorname{IsQuasiCompact}(j_W),\\
    \operatorname{IsOpenImmersion}(a_i)
      &\Longrightarrow\operatorname{IsIso}(\Delta_{a_i})
       \Longrightarrow\operatorname{IsQuasiCompact}(\Delta_{a_i})
       \Longleftrightarrow\operatorname{QuasiSeparated}(a_i),\\
    \operatorname{IsOpenImmersion}(j_W)
      &\Longrightarrow\operatorname{IsIso}(\Delta_{j_W})
       \Longrightarrow\operatorname{IsQuasiCompact}(\Delta_{j_W})
       \Longleftrightarrow\operatorname{QuasiSeparated}(j_W),\\
    \operatorname{IsQuasiCompact}(a_i)\ \wedge\
      \operatorname{QuasiSeparated}(a_i)
      &\Longleftrightarrow\operatorname{IsQCQS}(a_i),\\
    \operatorname{IsQuasiCompact}(j_W)\ \wedge\
      \operatorname{QuasiSeparated}(j_W)
      &\Longleftrightarrow\operatorname{IsQCQS}(j_W).
  \end{aligned}                                                    \tag{FD0q}
  \]
  By \Cref{lemma:nagata-two-open-dense-compactifications},
  $U_i\hookrightarrow X_i^0$ is scheme-theoretically dense.  Each morphism
  \[
    X_1=X_1^{\mathrm{cmp}}
      \xrightarrow{b_1^{\mathrm{cmp}}}X_1^{\mathrm{prep}}
      \xrightarrow{b_1^{\mathrm{prep}}}X_1^{\mathrm{gr}}
      \xrightarrow{b_1^{\mathrm{gr}}}X_1^0,
  \]
  \[
    X_2=X_2^{\mathrm{prep}}
      \xrightarrow{b_2^{\mathrm{prep}}}X_2^{\mathrm{gr}}
      \xrightarrow{b_2^{\mathrm{gr}}}X_2^0
  \]
  is a sequence of admissible blow-ups whose defining ideals restrict to the
  unit ideal on $U_i$.  Repeated application of
  \Cref{lemma:nagata-dense-admissible-blowup} is, for $i=1,2$, the equality
  \[
    \ker\bigl(
      \mathcal O_{X_i}\longrightarrow(j_i)_*\mathcal O_{U_i}
    \bigr)=0.                                                     \tag{36a}
  \]
  Restriction of (36a) along $k_i\colon X_i^\circ\hookrightarrow X_i$ and
  the equality $k_i\circ j_i^\circ=j_i$ are
  \[
    \ker\bigl(
      \mathcal O_{X_i^\circ}
        \longrightarrow(j_i^\circ)_*\mathcal O_{U_i}
    \bigr)=0.                                                     \tag{36b}
  \]
  Define
  \[
    \mathcal K_i
      =\ker\bigl(
          \mathcal O_{W_i}\longrightarrow(a_i)_*\mathcal O_U
        \bigr).
  \]
  On the open cover $W_i=a_i(U)\cup c_i(X_i^\circ)$, the gluing equalities
  in \Cref{def:nagata-two-open-intermediate-schemes} identify the two
  restrictions as
  \[
    \mathcal K_i|_{a_i(U)}=0,
    \qquad
    c_i^*\mathcal K_i
      =\ker\bigl(
          \mathcal O_{X_i^\circ}
            \longrightarrow(j_i^\circ)_*\mathcal O_{U_i}
        \bigr)=0                                                  \tag{37}
  \]
  by (36b).  The sheaf condition on this open cover is
  \[
    \mathcal K_i=0,
    \qquad
    [\,\operatorname{SIm}_{W_i}(a_i)\xhookrightarrow{i_{a_i}}W_i\,]
      =[\,W_i\xhookrightarrow{\operatorname{id}_{W_i}}W_i\,]
      \quad\text{in }\operatorname{ClosedSub}(W_i).              \tag{38}
  \]
  By \Cref{lemma:nagata-two-open-final-modification}, $b_i'$ is an
  $a_i(U)$-admissible blow-up and
  $b_i'\circ a_i'=a_i$.  Applying
  \Cref{lemma:nagata-dense-admissible-blowup} to (38) is
  \[
    \ker\bigl(
      \mathcal O_{W_i'}\longrightarrow(a_i')_*\mathcal O_U
    \bigr)=0.                                                     \tag{39}
  \]
  Define
  \[
    \mathcal K
      =\ker\bigl(
          \mathcal O_W\longrightarrow(j_W)_*\mathcal O_U
        \bigr).
  \]
  The equalities $j_W=o_i\circ a_i'$ from (34), the open-cover equality
  $W=o_1(W_1')\cup o_2(W_2')$, and (39) are the restriction equalities
  \[
    o_i^*\mathcal K
      =\ker\bigl(
          \mathcal O_{W_i'}\longrightarrow(a_i')_*\mathcal O_U
        \bigr)=0
    \quad(i=1,2),
    \qquad
    \mathcal K=0.                                                 \tag{39a}
  \]
  By \Cref{def:nagata-scheme-theoretic-image}, (39a) is the following class
  equality and morphism equality:
  \[
    [\,\operatorname{SIm}_W(j_W)\xhookrightarrow{i_{j_W}}W\,]
      =[\,W\xhookrightarrow{\operatorname{id}_W}W\,]
      \quad\text{in }\operatorname{ClosedSub}(W),
    \qquad
    \overline{|j_W|(|U|)}^{\,|W|}=|W|.                            \tag{39b}
  \]
  The assumed equality
  $\overline{|e|(|U_{12}|)}^{\,|U|}=|U|$ is the inclusion
  \[
    |j_W|(|U|)
      \subseteq\overline{|j_W|(|e|(|U_{12}|))}^{\,|W|}
      =\overline{|d|(|U_{12}|)}^{\,|W|}.
  \]
  Taking closures and applying (39b),
  \[
    |W|
      \subseteq\overline{|d|(|U_{12}|)}^{\,|W|}
      \subseteq|W|,
  \]
  proving density.

  The scheme $U_{12}$ is open in the Noetherian scheme $U$, hence is
  Noetherian.  For every affine open $Q\subseteq W$, the inverse image
  $d^{-1}(Q)$ is open in $U_{12}$ and therefore quasi-compact.  Thus $d$ is
  quasi-compact.
\end{proof}

\begin{lemma}[Valuative lifting for the two-open candidate]
  \label{lemma:nagata-two-open-valuative-lifting}
  \uses{def:nagata-two-open-compactification-candidate}
  \lean{}

  Let $A$ be a valuation ring, $K=\operatorname{Frac}(A)$, and
  $\nu_A\colon\operatorname{Spec}(K)\to\operatorname{Spec}(A)$.  For every
  \[
    \gamma\colon\operatorname{Spec}(K)\to U_{12},\qquad
    s_A\colon\operatorname{Spec}(A)\to S
  \]
  satisfying $s\circ e\circ\gamma=s_A\circ\nu_A$, there is a unique
  $y_A\colon\operatorname{Spec}(A)\to W$ satisfying
  \[
    y_A\circ\nu_A=j_W\circ e\circ\gamma,\qquad
    w\circ y_A=s_A.
  \]
\end{lemma}

\begin{proof}
  \uses{lemma:nagata-one-valuation-avoids-boundary,
    lemma:nagata-blowup-proper,
    lemma:nagata-two-open-final-modification,
    def:nagata-morphism-predicate-notation}
  For $i=1,2$, put
  \[
    \gamma_i=j_i\circ e_i\circ\gamma.
  \]
  Then
  \[
    q_i\circ\gamma_i
      =s\circ e\circ\gamma
      =s_A\circ\nu_A.
  \]
  Properness of $q_i$ and the valuative criterion
  (\cite[\href{https://stacks.math.columbia.edu/tag/0BX5}{Tag 0BX5}]
  {stacks-project}) define
  \[
    g_i\colon\operatorname{Spec}(A)\to X_i,\qquad
    g_i\circ\nu_A=\gamma_i,\qquad q_i\circ g_i=s_A.                \tag{40}
  \]
  Apply \Cref{lemma:nagata-one-valuation-avoids-boundary} to (40).  Its
  conclusion is an index $i$ and a morphism
  \[
    h_i\colon\operatorname{Spec}(A)\to W_i,\qquad
    h_i\circ\nu_A=a_i\circ e\circ\gamma,\qquad
    w_i\circ h_i=s_A.                                              \tag{41}
  \]
  The map $b_i'\colon W_i'\to W_i$ is proper by
  \Cref{lemma:nagata-blowup-proper}, and
  \[
    b_i'\circ a_i'\circ e\circ\gamma
      =a_i\circ e\circ\gamma
      =h_i\circ\nu_A.
  \]
  Its valuative criterion defines
  \[
    h_i'\colon\operatorname{Spec}(A)\to W_i',\qquad
    b_i'\circ h_i'=h_i,\qquad
    h_i'\circ\nu_A=a_i'\circ e\circ\gamma.                         \tag{42}
  \]
  Define $y_A=o_i\circ h_i'$.  Then
  \[
  \begin{aligned}
    y_A\circ\nu_A
      &=o_i\circ a_i'\circ e\circ\gamma
       =j_W\circ e\circ\gamma,\\
    w\circ y_A
      &=w_i\circ b_i'\circ h_i'
       =w_i\circ h_i
       =s_A.
  \end{aligned}
  \]
  Let $y_A'$ be a second lift and define its scheme-theoretic equalizer by
  \[
  \begin{gathered}
    E_A:=\operatorname{Spec}(A)
      \times_{W\times_SW,\,(y_A,y_A'),\,\Delta_w}W,\\
    e_A:=\operatorname{pr}_{\operatorname{Spec}(A)}
      \colon E_A\to\operatorname{Spec}(A),\qquad
    r_A:=\operatorname{pr}_W\colon E_A\to W,\\
    (y_A,y_A')\circ e_A=\Delta_w\circ r_A.
  \end{gathered}                                                   \tag{42a}
  \]
  By \Cref{lemma:nagata-two-open-final-modification},
  \[
    \operatorname{IsSeparated}(w),
    \qquad
    \operatorname{IsSeparated}(w)
      \Longleftrightarrow\operatorname{IsClosedImmersion}(\Delta_w),
    \qquad
    \operatorname{IsClosedImmersion}(\Delta_w).
  \]
  Base change along $(y_A,y_A')$ is the predicate
  \[
    \operatorname{IsClosedImmersion}(e_A).
  \]
  The affine closed-immersion description therefore defines the unique ideal
  $J_A\triangleleft A$, the quotient map
  $\rho_A\colon A\twoheadrightarrow A/J_A$, and the unique isomorphism over
  $\operatorname{Spec}(A)$
  \[
  \begin{gathered}
    \epsilon_A\colon E_A\xrightarrow{\sim}\operatorname{Spec}(A/J_A),
    \qquad
    \operatorname{Spec}(\rho_A)\circ\epsilon_A=e_A,\\
    |e_A|(|E_A|)=V(J_A).                                        
  \end{gathered} \tag{42b}
  \]
  The two generic equalities for $y_A$ and $y_A'$ are
  \[
    (y_A,y_A')\circ\nu_A
      =\Delta_w\circ(j_W\circ e\circ\gamma).
  \]
  The fibre-product universal property in (42a) defines
  \[
    \widetilde\nu_A\colon\operatorname{Spec}(K)\to E_A,
    \qquad e_A\circ\widetilde\nu_A=\nu_A.                       \tag{42c}
  \]
  If $\eta_A=(0)\in\operatorname{Spec}(A)$ is the generic point, (42c) and
  closedness of $|e_A|(|E_A|)=V(J_A)$ are the inclusions
  \[
    \eta_A\in V(J_A),
    \qquad
    \operatorname{Spec}(A)=\overline{\{\eta_A\}}
      \subseteq V(J_A)\subseteq\operatorname{Spec}(A),
  \]
  hence
  \[
    V(J_A)=\operatorname{Spec}(A),
    \qquad J_A\subseteq\sqrt{(0)}.
  \]
  A valuation ring is a domain, so $\sqrt{(0)}=(0)$; therefore
  \[
  \begin{gathered}
    J_A=(0),\qquad
    \rho_A=\operatorname{id}_A,\qquad
    \operatorname{Spec}(\rho_A)=
      \operatorname{id}_{\operatorname{Spec}(A)},\\
    e_A=\operatorname{Spec}(\rho_A)\circ\epsilon_A=\epsilon_A,
    \qquad
    e_A^{-1}:=\epsilon_A^{-1}\colon\operatorname{Spec}(A)\to E_A,\\
    e_A^{-1}\circ e_A=\operatorname{id}_{E_A},\qquad
    e_A\circ e_A^{-1}=\operatorname{id}_{\operatorname{Spec}(A)},\qquad
    \operatorname{IsIso}(e_A).                                  
  \end{gathered} \tag{42d}
  \]
  Define
  \[
    \overline r_A:=r_A\circ e_A^{-1}
      \colon\operatorname{Spec}(A)\longrightarrow W.
  \]
  Compose the equality in (42a) on the right with $e_A^{-1}$ and substitute
  (42d).  The resulting equality is
  \[
  \begin{aligned}
    (y_A,y_A')
      &=(y_A,y_A')\circ e_A\circ e_A^{-1}\\
      &=\Delta_w\circ r_A\circ e_A^{-1}
       =\Delta_w\circ\overline r_A.
  \end{aligned}
  \]
  Applying the two projections
  $\operatorname{pr}_1,\operatorname{pr}_2\colon W\times_SW\to W$ is the
  pair of equalities
  \[
  \begin{aligned}
    y_A&=\operatorname{pr}_1\circ\Delta_w\circ\overline r_A
      =\overline r_A,\\
    y_A'&=\operatorname{pr}_2\circ\Delta_w\circ\overline r_A
      =\overline r_A,
  \end{aligned}
  \]
  and hence $y_A=y_A'$.
\end{proof}

\begin{lemma}[Properness of the two-open candidate]
  \label{lemma:nagata-two-open-candidate-proper}
  \uses{def:nagata-two-open-compactification-candidate}
  \lean{}

  The morphism $w\colon W\to S$ is proper.
\end{lemma}

\begin{proof}
  \uses{lemma:nagata-final-density,
    lemma:nagata-two-open-valuative-lifting,
    lemma:nagata-two-open-final-modification,
    lemma:nagata-dense-source-valuative}
  Put $d=j_W\circ e\colon U_{12}\to W$.  By
  \Cref{lemma:nagata-final-density},
  \[
    \operatorname{IsQuasiCompact}(d),\qquad
    \overline{|d|(|U_{12}|)}=|W|.
  \]
  The open-immersion predicates for $e\colon U_{12}\hookrightarrow U$ and
  $j_W\colon U\hookrightarrow W$, composition, and the diagonal definition
  in \Cref{def:nagata-morphism-predicate-notation} are the following
  predicate chain:
  \[
  \begin{aligned}
    \operatorname{IsOpenImmersion}(e)\wedge
      \operatorname{IsOpenImmersion}(j_W)
      &\Longrightarrow\operatorname{IsOpenImmersion}(d)
       \Longrightarrow\operatorname{QuasiSeparated}(d),\\
    \operatorname{IsQuasiCompact}(d)\wedge
      \operatorname{QuasiSeparated}(d)
      &\Longrightarrow\operatorname{IsQCQS}(d).
  \end{aligned}
  \]
  By \Cref{lemma:nagata-two-open-valuative-lifting}, every valuation datum for
  $(w,d)$ has a unique lift.  Since $w$ is of finite type, it is
  quasi-compact.  The conclusion of
  \Cref{lemma:nagata-dense-source-valuative} under the substitution
  $(f,h)=(w,d)$ is
  \[
    w\text{ is universally closed}.                               \tag{43}
  \]
  By \Cref{lemma:nagata-two-open-final-modification},
  \[
    w\text{ is separated and of finite type}.                     \tag{44}
  \]
  Equations (43)--(44) are the defining conditions for $w$ to be proper.
\end{proof}

\begin{theorem}[{\cite[\href{https://stacks.math.columbia.edu/tag/0F40}{Tag 0F40}]{stacks-project}}]
  \label{thm:nagata-two-open}
  \uses{def:nagata-compactification,def:nagata-two-open-compactification-candidate}
  \lean{}

  For every
  \[
    D_i\in\operatorname{Comp}(s_i)\qquad(i=1,2),
  \]
  the factorization
  \[
    \mathsf C_{12}(U/S;D_1,D_2)
      =
      \bigl(
        U\xrightarrow{j_W(D_1,D_2)}W(D_1,D_2)
          \xrightarrow{w(D_1,D_2)}S,\,
        w(D_1,D_2)\circ j_W(D_1,D_2)=s
      \bigr)
  \]
  is a compactification of $s\colon U\to S$.  Consequently
  $\operatorname{Comp}(s)\neq\varnothing$.
\end{theorem}

\begin{proof}
  \uses{lemma:nagata-two-open-final-modification,
    lemma:nagata-two-open-candidate-proper,
    lemma:nagata-open-immersion-qc-over-separated-target}
  Fix $D_i\in\operatorname{Comp}(s_i)$ and perform the substitution
  $(C_1^0,C_2^0):=(D_1,D_2)$ from
  \Cref{def:nagata-two-open-compactification-candidate}.  Abbreviate
  \[
    (W_D,j_D,w_D)
      =(W(D_1,D_2),j_W(D_1,D_2),w(D_1,D_2)).
  \]
  Under this substitution,
  \Cref{lemma:nagata-two-open-final-modification} states
  \[
    j_D\colon U\hookrightarrow W_D\text{ is an open immersion},\qquad
    w_D\circ j_D=s.                                                \tag{45}
  \]
  and \Cref{lemma:nagata-two-open-candidate-proper} states
  \[
    w_D\colon W_D\to S\text{ is proper},                          \tag{46}
  \]
  hence separated.  Apply
  \Cref{lemma:nagata-open-immersion-qc-over-separated-target} to (45), using
  that $s$ is of finite type and $w_D$ is separated.  Its conclusion is
  \[
    j_D\text{ is quasi-compact}.                                  \tag{47}
  \]
  Equations (45)--(47) are exactly the conditions of
  \Cref{def:nagata-compactification}.  Thus
  \[
    \mathsf C_{12}(U/S;D_1,D_2)\in\operatorname{Comp}(s).
  \]
  Taking $(D_1,D_2)=(C_1^0,C_2^0)$ proves
  $\operatorname{Comp}(s)\neq\varnothing$.
\end{proof}
\section{The Noetherian case}

\begin{definition}[Affine-space immersion attached to a quasi-affine scheme]
  \label{def:nagata-quasi-affine-immersion-affine-space}
  \lean{}

  Let $\pi\colon X\to S$ be locally of finite type and let $X$ be
  quasi-affine.  Put
  \[
    A=\Gamma(X,\mathcal O_X).
  \]
  For $f\in A$, define
  \[
    X_f=\{x\in X\mid f_x\in\mathcal O_{X,x}^{\times}\},
    \qquad u_f\colon X_f\hookrightarrow X.
  \]
  Choose $m\geq1$, elements $f_1,\ldots,f_m\in A$, and affine open
  immersions
  $v_j\colon V_j=\operatorname{Spec}(R_j)\hookrightarrow S$ such that
  \[
    X=\bigcup_{j=1}^mX_{f_j},\qquad
    X_{f_j}=\operatorname{Spec}(A_{f_j}),\qquad
    \pi(X_{f_j})\subseteq v_j(V_j).
  \]
  Define $\pi_j\colon X_{f_j}\to V_j$ and $\theta_j$ by
  \[
    v_j\circ\pi_j=\pi\circ u_{f_j},
    \qquad
    \theta_j=\pi_j^\#\colon R_j\to A_{f_j}.                       \tag{47a}
  \]
  Finite type of $\theta_j$ defines integers $q_j\geq0$ and elements
  $b_{j,\nu}\in A_{f_j}$ satisfying
  \[
    A_{f_j}=\theta_j(R_j)[b_{j,1},\ldots,b_{j,q_j}].              \tag{47b}
  \]
  Choose $c_{j,\nu}\in A$ and $e_{j,\nu}\geq0$ satisfying
  \[
    b_{j,\nu}=c_{j,\nu}/f_j^{e_{j,\nu}}\quad\text{in }A_{f_j}.    \tag{47c}
  \]
  Enumerate the finite set
  $\{c_{j,\nu}\mid1\leq j\leq m,\ 1\leq\nu\leq q_j\}$ as
  $\{a_1,\ldots,a_N\}$.  Substitution of (47c) into (47b) is the equality
  \[
    A_{f_j}
      =\theta_j(R_j)\!\left[
          \frac{a_1}{1},\ldots,\frac{a_N}{1},\frac1{f_j}
        \right]\qquad(1\leq j\leq m).                              \tag{48}
  \]
  Put $n=m+N$.  If $t_1,\ldots,t_n$ are the coordinates of
  $\mathbb A^n_S$, define
  \[
    \iota_X=(\pi;f_1,\ldots,f_m,a_1,\ldots,a_N)
      \colon X\to\mathbb A^n_S                                   \tag{49}
  \]
  by the equalities
  \[
    \iota_X^*(t_j)=f_j\ (1\leq j\leq m),\qquad
    \iota_X^*(t_{m+k})=a_k\ (1\leq k\leq N).
  \]
\end{definition}

\begin{lemma}[Property of the quasi-affine affine-space map,
  {\cite[\href{https://stacks.math.columbia.edu/tag/04II}{Tag 04II}]{stacks-project}}]
  \label{lemma:nagata-quasi-affine-immersion-affine-space}
  \uses{def:nagata-quasi-affine-immersion-affine-space}
  \lean{}

  The morphism $\iota_X\colon X\to\mathbb A^n_S$ defined by (49) is an
  immersion.
\end{lemma}

\begin{proof}
  For $1\leq j\leq m$, define the open subscheme and its open immersion
  \[
    T_j=D(t_j)\times_SV_j
      =\operatorname{Spec}
        \bigl(R_j[t_1,\ldots,t_n,t_j^{-1}]\bigr)
      \subseteq\mathbb A^n_S,
    \qquad
    w_j\colon T_j\hookrightarrow\mathbb A^n_S.
  \]
  By (47a) and (49), the inverse-image open is
  \[
  \begin{aligned}
    \iota_X^{-1}(T_j)
      &=\iota_X^{-1}\bigl(D(t_j)\bigr)\cap\pi^{-1}(V_j)\\
      &=X_{f_j}\cap\pi^{-1}(V_j)\\
      &=X_{f_j}.
  \end{aligned}                                                   \tag{50}
  \]
  Define the pullback morphism
  $\iota_{X,j}=\iota_X|_{X_{f_j}}\colon X_{f_j}\to T_j$; it satisfies
  \[
    w_j\circ\iota_{X,j}=\iota_X\circ u_{f_j}.
  \]
  On coordinate rings, $\iota_{X,j}$ is the morphism corresponding to
  \[
  \begin{aligned}
    \Theta_j\colon
      R_j[t_1,\ldots,t_n,t_j^{-1}]
        &\longrightarrow A_{f_j},\\
    t_\ell&\longmapsto
      \begin{cases}
        f_\ell/1,&1\leq\ell\leq m,\\
        a_{\ell-m}/1,&m<\ell\leq n,
      \end{cases}\\
    t_j^{-1}&\longmapsto f_j^{-1}.
  \end{aligned}
  \]
  By (48), its image satisfies the chain of inclusions
  \[
  \begin{aligned}
    A_{f_j}
      &=\theta_j(R_j)
        [a_1/1,\ldots,a_N/1,f_j^{-1}]\\
      &\subseteq\operatorname{im}(\Theta_j)
       \subseteq A_{f_j},
  \end{aligned}
  \]
  Antisymmetry of the two outer inclusions is
  \[
    \operatorname{im}(\Theta_j)=A_{f_j},\qquad
    \iota_{X,j}\colon X_{f_j}\hookrightarrow T_j
      \text{ is a closed immersion}.                             \tag{50a}
  \]
  Define the open subscheme and its inclusion by
  \[
    T=\bigcup_{j=1}^mT_j,
    \qquad w_T\colon T\hookrightarrow\mathbb A^n_S.
  \]
  Equation (50) and the cover equality in
  \Cref{def:nagata-quasi-affine-immersion-affine-space} form the chain
  \[
    \iota_X^{-1}(T)
      =\bigcup_{j=1}^m\iota_X^{-1}(T_j)
      =\bigcup_{j=1}^mX_{f_j}
      =X.                                                         \tag{50b}
  \]
  The universal property of the open immersion $w_T$, applied to (50b),
  defines the unique morphism $\iota_T\colon X\to T$ satisfying
  \[
    w_T\circ\iota_T=\iota_X.                                    \tag{50c}
  \]
  For every $j$, (50), (50a), and (50c) identify the base change of
  $\iota_T$ along $T_j\hookrightarrow T$ as
  \[
    X\times_TT_j=X_{f_j},\qquad
    \iota_T\times_TT_j=\iota_{X,j}
      \colon X_{f_j}\hookrightarrow T_j.                         \tag{50d}
  \]
  Since $T=\bigcup_jT_j$ and every morphism in (50d) is a closed immersion,
  \cite[\href{https://stacks.math.columbia.edu/tag/01HL}{Tag 01HL}]
  {stacks-project}, applied to $\iota_T$, asserts
  \[
    \iota_T\colon X\hookrightarrow T\quad\text{a closed immersion}.
  \]
  Finally, (50c) is the factorization
  \[
    \iota_X
      =w_T\circ\iota_T
      \colon X\hookrightarrow T\hookrightarrow\mathbb A^n_S,
  \]
  which is an immersion by
  \cite[\href{https://stacks.math.columbia.edu/tag/01IO}{Tag 01IO}]
  {stacks-project}.
\end{proof}

\begin{definition}[Affine compactification construction]
  \label{def:nagata-affine-compactification}
  \uses{def:nagata-quasi-affine-immersion-affine-space,
    def:nagata-immersion-compactification,
    def:nagata-scheme-theoretic-image,
    def:nagata-scheme-image-factor,
    lemma:nagata-quasi-affine-immersion-affine-space,
    lemma:nagata-finite-type-to-separated-target}
  \lean{}

  Let $s_U\colon U\to S$ be separated and of finite type, with $U$ affine.
  Let $p_n$ and $\rho_n$ be the structural projections, and let $\alpha_n$ be
  the standard open immersion defined by the homogeneous coordinate
  $T_0\neq0$:
  \[
    p_n\colon\mathbb A^n_S\to S,
    \qquad
    \rho_n\colon\mathbb P^n_S\to S,
    \qquad
    \alpha_n\colon\mathbb A^n_S\hookrightarrow\mathbb P^n_S,
    \qquad
    \rho_n\circ\alpha_n=p_n.
  \]
  Let $\iota_U\colon U\to\mathbb A^n_S$ be the morphism of
  \Cref{def:nagata-quasi-affine-immersion-affine-space}, with structural
  equality $p_n\circ\iota_U=s_U$, and define
  \[
    e_U=\alpha_n\circ\iota_U\colon U\to\mathbb P^n_S.              \tag{51}
  \]
  By \Cref{lemma:nagata-quasi-affine-immersion-affine-space,lemma:nagata-finite-type-to-separated-target},
  the domain predicate
  required for the scheme-theoretic image is the chain
  \[
  \begin{aligned}
    \operatorname{Projective}(\rho_n)
      &\Longrightarrow\operatorname{Proper}(\rho_n)
       \Longrightarrow\operatorname{IsSeparated}(\rho_n)
       &&\text{by \cite[Tag 01WC]{stacks-project}},\\
    \operatorname{FiniteType}(s_U)\wedge
      \operatorname{IsSeparated}(\rho_n)\wedge
      (\rho_n\circ e_U=s_U)
      &\Longrightarrow\operatorname{FiniteType}(e_U)
       \Longrightarrow\operatorname{IsQuasiCompact}(e_U),\\
    \operatorname{IsImmersion}(\iota_U)\wedge
      \operatorname{IsOpenImmersion}(\alpha_n)\wedge
      (e_U=\alpha_n\circ\iota_U)
      &\Longrightarrow\operatorname{IsImmersion}(e_U)
       \Longrightarrow\operatorname{IsSeparated}(e_U)\\
      &\Longrightarrow\operatorname{QuasiSeparated}(e_U),\\
    \operatorname{IsQuasiCompact}(e_U)\wedge
      \operatorname{QuasiSeparated}(e_U)
      &\Longrightarrow\operatorname{IsQCQS}(e_U).
  \end{aligned}
  \]
  With the notation of
  \Cref{def:nagata-scheme-theoretic-image,def:nagata-scheme-image-factor},
  define
  \[
    h_U:=i_{e_U}\colon\operatorname{SIm}(e_U)\hookrightarrow
      \mathbb P^n_S,
    \qquad
    j_U:=e_U^{\mathrm{im}}\colon U\to\operatorname{SIm}(e_U),
    \qquad
    h_U\circ j_U=e_U.                                             \tag{51a}
  \]
  Define
  \[
    p_U:=\rho_n\circ h_U
      \colon\operatorname{SIm}(e_U)\to S.                         \tag{51b}
  \]
  Define
  \[
    \mathsf C_{\mathrm{aff}}(U/S)
      =\mathsf C(e_U)
      =
      \left(
        U\xrightarrow{j_U=e_U^{\mathrm{im}}}
        \operatorname{SIm}(e_U)
        \xrightarrow{p_U}S
      \right),                                                     \tag{52}
  \]
  where $\mathsf C(e_U)$ is the construction of
  \Cref{def:nagata-immersion-compactification}; explicitly,
  (51a)--(51b) identify its morphisms as
  $j_U=e_U^{\mathrm{im}}$ and $p_U=\rho_n\circ i_{e_U}$.
\end{definition}

\begin{proposition}[Property of the affine compactification construction]
  \label{prop:nagata-affine-compactification}
  \uses{def:nagata-affine-compactification}
  \lean{}

  The factorization $\mathsf C_{\mathrm{aff}}(U/S)$ in (52) is a
  compactification of $s_U\colon U\to S$.
\end{proposition}

\begin{proof}
  \uses{lemma:nagata-quasi-affine-immersion-affine-space,
    lemma:nagata-immersion-into-proper,
    lemma:nagata-finite-type-to-separated-target}
  By \Cref{lemma:nagata-quasi-affine-immersion-affine-space}, the predicates
  \[
    \operatorname{IsImmersion}(\iota_U),
    \qquad
    \operatorname{IsOpenImmersion}(\alpha_n)
      \Longrightarrow\operatorname{IsImmersion}(\alpha_n)
  \]
  hold.  Stability of immersions under composition, applied to (51), is the
  deduction
  \[
    \operatorname{IsImmersion}(\iota_U)
      \wedge\operatorname{IsImmersion}(\alpha_n)
      \Longrightarrow
    \operatorname{IsImmersion}(e_U).                              \tag{52a}
  \]
  From (51) and the structural equalities in
  \Cref{def:nagata-affine-compactification}, one has
  \[
  \begin{aligned}
    \rho_n\circ e_U
      &=\rho_n\circ(\alpha_n\circ\iota_U)\\
      &=(\rho_n\circ\alpha_n)\circ\iota_U\\
      &=p_n\circ\iota_U\\
      &=s_U.
  \end{aligned}                                                   \tag{52b}
  \]
  The hypotheses of
  \Cref{lemma:nagata-finite-type-to-separated-target}, instantiated by
  \[
    (f,g,h)=(e_U,\rho_n,s_U),
  \]
  are
  \[
    \rho_n\circ e_U=s_U,\qquad
    \operatorname{FiniteType}(s_U),\qquad
    \operatorname{Separated}(\rho_n).
  \]
  Substitution of the preceding three predicates into that lemma, together
  with the defining implication (P1), is the chain
  \[
    \operatorname{FiniteType}(e_U),
    \qquad
    \operatorname{FiniteType}(e_U)
      \Longrightarrow\operatorname{QuasiCompact}(e_U),
    \qquad
    \operatorname{QuasiCompact}(e_U).                             \tag{52c}
  \]
  The projective-space projection satisfies
  \[
    \operatorname{Projective}(\rho_n)
      \Longrightarrow\operatorname{Proper}(\rho_n)               \tag{52d}
  \]
  by
  \cite[\href{https://stacks.math.columbia.edu/tag/01WC}{Tag 01WC}]
  {stacks-project}.  Instantiate
  \Cref{lemma:nagata-immersion-into-proper} with
  \[
    (g,p_Y,f)=(e_U,\rho_n,s_U).
  \]
  Equations (52a)--(52d) verify, respectively, that $g$ is an immersion,
  that $p_Y\circ g=f$, that $g$ is quasi-compact, and that $p_Y$ is proper.
  The membership statement is
  \[
    \mathsf C(e_U)=\mathsf C_{\mathrm{aff}}(U/S)
      \in\operatorname{Comp}(s_U).
  \]
\end{proof}

\begin{definition}[Affine-neighbourhood candidate set]
  \label{def:nagata-affine-neighbourhood-generic-points}
  \lean{}

  Let $X$ be quasi-separated with distinct irreducible components
  $X^{(1)},\ldots,X^{(r)}$, let $\eta_i$ be the generic point of $X^{(i)}$,
  and let $x\in|X|$.  Define
  \[
    \mathscr A_X(x;\eta_1,\ldots,\eta_r)
      =
      \left\{
        V\subseteq X\ \middle|\
        \begin{array}{l}
          V\text{ is an affine open subscheme},\\
          x,\eta_1,\ldots,\eta_r\in|V|
        \end{array}
      \right\}.                                                    \tag{53}
  \]
\end{definition}

\begin{lemma}[Existence of the affine neighbourhood,
  {\cite[\href{https://stacks.math.columbia.edu/tag/01ZX}{Tag 01ZX}]{stacks-project}}]
  \label{lemma:nagata-affine-neighbourhood-generic-points}
  \uses{def:nagata-affine-neighbourhood-generic-points}
  \lean{}

  Let \(X,X^{(1)},\ldots,X^{(r)},\eta_1,\ldots,\eta_r,x\) satisfy the
  hypotheses of \Cref{def:nagata-affine-neighbourhood-generic-points}.
  The set in (53) is nonempty:
  \[
    \mathscr A_X(x;\eta_1,\ldots,\eta_r)\neq\varnothing.
  \]
\end{lemma}

\begin{proof}
  Define
  $I_x=\{i\in\{1,\ldots,r\}\mid x\in X^{(i)}\}$.  After reindexing,
  write $I_x=\{1,\ldots,m\}$, where $0\leq m\leq r$; equivalently,
  \[
    x\in\bigcap_{i=1}^mX^{(i)},\qquad
    x\notin\bigcup_{i=m+1}^rX^{(i)}.                              \tag{54}
  \]
  Define the open subscheme
  \[
    O_x:=X\setminus\bigcup_{i=m+1}^rX^{(i)}.
  \]
  Equation (54) is $x\in|O_x|$.  Choose an affine open immersion
  $w\colon W\hookrightarrow O_x\hookrightarrow X$ satisfying
  \[
    x\in|W|,\qquad
    W\cap X^{(i)}=\varnothing\quad(i>m).
  \]
  For $1\leq i\leq m$, the specialization relation is
  $\eta_i\leadsto x$ because
  $\overline{\{\eta_i\}}=|X^{(i)}|$ and $x\in X^{(i)}$.  Stability of
  $|W|$ under generization is precisely the displayed membership
  \[
    \eta_i\in|W|\quad(1\leq i\leq m).                             \tag{55}
  \]
  By
  \cite[\href{https://stacks.math.columbia.edu/tag/01ZV}{Tag 01ZV}]
  {stacks-project}, choose affine open immersions
  $q_i\colon Q_i\hookrightarrow X$ for $m<i\leq r$ satisfying
  \[
    \eta_i\in|Q_i|,
    \qquad Q_i\cap Q_k=\varnothing\quad(m<i<k\leq r).            \tag{55a}
  \]
  For each $i>m$, quasi-separatedness of $X$ and affineness of $W,Q_i$
  establish the first property below; the defining property of $W$
  establishes the second:
  \[
    \operatorname{QuasiCompact}(W\cap Q_i),
    \qquad
    \eta_i\notin|W\cap Q_i|.                                    \tag{55b}
  \]
  Apply
  \cite[\href{https://stacks.math.columbia.edu/tag/00EV}{Tag 00EV}]
  {stacks-project} to the quasi-compact open $W\cap Q_i\subseteq Q_i$ and
  the point $\eta_i\in|Q_i|\setminus|W\cap Q_i|$.  With
  $B_i:=\Gamma(Q_i,\mathcal O_{Q_i})$, choose $g_i\in B_i$ and define
  \[
    Q_i':=D_{Q_i}(g_i)=\operatorname{Spec}((B_i)_{g_i}),
    \qquad q_i'\colon Q_i'\hookrightarrow Q_i\hookrightarrow X,
  \]
  with the exact properties
  \[
    \eta_i\in|Q_i'|,\qquad Q_i'\cap W=\varnothing.                 \tag{56}
  \]
  Equations (55a) and (56) are the disjointness identities
  \[
    Q_i'\cap Q_k'=\varnothing\quad(m<i<k\leq r),
    \qquad
    W\cap Q_i'=\varnothing\quad(m<i\leq r).                     \tag{56a}
  \]
  Define the open subscheme and its inclusion by
  \[
    V:=W\cup\bigcup_{i=m+1}^rQ_i'\subseteq X,
    \qquad v\colon V\hookrightarrow X.
  \]
  Equation (56a) identifies this union with the coproduct
  \[
    V=W\amalg\coprod_{i=m+1}^rQ_i'.                              \tag{56b}
  \]
  The finite disjoint-union affine isomorphism
  (\cite[\href{https://stacks.math.columbia.edu/tag/01I5}{Tag 01I5}]
  {stacks-project}), applied to (56b), is
  \[
    V\cong
      \operatorname{Spec}\left(
        \Gamma(W,\mathcal O_W)
          \times\prod_{i=m+1}^r\Gamma(Q_i',\mathcal O_{Q_i'})
      \right),
    \qquad V\text{ is affine},
  \]
  where the second assertion is the definition of an affine scheme.  The
  point memberships in (55)--(56) expand to
  \[
  \begin{gathered}
    x\in|W|\subseteq|V|,\\
    \eta_i\in|W|\subseteq|V|\quad(1\leq i\leq m),\\
    \eta_i\in|Q_i'|\subseteq|V|\quad(m<i\leq r),
  \end{gathered}
  \]
  By Definition~\ref{def:nagata-affine-neighbourhood-generic-points},
  equation (53), these properties are exactly the symbolic statements
  \[
    V\in\mathscr A_X(x;\eta_1,\ldots,\eta_r),
    \qquad
    \mathscr A_X(x;\eta_1,\ldots,\eta_r)\neq\varnothing.
  \]
\end{proof}

\begin{definition}[Selected affine neighbourhood containing generic points]
  \label{def:nagata-selected-affine-neighbourhood-generic-points}
  \uses{def:nagata-choice-notation,
    def:nagata-affine-neighbourhood-generic-points}
  
  Let \(X\) be quasi-separated, let
  \(X^{(1)},\ldots,X^{(r)}\) be its distinct irreducible components, let
  \(\eta_i\in|X^{(i)}|\) be the generic point of \(X^{(i)}\), and let
  \(x\in|X|\).  Using
  \Cref{lemma:nagata-affine-neighbourhood-generic-points}, define
  \[
    A_X(x;\eta_1,\ldots,\eta_r)
      =
      \operatorname{ch}\bigl(
        \mathscr A_X(x;\eta_1,\ldots,\eta_r)
      \bigr)
      \in\mathscr A_X(x;\eta_1,\ldots,\eta_r).                    \tag{57}
  \]
  The final membership in (57) is the defining property
  $\operatorname{ch}(E)\in E$ of \Cref{def:nagata-choice-notation}, with
  $E=\mathscr A_X(x;\eta_1,\ldots,\eta_r)$.  Unpacking that membership by
  (53) defines the open immersion
  \[
    a_{X,x}\colon A_X(x;\eta_1,\ldots,\eta_r)\hookrightarrow X
  \]
  and records the exact properties
  \[
  \begin{gathered}
    \operatorname{IsOpenImmersion}(a_{X,x}),\qquad
    A_X(x;\eta_1,\ldots,\eta_r)\text{ is affine},\\
    \{x,\eta_1,\ldots,\eta_r\}
      \subseteq|A_X(x;\eta_1,\ldots,\eta_r)|.
  \end{gathered}
  \]
\end{definition}

\begin{definition}[Finite affine cover with selected dense intersection]
  \label{def:nagata-dense-affine-cover}
  \uses{def:nagata-selected-affine-neighbourhood-generic-points}
  \lean{}

  Let $X$ be quasi-compact and quasi-separated with irreducible components
  $X^{(1)},\ldots,X^{(r)}$ and generic points
  $\eta_1,\ldots,\eta_r$.  Define
  \[
    A_x:=A_X(x;\eta_1,\ldots,\eta_r),
    \qquad a_x:=a_{X,x}\colon A_x\hookrightarrow X
    \quad(x\in|X|).
  \]
  Equation (57), unpacked by (53), is the membership
  $x\in|A_x|$ for every $x\in|X|$; equivalently,
  \[
    X=\bigcup_{x\in|X|}A_x.
  \]
  Define
  \[
    \mathscr F_X
      :=\coprod_{n\geq1}
        \left\{
          (x_1,\ldots,x_n)\in|X|^n\ \middle|\
          X=\bigcup_{\ell=1}^nA_{x_\ell}
        \right\}.
  \]
  Quasi-compactness of $X$, applied to the displayed open cover, is the
  existence statement
  \[
    \exists n\geq1\ \exists(x_1,\ldots,x_n)\in|X|^n,\qquad
      X=\bigcup_{\ell=1}^nA_{x_\ell};
  \]
  by the definition of $\mathscr F_X$, this is
  $\mathscr F_X\neq\varnothing$.  Define
  \[
    \bigl(n,(x_1,\ldots,x_n)\bigr)
      :=\operatorname{ch}(\mathscr F_X),\qquad
    U_\ell:=A_{x_\ell},\qquad
    u_\ell:=a_{x_\ell}\colon U_\ell\hookrightarrow X,
  \]
  and define the open subscheme and its inclusion by
  \[
    D_X:=U_1\times_X\cdots\times_XU_n
      =\bigcap_{\ell=1}^nU_\ell,
    \qquad d_X\colon D_X\hookrightarrow X.                       \tag{58}
  \]
  Write
  \[
    \mathcal U_X=(U_1,\ldots,U_n).
  \]
\end{definition}

\begin{lemma}[Properties of the selected dense affine cover]
  \label{lemma:nagata-dense-affine-cover}
  \uses{def:nagata-dense-affine-cover}
  \lean{}

  The tuple $\mathcal U_X=(U_1,\ldots,U_n)$ satisfies
  \[
    X=\bigcup_{\ell=1}^nU_\ell,\qquad
    U_\ell\text{ is affine open},
    \qquad
    \overline{|D_X|}=|X|.
  \]
\end{lemma}

\begin{proof}
  \uses{lemma:nagata-affine-neighbourhood-generic-points}
  By \Cref{def:nagata-choice-notation} and (58),
  \[
    \bigl(n,(x_1,\ldots,x_n)\bigr)\in\mathscr F_X.
  \]
  Substitution in the defining coproduct of $\mathscr F_X$ is
  \[
    X=\bigcup_{\ell=1}^nA_{x_\ell}
      =\bigcup_{\ell=1}^nU_\ell.                                 \tag{58a}
  \]
  For every $(i,\ell)\in\{1,\ldots,r\}\times\{1,\ldots,n\}$,
  equations (53), (57), and (58) are the properties
  \[
    U_\ell=A_X(x_\ell;\eta_1,\ldots,\eta_r)
      \text{ is affine open},
    \qquad
    \eta_i\in|U_\ell|.                                          \tag{58b}
  \]
  The definition of $D_X$ in (58) and (58b) are, for every $i$,
  \[
    \eta_i\in\bigcap_{\ell=1}^n|U_\ell|=|D_X|,
    \qquad
    \{\eta_1,\ldots,\eta_r\}\subseteq|D_X|.                    \tag{58c}
  \]
  Since $\overline{\{\eta_i\}}=|X^{(i)}|$ and
  $X=\bigcup_{i=1}^rX^{(i)}$,
  \[
    |X|
      =\bigcup_{i=1}^r\overline{\{\eta_i\}}
      \subseteq\overline{|D_X|}
      \subseteq|X|.
  \]
  Antisymmetry of inclusion applied to the first and last terms is the
  equality
  \[
    \overline{|D_X|}=|X|.
  \]
\end{proof}

\begin{definition}[Partial unions of the selected affine cover]
  \label{def:nagata-partial-affine-unions}
  \uses{def:nagata-dense-affine-cover}
  \lean{}

  Let $s\colon X\to S$ be a morphism, and let
  $\mathcal U_X=(U_1,\ldots,U_n)$ be the selected cover of
  \Cref{def:nagata-dense-affine-cover}, with
  $u_\ell\colon U_\ell\hookrightarrow X$ as in (58).  For
  $1\leq r\leq n$, define the open subscheme
  \[
    X_r:=\bigcup_{\ell=1}^rU_\ell\subseteq X,
    \qquad v_r\colon X_r\hookrightarrow X,
  \]
  and define
  \[
    s_r:=s\circ v_r\colon X_r\to S.
  \]
  For $2\leq r\leq n$, define the overlap and its projections by
  \[
    W_r:=X_{r-1}\times_XU_r=X_{r-1}\cap U_r,
    \qquad
    w_{r,1}\colon W_r\hookrightarrow X_{r-1},
    \qquad
    w_{r,2}\colon W_r\hookrightarrow U_r.
  \]
\end{definition}

\begin{lemma}[Density needed for induction on the affine cover]
  \label{lemma:nagata-partial-union-overlap-dense}
  \uses{def:nagata-partial-affine-unions}
  \lean{}

  For $2\leq r\leq n$,
  \[
    \overline{|W_r|}^{\,|X_r|}
      =\overline{|X_{r-1}\cap U_r|}^{\,|X_r|}=|X_r|.
  \]
\end{lemma}

\begin{proof}
  \uses{lemma:nagata-dense-affine-cover}
  The definitions (58) and
  $U_1\subseteq X_{r-1}=\bigcup_{\ell=1}^{r-1}U_\ell$ form the inclusion chain
  \[
    D_X=\bigcap_{\ell=1}^nU_\ell
      \subseteq U_1\cap U_r
      \subseteq X_{r-1}\cap U_r=W_r.                             \tag{59}
  \]
  By \Cref{lemma:nagata-dense-affine-cover},
  $\overline{|D_X|}^{\,|X|}=|X|$.  Since $X_r$ is open in $X$ and
  $D_X\subseteq X_r$,
  \[
    \overline{|D_X|}^{\,|X_r|}
      =|X_r|\cap\overline{|D_X|}^{\,|X|}
      =|X_r|.                                                      \tag{60}
  \]
  Equations (59)--(60) form the inclusion chain
  \[
    |X_r|
      =\overline{|D_X|}^{\,|X_r|}
      \subseteq
        \overline{|X_{r-1}\cap U_r|}^{\,|X_r|}
      \subseteq|X_r|.
  \]
  Antisymmetry of the two outer inclusions is the equality
  \[
    \overline{|W_r|}^{\,|X_r|}
      =\overline{|X_{r-1}\cap U_r|}^{\,|X_r|}
      =|X_r|.
  \]
\end{proof}

\begin{definition}[Recursive Noetherian compactification construction]
  \label{def:nagata-noetherian-compactification}
  \uses{def:nagata-choice-notation,def:nagata-affine-compactification,
    def:nagata-dense-affine-cover,
    def:nagata-partial-affine-unions,
    def:nagata-two-open-compactification-candidate}
  \lean{}

  Let $S$ be Noetherian and let $s\colon X\to S$ be separated and of finite
  type.  If $|X|=\varnothing$, then for every morphism
  $\tau\colon T\to S$ define
  \[
    X_T:=X\times_{S,(s,\tau)}T,\qquad
    p_T:=\operatorname{pr}_T\colon X_T\to T.
  \]
  Define
  \[
  \begin{gathered}
    C_\varnothing(s):=
      (X,\operatorname{id}_X,s,e_\varnothing),\qquad
    e_\varnothing\colon s\circ\operatorname{id}_X=s,\\
    \operatorname{IsQuasiCompactOpenImmersion}(\operatorname{id}_X),
    \qquad
    (\forall\,V=\operatorname{Spec}(A)\hookrightarrow S)\quad
      s^{-1}(V)=\operatorname{Spec}(0),\quad
      0=A/(1),\\
    \bigl[(\forall\,V=\operatorname{Spec}(A)\hookrightarrow S)\quad
      \Gamma(s^{-1}(V),\mathcal O_X)=0=A/(1)\bigr]
      \Longrightarrow\operatorname{FinitePresentation}(s),\\
    \operatorname{FinitePresentation}(s)
      \Longrightarrow\operatorname{FiniteType}(s),\qquad
    \Delta_s\colon X\to X\times_SX
      =\operatorname{id}_{\varnothing},\qquad
    \operatorname{IsClosedImmersion}(\Delta_s),\\
    \operatorname{IsSeparated}(s),\qquad
    (\forall\tau\colon T\to S)\quad |X_T|=\varnothing,\\
    (\forall\tau\colon T\to S)\quad
      \operatorname{Im}(|p_T|)=\varnothing
      \text{ is closed in }|T|,\\
    \operatorname{UniversallyClosed}(s),\qquad
    \bigl(\operatorname{FiniteType}(s)\wedge
      \operatorname{IsSeparated}(s)\wedge
      \operatorname{UniversallyClosed}(s)\bigr)
      \Longleftrightarrow\operatorname{IsProper}(s),\\
    \operatorname{IsProper}(s),\qquad
    \chi_\varnothing(s)\colon
      C_\varnothing(s)\in\operatorname{Compactification}(s).
  \end{gathered}                                                    \tag{NC0}
  \]

  Assume for the remainder of this definition that $|X|\neq\varnothing$.
  The Noetherian finite-type theorem and the topological properties of a
  Noetherian scheme are the following implication chain:
  \[
  \begin{aligned}
    &\operatorname{Noetherian}(S)\wedge\operatorname{FiniteType}(s)\\
    &\qquad\Longrightarrow \operatorname{Noetherian}(X)\\
    &\qquad\Longrightarrow
      \bigl(\operatorname{QuasiCompact}(X)\wedge
        \operatorname{QuasiSeparated}(X)\wedge
        1\leq\#\operatorname{Irr}(X)<\infty\bigr).
  \end{aligned}                                                     \tag{NC1}
  \]
  Define the set of enumerated component/generic-point data
  \[
    \mathscr I_X:=
      \coprod_{r\geq1}
      \left\{(Z_i,\eta_i)_{i=1}^r\ \middle|\
      \begin{array}{l}
        Z_i\in\operatorname{Irr}(X),\quad
        \eta_i\in|Z_i|,\quad
        \overline{\{\eta_i\}}=|Z_i|,\\
        i\neq j\Longrightarrow Z_i\neq Z_j,\quad
        \operatorname{Irr}(X)=\{Z_1,\ldots,Z_r\}
      \end{array}\right\}.                                      \tag{NC2}
  \]
  Finiteness in (NC1), together with existence and uniqueness of the generic
  point of each irreducible closed subset of a scheme, is the assertion
  \[
    \mathscr I_X\neq\varnothing.                                  \tag{NC3}
  \]
  By \Cref{def:nagata-choice-notation}, define
  \[
    \mathsf I_X:=\operatorname{ch}(\mathscr I_X)
      =\bigl(X^{(i)},\eta_i\bigr)_{i=1}^{r_X}
      \in\mathscr I_X.                                            \tag{NC4}
  \]
  Membership in (NC4) expands to
  \[
  \begin{gathered}
    \operatorname{Irr}(X)=\{X^{(1)},\ldots,X^{(r_X)}\},\qquad
    \eta_i\in|X^{(i)}|,\qquad
    \overline{\{\eta_i\}}=|X^{(i)}|,\\
    i\neq j\Longrightarrow X^{(i)}\neq X^{(j)}.
  \end{gathered}                                                   \tag{NC5}
  \]
  Equations (NC1) and (NC5) instantiate
  \Cref{def:nagata-dense-affine-cover}.  Use
  $\mathcal U_X=(U_1,\ldots,U_n)$ and
  $u_r\colon U_r\hookrightarrow X$ from
  \Cref{def:nagata-dense-affine-cover}.  Use $X_r$, $v_r$, and
  $s_r=s\circ v_r$ from \Cref{def:nagata-partial-affine-unions}, and define
  \[
    s_{U_r}:=s\circ u_r\colon U_r\to S\qquad(1\leq r\leq n).
  \]
  The cited implications
  \cite[\href{https://stacks.math.columbia.edu/tag/01T6}{Tag 01T6}]
  {stacks-project},
  \cite[\href{https://stacks.math.columbia.edu/tag/01OX}{Tag 01OX}]
  {stacks-project}, and
  \cite[\href{https://stacks.math.columbia.edu/tag/01TU}{Tag 01TU}]
  {stacks-project}, together with the definitions of a Noetherian scheme and
  a quasi-compact morphism, are, for every $r$, the implications
  \[
  \begin{gathered}
    S\text{ Noetherian}\ \wedge\ \operatorname{FiniteType}(s)
      \Longrightarrow X\text{ Noetherian},\\
    X\text{ Noetherian}
      \Longrightarrow
      \bigl[\forall W\subseteq X\text{ open},\ W\text{ quasi-compact}\bigr],\\
    \operatorname{IsOpenImmersion}(u_r)\ \wedge\ X\text{ Noetherian}
      \Longrightarrow
      \bigl[\forall W\subseteq X\text{ open and quasi-compact},\quad
        u_r^{-1}(W)=U_r\cap W\text{ is quasi-compact}\bigr]
      \Longrightarrow\operatorname{QuasiCompact}(u_r)
      \Longrightarrow\operatorname{FinitePresentation}(u_r)
      \Longrightarrow\operatorname{FiniteType}(u_r),\\
    \operatorname{FiniteType}(s)\ \wedge\operatorname{FiniteType}(u_r)
      \Longrightarrow\operatorname{FiniteType}(s_{U_r}),\\
    \operatorname{Separated}(s)\ \wedge\operatorname{Separated}(u_r)
      \Longrightarrow\operatorname{Separated}(s_{U_r}).
  \end{gathered}                                                  \tag{60c}
  \]
  For $2\leq r\leq n$, define the open immersions
  \[
    \kappa_{r-1,r}\colon X_{r-1}\hookrightarrow X_r,
    \qquad
    \lambda_r\colon U_r\hookrightarrow X_r
  \]
  by the factorization equalities
  \[
    v_r\circ\kappa_{r-1,r}=v_{r-1},
    \qquad
    v_r\circ\lambda_r=u_r.                                      \tag{60a}
  \]
  Their structural and cover equalities are
  \[
    s_r\circ\kappa_{r-1,r}=s_{r-1},
    \qquad
    s_r\circ\lambda_r=s_{U_r},
    \qquad
    X_r=X_{r-1}\cup U_r.                                        \tag{60b}
  \]
  For $1\leq r\leq n$, define, by
  \Cref{def:nagata-compactification},
  \[
    \operatorname{Comp}(s_r):=\operatorname{Compactification}(s_r),
    \qquad
    \operatorname{Comp}(s_{U_r})
      :=\operatorname{Compactification}(s_{U_r}).
  \]
  For $r\geq2$ and
  \[
    C\in\operatorname{Comp}(s_{r-1}),\qquad
    D\in\operatorname{Comp}(s_{U_r}),
  \]
  define
  \[
    \mathfrak T_r(C,D)
      :=\mathsf C_{12}(X_r/S;C,D),                                \tag{61}
  \]
  where the symbols of
  \Cref{def:nagata-two-open-compactification-candidate} are instantiated by
  \[
  \begin{gathered}
    (U,U_1,U_2,U_{12},s,s_1,s_2)
      =(X_r,X_{r-1},U_r,W_r,s_r,s_{r-1},s_{U_r}),\\
    (U_1\hookrightarrow U,U_2\hookrightarrow U)
      =(\kappa_{r-1,r},\lambda_r).
  \end{gathered}                                                  \tag{61a}
  \]
  Define the density witness by the equality in
  \Cref{lemma:nagata-partial-union-overlap-dense}:
  \[
    \delta_r\colon\overline{|W_r|}^{\,|X_r|}=|X_r|.              \tag{61b}
  \]
  Under (61a), the substitution $(D_1,D_2)=(C,D)$ and the witness
  $\delta_r$ instantiate \Cref{thm:nagata-two-open} as
  \[
    \tau_r(C,D)\colon
      \mathfrak T_r(C,D)\in\operatorname{Comp}(s_r),              \tag{62}
  \]
  which defines the membership witness $\tau_r(C,D)$.
  For $1\leq r\leq n$, define
  \[
    C_{\mathrm{aff},r}:=\mathsf C_{\mathrm{aff}}(U_r/S).
  \]
  The substitution $(U,s_U)=(U_r,s_{U_r})$ in
  \Cref{prop:nagata-affine-compactification} is
  \[
    \chi_{\mathrm{aff},r}\colon
      C_{\mathrm{aff},r}\in\operatorname{Comp}(s_{U_r}),        \tag{62a}
  \]
  which defines the membership witness $\chi_{\mathrm{aff},r}$.  Define pairs
  $(\mathsf C_r,\chi_r)$ recursively by
  \[
    (\mathsf C_1,\chi_1)
      =(C_{\mathrm{aff},1},\chi_{\mathrm{aff},1}),                 \tag{63}
  \]
  and, for $2\leq r\leq n$, by
  \[
    \mathsf C_r
      =\mathfrak T_r(\mathsf C_{r-1},C_{\mathrm{aff},r}),\qquad
    \chi_r
      =\tau_r(\mathsf C_{r-1},C_{\mathrm{aff},r}).                 \tag{64}
  \]
  Combining the empty case (NC0) with the nonempty recursion (63)--(64),
  define
  \[
    \mathsf C_{\mathrm{Noeth}}(X/S):=
    \begin{cases}
      C_\varnothing(s),&|X|=\varnothing,\\
      \mathsf C_n,&|X|\neq\varnothing,
    \end{cases}
    \qquad
    \chi_{\mathrm{Noeth}}(s):=
    \begin{cases}
      \chi_\varnothing(s),&|X|=\varnothing,\\
      \chi_n,&|X|\neq\varnothing.
    \end{cases}                                                    \tag{65}
  \]
\end{definition}

\begin{theorem}[Nagata compactification over a Noetherian base]
  \label{thm:nagata-noetherian}
  \uses{def:nagata-compactification,def:nagata-noetherian-compactification}
  \lean{}

  Let $S$ be Noetherian and let $s\colon X\to S$ be separated and of finite
  type.  Then
  \[
    \mathsf C_{\mathrm{Noeth}}(X/S)
      \in\operatorname{Comp}(s).
  \]
  In particular,
  \[
    \operatorname{Compactifiable}(s).
  \]
\end{theorem}

\begin{proof}
  \uses{prop:nagata-affine-compactification,
    lemma:nagata-dense-affine-cover,
    lemma:nagata-partial-union-overlap-dense,
    thm:nagata-two-open}
  If $|X|=\varnothing$, (NC0) and the empty branch of (65) are the
  membership judgment
  \[
    \chi_{\mathrm{Noeth}}(s)=\chi_\varnothing(s)\colon
      \mathsf C_{\mathrm{Noeth}}(X/S)=C_\varnothing(s)
        \in\operatorname{Comp}(s).
  \]
  Thus assume $|X|\neq\varnothing$.  Equations (NC1)--(NC5) are the typed
  input
  \[
    \operatorname{QuasiCompact}(X)\wedge
      \operatorname{QuasiSeparated}(X),\qquad
    \operatorname{Irr}(X)=\{X^{(1)},\ldots,X^{(r_X)}\},\qquad
    \eta_i\in|X^{(i)}|,\qquad
    \overline{\{\eta_i\}}=|X^{(i)}|.                              \tag{65a}
  \]
  Substitution of (65a) into
  \Cref{def:nagata-dense-affine-cover}, together with
  \Cref{lemma:nagata-dense-affine-cover}, is the following cover and closure
  tuple:
  \[
    X=\bigcup_{\ell=1}^nU_\ell,\qquad
    U_\ell\text{ affine open},\qquad
    \overline{|D_X|}=|X|.                                       \tag{65b}
  \]

  We prove by induction on $r\in\{1,\ldots,n\}$ the membership judgment
  \[
    \chi_r\colon\mathsf C_r\in\operatorname{Comp}(s_r)
      \quad(1\leq r\leq n).                                      \tag{66}
  \]
  For $r=1$, the definitions of the partial union and of $s_{U_1}$ are
  \[
    X_1=U_1,\qquad v_1=u_1,\qquad s_1=s\circ v_1=s\circ u_1=s_{U_1}.
  \]
  The preceding equality together with (62a) and (63) is the identity
  \[
    \chi_1=\chi_{\mathrm{aff},1}\colon
    \mathsf C_1=C_{\mathrm{aff},1}
      \in\operatorname{Comp}(s_1).                               \tag{66a}
  \]
  Fix $2\leq r\leq n$ and assume
  \[
    \chi_{r-1}\colon
      \mathsf C_{r-1}\in\operatorname{Comp}(s_{r-1}).            \tag{66b}
  \]
  Equation (62a) supplies the second input
  \[
    \chi_{\mathrm{aff},r}\colon
      C_{\mathrm{aff},r}\in\operatorname{Comp}(s_{U_r}),        \tag{66c}
  \]
  and (61b) is the density input
  \[
    \delta_r\colon\overline{|W_r|}^{\,|X_r|}=|X_r|.              \tag{66d}
  \]
  Substitute $(C,D)=(\mathsf C_{r-1},C_{\mathrm{aff},r})$ and
  (66b)--(66d) into (62).  The judgment is
  \[
    \tau_r(\mathsf C_{r-1},C_{\mathrm{aff},r})\colon
    \mathfrak T_r(\mathsf C_{r-1},C_{\mathrm{aff},r})
      \in\operatorname{Comp}(s_r).                              \tag{66e}
  \]
  The two substitutions from (64),
  \[
    \mathsf C_r
      =\mathfrak T_r(\mathsf C_{r-1},C_{\mathrm{aff},r}),
    \qquad
    \chi_r=\tau_r(\mathsf C_{r-1},C_{\mathrm{aff},r}),
  \]
  transform (66e) into (66).  This completes the induction.

  By (65b) and \Cref{def:nagata-partial-affine-unions},
  \[
    X_n=\bigcup_{\ell=1}^nU_\ell=X,\qquad
    v_n=\operatorname{id}_X,\qquad
    s_n=s\circ v_n=s.                                            \tag{67}
  \]
  Substitute $r=n$ in (66), then substitute the nonempty branches of (65)
  and (67):
  \[
  \begin{aligned}
    \chi_n&\colon\mathsf C_n\in\operatorname{Comp}(s_n),\\
    \chi_{\mathrm{Noeth}}(s)=\chi_n&\colon
    \mathsf C_{\mathrm{Noeth}}(X/S)
      &=\mathsf C_n
       \in\operatorname{Comp}(s_n)
       =\operatorname{Comp}(s).
  \end{aligned}
  \]
  In particular,
  \[
    \mathsf C_{\mathrm{Noeth}}(X/S)\in\operatorname{Comp}(s)
      \Longrightarrow\operatorname{Comp}(s)\neq\varnothing.
  \]
  By \Cref{def:nagata-compactifiable}, the final predicate is exactly
  \[
    \operatorname{Comp}(s)\neq\varnothing
      \Longleftrightarrow\operatorname{Compactifiable}(s).
  \]
\end{proof}

\section{Finite-presentation reduction and approximation}

\subsection{Finite data and absolute Noetherian models}

\begin{definition}[Selected upper bound of a finite set]
  \label{def:nagata-directed-finite-upper-bound}
  \uses{def:nagata-choice-notation}
  \lean{}

  Let $(I,\leq)$ be a nonempty directed preorder, and let
  $\operatorname{Fin}(I)$ denote the set of finite subsets of $I$.  For
  $F\in\operatorname{Fin}(I)$ define
  \[
    \operatorname{UB}(F)
      =\{k\in I\mid (\forall i\in F)\ i\leq k\}.
  \]
  For $F=\varnothing$ one has
  \[
    \operatorname{UB}(F)=I\neq\varnothing.                       \tag{UB0}
  \]
  If $F\neq\varnothing$, choose an enumeration
  $F=\{i_1,\ldots,i_n\}$ with $n\geq1$.  Define
  \[
    k_1:=i_1,\qquad
    D_r:=\{k\in I\mid k_r\leq k\ \wedge\ i_{r+1}\leq k\}
      \quad(1\leq r<n).                                          \tag{UB1}
  \]
  Directedness is precisely
  \[
    D_r\neq\varnothing\quad(1\leq r<n).
  \]
  Recursively define
  \[
    k_{r+1}:=\operatorname{ch}(D_r)\in D_r
      \quad(1\leq r<n).                                          \tag{UB2}
  \]
  The defining membership in (UB2) and induction on $r$ are the chain
  \[
    (\forall j\leq r)\ i_j\leq k_r,\qquad
    k_r\leq k_{r+1},\quad i_{r+1}\leq k_{r+1}
      \Longrightarrow
    (\forall j\leq r+1)\ i_j\leq k_{r+1}.                       \tag{UB3}
  \]
  Starting from $i_1=k_1$, (UB3) at $r=n-1$ states
  \[
    k_n\in\operatorname{UB}(F),\qquad
    \operatorname{UB}(F)\neq\varnothing.                        \tag{UB4}
  \]
  Equations (UB0) and (UB4) define, by
  \Cref{def:nagata-choice-notation}, the function
  \[
    u_I\colon\operatorname{Fin}(I)\longrightarrow I,
    \qquad
    u_I(F):=\operatorname{ch}(\operatorname{UB}(F)),
    \qquad
    u_I(F)\in\operatorname{UB}(F).                               \tag{UB5}
  \]
  When the ambient preorder is fixed, write $u:=u_I$.

  \paragraph{Mathlib relation.}
  Existence of a common upper bound for a finite set is supplied by
  \texttt{\detokenize{Finset.exists_le}} (and, in directed form, by
  \texttt{\detokenize{Directed.finset_le}}).  The deterministic selected
  function $u_I$, including its value on the empty set, is the
  blueprint-local choice adapter (UB5); the proposed name
  \texttt{\detokenize{Order.Directed.finsetUpperBound}} is not an existing
  Mathlib declaration.
\end{definition}

\begin{lemma}[Property of the selected finite upper bound]
  \label{lemma:nagata-directed-finite-upper-bound}
  \uses{def:nagata-directed-finite-upper-bound}
  \lean{}

  Let \((I,\leq)\) be a nonempty directed preorder, let
  \(F\in\operatorname{Fin}(I)\), and let
  \(u_I\colon\operatorname{Fin}(I)\to I\) be the function (UB5) of
  \Cref{def:nagata-directed-finite-upper-bound}.  Then
  \[
    (\forall i\in F)\quad i\leq u_I(F).
  \]

  \paragraph{Mathlib relation.}
  The existential fact used to define $u_I(F)$ is exactly
  \texttt{\detokenize{Finset.exists_le}}.  This lemma records the defining
  property of the locally selected witness and is not a renamed Mathlib
  theorem.
\end{lemma}

\begin{proof}
  By \Cref{def:nagata-directed-finite-upper-bound},
  \[
    F=\varnothing\Longrightarrow
      (\forall i\in F)\ i\leq u_I(F)
  \]
  is the empty universal statement.  If $F\neq\varnothing$, then
  \[
    u_I(F)\in\operatorname{UB}(F)
      =\{k\in I\mid(\forall i\in F)\ i\leq k\}.
  \]
  In both cases,
  \[
    (\forall i\in F)\quad i\leq u_I(F).
  \]
\end{proof}

\begin{lemma}[A ring is the filtered colimit of its finite-type
  $\mathbf Z$-subalgebras]
  \label{lemma:nagata-ring-finite-type-subalgebras}
  \lean{}

  For a ring $A$, define the ordered set
  \[
    \mathcal A_A
      =\{B\subseteq A\mid B\text{ is a finitely generated }
        \mathbf Z\text{-subalgebra}\},
    \qquad
    B\leq C\Longleftrightarrow B\subseteq C,
  \]
  with transition homomorphism
  $\iota_{CB}\colon B\hookrightarrow C$ for $B\subseteq C$.  Define
  \[
    \Phi_A\colon\varinjlim_{B\in\mathcal A_A}B\longrightarrow A,
    \qquad
    \Phi_A([B,b])=b.
  \]
  Then $\mathcal A_A$ is directed and $\Phi_A$ is an isomorphism.  In
  particular, for every finite subset $E\subseteq A$, the algebra
  $\mathbf Z[E]\in\mathcal A_A$ contains $E$.
\end{lemma}

\begin{proof}
  If \([B,b]=[C,c]\) in the filtered colimit, its defining equivalence
  relation is the following predicate:
  \[
    \exists D\in\mathcal A_A:\quad
    B,C\subseteq D\ \wedge\
    \iota_{DB}(b)=\iota_{DC}(c).                                   \tag{67a}
  \]
  Because \(\iota_{DB},\iota_{DC}\), and \(D\hookrightarrow A\) are
  inclusions, (67a) is the equality chain
  \[
    b=\iota_{DB}(b)=\iota_{DC}(c)=c\quad\text{in }A.                \tag{67b}
  \]
  Thus \(\Phi_A([B,b]):=b\) is independent of the representative.
  Let $B=\mathbf Z[E_B]$ and $C=\mathbf Z[E_C]$ be elements of
  $\mathcal A_A$, where $E_B,E_C\subseteq A$ are finite.  Put
  \[
    D=\mathbf Z[E_B\cup E_C].
  \]
  Then $D\in\mathcal A_A$ and
  \[
    B\subseteq D,
    \qquad
    C\subseteq D,
  \]
  so $\mathcal A_A$ is directed.  The equality
  \[
    \Phi_A([\mathbf Z[a],a])=a\qquad(a\in A)
  \]
  proves that $\Phi_A$ is surjective.  If
  $\Phi_A([B,b])=\Phi_A([C,c])$, then $b=c$ in $A$.  With
  $D=\mathbf Z[E_B\cup E_C]\in\mathcal A_A$, the transition maps are
  inclusions, so
  \[
    \iota_{DB}(b)=b=c=\iota_{DC}(c)\quad\text{in }D.
  \]
  The defining equivalence relation of the filtered colimit is
  \[
    \bigl(\exists D\geq B,C,\ 
      \iota_{DB}(b)=\iota_{DC}(c)\bigr)
      \Longrightarrow [B,b]=[C,c].
  \]
  The witness $D=\mathbf Z[E_B\cup E_C]$ and the preceding equality prove
  $[B,b]=[C,c]$, so $\Phi_A$ is injective.  Finally,
  \[
    E\subseteq\mathbf Z[E]\subseteq A,
    \qquad
    \mathbf Z[E]\in\mathcal A_A
  \]
  for every finite $E\subseteq A$.
\end{proof}

\begin{definition}[Explicit limit datum]
  \label{def:nagata-limit-datum}
  \lean{}

  Let $J$ be a small category and let
  $F\colon J\to\mathbf{Sch}$ be a diagram.  Define
  \[
    \operatorname{LimitDatum}(F)
      :=\left\{(c,h)\ \middle|\
        c\in\operatorname{Cone}(F),\quad
        h\in\operatorname{IsLimit}(c)\right\}.                  \tag{MLL0}
  \]
  For $D=(c,h)\in\operatorname{LimitDatum}(F)$, write
  \[
  \begin{gathered}
    \lim(D):=c.\operatorname{pt},\qquad
    \pi_j^D:=c.\pi_j\colon\lim(D)\to F(j),\\
    \operatorname{lift}_D(s)\colon T\to\lim(D)
      \quad\text{for the lift of a cone }s\text{ with vertex }T.
  \end{gathered}                                                   \tag{MLL1}
  \]
  Henceforth every displayed object
  $\varprojlim_jF(j)$ and every phrase ``the limit projection'' in this
  chapter abbreviates a specified member $D$ of
  $\operatorname{LimitDatum}(F)$ and the notation (MLL1).  It never asserts
  an unqualified general limit-existence instance for schemes.

  \paragraph{Mathlib relation.}
  Formula (MLL0) is a notational wrapper around the exact declarations
  \texttt{\detokenize{CategoryTheory.Limits.Cone}} and
  \texttt{\detokenize{CategoryTheory.Limits.IsLimit}}.  The selected lift,
  factorization, and uniqueness operations are
  \texttt{\detokenize{Limits.IsLimit.lift}},
  \texttt{\detokenize{Limits.IsLimit.fac}}, and
  \texttt{\detokenize{Limits.IsLimit.uniq}}.  Mathlib's
  \texttt{\detokenize{AlgebraicGeometry.AffineTransitionLimit}} interface is
  conditional on such a cone and limit witness; it does not provide limits
  for arbitrary scheme diagrams.
\end{definition}

\begin{lemma}[Universal interface of explicit limit data]
  \label{lemma:nagata-limit-datum}
  \uses{def:nagata-limit-datum}

  If $D\in\operatorname{LimitDatum}(F)$ and $s$ is a cone on $F$ with vertex
  $T$, then
  \[
    \pi_j^D\circ\operatorname{lift}_D(s)=s.\pi_j
      \quad(j\in J),                                             \tag{MLL2}
  \]
  and any $u\colon T\to\lim(D)$ satisfying these equations equals
  $\operatorname{lift}_D(s)$.  Any two members $D,D'$ of
  $\operatorname{LimitDatum}(F)$ therefore determine a unique isomorphism
  \[
    \lim(D)\xrightarrow{\sim}\lim(D')                           \tag{MLL3}
  \]
  commuting with every projection.  All three assertions remain valid
  after reindexing whenever the reindexed cone is supplied with an
  $\operatorname{IsLimit}$ witness.

  \paragraph{Mathlib relation.}
  Equations (MLL2) and their uniqueness are
  \texttt{\detokenize{CategoryTheory.Limits.IsLimit.fac}} and
  \texttt{\detokenize{CategoryTheory.Limits.IsLimit.uniq}}; (MLL3) is
  \texttt{\detokenize{CategoryTheory.Limits.IsLimit.conePointUniqueUpToIso}}.
\end{lemma}

\begin{proof}
  Equations (MLL2) and their uniqueness are respectively the factorization
  and uniqueness fields of $h\in\operatorname{IsLimit}(c)$.  Apply the lift
  operation in both directions to two limit cones; the uniqueness field
  identifies the composites with the identities and proves (MLL3).  The
  reindexed assertion is the same argument for the supplied witness.
\end{proof}

\begin{lemma}[Existence of limits with affine transition morphisms]
  \label{lemma:nagata-affine-transition-limit-datum-exists}
  \uses{def:nagata-choice-notation,
    def:nagata-limit-datum}

  Let $J$ be a small cofiltered category and let
  $F\colon J\to\mathbf{Sch}$ be an inverse diagram whose transition
  morphisms are affine.  Then
  \[
    \operatorname{LimitDatum}(F)\neq\varnothing.                 \tag{MLL3a}
  \]
  Define the selected affine-transition limit datum and its notation by
  \[
    D_F^{\mathrm{aff}}
      :=\operatorname{ch}(\operatorname{LimitDatum}(F)),\qquad
    \varprojlim_{j\in J}F(j):=\lim(D_F^{\mathrm{aff}}),\qquad
    \pi_j:=\pi_j^{D_F^{\mathrm{aff}}}.                           \tag{MLL3b}
  \]
  Thus every later raw inverse-limit symbol for an affine-transition system
  has the explicit witness selected in (MLL3b), unless a different member of
  $\operatorname{LimitDatum}(F)$ is named.

  \paragraph{Mathlib relation.}
  The cone and witness produced here use the exact categorical declarations
  in \Cref{def:nagata-limit-datum}.  Mathlib's
  \texttt{\detokenize{Mathlib.AlgebraicGeometry.AffineTransitionLimit}}
  module provides results conditional on such a witness, but does not itself
  expose this general existence statement.  Consequently (MLL3a) is a
  blueprint-local construction theorem.
\end{lemma}

\begin{proof}
  This is the affine-transition inverse-limit construction
  \cite[\href{https://stacks.math.columbia.edu/tag/01ZM}{Tag 01ZM}]
  {stacks-project}.  On affine charts, take spectra of the filtered colimits
  of the contravariant coordinate-ring diagrams.  The localization
  comparisons identify the constructions on overlaps; scheme gluing gives a
  cone, and the affine universal properties plus morphism gluing prove that
  it is limiting.  This supplies (MLL3a), after which (MLL3b) is
  \Cref{def:nagata-choice-notation}.
\end{proof}

\begin{lemma}[Finite base change of an affine-transition limit datum]
  \label{lemma:nagata-affine-transition-limit-base-change}
  \uses{def:nagata-limit-datum,
    lemma:nagata-affine-transition-limit-datum-exists,
    lemma:nagata-limit-datum}

  Let $(X_\alpha,t_{\beta\alpha})$ be a cofiltered inverse system of qcqs
  schemes with affine transition morphisms, and let
  $D\in\operatorname{LimitDatum}(X_\bullet)$.  Any inverse system obtained
  levelwise from $X_\alpha$ by a finite succession of products over a fixed
  base and fibre products along fixed morphisms carries a canonical limit
  datum.  Its limit object is canonically isomorphic to the scheme obtained
  by applying the same finite-limit construction to $\lim(D)$, and the
  comparison commutes with every projection.\hfill\textup{(MLL4)}

  The assertion remains true after restriction to a cofinal full
  subcategory, provided the restricted cone is equipped with the induced
  $\operatorname{IsLimit}$ witness.

  \paragraph{Mathlib relation.}
  Mathlib's
  \texttt{\detokenize{Mathlib.AlgebraicGeometry.AffineTransitionLimit}}
  module supplies the conditional affine-transition comparison API once a
  cone and \texttt{\detokenize{CategoryTheory.Limits.IsLimit}} witness are
  given.  The finite iteration and reindexing package stated here is a
  blueprint-local lemma used for the systems in (FS4).
\end{lemma}

\begin{proof}
  Affine-locally on the fixed base, inverse limits with affine transitions
  are spectra of filtered colimits of rings.  Tensor products and quotient
  presentations give the displayed finite products and fibre products, and
  their universal properties commute with those filtered colimits.  The
  affine comparison isomorphisms commute with localization, so they glue.
  Projection compatibility and uniqueness are the corresponding
  \Cref{lemma:nagata-limit-datum} equations.  Cofinal reindexing preserves
  the same universal property.
\end{proof}

\begin{definition}[Absolute Noetherian approximation datum]
  \label{def:nagata-absolute-noetherian-approximation-datum}
  \uses{def:nagata-choice-notation,
    def:nagata-limit-datum,
    lemma:nagata-limit-datum}
  \lean{}

  Let \(S\) be a quasi-compact and quasi-separated scheme.  An
  \emph{absolute Noetherian approximation datum} for \(S\) is a tuple
  \[
    \mathsf A=
    \bigl(I,\leq,(S_\lambda)_{\lambda\in I},
      (t_{\mu\lambda})_{\lambda\leq\mu},
      (q_\lambda)_{\lambda\in I},
      L,(\pi_\lambda)_{\lambda\in I},\kappa,\eta\bigr)              \tag{68}
  \]
  with the explicitly typed entries
  \[
  \begin{gathered}
    I\text{ a set},\qquad {\leq}\ \subseteq I\times I,\qquad
    S_\lambda\text{ a scheme}\quad(\lambda\in I),\qquad
    L\text{ a scheme},\\
    \pi_\lambda\colon L\to S_\lambda,\quad
    \kappa\colon S\to L,\quad \eta\colon L\to S,\\
    (\forall\lambda\in I)\ \lambda\leq\lambda,\qquad
    (\forall\lambda\leq\mu\leq\nu)\ \lambda\leq\nu .
  \end{gathered}                                                     \tag{68a}
  \]
  Thus \((I,\leq)\) is a preorder.  Require the following additional
  symbolic conditions:
  \begin{align}
    &I\neq\varnothing,\qquad
    (\forall\lambda,\mu\in I)(\exists\nu\in I)\ 
      (\lambda\leq\nu\ \wedge\ \mu\leq\nu),                         \tag{69}\\
    &\operatorname{FiniteType}(s_{S_\lambda})
      \quad(\lambda\in I),                                          \tag{70}\\
    &t_{\mu\lambda}\colon S_\mu\to S_\lambda,\qquad
      \operatorname{IsAffine}(t_{\mu\lambda})
      \quad(\lambda\leq\mu),                                        \tag{71}\\
    &t_{\lambda\lambda}=\operatorname{id}_{S_\lambda},\qquad
    t_{\nu\lambda}=t_{\mu\lambda}\circ t_{\nu\mu}
      \quad(\lambda\leq\mu\leq\nu),                                 \tag{72}\\
    &q_\lambda\colon S\to S_\lambda,\qquad
    t_{\mu\lambda}\circ q_\mu=q_\lambda
      \quad(\lambda\leq\mu).                                        \tag{73}
  \end{align}
  Require the projection cone equalities
  \[
    t_{\mu\lambda}\circ\pi_\mu=\pi_\lambda
      \qquad(\lambda\leq\mu).                                      \tag{73a}
  \]
  For a scheme \(T\), define
  \[
    \mathcal C_{\mathsf A}(T):=
      \left\{(g_\lambda)_{\lambda\in I}\in
        \prod_{\lambda\in I}\operatorname{Hom}(T,S_\lambda)
        \ \middle|\
        (\forall\lambda\leq\mu)\
          t_{\mu\lambda}\circ g_\mu=g_\lambda\right\}.              \tag{74}
  \]
  The assertion that \((L,(\pi_\lambda)_\lambda)\) is the inverse limit
  of \((S_\lambda,t_{\mu\lambda})\) means that, for every scheme \(T\),
  the function
  \[
    \Lambda_{\mathsf A,T}\colon
      \operatorname{Hom}(T,L)\xrightarrow{\sim}
      \mathcal C_{\mathsf A}(T),\qquad
      h\longmapsto(\pi_\lambda\circ h)_{\lambda\in I}               \tag{75}
  \]
  is a bijection.  The remaining entries in (68) satisfy
  \begin{align}
    &\kappa:=\Lambda_{\mathsf A,S}^{-1}
      ((q_\lambda)_{\lambda\in I})\colon S\to L,\qquad
    \pi_\lambda\circ\kappa=q_\lambda,                               \tag{76}\\
    &\eta\colon L\to S,\qquad
    \eta\circ\kappa=\operatorname{id}_S,\qquad
    \kappa\circ\eta=\operatorname{id}_L.                            \tag{77}
  \end{align}
  Define
  \[
    \operatorname{ANApp}(S):=
      \{\mathsf A\text{ of the form (68) satisfying
        (68a), (69)--(73), (73a), and (74)--(77)}\}.                 \tag{78}
  \]

  \paragraph{Mathlib relation.}
  Conditions (74)--(75) are exactly the cone and
  \texttt{\detokenize{CategoryTheory.Limits.IsLimit}} component of
  \Cref{def:nagata-limit-datum}, written out on Hom-sets.  The directed
  Noetherian system, affine transitions, and comparison isomorphism are
  additional blueprint data; Mathlib does not currently package
  $\operatorname{ANApp}(S)$ or its existence theorem.
\end{definition}

\begin{lemma}[Existence of an absolute Noetherian approximation datum]
  \label{lemma:nagata-absolute-noetherian-approximation-exists}
  \uses{def:nagata-absolute-noetherian-approximation-datum,
    lemma:nagata-affine-transition-limit-datum-exists}
  \lean{}

  For every quasi-compact and quasi-separated scheme \(S\),
  \[
    \operatorname{ANApp}(S)\neq\varnothing.                         \tag{79}
  \]
\end{lemma}

\begin{proof}
  Proposition
  \cite[\href{https://stacks.math.columbia.edu/tag/01ZA}{Tag 01ZA}]
    {stacks-project}
  supplies entries for which the exact predicates are
  \[
  \begin{gathered}
    I\neq\varnothing,\qquad
    (\forall\lambda\in I)\ \lambda\leq\lambda,\qquad
    (\forall\lambda\leq\mu\leq\nu)\ \lambda\leq\nu,\\
    (\forall\lambda,\mu\in I)(\exists\nu\in I)\
      (\lambda\leq\nu\wedge\mu\leq\nu),\\
    \operatorname{FiniteType}(s_{S_\lambda})\quad(\lambda\in I),
    \qquad
    \operatorname{IsAffine}(t_{\mu\lambda})\quad(\lambda\leq\mu).
  \end{gathered}                                                     \tag{79a}
  \]
  The same proposition supplies the isomorphism
  \[
    \alpha_S\colon
      S\xrightarrow{\sim}
      L:=\varprojlim_{\lambda\in I}S_\lambda.                       \tag{80}
  \]
  Let \(\pi_\lambda\colon L\to S_\lambda\) be the limit projections and
  define
  \[
    q_\lambda:=\pi_\lambda\circ\alpha_S,\qquad
    \kappa:=\alpha_S,\qquad
    \eta:=\alpha_S^{-1}.                                            \tag{81}
  \]
  The remaining typed entries and the inverse-system identities are
  \begin{align}
    &L\text{ is a scheme},\qquad
      \pi_\lambda\colon L\to S_\lambda,\qquad
      \kappa\colon S\to L,\qquad\eta\colon L\to S,                  \tag{80a}\\
    &t_{\lambda\lambda}=\operatorname{id}_{S_\lambda},\qquad
    t_{\nu\lambda}=t_{\mu\lambda}\circ t_{\nu\mu},\qquad
    t_{\mu\lambda}\circ\pi_\mu=\pi_\lambda,                         \tag{82}\\
    &t_{\mu\lambda}\circ q_\mu
      =t_{\mu\lambda}\circ\pi_\mu\circ\alpha_S
      =\pi_\lambda\circ\alpha_S=q_\lambda,                          \tag{83}\\
    &\eta\circ\kappa
      =\alpha_S^{-1}\circ\alpha_S=\operatorname{id}_S,\qquad
    \kappa\circ\eta
      =\alpha_S\circ\alpha_S^{-1}=\operatorname{id}_L.              \tag{84}
  \end{align}
  The universal property of \(L=\varprojlim_\lambda S_\lambda\) is exactly
  the bijection (75).  Thus the tuple in (68), with entries (80)--(81),
  satisfies (69)--(77), and hence
  \[
    \mathsf A\in\operatorname{ANApp}(S).
  \]
  This membership is a witness for (79).
\end{proof}

\begin{definition}[Selected Noetherian inverse system]
  \label{def:nagata-noetherian-inverse-system}
  \uses{def:nagata-choice-notation,
    lemma:nagata-absolute-noetherian-approximation-exists}
  \lean{}

  Let \(S\) be a scheme satisfying
  \(\operatorname{QuasiCompact}(S)\wedge
    \operatorname{QuasiSeparated}(S)\).
  Define
  \[
  \begin{aligned}
    \mathsf A_S
      &:=\operatorname{ch}(\operatorname{ANApp}(S))\\
      &=\bigl(I_S,\leq_S,(S_\lambda)_{\lambda\in I_S},
          (t_{\mu\lambda})_{\lambda\leq_S\mu},
          (q_\lambda)_{\lambda\in I_S},
          L_S,(\pi_\lambda)_{\lambda\in I_S},\kappa_S,\eta_S\bigr).
  \end{aligned}                                                     \tag{85}
  \]
  The choice is defined because
  \(\operatorname{ANApp}(S)\neq\varnothing\) by
  \Cref{lemma:nagata-absolute-noetherian-approximation-exists}.
  Whenever \(\lambda,\mu\in I_S\), write
  \[
    \lambda\leq\mu\quad:\Longleftrightarrow\quad
    \lambda\leq_S\mu.                                               \tag{86}
  \]
  Define the inverse-system tuple
  \[
    \mathsf N(S):=
      \bigl(I_S,\leq_S,(S_\lambda)_{\lambda\in I_S},
        (t_{\mu\lambda})_{\lambda\leq\mu}\bigr).                    \tag{87}
  \]
  Membership \(\mathsf A_S\in\operatorname{ANApp}(S)\) expands to
  \begin{align}
    &{\leq}\ \subseteq I_S\times I_S,\qquad
      (\forall\lambda\in I_S)\ \lambda\leq\lambda,\qquad
      (\forall\lambda\leq\mu\leq\nu)\ \lambda\leq\nu,\qquad
      I_S\neq\varnothing,\qquad
    (\forall\lambda,\mu\in I_S)(\exists\nu\in I_S)\
      (\lambda\leq\nu\ \wedge\ \mu\leq\nu),                     \tag{88}\\
    &(\forall\lambda\in I_S)\ 
      \operatorname{FiniteType}(s_{S_\lambda}),\qquad
      (\forall\lambda,\mu\in I_S)\bigl(
        \lambda\leq\mu\Longrightarrow
        \operatorname{IsAffine}(t_{\mu\lambda})\bigr),             \tag{89}\\
    &t_{\lambda\lambda}=\operatorname{id}_{S_\lambda},\qquad
    t_{\nu\lambda}=t_{\mu\lambda}\circ t_{\nu\mu},\qquad
    t_{\mu\lambda}\circ q_\mu=q_\lambda.                            \tag{90}
  \end{align}
\end{definition}

\begin{definition}[Comparison from the original scheme to the inverse limit]
  \label{def:nagata-glued-limit-comparison}
  \uses{def:nagata-noetherian-inverse-system}
  \lean{}

  Let \(S\) satisfy the hypotheses of
  \Cref{def:nagata-noetherian-inverse-system}, and let
  \(\mathsf A_S\) be the tuple (85).  By (85), the limit object and its
  projections are
  \[
    L_S=\varprojlim_{\lambda\in I_S}S_\lambda,\qquad
    \pi_\lambda\colon L_S\to S_\lambda,\qquad
    t_{\mu\lambda}\circ\pi_\mu=\pi_\lambda.                         \tag{91}
  \]
  Define
  \[
    \kappa_S
      :=\Lambda_{\mathsf A_S,S}^{-1}
          ((q_\lambda)_{\lambda\in I_S})
      \colon S\longrightarrow L_S,\qquad
    \pi_\lambda\circ\kappa_S=q_\lambda.                             \tag{92}
  \]
  The selected inverse in (85) is
  \[
    \eta_S\colon L_S\to S,\qquad
    \eta_S\circ\kappa_S=\operatorname{id}_S,\qquad
    \kappa_S\circ\eta_S=\operatorname{id}_{L_S}.                    \tag{93}
  \]
\end{definition}

\begin{lemma}[Recovery of the selected inverse limit]
  \label{lemma:nagata-glued-limit-recovery}
  \uses{def:nagata-glued-limit-comparison}
  \lean{}

  Let $S$ be a quasi-compact and quasi-separated scheme, and let
  $\mathsf A_S$ and $\mathsf N(S)$ be the tuples of
  \Cref{def:nagata-noetherian-inverse-system}.  The morphism
  \[
    \kappa_S\colon S\xrightarrow{\sim}
      \varprojlim_{\lambda\in I_S}S_\lambda
  \]
  is an isomorphism with inverse \(\eta_S\).
\end{lemma}

\begin{proof}
  Equations (91)--(93) are
  \[
    \eta_S\circ\kappa_S=\operatorname{id}_S,\qquad
    \kappa_S\circ\eta_S
      =\operatorname{id}_{\varprojlim_{\lambda\in I_S}S_\lambda}.
  \]
  These two equations are the inverse criterion for
  \(\kappa_S\colon S\to\varprojlim_\lambda S_\lambda\).
\end{proof}

\begin{proposition}[{\cite[\href{https://stacks.math.columbia.edu/tag/01ZA}{Tag 01ZA}]{stacks-project}}]
  \label{prop:nagata-noetherian-inverse-system}
  \uses{def:nagata-noetherian-inverse-system,
    lemma:nagata-glued-limit-recovery}
  \lean{}

  Let $S$ be a quasi-compact and quasi-separated scheme, and let
  $\mathsf A_S$ and $\mathsf N(S)$ be the tuples of
  \Cref{def:nagata-noetherian-inverse-system}.  The system
  \(\mathsf N(S)\) satisfies
  \[
  \begin{gathered}
    \operatorname{Noetherian}(S_\lambda)
      \ \wedge\ \operatorname{FiniteType}(s_{S_\lambda})
      \quad(\lambda\in I_S),\\
    t_{\mu\lambda}\colon S_\mu\to S_\lambda,\qquad
      \operatorname{IsAffine}(t_{\mu\lambda})
      \quad(\lambda\leq\mu),\\
    \kappa_S\colon S\xrightarrow{\sim}
      \varprojlim_{\lambda\in I_S}S_\lambda.
  \end{gathered}                                                     \tag{94}
  \]
\end{proposition}

\begin{proof}
  Equation (89) contains the first finite-type assertion and the affineness
  assertion.  Since \(\mathbf Z\) is Noetherian, the finite-type
  preservation theorem for Noetherian schemes is the implication
  \[
    \operatorname{FiniteType}(s_{S_\lambda})
      \Longrightarrow\operatorname{Noetherian}(S_\lambda).         \tag{95}
  \]
  By \Cref{lemma:nagata-glued-limit-recovery}, the two inverse equalities in
  (93) are
  \[
    \eta_S\circ\kappa_S=\operatorname{id}_S,
    \qquad
    \kappa_S\circ\eta_S
      =\operatorname{id}_{\varprojlim_{\lambda\in I_S}S_\lambda},
  \]
  equivalently $\operatorname{IsIso}(\kappa_S)$.  Conjunction of (89), (95),
  and this predicate is exactly the three-line assertion (94).
\end{proof}
\subsection{Embedding a finite-type scheme into a finite-presentation model}

\begin{lemma}[Finite affine presentation subordinate to a locally finite-type
  morphism]
  \label{lemma:nagata-finite-affine-presentation-locally-finite-type}
  \lean{}

  Let $f\colon X\to S$ be locally of finite type and let $X$ be
  quasi-compact.  There exist an integer $r\geq1$ and, for
  $A_X=\{1,\ldots,r\}$, affine open immersions
  \[
    v_a\colon V_a=\operatorname{Spec}(B_a)\hookrightarrow X,
    \qquad
    \tau_a\colon T_a=\operatorname{Spec}(R_a)\hookrightarrow S,
  \]
  morphisms $f_a\colon V_a\to T_a$, integers $m_a\geq0$, and elements
  $b_{a,j}\in B_a$ such that, with $\varphi_a=f_a^\#$,
  \[
  \begin{gathered}
    X=\bigcup_{a\in A_X}v_a(V_a),\qquad
    \tau_a\circ f_a=f\circ v_a,\\
    B_a=\varphi_a(R_a)[b_{a,1},\ldots,b_{a,m_a}]
      \quad(a\in A_X).
  \end{gathered}
  \]
\end{lemma}

\begin{proof}
  If $X=\varnothing$, put
  \[
    r=1,\qquad
    V_1=X=\operatorname{Spec}(0),\qquad
    B_1=0,\qquad m_1=0.
  \]
  Choose an affine open $T_1=\operatorname{Spec}(R_1)\subseteq S$ when
  $S\neq\varnothing$, and put $T_1=S=\operatorname{Spec}(0)$ when
  $S=\varnothing$.  For the unique $f_1\colon V_1\to T_1$,
  \[
    X=V_1,\qquad
    \tau_1\circ f_1=f\circ v_1,\qquad
    B_1=f_1^\#(R_1),
  \]
  which is the asserted tuple.  Assume henceforth that $X\neq\varnothing$.
  For every $x\in X$, the affine-local definition of locally finite type
  supplies affine open immersions and a factorization
  \[
  \begin{gathered}
    v_x\colon V_x=\operatorname{Spec}(B_x)\hookrightarrow X,\qquad
    x\in v_x(V_x),\\
    \tau_x\colon T_x=\operatorname{Spec}(R_x)\hookrightarrow S,\qquad
    f_x\colon V_x\to T_x,\qquad
    \tau_x\circ f_x=f\circ v_x,\\
    E_x\subseteq B_x\text{ finite},\qquad
    B_x=f_x^\#(R_x)[E_x].
  \end{gathered}                                                   \tag{FAP0a}
  \]
  The equality
  \[
    X=\bigcup_{x\in X}v_x(V_x)
  \]
  and quasi-compactness of $X$ define points
  $x_1,\ldots,x_r\in X$ satisfying
  \[
    X=\bigcup_{a=1}^r v_{x_a}(V_{x_a}).                           \tag{80b}
  \]
  For $a\in A_X$, define
  \[
  \begin{gathered}
    (V_a,B_a,v_a,T_a,R_a,\tau_a,f_a)
      =(V_{x_a},B_{x_a},v_{x_a},T_{x_a},R_{x_a},\tau_{x_a},f_{x_a}),\\
    m_a=|E_{x_a}|,\qquad
    E_{x_a}=\{b_{a,1},\ldots,b_{a,m_a}\},\qquad
    \varphi_a=f_a^\#.
  \end{gathered}                                                   \tag{80c}
  \]
  Substitution of (80c) into (FAP0a)--(80b) is the list of equalities
  \[
  \begin{gathered}
    X=\bigcup_{a=1}^{r}v_a(V_a),
    \qquad
    \tau_a\circ f_a=f\circ v_a,\\
    E_{x_a}=\{b_{a,1},\ldots,b_{a,m_a}\},
    \qquad
    B_a=\varphi_a(R_a)[b_{a,1},\ldots,b_{a,m_a}]
      \quad(a\in A_X),
  \end{gathered}
  \]
  which is the conclusion of the lemma.
\end{proof}

\begin{definition}[Finite-presentation immersion model]
  \label{def:nagata-immersion-into-fp-model}
  \uses{def:nagata-choice-notation}
  \lean{}

  Let \(f\colon X\to S\) be locally of finite type, with \(X\) quasi-compact
  and quasi-separated.  Define the set
  \[
    \mathscr M_{\mathrm{imm}}(f):=
      \left\{(Y,q,i,e)\ \middle|\
      \begin{array}{l}
        Y\text{ is a scheme},\\
        q\colon Y\to S\text{ is of finite presentation},\\
        i\colon X\to Y\text{ is an immersion},\\
        e\colon q\circ i=f
      \end{array}\right\}.                                        \tag{IMM1}
  \]
  Lemma
  \cite[\href{https://stacks.math.columbia.edu/tag/01ZE}{Tag 01ZE}]
    {stacks-project}, applied to \(f\), is the existence assertion
  \[
    \mathscr M_{\mathrm{imm}}(f)\neq\varnothing.                   \tag{IMM2}
  \]
  By \Cref{def:nagata-choice-notation}, define
  \[
    \bigl(Y^{\mathrm{imm}}(X/S),q^{\mathrm{imm}},
      i^{\mathrm{imm}},e^{\mathrm{imm}}\bigr)
      :=\operatorname{ch}(\mathscr M_{\mathrm{imm}}(f)).           \tag{IMM3}
  \]
  Thus the defining membership in (IMM3) expands to
  \[
  \begin{gathered}
    q^{\mathrm{imm}}\colon Y^{\mathrm{imm}}(X/S)\to S
      \text{ is of finite presentation},\\
    i^{\mathrm{imm}}\colon X\to Y^{\mathrm{imm}}(X/S)
      \text{ is an immersion},\qquad
    e^{\mathrm{imm}}\colon q^{\mathrm{imm}}\circ i^{\mathrm{imm}}=f.
  \end{gathered}                                                     \tag{IMM4}
  \]
\end{definition}

\begin{lemma}[Properties of the finite-presentation immersion model,
  {\cite[\href{https://stacks.math.columbia.edu/tag/01ZE}{Tag 01ZE}]{stacks-project}}]
  \label{lemma:nagata-immersion-into-fp-model}
  \uses{def:nagata-immersion-into-fp-model}
  \lean{}

  The selected tuple in
  \Cref{def:nagata-immersion-into-fp-model} satisfies
  \[
  \begin{gathered}
    q^{\mathrm{imm}}\colon Y^{\mathrm{imm}}(X/S)\to S
      \text{ is of finite presentation},\\
    i^{\mathrm{imm}}\colon X\to Y^{\mathrm{imm}}(X/S)
      \text{ is an immersion},\qquad
    q^{\mathrm{imm}}\circ i^{\mathrm{imm}}=f.
  \end{gathered}
  \]
\end{lemma}

\begin{proof}
  Equation (IMM3) and
  \(\operatorname{ch}(\mathscr M_{\mathrm{imm}}(f))
    \in\mathscr M_{\mathrm{imm}}(f)\)
  are the membership
  \[
    (Y^{\mathrm{imm}}(X/S),q^{\mathrm{imm}},
      i^{\mathrm{imm}},e^{\mathrm{imm}})
      \in\mathscr M_{\mathrm{imm}}(f).
  \]
  Expanding the predicate in (IMM1) is exactly the three assertions in (IMM4).
\end{proof}
\begin{definition}[Quasi-compact open on which an immersion is closed]
  \label{def:nagata-immersion-shrink-to-closed}
  \uses{def:nagata-choice-notation}
  \lean{}

  Let \(i\colon X\to Y\) be an immersion, with \(X\) quasi-compact and
  \(Y\) quasi-separated.  Define
  \[
    \mathscr N_{\mathrm{cl}}(i):=
      \left\{(W,w,i_W,e_W)\ \middle|\
      \begin{array}{l}
        W\text{ is a quasi-compact scheme},\\
        w\colon W\hookrightarrow Y\text{ is a quasi-compact open immersion},\\
        i_W\colon X\hookrightarrow W\text{ is a closed immersion},\\
        e_W\colon w\circ i_W=i
      \end{array}\right\}.                                        \tag{105}
  \]
  We now define a member of this set.  If \(X=\varnothing\), put
  \[
    (W_0,w_0,i_0,e_0):=
      (\varnothing,\varnothing\hookrightarrow Y,
        \operatorname{id}_{\varnothing},
        (\varnothing\hookrightarrow Y)\circ
          \operatorname{id}_{\varnothing}=i).                     \tag{106}
  \]
  Assume \(X\neq\varnothing\).  For every \(x\in X\), the definition of an
  immersion supplies an affine open immersion
  \[
    w_x\colon W_x\hookrightarrow Y,\qquad i(x)\in w_x(W_x),
    \qquad
    i_x\colon X\times_YW_x\hookrightarrow W_x
      \text{ a closed immersion}.                                 \tag{107}
  \]
  Quasi-compactness of \(X\) supplies a finite nonempty set \(A\subseteq X\)
  satisfying
  \[
    X=\bigcup_{x\in A}i^{-1}(w_x(W_x)).                            \tag{108}
  \]
  Define
  \[
    W_0:=\bigcup_{x\in A}w_x(W_x)\subseteq Y,\qquad
    w_0\colon W_0\hookrightarrow Y,\qquad
    i_0\colon X\to W_0,\qquad
    w_0\circ i_0=i.                                                \tag{109}
  \]
  The last morphism is the unique factorization through \(w_0\), whose
  existence is equivalent to (108).  Since \(A\) is finite and every \(W_x\)
  is affine,
  \[
    W_0=\bigcup_{x\in A}W_x\quad\Longrightarrow\quad
    W_0\text{ is quasi-compact}.                                  \tag{110}
  \]
  For a quasi-compact open \(V\subseteq Y\), quasi-separatedness of \(Y\)
  and finiteness of \(A\) are the two hypotheses in the chain
  \[
    w_0^{-1}(V)=W_0\cap V
      =\bigcup_{x\in A}(W_x\cap V),\qquad
    W_x\cap V\text{ is quasi-compact},                             \tag{111}
  \]
  Consequently the defining quantifier in
  \Cref{def:nagata-morphism-predicate-notation} is
  \[
    \forall V\subseteq Y\text{ quasi-compact open},\quad
      w_0^{-1}(V)\text{ is quasi-compact},
    \qquad\text{i.e.}\qquad
    \operatorname{IsQuasiCompact}(w_0).
  \]
  The cartesian restrictions are
  \[
    i_0\times_{W_0}W_x
      =i_x\colon X\times_YW_x\hookrightarrow W_x\quad(x\in A).    \tag{112}
  \]
  Target-locality of closed immersions on the cover
  \(W_0=\bigcup_{x\in A}W_x\), applied to (112), is the predicate
  \[
    \operatorname{IsClosedImmersion}(i_0\colon X\to W_0).          \tag{113}
  \]
  Equations (106), respectively (109)--(113), prove
  \[
    (W_0,w_0,i_0,e_0)\in\mathscr N_{\mathrm{cl}}(i),\qquad
    \mathscr N_{\mathrm{cl}}(i)\neq\varnothing,                    \tag{114}
  \]
  where in the nonempty case
  \(e_0\) denotes the equality \(w_0\circ i_0=i\) in (109).
  By \Cref{def:nagata-choice-notation}, define
  \[
    \bigl(Y^{\mathrm{cl}}(i),w_i^{\mathrm{cl}},
      i^{\mathrm{cl}},e_i^{\mathrm{cl}}\bigr)
      :=\operatorname{ch}(\mathscr N_{\mathrm{cl}}(i)).            \tag{115}
  \]
\end{definition}

\begin{lemma}[Properties of the closed-immersion neighbourhood,
  {\cite[\href{https://stacks.math.columbia.edu/tag/01ZG}{Tag 01ZG}]{stacks-project}}]
  \label{lemma:nagata-immersion-shrink-to-closed}
  \uses{def:nagata-immersion-shrink-to-closed}
  \lean{}

  The tuple selected in
  \Cref{def:nagata-immersion-shrink-to-closed} satisfies
  \[
  \begin{gathered}
    Y^{\mathrm{cl}}(i)\text{ is quasi-compact},\\
    w_i^{\mathrm{cl}}\colon Y^{\mathrm{cl}}(i)\hookrightarrow Y
      \text{ is a quasi-compact open immersion},\\
    i^{\mathrm{cl}}\colon X\hookrightarrow Y^{\mathrm{cl}}(i)
      \text{ is a closed immersion},\qquad
    w_i^{\mathrm{cl}}\circ i^{\mathrm{cl}}=i.
  \end{gathered}
  \]
\end{lemma}

\begin{proof}
  The defining property of the choice function and (115) are
  \[
    (Y^{\mathrm{cl}}(i),w_i^{\mathrm{cl}},
      i^{\mathrm{cl}},e_i^{\mathrm{cl}})
      \in\mathscr N_{\mathrm{cl}}(i).
  \]
  Expanding the four predicates in (105) is exactly the four assertions in
  the statement.
\end{proof}
\begin{definition}[Finite-presentation closed-subscheme system]
  \label{def:nagata-closed-subscheme-fp-system}
  \uses{def:nagata-direct-image-structure-map,
    def:nagata-limit-datum,
    def:nagata-presented-relative-spec,
    lemma:nagata-affine-transition-limit-datum-exists,
    lemma:nagata-closed-immersion-relative-spec-recovery,
    lemma:nagata-limit-datum}
  \lean{}

  Let $i\colon X\hookrightarrow Y$ be a closed immersion with $Y$ qcqs.
  By \Cref{def:nagata-direct-image-structure-map}, define
  \[
  \begin{gathered}
    \pi_i:=\varepsilon_i:=i^\flat\colon
      \mathcal O_Y\twoheadrightarrow i_*\mathcal O_X,
    \qquad
    \mathcal I:=\ker(\pi_i),
    \qquad
    q_{\mathcal I}\colon\mathcal O_Y\twoheadrightarrow
      \mathcal O_Y/\mathcal I,\\
    \overline\pi_i\colon\mathcal O_Y/\mathcal I
      \xrightarrow{\sim}i_*\mathcal O_X,
    \qquad
    \overline\pi_i\circ q_{\mathcal I}=\pi_i.                 
  \end{gathered} \tag{CSFP0}
  \]
  Let
  \[
    \rho_i\colon X\xrightarrow{\sim}
      \operatorname{Spec}_Y
        (i_*\mathcal O_X;\mathcal I,\overline\pi_i)
  \]
  be the $Y$-isomorphism of
  \Cref{lemma:nagata-closed-immersion-relative-spec-recovery}; equivalently,
  \[
    \operatorname{Spec}_Y(\pi_i)\circ\rho_i=i.
  \]
  Define
  \[
  \begin{gathered}
    X_{\mathcal I}:=\operatorname{Spec}_Y(\mathcal O_Y/\mathcal I),
    \qquad
    k_{\mathcal I}:=\operatorname{Spec}_Y(q_{\mathcal I})\colon
      X_{\mathcal I}\hookrightarrow Y,\\
    \sigma_i:=\operatorname{Spec}_Y(\overline\pi_i)\circ\rho_i
      \colon X\xrightarrow{\sim}X_{\mathcal I},\\
    k_{\mathcal I}\circ\sigma_i
      =\operatorname{Spec}_Y(q_{\mathcal I})
         \circ\operatorname{Spec}_Y(\overline\pi_i)\circ\rho_i
      =\operatorname{Spec}_Y(\overline\pi_i\circ q_{\mathcal I})
         \circ\rho_i
      =\operatorname{Spec}_Y(\pi_i)\circ\rho_i=i.              
  \end{gathered} \tag{115a}
  \]
  Define
  \[
    A_{X/Y}:=\left\{(\mathcal J,\iota_{\mathcal J})\ \middle|\
      \begin{array}{l}
        \mathcal J\subseteq\mathcal I\text{ is a finite-type
          quasi-coherent ideal of }\mathcal O_Y,\\
        \iota_{\mathcal J}\colon\mathcal J\hookrightarrow\mathcal I
          \text{ is the inclusion}
      \end{array}\right\}.                                      \tag{116}
  \]
  For
  $\alpha=(\mathcal J,\iota_{\mathcal J})\in A_{X/Y}$, define
  \[
    \mathcal I_\alpha:=\mathcal J,
    \qquad
    \iota_\alpha:=\iota_{\mathcal J}.
  \]
  For $\alpha,\beta\in A_{X/Y}$, define
  \[
  \begin{aligned}
    \alpha\leq\beta
    &\Longleftrightarrow \mathcal I_\alpha\subseteq\mathcal I_\beta\\
    &\Longleftrightarrow
      \bigl(\exists\,(\jmath_{\beta\alpha}\colon
        \mathcal I_\alpha\hookrightarrow\mathcal I_\beta)\bigr)
      \bigl[
        (\jmath_{\beta\alpha}\text{ is the inclusion})
        \ \wedge\
        (\iota_\beta\circ\jmath_{\beta\alpha}=\iota_\alpha)
      \bigr],
  \end{aligned}                                                   \tag{116a}
  \]
  For $\alpha\leq\beta$, define the canonical inclusion occurring in
  (116a) by
  \[
    \jmath_{\beta\alpha}\colon
      \mathcal I_\alpha\hookrightarrow\mathcal I_\beta,
    \qquad
    (\jmath_{\beta\alpha})_V(s):=s
      \quad(V\subseteq Y\text{ open},\
        s\in\mathcal I_\alpha(V)),
    \qquad
    \iota_\beta\circ\jmath_{\beta\alpha}=\iota_\alpha.
  \]
  Define
  \[
    0_{X/Y}:=(0,\iota_0)\in A_{X/Y},
    \qquad
    \iota_0\colon0\hookrightarrow\mathcal I.
  \]
  For $\alpha,\beta\in A_{X/Y}$, define the sum and its inclusion by
  \[
  \begin{gathered}
    \mathcal I_\alpha+\mathcal I_\beta
      :=\operatorname{im}\!\left(
        \mathcal I_\alpha\oplus\mathcal I_\beta
          \xrightarrow{\,\iota_\alpha+\iota_\beta\,}\mathcal I\right),\\
    \iota_{\alpha\beta}\colon
      \mathcal I_\alpha+\mathcal I_\beta\hookrightarrow\mathcal I,
    \qquad
    \gamma(\alpha,\beta):=
      (\mathcal I_\alpha+\mathcal I_\beta,
        \iota_{\alpha\beta}).                                  
  \end{gathered} \tag{116b}
  \]
  On an affine open $W\subseteq Y$, if
  \[
    \Gamma(W,\mathcal I_\alpha)=(a_1,\ldots,a_m),
    \qquad
    \Gamma(W,\mathcal I_\beta)=(b_1,\ldots,b_n),
  \]
  then the quasi-coherent-ideal sum identity is
  \[
    \Gamma(W,\mathcal I_\alpha+\mathcal I_\beta)
      =\Gamma(W,\mathcal I_\alpha)+\Gamma(W,\mathcal I_\beta)
      =(a_1,\ldots,a_m,b_1,\ldots,b_n).
  \]
  Hence
  \[
    \gamma(\alpha,\beta)\in A_{X/Y},
    \qquad
    \alpha\leq\gamma(\alpha,\beta),
    \qquad
    \beta\leq\gamma(\alpha,\beta),                            \tag{116c}
  \]
  so (116), (116b), and (116c) verify nonemptiness and directedness.

  For $\alpha\in A_{X/Y}$, let
  \[
    \pi_\alpha^Y\colon\mathcal O_Y\twoheadrightarrow
      \mathcal O_Y/\mathcal I_\alpha
  \]
  be the quotient.  For $\alpha\in A_{X/Y}$, define
  $\sigma_{\infty\alpha}$ uniquely by the first line below; for
  $\alpha\leq\beta$, define $\sigma_{\beta\alpha}$ uniquely by the second:
  \[
  \begin{gathered}
    \sigma_{\infty\alpha}\colon
      \mathcal O_Y/\mathcal I_\alpha
        \twoheadrightarrow\mathcal O_Y/\mathcal I,
    \qquad
    \sigma_{\infty\alpha}\circ\pi_\alpha^Y=q_{\mathcal I},\\
    \sigma_{\beta\alpha}\colon
      \mathcal O_Y/\mathcal I_\alpha
        \twoheadrightarrow\mathcal O_Y/\mathcal I_\beta,
    \qquad
    \sigma_{\beta\alpha}\circ\pi_\alpha^Y=\pi_\beta^Y.
  \end{gathered} \tag{117}
  \]
  For every $\alpha$, the morphism $\pi_\alpha^Y$ is an epimorphism.  The
  following computations and epimorphism cancellation
  \[
  \begin{aligned}
    (\sigma_{\infty\beta}\circ\sigma_{\beta\alpha})
      \circ\pi_\alpha^Y
      &=\sigma_{\infty\beta}\circ\pi_\beta^Y
       =q_{\mathcal I}
       =\sigma_{\infty\alpha}\circ\pi_\alpha^Y,\\
    (\sigma_{\gamma\beta}\circ\sigma_{\beta\alpha})
      \circ\pi_\alpha^Y
      &=\sigma_{\gamma\beta}\circ\pi_\beta^Y
       =\pi_\gamma^Y
       =\sigma_{\gamma\alpha}\circ\pi_\alpha^Y
  \end{aligned}
  \]
  are equivalent to the equalities
  \[
    \sigma_{\infty\beta}\circ\sigma_{\beta\alpha}
      =\sigma_{\infty\alpha},
    \qquad
    \sigma_{\gamma\beta}\circ\sigma_{\beta\alpha}
      =\sigma_{\gamma\alpha}
      \quad(\alpha\leq\beta\leq\gamma).                       \tag{117a}
  \]
  Define
  \[
    X_\alpha:=\operatorname{Spec}_Y
      (\mathcal O_Y/\mathcal I_\alpha),
    \qquad
    k_\alpha:=\operatorname{Spec}_Y(\pi_\alpha^Y)\colon
      X_\alpha\hookrightarrow Y,                                \tag{118}
  \]
  and, for $\alpha\leq\beta$, define
  \[
  \begin{gathered}
    t_{\beta\alpha}:=\operatorname{Spec}_Y(\sigma_{\beta\alpha})
      \colon X_\beta\hookrightarrow X_\alpha,\\
    i_\alpha:=\operatorname{Spec}_Y(\sigma_{\infty\alpha})
      \circ\sigma_i\colon X\hookrightarrow X_\alpha.          
  \end{gathered} \tag{119}
  \]
  Contravariant functoriality of $\operatorname{Spec}_Y$ and
  (115a), (117), and (117a) are the following explicit equality chains:
  \[
  \begin{aligned}
    k_\alpha\circ t_{\beta\alpha}
      &=\operatorname{Spec}_Y(\pi_\alpha^Y)
        \circ\operatorname{Spec}_Y(\sigma_{\beta\alpha})
       =\operatorname{Spec}_Y(
          \sigma_{\beta\alpha}\circ\pi_\alpha^Y)=k_\beta,\\
    t_{\beta\alpha}\circ i_\beta
      &=\operatorname{Spec}_Y(\sigma_{\beta\alpha})
        \circ\operatorname{Spec}_Y(\sigma_{\infty\beta})
        \circ\sigma_i
       =\operatorname{Spec}_Y(\sigma_{\infty\alpha})
        \circ\sigma_i=i_\alpha,\\
    k_\alpha\circ i_\alpha
      &=\operatorname{Spec}_Y(\pi_\alpha^Y)
        \circ\operatorname{Spec}_Y(\sigma_{\infty\alpha})
        \circ\sigma_i
       =\operatorname{Spec}_Y(q_{\mathcal I})\circ\sigma_i
       =k_{\mathcal I}\circ\sigma_i=i,\\
    t_{\beta\alpha}\circ t_{\gamma\beta}
      &=\operatorname{Spec}_Y(\sigma_{\beta\alpha})
        \circ\operatorname{Spec}_Y(\sigma_{\gamma\beta})
       =\operatorname{Spec}_Y(\sigma_{\gamma\alpha})
       =t_{\gamma\alpha}.
  \end{aligned}                                                   \tag{120}
  \]
  If
  $\pi_\alpha\colon\varprojlim_{\gamma\in A_{X/Y}}X_\gamma\to X_\alpha$
  is the limit projection, define
  \[
    \kappa_{X/Y}\colon X\longrightarrow
      \varprojlim_{\alpha\in A_{X/Y}}X_\alpha
  \]
  to be the unique morphism satisfying
  \[
    \pi_\alpha\circ\kappa_{X/Y}=i_\alpha,
    \qquad
    t_{\beta\alpha}\circ i_\beta=i_\alpha
      \quad(\alpha\leq\beta).                                  \tag{121}
  \]
  Denote this specified inverse system by
  \[
    \mathsf C_{\mathrm{fp}}(X/Y)
      =\bigl(A_{X/Y},(X_\alpha,k_\alpha)_\alpha,
        (t_{\beta\alpha})_{\alpha\leq\beta},(i_\alpha)_\alpha,
        \kappa_{X/Y}\bigr).                                     \tag{122}
  \]

  \paragraph{Mathlib relation.}
  The displayed inverse limit is interpreted by
  (MLL3b) of
  \Cref{lemma:nagata-affine-transition-limit-datum-exists}; its cone and
  \texttt{\detokenize{CategoryTheory.Limits.IsLimit}} witness are the
  globally selected datum $D_{F_{X/Y}}^{\mathrm{aff}}$.  The affine chart
  description is proved below.  No global Mathlib
  \texttt{\detokenize{HasLimit}} instance for arbitrary scheme diagrams is
  assumed.
\end{definition}

\begin{lemma}[Localization commutes with a filtered colimit of rings]
  \label{lemma:nagata-localization-filtered-colimit}
  \lean{}

  Let $(A_\alpha,\varphi_{\beta\alpha})$ be a filtered direct system of
  commutative rings.  Let $s_\alpha\in A_\alpha$ be compatible elements,
  so $\varphi_{\beta\alpha}(s_\alpha)=s_\beta$, and write
  \[
    A:=\varinjlim_\alpha A_\alpha,
    \qquad s:=[\alpha,s_\alpha]\in A.
  \]
  There is a natural ring isomorphism
  \[
    \delta_s\colon
      A_s=\left(\varinjlim_\alpha A_\alpha\right)_s
        \xrightarrow{\sim}
      \varinjlim_\alpha(A_\alpha)_{s_\alpha},                   \tag{MLFC0}
  \]
  characterized on representatives by
  \[
    \delta_s\!\left(\frac{[\alpha,a]}{s^n}\right)
      =\left[\alpha,\frac{a}{s_\alpha^n}\right].                \tag{MLFC1}
  \]
  It is compatible with every morphism of filtered systems carrying
  $s_\alpha$ to the corresponding selected element.

  \paragraph{Mathlib relation.}
  The two universal constructions are the exact Mathlib APIs
  \texttt{\detokenize{CategoryTheory.Limits.colimit}} and
  \texttt{\detokenize{Localization.Away}}.  The natural isomorphism
  (MLFC0), in the representative-level form required below, is a
  blueprint-local lemma; it is not claimed to be a single existing Mathlib
  declaration.
\end{lemma}

\begin{proof}
  The compatible maps
  $A_\alpha\to\varinjlim_\gamma(A_\gamma)_{s_\gamma}$ send
  $s_\alpha$ to a unit.  The colimit universal property first gives a map
  from $A$, and the localization universal property extends it uniquely to
  the forward map in (MLFC0).  Conversely, the maps
  $(A_\alpha)_{s_\alpha}\to A_s$ induced by $A_\alpha\to A$ are compatible,
  so the colimit universal property gives the inverse.  Formula (MLFC1)
  follows from the two defining factorization equations.  Those equations
  also show that both composites fix the canonical maps, hence are
  identities by uniqueness.  The same uniqueness proves naturality.
\end{proof}

\begin{lemma}[Limit datum for the finite-presentation closed-subscheme system]
  \label{lemma:nagata-closed-subscheme-fp-limit-datum}
  \uses{def:nagata-closed-subscheme-fp-system,
    def:nagata-limit-datum,
    lemma:nagata-affine-transition-limit-datum-exists,
    lemma:nagata-limit-datum,
    lemma:nagata-localization-filtered-colimit}

  For the inverse system
  $\mathsf C_{\mathrm{fp}}(X/Y)$ of
  \Cref{def:nagata-closed-subscheme-fp-system}, let
  $F_{X/Y}$ be the functor sending $\alpha$ to $X_\alpha$ and an inequality
  $\alpha\leq\beta$ to $t_{\beta\alpha}$.  Then
  \[
    \operatorname{LimitDatum}(F_{X/Y})\neq\varnothing.           \tag{MLCSL0}
  \]
  Put
  $D_{X/Y}:=D_{F_{X/Y}}^{\mathrm{aff}}$, the globally selected datum of
  (MLL3b).  For every affine open
  $U=\operatorname{Spec}(A)\subseteq Y$, its restriction is
  \[
    \lim(D_{X/Y})\times_YU
      \cong
    \operatorname{Spec}\!\left(
      \varinjlim_\alpha A/I_{\alpha,U}\right),                  \tag{MLCSL1}
  \]
  and its projections are induced by the canonical maps
  $A/I_{\alpha,U}\to\varinjlim_\beta A/I_{\beta,U}$.

  \paragraph{Mathlib relation.}
  The result produces the explicit
  \texttt{\detokenize{CategoryTheory.Limits.Cone}} and
  \texttt{\detokenize{CategoryTheory.Limits.IsLimit}} required by
  \Cref{def:nagata-limit-datum}.  It is a project lemma for this particular
  affine-transition system, not a general Mathlib limit instance.
\end{lemma}

\begin{proof}
  Nonemptiness and the selected datum follow first from
  \Cref{lemma:nagata-affine-transition-limit-datum-exists}, since the
  transitions $t_{\beta\alpha}$ are affine.  It remains to record the affine
  model of the selected datum.
  On an affine open $U=\operatorname{Spec}(A)$ the transition morphisms are
  affine, so contravariance of affine spectra turns the system into the
  filtered direct system of quotient rings $A/I_{\alpha,U}$.  Its colimit
  gives (MLCSL1) and the affine cone.  On a principal subopen $D(s)$,
  \Cref{lemma:nagata-localization-filtered-colimit} identifies the
  localization of that colimit with the colimit of the localized quotient
  rings.  Hence the affine cones agree on overlaps and glue to a scheme cone
  $D'$.
  The affine-spectrum/filtered-colimit universal property proves
  $\operatorname{IsLimit}$ locally; morphisms into schemes glue uniquely, so
  the global cone has the same universal property.  The canonical comparison
  isomorphism (MLL3) from $D_{X/Y}$ to $D'$ transports the affine chart
  description to the globally selected datum and proves (MLCSL1).
\end{proof}

\begin{lemma}[Properties of the closed-subscheme system,
  {\cite[\href{https://stacks.math.columbia.edu/tag/09ZP}{Tag 09ZP}]{stacks-project}}]
  \label{lemma:nagata-closed-subscheme-fp-system}
  \uses{def:nagata-closed-subscheme-fp-system,
    lemma:nagata-closed-subscheme-fp-limit-datum,
    lemma:nagata-localization-filtered-colimit}
  \lean{}

  Every $X_\alpha\hookrightarrow Y$ is of finite presentation, all
  transitions are closed immersions, and the morphism
  $\kappa_{X/Y}$ of \Cref{def:nagata-closed-subscheme-fp-system} is an
  isomorphism.
\end{lemma}

\begin{proof}
  By \cite[\href{https://stacks.math.columbia.edu/tag/09ZP}{Tag 09ZP}]
  {stacks-project},
  \[
    \mathcal I=\varinjlim_{\alpha\in A_{X/Y}}\mathcal I_\alpha
      =\bigcup_{\alpha\in A_{X/Y}}\mathcal I_\alpha,
    \qquad
    (A_{X/Y},\leq)\quad\text{is directed}.                     \tag{123}
  \]
  Put $D_{X/Y}:=D_{F_{X/Y}}^{\mathrm{aff}}$ as in
  \Cref{lemma:nagata-closed-subscheme-fp-limit-datum}, and define
  \[
    L:=\lim(D_{X/Y})=\varprojlim_{\alpha\in A_{X/Y}}X_\alpha,\qquad
    \pi_\alpha:=\pi_\alpha^{D_{X/Y}}\colon L\to X_\alpha,\qquad
    k_L:=k_{0_{X/Y}}\circ\pi_{0_{X/Y}}\colon L\to Y.             \tag{123c}
  \]
  The inequalities $0_{X/Y}\leq\alpha$ for
  $\alpha\in A_{X/Y}$, equations (120), and the limit projection identity
  are the following equality chain:
  \[
    k_L
      =k_{0_{X/Y}}\circ\pi_{0_{X/Y}}
      =k_{0_{X/Y}}\circ t_{\alpha,0_{X/Y}}\circ\pi_\alpha
      =k_\alpha\circ\pi_\alpha.                                  \tag{123d}
  \]
  For every $\alpha$, equations (121), (120), and (123d) are the following
  equality chain:
  \[
  \begin{aligned}
    k_L\circ\kappa_{X/Y}
      &=k_\alpha\circ\pi_\alpha\circ\kappa_{X/Y}\\
      &=k_\alpha\circ i_\alpha\\
      &=i.
  \end{aligned}                                                   \tag{123e}
  \]
  Fix an affine open immersion $j_U\colon U=\operatorname{Spec}(A)\hookrightarrow
  Y$, and define the ideals
  \[
    I_U:=\Gamma(U,\mathcal I|_U),
    \qquad
    I_{\alpha,U}:=\Gamma(U,\mathcal I_\alpha|_U)
      \quad(\alpha\in A_{X/Y}).                                  \tag{123a}
  \]
  The restriction of $\pi_\alpha^Y$ in (117) and the affine
  finite-presentation criterion are the following predicate chain:
  \[
  \begin{gathered}
    \Gamma\!\left(U,
      (\mathcal O_Y/\mathcal I_\alpha)|_U\right)=A/I_{\alpha,U},\\
    \operatorname{Spec}(A/I_{\alpha,U})\to\operatorname{Spec}(A)
      \text{ is of finite presentation}
      \Longleftrightarrow
    I_{\alpha,U}\text{ is finitely generated}.                  
  \end{gathered} \tag{123b}
  \]
  Equation (116), the affine criterion (123b), and target-locality of finite
  presentation are the following predicate chain:
  \[
  \begin{aligned}
    \operatorname{FiniteTypeModule}_Y(\mathcal I_\alpha)
      &\Longrightarrow
        (\forall U\subseteq Y\text{ affine})\ 
        \operatorname{FinitePresentation}
          (k_\alpha|_{X_\alpha\times_YU})\\
      &\Longrightarrow
        \operatorname{FinitePresentation}(k_\alpha),\\
    k_\alpha&\colon
      X_\alpha=\operatorname{Spec}_Y
        (\mathcal O_Y/\mathcal I_\alpha)\hookrightarrow Y.
  \end{aligned}                                                   \tag{124}
  \]
  If $\alpha\leq\beta$, (117) states that
  $\sigma_{\beta\alpha}$ is surjective.  The affine criterion
  \[
    \operatorname{Spec}_Y(\sigma_{\beta\alpha})
      \text{ is a closed immersion}
      \Longleftrightarrow
    \sigma_{\beta\alpha}\text{ is surjective}
  \]
  identifies
  \[
    t_{\beta\alpha}
      =\operatorname{Spec}_Y(\sigma_{\beta\alpha})\colon
        X_\beta\hookrightarrow X_\alpha
      \quad\text{is a closed immersion}.                       \tag{125}
  \]

  Restriction of (123) to the affine open $U$ and the equivalence between
  quasi-coherent $\mathcal O_U$-modules and $A$-modules are the following
  equality chain:
  \[
    I_U=\varinjlim_{\alpha\in A_{X/Y}}I_{\alpha,U}
       =\bigcup_{\alpha\in A_{X/Y}}I_{\alpha,U}.                \tag{125a}
  \]
  Define
  \[
    \chi_U\colon
      \varinjlim_{\alpha\in A_{X/Y}}A/I_{\alpha,U}\longrightarrow A/I_U,
    \qquad
    \chi_U([\alpha,a+I_{\alpha,U}]):=a+I_U.                    \tag{126}
  \]
  If
  $a+I_{\alpha,U}=a'+I_{\alpha,U}$, then
  $a-a'\in I_{\alpha,U}\subseteq I_U$, and hence
  \[
    a+I_U=a'+I_U.
  \]
  If $\alpha\leq\beta$, then
  \[
    \chi_U([\beta,a+I_{\beta,U}])=a+I_U
      =\chi_U([\alpha,a+I_{\alpha,U}]),
  \]
  so (126) is a well-defined homomorphism from the indicated colimit.  Since
  $I_{0_{X/Y},U}=0$, for every $a+I_U\in A/I_U$ one has
  \[
    \chi_U([0_{X/Y},a+I_{0_{X/Y},U}])=a+I_U,
  \]
  and therefore
  \[
    \operatorname{im}(\chi_U)
      =\{a+I_U\mid a\in A\}=A/I_U.                             \tag{126a}
  \]
  If $\chi_U([\alpha,a+I_{\alpha,U}])=0$, then
  $a\in I_U$.  By (125a), choose $\beta$ such that
  $a\in I_{\beta,U}$, and by directedness choose
  $\gamma\geq\alpha,\beta$.  The transition maps are the following
  equality:
  \[
    a+I_{\alpha,U}\longmapsto a+I_{\gamma,U}=0
      \quad\text{in }A/I_{\gamma,U},                           \tag{127}
  \]
  whence
  \[
    [\alpha,a+I_{\alpha,U}]=0,
    \qquad
    \ker(\chi_U)=0.
  \]
  The ring-homomorphism criterion, (126a), and the preceding kernel equality
  are the predicate chain
  \[
  \begin{aligned}
    \bigl(\ker\chi_U=0\ \wedge\
      \operatorname{im}\chi_U=A/I_U\bigr)
      &\Longleftrightarrow \operatorname{IsIso}(\chi_U),\\
    \ker\chi_U=0\ \wedge\
      \operatorname{im}\chi_U=A/I_U
      &\Longrightarrow
      \chi_U\colon\varinjlim_\alpha A/I_{\alpha,U}
        \xrightarrow{\sim}A/I_U.
  \end{aligned}                                                     \tag{128}
  \]
  Define
  \[
  \begin{gathered}
    X_{\alpha,U}:=X_\alpha\times_YU
      =\operatorname{Spec}(A/I_{\alpha,U}),\\
    L_U:=L\times_YU=\varprojlim_\alpha X_{\alpha,U}
      =\operatorname{Spec}\!\left(
        \varinjlim_\alpha A/I_{\alpha,U}\right),\\
    \sigma_{i,U}:=\sigma_i\times_Y\operatorname{id}_U\colon
      X\times_YU\xrightarrow{\sim}
      X_{\mathcal I}\times_YU=\operatorname{Spec}(A/I_U).       
  \end{gathered} \tag{128a}
  \]
  Let $\lambda_{\alpha,U}\colon L_U\to X_{\alpha,U}$ denote the limit
  projection and put
  \[
    \kappa_U:=\kappa_{X/Y}\times_Y\operatorname{id}_U\colon
      X\times_YU\longrightarrow L_U.
  \]
  By (119), (121), (126), and the affine relative-spectrum equivalence, for
  every $\alpha$ the two composites satisfy
  \[
    \lambda_{\alpha,U}\circ\kappa_U
      =i_\alpha\times_Y\operatorname{id}_U
      =\lambda_{\alpha,U}\circ\operatorname{Spec}(\chi_U)
         \circ\sigma_{i,U}.                                    \tag{129}
  \]
  The uniqueness clause in the universal property of $L_U$ applied to
  (129) is the following equality:
  \[
    \kappa_U=\operatorname{Spec}(\chi_U)\circ\sigma_{i,U}
      \colon X\times_YU\xrightarrow{\sim}L_U.                  \tag{130}
  \]
  For $s\in A$, set $U_s:=D_U(s)=\operatorname{Spec}(A_s)$.  The
  quasi-coherent sheaf--module correspondence defines the localization
  isomorphisms
  \[
  \begin{gathered}
    \lambda_{\mathcal I,s}\colon (I_U)_s\xrightarrow{\sim}
      \Gamma(U_s,\mathcal I|_{U_s}),\qquad
    \lambda_{\mathcal I,s}(x/s^n)
      :=(x|_{U_s})(s|_{U_s})^{-n},\\
    \lambda_{\mathcal I_\alpha,s}\colon
      (I_{\alpha,U})_s\xrightarrow{\sim}
      \Gamma(U_s,\mathcal I_\alpha|_{U_s}),\qquad
    \lambda_{\mathcal I_\alpha,s}(x/s^n)
      :=(x|_{U_s})(s|_{U_s})^{-n},\\
    \mu_{\mathcal I,s}\colon
      (A/I_U)_s\xrightarrow{\sim}A_s/(I_U)_s,\\
    \mu_{\mathcal I,s}\!\left(\frac{a+I_U}{s^n}\right)
      :=\frac a{s^n}+(I_U)_s,\\
    \mu_{\alpha,s}\colon
      (A/I_{\alpha,U})_s\xrightarrow{\sim}
        A_s/(I_{\alpha,U})_s,\qquad
    \mu_{\alpha,s}\!\left(\frac{a+I_{\alpha,U}}{s^n}\right)
      :=\frac a{s^n}+(I_{\alpha,U})_s.
  \end{gathered} \tag{131}
  \]
  Define the filtered-colimit localization isomorphisms
  \[
  \begin{gathered}
    \delta_s\colon
      \left(\varinjlim_\alpha A/I_{\alpha,U}\right)_s
        \xrightarrow{\sim}
      \varinjlim_\alpha(A/I_{\alpha,U})_s,\\
    \overline\mu_s:=
      \varinjlim_\alpha\mu_{\alpha,s}\colon
      \varinjlim_\alpha(A/I_{\alpha,U})_s
        \xrightarrow{\sim}
      \varinjlim_\alpha A_s/(I_{\alpha,U})_s,\\
    \delta_s^{-1}\colon
      \varinjlim_\alpha(A/I_{\alpha,U})_s
        \xrightarrow{\sim}
      \left(\varinjlim_\alpha A/I_{\alpha,U}\right)_s,\\
    \overline\mu_s^{-1}:=
      \varinjlim_\alpha\mu_{\alpha,s}^{-1}\colon
      \varinjlim_\alpha A_s/(I_{\alpha,U})_s
        \xrightarrow{\sim}
      \varinjlim_\alpha(A/I_{\alpha,U})_s,\\
    \nu_s:=\overline\mu_s\circ\delta_s,\qquad
    \nu_s^{-1}:=\delta_s^{-1}\circ\overline\mu_s^{-1}.            
  \end{gathered} \tag{131a}
  \]
  Transport the localization $(\chi_U)_s$ through the named isomorphisms by
  defining
  \[
    \chi_{U_s}:=
      \mu_{\mathcal I,s}\circ(\chi_U)_s\circ\nu_s^{-1}\colon
      \varinjlim_\alpha A_s/(I_{\alpha,U})_s
        \longrightarrow A_s/(I_U)_s.                             \tag{131b}
  \]
  For every representative, direct substitution into (131b) is the following
  equality:
  \[
    \chi_{U_s}\!\left(
      [\alpha,a/s^n+(I_{\alpha,U})_s]\right)
      =a/s^n+(I_U)_s,
  \]
  Composition of (131b) with $\nu_s$ is the following exact transport
  equation:
  \[
    \chi_{U_s}\circ\nu_s
      =\mu_{\mathcal I,s}\circ(\chi_U)_s.                         \tag{131c}
  \]
  Contravariant application of $\operatorname{Spec}$ to (131c) is the
  following equality:
  \[
    \operatorname{Spec}(\nu_s)\circ\operatorname{Spec}(\chi_{U_s})
      =\operatorname{Spec}((\chi_U)_s)
         \circ\operatorname{Spec}(\mu_{\mathcal I,s}).            \tag{131d}
  \]
  Define the restrictions
  \[
    \sigma_{i,U_s}:=\sigma_i\times_Y\operatorname{id}_{U_s},
    \qquad
    \kappa_{U_s}:=\kappa_{X/Y}\times_Y\operatorname{id}_{U_s}.
  \]
  Their base-change definitions are the following equalities:
  \[
    \sigma_{i,U_s}
      =\sigma_{i,U}|_{X\times_YU_s},
    \qquad
    \kappa_{U_s}
      =\kappa_U|_{X\times_YU_s}.                                  \tag{131e}
  \]
  The inverse of the isomorphism $\operatorname{Spec}(\nu_s)$ and (131d)
  are the following typed equality:
  \[
    \operatorname{Spec}(\chi_{U_s})
      =\operatorname{Spec}(\nu_s)^{-1}
        \circ\operatorname{Spec}((\chi_U)_s)
        \circ\operatorname{Spec}(\mu_{\mathcal I,s}).             \tag{131d'}
  \]
  Equations (130) and (131d)--(131e) are the following equality chain:
  \[
    \kappa_{U_s}
      =\operatorname{Spec}(\chi_{U_s})\circ\sigma_{i,U_s}
      =\kappa_U|_{X\times_YU_s}.
  \]
  Finally, for every affine open $U\subseteq Y$,
  \[
    L\times_YU=L_U,
    \qquad
    \kappa_{X/Y}\times_Y\operatorname{id}_U=\kappa_U,
    \qquad
    \operatorname{IsIso}(\kappa_U)                              \tag{131f}
  \]
  by (130).  The cover by the opens $L\times_YU$, (131f), and target-locality
  of isomorphisms are the following predicate chain:
  \[
  \begin{aligned}
    (\forall U\subseteq Y\text{ affine})\quad
      \operatorname{IsIso}
        (\kappa_{X/Y}|_{X\times_YU})
      &\Longrightarrow \operatorname{IsIso}(\kappa_{X/Y}),\\
    \kappa_{X/Y}&\colon X\xrightarrow{\sim}
      \varprojlim_{\alpha\in A_{X/Y}}X_\alpha.
  \end{aligned}                                                   \tag{132}
  \]
\end{proof}

\begin{definition}[Finite-presentation model system of a finite-type morphism]
  \label{def:nagata-finite-type-fp-model-system}
  \uses{def:nagata-immersion-into-fp-model,
    def:nagata-immersion-shrink-to-closed,
    def:nagata-closed-subscheme-fp-system,
    lemma:nagata-immersion-into-fp-model,
    lemma:nagata-immersion-shrink-to-closed}
  \lean{}

  Let $f\colon X\to S$ be locally of finite type, with $X$ qcqs and $S$
  quasi-separated.  Define the abbreviation and the selected morphisms
  \[
    Y^{\mathrm{imm}}:=Y^{\mathrm{imm}}(X/S),
    \qquad
    i^{\mathrm{imm}}\colon X\to Y^{\mathrm{imm}},
    \qquad
    q^{\mathrm{imm}}\colon Y^{\mathrm{imm}}\to S                \tag{132a}
  \]
  by \Cref{def:nagata-immersion-into-fp-model}, equation (IMM3).  By
  \Cref{lemma:nagata-immersion-into-fp-model}, $q^{\mathrm{imm}}$ is of finite
  presentation.  The standard implication from finite presentation to
  quasi-separatedness is
  \[
    q^{\mathrm{imm}}\text{ is of finite presentation}
      \Longrightarrow
    \Delta_{q^{\mathrm{imm}}}\text{ is quasi-compact}
      \Longleftrightarrow
    q^{\mathrm{imm}}\text{ is quasi-separated}.
  \]
  Let $p_1,p_2\colon
    Y^{\mathrm{imm}}\times_SY^{\mathrm{imm}}\to Y^{\mathrm{imm}}$
  be the fibre-product projections and let
  $\operatorname{pr}_1,\operatorname{pr}_2\colon
    Y^{\mathrm{imm}}\times Y^{\mathrm{imm}}\to Y^{\mathrm{imm}}$
  be the product projections.  Define
  \[
  \begin{gathered}
    k_{Y/S}:=(p_1,p_2)\colon
      Y^{\mathrm{imm}}\times_SY^{\mathrm{imm}}\longrightarrow
      Y^{\mathrm{imm}}\times Y^{\mathrm{imm}},
    \qquad
    \operatorname{pr}_j\circ k_{Y/S}=p_j\quad(j=1,2),\\
    q_{12}:=q^{\mathrm{imm}}\circ p_1
      =q^{\mathrm{imm}}\circ p_2\colon
      Y^{\mathrm{imm}}\times_SY^{\mathrm{imm}}\to S.
  \end{gathered}
  \]
  Define the fibre product and its projections
  \[
  \begin{gathered}
    P_{Y/S}:=(Y^{\mathrm{imm}}\times Y^{\mathrm{imm}})
      \times_{q^{\mathrm{imm}}\times q^{\mathrm{imm}},\,
        S\times S,\,\Delta_S}S,\\
    a_{Y/S}\colon P_{Y/S}\to
      Y^{\mathrm{imm}}\times Y^{\mathrm{imm}},
    \qquad
    b_{Y/S}\colon P_{Y/S}\to S,\\
    (q^{\mathrm{imm}}\times q^{\mathrm{imm}})\circ a_{Y/S}
      =\Delta_S\circ b_{Y/S}.                                   
  \end{gathered} \tag{132b}
  \]
  Composition of (132b) with the two projections $S\times S\to S$ is the
  following equality:
  \[
    q^{\mathrm{imm}}\circ\operatorname{pr}_1\circ a_{Y/S}
      =b_{Y/S}
      =q^{\mathrm{imm}}\circ\operatorname{pr}_2\circ a_{Y/S}.
  \]
  Hence the two fibre-product universal properties define
  \[
  \begin{gathered}
    c_{Y/S}:=(k_{Y/S},q_{12})\colon
      Y^{\mathrm{imm}}\times_SY^{\mathrm{imm}}\to P_{Y/S},\\
    r_{Y/S}:=(\operatorname{pr}_1\circ a_{Y/S},
      \operatorname{pr}_2\circ a_{Y/S})\colon
      P_{Y/S}\to Y^{\mathrm{imm}}\times_SY^{\mathrm{imm}}.     
  \end{gathered} \tag{132c}
  \]
  Their projection equations are the following equality chains:
  \[
  \begin{aligned}
    r_{Y/S}\circ c_{Y/S}
      &=(p_1,p_2)=\operatorname{id}_{
        Y^{\mathrm{imm}}\times_SY^{\mathrm{imm}}},\\
    a_{Y/S}\circ c_{Y/S}\circ r_{Y/S}
      &=k_{Y/S}\circ r_{Y/S}
       =(\operatorname{pr}_1\circ a_{Y/S},
          \operatorname{pr}_2\circ a_{Y/S})=a_{Y/S},\\
    b_{Y/S}\circ c_{Y/S}\circ r_{Y/S}
      &=q_{12}\circ r_{Y/S}
       =q^{\mathrm{imm}}\circ\operatorname{pr}_1\circ a_{Y/S}
       =b_{Y/S}.
  \end{aligned}
  \]
  Uniqueness in the universal property of $P_{Y/S}$ is the following inverse
  equality:
  \[
    c_{Y/S}\circ r_{Y/S}=\operatorname{id}_{P_{Y/S}},
    \qquad
    c_{Y/S}\colon
      Y^{\mathrm{imm}}\times_SY^{\mathrm{imm}}
        \xrightarrow{\sim}P_{Y/S}.                              \tag{132d}
  \]
  Thus the square
  \[
  \begin{tikzcd}
    Y^{\mathrm{imm}}\times_SY^{\mathrm{imm}}
      \arrow[r,"k_{Y/S}"] \arrow[d,"q_{12}"'] &
      Y^{\mathrm{imm}}\times Y^{\mathrm{imm}}
      \arrow[d,"q^{\mathrm{imm}}\times q^{\mathrm{imm}}"]\\
    S \arrow[r,"\Delta_S"'] & S\times S
  \end{tikzcd}
  \]
  is cartesian, with comparison isomorphism $c_{Y/S}$ in (132d).
  Quasi-separatedness of $S$ is
  $\operatorname{IsQuasiCompact}(\Delta_S)$; stability under base change in
  this cartesian square is the following base-change implication:
  \[
    \operatorname{IsQuasiCompact}(\Delta_S)
      \Longrightarrow
      \operatorname{IsQuasiCompact}(k_{Y/S}).                   \tag{132e}
  \]
  The projection equalities
  \[
    \operatorname{pr}_j\circ k_{Y/S}\circ
      \Delta_{q^{\mathrm{imm}}}
      =p_j\circ\Delta_{q^{\mathrm{imm}}}
      =\operatorname{id}_{Y^{\mathrm{imm}}}
      =\operatorname{pr}_j\circ\Delta_{Y^{\mathrm{imm}}}
      \quad(j=1,2)
  \]
  and the product universal property are the following equality:
  \[
    \Delta_{Y^{\mathrm{imm}}}
      =k_{Y/S}\circ\Delta_{q^{\mathrm{imm}}}.                   \tag{133}
  \]
  Quasi-separatedness of $q^{\mathrm{imm}}$ is
  $\operatorname{IsQuasiCompact}(\Delta_{q^{\mathrm{imm}}})$.
  Equations (132e), (133), and stability under composition are the following
  predicate chain:
  \[
  \begin{aligned}
    \operatorname{IsQuasiCompact}(k_{Y/S})\ \wedge\
      \operatorname{IsQuasiCompact}(\Delta_{q^{\mathrm{imm}}})
      &\Longrightarrow
        \operatorname{IsQuasiCompact}
          (k_{Y/S}\circ\Delta_{q^{\mathrm{imm}}})\\
      &\Longleftrightarrow
        \operatorname{IsQuasiCompact}(\Delta_{Y^{\mathrm{imm}}})\\
      &\Longleftrightarrow
        \operatorname{QuasiSeparated}(Y^{\mathrm{imm}}).
  \end{aligned}                                                   \tag{134}
  \]
  Since $X$ is quasi-compact, (134) verifies the hypotheses of
  \Cref{def:nagata-immersion-shrink-to-closed}.  Define
  \[
    Y^\circ:=Y^{\mathrm{cl}}(i^{\mathrm{imm}}),
    \qquad
    u^\circ:=w_{i^{\mathrm{imm}}}^{\mathrm{cl}}
      \colon Y^\circ\hookrightarrow Y^{\mathrm{imm}},
    \qquad
    q^\circ:=q^{\mathrm{imm}}\circ u^\circ\colon Y^\circ\to S,
    \qquad
    i^\circ:=(i^{\mathrm{imm}})^{\mathrm{cl}}\colon X\hookrightarrow Y^\circ.
                                                                    \tag{135}
  \]
  The final equality in
  \Cref{lemma:nagata-immersion-shrink-to-closed}, instantiated with
  $i=i^{\mathrm{imm}}$, is
  \[
    w_{i^{\mathrm{imm}}}^{\mathrm{cl}}
      \circ(i^{\mathrm{imm}})^{\mathrm{cl}}=i^{\mathrm{imm}}.
  \]
  Substitution of the definitions in (135) is the following equality chain:
  \[
  \begin{aligned}
    u^\circ\circ i^\circ&=i^{\mathrm{imm}},\\
    q^\circ\circ i^\circ
      &=(q^{\mathrm{imm}}\circ u^\circ)\circ i^\circ\\
      &=q^{\mathrm{imm}}\circ(u^\circ\circ i^\circ)\\
      &=q^{\mathrm{imm}}\circ i^{\mathrm{imm}}\\
      &=f
        &&\text{by \Cref{def:nagata-immersion-into-fp-model}}.
  \end{aligned}                                                  \tag{136}
  \]
  By \Cref{lemma:nagata-immersion-shrink-to-closed}, $Y^\circ$ is
  quasi-compact, $u^\circ\colon Y^\circ\hookrightarrow Y^{\mathrm{imm}}$ is
  an open immersion, and $i^\circ$ is a closed immersion.  Let
  $\operatorname{pr}_1^\circ,\operatorname{pr}_2^\circ\colon
  Y^\circ\times Y^\circ\to Y^\circ$ be the projections, and define
  \[
  \begin{gathered}
    D^\circ:=Y^{\mathrm{imm}}
      \times_{\Delta_{Y^{\mathrm{imm}}},\,
        Y^{\mathrm{imm}}\times Y^{\mathrm{imm}},\,
        u^\circ\times u^\circ}(Y^\circ\times Y^\circ),\\
    a^\circ\colon D^\circ\to Y^{\mathrm{imm}},
    \qquad
    b^\circ\colon D^\circ\to Y^\circ\times Y^\circ,\\
    \Delta_{Y^{\mathrm{imm}}}\circ a^\circ
      =(u^\circ\times u^\circ)\circ b^\circ,
    \qquad
    b_j^\circ:=\operatorname{pr}_j^\circ\circ b^\circ
      \quad(j=1,2).                                               
  \end{gathered} \tag{136a}
  \]
  Composition of (136a) with the two projections
  $Y^{\mathrm{imm}}\times Y^{\mathrm{imm}}\to Y^{\mathrm{imm}}$
  is the following pair of equalities:
  \[
    a^\circ=u^\circ\circ b_1^\circ,
    \qquad
    a^\circ=u^\circ\circ b_2^\circ.
  \]
  The implication from the open-immersion predicate to left cancellation,
  together with the preceding component equalities, is
  \[
  \begin{aligned}
    \operatorname{IsOpenImmersion}(u^\circ)
      &\Longrightarrow
      \Bigl[
        (\forall r_1,r_2\colon D^\circ\to Y^\circ)\,
        \bigl(u^\circ\circ r_1=u^\circ\circ r_2
          \Longrightarrow r_1=r_2\bigr)
      \Bigr],\\
    u^\circ\circ b_1^\circ=u^\circ\circ b_2^\circ
      &\Longrightarrow b_1^\circ=b_2^\circ.
  \end{aligned}                                                   \tag{136b}
  \]
  The equality
  $\Delta_{Y^{\mathrm{imm}}}\circ u^\circ
    =(u^\circ\times u^\circ)\circ\Delta_{Y^\circ}$ defines
  \[
    c^\circ:=(u^\circ,\Delta_{Y^\circ})\colon Y^\circ\to D^\circ,
    \qquad
    r^\circ:=b_1^\circ\colon D^\circ\to Y^\circ.                 \tag{136c}
  \]
  Equations (136a)--(136c) are the following inverse equality chains:
  \[
  \begin{aligned}
    r^\circ\circ c^\circ
      &=\operatorname{pr}_1^\circ\circ\Delta_{Y^\circ}
       =\operatorname{id}_{Y^\circ},\\
    a^\circ\circ c^\circ\circ r^\circ
      &=u^\circ\circ b_1^\circ=a^\circ,\\
    b^\circ\circ c^\circ\circ r^\circ
      &=\Delta_{Y^\circ}\circ b_1^\circ
       =(b_1^\circ,b_1^\circ)
       =(b_1^\circ,b_2^\circ)=b^\circ.
  \end{aligned}
  \]
  Uniqueness in the universal property of $D^\circ$ is the following inverse
  equality:
  \[
    c^\circ\circ r^\circ=\operatorname{id}_{D^\circ},
    \qquad
    c^\circ\colon Y^\circ\xrightarrow{\sim}D^\circ.              \tag{136d}
  \]
  Thus the commutative square
  \[
  \begin{tikzcd}
    Y^\circ \arrow[r,"\Delta_{Y^\circ}"] \arrow[d,"u^\circ"'] &
      Y^\circ\times Y^\circ
        \arrow[d,"u^\circ\times u^\circ"]\\
    Y^{\mathrm{imm}} \arrow[r,"\Delta_{Y^{\mathrm{imm}}}"'] &
      Y^{\mathrm{imm}}\times Y^{\mathrm{imm}}
  \end{tikzcd}
  \]
  is cartesian, with comparison isomorphism $c^\circ$ in (136d).  In the
  defining fibre-product square (136a), stability under base change and
  (134) are the following predicate:
  \[
    \operatorname{IsQuasiCompact}(b^\circ).
  \]
  The equality
  \[
    b^\circ\circ c^\circ=\Delta_{Y^\circ},
  \]
  the isomorphism predicate for $c^\circ$, and stability under composition
  are the following predicate chain:
  \[
  \begin{aligned}
    \operatorname{IsQuasiCompact}(b^\circ)\ \wedge\
      \operatorname{IsIso}(c^\circ)
      &\Longrightarrow
        \operatorname{IsQuasiCompact}(b^\circ\circ c^\circ)\\
      &\Longleftrightarrow
        \operatorname{IsQuasiCompact}(\Delta_{Y^\circ}).
  \end{aligned}
  \]
  Together with the quasi-compactness of $Y^\circ$, these equalities are
  \[
    Y^\circ\text{ is quasi-compact},
    \qquad
    \Delta_{Y^\circ}\text{ is quasi-compact},
    \qquad
    Y^\circ\text{ is quasi-separated}.                          \tag{137}
  \]
  The first and third assertions in (137) are exactly that $Y^\circ$ is
  qcqs.  Apply
  \Cref{def:nagata-closed-subscheme-fp-system} to the specified closed
  immersion $i^\circ\colon X\hookrightarrow Y^\circ$, and set
  \[
  \begin{gathered}
    \pi_{i^\circ}:=\varepsilon_{i^\circ}:=(i^\circ)^\flat\colon
      \mathcal O_{Y^\circ}\twoheadrightarrow i^\circ_*\mathcal O_X,
    \qquad
    \mathcal I^\circ:=\ker(\pi_{i^\circ}),
    \qquad
    A_{X/S}:=A_{X/Y^\circ},\\
    q_{\mathcal I^\circ}\colon\mathcal O_{Y^\circ}
      \twoheadrightarrow\mathcal O_{Y^\circ}/\mathcal I^\circ,
    \qquad
    \sigma_{i^\circ}\colon X\xrightarrow{\sim}
      \operatorname{Spec}_{Y^\circ}
        (\mathcal O_{Y^\circ}/\mathcal I^\circ)
  \end{gathered}                                                  \tag{138a}
  \]
  with $\sigma_{i^\circ}$ equal to the morphism $\sigma_i$ of (115a)
  instantiated at $i=i^\circ$.  For
  $\alpha=(\mathcal J_\alpha,\iota_{\mathcal J_\alpha})\in A_{X/S}$,
  and for $\alpha\leq\beta$, define
  \[
  \begin{gathered}
    \mathcal I_\alpha:=\mathcal J_\alpha,\qquad
    \pi_\alpha^{Y^\circ}\colon\mathcal O_{Y^\circ}
      \twoheadrightarrow\mathcal O_{Y^\circ}/\mathcal I_\alpha,\\
    \sigma_{\beta\alpha}\colon
      \mathcal O_{Y^\circ}/\mathcal I_\alpha
        \twoheadrightarrow\mathcal O_{Y^\circ}/\mathcal I_\beta,
    \qquad
    \sigma_{\infty\alpha}\colon
      \mathcal O_{Y^\circ}/\mathcal I_\alpha
        \twoheadrightarrow\mathcal O_{Y^\circ}/\mathcal I^\circ,\\
    \sigma_{\beta\alpha}\circ\pi_\alpha^{Y^\circ}
      =\pi_\beta^{Y^\circ},
    \qquad
    \sigma_{\infty\alpha}\circ\pi_\alpha^{Y^\circ}
      =q_{\mathcal I^\circ},\\
    X_\alpha:=\operatorname{Spec}_{Y^\circ}
      (\mathcal O_{Y^\circ}/\mathcal I_\alpha),\qquad
    k_\alpha:=\operatorname{Spec}_{Y^\circ}(\pi_\alpha^{Y^\circ})
      \colon X_\alpha\hookrightarrow Y^\circ,\\
    t_{\beta\alpha}:=\operatorname{Spec}_{Y^\circ}
      (\sigma_{\beta\alpha})\colon X_\beta\hookrightarrow X_\alpha,\\
    i_\alpha:=\operatorname{Spec}_{Y^\circ}
      (\sigma_{\infty\alpha})\circ\sigma_{i^\circ}
      \colon X\hookrightarrow X_\alpha,\qquad
    f_\alpha:=q^\circ\circ k_\alpha\colon X_\alpha\to S.
  \end{gathered}                                                  \tag{138}
  \]
  Equations (120), (136), and (138) expand to
  \[
  \begin{aligned}
    k_\alpha\circ t_{\beta\alpha}=k_\beta,
    &\qquad
    t_{\beta\alpha}\circ i_\beta=i_\alpha,
    \qquad
    k_\alpha\circ i_\alpha=i^\circ,\\
    f_\alpha\circ t_{\beta\alpha}
      &=(q^\circ\circ k_\alpha)\circ t_{\beta\alpha}
       =q^\circ\circ(k_\alpha\circ t_{\beta\alpha})
       =q^\circ\circ k_\beta=f_\beta.
  \end{aligned}                                                 \tag{139}
  \]
  Moreover,
  \[
    f_\alpha\circ i_\alpha
      =(q^\circ\circ k_\alpha)\circ i_\alpha
      =q^\circ\circ(k_\alpha\circ i_\alpha)
      =q^\circ\circ i^\circ=f.                                \tag{140}
  \]
  Finally, let $\kappa_{X/Y^\circ}$ be the morphism (121), and define the
  finite-presentation model system to be the full tuple
  \[
    \mathsf{FP}(X/S)
      =\bigl(A_{X/S},(X_\alpha,f_\alpha,k_\alpha)_\alpha,
        (t_{\beta\alpha})_{\alpha\leq\beta},(i_\alpha)_\alpha,
        \kappa_{X/Y^\circ}\bigr).                               \tag{141}
  \]
\end{definition}

\begin{lemma}[Properties of the finite-presentation model system,
  {\cite[\href{https://stacks.math.columbia.edu/tag/09ZQ}{Tag 09ZQ}]{stacks-project}}]
  \label{lemma:nagata-finite-type-fp-model-system}
  \uses{def:nagata-finite-type-fp-model-system}
  \lean{}

  Every $X_\alpha\to S$ is of finite presentation, every transition map is a
  closed immersion, every $i_\alpha$ is a closed immersion, and
  $X\cong\varprojlim_\alpha X_\alpha$ over $S$.
\end{lemma}

\begin{proof}
  \uses{lemma:nagata-immersion-into-fp-model,
    lemma:nagata-immersion-shrink-to-closed,
    lemma:nagata-closed-subscheme-fp-system}
  Use the abbreviations
  \[
    Y^{\mathrm{imm}}:=Y^{\mathrm{imm}}(X/S),
    \qquad
    Y^\circ:=Y^{\mathrm{cl}}(i^{\mathrm{imm}})
  \]
  of (132a) and (135).  By
  \Cref{lemma:nagata-immersion-into-fp-model},
  \[
    q^{\mathrm{imm}}\colon Y^{\mathrm{imm}}\longrightarrow S
      \quad\text{is of finite presentation}.                   \tag{142}
  \]
  By \Cref{lemma:nagata-immersion-shrink-to-closed}, the inclusion
  $u^\circ\colon Y^\circ\hookrightarrow Y^{\mathrm{imm}}$ is a
  quasi-compact open immersion.  The finite-presentation criterion for an
  open immersion is the equivalence
  \[
    u^\circ\text{ is of finite presentation}
      \Longleftrightarrow
    u^\circ\text{ is quasi-compact}.
  \]
  Its application to $u^\circ$ records
  \[
    u^\circ\quad\text{is of finite presentation}.              \tag{143}
  \]
  Stability of finite presentation under composition, instantiated at
  (135), is
  \[
    \bigl(q^{\mathrm{imm}}\text{ is of finite presentation}\bigr)
      \wedge
    \bigl(u^\circ\text{ is of finite presentation}\bigr)
      \Longrightarrow
    q^{\mathrm{imm}}\circ u^\circ\text{ is of finite presentation}.
  \]
  Applying this statement to (142)--(143) records
  \[
    q^\circ=q^{\mathrm{imm}}\circ u^\circ\colon Y^\circ\to S
      \quad\text{is of finite presentation}.                   \tag{144}
  \]
  For the system of
  \Cref{def:nagata-closed-subscheme-fp-system},
  \Cref{lemma:nagata-closed-subscheme-fp-system} asserts the exact
  properties
  \[
    k_\alpha\colon X_\alpha\hookrightarrow Y^\circ
      \text{ of finite presentation},
    \qquad
    t_{\beta\alpha}\colon X_\beta\hookrightarrow X_\alpha
      \text{ a closed immersion},                              \tag{145}
  \]
  and
  \[
    \kappa_{X/Y^\circ}\colon X\xrightarrow{\sim}
      \varprojlim_\alpha X_\alpha.                             \tag{146}
  \]
  Stability under composition applied to
  $f_\alpha=q^\circ\circ k_\alpha$ in (138) and the two
  finite-presentation assertions (144)--(145) records
  \[
    f_\alpha=q^\circ\circ k_\alpha\colon X_\alpha\to S
      \quad\text{of finite presentation}.                      \tag{147}
  \]
  The isomorphism (115a) and quotient map (117), instantiated with
  $Y=Y^\circ$ and $i=i^\circ$, are
  \[
  \begin{gathered}
    \sigma_{i^\circ}\colon X\xrightarrow{\sim}
      \operatorname{Spec}_{Y^\circ}
        (\mathcal O_{Y^\circ}/\mathcal I^\circ),\\
    \sigma_{\infty\alpha}\colon
      \mathcal O_{Y^\circ}/\mathcal I_\alpha
      \twoheadrightarrow\mathcal O_{Y^\circ}/\mathcal I^\circ.
  \end{gathered}                                                  \tag{148}
  \]
  Equation (148) states that $\sigma_{\infty\alpha}$ is surjective.  The
  relative affine criterion is the equivalence
  \[
    \operatorname{Spec}_{Y^\circ}(\sigma_{\infty\alpha})
      \text{ is a closed immersion}
      \Longleftrightarrow
    \sigma_{\infty\alpha}\text{ is surjective}.
  \]
  Under (119), the exact factorization is
  \[
    i_\alpha
      =\operatorname{Spec}_{Y^\circ}(\sigma_{\infty\alpha})
         \circ\sigma_{i^\circ}\colon X\hookrightarrow X_\alpha. \tag{149}
  \]
  The first factor in (149) is a closed immersion by the displayed
  equivalence, and $\sigma_{i^\circ}$ is an isomorphism by (148).  Hence
  stability under composition gives
  \[
    \operatorname{IsClosedImmersion}(i_\alpha).
  \]

  Put $L=\varprojlim_\alpha X_\alpha$, and let
  $\pi_\alpha\colon L\to X_\alpha$ be its projections.  For
  $\alpha,\beta\in A_{X/S}$, choose
  $\gamma\in A_{X/S}$ with $\alpha,\beta\leq\gamma$.  The projection
  identities and (139) form the chain
  \[
  \begin{aligned}
    f_\alpha\circ\pi_\alpha
      &=f_\alpha\circ t_{\gamma\alpha}\circ\pi_\gamma
       =f_\gamma\circ\pi_\gamma\\
      &=f_\beta\circ t_{\gamma\beta}\circ\pi_\gamma
       =f_\beta\circ\pi_\beta.
  \end{aligned}
  \]
  Define the morphism represented by this common composite by
  \[
    f_L=f_\alpha\circ\pi_\alpha\colon L\to S                    \tag{150}
  \]
  for any $\alpha\in A_{X/S}$.  Substitution of (121) and (140) into (150)
  is the chain
  \[
    f_L\circ\kappa_{X/Y^\circ}
      =f_\alpha\circ\pi_\alpha\circ\kappa_{X/Y^\circ}
      =f_\alpha\circ i_\alpha
      =f.                                                        \tag{151}
  \]
  By definition, $\kappa_{X/Y^\circ}$ is an $S$-morphism precisely when
  $f_L\circ\kappa_{X/Y^\circ}=f$; equation (151) verifies this equality.
  Equations (145), (147), (149),
  and (151) are the four assertions of the lemma.
\end{proof}

\begin{definition}[Selected stage for fixed-base separatedness]
  \label{def:nagata-fixed-base-separated-stage}
  \uses{def:nagata-choice-notation,
    def:nagata-directed-finite-upper-bound,
    def:nagata-limit-datum,
    lemma:nagata-affine-transition-limit-base-change,
    lemma:nagata-limit-datum}
  \lean{}

  Let $(A,\leq)$ be a nonempty directed preorder, and let
  \[
    \mathsf X
      =\bigl((X_\alpha,f_\alpha)_{\alpha\in A},
        (t_{\beta\alpha})_{\alpha\leq\beta},
        X,f,L_{\mathsf X},(\pi_\alpha)_{\alpha\in A},
        (p_\alpha)_{\alpha\in A},\kappa_{\mathsf X}\bigr)         \tag{152}
  \]
  be an inverse system over a quasi-separated scheme $S$ such that
  \[
  \begin{gathered}
    f_\alpha\colon X_\alpha\to S,\qquad f\colon X\to S,\\
    t_{\beta\alpha}\colon X_\beta\to X_\alpha,\qquad
    \operatorname{IsAffine}(t_{\beta\alpha}),\qquad
    \operatorname{QuasiCompact}(X_\alpha)
      \wedge\operatorname{QuasiSeparated}(X_\alpha),\\
    p_\alpha\colon X\to X_\alpha,\qquad
    t_{\alpha\alpha}=\operatorname{id}_{X_\alpha},\qquad
    t_{\gamma\alpha}
      =t_{\beta\alpha}\circ t_{\gamma\beta}
      \quad(\alpha\leq\beta\leq\gamma),\\
    f_\alpha\circ t_{\beta\alpha}=f_\beta,\qquad
    t_{\beta\alpha}\circ p_\beta=p_\alpha,\qquad
    f_\alpha\circ p_\alpha=f,\qquad
    L_{\mathsf X}:=\varprojlim_{\alpha\in A}X_\alpha,\qquad
    \pi_\alpha\colon L_{\mathsf X}\to X_\alpha,\qquad
    t_{\beta\alpha}\circ\pi_\beta=\pi_\alpha,\\
    \kappa_{\mathsf X}\colon X\to L_{\mathsf X},\qquad
    \operatorname{IsIso}(\kappa_{\mathsf X}),\qquad
    \pi_\alpha\circ\kappa_{\mathsf X}=p_\alpha .
  \end{gathered}                                                 \tag{153}
  \]
  For every scheme \(T\), define
  \[
    \mathcal C_{\mathsf X}(T):=
      \left\{(g_\alpha)_\alpha\in
        \prod_{\alpha\in A}\operatorname{Hom}(T,X_\alpha)
      \ \middle|\
        (\forall\alpha\leq\beta)\
          t_{\beta\alpha}\circ g_\beta=g_\alpha\right\}.           \tag{153a}
  \]
  The symbol \(L_{\mathsf X}=\varprojlim_\alpha X_\alpha\) in (153) means
  that, for every scheme \(T\), the function
  \[
    \Lambda_{\mathsf X,T}\colon
      \operatorname{Hom}(T,L_{\mathsf X})
        \xrightarrow{\sim}\mathcal C_{\mathsf X}(T),\qquad
    h\longmapsto(\pi_\alpha\circ h)_\alpha                         \tag{153b}
  \]
  is bijective.
  Assume \(\operatorname{IsSeparated}(f)\).  Fix $\alpha_0\in A$, an integer $m\geq1$,
  and affine open immersions
  $e_r\colon P_r\hookrightarrow X_{\alpha_0}\times_SX_{\alpha_0}$ satisfying
  \[
    \left|X_{\alpha_0}\times_SX_{\alpha_0}\right|
      =\bigcup_{r=1}^m |e_r|(|P_r|).                            \tag{154}
  \]
  For $\alpha\geq\alpha_0$, define
  \[
  \begin{gathered}
    t^{(2)}_{\alpha,\alpha_0}
      =t_{\alpha,\alpha_0}\times_St_{\alpha,\alpha_0}\colon
      X_\alpha\times_SX_\alpha
        \longrightarrow X_{\alpha_0}\times_SX_{\alpha_0},\\
    P_{r,\alpha}
      =P_r\times_{X_{\alpha_0}\times_SX_{\alpha_0}}
        (X_\alpha\times_SX_\alpha),
    \qquad
    u_{r,\alpha}\colon P_{r,\alpha}\to P_r,\qquad
    e_{r,\alpha}\colon P_{r,\alpha}\hookrightarrow
      X_\alpha\times_SX_\alpha,\qquad
    e_r\circ u_{r,\alpha}
      =t^{(2)}_{\alpha,\alpha_0}\circ e_{r,\alpha}.
  \end{gathered}                                                 \tag{155}
  \]
  Define also
  \[
    D_{r,\alpha}
      =X_\alpha\times_{\Delta_{f_\alpha},\,
          X_\alpha\times_SX_\alpha,\,e_{r,\alpha}}P_{r,\alpha},
    \qquad
    d_{r,\alpha}\colon D_{r,\alpha}\to X_\alpha,\qquad
    \delta_{r,\alpha}\colon D_{r,\alpha}\to P_{r,\alpha},
    \qquad
    \Delta_{f_\alpha}\circ d_{r,\alpha}
      =e_{r,\alpha}\circ\delta_{r,\alpha}.                     \tag{156}
  \]
  For $\alpha_0\leq\alpha\leq\beta$, define the transition morphisms
  $s^P_{r,\beta\alpha}$ and $s^D_{r,\beta\alpha}$ by the fibre-product
  equations
  \[
  \begin{gathered}
    s^P_{r,\beta\alpha}\colon P_{r,\beta}\to P_{r,\alpha},\qquad
    u_{r,\alpha}\circ s^P_{r,\beta\alpha}=u_{r,\beta},\qquad
    e_{r,\alpha}\circ s^P_{r,\beta\alpha}
      =(t_{\beta\alpha}\times_St_{\beta\alpha})\circ e_{r,\beta},\\
    s^D_{r,\beta\alpha}\colon D_{r,\beta}\to D_{r,\alpha},\qquad
    d_{r,\alpha}\circ s^D_{r,\beta\alpha}
      =t_{\beta\alpha}\circ d_{r,\beta},\qquad
    \delta_{r,\alpha}\circ s^D_{r,\beta\alpha}
      =s^P_{r,\beta\alpha}\circ\delta_{r,\beta}.
  \end{gathered}                                                  \tag{FS0}
  \]
  Define
  \[
  \begin{gathered}
    p_{\alpha_0}^{(2)}:=p_{\alpha_0}\times_Sp_{\alpha_0},\qquad
    P_{r,\infty}:=
      P_r\times_{X_{\alpha_0}\times_SX_{\alpha_0},\,
        (e_r,p_{\alpha_0}^{(2)})}(X\times_SX),\\
    u_{r,\infty}\colon P_{r,\infty}\to P_r,\qquad
    e_{r,\infty}\colon P_{r,\infty}\to X\times_SX,\qquad
    e_r\circ u_{r,\infty}=p_{\alpha_0}^{(2)}\circ e_{r,\infty},\\
    D_{r,\infty}:=
      X\times_{\Delta_f,\,X\times_SX,\,e_{r,\infty}}P_{r,\infty},
    \qquad
    d_{r,\infty}\colon D_{r,\infty}\to X,
    \qquad
    \delta_{r,\infty}\colon D_{r,\infty}\to P_{r,\infty},\\
    \Delta_f\circ d_{r,\infty}
      =e_{r,\infty}\circ\delta_{r,\infty}.
  \end{gathered}                                                  \tag{FS1}
  \]
  Since $\delta_{r,\infty}$ is the base change of $\Delta_f$ along
  $e_{r,\infty}$, separatedness of $f$ and base-change stability of closed
  immersions give the predicate chain
  \[
    \operatorname{IsSeparated}(f)
      \Longleftrightarrow
    \operatorname{IsClosedImmersion}(\Delta_f)
      \Longrightarrow
    \operatorname{IsClosedImmersion}(\delta_{r,\infty}).         \tag{FS2}
  \]
  For $\alpha\geq\alpha_0$, let
  $p^{(2)}_\alpha:=p_\alpha\times_Sp_\alpha$ and define the cone maps
  \[
  \begin{gathered}
    \rho^P_{r,\alpha}\colon P_{r,\infty}\to P_{r,\alpha},\qquad
    u_{r,\alpha}\circ\rho^P_{r,\alpha}=u_{r,\infty},\qquad
    e_{r,\alpha}\circ\rho^P_{r,\alpha}
      =p^{(2)}_\alpha\circ e_{r,\infty},\\
    \rho^D_{r,\alpha}\colon D_{r,\infty}\to D_{r,\alpha},\qquad
    d_{r,\alpha}\circ\rho^D_{r,\alpha}
      =p_\alpha\circ d_{r,\infty},\qquad
    \delta_{r,\alpha}\circ\rho^D_{r,\alpha}
      =\rho^P_{r,\alpha}\circ\delta_{r,\infty}.
  \end{gathered}                                                  \tag{FS3}
  \]
  By
  \Cref{lemma:nagata-affine-transition-limit-base-change}, the reindexed
  systems in (FS0) carry limit data.  Their limit/base-change isomorphisms
  are the unique morphisms
  \[
  \begin{gathered}
    \kappa^P_r\colon P_{r,\infty}
      \xrightarrow{\sim}\varprojlim_{\alpha\geq\alpha_0}P_{r,\alpha},
    \qquad \operatorname{pr}^P_{r,\alpha}\circ\kappa^P_r
      =\rho^P_{r,\alpha},\\
    \kappa^D_r\colon D_{r,\infty}
      \xrightarrow{\sim}\varprojlim_{\alpha\geq\alpha_0}D_{r,\alpha},
    \qquad \operatorname{pr}^D_{r,\alpha}\circ\kappa^D_r
      =\rho^D_{r,\alpha},\\
    \left(\varprojlim_{\alpha\geq\alpha_0}\delta_{r,\alpha}\right)
      \circ\kappa^D_r
      =\kappa^P_r\circ\delta_{r,\infty}.
  \end{gathered}                                                  \tag{FS4}
  \]
  For \(r\in\{1,\ldots,m\}\), define the witness set
  \[
    E_r:=\left\{\beta\in A\ \middle|\
      \alpha_0\leq\beta\ \wedge\
      (\forall\alpha\in A)\
      \bigl(\beta\leq\alpha\Longrightarrow
        \operatorname{IsClosedImmersion}(\delta_{r,\alpha})\bigr)
      \right\}.                                                   \tag{156a}
  \]
  Applied to (FS0)--(FS4), the finite-presentation limit theorem for a
  closed immersion
  (\cite[\href{https://stacks.math.columbia.edu/tag/01ZH}{Tag 01ZH}]
  {stacks-project}) is the implication
  \[
    \operatorname{IsClosedImmersion}(\delta_{r,\infty})
      \Longrightarrow
    (\exists\beta\geq\alpha_0)(\forall\alpha\geq\beta)\,
      \operatorname{IsClosedImmersion}(\delta_{r,\alpha}),       \tag{FS5}
  \]
  and (FS2), (FS5), and the definition (156a) form the chain
  \[
    \operatorname{IsClosedImmersion}(\delta_{r,\infty})
      \Longrightarrow (\exists\beta\in E_r)
      \Longrightarrow E_r\neq\varnothing
      \qquad(1\leq r\leq m).                                     \tag{156b}
  \]
  By \Cref{def:nagata-choice-notation}, define
  \[
    \alpha_r:=\operatorname{ch}(E_r),\qquad
    \alpha_r\in E_r,
  \]
  whose membership expands to
  \[
    \alpha_0\leq\alpha_r,\qquad
    \alpha_r\leq\alpha
      \Longrightarrow
      \operatorname{IsClosedImmersion}(\delta_{r,\alpha}).        \tag{157}
  \]
  Let $u_A\colon\operatorname{Fin}(A)\to A$ be the function of
  \Cref{def:nagata-directed-finite-upper-bound}, and define
  \[
  \begin{gathered}
    F_{\mathrm{sep}}(\mathsf X,\alpha_0,(P_r)_r)
      =\{\alpha_0,\alpha_1,\ldots,\alpha_m\}\subseteq A,\\
    \alpha_{\mathrm{sep}}
      =u_A\bigl(F_{\mathrm{sep}}(\mathsf X,\alpha_0,(P_r)_r)\bigr).
  \end{gathered}                                                 \tag{158}
  \]
  Define the witness tuple
  \[
    \mathcal W_{\mathrm{sep}}(\mathsf X)
      =\bigl(\alpha_0,m,(P_r,e_r)_{r=1}^m,
        (\alpha_r)_{r=1}^m,F_{\mathrm{sep}},
        \alpha_{\mathrm{sep}}\bigr)                              \tag{159}
  \]
  whose components are exactly the objects defined in (154), (157), and
  (158).
\end{definition}

\begin{lemma}[Eventual separatedness over a fixed base,
  {\cite[\href{https://stacks.math.columbia.edu/tag/01ZH}{Tag 01ZH}]{stacks-project}}]
  \label{lemma:nagata-fixed-base-separated-eventually}
  \uses{def:nagata-fixed-base-separated-stage}
  \lean{}

  Let $S$ be quasi-separated, let $(A,\leq)$ be a nonempty directed
  preorder, and let $\mathsf X$ and
  $\mathcal W_{\mathrm{sep}}(\mathsf X)$ be the tuples (152) and (159) of
  \Cref{def:nagata-fixed-base-separated-stage}.  For every
  $\alpha\geq\alpha_{\mathrm{sep}}$,
  \[
    \operatorname{IsSeparated}(f_\alpha).
  \]
\end{lemma}

\begin{proof}
  Use the symbols (152)--(159), and fix
  $\alpha\in A$ with $\alpha_{\mathrm{sep}}\leq\alpha$.  By
  \Cref{lemma:nagata-directed-finite-upper-bound}, for every
  $r\in\{1,\ldots,m\}$,
  \[
    \alpha_r\leq\alpha_{\mathrm{sep}}\leq\alpha.                 \tag{160}
  \]
  Equations (157) and (160) are the assertion
  \[
    \operatorname{IsClosedImmersion}(\delta_{r,\alpha})
      \qquad(1\leq r\leq m).
                                                                    \tag{161}
  \]
  The cartesian definition (155) gives the equality of underlying subsets
  \[
    |e_{r,\alpha}|(|P_{r,\alpha}|)
      =|t^{(2)}_{\alpha,\alpha_0}|^{-1}\bigl(|e_r|(|P_r|)\bigr).
  \]
  Taking inverse images in (154) and distributing inverse image over the
  finite union gives the chain
  \[
  \begin{aligned}
    |X_\alpha\times_SX_\alpha|
      &=|t^{(2)}_{\alpha,\alpha_0}|^{-1}
          \left(\left|X_{\alpha_0}\times_SX_{\alpha_0}\right|\right)\\
      &=|t^{(2)}_{\alpha,\alpha_0}|^{-1}
          \left(\bigcup_{r=1}^m|e_r|(|P_r|)\right)\\
      &=\bigcup_{r=1}^m
          |t^{(2)}_{\alpha,\alpha_0}|^{-1}\bigl(|e_r|(|P_r|)\bigr)\\
      &=\bigcup_{r=1}^m|e_{r,\alpha}|(|P_{r,\alpha}|).
  \end{aligned}                                                   \tag{162}
  \]
  By (156), the base change of $\Delta_{f_\alpha}$ by
  $e_{r,\alpha}$ is $\delta_{r,\alpha}$.  The target-locality equivalence for
  closed immersions, instantiated with the cover (162), is
  \[
    \operatorname{IsClosedImmersion}(\Delta_{f_\alpha})
      \Longleftrightarrow
    (\forall r)\ \operatorname{IsClosedImmersion}(\delta_{r,\alpha}).
  \]
  Its right-hand side is (161), so
  \[
    \operatorname{IsClosedImmersion}(\Delta_{f_\alpha}).          \tag{163}
  \]
  By the diagonal criterion, (163) is equivalent to
  \[
    \operatorname{IsSeparated}(f_\alpha).
  \]
\end{proof}

\begin{lemma}[A separated finite-type morphism over a qcqs base has qcqs source]
  \label{lemma:nagata-qcqs-source-separated-finite-type}
  \lean{}

  Let \(f\colon X\to S\) satisfy
  \[
    \operatorname{IsSeparated}(f)\ \wedge\
    \operatorname{FiniteType}(f)\ \wedge\
    \operatorname{QuasiCompact}(S)\ \wedge\
    \operatorname{QuasiSeparated}(S).
  \]
  Substitution of the displayed expansion of $z$ into the fraction $z/b^n$
  is the equality
  \[
    \operatorname{QuasiCompact}(X)\ \wedge\
    \operatorname{QuasiSeparated}(X).
  \]
\end{lemma}

\begin{proof}
  Let \(s_X\colon X\to\operatorname{Spec}(\mathbf Z)\) and
  \(s_S\colon S\to\operatorname{Spec}(\mathbf Z)\) be the structural
  morphisms of (P1a).  Expanding (P1) and applying (PC3) of
  \Cref{lemma:nagata-morphism-property-calculus} gives
  \[
  \begin{gathered}
    \operatorname{FiniteType}(f)
      \Longrightarrow\operatorname{IsQuasiCompact}(f),\\
    \operatorname{QuasiCompact}(S)
      \Longleftrightarrow\operatorname{IsQuasiCompact}(s_S),\\
    s_X=s_S\circ f,\qquad
    \operatorname{IsQuasiCompact}(s_S)
      \wedge\operatorname{IsQuasiCompact}(f)
      \Longrightarrow\operatorname{IsQuasiCompact}(s_X),\\
    \operatorname{IsQuasiCompact}(s_X)
      \Longleftrightarrow\operatorname{QuasiCompact}(X).
  \end{gathered}                                                     \tag{163a}
  \]
  Define
  \[
    k=(\operatorname{pr}_1,\operatorname{pr}_2)\colon
      X\times_SX\longrightarrow X\times X.
  \]
  For every scheme $Z$, define
  \[
  \begin{aligned}
    \Theta_Z\colon
    \operatorname{Hom}\!\left(
      Z,(X\times X)\times_{S\times S,\Delta_S}S\right)
      &\longrightarrow\operatorname{Hom}(Z,X\times_SX),\\
    (x_1,x_2,s)&\longmapsto(x_1,x_2).
  \end{aligned}                                                   \tag{164}
  \]
  The source and target of (164) are the explicitly equal sets
  \[
  \begin{aligned}
    &\{(x_1,x_2,s)\mid f\circ x_1=s=f\circ x_2\}
      \xrightarrow[\ (x_1,x_2)\mapsto
          (x_1,x_2,f\circ x_1)\ ]{\ (x_1,x_2,s)\mapsto(x_1,x_2)\ }
      \{(x_1,x_2)\mid f\circ x_1=f\circ x_2\}.
  \end{aligned}                                                  \tag{165}
  \]
  The two displayed arrows in (165) are mutually inverse.  Naturality in
  $Z$ and Yoneda define mutually inverse morphisms
  \[
  \begin{gathered}
    \vartheta\colon
      (X\times X)\times_{S\times S,\Delta_S}S
        \xrightarrow{\sim}X\times_SX,\\
    \vartheta^{-1}=(k,f\circ\operatorname{pr}_1)\colon
      X\times_SX
        \xrightarrow{\sim}
      (X\times X)\times_{S\times S,\Delta_S}S,\\
    \vartheta\circ\vartheta^{-1}=\operatorname{id}_{X\times_SX},
    \qquad
    \vartheta^{-1}\circ\vartheta
      =\operatorname{id}_{(X\times X)\times_{S\times S,\Delta_S}S}.
  \end{gathered}                                                   \tag{166}
  \]
  The defining equality for $X\times_SX$ gives
  \[
    (f\times f)\circ k
      =(f\circ\operatorname{pr}_1,f\circ\operatorname{pr}_2)
      =(f\circ\operatorname{pr}_1,f\circ\operatorname{pr}_1)
      =\Delta_S\circ f\circ\operatorname{pr}_1.
  \]
  This equality and the comparison isomorphism
  $\vartheta^{-1}$ in (166) identify the square
  \[
  \begin{tikzcd}
    X\times_SX \arrow[r,"k"]
      \arrow[d,"f\circ\operatorname{pr}_1"'] &
      X\times X \arrow[d,"f\times f"] \\
    S \arrow[r,"\Delta_S"'] & S\times S
  \end{tikzcd}
  \]
  as cartesian.
  By (P1b),
  \(\operatorname{QuasiSeparated}(S)
    \Longleftrightarrow\operatorname{IsQuasiCompact}(\Delta_S)\).
  Stability under the displayed base change is the implication
  \[
    \operatorname{IsQuasiCompact}(\Delta_S)
      \Longrightarrow\operatorname{IsQuasiCompact}(k).             \tag{166a}
  \]
  By (P0)--(P0a), the homeomorphism
  \(|\Delta_f|_{\operatorname{Im}}\) identifies
  \(|\Delta_f|^{-1}(W)\) with the closed subset
  \(W\cap\operatorname{Im}(|\Delta_f|)\) of every quasi-compact open
  \(W\subseteq|X\times_SX|\).  Thus
  \[
  \begin{aligned}
    \operatorname{IsSeparated}(f)
      \Longleftrightarrow\operatorname{IsClosedImmersion}(\Delta_f)
      &\Longrightarrow
      \bigl[(\forall W\subseteq|X\times_SX|)\
        (W\text{ open and quasi-compact})\\
      &\hspace{6em}\Longrightarrow
        |\Delta_f|^{-1}(W)\text{ quasi-compact}\bigr]\\
      &\Longleftrightarrow\operatorname{IsQuasiCompact}(\Delta_f).
  \end{aligned}                                                     \tag{166b}
  \]
  The absolute diagonal has the explicit factorization determined by
  \[
    \operatorname{pr}_1\circ k\circ\Delta_f
      =\operatorname{id}_X
      =\operatorname{pr}_2\circ k\circ\Delta_f,
    \qquad
    \Delta_X=k\circ\Delta_f\colon X\longrightarrow X\times X.   \tag{167}
  \]
  Stability of quasi-compact morphisms under composition, applied to (167),
  is
  \[
    \operatorname{IsQuasiCompact}(k)
      \wedge\operatorname{IsQuasiCompact}(\Delta_f)
      \Longrightarrow\operatorname{IsQuasiCompact}(\Delta_X).
  \]
  By (P1b),
  \[
    \operatorname{IsQuasiCompact}(\Delta_X)
      \Longleftrightarrow\operatorname{QuasiSeparated}(X).
  \]
\end{proof}

\begin{definition}[Separated finite-presentation envelope]
  \label{def:nagata-finite-type-closed-in-fp}
  \uses{def:nagata-choice-notation,def:nagata-finite-type-fp-model-system,
    def:nagata-fixed-base-separated-stage,
    lemma:nagata-qcqs-source-separated-finite-type,
    lemma:nagata-finite-type-fp-model-system}
  \lean{}

  Let \(f\colon X\to S\) satisfy
  \[
    \operatorname{QuasiCompact}(S)\ \wedge\
    \operatorname{QuasiSeparated}(S)\ \wedge\
    \operatorname{IsSeparated}(f)\ \wedge\
    \operatorname{FiniteType}(f).                                 \tag{167a}
  \]
  By \Cref{lemma:nagata-qcqs-source-separated-finite-type},
  \[
    \operatorname{QuasiCompact}(X)
      \ \wedge\ \operatorname{QuasiSeparated}(X).                 \tag{168}
  \]
  Equation (168) verifies the qcqs hypothesis of
  \Cref{def:nagata-finite-type-fp-model-system}; denote its output tuple by
  \[
    \mathsf{FP}(X/S)
      =\bigl(A_{X/S},(X_\alpha,f_\alpha,k_\alpha)_\alpha,
        (t_{\beta\alpha})_{\alpha\leq\beta},(i_\alpha)_\alpha,
        \kappa_{X/Y^\circ}\bigr).                               \tag{169}
  \]
  By \Cref{lemma:nagata-finite-type-fp-model-system}, for every
  $\alpha\leq\beta$,
  \[
  \begin{gathered}
    \operatorname{FinitePresentation}(f_\alpha),\qquad
    \operatorname{IsClosedImmersion}(t_{\beta\alpha}),\\
    \operatorname{IsClosedImmersion}(i_\alpha),\qquad
    \operatorname{IsIso}(\kappa_{X/Y^\circ}),\qquad
    \kappa_{X/Y^\circ}\colon X\to\varprojlim_\alpha X_\alpha.
  \end{gathered}                                                 \tag{170}
  \]
  Finite presentation makes $f_\alpha$ quasi-compact and
  quasi-separated.  Since $S$ is quasi-compact,
  \[
    \operatorname{QuasiCompact}(X_\alpha).                         \tag{171}
  \]
  Define
  \[
    c_\alpha=(\operatorname{pr}_1,\operatorname{pr}_2)\colon
      X_\alpha\times_SX_\alpha\to X_\alpha\times X_\alpha.
  \]
  The map $c_\alpha$ is the base change of $\Delta_S$, and
  \[
    \Delta_{X_\alpha}=c_\alpha\circ\Delta_{f_\alpha}.            \tag{172}
  \]
  Both factors on the right of (172) are quasi-compact: $c_\alpha$ because
  it is the base change of the quasi-compact morphism $\Delta_S$, and
  $\Delta_{f_\alpha}$ because $f_\alpha$ is quasi-separated.  Stability
  under composition applied to (172) is
  \[
    \operatorname{IsQuasiCompact}(\Delta_{X_\alpha}),\qquad
    \operatorname{QuasiSeparated}(X_\alpha).                       \tag{173}
  \]
  Equations (171) and (173) state that $X_\alpha$ is qcqs.
  Closed immersions are affine, so (170) gives
  \[
    \operatorname{IsClosedImmersion}(t_{\beta\alpha})
      \Longrightarrow\operatorname{IsAffine}(t_{\beta\alpha}).     \tag{173a}
  \]
  Put
  \[
    L_{X/S}:=\varprojlim_{\alpha\in A_{X/S}}X_\alpha,\qquad
    \pi_\alpha\colon L_{X/S}\to X_\alpha,\qquad
    p_\alpha:=i_\alpha,\qquad
    \kappa_{X/S}:=\kappa_{X/Y^\circ}\colon X\to L_{X/S}.            \tag{173b}
  \]
  Equations (121) and (140) give
  \[
    t_{\beta\alpha}\circ\pi_\beta=\pi_\alpha,\qquad
    \pi_\alpha\circ\kappa_{X/S}=p_\alpha,\qquad
    f_\alpha\circ p_\alpha=f.                                      \tag{173c}
  \]
  Thus (140), (170)--(173c) are the conditions (152)--(153).  Define this
  fixed-base system by
  \[
    \mathsf X_{X/S}
      =\bigl((X_\alpha,f_\alpha)_\alpha,
        (t_{\beta\alpha})_{\alpha\leq\beta},
        X,f,L_{X/S},(\pi_\alpha)_\alpha,
        (i_\alpha)_\alpha,\kappa_{X/S}\bigr).                      \tag{174}
  \]

  Let $u_{X/S}\colon\operatorname{Fin}(A_{X/S})\to A_{X/S}$ be the selected
  upper-bound function of
  \Cref{def:nagata-directed-finite-upper-bound}, and define
  \[
    \alpha^{X/S}_0=u_{X/S}(\varnothing).                         \tag{175}
  \]
  Put $\alpha_0=\alpha^{X/S}_0$.  The product
  $X_{\alpha_0}\times X_{\alpha_0}$ is quasi-compact by (171), and
  $c_{\alpha_0}$ is quasi-compact by (172).  The inverse image of a
  quasi-compact open under a quasi-compact morphism is quasi-compact; here
  this assertion is
  \[
    |X_{\alpha_0}\times_SX_{\alpha_0}|
      =|c_{\alpha_0}|^{-1}(|X_{\alpha_0}\times X_{\alpha_0}|),
    \qquad
    \operatorname{QuasiCompact}(X_{\alpha_0}\times_SX_{\alpha_0}).
                                                                    \tag{175a}
  \]
  Define the affine-cover witness set
  \[
    \mathscr A_{X/S}:=
      \left\{\bigl(m,(R_r,e_r)_{r=1}^m\bigr)\ \middle|\
      \begin{array}{l}
        m\in\mathbf N,\ m\geq1,\\
        e_r\colon\operatorname{Spec}(R_r)\to
          X_{\alpha^{X/S}_0}\times_SX_{\alpha^{X/S}_0},\\
        \operatorname{IsOpenImmersion}(e_r)\quad(1\leq r\leq m),\\
        |X_{\alpha^{X/S}_0}\times_SX_{\alpha^{X/S}_0}|
          =\displaystyle\bigcup_{r=1}^m
             \operatorname{Im}(|e_r|)
      \end{array}\right\}.                                         \tag{175b}
  \]
  The affine-basis axiom for schemes and (175a) give the formal chain
  \[
  \begin{gathered}
    (\forall x\in|X_{\alpha_0}\times_SX_{\alpha_0}|)
      \ \exists(R_x\text{ a ring})\
      \exists\bigl(e_x\colon\operatorname{Spec}(R_x)\to
        X_{\alpha_0}\times_SX_{\alpha_0}\bigr):\
      x\in\operatorname{Im}(|e_x|)
      \ \wedge\ \operatorname{IsOpenImmersion}(e_x),\\
    \operatorname{QuasiCompact}(X_{\alpha_0}\times_SX_{\alpha_0})
      \Longrightarrow\mathscr A_{X/S}\neq\varnothing .
  \end{gathered}                                                     \tag{175c}
  \]
  By \Cref{def:nagata-choice-notation}, define
  \[
    \bigl(m_{X/S},(R^{X/S}_r,e^{X/S}_r)_{r=1}^{m_{X/S}}\bigr)
      :=\operatorname{ch}(\mathscr A_{X/S}),\qquad
    P^{X/S}_r:=\operatorname{Spec}(R^{X/S}_r).                     \tag{175d}
  \]
  Membership of the selected tuple in \(\mathscr A_{X/S}\) is
  \[
  \begin{gathered}
    m_{X/S}\geq1,\qquad
    e^{X/S}_r\colon P^{X/S}_r\to
      X_{\alpha^{X/S}_0}\times_SX_{\alpha^{X/S}_0},\\
    \operatorname{IsOpenImmersion}(e^{X/S}_r)
      \quad(1\leq r\leq m_{X/S}),\\
    |X_{\alpha^{X/S}_0}\times_SX_{\alpha^{X/S}_0}|
      =\bigcup_{r=1}^{m_{X/S}}\operatorname{Im}(|e^{X/S}_r|).
  \end{gathered}                                                     \tag{176}
  \]
  Apply \Cref{def:nagata-fixed-base-separated-stage} to the exact tuple
  $(\mathsf X_{X/S},\alpha^{X/S}_0,m_{X/S},
    (P^{X/S}_r,e^{X/S}_r)_r)$.  Denote its witness tuple (159) by
  \[
    \mathcal W^{X/S}_{\mathrm{sep}}
      =\bigl(\alpha^{X/S}_0,m_{X/S},
        (P^{X/S}_r,e^{X/S}_r)_r,
        (\alpha^{X/S}_r)_r,F^{X/S}_{\mathrm{sep}},
        \alpha^{X/S}_{\mathrm{sep}}\bigr),                       \tag{177}
  \]
  where, exactly as in (158),
  \[
    F^{X/S}_{\mathrm{sep}}
      =\{\alpha^{X/S}_0,\alpha^{X/S}_1,\ldots,
          \alpha^{X/S}_{m_{X/S}}\},
    \qquad
    \alpha^{X/S}_{\mathrm{sep}}
      =u_{X/S}(F^{X/S}_{\mathrm{sep}}).                          \tag{178}
  \]
  Define
  \[
  \begin{gathered}
    \qquad
    X^{\mathrm{fp}}=X_{\alpha^{X/S}_{\mathrm{sep}}},
    \qquad
    i^{\mathrm{fp}}_{X/S}
      =i_{\alpha^{X/S}_{\mathrm{sep}}}\colon X\hookrightarrow X^{\mathrm{fp}},
    \\
    q^{\mathrm{fp}}_{X/S}
      =f_{\alpha^{X/S}_{\mathrm{sep}}}\colon X^{\mathrm{fp}}\to S,
    \qquad
    q^{\mathrm{fp}}_{X/S}\circ i^{\mathrm{fp}}_{X/S}=f.
  \end{gathered}                                                 \tag{179}
  \]
  The last equality in (179) is the instance of (140) at
  $\alpha=\alpha^{X/S}_{\mathrm{sep}}$.  Define the envelope pair by
  \[
    \mathsf{Env}_{\mathrm{fp}}(X/S)
      =(X^{\mathrm{fp}},i^{\mathrm{fp}}_{X/S}).                  \tag{180}
  \]
\end{definition}

\begin{proposition}[Properties of the separated finite-presentation envelope,
  {\cite[\href{https://stacks.math.columbia.edu/tag/01ZJ}{Tag 01ZJ}]{stacks-project}}]
  \label{prop:nagata-finite-type-closed-in-fp}
  \uses{def:nagata-finite-type-closed-in-fp}
  \lean{}

  The morphism
  $q^{\mathrm{fp}}_{X/S}\colon X^{\mathrm{fp}}\to S$ of (179) is separated
  and of finite presentation, and
  $i^{\mathrm{fp}}_{X/S}\colon X\to X^{\mathrm{fp}}$ is a closed immersion
  satisfying $q^{\mathrm{fp}}_{X/S}\circ i^{\mathrm{fp}}_{X/S}=f$.
\end{proposition}

\begin{proof}
  \uses{lemma:nagata-finite-type-fp-model-system,
    lemma:nagata-fixed-base-separated-eventually}
  Put $\alpha:=\alpha^{X/S}_{\mathrm{sep}}$.  Equations (169), (179), and
  \Cref{lemma:nagata-finite-type-fp-model-system} are the assertions
  \[
    q^{\mathrm{fp}}_{X/S}
      =f_\alpha\colon X^{\mathrm{fp}}=X_\alpha\longrightarrow S
      \text{ of finite presentation},
    \qquad
    i^{\mathrm{fp}}_{X/S}=i_\alpha\colon X\hookrightarrow X_\alpha
      \text{ a closed immersion}.
  \]
  Equations (177)--(178) identify $\alpha$ with
  $\alpha_{\mathrm{sep}}$ in (158).  The conclusion of
  \Cref{lemma:nagata-fixed-base-separated-eventually}, instantiated at this
  $\alpha$, is
  \[
    q^{\mathrm{fp}}_{X/S}=f_\alpha\colon X^{\mathrm{fp}}\to S
      \quad\text{is separated}.
  \]
  Finally, the last equality in (179) is
  \[
    q^{\mathrm{fp}}_{X/S}\circ i^{\mathrm{fp}}_{X/S}=f.
  \]
\end{proof}


\subsection{Descent of relative finite-presentation data}

\begin{definition}[Limit datum for relative descent]
  \label{def:nagata-relative-limit-datum}
  \uses{def:nagata-limit-datum,
    lemma:nagata-limit-datum}

  Let $S$ be a scheme.  A candidate relative limit datum over $S$ is a tuple
  \[
    \mathsf B
      :=\bigl(
        I,\leq,(S_\lambda)_{\lambda\in I},
        (t_{\mu\lambda})_{\lambda\leq\mu},
        S,(s_\lambda)_{\lambda\in I},
        L,(\pi_\lambda)_{\lambda\in I},\kappa_{\mathsf B}
        \bigr),                                                    \tag{RL0}
  \]
  with component types
  \[
  \begin{gathered}
    I\text{ is a set},\qquad {\leq}\ \subseteq I\times I,\qquad
    S_\lambda\text{ is a scheme}\quad(\lambda\in I),\qquad
    L\text{ is a scheme},\\
    t_{\mu\lambda}\colon S_\mu\to S_\lambda
      \quad(\lambda,\mu\in I,\ \lambda\leq\mu),\qquad
    s_\lambda\colon S\to S_\lambda\quad(\lambda\in I),\\
    \pi_\lambda\colon L\to S_\lambda\quad(\lambda\in I),\qquad
    \kappa_{\mathsf B}\colon S\to L.
  \end{gathered}
  \]
  Define the first three structural predicates by
  \begin{align}
    \mathsf{RL}_1(\mathsf B)
      &:\Longleftrightarrow
      \bigl(I\neq\varnothing\bigr)\ \wedge\
      \bigl(\forall\lambda,\mu\in I\bigr)
      \bigl(\exists\nu\in I\bigr)
        \bigl((\lambda\leq\nu)\wedge(\mu\leq\nu)\bigr),             \tag{RL1}\\
    \mathsf{RL}_2(\mathsf B)
      &:\Longleftrightarrow
      \Bigl[
        \bigl(\forall\lambda\in I\bigr)\
          t_{\lambda\lambda}=\operatorname{id}_{S_\lambda}
      \Bigr]
      \ \wedge\
      \Bigl[
        \bigl(\forall\lambda,\mu,\nu\in I\bigr)
        \Bigl(
          \bigl((\lambda\leq\mu)\wedge(\mu\leq\nu)\bigr)
          \Longrightarrow
          t_{\mu\lambda}\circ t_{\nu\mu}=t_{\nu\lambda}
        \Bigr)
      \Bigr],                                                       \tag{RL2}\\
    \mathsf{RL}_3(\mathsf B)
      &:\Longleftrightarrow
      \bigl(\forall\lambda,\mu\in I\bigr)
      \Bigl(
        \bigl(\lambda\leq\mu\bigr)
        \Longrightarrow t_{\mu\lambda}\circ s_\mu=s_\lambda
      \Bigr).                                                       \tag{RL3}
  \end{align}
  For every scheme $Z$, define the set of cones from $Z$ to the typed diagram
  in (RL0) by
  \[
    \mathcal C_{\mathsf B}(Z)
      :=\left\{
        (u_\lambda)_{\lambda\in I}
          \in\prod_{\lambda\in I}
            \operatorname{Hom}_{\mathbf{Sch}}(Z,S_\lambda)
        \ \middle|\
        \begin{array}{c}
          (\forall\lambda,\mu\in I)\
          \bigl((\lambda\leq\mu)\Longrightarrow
            t_{\mu\lambda}\circ u_\mu=u_\lambda\bigr)
        \end{array}
      \right\},
  \]
  and define
  \[
    \Phi_{\mathsf B,Z}\colon
      \operatorname{Hom}_{\mathbf{Sch}}(Z,L)
        \longrightarrow
        \prod_{\lambda\in I}
          \operatorname{Hom}_{\mathbf{Sch}}(Z,S_\lambda),
    \qquad
    h\longmapsto(\pi_\lambda\circ h)_{\lambda\in I}.
  \]
  Define the remaining two structural predicates by
  \begin{align}
    \mathsf{RL}_4(\mathsf B)
      &:\Longleftrightarrow
      \Bigl[
        (\forall\lambda,\mu\in I)\
        \bigl((\lambda\leq\mu)\Longrightarrow
          t_{\mu\lambda}\circ\pi_\mu=\pi_\lambda\bigr)
      \Bigr]
      \ \wedge\
      \Bigl[
        (\forall Z)\
        \Bigl(
          Z\text{ is a scheme}\Longrightarrow
          \Bigl(
          \Bigl[
            \Phi_{\mathsf B,Z}
              \bigl(\operatorname{Hom}_{\mathbf{Sch}}(Z,L)\bigr)
              =\mathcal C_{\mathsf B}(Z)
          \Bigr]
          \ \wedge\
          \Bigl[
            (\forall h,h'\in
              \operatorname{Hom}_{\mathbf{Sch}}(Z,L))\
            \bigl(
              \Phi_{\mathsf B,Z}(h)=\Phi_{\mathsf B,Z}(h')
              \Longrightarrow h=h'
            \bigr)
          \Bigr]
          \Bigr)
        \Bigr)
      \Bigr],                                                       \tag{RL4}\\
    \mathsf{RL}_5(\mathsf B)
      &:\Longleftrightarrow
      \Bigl[\kappa_{\mathsf B}\colon S\xrightarrow{\sim}L\Bigr]
      \ \wedge\
      \Bigl[
        (\forall\lambda\in I)\
          \pi_\lambda\circ\kappa_{\mathsf B}=s_\lambda
      \Bigr].                                                       \tag{RL5}
  \end{align}
  Define
  \[
    \operatorname{RelLim}(S)
      :=\left\{
      \mathsf B=
        \bigl(I,\leq,(S_\lambda)_{\lambda\in I},
          (t_{\mu\lambda})_{\lambda\leq\mu},
          S,(s_\lambda)_{\lambda\in I},
          L,(\pi_\lambda)_{\lambda\in I},\kappa_{\mathsf B}\bigr)
      \ \middle|\
      \begin{array}{l}
        I\text{ is a set},\quad {\leq}\ \subseteq I\times I,\\
        (\forall\lambda\in I)\
          S_\lambda\text{ is a scheme},\quad L\text{ is a scheme},\\
        (\forall\lambda,\mu\in I)\
          (\lambda\leq\mu)\Longrightarrow
          t_{\mu\lambda}\colon S_\mu\to S_\lambda,\\
        (\forall\lambda\in I)\
          s_\lambda\colon S\to S_\lambda,\\
        (\forall\lambda\in I)\
          \pi_\lambda\colon L\to S_\lambda,\quad
          \kappa_{\mathsf B}\colon S\to L,\\
        (\forall\lambda\in I)\ \lambda\leq\lambda,\\
        (\forall\lambda,\mu,\nu\in I)\
          ((\lambda\leq\mu)\wedge(\mu\leq\nu))
            \Longrightarrow\lambda\leq\nu,\\
        (\forall\lambda\in I)\
          \bigl(
            \operatorname{QuasiCompact}(S_\lambda)\ \wedge\
            \operatorname{QuasiSeparated}(S_\lambda)
          \bigr),\\
        (\forall\lambda,\mu\in I)\
          (\lambda\leq\mu)\Longrightarrow
          \operatorname{IsAffine}(t_{\mu\lambda}),\\
        \mathsf{RL}_1(\mathsf B)\ \wedge\
        \mathsf{RL}_2(\mathsf B)\ \wedge\
        \mathsf{RL}_3(\mathsf B)\ \wedge\
        \mathsf{RL}_4(\mathsf B)\ \wedge\
        \mathsf{RL}_5(\mathsf B)
      \end{array}
      \right\}.                                                    \tag{RL5a}
  \]

  \paragraph{Mathlib relation.}
  Predicate $\mathsf{RL}_4$ is precisely a
  \texttt{\detokenize{CategoryTheory.Limits.Cone}} together with
  \texttt{\detokenize{CategoryTheory.Limits.IsLimit}}, expanded into
  projection equations, existence of lifts, and uniqueness.  Predicate
  $\mathsf{RL}_5$ adds the comparison with $S$.  Thus
  $\operatorname{RelLim}(S)$ is a blueprint-local structured wrapper around
  \Cref{def:nagata-limit-datum}, not a general Mathlib limit-existence
  theorem for schemes.
\end{definition}

\begin{definition}[Base-change notation for relative descent]
  \label{def:nagata-relative-base-change-notation}
  \uses{def:nagata-relative-limit-datum}

  Let \(\mathsf B\in\operatorname{RelLim}(S)\).  For
  \(\lambda\in I\) and \(r_\lambda\colon Y_\lambda\to S_\lambda\),
  define
  \[
  \begin{gathered}
    Y_{\lambda,S}^{r_\lambda}
      :=Y_\lambda\times_{S_\lambda,s_\lambda}S,\\
    p_{Y_\lambda}^{\lambda,r_\lambda}\colon
       Y_{\lambda,S}^{r_\lambda}\to Y_\lambda,\qquad
    p_S^{\lambda,r_\lambda}\colon
       Y_{\lambda,S}^{r_\lambda}\to S,\\
    r_\lambda\circ p_{Y_\lambda}^{\lambda,r_\lambda}
      =s_\lambda\circ p_S^{\lambda,r_\lambda}.
  \end{gathered}                                                    \tag{RL6}
  \]
  If \(\lambda\leq\mu\), define
  \[
    Y_\mu^{(\lambda,r_\lambda)}
      :=Y_\lambda\times_{S_\lambda,t_{\mu\lambda}}S_\mu,\qquad
    r_\mu^{(\lambda,r_\lambda)}
      :=\operatorname{pr}_{S_\mu}\colon
         Y_\mu^{(\lambda,r_\lambda)}\to S_\mu.                    \tag{RL7}
  \]
  The fibre-product equations in (RL6)--(RL7) and the equation
  \(t_{\mu\lambda}\circ s_\mu=s_\lambda\) in (RL3) give the
  compatibility chain
  \[
  \begin{aligned}
    &r_\lambda\circ\operatorname{pr}_{Y_\lambda}\circ
      p_{Y_\mu^{(\lambda,r_\lambda)}}^{
        \mu,r_\mu^{(\lambda,r_\lambda)}}\\
    &\quad=
      t_{\mu\lambda}\circ\operatorname{pr}_{S_\mu}\circ
      p_{Y_\mu^{(\lambda,r_\lambda)}}^{
        \mu,r_\mu^{(\lambda,r_\lambda)}}\\
    &\quad=
      t_{\mu\lambda}\circ r_\mu^{(\lambda,r_\lambda)}\circ
      p_{Y_\mu^{(\lambda,r_\lambda)}}^{
        \mu,r_\mu^{(\lambda,r_\lambda)}}\\
    &\quad=
      t_{\mu\lambda}\circ s_\mu\circ
      p_S^{\mu,r_\mu^{(\lambda,r_\lambda)}}
      =s_\lambda\circ
      p_S^{\mu,r_\mu^{(\lambda,r_\lambda)}}.
  \end{aligned}
  \]
  Define
  \[
    a_{\mu\lambda}^{Y_\lambda,r_\lambda}\colon
      \bigl(Y_\mu^{(\lambda,r_\lambda)}\bigr)_{
        \mu,S}^{r_\mu^{(\lambda,r_\lambda)}}
        \longrightarrow Y_{\lambda,S}^{r_\lambda}                 \tag{RL8}
  \]
  as the unique fibre-product morphism satisfying
  \[
  \begin{aligned}
    p_{Y_\lambda}^{\lambda,r_\lambda}
      \circ a_{\mu\lambda}^{Y_\lambda,r_\lambda}
      &=
      \operatorname{pr}_{Y_\lambda}\circ
      p_{Y_\mu^{(\lambda,r_\lambda)}}^{
        \mu,r_\mu^{(\lambda,r_\lambda)}},\\
    p_S^{\lambda,r_\lambda}
      \circ a_{\mu\lambda}^{Y_\lambda,r_\lambda}
      &=
      p_S^{\mu,r_\mu^{(\lambda,r_\lambda)}}.
  \end{aligned}                                                     \tag{RL9}
  \]
  Before defining the first component of the inverse, equations (RL3) and
  (RL6) give
  \[
    r_\lambda\circ p_{Y_\lambda}^{\lambda,r_\lambda}
      =s_\lambda\circ p_S^{\lambda,r_\lambda}
      =t_{\mu\lambda}\circ s_\mu\circ
        p_S^{\lambda,r_\lambda}.
  \]
  Hence the following fibre-product morphism is defined:
  \[
    \widehat b_{\lambda\mu}:=
      (p_{Y_\lambda}^{\lambda,r_\lambda},
        s_\mu\circ p_S^{\lambda,r_\lambda})
      \colon Y_{\lambda,S}^{r_\lambda}
        \longrightarrow Y_\mu^{(\lambda,r_\lambda)}.
  \]
  Its second projection and (RL7) give the compatibility equation
  \[
    r_\mu^{(\lambda,r_\lambda)}\circ\widehat b_{\lambda\mu}
      =\operatorname{pr}_{S_\mu}\circ\widehat b_{\lambda\mu}
      =s_\mu\circ p_S^{\lambda,r_\lambda}.
  \]
  Define
  \[
    b_{\lambda\mu}^{Y_\lambda,r_\lambda}:=
      (\widehat b_{\lambda\mu},p_S^{\lambda,r_\lambda})
      \colon Y_{\lambda,S}^{r_\lambda}\longrightarrow
      (Y_\mu^{(\lambda,r_\lambda)})_{
        \mu,S}^{r_\mu^{(\lambda,r_\lambda)}}.
                                                                    \tag{RL10}
  \]
  Equations (RL9)--(RL10) are the component identities
  \[
  \begin{aligned}
    p_{Y_\lambda}^{\lambda,r_\lambda}\circ
      a_{\mu\lambda}^{Y_\lambda,r_\lambda}\circ
      b_{\lambda\mu}^{Y_\lambda,r_\lambda}
      &=p_{Y_\lambda}^{\lambda,r_\lambda},\\
    p_S^{\lambda,r_\lambda}\circ
      a_{\mu\lambda}^{Y_\lambda,r_\lambda}\circ
      b_{\lambda\mu}^{Y_\lambda,r_\lambda}
      &=p_S^{\lambda,r_\lambda},
  \end{aligned}
  \]
  Fibre-product uniqueness applied to these two equalities is
  \[
    a_{\mu\lambda}^{Y_\lambda,r_\lambda}\circ
      b_{\lambda\mu}^{Y_\lambda,r_\lambda}
      =\operatorname{id}_{Y_{\lambda,S}^{r_\lambda}}.             \tag{RL11}
  \]
  The other composite has projection chains
  \[
  \begin{aligned}
    \operatorname{pr}_{Y_\lambda}\circ\widehat b_{\lambda\mu}\circ
      a_{\mu\lambda}^{Y_\lambda,r_\lambda}
      &=\operatorname{pr}_{Y_\lambda}\circ
        p_{Y_\mu^{(\lambda,r_\lambda)}}^
          {\mu,r_\mu^{(\lambda,r_\lambda)}},\\
    \operatorname{pr}_{S_\mu}\circ\widehat b_{\lambda\mu}\circ
      a_{\mu\lambda}^{Y_\lambda,r_\lambda}
      &=s_\mu\circ p_S^{\mu,r_\mu^{(\lambda,r_\lambda)}}
       =\operatorname{pr}_{S_\mu}\circ
        p_{Y_\mu^{(\lambda,r_\lambda)}}^
          {\mu,r_\mu^{(\lambda,r_\lambda)}},\\
    p_S^{\mu,r_\mu^{(\lambda,r_\lambda)}}\circ
      b_{\lambda\mu}^{Y_\lambda,r_\lambda}\circ
      a_{\mu\lambda}^{Y_\lambda,r_\lambda}
      &=p_S^{\mu,r_\mu^{(\lambda,r_\lambda)}}.
  \end{aligned}
  \]
  Two applications of fibre-product uniqueness are the equality
  \[
    b_{\lambda\mu}^{Y_\lambda,r_\lambda}\circ
      a_{\mu\lambda}^{Y_\lambda,r_\lambda}
      =\operatorname{id}_{
        (Y_\mu^{(\lambda,r_\lambda)})_{
          \mu,S}^{r_\mu^{(\lambda,r_\lambda)}}}.                  \tag{RL12}
  \]
  Equations (RL11)--(RL12) give
  \[
    a_{\mu\lambda}^{Y_\lambda,r_\lambda}\colon
      \bigl(Y_\mu^{(\lambda,r_\lambda)}\bigr)_{
        \mu,S}^{r_\mu^{(\lambda,r_\lambda)}}
        \xrightarrow{\sim}Y_{\lambda,S}^{r_\lambda},
    \qquad
    (a_{\mu\lambda}^{Y_\lambda,r_\lambda})^{-1}
      =b_{\lambda\mu}^{Y_\lambda,r_\lambda}.
  \]

  \paragraph{Mathlib relation.}
  The objects and projections in (RL6)--(RL7) are
  \texttt{\detokenize{CategoryTheory.Limits.pullback}},
  \texttt{\detokenize{CategoryTheory.Limits.pullback.fst}}, and
  \texttt{\detokenize{CategoryTheory.Limits.pullback.snd}}.  The comparison
  (RL8)--(RL12) specializes
  \texttt{\detokenize{CategoryTheory.Limits.pullbackAssoc}} and then uses
  the pullback-of-an-isomorphism identification supplied by
  \texttt{\detokenize{CategoryTheory.Limits.pullbackConeOfLeftIso}}.
\end{definition}

\subsection{Eventual separatedness of the descended model}

\begin{definition}[Separated stage of a finite-presentation morphism]
  \label{def:nagata-separated-descent-stage}
  \uses{def:nagata-choice-notation,
    def:nagata-relative-limit-datum,
    def:nagata-relative-base-change-notation}

  Let $S$ and $X$ be schemes, let
  \(\mathsf B\in\operatorname{RelLim}(S)\), and let
  \[
    q\colon X\longrightarrow S
  \]
  satisfy
  \[
    \operatorname{Separated}(q)\wedge
      \operatorname{FinitePresentation}(q).
  \]
  Define
  \[
  \begin{split}
    \mathscr D_{\mathrm{sep}}(q;\mathsf B)
      :=\Bigl\{&
       (\lambda,Y_\lambda,r_\lambda,\epsilon,e)\ \Bigm|\;
       \lambda\in I,\quad
       Y_\lambda\text{ is a scheme},\quad
       r_\lambda\colon Y_\lambda\to S_\lambda,\\
       &\operatorname{Separated}(r_\lambda)
          \wedge\operatorname{FinitePresentation}(r_\lambda),\\
       &\epsilon\colon
          Y_{\lambda,S}^{r_\lambda}\xrightarrow{\sim}X,\quad
       e\colon q\circ\epsilon=p_S^{\lambda,r_\lambda}
       \Bigr\}.
  \end{split}                                                       \tag{SD1}
  \]

  Relative finite-presentation descent
  \cite[\href{https://stacks.math.columbia.edu/tag/01ZM}{Tag 01ZM}]
       {stacks-project} is the quantified statement
  \[
  \begin{split}
    \exists\lambda_0\in I\ \exists
      \bigl(Y_{\lambda_0}^0\text{ a scheme}\bigr)\ \exists
      \bigl(r_{\lambda_0}^0\colon
        Y_{\lambda_0}^0\to S_{\lambda_0}\bigr)\ \exists
      \left(
        \epsilon_0\colon
          (Y_{\lambda_0}^0)_{\lambda_0,S}^{r_{\lambda_0}^0}
            \xrightarrow{\sim}X
      \right):\quad&
      \operatorname{FinitePresentation}(r_{\lambda_0}^0)\\
      &{}\wedge\
      \bigl[
        q\circ\epsilon_0
          =p_S^{\lambda_0,r_{\lambda_0}^0}
      \bigr].
  \end{split}                                                       \tag{SD2}
  \]
  Fix witnesses \((\lambda_0,Y_{\lambda_0}^0,
  r_{\lambda_0}^0,\epsilon_0)\) of (SD2); their complete list of types and
  equations is
  \[
  \begin{gathered}
    \lambda_0\in I,\qquad
    Y_{\lambda_0}^0\text{ is a scheme},\qquad
    r_{\lambda_0}^0\colon Y_{\lambda_0}^0\to S_{\lambda_0},\\
    \operatorname{FinitePresentation}(r_{\lambda_0}^0),\qquad
    \epsilon_0\colon
      (Y_{\lambda_0}^0)_{\lambda_0,S}^{r_{\lambda_0}^0}
        \xrightarrow{\sim}X,\qquad
    q\circ\epsilon_0=p_S^{\lambda_0,r_{\lambda_0}^0}.
  \end{gathered}
  \]
  For every \(\mu\in I\) satisfying \(\lambda_0\leq\mu\), define
  \[
  \begin{gathered}
    Y_\mu^0
      :=Y_\mu^{(\lambda_0,r_{\lambda_0}^0)}
      =Y_{\lambda_0}^0
          \times_{S_{\lambda_0},t_{\mu\lambda_0}}S_\mu
      \quad\text{(a scheme)},\\
    r_\mu^0
      :=r_\mu^{(\lambda_0,r_{\lambda_0}^0)}
      =\operatorname{pr}_{S_\mu}\colon Y_\mu^0\to S_\mu,\qquad
    a_\mu
      :=a_{\mu\lambda_0}^{Y_{\lambda_0}^0,r_{\lambda_0}^0}
      \colon (Y_\mu^0)_{\mu,S}^{r_\mu^0}
        \xrightarrow{\sim}
        (Y_{\lambda_0}^0)_{\lambda_0,S}^{r_{\lambda_0}^0},\\
    \epsilon_\mu:=\epsilon_0\circ a_\mu
      \colon (Y_\mu^0)_{\mu,S}^{r_\mu^0}\xrightarrow{\sim}X.
  \end{gathered}                                                    \tag{SD3}
  \]
  Equations (RL9), (SD2), and (SD3) form the chain
  \[
  \begin{aligned}
    q\circ\epsilon_\mu
      &=q\circ(\epsilon_0\circ a_\mu)\\
      &=(q\circ\epsilon_0)\circ a_\mu\\
      &=p_S^{\lambda_0,r_{\lambda_0}^0}\circ a_\mu\\
      &=p_S^{\mu,r_\mu^0}.
  \end{aligned}                                                     \tag{SD4}
  \]
  Stability under base change is the implication
  \[
    \operatorname{FinitePresentation}(r_{\lambda_0}^0)
      \Longrightarrow
    \operatorname{FinitePresentation}(r_\mu^0).                   \tag{SD5}
  \]
  Define the eventual-separatedness witness set
  \[
    \mathscr U_{\mathrm{sep}}
      (q;\mathsf B,\lambda_0,r_{\lambda_0}^0)
      :=\left\{
        \mu\in I\ \middle|\
        \lambda_0\leq\mu\ \wedge\
        \operatorname{Separated}(r_\mu^0)
      \right\}.
                                                                    \tag{SD5a}
  \]
  Eventual separatedness
  \cite[\href{https://stacks.math.columbia.edu/tag/01ZQ}{Tag 01ZQ}]
       {stacks-project}, applied to (SD2)--(SD5), and the hypothesis on $q$
  give the chain
  \[
  \begin{aligned}
    \operatorname{Separated}(q)
      &\Longrightarrow
      \mathscr U_{\mathrm{sep}}
        (q;\mathsf B,\lambda_0,r_{\lambda_0}^0)\neq\varnothing,
      &&\text{by Tag 01ZQ},\\
    \mathscr U_{\mathrm{sep}}
        (q;\mathsf B,\lambda_0,r_{\lambda_0}^0)
      &\neq\varnothing,
      &&\text{by the first line and the hypothesis},\\
    \mu_{\mathrm{sep}}^0
        (q;\mathsf B,\lambda_0,r_{\lambda_0}^0)
      &:=\operatorname{ch}\!\left(
        \mathscr U_{\mathrm{sep}}
          (q;\mathsf B,\lambda_0,r_{\lambda_0}^0)
        \right)
      \in
        \mathscr U_{\mathrm{sep}}
          (q;\mathsf B,\lambda_0,r_{\lambda_0}^0),
      &&\text{by (C0)},\\
    \mu_{\mathrm{sep}}^0
      &:=\mu_{\mathrm{sep}}^0
        (q;\mathsf B,\lambda_0,r_{\lambda_0}^0).
  \end{aligned}                                                     \tag{SD6}
  \]
  Membership
  \(\mu_{\mathrm{sep}}^0\in
    \mathscr U_{\mathrm{sep}}
      (q;\mathsf B,\lambda_0,r_{\lambda_0}^0)\)
  in (SD6), expanded by (SD5a), is
  \[
  \begin{gathered}
    \mu_{\mathrm{sep}}^0\in I,\qquad
    \lambda_0\leq\mu_{\mathrm{sep}}^0,\qquad
    \operatorname{Separated}(r_{\mu_{\mathrm{sep}}^0}^0).
  \end{gathered}
  \]
  Equations (SD2) and (SD5), evaluated at
  \(\mu=\mu_{\mathrm{sep}}^0\), form the implication chain
  \[
  \begin{gathered}
    \operatorname{FinitePresentation}(r_{\lambda_0}^0),\\
    \operatorname{FinitePresentation}(r_{\lambda_0}^0)
      \Longrightarrow
      \operatorname{FinitePresentation}(r_{\mu_{\mathrm{sep}}^0}^0),\\
    \operatorname{FinitePresentation}(r_{\mu_{\mathrm{sep}}^0}^0).
  \end{gathered}
  \]
  Substitution of \(\mu=\mu_{\mathrm{sep}}^0\) into (SD3)--(SD4), together
  with this last predicate and the expanded membership after (SD6), gives the
  explicit conjunction
  \[
  \begin{gathered}
    \mu_{\mathrm{sep}}^0\in I,\qquad
    Y_{\mu_{\mathrm{sep}}^0}^0\text{ is a scheme},\qquad
    r_{\mu_{\mathrm{sep}}^0}^0\colon
      Y_{\mu_{\mathrm{sep}}^0}^0\to S_{\mu_{\mathrm{sep}}^0},\\
    \operatorname{Separated}(r_{\mu_{\mathrm{sep}}^0}^0)
      \ \wedge\
    \operatorname{FinitePresentation}(r_{\mu_{\mathrm{sep}}^0}^0),\\
    \epsilon_{\mu_{\mathrm{sep}}^0}\colon
      (Y_{\mu_{\mathrm{sep}}^0}^0)_{
        \mu_{\mathrm{sep}}^0,S}^{r_{\mu_{\mathrm{sep}}^0}^0}
        \xrightarrow{\sim}X,\\
    q\circ\epsilon_{\mu_{\mathrm{sep}}^0}
      =p_S^{\mu_{\mathrm{sep}}^0,r_{\mu_{\mathrm{sep}}^0}^0}.
  \end{gathered}
  \]
  Substitution of this conjunction into the defining set equality (SD1) is
  \[
    \bigl(
      \mu_{\mathrm{sep}}^0,Y_{\mu_{\mathrm{sep}}^0}^0,
      r_{\mu_{\mathrm{sep}}^0}^0,\epsilon_{\mu_{\mathrm{sep}}^0},
      (q\circ\epsilon_{\mu_{\mathrm{sep}}^0}
        =p_S^{\mu_{\mathrm{sep}}^0,r_{\mu_{\mathrm{sep}}^0}^0})
    \bigr)
      \in\mathscr D_{\mathrm{sep}}(q;\mathsf B).                  \tag{SD7}
  \]
  Equation (SD7) supplies the explicit element in the nonemptiness assertion
  \[
    \mathscr D_{\mathrm{sep}}(q;\mathsf B)\neq\varnothing.        \tag{SD8}
  \]
  By \Cref{def:nagata-choice-notation}, define
  \[
  \begin{split}
    \bigl(&
      \lambda_{\mathrm{sep}}(q;\mathsf B),
      X_{\mathrm{sep}}(q;\mathsf B),
      q_{\mathrm{sep}}(q;\mathsf B),
      \epsilon_{\mathrm{sep}}(q;\mathsf B),
      e_{\mathrm{sep}}(q;\mathsf B)
    \bigr)
      :=\operatorname{ch}
        \bigl(\mathscr D_{\mathrm{sep}}(q;\mathsf B)\bigr).
  \end{split}                                                       \tag{SD9}
  \]
\end{definition}

\begin{lemma}[Property of the selected separated stage]
  \label{lemma:nagata-separated-eventually}
  \uses{def:nagata-separated-descent-stage}

  Let $S$ and $X$ be schemes, let
  \(\mathsf B\in\operatorname{RelLim}(S)\), and let
  \(q\colon X\to S\) satisfy
  \[
    \operatorname{Separated}(q)\ \wedge\
    \operatorname{FinitePresentation}(q).
  \]
  For the five components of (SD9),
  \[
  \begin{gathered}
    \lambda_{\mathrm{sep}}(q;\mathsf B)\in I,\qquad
    X_{\mathrm{sep}}(q;\mathsf B)\text{ is a scheme},\\
    q_{\mathrm{sep}}(q;\mathsf B)\colon
      X_{\mathrm{sep}}(q;\mathsf B)
        \to S_{\lambda_{\mathrm{sep}}(q;\mathsf B)},\\
    \operatorname{Separated}
      (q_{\mathrm{sep}}(q;\mathsf B))
      \wedge
    \operatorname{FinitePresentation}
      (q_{\mathrm{sep}}(q;\mathsf B)),\\
    \epsilon_{\mathrm{sep}}(q;\mathsf B)\colon
      \bigl(X_{\mathrm{sep}}(q;\mathsf B)\bigr)_{
        \lambda_{\mathrm{sep}}(q;\mathsf B),S}^{
          q_{\mathrm{sep}}(q;\mathsf B)}
        \xrightarrow{\sim}X,\\
    q\circ\epsilon_{\mathrm{sep}}(q;\mathsf B)
      =p_S^{
        \lambda_{\mathrm{sep}}(q;\mathsf B),
        q_{\mathrm{sep}}(q;\mathsf B)}.
  \end{gathered}                                                    \tag{SD10}
  \]
\end{lemma}

\begin{proof}
  Equations (SD8), (SD9), and
  \Cref{def:nagata-choice-notation} establish
  \[
  \begin{split}
    \bigl(&
      \lambda_{\mathrm{sep}}(q;\mathsf B),
      X_{\mathrm{sep}}(q;\mathsf B),
      q_{\mathrm{sep}}(q;\mathsf B),
      \epsilon_{\mathrm{sep}}(q;\mathsf B),
      e_{\mathrm{sep}}(q;\mathsf B)
    \bigr)
    \in\mathscr D_{\mathrm{sep}}(q;\mathsf B).
  \end{split}                                                       \tag{SD11}
  \]
  Substitution of the tuple in (SD11) into the set equality (SD1) is
  \[
  \begin{split}
    \text{\rm(SD11)}
    \quad\Longleftrightarrow\quad&
    \lambda_{\mathrm{sep}}(q;\mathsf B)\in I\\
    &{}\wedge\
    \left[
      X_{\mathrm{sep}}(q;\mathsf B)\text{ is a scheme}
    \right]\\
    &{}\wedge\
    \left[
      q_{\mathrm{sep}}(q;\mathsf B)\colon
      X_{\mathrm{sep}}(q;\mathsf B)
        \to S_{\lambda_{\mathrm{sep}}(q;\mathsf B)}
    \right]\\
    &{}\wedge\
    \left[
      \operatorname{Separated}(q_{\mathrm{sep}}(q;\mathsf B))
      \ \wedge\
      \operatorname{FinitePresentation}(q_{\mathrm{sep}}(q;\mathsf B))
    \right]\\
    &{}\wedge\
    \left[
      \epsilon_{\mathrm{sep}}(q;\mathsf B)\colon
      \bigl(X_{\mathrm{sep}}(q;\mathsf B)\bigr)_{
        \lambda_{\mathrm{sep}}(q;\mathsf B),S}^{
          q_{\mathrm{sep}}(q;\mathsf B)}
        \xrightarrow{\sim}X
    \right]\\
    &{}\wedge\
    \left[
      q\circ\epsilon_{\mathrm{sep}}(q;\mathsf B)
        =p_S^{
          \lambda_{\mathrm{sep}}(q;\mathsf B),
          q_{\mathrm{sep}}(q;\mathsf B)}
    \right].
  \end{split}                                                       \tag{SD12}
  \]
  The forward direction of (SD12), evaluated at (SD11), is the
  conjunction (SD10).
\end{proof}


\section{Lean interfaces and implementation order}

The basic bundled structure should expose all construction data rather than
only a proposition.

\begin{verbatim}
structure Compactification {X S : Scheme} (f : X -> S) where
  Xbar : Scheme
  j : X -> Xbar
  p : Xbar -> S
  fac : j >> p = f
  isOpenImmersion_j : IsOpenImmersion j
  quasiCompact_j : QuasiCompact j
  isProper_p : IsProper p

def Compactifiable {X S : Scheme} (f : X -> S) : Prop :=
  Nonempty (Compactification f)
\end{verbatim}

The implementation order forced by the dependency graph is:
\begin{enumerate}
  \item compactifications, scheme-theoretic images, restriction to opens, and
    closed-subscheme inheritance;
  \item admissible blow-ups, strict transforms, extension of centers, and
    domination by product ideals;
  \item boundary separation and elimination of indeterminacy;
  \item common modifications, graph extension, and separated gluing;
  \item the dense-source valuative criterion and the two-open theorem;
  \item affine compactification and the Noetherian theorem;
  \item finite-presentation reduction, Noetherian approximation, and the final
    assembly.
\end{enumerate}

Before any declaration receives a \verb|\leanok| marker, the following points
must be checked in Lean:
\begin{enumerate}
  \item every displayed composite equality is stated using the actual chosen
    pullback and gluing morphisms;
  \item every scheme-theoretic image is formed under the required
    quasi-compactness hypotheses;
  \item restriction of a scheme-theoretic image to an open target is available
    as an equality or a specified isomorphism of closed subschemes;
  \item every admissible blow-up carries a chosen lift and chosen isomorphism
    over the preserved open;
  \item the boundary-separation blow-ups preserve $V_i$, not only $U_i$;
  \item the common-modification proof uses properness of
    $X'_{12}\to X_1\times_SX_2$ and separatedness of the new product over the
    old product; it does not incorrectly assume the projections
    $X_{12}\to X_i$ are proper in the general common-modification lemma;
  \item the dense-source valuative criterion is invoked with uniqueness, not
    with existence alone;
  \item the pair of valuation extensions factors scheme-theoretically through
    $X_{12}$ by injectivity of $A\hookrightarrow K$;
  \item the finite affine cover has dense total intersection, and the induction
    uses density of each partial overlap;
  \item finite presentation, the base-change identification, and separatedness
    all descend at one common sufficiently large index.
\end{enumerate}

\section{Explicit final assembly}

In this final section let $S$ be quasi-compact and quasi-separated and let
$f\colon X\to S$ be separated and of finite type.  Every object used below is
the output of a previously defined construction; no new existential choice is
made inside the final theorem.

\begin{definition}[Finite-presentation envelope used in the final assembly]
  \label{def:nagata-final-fp-envelope}
  \uses{def:nagata-finite-type-closed-in-fp}
  \checked\lean{}

  Define
  \[
    (X_f^{\mathrm{fp}},i_f^{\mathrm{fp}})
      :=\mathsf{Env}_{\mathrm{fp}}(X/S)
      =(X^{\mathrm{fp}},i^{\mathrm{fp}}_{X/S}),
    \qquad
    q_f^{\mathrm{fp}}
      :=q^{\mathrm{fp}}_{X/S}
      \colon X_f^{\mathrm{fp}}\longrightarrow S,
  \]
  using the objects and morphisms in (179)--(180) of
  \Cref{def:nagata-finite-type-closed-in-fp}.  Set
  \[
    \mathsf E_f
      =(X_f^{\mathrm{fp}},i_f^{\mathrm{fp}},q_f^{\mathrm{fp}}).
  \]
\end{definition}

\begin{lemma}[Properties of the final finite-presentation envelope]
\checked
  \label{lemma:nagata-final-fp-envelope}
  \uses{def:nagata-final-fp-envelope}
  \lean{}

  The map $i_f^{\mathrm{fp}}\colon X\hookrightarrow X_f^{\mathrm{fp}}$ is a
  closed immersion over $S$, and
  $q_f^{\mathrm{fp}}\colon X_f^{\mathrm{fp}}\to S$ is separated and of
  finite presentation.  Explicitly,
  \[
    q_f^{\mathrm{fp}}\circ i_f^{\mathrm{fp}}=f.
  \]
\end{lemma}

\begin{proof}
  \uses{prop:nagata-finite-type-closed-in-fp}
  By \Cref{def:nagata-final-fp-envelope},
  \[
    (X_f^{\mathrm{fp}},i_f^{\mathrm{fp}})
      =\mathsf{Env}_{\mathrm{fp}}(X/S)
      =(X^{\mathrm{fp}},i^{\mathrm{fp}}_{X/S}),
    \qquad
    q_f^{\mathrm{fp}}=q^{\mathrm{fp}}_{X/S}
      \colon X_f^{\mathrm{fp}}\to S
  \]
  by \Cref{def:nagata-final-fp-envelope}.  The conclusion of
  \Cref{prop:nagata-finite-type-closed-in-fp}, evaluated at this specified
  pair, is the conjunction
  \[
  \begin{gathered}
    i_f^{\mathrm{fp}}\colon X\hookrightarrow X_f^{\mathrm{fp}}
      \text{ a closed immersion},\\
    q_f^{\mathrm{fp}}\text{ separated and of finite presentation},\\
    q_f^{\mathrm{fp}}\circ i_f^{\mathrm{fp}}=f.
  \end{gathered}
  \]
\end{proof}

\begin{definition}[Noetherian base system used in the final assembly]
  \label{def:nagata-final-noetherian-base-system}
  \uses{def:nagata-noetherian-inverse-system, def:nagata-glued-limit-comparison}
  \checked\lean{}

  Define
  \[
    \mathsf N_f=\mathsf N(S)
      =(I_S,\leq_S,(S_\lambda)_\lambda,
        (t_{\mu\lambda})_{\lambda\leq\mu})
  \]
  together with the morphisms $q_\lambda\colon S\to S_\lambda$ of
  \Cref{def:nagata-glued-limit-comparison}.  For each $\lambda\in I_S$, let
  \[
    \pi_\lambda\colon\varprojlim_{\nu\in I_S}S_\nu\longrightarrow S_\lambda
  \]
  be the limit projection; these morphisms are specified by
  \[
    t_{\mu\lambda}\circ\pi_\mu=\pi_\lambda
      \qquad(\lambda\leq\mu).
  \]
  Define
  \[
    \kappa_S\colon S\longrightarrow\varprojlim_\lambda S_\lambda
  \]
  by the equations
  \[
    \pi_\lambda\circ\kappa_S=q_\lambda,
    \qquad
    t_{\mu\lambda}\circ q_\mu=q_\lambda
      \quad(\lambda\leq\mu).
  \]
\end{definition}

\begin{lemma}[Properties of the final Noetherian base system]
  \label{lemma:nagata-final-noetherian-base-system}
  \uses{def:nagata-final-noetherian-base-system}
  \checked\lean{}

  For the system in
  \Cref{def:nagata-final-noetherian-base-system}, every $S_\lambda$ is
  Noetherian and of finite type over $\mathbf Z$, every
  $t_{\mu\lambda}\colon S_\mu\to S_\lambda$ is affine, and $\kappa_S$ is an
  isomorphism.
\end{lemma}

\begin{proof}
\checked
  \uses{prop:nagata-noetherian-inverse-system}
  The equality
  \[
    \mathsf N_f=\mathsf N(S)
  \]
  is \Cref{def:nagata-final-noetherian-base-system}.  The conclusion of
  \Cref{prop:nagata-noetherian-inverse-system}, after substitution of this
  equality, is the conjunction
  \[
  \begin{gathered}
    S_\lambda\text{ Noetherian and of finite type over }\mathbf Z,\\
    t_{\mu\lambda}\colon S_\mu\to S_\lambda\text{ affine}
      \quad(\lambda\leq\mu),\\
    \kappa_S\colon S\xrightarrow{\sim}\varprojlim_\lambda S_\lambda.
  \end{gathered}
  \]
\end{proof}

\begin{definition}[Separated descended model used in the final assembly]
  \label{def:nagata-final-descended-model}
  \uses{def:nagata-final-fp-envelope,
    def:nagata-final-noetherian-base-system,
    def:nagata-glued-limit-comparison,
    def:nagata-relative-limit-datum,
    def:nagata-separated-descent-stage,
    lemma:nagata-final-fp-envelope,
    lemma:nagata-final-noetherian-base-system}
  \checked\lean{}

  Define the relative limit datum
  \[
    \mathsf B_f
      :=\bigl(
        I_S,\leq,(S_\lambda)_{\lambda\in I_S},
        (t_{\mu\lambda})_{\lambda\leq\mu},
        S,(q_\lambda)_{\lambda\in I_S},
        L_S,(\pi_\lambda)_{\lambda\in I_S},\kappa_S
        \bigr).                                                    \tag{F0}
  \]
  The tuple definitions (85)--(87) of
  \Cref{def:nagata-noetherian-inverse-system}, equations (91)--(92) of
  \Cref{def:nagata-glued-limit-comparison}, and
  \Cref{def:nagata-final-noetherian-base-system} give the component types
  \[
  \begin{gathered}
    I_S\text{ is a set},\qquad
    {\leq}\ \subseteq I_S\times I_S,\qquad
    S_\lambda\text{ is a scheme}\quad(\lambda\in I_S),\qquad
    L_S\text{ is a scheme},\\
    t_{\mu\lambda}\colon S_\mu\to S_\lambda
      \quad(\lambda,\mu\in I_S,\ \lambda\leq\mu),\qquad
    q_\lambda\colon S\to S_\lambda\quad(\lambda\in I_S),\\
    \pi_\lambda\colon L_S\to S_\lambda\quad(\lambda\in I_S),\qquad
    \kappa_S\colon S\to L_S.
  \end{gathered}                                                    \tag{F0a}
  \]
  By \Cref{lemma:nagata-final-noetherian-base-system},
  \[
    (\forall\lambda\in I_S)\
      \operatorname{Noetherian}(S_\lambda).                       \tag{F0b}
  \]
  The defining implication for a Noetherian scheme and
  \cite[\href{https://stacks.math.columbia.edu/tag/01OY}{Tag 01OY}]
    {stacks-project} are, respectively,
  \[
  \begin{aligned}
    \operatorname{Noetherian}(T)
      &\Longrightarrow\operatorname{QuasiCompact}(T),\\
    \operatorname{Noetherian}(T)
      &\Longrightarrow\operatorname{QuasiSeparated}(T)
        &&\text{by Tag 01OY}.
  \end{aligned}
  \]
  Substitution of $T=S_\lambda$ and (F0b) gives the required stage predicate
  \[
    (\forall\lambda\in I_S)\
      \bigl(
        \operatorname{QuasiCompact}(S_\lambda)
        \ \wedge\
        \operatorname{QuasiSeparated}(S_\lambda)
      \bigr).                                                       \tag{F0c}
  \]
  Equations (88)--(90) give
  \[
  \begin{gathered}
    (\forall\lambda\in I_S)\ \lambda\leq\lambda,\qquad
    (\forall\lambda,\mu,\nu\in I_S)\
      ((\lambda\leq\mu)\wedge(\mu\leq\nu))
        \Longrightarrow\lambda\leq\nu,\\
    \mathsf{RL}_1(\mathsf B_f),\qquad
    \mathsf{RL}_2(\mathsf B_f),\qquad
    \mathsf{RL}_3(\mathsf B_f),\\
    (\forall\lambda,\mu\in I_S)\
      (\lambda\leq\mu)\Longrightarrow
      \operatorname{IsAffine}(t_{\mu\lambda}).
  \end{gathered}                                                    \tag{F0d}
  \]
  Equation (91) of \Cref{def:nagata-glued-limit-comparison} is
  \[
  \begin{gathered}
    L_S=\varprojlim_{\lambda\in I_S}S_\lambda,\\
    (\forall\lambda,\mu\in I_S)\
      \bigl((\lambda\leq\mu)\Longrightarrow
        t_{\mu\lambda}\circ\pi_\mu=\pi_\lambda\bigr).
  \end{gathered}
  \]
  Hence, for every scheme $Z$, the inverse-limit universal property in (91)
  is the explicit bijection
  \[
    \Phi_{\mathsf B_f,Z}\colon
      \operatorname{Hom}_{\mathbf{Sch}}(Z,L_S)
        \xrightarrow{\sim}\mathcal C_{\mathsf B_f}(Z),
    \qquad
    h\longmapsto(\pi_\lambda\circ h)_{\lambda\in I_S}.
  \]
  For every scheme $Z$, the set-theoretic definition of bijectivity is the
  equivalence
  \[
  \begin{aligned}
    &\Bigl[
      \Phi_{\mathsf B_f,Z}\colon
        \operatorname{Hom}_{\mathbf{Sch}}(Z,L_S)
          \xrightarrow{\sim}\mathcal C_{\mathsf B_f}(Z)
    \Bigr]\\
    &\quad\Longleftrightarrow
    \Bigl[
      \Phi_{\mathsf B_f,Z}
        \bigl(\operatorname{Hom}_{\mathbf{Sch}}(Z,L_S)\bigr)
        =\mathcal C_{\mathsf B_f}(Z)
    \Bigr]
    \ \wedge\
    \Bigl[
      (\forall h,h'\in\operatorname{Hom}_{\mathbf{Sch}}(Z,L_S))\
      \bigl(
        \Phi_{\mathsf B_f,Z}(h)=\Phi_{\mathsf B_f,Z}(h')
        \Longrightarrow h=h'
      \bigr)
    \Bigr].
  \end{aligned}
  \]
  Substitution of the projection equation and the right-hand side of this
  equivalence into (RL4) of
  \Cref{def:nagata-relative-limit-datum} is
  \[
  \begin{aligned}
    &\Bigl[
      (\forall\lambda,\mu\in I_S)\
      \bigl((\lambda\leq\mu)\Longrightarrow
        t_{\mu\lambda}\circ\pi_\mu=\pi_\lambda\bigr)
    \Bigr]\\
    &\qquad\wedge
    \Bigl[
      (\forall Z)\
      \Bigl(
        Z\text{ is a scheme}\Longrightarrow
        \Bigl(
        \Bigl[
          \Phi_{\mathsf B_f,Z}
            \bigl(\operatorname{Hom}_{\mathbf{Sch}}(Z,L_S)\bigr)
            =\mathcal C_{\mathsf B_f}(Z)
        \Bigr]
        \ \wedge\
        \Bigl[
          (\forall h,h'\in
            \operatorname{Hom}_{\mathbf{Sch}}(Z,L_S))\
          \bigl(
            \Phi_{\mathsf B_f,Z}(h)=\Phi_{\mathsf B_f,Z}(h')
            \Longrightarrow h=h'
          \bigr)
        \Bigr]
        \Bigr)
      \Bigr)
    \Bigr]
      =\mathsf{RL}_4(\mathsf B_f).
  \end{aligned}                                                     \tag{F0e}
  \]
  The isomorphism assertion in
  \Cref{lemma:nagata-final-noetherian-base-system}, equation (91), and
  equation (92) of \Cref{def:nagata-glued-limit-comparison} are
  \[
    \kappa_S\colon S\xrightarrow{\sim}L_S,\qquad
    (\forall\lambda\in I_S)\
      \pi_\lambda\circ\kappa_S=q_\lambda.
  \]
  Substitution of these two formulas into (RL5) of
  \Cref{def:nagata-relative-limit-datum} is
  \[
    \Bigl[\kappa_S\colon S\xrightarrow{\sim}L_S\Bigr]
      \ \wedge\
    \Bigl[(\forall\lambda\in I_S)\
      \pi_\lambda\circ\kappa_S=q_\lambda\Bigr]
      =\mathsf{RL}_5(\mathsf B_f).                                 \tag{F0f}
  \]
  Substitution of (F0a), (F0c)--(F0f) into the defining conjunction (RL5a)
  gives
  \[
    \mathsf B_f\in\operatorname{RelLim}(S).                        \tag{F1}
  \]
  By \Cref{lemma:nagata-final-fp-envelope}, the predicates required by
  \Cref{def:nagata-separated-descent-stage} are
  \[
    \operatorname{Separated}(q_f^{\mathrm{fp}})
      \wedge
    \operatorname{FinitePresentation}(q_f^{\mathrm{fp}}).          \tag{F1a}
  \]
  Substitute
  \[
    (q,\mathsf B)
      :=(q_f^{\mathrm{fp}},\mathsf B_f)                            \tag{F2}
  \]
  into the choice equation (SD9) of
  \Cref{def:nagata-separated-descent-stage}.  Define its selected
  components by
  \[
  \begin{gathered}
    \lambda_f
      :=\lambda_{\mathrm{sep}}(q_f^{\mathrm{fp}};\mathsf B_f),
    \qquad
    X_{f,\lambda_f}^{\mathrm{fp}}
      :=X_{\mathrm{sep}}(q_f^{\mathrm{fp}};\mathsf B_f),
    \qquad
    q_{f,\lambda_f}
      :=q_{\mathrm{sep}}(q_f^{\mathrm{fp}};\mathsf B_f)
      \colon X_{f,\lambda_f}^{\mathrm{fp}}\to S_{\lambda_f},
    \\
    \bigl(X_{f,\lambda_f}^{\mathrm{fp}}\bigr)_{
      \lambda_f,S}^{q_{f,\lambda_f}}
      =X_{f,\lambda_f}^{\mathrm{fp}}\times_{S_{\lambda_f}}S,
    \qquad
    \epsilon_f
      :=\epsilon_{\mathrm{sep}}(q_f^{\mathrm{fp}};\mathsf B_f)\colon
      X_{f,\lambda_f}^{\mathrm{fp}}\times_{S_{\lambda_f}}S
      \xrightarrow{\sim}X_f^{\mathrm{fp}}.
  \end{gathered}                                                    \tag{F3}
  \]
\end{definition}

\begin{lemma}[Properties of the final descended model]
  \label{lemma:nagata-final-descended-model}
  \uses{def:nagata-final-descended-model}
  \checked\lean{}

  The morphism
  $q_{f,\lambda_f}\colon X_{f,\lambda_f}^{\mathrm{fp}}\to S_{\lambda_f}$
  is separated and of
  finite presentation, $S_{\lambda_f}$ is Noetherian, and $\epsilon_f$ is an
  isomorphism satisfying
  \[
    q_f^{\mathrm{fp}}\circ\epsilon_f
      =p_S^{\lambda_f,q_{f,\lambda_f}}
      =\operatorname{pr}_S.
  \]
\end{lemma}

\begin{proof}
  \uses{lemma:nagata-separated-eventually,
    lemma:nagata-final-noetherian-base-system}
  Substitution of (F2)--(F3) into (SD10) of
  \Cref{lemma:nagata-separated-eventually} is the conjunction
  \[
  \begin{gathered}
    \lambda_f\in I_S,\\
    \operatorname{Separated}(q_{f,\lambda_f})\wedge
      \operatorname{FinitePresentation}(q_{f,\lambda_f}),\\
    \epsilon_f\colon
      X_{f,\lambda_f}^{\mathrm{fp}}\times_{S_{\lambda_f}}S
        \xrightarrow{\sim}X_f^{\mathrm{fp}},\\
    q_f^{\mathrm{fp}}\circ\epsilon_f
      =p_S^{\lambda_f,q_{f,\lambda_f}}
      =\operatorname{pr}_S.
  \end{gathered}
  \]
  The universal statement in
  \Cref{lemma:nagata-final-noetherian-base-system}, evaluated at
  the first conjunct \(\lambda_f\in I_S\), is
  \[
    \operatorname{Noetherian}(S_{\lambda_f}).
  \]
\end{proof}

\begin{definition}[Noetherian compactification used in the final assembly]
  \label{def:nagata-final-noetherian-compactification}
  \uses{def:nagata-final-descended-model,
    def:nagata-noetherian-compactification,
    lemma:nagata-final-descended-model}
  \checked\lean{}

  By \Cref{lemma:nagata-final-descended-model}, the exact input predicates are
  \[
    \operatorname{Noetherian}(S_{\lambda_f}),\qquad
    \operatorname{Separated}(q_{f,\lambda_f}),\qquad
    \operatorname{FinitePresentation}(q_{f,\lambda_f}).           \tag{FNC0}
  \]
  The implication
  \[
    \operatorname{FinitePresentation}(q_{f,\lambda_f})
      \Longrightarrow
    \operatorname{FiniteType}(q_{f,\lambda_f})                    \tag{FNC1}
  \]
  is
  \cite[\href{https://stacks.math.columbia.edu/tag/01TW}{Tag 01TW}]
  {stacks-project}.  Hence the tuple of hypotheses supplied to
  \Cref{def:nagata-noetherian-compactification} is
  \[
    \bigl(
      \operatorname{Noetherian}(S_{\lambda_f}),
      \operatorname{Separated}(q_{f,\lambda_f}),
      \operatorname{FiniteType}(q_{f,\lambda_f})
    \bigr).                                                        \tag{FNC2}
  \]
  Define
  \[
    C_{f,\lambda_f}
      =\mathsf C_{\mathrm{Noeth}}
        (q_{f,\lambda_f})
      =\bigl(
        \overline X_{f,\lambda_f}^{\mathrm{fp}},
        j_{f,\lambda_f},p_{f,\lambda_f},e_{f,\lambda_f}
        \bigr).
  \]
  The component types and equality are
  \[
  \begin{gathered}
    j_{f,\lambda_f}\colon X_{f,\lambda_f}^{\mathrm{fp}}
      \to\overline X_{f,\lambda_f}^{\mathrm{fp}},\qquad
    p_{f,\lambda_f}\colon\overline X_{f,\lambda_f}^{\mathrm{fp}}
      \to S_{\lambda_f},\\
    e_{f,\lambda_f}\colon
      p_{f,\lambda_f}\circ j_{f,\lambda_f}=q_{f,\lambda_f}.
  \end{gathered}
  \]
\end{definition}

\begin{lemma}[Property of the selected Noetherian compactification]
  \label{lemma:nagata-final-noetherian-compactification}
  \uses{def:nagata-final-noetherian-compactification}
  \checked\lean{}

  The constructed $C_{f,\lambda_f}$ is a compactification of
  $X_{f,\lambda_f}^{\mathrm{fp}}\to S_{\lambda_f}$.
\end{lemma}

\begin{proof}
\checked
  \uses{lemma:nagata-final-descended-model,
    thm:nagata-noetherian}
  By \Cref{lemma:nagata-final-descended-model},
  \[
    \operatorname{Noetherian}(S_{\lambda_f}),
    \qquad
    \operatorname{Separated}(q_{f,\lambda_f}),
    \qquad
    \operatorname{FinitePresentation}(q_{f,\lambda_f}).
  \]
  By
  \cite[\href{https://stacks.math.columbia.edu/tag/01TW}{Tag 01TW}]
  {stacks-project},
  \[
    \operatorname{FinitePresentation}(q_{f,\lambda_f})
      \Longrightarrow
    \operatorname{FiniteType}(q_{f,\lambda_f}).
  \]
  Apply
  \Cref{thm:nagata-noetherian} to the specified morphism
  $q_{f,\lambda_f}$ and the construction in
  \Cref{def:nagata-final-noetherian-compactification}.  The result is
  \[
  \begin{gathered}
    \operatorname{IsQuasiCompactOpenImmersion}(j_{f,\lambda_f}),\\
    \operatorname{Proper}(p_{f,\lambda_f}),\\
    p_{f,\lambda_f}\circ j_{f,\lambda_f}=q_{f,\lambda_f}.
  \end{gathered}
  \]
  By \Cref{def:nagata-compactification}, these equalities and properties say
  that $C_{f,\lambda_f}$ is a compactification of $q_{f,\lambda_f}$.
\end{proof}

\begin{definition}[Base-changed compactification of the finite-presentation envelope]
\checked
  \label{def:nagata-final-basechanged-compactification}
  \uses{def:nagata-compactification-base-change,
    def:nagata-final-noetherian-compactification,
    def:nagata-final-descended-model,
    lemma:nagata-final-descended-model}
  \lean{}

  Let $q_{\lambda_f}\colon S\to S_{\lambda_f}$ be the morphism in
  \Cref{def:nagata-final-noetherian-base-system}.  Define
  \[
  \begin{aligned}
    X_{f,S}^{\mathrm{fp}}
      &=X_{f,\lambda_f}^{\mathrm{fp}}
        \times_{S_{\lambda_f}}S,\\
    \overline X_f^{\mathrm{fp}}
      &=\overline X_{f,\lambda_f}^{\mathrm{fp}}
        \times_{S_{\lambda_f}}S,\\
    j_{f,S}
      &=j_{f,\lambda_f}\times_{S_{\lambda_f}}\operatorname{id}_S
        \colon X_{f,S}^{\mathrm{fp}}\to\overline X_f^{\mathrm{fp}},\\
    p_f^{\mathrm{fp}}
      &=\operatorname{pr}_S\colon\overline X_f^{\mathrm{fp}}\to S,\\
    j_f^{\mathrm{fp}}
      &=j_{f,S}\circ\epsilon_f^{-1}
        \colon X_f^{\mathrm{fp}}\to\overline X_f^{\mathrm{fp}}.
  \end{aligned}
  \]
  Define the base-change equality
  \[
    e_{f,S}\colon
      p_f^{\mathrm{fp}}\circ j_{f,S}=\operatorname{pr}_S
  \]
  to be the pullback of $e_{f,\lambda_f}$ along
  $q_{\lambda_f}\colon S\to S_{\lambda_f}$.  By
  \Cref{lemma:nagata-final-descended-model},
  \[
    q_f^{\mathrm{fp}}\circ\epsilon_f=\operatorname{pr}_S,
    \qquad
    \epsilon_f\circ\epsilon_f^{-1}
      =\operatorname{id}_{X_f^{\mathrm{fp}}}.
  \]
  Define the equality witness $e_f^{\mathrm{fp}}$ by the chain
  \[
  \begin{aligned}
    e_f^{\mathrm{fp}}\colon\quad
    p_f^{\mathrm{fp}}\circ j_f^{\mathrm{fp}}
      &=p_f^{\mathrm{fp}}\circ j_{f,S}\circ\epsilon_f^{-1}\\
      &=\operatorname{pr}_S\circ\epsilon_f^{-1}\\
      &=(q_f^{\mathrm{fp}}\circ\epsilon_f)\circ\epsilon_f^{-1}\\
      &=q_f^{\mathrm{fp}}\circ
        (\epsilon_f\circ\epsilon_f^{-1})\\
      &=q_f^{\mathrm{fp}}\circ
        \operatorname{id}_{X_f^{\mathrm{fp}}}\\
      &=q_f^{\mathrm{fp}},
  \end{aligned}
  \]
  where the second equality is $e_{f,S}$ and the third equality substitutes
  $q_f^{\mathrm{fp}}\circ\epsilon_f=\operatorname{pr}_S$.  Define
  \[
    C_f^{\mathrm{fp}}
      :=\bigl(\overline X_f^{\mathrm{fp}},j_f^{\mathrm{fp}},
        p_f^{\mathrm{fp}},e_f^{\mathrm{fp}}\bigr).
  \]
\end{definition}

\begin{lemma}[Property of the base-changed compactification]
  \label{lemma:nagata-final-basechanged-compactification}
  \uses{def:nagata-final-basechanged-compactification}
  \lean{}
\checked
  The constructed $C_f^{\mathrm{fp}}$ is a compactification of
  $X_f^{\mathrm{fp}}\to S$.
\end{lemma}

\begin{proof}
\checked
  \uses{lemma:nagata-final-noetherian-compactification,
    lemma:nagata-compactification-base-change,
    lemma:nagata-final-descended-model,
    lemma:nagata-morphism-property-calculus}
  By \Cref{lemma:nagata-final-noetherian-compactification},
  \[
    \operatorname{IsQuasiCompactOpenImmersion}(j_{f,\lambda_f}),
    \qquad
    \operatorname{IsProper}(p_{f,\lambda_f}),
    \qquad
    p_{f,\lambda_f}\circ j_{f,\lambda_f}=q_{f,\lambda_f}.
  \]
  The conclusion of \Cref{lemma:nagata-compactification-base-change} for
  $q_{\lambda_f}\colon S\to S_{\lambda_f}$ is
  \[
    \operatorname{IsQuasiCompactOpenImmersion}(j_{f,S}),
    \qquad
    \operatorname{IsProper}(p_f^{\mathrm{fp}}),
    \qquad
    p_f^{\mathrm{fp}}\circ j_{f,S}=\operatorname{pr}_S.
  \]
  By \Cref{lemma:nagata-final-descended-model},
  \[
    \operatorname{IsIso}(\epsilon_f),
    \qquad
    \operatorname{IsIso}(\epsilon_f^{-1}).
  \]
  Its inverse satisfies
  $\epsilon_f\circ\epsilon_f^{-1}
    =\operatorname{id}_{X_f^{\mathrm{fp}}}$, and the equality
  $q_f^{\mathrm{fp}}\circ\epsilon_f=\operatorname{pr}_S$ produces the chain
  \[
  \begin{aligned}
    \operatorname{pr}_S\circ\epsilon_f^{-1}
      &=(q_f^{\mathrm{fp}}\circ\epsilon_f)
        \circ\epsilon_f^{-1}\\
      &=q_f^{\mathrm{fp}}\circ
        (\epsilon_f\circ\epsilon_f^{-1})\\
      &=q_f^{\mathrm{fp}}\circ
        \operatorname{id}_{X_f^{\mathrm{fp}}}\\
      &=q_f^{\mathrm{fp}}.
  \end{aligned}
  \]
  Expanding (P0)--(P0a) in
  \Cref{def:nagata-morphism-predicate-notation} and applying (PC3) in
  \Cref{lemma:nagata-morphism-property-calculus} is the predicate chain
  \[
  \begin{aligned}
    \operatorname{IsIso}(\epsilon_f^{-1})
      &\Longrightarrow
        \operatorname{IsQuasiCompactOpenImmersion}(\epsilon_f^{-1}),\\
    \operatorname{IsQuasiCompactOpenImmersion}(\epsilon_f^{-1})
      \wedge\operatorname{IsQuasiCompactOpenImmersion}(j_{f,S})
      &\Longrightarrow
        \operatorname{IsQuasiCompactOpenImmersion}
          (j_{f,S}\circ\epsilon_f^{-1})\\
      &\Longleftrightarrow
        \operatorname{IsQuasiCompactOpenImmersion}(j_f^{\mathrm{fp}}).
  \end{aligned}
  \]
  The factorization equality is the chain
  \[
  \begin{aligned}
    p_f^{\mathrm{fp}}\circ j_f^{\mathrm{fp}}
      &=p_f^{\mathrm{fp}}\circ j_{f,S}\circ\epsilon_f^{-1}\\
      &=\operatorname{pr}_S\circ\epsilon_f^{-1}\\
      &=q_f^{\mathrm{fp}}.
  \end{aligned}
  \]
  Substitution into \Cref{def:nagata-compactification} is the membership
  equality
  \[
    C_f^{\mathrm{fp}}
      =(\overline X_f^{\mathrm{fp}},j_f^{\mathrm{fp}},
        p_f^{\mathrm{fp}},e_f^{\mathrm{fp}})
      \in\operatorname{Compactification}(q_f^{\mathrm{fp}}).
  \]
\end{proof}

\begin{definition}[Final Nagata compactification candidate]
\checked
  \label{def:nagata-final-compactification}
  \uses{def:nagata-closed-subscheme-compactification,
    def:nagata-scheme-theoretic-image,
    def:nagata-scheme-image-factor,
    def:nagata-final-fp-envelope,
    def:nagata-final-basechanged-compactification,
    lemma:nagata-final-fp-envelope,
    lemma:nagata-final-basechanged-compactification,
    lemma:nagata-morphism-property-calculus}
  \lean{}

  With the symbols of
  \Cref{def:nagata-final-fp-envelope,def:nagata-final-basechanged-compactification},
  define
  \[
    g_f=j_f^{\mathrm{fp}}\circ i_f^{\mathrm{fp}}
      \colon X\longrightarrow\overline X_f^{\mathrm{fp}},
  \]
  By \Cref{lemma:nagata-final-fp-envelope,lemma:nagata-final-basechanged-compactification,lemma:nagata-morphism-property-calculus},
  the exact predicate chain is
  \[
  \begin{aligned}
    \operatorname{IsClosedImmersion}(i_f^{\mathrm{fp}})
      &\Longrightarrow
        \bigl(
          \operatorname{IsQuasiCompact}(i_f^{\mathrm{fp}})\wedge
          \operatorname{QuasiSeparated}(i_f^{\mathrm{fp}})
        \bigr)
       \Longrightarrow\operatorname{IsQCQS}(i_f^{\mathrm{fp}}),\\
    \operatorname{IsQuasiCompactOpenImmersion}(j_f^{\mathrm{fp}})
      &\Longrightarrow
        \bigl(
          \operatorname{IsQuasiCompact}(j_f^{\mathrm{fp}})\wedge
          \operatorname{QuasiSeparated}(j_f^{\mathrm{fp}})
        \bigr)
       \Longrightarrow\operatorname{IsQCQS}(j_f^{\mathrm{fp}}),\\
    \operatorname{IsQCQS}(i_f^{\mathrm{fp}})\wedge
      \operatorname{IsQCQS}(j_f^{\mathrm{fp}})
      &\Longrightarrow
        \operatorname{IsQCQS}
          (j_f^{\mathrm{fp}}\circ i_f^{\mathrm{fp}})
       =\operatorname{IsQCQS}(g_f).
  \end{aligned}
  \]
  Only after this last predicate has been established, define
  \[
    \overline X_f
      :=\operatorname{SIm}_{\overline X_f^{\mathrm{fp}}}(g_f).
  \]
  By \Cref{def:nagata-scheme-theoretic-image}, define the closed immersion
  \[
    \iota_f:=i_{g_f}\colon\overline X_f\hookrightarrow
      \overline X_f^{\mathrm{fp}}
  \]
  by the ideal
  \[
    \varepsilon_{g_f}:=g_f^\flat\colon
      \mathcal O_{\overline X_f^{\mathrm{fp}}}
        \longrightarrow(g_f)_*\mathcal O_X,
    \qquad
    \mathcal K_{g_f}:=\ker(\varepsilon_{g_f}).
  \]
  By
  \Cref{def:nagata-scheme-image-factor}, define
  \[
    j_f=g_f^{\mathrm{im}}\colon X\longrightarrow\overline X_f,
    \qquad
    \iota_f\circ j_f=g_f.
  \]
  Define
  \[
    p_f=p_f^{\mathrm{fp}}\circ\iota_f
      \colon\overline X_f\longrightarrow S.
  \]
  The equalities in
  \Cref{def:nagata-scheme-image-factor,lemma:nagata-final-basechanged-compactification,lemma:nagata-final-fp-envelope}
  define the equality witness
  \[
  \begin{aligned}
    e_f\colon\quad p_f\circ j_f
      &=(p_f^{\mathrm{fp}}\circ\iota_f)\circ j_f\\
      &=p_f^{\mathrm{fp}}\circ(\iota_f\circ j_f)\\
      &=p_f^{\mathrm{fp}}\circ g_f\\
      &=p_f^{\mathrm{fp}}\circ
        (j_f^{\mathrm{fp}}\circ i_f^{\mathrm{fp}})\\
      &=(p_f^{\mathrm{fp}}\circ j_f^{\mathrm{fp}})
        \circ i_f^{\mathrm{fp}}\\
      &=q_f^{\mathrm{fp}}\circ i_f^{\mathrm{fp}}\\
      &=f.
  \end{aligned}
  \]
  Finally, define
  \[
    \mathsf C_{\mathrm{Nag}}(f)
      :=\mathsf{Closed}(i_f^{\mathrm{fp}},C_f^{\mathrm{fp}})
      =\bigl(\overline X_f,j_f,p_f,e_f\bigr).
  \]
\end{definition}

\begin{lemma}[The final ambient morphism is proper]
\checked
  \label{lemma:nagata-final-ambient-proper}
  \uses{def:nagata-final-compactification}
  \lean{}

  The constructed morphism $p_f\colon\overline X_f\to S$ satisfies
  \[
    \operatorname{IsProper}(p_f).
  \]
\end{lemma}

\begin{proof}
\checked
  \uses{prop:nagata-closed-subscheme-compactifiable}
  By \Cref{def:nagata-final-compactification},
  \[
    \mathsf C_{\mathrm{Nag}}(f)
      =\mathsf{Closed}(i_f^{\mathrm{fp}},C_f^{\mathrm{fp}})
      =(\overline X_f,j_f,p_f,e_f).
  \]
  The conclusion of
  \Cref{prop:nagata-closed-subscheme-compactifiable}, instantiated at
  $(i,C')=(i_f^{\mathrm{fp}},C_f^{\mathrm{fp}})$, is
  \[
    \mathsf{Closed}(i_f^{\mathrm{fp}},C_f^{\mathrm{fp}})
      \in\operatorname{Compactification}(f).
  \]
  Substitution of the tuple equality into
  \Cref{def:nagata-compactification} and projection to its properness
  component is the implication
  \[
    (\overline X_f,j_f,p_f,e_f)
      \in\operatorname{Compactification}(f)
      \Longrightarrow\operatorname{IsProper}(p_f).
  \]
\end{proof}

\begin{lemma}[The final source map is a quasi-compact open immersion]
  \label{lemma:nagata-final-open-immersion}
  \uses{def:nagata-final-compactification}
  \checked\lean{}

  The constructed morphism $j_f\colon X\to\overline X_f$ satisfies
  \[
    \operatorname{IsQuasiCompactOpenImmersion}(j_f).
  \]
\end{lemma}

\begin{proof}
\checked
  \uses{def:nagata-closed-subscheme-class,
    def:nagata-direct-image-structure-map,
    def:nagata-final-compactification,
    def:nagata-scheme-image-comparison,
    lemma:nagata-final-fp-envelope,
    lemma:nagata-final-basechanged-compactification,
    lemma:nagata-scheme-image-restrict-open,
    lemma:nagata-morphism-property-calculus}
  By \Cref{def:nagata-final-compactification} and the notation of
  \Cref{lemma:nagata-scheme-image-restrict-open}, define the fibre product and
  both of its projections by
  \[
  \begin{gathered}
    P_f:=P_{g_f,j_f^{\mathrm{fp}}}
      :=\overline X_f\times_{\overline X_f^{\mathrm{fp}}}
        X_f^{\mathrm{fp}},\\
    \pi_{\overline X,f}:=\operatorname{pr}_{\overline X_f}
      \colon P_f\to\overline X_f,
    \qquad
    \pi_{X,f}:=\operatorname{pr}_{X_f^{\mathrm{fp}}}
      \colon P_f\to X_f^{\mathrm{fp}},\\
    \iota_f\circ\pi_{\overline X,f}
      =j_f^{\mathrm{fp}}\circ\pi_{X,f}.
  \end{gathered}                                                   \tag{FO0}
  \]
  The base-change implication in
  \Cref{lemma:nagata-morphism-property-calculus} is
  \[
    \operatorname{IsClosedImmersion}(\iota_f)
      \Longrightarrow
    \operatorname{IsClosedImmersion}(\pi_{X,f}).                 \tag{FO1}
  \]

  Define the base change of $g_f$ along $j_f^{\mathrm{fp}}$ by
  \[
  \begin{gathered}
    X_{f,U}:=X\times_{\overline X_f^{\mathrm{fp}}}
      X_f^{\mathrm{fp}},\\
    r_X:=\operatorname{pr}_X\colon X_{f,U}\to X,
    \qquad
    r_U:=\operatorname{pr}_{X_f^{\mathrm{fp}}}
      \colon X_{f,U}\to X_f^{\mathrm{fp}},\\
    g_{f,U}:=r_U,
    \qquad
    g_f\circ r_X=j_f^{\mathrm{fp}}\circ g_{f,U}.
  \end{gathered}                                                    \tag{FO2}
  \]
  By
  \Cref{lemma:nagata-final-fp-envelope,lemma:nagata-final-basechanged-compactification},
  the morphism-property implications and base-change stability are
  \[
  \begin{gathered}
    \operatorname{IsClosedImmersion}(i_f^{\mathrm{fp}})
      \Longrightarrow\operatorname{IsQCQS}(i_f^{\mathrm{fp}}),\\
    \operatorname{IsQuasiCompactOpenImmersion}(j_f^{\mathrm{fp}})
      \Longrightarrow\operatorname{IsQCQS}(j_f^{\mathrm{fp}}),\\
    \operatorname{IsQCQS}(i_f^{\mathrm{fp}})
      \ \wedge\ \operatorname{IsQCQS}(j_f^{\mathrm{fp}})
      \ \wedge\ (g_f=j_f^{\mathrm{fp}}\circ i_f^{\mathrm{fp}})
      \Longrightarrow\operatorname{IsQCQS}(g_f)
      \Longrightarrow\operatorname{IsQCQS}(g_{f,U}).
  \end{gathered}                                                    \tag{FO2a}
  \]
  Only after (FO2a), define
  \[
  \begin{gathered}
    Z_{f,U}:=\operatorname{SIm}_{X_f^{\mathrm{fp}}}(g_{f,U}),\\
    c_{f,U}:=i_{g_{f,U}}\colon Z_{f,U}\hookrightarrow X_f^{\mathrm{fp}}.
  \end{gathered}                                                    \tag{FO2b}
  \]
  With the substitutions
  \[
    (g,w,W,g_W,P_{g,w})
      =(g_f,j_f^{\mathrm{fp}},X_f^{\mathrm{fp}},g_{f,U},P_f),
  \]
  \Cref{lemma:nagata-scheme-image-restrict-open} defines
  \[
    \delta_f:=\delta_{g_f,j_f^{\mathrm{fp}}}
      \colon Z_{f,U}\xrightarrow{\sim}P_f,
    \qquad
    \pi_{X,f}\circ\delta_f=c_{f,U}.                              \tag{FO3}
  \]

  The equality
  $g_f=j_f^{\mathrm{fp}}\circ i_f^{\mathrm{fp}}$ in
  \Cref{def:nagata-final-compactification} defines
  \[
  \begin{gathered}
    a_f:=(\operatorname{id}_X,i_f^{\mathrm{fp}})
      \colon X\to X_{f,U},
    \qquad
    b_f:=r_X\colon X_{f,U}\to X,\\
    r_X\circ a_f=\operatorname{id}_X,
    \qquad
    r_U\circ a_f=i_f^{\mathrm{fp}},
    \qquad
    g_{f,U}\circ a_f=i_f^{\mathrm{fp}}.
  \end{gathered}                                                    \tag{FO4}
  \]
  Since $j_f^{\mathrm{fp}}$ is a monomorphism, the defining fibre-product
  equality in (FO2) and the factorization of $g_f$ give
  \[
  \begin{aligned}
    j_f^{\mathrm{fp}}\circ r_U
      &=g_f\circ r_X
       =j_f^{\mathrm{fp}}\circ i_f^{\mathrm{fp}}\circ r_X,\\
    r_U&=i_f^{\mathrm{fp}}\circ r_X.
  \end{aligned}                                                     \tag{FO4a}
  \]
  Equations (FO4) and (FO4a) give both projection computations
  \[
  \begin{aligned}
    r_X\circ(a_f\circ b_f)
      &=r_X
       =r_X\circ\operatorname{id}_{X_{f,U}},\\
    r_U\circ(a_f\circ b_f)
      &=i_f^{\mathrm{fp}}\circ r_X
       =r_U
       =r_U\circ\operatorname{id}_{X_{f,U}}.
  \end{aligned}
  \]
  Fibre-product uniqueness and (FO4) are therefore the inverse equations
  \[
    a_f\circ b_f=\operatorname{id}_{X_{f,U}},
    \qquad
    b_f\circ a_f=\operatorname{id}_X,
    \qquad
    \operatorname{IsIso}(a_f),\quad a_f^{-1}=b_f.                 \tag{FO4b}
  \]

  By \Cref{def:nagata-direct-image-structure-map}, define every structure map
  used below by its adjoint structure morphism and define both ideals as its
  kernel:
  \[
  \begin{gathered}
    \varepsilon_{a_f}:=a_f^\flat\colon
      \mathcal O_{X_{f,U}}\to(a_f)_*\mathcal O_X,\\
    \varepsilon_{g_{f,U}}:=g_{f,U}^\flat\colon
      \mathcal O_{X_f^{\mathrm{fp}}}
        \to(g_{f,U})_*\mathcal O_{X_{f,U}},\\
    \varepsilon_{i_f^{\mathrm{fp}}}:=(i_f^{\mathrm{fp}})^\flat\colon
      \mathcal O_{X_f^{\mathrm{fp}}}
        \to(i_f^{\mathrm{fp}})_*\mathcal O_X,\\
    \mathcal K_{f,U}:=\ker(\varepsilon_{g_{f,U}}),
    \qquad
    \mathcal K_{f,i}:=\ker(\varepsilon_{i_f^{\mathrm{fp}}}).
  \end{gathered}                                                    \tag{FO5}
  \]
  Define the comparison map of $\mathcal O_{X_f^{\mathrm{fp}}}$-algebras by
  \[
  \begin{aligned}
    \vartheta_f
      &:=(g_{f,U})_*\varepsilon_{a_f}\colon
        (g_{f,U})_*\mathcal O_{X_{f,U}}
          \xrightarrow{\sim}
        (g_{f,U})_*(a_f)_*\mathcal O_X\\
      &=(g_{f,U}\circ a_f)_*\mathcal O_X
       =(i_f^{\mathrm{fp}})_*\mathcal O_X.
  \end{aligned}                                                     \tag{FO5a}
  \]
  The composition identity in
  \Cref{def:nagata-direct-image-structure-map} gives the equality chain
  \[
  \begin{aligned}
    \vartheta_f\circ\varepsilon_{g_{f,U}}
      &=\bigl((g_{f,U})_*\varepsilon_{a_f}\bigr)
        \circ\varepsilon_{g_{f,U}}\\
      &=\varepsilon_{g_{f,U}\circ a_f}
       =\varepsilon_{i_f^{\mathrm{fp}}}.
  \end{aligned}                                                     \tag{FO5b}
  \]
  Since (FO4b) makes $\varepsilon_{a_f}$ and hence $\vartheta_f$
  isomorphisms, kernel invariance under postcomposition by $\vartheta_f$ is
  the formal chain
  \[
  \begin{aligned}
    \mathcal K_{f,U}
      &=\ker(\varepsilon_{g_{f,U}})
       =\ker(\vartheta_f\circ\varepsilon_{g_{f,U}})\\
      &=\ker(\varepsilon_{i_f^{\mathrm{fp}}})
       =\mathcal K_{f,i}.
  \end{aligned}                                                     \tag{FO6}
  \]

  By \Cref{lemma:nagata-final-fp-envelope}, the predicate chain
  \[
    \operatorname{IsClosedImmersion}(i_f^{\mathrm{fp}})
      \Longrightarrow\operatorname{IsQCQS}(i_f^{\mathrm{fp}})    \tag{FO6a}
  \]
  holds.  Only after (FO6a), put
  \[
    Z_{f,i}:=\operatorname{SIm}_{X_f^{\mathrm{fp}}}
      (i_f^{\mathrm{fp}}),
    \qquad
    c_{f,i}:=i_{i_f^{\mathrm{fp}}}\colon
      Z_{f,i}\hookrightarrow X_f^{\mathrm{fp}}.                  \tag{FO7}
  \]
  Equality (FO6) defines the quotient-algebra isomorphism and its inverse by
  \[
  \begin{gathered}
    \Lambda_f\colon
      \mathcal O_{X_f^{\mathrm{fp}}}/\mathcal K_{f,i}
        \xrightarrow{\sim}
      \mathcal O_{X_f^{\mathrm{fp}}}/\mathcal K_{f,U},
    \qquad
    \Lambda_f([s]_{\mathcal K_{f,i}}):=[s]_{\mathcal K_{f,U}},\\
    \Lambda_f^{-1}([s]_{\mathcal K_{f,U}}):=[s]_{\mathcal K_{f,i}},
    \qquad
    \Lambda_f^{-1}\circ\Lambda_f=\operatorname{id},
    \qquad
    \Lambda_f\circ\Lambda_f^{-1}=\operatorname{id}.
  \end{gathered}                                                    \tag{FO7a}
  \]
  By the relative-spectrum construction in
  \Cref{def:nagata-scheme-theoretic-image}, define
  \[
  \begin{gathered}
    \lambda_f:=\operatorname{Spec}_{X_f^{\mathrm{fp}}}(\Lambda_f)
      \colon Z_{f,U}\xrightarrow{\sim}Z_{f,i},\\
    c_{f,i}\circ\lambda_f=c_{f,U},
    \qquad
    c_{f,U}\circ\lambda_f^{-1}
      =c_{f,i}\circ\lambda_f\circ\lambda_f^{-1}=c_{f,i}.
  \end{gathered}                                                    \tag{FO7b}
  \]
  Apply \Cref{def:nagata-scheme-image-comparison} with
  \[
    (g,h,z,\mathcal K_z,r_z)=
      (i_f^{\mathrm{fp}},\operatorname{id}_X,i_f^{\mathrm{fp}},
        \mathcal K_{f,i},c_{f,i})
  \]
  and define
  \[
    \sigma_f:=\sigma_{i_f^{\mathrm{fp}}}\colon X\xrightarrow{\sim}Z_{f,i},
    \qquad
    c_{f,i}\circ\sigma_f=i_f^{\mathrm{fp}}.                      \tag{FO8}
  \]

  Equations (FO3), (FO7b), and (FO8), substituted into (CSub3) of
  \Cref{def:nagata-closed-subscheme-class}, give the typed chain
  \[
  \begin{aligned}
    [\,X\xhookrightarrow{i_f^{\mathrm{fp}}}X_f^{\mathrm{fp}}\,]
      &\overset{\sigma_f}{=}
        [\,Z_{f,i}\xhookrightarrow{c_{f,i}}X_f^{\mathrm{fp}}\,]\\
      &\overset{\lambda_f^{-1}}{=}
        [\,Z_{f,U}\xhookrightarrow{c_{f,U}}X_f^{\mathrm{fp}}\,]\\
      &\overset{\delta_f}{=}
        [\,P_f\xhookrightarrow{\pi_{X,f}}X_f^{\mathrm{fp}}\,]
      \quad\text{in }\operatorname{ClosedSub}(X_f^{\mathrm{fp}}).
  \end{aligned}                                                     \tag{FO9}
  \]
  Define the comparison set
  \[
    \mathscr P_f:=
      \left\{\chi\in\operatorname{Hom}(X,P_f)\ \middle|\
        \operatorname{IsIso}(\chi)\ \wedge\
        \pi_{X,f}\circ\chi=i_f^{\mathrm{fp}}\right\}.            \tag{FO10}
  \]
  The first and last members of (FO9) and (CSub3) give the equivalences
  \[
  \begin{aligned}
    &[\,X\xhookrightarrow{i_f^{\mathrm{fp}}}X_f^{\mathrm{fp}}\,]
      =[\,P_f\xhookrightarrow{\pi_{X,f}}X_f^{\mathrm{fp}}\,]\\
    &\quad\Longleftrightarrow
      \exists\chi\in\operatorname{Hom}(X,P_f):
        \operatorname{IsIso}(\chi)\ \wedge\
        \pi_{X,f}\circ\chi=i_f^{\mathrm{fp}}\\
    &\quad\Longleftrightarrow\mathscr P_f\neq\varnothing.
  \end{aligned}                                                     \tag{FO10a}
  \]
  Define the explicit comparison morphism by the three witnesses in (FO9):
  \[
  \begin{aligned}
    \chi_f
      &:=\delta_f\circ\lambda_f^{-1}\circ\sigma_f
        \colon X\longrightarrow P_f,\\
    \operatorname{IsIso}(\delta_f)
      \ \wedge\ \operatorname{IsIso}(\lambda_f^{-1})
      \ \wedge\ \operatorname{IsIso}(\sigma_f)
      &\Longrightarrow\operatorname{IsIso}(\chi_f),\\
    \pi_{X,f}\circ\chi_f
      &=\pi_{X,f}\circ\delta_f\circ\lambda_f^{-1}\circ\sigma_f\\
      &=c_{f,U}\circ\lambda_f^{-1}\circ\sigma_f
       =c_{f,i}\circ\sigma_f
       =i_f^{\mathrm{fp}},\\
    \operatorname{IsIso}(\chi_f)
      \ \wedge\ (\pi_{X,f}\circ\chi_f=i_f^{\mathrm{fp}})
      &\Longleftrightarrow\chi_f\in\mathscr P_f
       \Longrightarrow\mathscr P_f\neq\varnothing.
  \end{aligned}                                                     \tag{FO11}
  \]
  Equations (FO0), (FO11), and the two factorization equalities in
  \Cref{def:nagata-final-compactification} give
  \[
  \begin{aligned}
    \iota_f\circ\pi_{\overline X,f}\circ\chi_f
      &=j_f^{\mathrm{fp}}\circ\pi_{X,f}\circ\chi_f\\
      &=j_f^{\mathrm{fp}}\circ i_f^{\mathrm{fp}}
       =g_f
       =\iota_f\circ j_f.
  \end{aligned}
  \]
  Cancellation of the closed immersion $\iota_f$ is the second projection
  equation
  \[
    \pi_{\overline X,f}\circ\chi_f=j_f.                          \tag{FO12}
  \]

  The equality
  $\iota_f\circ j_f=g_f=j_f^{\mathrm{fp}}\circ i_f^{\mathrm{fp}}$
  defines the canonical fibre-product comparison map
  \[
  \begin{gathered}
    \kappa_f:=(j_f,i_f^{\mathrm{fp}})\colon X\to P_f,\\
    \pi_{\overline X,f}\circ\kappa_f=j_f,
    \qquad
    \pi_{X,f}\circ\kappa_f=i_f^{\mathrm{fp}}.
  \end{gathered}                                                    \tag{FO13}
  \]
  The two projection comparisons are explicitly
  \[
  \begin{aligned}
    \pi_{\overline X,f}\circ\kappa_f
      &=j_f=\pi_{\overline X,f}\circ\chi_f,\\
    \pi_{X,f}\circ\kappa_f
      &=i_f^{\mathrm{fp}}=\pi_{X,f}\circ\chi_f.
  \end{aligned}
  \]
  Fibre-product uniqueness and (FO11) therefore give
  \[
    \kappa_f=\chi_f,
    \qquad
    \operatorname{IsIso}(\kappa_f).                              \tag{FO13a}
  \]
  Let (FO14) denote the commutative square
  \[
  \begin{tikzcd}
    X \arrow[r,"i_f^{\mathrm{fp}}"] \arrow[d,"j_f"'] &
      X_f^{\mathrm{fp}} \arrow[d,"j_f^{\mathrm{fp}}"] \\
    \overline X_f \arrow[r,"\iota_f"'] & \overline X_f^{\mathrm{fp}}.
  \end{tikzcd}                                                       \tag{FO14}
  \]
  \[
  \begin{gathered}
    \pi_{\overline X,f}\circ\kappa_f=j_f,
    \qquad
    \pi_{X,f}\circ\kappa_f=i_f^{\mathrm{fp}},\\
    \bigl[\textup{(FO14) is cartesian}\bigr]
      \Longleftrightarrow\operatorname{IsIso}(\kappa_f).
  \end{gathered}                                                    \tag{FO14a}
  \]
  Equation (FO13a) therefore proves the left-hand member of (FO14a).

  By \Cref{lemma:nagata-final-basechanged-compactification} and base-change
  stability in \Cref{lemma:nagata-morphism-property-calculus},
  \[
  \begin{aligned}
    \operatorname{IsQuasiCompactOpenImmersion}(j_f^{\mathrm{fp}})
      &\Longrightarrow
        \operatorname{IsQuasiCompactOpenImmersion}
          ((j_f^{\mathrm{fp}})_{\iota_f}),\\
    (j_f^{\mathrm{fp}})_{\iota_f}
      &:=\pi_{\overline X,f}\colon P_f\to\overline X_f.
  \end{aligned}                                                     \tag{FO15}
  \]
  Finally, (FO11), (FO12), and composition stability give
  \[
  \begin{aligned}
    \operatorname{IsIso}(\chi_f)
      &\Longrightarrow
        \operatorname{IsQuasiCompactOpenImmersion}(\chi_f),\\
    \operatorname{IsQuasiCompactOpenImmersion}(\chi_f)
      \ \wedge\
    \operatorname{IsQuasiCompactOpenImmersion}(\pi_{\overline X,f})
      &\Longrightarrow
        \operatorname{IsQuasiCompactOpenImmersion}
          (\pi_{\overline X,f}\circ\chi_f)\\
      &=\operatorname{IsQuasiCompactOpenImmersion}(j_f).
  \end{aligned}                                                     \tag{FO16}
  \]
\end{proof}

\begin{lemma}[Factorization equality for the final candidate]
  \label{lemma:nagata-final-factorization}
  \uses{def:nagata-final-compactification}
  \lean{}
    \checked
  The maps in \Cref{def:nagata-final-compactification} satisfy
  \[
    p_f\circ j_f=f.
  \]
\end{lemma}

\begin{proof}
  \uses{lemma:nagata-scheme-image-universal,
    lemma:nagata-final-basechanged-compactification,
    lemma:nagata-final-fp-envelope}
  By \Cref{def:nagata-final-compactification},
  \[
    p_f=p_f^{\mathrm{fp}}\circ\iota_f,
    \qquad
    g_f=j_f^{\mathrm{fp}}\circ i_f^{\mathrm{fp}}.
  \]
  By the factorization equality in
  \Cref{lemma:nagata-scheme-image-universal}, applied to $g_f$, one has
  \[
    \iota_f\circ j_f=g_f.
  \]
  By \Cref{lemma:nagata-final-basechanged-compactification} and
  \Cref{lemma:nagata-final-fp-envelope}, respectively,
  \[
    p_f^{\mathrm{fp}}\circ j_f^{\mathrm{fp}}=q_f^{\mathrm{fp}},
    \qquad
    q_f^{\mathrm{fp}}\circ i_f^{\mathrm{fp}}=f.
  \]
  Successive substitution of these four equalities, with associativity of
  composition at the second and fifth lines, is the chain
  \[
  \begin{aligned}
    p_f\circ j_f
      &=(p_f^{\mathrm{fp}}\circ\iota_f)\circ j_f\\
      &=p_f^{\mathrm{fp}}\circ(\iota_f\circ j_f)\\
      &=p_f^{\mathrm{fp}}\circ g_f\\
      &=p_f^{\mathrm{fp}}\circ
        (j_f^{\mathrm{fp}}\circ i_f^{\mathrm{fp}})\\
      &=(p_f^{\mathrm{fp}}\circ j_f^{\mathrm{fp}})
        \circ i_f^{\mathrm{fp}}\\
      &=q_f^{\mathrm{fp}}\circ i_f^{\mathrm{fp}}\\
      &=f.
  \end{aligned}
  \]
\end{proof}

\begin{theorem}[Nagata compactification,
{\cite[\href{https://stacks.math.columbia.edu/tag/0F3T}{Tag 0F3T}]{stacks-project}}]
  \label{thm:nagata-compactification}
  \uses{def:nagata-compactification,def:nagata-final-compactification}
  \lean{}
\checked
  For the morphism $f\colon X\to S$ fixed in this section,
  \[
    \mathsf C_{\mathrm{Nag}}(f)
      \in\operatorname{Compactification}(f),
    \qquad
    \operatorname{Compactifiable}(f).
  \]
\end{theorem}

\begin{proof}
    \checked
  \uses{lemma:nagata-final-ambient-proper,
    lemma:nagata-final-open-immersion,
    lemma:nagata-final-factorization,
    def:nagata-compactifiable}
  Definition \ref{def:nagata-final-compactification} is the typed
  quadruple
  \[
    \mathsf C_{\mathrm{Nag}}(f)
      =\bigl(\overline X_f,j_f,p_f,e_f\bigr),
    \qquad e_f\colon p_f\circ j_f=f.
  \]
  Evaluation of the three cited lemmas at the components of
  $\mathsf C_{\mathrm{Nag}}(f)$ is the conjunction
  \[
    \operatorname{IsQuasiCompactOpenImmersion}(j_f)
      \ \wedge\ 
    \operatorname{IsProper}(p_f)
      \ \wedge\ 
    (p_f\circ j_f=f).
  \]
  where the three conjuncts are respectively
  \Cref{lemma:nagata-final-open-immersion,lemma:nagata-final-ambient-proper,lemma:nagata-final-factorization}.
  Substitution of these data into
  \Cref{def:nagata-compactification} defines the element
  \[
    (\overline X_f,j_f,p_f,e_f)
      \in\operatorname{Compactification}(f).
  \]
  The projection to nonemptiness and the definition of the compactifiability
  predicate are the formal chain
  \[
    \mathsf C_{\mathrm{Nag}}(f)
      \in\operatorname{Compactification}(f)
      \Longrightarrow
    \operatorname{Compactification}(f)\neq\varnothing
      \Longleftrightarrow
    \operatorname{Compactifiable}(f),
  \]
  by \Cref{def:nagata-compactifiable}.
\end{proof}
